\section{NN-A001: exact arithmetic label and a certified correction}
\label{sec:NN-A001}

\paragraph{Inputs.}
A supplies RH.  The analytic inputs are the Riemann--von Mangoldt identity
$N(t)=\theta(t)/\pi+1+S(t)$ with midpoint convention at a zero, and the
Hurwitz derivative identity DLMF 25.11.17.  The supplied consult
\path{OURS_CONSULT_ad855737_ZETA_LABELS.md}, Q1, Proposition 1, already
displays the correction term in its proof, but omits it from its claimed
equivalence.  We retain that proof's term and independently certify a
counterexample to the simplified assertion.

\begin{theorem}[NN-A001]
Let $\rho_n=1/2+i\gamma_n$ be a simple positive-ordinate zero, indexed with
multiplicity among all positive-ordinate zeros, and set
\[
 Z(t)=e^{i\theta(t)}\zeta(1/2+it),\quad
 \phi_n=\arg\zeta(3/2+i\gamma_n),\quad
 W_n=S(\gamma_n)-\phi_n/\pi.
\]
Choose the argument on the absolutely convergent Euler-product branch.
For the deformation $\zeta(s,1+t)$ put
$v_n=\rho_n\zeta(\rho_n+1)/\zeta'(\rho_n)$ and
$\delta_n=\arctan(1/(2\gamma_n))$.  Then
\[
 \operatorname{Re}v_n=
 \frac{|\zeta(\rho_n+1)|}{|Z'(\gamma_n)|}
 \left(\gamma_n\sin\pi W_n+\frac12\cos\pi W_n\right).       \tag{A001.1}
\]
Thus the exact right-mover criterion is
$\sin(\pi W_n+\delta_n)>0$, equivalently
$W_n+\delta_n/\pi\bmod2\in(0,1)$.  Tangency is exactly equality to an
integer modulo one after this correction.  The uncorrected criterion in
the prompt is false: the indexed zero $n=1169$ is a certified right-mover
although $W_{1169}<0$ and $\sin(\pi W_{1169})<0$.
\end{theorem}

\begin{proof}
Implicit differentiation using DLMF 25.11.17 gives the stated $v_n$.
Differentiating $Z(t)$ at a zero gives
$\zeta'(\rho_n)=-i e^{-i\theta(\gamma_n)}Z'(\gamma_n)$.
Since $Z(0)<0$, the sign of $Z'$ at this simple zero is
$(-1)^{n+1}$; preceding zeros of even multiplicity change neither this
sign nor the parity of the total zero count.  The midpoint convention
gives $\theta(\gamma_n)/\pi=n-3/2-S(\gamma_n)$.  Hence
\[
 \cos(\theta+\phi)=(-1)^n\sin\pi W_n,\qquad
 \sin(\theta+\phi)=(-1)^n\cos\pi W_n.
\]
Taking the real part of
$i\rho_n e^{i\theta}\zeta(\rho_n+1)/Z'(\gamma_n)$ proves (A001.1).
The positive factor $\sqrt{\gamma_n^2+1/4}$ converts its last parenthesis
to $\sqrt{\gamma_n^2+1/4}\sin(\pi W_n+\delta_n)$.

The included program \path{check_label_1169.py} calls the indexed
FLINT/Arb zero isolator and evaluates the derivative through a complex
ball power series, at 60 decimal digits.  It certifies
\[
 \gamma_{1169}=1613.0788554510274235\ldots,\quad
 W_{1169}=-0.0000606299724187599509\ldots,
\]
\[
 \delta_{1169}/\pi=0.0000986653178528222938\ldots,
 \qquad \operatorname{Re}v_{1169}=0.025925651967163814\ldots>0.
\]
The saved ball intervals in \path{label_certificate.json} exclude zero
for both sine signs and for the velocity's real part.  The program asserts
the relevant inequalities using ball comparisons, rather than midpoint
rounding.  This supplies the counterexample claimed in the theorem.
\end{proof}

\paragraph{Interpretation and geometry.}
The formula is a real normal component of a tangent vector to a moving
analytic Cartier divisor.  Its sign is a Boolean-valued statistic where
the normal component is nonzero; it is not a homomorphism of the compact
circle into a two-element group.  Indeed $\operatorname{sgn}\sin\pi x$ is
a discontinuous square wave, and even as a set map is not multiplicative
under addition: two small positive $x,y$ give $+,+,+$, while suitable
$x,y$ give $+,+,-$.  In particular the consult's description as a sign
character of a quotient circle is not a representation-theoretic identity.
The weight-zero Boolean skeleton in the split-zero construction contains
no value of $\zeta'(\rho)$ or of this normal vector.  This corrected label
therefore supplies no SAT, parity, or multiplicity-positivity algorithm.

\section{NN-A002: the exact statistical gate and why a CLT does not close it}
\label{sec:NN-A002}

\begin{theorem}[NN-A002]
A central limit theorem for $S(\gamma_n)$ after division by a standard
deviation tending to infinity, even together with boundedness of
$\phi_n$, does not imply that $W_n\bmod2$ equidistributes, nor that half
the labels in NN-A001 are positive.

A sufficient arithmetic gate for the half-proportion, on the chosen
simple-zero subsequence, is the following pair of assertions.  For every
odd integer $k$, the empirical means of $e^{i\pi kW_n}$ tend to zero;
and
\[
 \lim_{\epsilon\downarrow0}\limsup_{T\to\infty}
 \frac{\#\{n:\gamma_n\leq T,\ n\text{ simple},\
                 \operatorname{dist}(W_n,\mathbb Z)<\epsilon\}}
      {\#\{n:\gamma_n\leq T,\ n\text{ simple}\}}=0.       \tag{A002.1}
\]
All nonzero Fourier modes tending to zero is the stronger usual Weyl
equidistribution gate.  The required odd modes have the exact form
\[
 e^{i\pi kW_n}=e^{-3\pi i k/2}(-1)^n
          e^{-ik\theta(\gamma_n)}e^{-ik\phi_n}\qquad(k\text{ odd}).
                                                               \tag{A002.2}
\]
\end{theorem}

\begin{proof}
Let $G$ be standard normal, let $\sigma_j\to\infty$, and set
$X_j=2\lfloor\sigma_jG/2\rfloor+1/4$.
Then $|X_j-\sigma_jG|\leq7/4$, so $X_j/\sigma_j$ converges almost surely
and in every fixed $L^p$ to $G$.  In particular the normalized central
limit law and all its fixed normalized moment consequences hold.
Nevertheless $X_j\bmod2=1/4$ identically, so every square-wave label is
positive.  Equivalently, an exactly Gaussian variable $\sigma_jG$ can
be altered by a uniformly bounded, correlated correction to have this
completely nonuniform residue class.  This proves the failure of the
implication invoked in the supplied consult, Q2, Proposition 4.

For the sufficient gate, approximate the square wave
$f(x)=\operatorname{sgn}\sin\pi x$ by its Fejer trigonometric means.
It has only odd Fourier modes, because $f(x+1)=-f(x)$, and zero constant
term.  On any closed set a positive distance from $\mathbb Z$, these
means converge uniformly to $f$, and they are bounded in absolute value
by one.  For fixed truncation the empirical mean tends to zero by the
odd-mode hypothesis.  Equation (A002.1) controls the omitted boundary
neighborhood; then let the truncation tend to infinity and that
neighborhood shrink.  Thus the empirical mean of $f(W_n)$ tends to zero.
Since $\delta_n\to0$, (A002.1) also shows that replacing $f(W_n)$ by
$f(W_n+\delta_n/\pi)$ changes an asymptotically negligible proportion.
This proves the half-proportion for the exact labels.  Finally substitute
$W_n=n-3/2-\theta(\gamma_n)/\pi-\phi_n/\pi$ to get (A002.2).
\end{proof}

\paragraph{A primary CLT locator and its frequency range.}
The previously missing source can be made precise: Korolev,
\emph{On the Gram's Law in the Theory of Riemann Zeta Function},
arXiv:1011.3997, Theorem~14, printed pp.33--34, gives a Gaussian
law for $S(\gamma_n\pm0)/\sigma_N$, where
$\sigma_N=\sqrt{L/2}/\pi$ and $L=\log\log N$, on the stated
index blocks. Its proof invokes Theorem~10, pp.21--22. There the
characteristic-function argument is truncated at normalized frequency
$\lambda\asymp\sqrt{\log L}$. In the original $S$ coordinate
this reaches only $\lambda/\sigma_N
=O(\sqrt{\log L/L})\to0$, not the fixed nonzero frequencies
$\pi k$ required by NN-A002. The theorem also does not insert the
correlated exterior phase. Thus the concrete classical CLT, and not
merely an abstract CLT, leaves the identified fixed-frequency gate open.

\paragraph{Located progress beyond the invalid CLT route.}
Arias de Reyna, \emph{On the distribution (mod 1) of the normalized zeros
of the Riemann zeta-function}, arXiv:1403.7126, Theorem 3.1 and Corollary
3.2, prove cancellation of
$N(T)^{-1}\sum_{0<\gamma\leq T}e^{2i\kappa\theta(\gamma)}$ for
$0<\kappa<6/5$ once A identifies the ordinates with the real arguments
in that theorem.  The theorem itself is formulated without RH using
complex ordinates.  Its statement does not supply the alternating index
factor and exterior phase in (A002.2), does not select the simple-zero
subsequence, and does not establish all the required modes.  These are
specific additional gates.  The weighted first moment of velocities is
also insufficient: cancellation of a sum of real velocities controls
their sizes with signs, whereas the question here counts signs without
those weights.

\paragraph{Dictionary and complexity impact.}
The Fourier gate is a statement about characters $e^{i\pi kx}$ of a
compact real torus and a discontinuous observable on it.  The sample
points are analytic divisor points; no finite-type scheme or Frobenius
action on these phase values has been constructed.  This statement
neither establishes a nonstandard Weil purity theorem nor changes the
complexity of support, parity, or satisfiability.


\section{NN-A003: a CUE theorem with the exterior phase retained}
\label{sec:NN-A003}

\paragraph{Inputs and status.}
This result uses no instance of A, M, E, or a GCT hypothesis.  Its inputs are
the Haar eigenangle density for $U(N)$ (Weyl integration), the elementary
determinantal integral identity proved below, and elementary Fourier analysis.
The Haar character orthogonality used in the final phase estimate also appears
in the supplied paper 12, Theorem \texttt{thm:first-deriv-mean}.  This is a
proved form within the specified CUE model, not a transfer theorem for zeta
zeros.  In particular the arithmetic phase is not discarded on the grounds
that it is bounded.

\begin{theorem}[NN-A003: marked CUE labels and their correlations]
Let $U_N$ be Haar distributed in $U(N)$, let
$0\leq\theta_1<\cdots<\theta_N<2\pi$ be its eigenangles, and choose an index
$I$ uniformly from $\{1,\ldots,N\}$, independently after sampling $U_N$.
Fix $0\leq r<1$ and $a\in\mathbb R$.  Define the continuous exterior phase
\[
 \phi_{N,r}(x)=\operatorname{Im}\sum_{j=1}^N
       \log(1-r e^{i(x-\theta_j)}),
 \qquad
 W_{N,i}=i-\frac{N\theta_i}{2\pi}-a-
                 \frac{\phi_{N,r}(\theta_i)}{\pi}.
\]
The logarithm is its power-series branch in $r$, so there is no branch
ambiguity.  Set $\eta_{N,i}=\operatorname{sgn}\sin(\pi W_{N,i})$; a convention
at zero has no effect on the limits.  For fixed $j\geq1$, let
\[
 D_{j,N}=\frac{N}{2\pi}(\theta_{I+j}-\theta_I),
 \qquad \theta_{i+N}=\theta_i+2\pi.
\]
Then $(W_{N,I}\bmod2,D_{j,N})$ converges in distribution to $(V,D_j)$,
where $V$ is uniform on $\mathbb R/2\mathbb Z$, independent of $D_j$, and
$D_j$ is the distance to the $j$th point to the right of zero in the reduced
Palm sine process.  The latter process has kernel
\[
 K^{0}(x,y)=\frac{\sin\pi(x-y)}{\pi(x-y)}-
      \frac{\sin\pi x}{\pi x}\frac{\sin\pi y}{\pi y}.
\]
Consequently
\[
 \lim_{N\to\infty}\mathbb P(\eta_{N,I}=1)=\frac12,
 \qquad
 \lim_{N\to\infty}\mathbb E[\eta_{N,I}\eta_{N,I+j}]
       =\mathbb E[1-2d_2(j-D_j)],                 \tag{A003.1}
\]
where $d_2(x)=\min_{m\in\mathbb Z}|x-2m|\in[0,1]$.
In particular the adjacent agreement probability tends to
$1-\mathbb E d_2(1-D_1)$.  All assertions remain valid for $r=r_N$ in a fixed
compact subinterval of $[0,1)$, provided a subsequence is used if a limiting
phase law itself is requested; the displayed label and gap limits are
independent of that phase law.
\end{theorem}

\begin{proof}
Normalize circle measure to $d\theta/(2\pi)$ and let $P_N$ be the orthogonal
projection onto the span of $1,e^{i\theta},\ldots,e^{i(N-1)\theta}$.
The squared Vandermonde density is the determinant of the kernel of $P_N$.
Expanding its two determinants and integrating each coordinate shows directly
that, for a bounded multiplier $u$,
\[
 \mathbb E\prod_i u(\theta_i)
       =\det(P_NuP_N|_{\operatorname{ran}P_N}).       \tag{A003.2}
\]
Indeed each term in the determinant expansion is a product of the matrix
entries $\int u(\theta)e^{i(k-l)\theta}d\theta/(2\pi)$, and the permutation
signs are precisely the determinant signs.  Conditioning one uniformly
marked point to be $x$ gives the same formula with the rank $N-1$ projection
\[
 Q_{N,x}=P_N-|v_x\rangle\langle v_x|,
 \qquad v_x(\theta)=P_N(\theta,x)/\sqrt N.
\]
This conditional kernel follows either by the Schur complement of the
determinant with the marked row and column fixed, or by expanding that row
and column in (A003.2).  The marked point itself is uniform on the circle.

For any finite-rank projection $Q$ and multiplier $u$ with $|u|=1$, put
$A=QuQ|_{\operatorname{ran}Q}$.  Then
\[
 |\det A|^2=\det\bigl(I_Q-Qu^*(1-Q)uQ\bigr)
       \leq \exp\bigl(-\|(1-Q)uQ\|_{\rm HS}^2\bigr),       \tag{A003.3}
\]
because the eigenvalues of the positive contraction subtracted from $I_Q$
are in $[0,1]$, and $1-y\leq e^{-y}$.  For $Q=P_N$ its exponent equals
\[
 F_N(u):=\|(1-P_N)uP_N\|_{\rm HS}^2
       =\sum_{n\in\mathbb Z}\min(N,|n|)|\widehat u(n)|^2.   \tag{A003.4}
\]
This follows by applying $u$ to each Fourier basis vector and counting how
many translates by $n$ leave $\{0,\ldots,N-1\}$.  Removing one vector from
the projection changes this leakage by at most one.  To see this write
$P=Q+R$, where $R$ has rank one.  Orthogonality gives
\[
 F(P)-F(Q)=-\|RuQ\|_{\rm HS}^2+
                 \|(1-P)uR\|_{\rm HS}^2\in[-1,1].       \tag{A003.5}
\]

Consider a nonzero Fourier mode $k\in\mathbb Z$ of $W\bmod2$.
When $k$ is even, $e^{i\pi kI}=1$.  The exterior phase at the marked
eigenangle and all cyclic gaps are invariant under simultaneous rotation of
all eigenangles.  Averaging the factor $e^{-ikNx/2}$ over that rotation
therefore makes the Fourier mode exactly zero, also after multiplication
by any bounded function of cyclic gaps.

When $k$ is odd, condition on the marked point $x\in(0,2\pi)$.  Its index
is $1+$ the number of remaining points in $(0,x)$.  Up to a scalar of
modulus one, the conditional Fourier mode is (A003.2) for $Q_{N,x}$ and
\[
 u_x(\theta)=(-1)^{\mathbf1_{(0,x)}(\theta)}
        \exp\{-ik\,\operatorname{Im}\log(1-r e^{i(x-\theta)})\}.
\]
This function is smooth except for two jumps, at $0$ and $x$, each of
absolute size two.  Piecewise integration by parts twice gives
\[
 \widehat u_x(n)=\frac{J_0+J_x e^{-inx}}{2\pi i n}+O_{x,k,r}(n^{-2}),
 \qquad |J_0|=|J_x|=2.
\]
For fixed $x\notin2\pi\mathbb Z$ the cross sum
$\sum_{n\leq N}e^{inx}/n$ is bounded.  Thus (A003.4), including positive
and negative frequencies, gives
\[
 F_N(u_x)\geq\frac4{\pi^2}\log N+O_{x,k,r}(1).
\]
Equations (A003.3)--(A003.5) imply that the conditional odd mode tends to
zero.  Its absolute value is at most one, so dominated convergence over
the marked angle proves marginal equidistribution.  Derivative bounds of
the smooth factor are uniform when $r$ stays bounded away from one.

The same estimate proves independence from local gaps, which does not
follow from marginal equidistribution alone.  Fix $s>0$ and let $J$ be the
arc of length $2\pi s/N$ immediately after $x$.  For $|z|\leq1$ replace
$u_x$ by $w=u_x z^{\mathbf1_J}$.  This is a contraction multiplier, and
\[
 \|(w-u_x)Q_{N,x}\|_{\rm HS}^2
      \leq 4\operatorname{Tr}(Q_{N,x}\mathbf1_JQ_{N,x})\leq4s.
\]
For a contraction $w$, $I_Q-(QwQ)^*(QwQ)$ is the sum of
$Q(1-|w|^2)Q$ and $Qw^*(1-Q)wQ$, both positive.  Hence its determinant
bound has exponent at least $\|(1-Q)wQ\|_{\rm HS}^2$.  By the triangle
inequality this is at least
\[
 \bigl(\|(1-Q)u_xQ\|_{\rm HS}-2\sqrt s\bigr)_+^2,
\]
which still tends to infinity.  It follows uniformly on every circle
$|z|=b<1$ that the mixed conditional generating function of the odd
Fourier mode and $z^{\#J}$ tends to zero.  Cauchy coefficient extraction
then gives the same assertion with $\mathbf1_{\#J=l}$ for every fixed
$l$.  Since $D_{j,N}>s$ means $\#J\leq j-1$, this proves vanishing of
all mixed Fourier modes against these gap events.  The even modes already
vanish by rotation.  The gaps are tight, since their mean is exactly $j$
by cyclic summation of all gaps.

After scaling by $N/(2\pi)$ and removing a harmless phase, the conditional
kernel tends uniformly on each compact interval to $K^0$ displayed in the
theorem.  Determinantal generating functions on such an interval converge:
their Fredholm series converge term by term, and Hadamard's bound for
positive kernels with diagonal at most one bounds the $m$th term by
$ (2s)^m/m!$. The limiting $D_j$ has no atoms: a determinantal process
with continuous intensity has probability zero of a point at any
prescribed location, and $D_j=s$ requires such a point at $s$.
The local count formulas and tightness therefore give convergence of
the scaled $j$th gaps to $D_j$ at every survival threshold. The mixed Fourier
identities just proved characterize the joint limit as the product of
Haar measure modulo two and this gap law.

It remains to compare the phase at neighboring eigenangles.  The convergent
Fourier expansion gives
\[
 \sup_x|\phi_{N,r}'(x)|\leq
          \sum_{m\geq1}r^m|\operatorname{Tr}(U_N^m)|.
\]
Direct Fourier integration of the CUE two-point determinant gives
$\mathbb E|\operatorname{Tr}(U_N^m)|^2=\min(m,N)$; the same identity
follows by substituting $e^{im\theta}$ into (A003.4).  Cauchy--Schwarz
therefore bounds the expectation of this supremum by
$\sum_{m\geq1}r^m\sqrt m<\infty$, uniformly in $N$.
Tightness of $D_{j,N}$ shows that
$\phi_{N,r}(\theta_{I+j})-\phi_{N,r}(\theta_I)$ tends to zero in
probability.  With unwrapped indices and angles the exact difference is
\[
 W_{N,I+j}-W_{N,I}=j-D_{j,N}
       -\frac{\phi_{N,r}(\theta_{I+j})-\phi_{N,r}(\theta_I)}{\pi}.
\]
Finally, for uniform $V\bmod2$, the two square-wave signs at $V$ and
$V+v$ agree on a set of normalized measure $1-d_2(v)$.  Their correlation
is $1-2d_2(v)$.  The discontinuity boundaries have zero limiting Haar
measure, so the joint convergence proves (A003.1).
\end{proof}

\paragraph{What this changes.}
The CUE half-proportion and correlation formulas are justified within an
explicit model containing the correlated exterior phase.  The earlier
nonperiodic expression $1-\mathbb E|1-D_1|$ for adjacent agreement must be
replaced by $1-\mathbb E d_2(1-D_1)$; gaps larger than two cannot be silently
omitted.  To apply this to the arithmetic labels still requires a joint
arithmetic limit theorem for the label phase and local zero gaps.
Montgomery's restricted pair-correlation theorem supplies neither the odd
label modes nor the nearest-neighbor law.  This finding supplies no
algorithm for a Boolean counting problem and no separation of complexity
classes.

\paragraph{Geometric and representation-theoretic dictionary.}
At finite $N$, (A003.2) is an exterior-algebra determinant of a compressed
operator, derived from Weyl integration on the compact real form of
$\mathrm{GL}_N$.  The rank-one Palm operation is removal of the evaluation
line, and the leakage in (A003.3) is a Hilbert--Schmidt norm.  These are
precise Hilbert-space and representation-theoretic constructions; they do
not assert a scheme over $\mathbb F_q$, a finite Weil object, or Frobenius
purity.  The limiting sine projection has infinite rank.  No Grothendieck
archive cote is claimed for this new analytic argument.

\paragraph{Exact scope of the classical arithmetic pair theorem.}
The primary locator is Montgomery, \emph{The pair correlation of zeros
of the zeta function} (1973), Section~1, Theorem and equation~(2),
printed p.181; the Fourier identity (3)--(4) and explicit support
restriction are on p.182. Under RH, for each fixed $\epsilon>0$,
\[
 F(\alpha,T)=(1+o(1))T^{-2|\alpha|}\log T+|\alpha|+o(1)
 \quad (|\alpha|\le1-\epsilon),
\]
with the stated uniformity. The full continuation beyond this range
is presented as conjecture (11)--(12), pp.183--184. This controls
band-limited pair observables, while a neighboring-gap event requires
an empty interval and thus higher count information. The arithmetic
derivative and exterior phase are additional marks. Neither of these
missing pieces is supplied by that restricted pair theorem. This is
existing background, not a repeated consequence from the prior catalogue.


\section{Quantitative empirical balance of the CUE labels}
\label{sec:analytic-cue-empirical-balance}

NN-A003 concerns a uniformly marked eigenangle. The following
statement controls the actual fraction of labels in a single
sample, while retaining the same correlated exterior phase.

\begin{theorem}[NN-A031: empirical half proportions in mean square]
\label{thm:analytic-cue-empirical-balance}
Fix $r_0<1$. For each $N$, use the CUE labels of NN-A003 with
arbitrary $0\le r_N\le r_0$ and arbitrary offset $a_N\in\mathbb R$:
\[
 W_{N,i}=i-\frac{N\theta_i}{2\pi}-a_N
          -\frac1\pi\operatorname{Im}\sum_j
                   \log(1-r_Ne^{i(\theta_i-\theta_j)}),
 \qquad \eta_{N,i}=\operatorname{sgn}\sin(\pi W_{N,i}).
\]
Almost surely none of these labels is zero. There is a constant
$C_{r_0}$, independent of $N,r_N,a_N$, such that for $N\ge3$,
\begin{equation}\label{eq:analytic-cue-empirical-rate}
 \mathbb E\left[\left(\frac1N\sum_{i=1}^N\eta_{N,i}\right)^2\right]
                          \le\frac{C_{r_0}}{\log N}.
\end{equation}
In particular, the observed positive-label fraction $L_N$ satisfies
\[
 \mathbb E[(L_N-1/2)^2]\le\frac{C_{r_0}}{4\log N},\qquad
 \mathbb P(|L_N-1/2|>\varepsilon)
                   \le\frac{C_{r_0}}{4\varepsilon^2\log N}.
\]
Thus $L_N\to1/2$ in $L^2$ and in probability. No optimality of
the logarithmic error bound, almost-sure convergence along all
$N$, or arithmetic transfer is asserted.
\end{theorem}

\begin{proof}
Write $P=P_N$ for the CUE projection in
$L^2(S^1,d\theta/(2\pi))$, and let $I,J$ be two distinct ordered
uniform indices sampled after the matrix. All constants below
are independent of the offset and of $r\in[0,r_0]$.

\emph{An exact uniform marginal modulo one.}
Under simultaneous rotation by $\alpha$, the exterior phase at
a carried marked point is unchanged. Its coordinate changes by
$\alpha$ modulo $2\pi$, and its new index differs from its old
index by an integer. Consequently its $W$ coordinate modulo one
changes by $-N\alpha/(2\pi)$. The marked Haar ensemble is
rotation invariant, so for every nonzero integer $h$,
\[
 \mathbb E e^{2\pi ihW_{N,I}}
 =e^{-ihN\alpha}\mathbb E e^{2\pi ihW_{N,I}}
 \quad\hbox{for every }\alpha.
\]
These expectations vanish. Hence $W_{N,I}\bmod1$ is exactly
uniform for every $N$. In particular a uniformly marked label
is almost surely nonzero, and a finite union gives the same
assertion for all labels in a sample. For any integrable
one-periodic function $g$ we have the exact identity
\begin{equation}\label{eq:analytic-cue-exact-modone}
             \mathbb E g(W_{N,I})=\int_0^1g(u)\,du.
\end{equation}

\emph{Two-Palm odd Fourier modes.}
Condition on the marked eigenangles $x\ne y$. The remaining
points form a projection determinantal process with projection
$Q_{N,x,y}$ obtained by removing the two-dimensional span of
the evaluation vectors at $x,y$ from $\operatorname{ran}P$.
This follows by the Schur complement of their two rows and
columns in the CUE density; the two evaluation vectors are
linearly independent for $N\ge2$ and $x\ne y$.

Let $k,l$ be odd integers. Their index parities combine to the
parity of the number of remaining eigenangles in the arc
between $x$ and $y$ that does not cross the coordinate origin.
For example, when $x<y$, the two counts in $(0,x)$ and $(0,y)$
have sum congruent modulo two to the count in $(x,y)$; the
contribution from the two marked points is a constant sign.
Define
\[
 j_{x,y}(\theta)=(-1)^{\mathbf1_{(\min(x,y),\max(x,y))}(\theta)},
 \qquad g_x(\theta)=\operatorname{Im}\log(1-re^{i(x-\theta)}),
\]
and
\[
 v(\theta)=e^{-ikg_x(\theta)-ilg_y(\theta)},\qquad u=j_{x,y}v.
\]
The self- and cross-contributions of the marked points to the
exterior phases are fixed scalars of modulus one. The determinant
identity in NN-A003 therefore gives
\begin{equation}\label{eq:analytic-cue-two-palm-mode}
 \left|\mathbb E\left[e^{i\pi(kW_{N,I}+lW_{N,J})}
                     \mid x,y\right]\right|
   =\left|\det(Q_{N,x,y}uQ_{N,x,y})\right|.
\end{equation}
The determinant is on the range of $Q_{N,x,y}$.

For a unimodular multiplier $w$, retain the notation
\[
 F_N(w)=\|(1-P)wP\|_{\rm HS}^2
       =\sum_{n\in\mathbb Z}\min(N,|n|)|\widehat w(n)|^2.
\]
Expanding the two off-diagonal blocks of the commutator gives
$\|[P,w]\|_{\rm HS}^2=2F_N(w)$: both blocks have squared norm
$N-\|PwP\|_{\rm HS}^2$. Since
$[P,jv]=[P,j]v+j[P,v]$, the triangle inequality gives
\begin{equation}\label{eq:analytic-cue-smooth-multiplier}
       \sqrt{F_N(jv)}\ge\sqrt{F_N(j)}-\sqrt{F_N(v)}.
\end{equation}
The smooth multiplier has derivative bounded by
\[
 \|v'\|_2\le C_0(|k|+|l|),\qquad C_0=\frac{r_0}{1-r_0}.
\]
By Cauchy--Schwarz and Parseval,
\begin{equation}\label{eq:analytic-cue-smooth-leakage}
 F_N(v)\le\sum_n|n||\widehat v(n)|^2
     \le\|v'\|_2\|v\|_2
     \le C_0(|k|+|l|).
\end{equation}

Let $d(x,y)\in(0,\pi]$ be their circular separation and let
$\ell=|x-y|$. The arc multiplier has the exact coefficients
\[
 |\widehat j_{x,y}(n)|^2
       =\frac{2-2\cos(n\ell)}{\pi^2n^2}\quad(n\ne0).
\]
Consequently
\[
 F_N(j_{x,y})\ge\frac4{\pi^2}
       \left(H_N-\sum_{n=1}^N\frac{\cos(n\ell)}n\right).
\]
Summation by parts bounds the tail of
$\sum_{n\ge1}e^{in\ell}/n$ by
$2/(N|1-e^{i\ell}|)$. Its real infinite sum is
$-\log|1-e^{i\ell}|$. Since
$|1-e^{i\ell}|\ge2d(x,y)/\pi$, for
$d(x,y)\ge N^{-1/2}$ we obtain, uniformly,
\[
 F_N(j_{x,y})\ge\frac2{\pi^2}\log N-O(1)
                         \ge\frac1{\pi^2}\log N
\]
for all sufficiently large $N$.

Put
\[
 c_* =\frac1{64\pi^2(1+C_0)},\qquad
 M=\lfloor c_*\log N\rfloor,
\]
and assume $N$ is sufficiently large that $M\ge2$. If $k,l$ are
odd and $|k|,|l|\le2M$, then
\[
 F_N(v)\le4C_0M\le\frac{\log N}{16\pi^2}
                         \le\frac1{16}F_N(j_{x,y}).
\]
Equation \eqref{eq:analytic-cue-smooth-multiplier} yields
$F_N(u)\ge(9/16)F_N(j_{x,y})\ge(\log N)/(2\pi^2)$.
Removing a rank-one projection changes the leakage by at most
one, as proved in NN-A003; hence the two-Palm leakage is at
least $F_N(u)-2$. The determinant contraction from NN-A003 gives
\[
 |\det(Q_{N,x,y}uQ_{N,x,y})|
                 \le eN^{-\kappa},\qquad
                      \kappa=\frac1{4\pi^2},
\]
whenever $d(x,y)\ge N^{-1/2}$.

The joint density of the ordered marked pair, relative to
$dx\,dy/(2\pi)^2$, is
\[
 \frac{N^2-|P_N(x,y)|^2}{N(N-1)}\le\frac{N}{N-1}\le2.
\]
The excluded region $d(x,y)<N^{-1/2}$ has measure
$O(N^{-1/2})$, and a conditional characteristic function has
absolute value at most one there. Integration therefore proves
the uniform estimate
\begin{equation}\label{eq:analytic-cue-uniform-odd-pair}
 \left|\mathbb E e^{i\pi(kW_{N,I}+lW_{N,J})}\right|
                  \le C N^{-\kappa}
 \quad\left(\begin{array}{c}k,l\text{ odd},\\|k|,|l|\le2M\end{array}\right).
\end{equation}
No independence of the exterior phases has been used.

\emph{A positive finite Fourier approximation to the labels.}
On $\mathbb R/2\mathbb Z$, with normalized measure $du/2$, let
$f(u)=\operatorname{sgn}\sin(\pi u)$ and normalize the positive
Jackson kernel
\[
 J_M(u)=\frac1{Z_M}
       \left|\frac{\sin(M\pi u/2)}{\sin(\pi u/2)}\right|^4,
 \qquad \int_{-1}^1J_M(u)\frac{du}{2}=1.
\]
The quotient has its continuous value $M$ at zero. Elementary
sine bounds give $Z_M\asymp M^3$ and
\[
 J_M(u)\le CM\min\{1,(M|u|)^{-4}\},\qquad
 \int_{-1}^1|u|J_M(u)\frac{du}{2}\le\frac CM.
\]
For the lower bound on $Z_M$, restrict to $|u|\le1/M$ and use
$|\sin(M\pi u/2)|\ge M|u|$ and
$|\sin(\pi u/2)|\le\pi|u|/2$. The upper and moment bounds follow
by splitting at $1/M$.

Set $q_M=f*J_M$. This is a real trigonometric polynomial of
degree at most $2(M-1)$, with $|q_M|\le1$. Because
$f(u+1)=-f(u)$, it has only odd Fourier modes. Moreover,
\begin{equation}\label{eq:analytic-cue-jackson-error}
 \|f-q_M\|_{L^1(\mathbb R/2\mathbb Z)}\le\frac CM,
 \qquad \sum_k|\widehat q_M(k)|\le C\log(M+1).
\end{equation}
The first estimate follows by averaging
$\|f(\cdot-u)-f\|_1\le2|u|$ against $J_M$. For the second use
$|\widehat J_M(k)|\le1$ and
$|\widehat f(k)|=2/(\pi|k|)$ for odd $k$.

The function $|f-q_M|$ is one-periodic. Thus
\eqref{eq:analytic-cue-exact-modone} makes its expected value at
the marked $W$ exactly its circle integral, without an
asymptotic equidistribution error. Since all labels and $q_M$
values have absolute value at most one,
\[
 \left|\mathbb E[\eta_{N,I}\eta_{N,J}]
       -\mathbb E[q_M(W_{N,I})q_M(W_{N,J})]\right|
                                             \le\frac CM.
\]
Using the odd-mode estimate
\eqref{eq:analytic-cue-uniform-odd-pair} in the finite Fourier
expansion gives
\[
 |\mathbb E[\eta_{N,I}\eta_{N,J}]|
       \le\frac CM+C N^{-\kappa}\log^2(M+1).
\]
Finally the diagonal terms of the empirical square contribute
$1/N$, so
\[
 \mathbb E\left[\left(\frac1N\sum_i\eta_{N,i}\right)^2\right]
 \le\frac1N+\frac CM+C N^{-\kappa}\log^2(M+1)
                              =O_{r_0}(1/\log N).
\]
Increasing the constant covers the finitely many smaller $N$.
The assertions about $L_N=(1+N^{-1}\sum_i\eta_{N,i})/2$ follow
from this identity and Markov's inequality.
\end{proof}

\paragraph{Geometric translation.}
Identify $\operatorname{ran}P_N$ with the restrictions to the unit
circle of $H^0(\mathbb P^1,\mathcal O(N-1))$, carrying its Haar
Hermitian norm. At two distinct marked points $z_x,z_y$, polynomial
division gives the exact evaluation sequence
\[
 0\longrightarrow (z-z_x)(z-z_y)\mathbb C[z]_{\le N-3}
 \longrightarrow\mathbb C[z]_{\le N-1}
 \longrightarrow\mathbb C^2\longrightarrow0.
\]
The two-Palm projection is the orthogonal projection onto this
vanishing subspace. Its exterior determinant controls the marked
Fourier statistics; the rank-two correction and the commutator
norm make that control quantitative. The representation-ring
exterior-power formalism is SGA~6, Expos\'e~V, Example~2.2(b) and
Lemma~2.2.1. These actual evaluation spaces and Haar projections
give a probabilistic balance law. They do not provide an
arithmetic comparison map or a Weil purity argument.


\section{A local gap limit and distant-label decorrelation in the CUE limit}

\begin{theorem}[NN-A016: periodic gap displacement and pair decorrelation]
Let $D_j$ be the $j$th positive point of the reduced Palm sine process
at zero, with kernel $K^0$ as in NN-A003. It has a density $f_j$.
For $j\ge2$ put $\sigma_j=\sqrt{\log j}/\pi$ and extend $f_j(s)$
by zero for $s\le0$. Then
\[
 \int_{\mathbb R}\left|
 \sigma_j f_j(j+\sigma_j y)
       -\frac{e^{-y^2/2}}{\sqrt{2\pi}}\right|dy\longrightarrow0.
                                                        \tag{A016.1}
\]
In particular, for every fixed nonzero real $a$,
\[
 \mathbb E e^{ia(D_j-j)}\longrightarrow0,               \tag{A016.2}
\]
and $(D_j-j)\bmod2$ tends to Haar measure on $\mathbb R/2\mathbb Z$.
Consequently the limiting CUE label correlations in NN-A003 satisfy
\[
 \mathbb E[1-2d_2(j-D_j)]\longrightarrow0.
\]
Thus two labels far apart in the limiting indexed model become
independent fair signs. The order of limits is first $N\to\infty$
for fixed $j$ in NN-A003, then $j\to\infty$ here. No uniform
simultaneous-$j,N$ claim or arithmetic transfer is asserted.
\end{theorem}

\begin{proof}
We include the local-limit step, because an ordinary central limit
theorem alone would not imply (A016.2).

Write $k(x)=\sin(\pi x)/(\pi x)$. The sine kernel $K(x,y)=k(x-y)$
is the orthogonal projection in $L^2(\mathbb R)$ onto Fourier
frequencies $[-1/2,1/2]$. Its compression $A_s$ to $(0,s)$ is a
positive trace-class contraction with trace $s$. The determinant
generating function gives
\[
 \mathbb E z^{N(0,s)}=\det(I+(z-1)A_s)
       =\prod_\ell(1-\lambda_\ell+\lambda_\ell z).
\]
Thus this count has the distribution of a sum of independent Bernoulli
variables with parameters the eigenvalues $\lambda_\ell\in[0,1]$.
The product converges because $\sum\lambda_\ell=s<\infty$;
its probability generating function identifies the distribution.
Direct integration gives
\[
 \operatorname{Var}N(0,s)
 =s-2\int_0^s(s-u)k(u)^2du
 =\frac{\log s}{\pi^2}+O(1).                          \tag{A016.3}
\]
Indeed $\int_0^\infty k(u)^2du=1/2$; the remaining tail
$2s\int_s^\infty k(u)^2du$ is bounded, while
$2\int_0^s uk(u)^2du=(\log s)/\pi^2+O(1)$ by
$\sin^2(\pi u)=(1-\cos(2\pi u))/2$.

Conditioning on a fixed point removes its kernel evaluation line
from the projection. Conditioning on two distinct points $0,s$
removes their two-dimensional span. This follows either directly
from the Schur complement of the determinant correlation functions
or from the Gram determinant formula for those evaluation vectors.
Let $B_s$ be the restriction to $(0,s)$ of this two-point reduced
Palm projection. Then $A_s-B_s\ge0$ has trace at most two and rank
at most two. Hence
\[
 \operatorname{Tr}B_s=s+O(1),\qquad
 \operatorname{Tr}(B_s-B_s^2)
       =\frac{\log s}{\pi^2}+O(1).                    \tag{A016.4}
\]
For clarity, if $R=A_s-B_s$, cyclicity of trace gives the variance
difference $\operatorname{Tr}R-\operatorname{Tr}(R(A_s+B_s))$.
Since $0\le A_s+B_s\le2I$, its absolute value is at most
$\operatorname{Tr}R\le2$. The same assertions, with rank one,
hold for the Palm process at zero alone. All these bounds are
uniform as $s\to\infty$.

Here is the needed local central limit lemma. If $X$ is any finite
or countable sum of independent Bernoulli variables, has finite
mean $\mu$, and variance $v\to\infty$ in a sequence of such sums,
then
\[
 \sup_{k\in\mathbb Z}\left|
 \sqrt v\,\mathbb P(X=k)
 -\frac1{\sqrt{2\pi}}e^{-(k-\mu)^2/(2v)}
 \right|\longrightarrow0.                            \tag{A016.5}
\]
This statement is uniform over the Bernoulli parameters once $v$
is sufficiently large. To prove it, the centered characteristic
function $\psi$ satisfies, for $|\theta|\le\pi$,
\[
 |\psi(\theta)|\le
 \exp[-2v\sin^2(\theta/2)]
 \le\exp[-2v\theta^2/\pi^2].
\]
For fixed $u$, Taylor expansion of the Bernoulli logarithms gives
$\log\psi(u/\sqrt v)=-u^2/2+O(|u|^3/\sqrt v)$.
The remainder is uniform because the third derivative of each
centered Bernoulli logarithm, for small argument, is bounded by a
constant times $p(1-p)$. Fourier inversion, scaling by $\sqrt v$,
and dominated convergence on the interval
$[-\pi\sqrt v,\pi\sqrt v]$ prove (A016.5): the uniform error
in $k$ is at most the $L^1$ difference between the scaled
characteristic function, extended by zero, and $e^{-u^2/2}$.
The same argument applies to the countable product by finite
truncation. This proves the lemma without a local-limit assumption.

The Palm Campbell formula gives a density for the $j$th positive
point:
\[
 f_j(s)=K^0(s,s)\,
 \mathbb P^{0,s}\{N(0,s)=j-1\},\qquad s>0,            \tag{A016.6}
\]
where $\mathbb P^{0,s}$ is the two-point reduced Palm law.
Indeed summing a test function over points $s$ with exactly $j-1$
points between $0$ and $s$ selects precisely the $j$th point;
conditioning that summand at $s$ gives the right side.
The count in (A016.6) has the Bernoulli eigenvalue description for
$B_s$. Also $K^0(s,s)=1-k(s)^2\to1$.
For $s=j+y\sigma_j$ with $y$ in any fixed compact set, (A016.4)
and (A016.5) therefore give, uniformly in $y$,
\[
 \sigma_j f_j(j+y\sigma_j)
       \longrightarrow\frac{e^{-y^2/2}}{\sqrt{2\pi}}. \tag{A016.7}
\]
Only an $O(1)$ error in the Palm mean is needed, since
$\sigma_j\to\infty$.

To control tails, apply Chebyshev to the one-point Palm count,
whose mean is $s+O(1)$ and variance is (A016.3) up to $O(1)$.
For each fixed large $M$, the identities
$D_j>s\iff N(0,s)\le j-1$ and their lower-tail counterparts give
\[
 \limsup_{j\to\infty}
 \mathbb P(|D_j-j|>M\sigma_j)\le C/M^2.
\]
There are infinitely many positive points almost surely: the same
variance estimate and Chebyshev at dyadic $s$ give this by
Borel--Cantelli. Thus (A016.6) has total integral one.
Compact-uniform density convergence (A016.7), the displayed tail
bound, and Gaussian tail decay prove (A016.1).

Finally write $g_j(y)=\sigma_j f_j(j+\sigma_j y)$ and
$\varphi(y)=e^{-y^2/2}/\sqrt{2\pi}$. For $a\ne0$,
\[
 \left|\mathbb E e^{ia(D_j-j)}-e^{-a^2\sigma_j^2/2}\right|
 \le\|g_j-\varphi\|_{L^1}\longrightarrow0.
\]
This proves (A016.2), and taking $a=\pi k$ proves all nonconstant
Fourier coefficients of the residue-class law vanish. The continuous
periodic function $1-2d_2(x)$ has Haar mean zero, so its expectation
tends to zero. In NN-A003 both signs have zero mean and the joint
law is unchanged by flipping both, because the uniform phase can
be shifted by one. Their four joint probabilities are therefore
$(1\pm C_j)/4$, where $C_j$ is that correlation, and tend to $1/4$.
\end{proof}

The geometry in this result is again the sine projection, finite-rank
Palm modifications, and positive trace-class contractions on intervals.
The full Fredholm generating function supplies local count information
that a pair-correlation statement does not contain. The proof illustrates
the specific strengthening needed to pass from a Gaussian scaling law
to a fixed-frequency periodic statistic. It asserts neither the
arithmetic counterpart nor mixing of every finite label block.


\section{An exact series and certified value for adjacent CUE agreement}

\begin{theorem}[NN-A017: adjacent agreement from integer gap determinants]
Let $K_s$ be the sine-kernel operator on an interval of length $s$,
and put $E(s)=\det(I-K_s)$, the empty-interval probability for the
unit-intensity stationary sine process. The limiting adjacent-label
agreement probability in NN-A003 is
\[
 A_{\rm adj}=1-2\sum_{n=1}^\infty(-1)^{n+1}E(n).
                                                        \tag{A017.1}
\]
The series is absolutely convergent and its alternating partial sums
give certified two-sided bounds. In the specified model, for every
fixed exterior radius $0\le r<1$ and offset,
\[
 0.66654660183721072003
 < A_{\rm adj} <
 0.66654660183721072008.                                \tag{A017.2}
\]
The interval is certified using degree-$100$ Taylor features and
Arb ball arithmetic at $160$ decimal digits; no quadrature or
uncertified extrapolation is used. It is a CUE-model value, not a
proved proportion for arithmetic zero labels.
\end{theorem}

\begin{proof}
Let $D_1$ be the first positive point under the reduced Palm sine law.
The endpoint derivative of a void probability gives
\[
 -E'(s)=\mathbb P^0(D_1>s),\qquad E''(s)=f_1(s),
                                                        \tag{A017.3}
\]
where the second equality may first be read distributionally.
Indeed, enlarging $(0,s)$ by $ds$ loses precisely configurations
with a point in that added strip and no point in the old interval;
the Campbell formula gives the reduced Palm void probability at
the right endpoint. Translation and reflection symmetry identify
it with the probability on the right side of (A017.3).
The chance of two added points is $O(ds^2)$, since the two-point
intensity is bounded by one. Formula (A016.6) supplies its density;
the two-point Palm kernel varies continuously in trace norm after
rescaling each interval to a fixed unit interval, on every compact
set of $s>0$, so that formula also gives continuity. The distributional
form alone is enough below. In particular $E(0)=1$ and $E'(0)=-1$.

Write $g(s)=d_2(1-s)$ for $s\ge0$. It is the triangular function
with $g(0)=1$, right derivative $g'(0)=-1$, and distributional
second derivative
\[
 g''=2\sum_{n\ge1}(-1)^{n+1}\delta_n.
\]
All eigenvalues of $K_s$ lie in $[0,1]$ and their sum is $s$;
therefore $0\le E(s)\le e^{-s}$. Also $E'(s)\to0$ because
$D_1$ is finite almost surely. Twice integrating by parts, first on
a finite interval and then passing to infinity, gives
\[
 \mathbb E d_2(1-D_1)
 =\int_0^\infty g(s)E''(s)ds
 =2\sum_{n\ge1}(-1)^{n+1}E(n).
\]
The two terms at zero cancel: $-g(0)E'(0)=1$ and
$g'(0)E(0)=-1$. The boundary terms at infinity vanish; the series
converges absolutely by $E(n)\le e^{-n}$. NN-A003 now gives
(A017.1). Since void probabilities decrease with interval length,
the alternating remainder after $m$ terms has absolute value at
most $E(m+1)$ and the usual alternating sign.

Here are full error bounds for the supplied numerical certificate.
Center the spatial interval and write $x=sa$, with
$a,\xi\in[-1/2,1/2]$. The Fourier-feature map has kernel
\[
 T(a,\xi)=\sqrt s\,e^{2\pi is a\xi},\qquad K_s\simeq TT^*.
\]
Replace the exponential by its Taylor polynomial of degree $p$,
defining $T_p$. With $h=\pi s/2$ and
$R=e^h h^{p+1}/(p+1)!$, the uniform remainder gives
\[
 \|T-T_p\|_{\rm HS}\le\sqrt s R,\quad
 \|T\|_{\rm HS}=\sqrt s,\quad
 \|T_p\|_{\rm HS}\le\sqrt s(1+R).
\]
Thus, for $A=TT^*$ and $B=T_pT_p^*$,
\[
 \|A-B\|_1\le sR(2+R)=:\delta,
 \qquad \operatorname{Tr}B\le s(1+R)^2.               \tag{A017.4}
\]
For positive finite matrices $A,B$, differentiation along their
positive linear interpolation proves
\[
 |\det(I-A)-\det(I-B)|
 \le\|A-B\|_1e^{\max(\operatorname{Tr}A,
                                    \operatorname{Tr}B)}.
\]
To see this even when the interpolation passes through eigenvalue
one, use the polynomial adjugate derivative. If its positive
eigenvalues are $\lambda_i$, the adjugate operator norm is bounded
by $\prod_i(1+\lambda_i)\le e^{\sum_i\lambda_i}$, since
$|1-\lambda_i|\le1+\lambda_i$. Pairing with $A-B$ uses its
trace norm. For positive trace-class operators, compress by common
finite-rank projections and pass to the limit. The determinant
convergence follows from its absolutely convergent exterior-power
series, bounded by $\sum_k\|A\|_1^k/k!$; trace-norm convergence
of the compressions follows first for finite-rank approximants and
then by approximation. Consequently the certificate's more generous
error bound
\[
 \delta\exp(1+s+s(1+R)^2)                             \tag{A017.5}
\]
is valid.

The finite determinant is evaluated without quadrature. Let
\[
 G_{kl}=\int_{-1/2}^{1/2}u^{k+l}du
 =\begin{cases}
  [2^{k+l}(k+l+1)]^{-1},&k+l\ \text{even},\\
 0,&k+l\ \text{odd},
 \end{cases}
 \quad D_{kk}=\frac{(2\pi is)^k}{k!},\quad0\le k,l\le p.
\]
Sylvester's finite-rank determinant identity gives
\[
 \det(I-B)=\det(I-sDGD^*G).
\]
The even and odd index blocks are real and separate. In those
blocks $(DGD^*)_{kl}=(-1)^{(k-l)/2}
 (2\pi s)^k(2\pi s)^lG_{kl}/(k!l!)$.
The program uses exactly these rational-moment matrices with Arb
balls for $\pi$, and enlarges each determinant ball by (A017.5).
For $p=100$ it certifies the six integer gap probabilities, beginning
with
\begin{align*}
 E(1)&=0.17021742137918523073265305256\ldots,\\
 E(2)&=0.00349732514916909769841342811\ldots.
\end{align*}
and ending with
\[
 E(5)=1.55338899144701217136\ldots\,10^{-14},\qquad
 E(6)=1.89553739144083908383\ldots\,10^{-20}.
\]
Their complete intervals and analytic approximation errors are in
\path{cue_agreement_certificate.json}. The odd partial sum
through five gives a lower bound for $A_{\rm adj}$, and the even
partial sum through six an upper bound. Outward rounding gives
exactly (A017.2). The included program asserts both displayed
decimal inequalities as ball comparisons.
\end{proof}

The determinant is the exterior-algebra trace of the sine projection
restricted to an interval. The finite certificate uses an actual
polynomial feature space and its two Gram matrices, not a finite
field or a claimed arithmetic motive. This supplies a rigorously
computable model statistic, with an analytic approximation error;
it does not certify that arithmetic labels have this limit.


\section{Actual motion, ramified branches, and contour traces}

\begin{theorem}[NN-A004: all shifted-series jets and factorization obstruction]
Let $L(s)=\sum_{n\ge1}a_n n^{-s}$ converge absolutely for $\Re s>1$
and continue holomorphically near $\rho$, where $\Re\rho>0$ and
$L(\rho)=0$. In a sufficiently small neighborhood of $(\rho,0)$, set
\[
 L_t(s)=\sum_{k\ge0}\frac{(-1)^k(s)_k}{k!}L(s+k)t^k.
\]
This is holomorphic and agrees with $\sum a_n(n+t)^{-s}$ in the initial
half-plane. If $\rho$ is simple, its unique analytic zero branch has
\[
 v=\rho'(0)=\frac{\rho L(\rho+1)}{L'(\rho)},
\quad
 \rho''(0)=-\frac{L''(\rho)v^2
 -2\{L(\rho+1)+\rho L'(\rho+1)\}v
 +\rho(\rho+1)L(\rho+2)}{L'(\rho)}.
\]
All higher coefficients $\rho(t)=\rho+\sum_{n\ge1}c_nt^n$ are determined
without a new existence hypothesis by
\[
 c_n=-\frac{[t^n]L_t(\rho+\sum_{j<n}c_jt^j)}{L'(\rho)}.
\]
For real $t\downarrow0$, the first nonzero $\Re c_n$ determines the
normal side of the critical line. If every $\Re c_n=0$, the branch stays
on that line near zero. A nonzero complex velocity alone does not exclude
normal tangency.

If $L=FG$, $F(\rho)=0$, $F'(\rho)G(\rho)\ne0$, then the shifts of the
Dirichlet coefficients satisfy
\[
 v_L=v_F\,\frac{G(\rho+1)}{G(\rho)}.
\]
Consequently factorwise GRH does not give a factorwise rule for the motion
labels: the nonvanishing factor has a retained complex multiplier.
\end{theorem}
\begin{proof}
For $s$ in a fixed compact neighborhood of $\rho$, choose $k_0$ such that
$\Re(s+k_0)>1$. Absolute convergence gives a uniform bound for $L(s+k)$
for $k\ge k_0$, while $(s)_k/k!=O(k^C)$ uniformly on this compact set.
The displayed series and its local holomorphic derivatives converge for
$|t|<1$. There are only finitely many lower $k$, and their neighborhoods
can avoid any poles. Expanding $(n+t)^{-s}$ by the binomial theorem and
using absolute convergence proves equality in the initial half-plane.
The implicit function theorem applies because $L'(\rho)\ne0$.
The first two derivatives of $L_t(\rho(t))=0$ use
$\partial_tL_t|_0=-sL(s+1)$ and
$\partial_t^2L_t|_0=s(s+1)L(s+2)$, and give the stated formulas.
Adding $c_nt^n$ to the already constructed truncation changes its degree
$n$ coefficient by exactly $L'(\rho)c_n$, proving the recursion. A
real analytic function is either identically zero locally or has the sign
of its first nonzero coefficient for small positive argument. Finally
$L'(\rho)=F'(\rho)G(\rho)$ and
$L(\rho+1)=F(\rho+1)G(\rho+1)$ prove the product formula.
\end{proof}

This applies to $\zeta$, Dirichlet $L$-functions, and Dedekind $\zeta_K$.
A supplies the critical line for $\zeta$ and Dirichlet factors; for abelian
$K/\mathbb Q$ the standard Dirichlet factorization transfers that location
to $\zeta_K$. General $K$ requires the stipulated Hecke input M, or a
separate GRH hypothesis. None of these inputs supplies simplicity. At a
nontrivial critical zero, the Euler product at $\rho+1$ shows $v\ne0$
when the zero is simple; it gives no sign for $\Re v$.
The geometric object here is the analytic zero divisor in the two
coordinates $(s,t)$ and its sections over the parameter disc. It is not
an algebraic family over a finite field, and its complex multipliers have
not been identified with Frobenius eigenvalues.

\begin{theorem}[NN-A005: multiple-zero branches and the residue trace]
Suppose instead that $L$ has multiplicity $m\ge2$ at $\rho$ and
$\rho L(\rho+1)\ne0$. Put
\[
 C=\frac{m!\rho L(\rho+1)}{L^{(m)}(\rho)}.
\]
The $m$ local branches on the ramified parameter $u=t^{1/m}$ satisfy
\[
 s_j(u)=\rho+c\omega_m^j u+O(u^2),\qquad c^m=C.
\]
In particular they are not individually differentiable as holomorphic
functions of $t$ at zero. Their symmetric trace is analytic. More
generally, let $\Gamma$ be a positively oriented finite union of simple
contours, bounding a region where $L_t$ is holomorphic and has no boundary
zeros for small $t$. If $P$ is holomorphic on that region, then
\[
 \left.\frac d{dt}\sum_{L_t(z)=0}P(z)\right|_{t=0}
 =\frac1{2\pi i}\int_\Gamma
 P'(s)\frac{sL(s+1)}{L(s)}\,ds.
\]
Zeros in this formula are counted with multiplicity. For $P(s)=s$, the
contribution of a multiplicity-$m$ zero is
\[
 \operatorname{Res}_{s=\rho}\frac{sL(s+1)}{L(s)}
 =\frac1{(m-1)!}
 \left.\left(\frac d{ds}\right)^{m-1}
 \frac{sL(s+1)}{L(s)/(s-\rho)^m}\right|_{s=\rho}.
\]
Thus a contour mean remains meaningful without simple zeros, whereas a
sum of the simple-zero velocities does not.
\end{theorem}
\begin{proof}
With $a=L^{(m)}(\rho)/m!$, the quotient
$u^{-m}L_{u^m}(\rho+uy)$ extends holomorphically to $u=0$ and there
equals $ay^m-\rho L(\rho+1)$. Its $m$ roots are nonzero and simple.
The implicit function theorem gives the $m$ analytic $y_j(u)$, and
Rouche's theorem counts all roots in the local cluster. The nonzero
leading terms $c\omega_m^j t^{1/m}$ exclude holomorphic individual
branches in $t$.
The argument principle writes the symmetric sum as
$(2\pi i)^{-1}\int_\Gamma P(s)\partial_s\log L_t(s)\,ds$.
The quotient $\partial_tL_t/L_t$ is analytic in $t$ on the compact
contour, so differentiation under the integral and integration by parts
give $-(2\pi i)^{-1}\int P'(s)\partial_tL_t/L_t$. Insert
$\partial_tL_t|_0=-sL(s+1)$. The residue expression is the ordinary
formula for a pole of order at most $m$.
\end{proof}

For $L=\zeta$, the right-hand contour integrand has the fully explicit
Dirichlet series, initially on $\Re s>1$,
\[
 \frac{s\zeta(s+1)}{\zeta(s)}
 =s\sum_{n\ge1}\frac{b(n)}{n^s},\qquad
 b(n)=\frac1n\prod_{p\mid n}(1-p).
\]
Indeed convolution with $1/\zeta(s)=\sum\mu(n)n^{-s}$ gives
$b(n)=n^{-1}\sum_{d\mid n}d\mu(d)$. For general $L$, replace
$b(n)$ by $\sum_{d\mid n}a_d d^{-1}a^{-1}_{n/d}$ whenever the
Dirichlet inverse converges absolutely. This is an exact residue
formula, not an unproved asymptotic on contours passing through zeros.
In the zero-divisor geometry, symmetric sums are functions on the relative
degree-$m$ divisor (or its symmetric power); individual sections require
the stated ramified base change. No complexity bound follows from the
existence of the contour integral.

\begin{proposition}[NN-A010: the completed Hurwitz deformation has extra poles]
Let $\Lambda_t(s)=\pi^{-s/2}\Gamma(s/2)\zeta(s,1+t)$. For $m\ge1$,
\[
 \operatorname{Res}_{s=-2m}\Lambda_t(s)
 =-\frac{2(-1)^m\pi^m}{m!(2m+1)}B_{2m+1}(1+t).
\]
At $t=0$ these residues vanish; their first $t$ derivatives are
$-2(-1)^m\pi^mB_{2m}/m!\ne0$. At $t=1$, every such residue is
$-2(-1)^m\pi^m/m!$. Consequently the rational principal-part
bookkeeping using the poles at zero and one is not the complete
meromorphic divisor of $\Lambda_t$. It does not provide a meromorphic
function on $\mathbb P^1$ for the full completion.
\end{proposition}
\begin{proof}
The residue of $\Gamma(s/2)$ at $s=-2m$ is $2(-1)^m/m!$.
Use the exact Hurwitz value
$\zeta(-2m,a)=-B_{2m+1}(a)/(2m+1)$ (DLMF 25.11.14).
Bernoulli differentiation and $B_{2m}(1)=B_{2m}$ give the derivative.
For $t=1$, $\zeta(s,2)=\zeta(s)-1$, so its value at every negative
even integer is $-1$. A meromorphic function on compact
$\mathbb P^1$ has finitely many poles, whereas this example has
infinitely many. This proves the asserted scope restriction.
\end{proof}


\section{Central Dedekind branches and an unconditional arithmetic imbalance}
\label{sec:analytic-central-puiseux}

\begin{theorem}[NN-A018: the second central Puiseux coefficient]
Let $K$ have degree $n=r_1+2r_2$ and absolute discriminant $D_K$.
Write $L=\zeta_K$ and suppose its zero at $\rho=1/2$ has order
$m\ge2$. Put
\[
 Q_K(s)=s(s-1)D_K^{s/2}\Gamma_{\mathbb R}(s)^{r_1}
           \Gamma_{\mathbb C}(s)^{r_2},\qquad
 \Gamma_{\mathbb R}(s)=\pi^{-s/2}\Gamma(s/2),\quad
 \Gamma_{\mathbb C}(s)=(2\pi)^{-s}\Gamma(s),
\]
so that $\xi_K=Q_KL$ is entire and satisfies
$\xi_K(s)=\xi_K(1-s)$. The harmless constant often inserted in
$\Gamma_{\mathbb C}$ changes none of the formulas below. Set
\[
 a_m=\frac{L^{(m)}(1/2)}{m!},\qquad
 C_K=2+\frac{\zeta_K'}{\zeta_K}(3/2)
                 +\frac{Q_K'}{Q_K}(1/2).
\]
The shifted Dirichlet-series deformation from NN-A004 has exactly $m$
local branches, counted with multiplicity. For positive real $t$ they
are distinct for all sufficiently small $t$ and have expansions
\begin{equation}\label{eq:central-second-jet}
 s_j(t)=\frac12+c_jt^{1/m}
              +\frac{C_K}{m}c_j^2t^{2/m}+O(t^{3/m}),\qquad
 c_j^m=\frac{\zeta_K(3/2)}{2a_m}.
\end{equation}
The central order $m$ is even. If $a_m<0$, let
$a=(\zeta_K(3/2)/(2|a_m|))^{1/m}$ and choose
$c_j=a\exp((2j+1)\pi i/m)$. If $m\equiv0\pmod4$, exactly $m/2$
branches lie on each side of the critical line for small positive $t$.
If $m\equiv2\pmod4$, precisely two leading directions are vertical.
When $C_K>0$ these two branches bend left, and the left/right counts
are $(m+2)/2,(m-2)/2$; when $C_K<0$ both bend right and these counts
are reversed. No conclusion about their normal direction is made
from this coefficient when $C_K=0$.

The hypothesis $a_m<0$ follows from GRH for $\zeta_K$, hence from A
for abelian $K$ and conditionally from M in general. It can also be
established without GRH by a local factorization, as in the example
below.
\end{theorem}

\begin{proof}
Evenness of $\xi_K(1/2+z)$ and the nonvanishing of $Q_K(1/2)$ imply
that $m$ is even. Write
\[
 L(1/2+z)=a_mz^m+a_{m+1}z^{m+1}+O(z^{m+2}).
\]
The coefficient of $z^{m+1}$ in $Q_K(1/2+z)L(1/2+z)$ vanishes, so
\begin{equation}\label{eq:central-unit-jet}
 \frac{a_{m+1}}{a_m}=-\frac{Q_K'}{Q_K}(1/2).
\end{equation}
The deformation has the local expansion
\[
 L_t(1/2+z)=a_mz^m+a_{m+1}z^{m+1}
       -t(B_0+B_1z)+O(z^{m+2})+O(tz^2)+O(t^2),
\]
where $B_0=L(3/2)/2>0$ and
$B_1=L(3/2)+L'(3/2)/2$. By NN-A005, the base change $t=u^m$
produces $m$ holomorphic branches $z=c_ju+d_ju^2+O(u^3)$,
with distinct $c_j$ and $a_mc_j^m=B_0$. The coefficient of $u^{m+1}$
is
\[
 ma_mc_j^{m-1}d_j+a_{m+1}c_j^{m+1}-B_1c_j=0.
\]
Because $2m\ge m+2$, the $O(t^2)$ term does not contribute.
Using $a_mc_j^m=B_0$ and \eqref{eq:central-unit-jet} gives
\[
 d_j=\frac{c_j^2}{m}
       \left(\frac{B_1}{B_0}-\frac{a_{m+1}}{a_m}\right)
     =\frac{C_K}{m}c_j^2.
\]
This proves \eqref{eq:central-second-jet} and the simplicity of
the branches at small nonzero $t$.

For $a_m<0$ the displayed $m$ directions are the roots of a negative
real number. A direction is vertical exactly when $m\equiv2\pmod4$,
and then the two vertical directions are $\pm ia$. Their second
normal coefficient is $-a^2C_K/m$, the same for both. The other
directions pair across the imaginary axis. This gives all the
claimed counts. The $O(t^{3/m})$ error is smaller than the nonzero
first or second normal term, so the statements hold on a common
positive interval for all $m$ branches.

Under GRH, $\xi_K$ has no real zero in $(1/2,\infty)$ and is positive
on $(1,\infty)$. Its value at $1$ is the positive Dedekind residue
times positive factors; thus it remains positive on $(1/2,\infty)$.
Its first nonzero coefficient at $1/2$ is consequently positive.
Since $Q_K(1/2)<0$, this gives $a_m<0$.
\end{proof}

\begin{proposition}[An effective conductor criterion]
Let
\[
 a_{\mathbb R}=\frac12\psi(1/4)-\frac12\log\pi,
 \qquad C_*=-a_{\mathbb R}-\frac{\zeta'}{\zeta}(3/2).
\]
Then
\begin{equation}\label{eq:central-C-bound}
 C_K\ge2+\frac12\log D_K-nC_*+\frac\pi2r_2.
\end{equation}
In particular, root discriminant $D_K^{1/n}\ge4371$ implies $C_K>0$.
Consequently, whenever $a_m<0$ and $m\equiv2\pmod4$, the two
initially vertical branches bend left for such a field.
\end{proposition}

\begin{proof}
The logarithmic derivative of $s(s-1)$ vanishes at $1/2$.
The gamma reflection and duplication formulas give
\[
 \frac{Q_K'}{Q_K}(1/2)=\frac12\log D_K
                   +na_{\mathbb R}+\frac\pi2r_2.
\]
For a rational prime $p$, write $f_{\mathfrak p}$ and
$e_{\mathfrak p}$ for the residue and ramification degrees.
For $\sigma>1$,
\[
 -\frac{\zeta_K'}{\zeta_K}(\sigma)
 =\sum_p\log p\sum_{\mathfrak p\mid p}
       \frac{f_{\mathfrak p}}{p^{f_{\mathfrak p}\sigma}-1}
 \le n\sum_p\frac{\log p}{p^\sigma-1}
 =-n\frac{\zeta'}{\zeta}(\sigma).
\]
Here $p^{f\sigma}-1\ge p^\sigma-1$ and
$\sum f_{\mathfrak p}\le\sum e_{\mathfrak p}f_{\mathfrak p}=n$.
This proves \eqref{eq:central-C-bound}. The saved Arb certificate
gives $e^{2C_*}<4371$ by a strict interval comparison; hence the
right side of \eqref{eq:central-C-bound} is positive at the stated
root discriminant. One may instead use the exact threshold
$\log D_K>2nC_*-\pi r_2-4$ without decimal constants.
\end{proof}

\begin{corollary}[An unconditional arithmetic example with two excess right branches]
Let
\[
 F=\mathbb Q(\sqrt5,\sqrt{41}),\qquad
 N=F\left(\sqrt{\frac{(5+\sqrt5)(41+\sqrt{205})}{4}}\right).
\]
The central order $m=\operatorname{ord}_{s=1/2}\zeta_N(s)$ is
positive and congruent to $2$ modulo $4$. For the shifted
Dirichlet-series deformation of $\zeta_N$, and every sufficiently
small positive $t$, its $m$ branches from the central zero consist
of exactly $(m+2)/2$ right branches and $(m-2)/2$ left branches.
The two initially vertical branches both bend right. This assertion
uses neither A nor M and does not assume $m=2$.
\end{corollary}

\begin{proof}
The field, its totally real tame $Q_8$ Galois group, ramification
only at $5,41$, and the nonfreeness of $\mathcal O_N$ over
$\mathbb Z[Q_8]$ are Martinet's Example~3, printed p.~407, with
Theorem~IV.1 on p.~406. They are classical inputs, already used for
central existence in prior MT-ART10; central existence is not the
new claim here. Fr\"ohlich's Theorem~1, printed p.~82 of his 1974
paper, identifies this nontrivial module class with root number
$-1$ for the unique irreducible character $\chi$ of degree two.

Choose a cyclic subgroup of order four. Its faithful linear
character induces $\chi$, so Artin induction identifies $L(s,\chi)$
with a nontrivial finite-order Hecke $L$-function over the quadratic
fixed field. It is entire. Its selfdual functional equation with
sign $-1$ makes its central order $\ell$ positive and odd. Its
Dirichlet coefficients are real, so its leading central coefficient
$b_\ell$ is a nonzero real number. The regular representation gives
the exact all-prime factorization
\[
 \zeta_N(s)=\zeta(s)L(s,\chi_5)L(s,\chi_{41})L(s,\chi_{205})
                     L(s,\chi)^2.
\]
The included 70-digit Arb certificate, using the finite Hurwitz
formula, proves
\[
 0.23<L(1/2,\chi_5)<0.24,\quad
 1.10<L(1/2,\chi_{41})<1.11,\quad
 1.17<L(1/2,\chi_{205})<1.18.
\]
Also $\zeta(1/2)<0$: the alternating eta series is positive by
pairing successive terms and $1-2^{1-s}<0$ for $0<s<1$.
It follows that $m=2\ell\equiv2\pmod4$ and the leading coefficient
of $\zeta_N$ is
$a_m=\zeta_F(1/2)b_\ell^2<0$, without a zero-location hypothesis.

The discriminant is $D_N=205^6$. For completeness, at either
$p=5$ or $41$ the inertia image in
$\operatorname{Gal}(F/\mathbb Q)=Q_8/\{\pm1\}$ is nontrivial.
Tame inertia in $Q_8$ is cyclic; any element with nontrivial image
in that quotient has order four. Thus the inertia order is four
and the tame discriminant exponent is $8(1-1/4)=6$. There are no
other ramified primes. This also agrees with Martinet's
Propositions~III.1--III.2, printed pp.~404--405.

The Euler product gives $\zeta_N'/\zeta_N(3/2)<0$ strictly.
Since $N$ is totally real of degree eight, the same certificate gives
\[
 C_N<2+3\log205+8a_{\mathbb R}
       <-3.5197<0.
\]
The theorem now gives the stated right/left branch counts.
\end{proof}

\paragraph{Source and verification locators.}
Martinet, \emph{Modules sur l'alg\`ebre du groupe quaternionien},
Ann. Sci. ENS~4 (1971), 399--408,
DOI~10.24033/asens.1216, Propositions~III.1--III.2,
Theorem~IV.1 and Example~3. The original p.~407 image was checked
again here: its second radical is $\sqrt{5\cdot41}$.
Fr\"ohlich, \emph{The Galois module structure of algebraic integer
rings in fields with generalised quaternion group},
M\'em. SMF~37 (1974), 81--86, DOI~10.24033/msmf.132,
Theorem~1 p.~82. This is an imported classical theorem; its
referenced 1972 proof is not claimed to have been re-proved here.
The short local deduction above is complete.
The script \path{check_central_jets.py} uses exact Kronecker
characters and
\[
 L(s,\chi_d)=d^{-s}\sum_{a=1}^d\chi_d(a)\zeta(s,a/d)
\]
with complex-ball Hurwitz evaluation, and writes every sign interval
to \path{central_jets_certificate.json}.

\paragraph{Geometric content and limitation.}
The central zero has local divisor length $m$; after the ramified
base change $t=u^m$, the normalization has $m$ sections. The
coefficient of $u^2$ depends on the completion unit and the shifted
Euler value, data not retained by the reduced support. On the chosen
positive real parameter ray the theorem assigns a canonical
multiset of sides. It does not assign a single unramified Boolean
label to the central support point, nor a colouring invariant under
complex monodromy. In the example, a nontrivial integral
$\mathbb Z[Q_8]$-module class supplies the odd Artin order, while
the analytic second jet supplies the excess of two right branches.
Those are precise distinct steps; no finite pure motive or
complexity lower bound is inferred.


\section{A uniform central Dirichlet second-jet sign under A}
\label{sec:analytic-central-dirichlet-sign}

Let $\chi$ be a nonprincipal real primitive Dirichlet character of
conductor $q$, and put $a\in\{0,1\}$ for its parity,
$\chi(-1)=(-1)^a$. Throughout this section the completion is
\begin{equation}\label{eq:analytic-dirichlet-completion}
 Q_\chi(s)=\left(\frac q\pi\right)^{(s+a)/2}
                \Gamma\left(\frac{s+a}{2}\right),\qquad
 \Lambda_\chi(s)=Q_\chi(s)L(s,\chi).
\end{equation}
This fixes every gamma and conductor factor used in the local
coefficient. The principal primitive character, whose completion
has poles, is excluded.

\begin{theorem}[NN-A032: positive central coefficient and exact side counts]
\label{thm:analytic-central-dirichlet-sign}
Let $m=\operatorname{ord}_{s=1/2}L(s,\chi)\ge0$. It is even.
The coefficient $C_\chi$ defined below satisfies, unconditionally,
$C_\chi>m+2-\pi/2$ for $a=0$ and
$C_\chi>m+\pi/2$ for $a=1$.
Under A, list the positive ordinates of all noncentral zeros
with their multiplicities. Then the exact identity is
\begin{equation}\label{eq:analytic-dirichlet-C-positive}
 \begin{split}
 C_\chi&:=2+\frac{L'}L(3/2,\chi)
                         +\frac{Q_\chi'}{Q_\chi}(1/2)\\
 &=m+\sum_{\gamma>0}\frac2{1+\gamma^2}
     +\begin{cases}2-\pi/2,&a=0,\\ \pi/2,&a=1.
       \end{cases}
 \end{split}
\end{equation}
The sum converges absolutely. In particular,
\[
 C_\chi\ge m+2-\pi/2>0\quad(a=0),\qquad
 C_\chi\ge m+\pi/2>0\quad(a=1).
\]

Assume A for the rest of the theorem. If $m>0$, write
$A_m=L^{(m)}(1/2,\chi)/m!$. Then $A_m>0$. For the shifted
Dirichlet deformation of NN-A004,
\[
 L_t(s,\chi)=\sum_{n\ge1}\frac{\chi(n)}{(n+t)^s},
\]
continued analytically in $s$ and locally in $t$, its $m$ zero
branches near the center have, for positive $t$, the expansions
\begin{equation}\label{eq:analytic-dirichlet-central-branches}
 s_j(t)=\frac12+c_jt^{1/m}
                  +\frac{C_\chi}{m}c_j^2t^{2/m}
                  +O(t^{3/m}),\qquad
 c_j^m=\frac{L(3/2,\chi)}{2A_m}>0.
\end{equation}
For all sufficiently small $t>0$, the exact counts are
\[
 \begin{array}{c|cc}
 &\Re s_j(t)<1/2&\Re s_j(t)>1/2\\ \hline
 m\equiv2\pmod4&m/2&m/2\\
 m\equiv0\pmod4,\ m\ge4&m/2+1&m/2-1.
 \end{array}
\]
In the second row the two branches with purely imaginary first
coefficient both bend strictly left at second order. No central
zero is asserted to exist; the side-count statement is
conditional on $m>0$. Neither quadratic central nonvanishing nor
simplicity follows from this theorem.
\end{theorem}

\begin{proof}
\emph{The exact functional-equation sign and entire completion.}
A nonprincipal real primitive character is quadratic. Its
quadratic Gauss sum has the sign
$\tau(\chi)=i^a\sqrt q$, so the root number in the convention
\eqref{eq:analytic-dirichlet-completion} is $+1$. A primary
locator, including the sign rather than just its square, is
B.~H.~Gross, \emph{Some remarks on signs in functional equations},
section~1, equations~(1.3)--(1.6), p.~2. There the sign is proved
by writing the completed quadratic $L$-function as the quotient
of the completed zeta functions of the quadratic field and
$\mathbb Q$. DLMF~25.15.5 gives the primitive functional equation
with the Gauss sum explicitly.

For completeness, the usual order-one analytic input can be
seen directly with these normalizations. Put
\[
 \theta_\chi(x)=\sum_{n\in\mathbb Z}n^a\chi(n)e^{-\pi n^2x/q},
 \qquad x>0,
\]
where the $n=0$ term is zero. Poisson summation in the residue
classes modulo $q$, using
\[
 \sum_{b\bmod q}\chi(b)e^{2\pi inb/q}=\tau(\chi)\chi(n),
\]
and the Gaussian Fourier transform gives
\[
 \theta_\chi(x)=\frac{(-i)^a\tau(\chi)}{\sqrt q}
        x^{-a-1/2}\theta_\chi(1/x)
              =x^{-a-1/2}\theta_\chi(1/x).
\]
For $a=1$ the needed transform follows by differentiating the
Gaussian transform once; this accounts for the factor $-i$.
Termwise Mellin integration for $\Re s>1$, followed by the
substitution $x\mapsto1/x$ on $(0,1)$, yields
\begin{equation}\label{eq:analytic-dirichlet-theta-entire}
 \Lambda_\chi(s)=\frac12\int_1^\infty\theta_\chi(x)
       \left(x^{(s+a)/2}+x^{(1-s+a)/2}\right)\frac{dx}{x}.
\end{equation}
The theta series decays exponentially for $x\ge1$; hence this
formula is entire, is invariant under $s\mapsto1-s$, and gives
$\log\max_{|s|\le R}|\Lambda_\chi(s)|=O_q(R\log(R+2))$ by
comparison with a gamma integral. Thus the entire function has
order at most one. Its zeros are precisely the nontrivial
$L$-zeros: the gamma poles are canceled by the trivial zeros,
with nonzero resulting values by the functional equation and
the Euler product on the opposite half-plane. At the endpoint
$s=0$ for an even character, the reflected point is $s=1$;
here one also uses Dirichlet's nonvanishing theorem
$L(1,\chi)\ne0$ (DLMF~25.15.9, with source Apostol, 1976,
p.~149), rather than the Euler product in its open half-plane.

\emph{Unconditional sign of the second-jet coefficient.}
Write $F(z)=\Lambda_\chi(1/2+z)$. All its zeros
$\alpha=\rho-1/2$ lie in $|\Re\alpha|\le1/2$ by the classical
critical-strip theorem, independently of A. The paired
Hadamard argument, exactly as in NN-A033, gives
\[
 \frac{F'(1)}{F(1)}
       =m+\sum_{\{\alpha,-\alpha\},\ \alpha\ne0}
                         \frac2{1-\alpha^2}.
\]
The series is absolutely convergent and is grouped under
conjugation. If $\alpha=b+ic$, then
$\Re(1-\alpha^2)=1-b^2+c^2>0$, whence
$\Re(1-\alpha^2)^{-1}>0$. Purely real and purely imaginary
pairs also contribute positive terms. There are noncentral
zeros, since the primitive $L$-function has infinitely many
zeros in the strip (DLMF~25.15(ii)). Hence $F'(1)/F(1)>m$.
The completion identity and the gamma differences calculated
below give, with no zero-location assumption,
\[
 C_\chi=\frac{F'(1)}{F(1)}+
       \begin{cases}2-\pi/2,&a=0,\\ \pi/2,&a=1.
       \end{cases}
\]
This proves the unconditional strict inequalities. It does not
fix the sign of the central leading coefficient without A.

\emph{Real positivity and the paired zero product under A.}
For real $\sigma>1$, every factor of the Euler product
\[
            L(\sigma,\chi)=\prod_p(1-\chi(p)p^{-\sigma})^{-1}
\]
is positive; this is DLMF~25.15.2. The gamma and conductor
factors are also positive. Assumption A excludes real zeros in
$(1/2,\infty)$, so continuity gives
$\Lambda_\chi(\sigma)>0$ throughout this interval. The
sign-$+1$ functional equation
makes $F$ even, so its central order $m$ is even. Its central
leading coefficient is positive, by its positivity for $z>0$.
Dividing by the positive number $Q_\chi(1/2)$ gives $A_m>0$
when $m>0$.

Under A all noncentral zeros of $F$ are $\pm i\gamma$. The
paired genus-one Hadamard product is
\begin{equation}\label{eq:analytic-dirichlet-positive-product}
 F(z)=Bz^m\prod_{\gamma>0}\left(1+\frac{z^2}{\gamma^2}\right),
                         \qquad B>0.
\end{equation}
Indeed the order-one factorization initially permits an
exponential $e^{b z}$; evenness forces $b=0$, and the paired
genus-one exponentials cancel. The sum of reciprocal squared
ordinates converges, so the product and its logarithmic
derivative converge normally on compact sets avoiding zeros.
At $z=1$ this gives the exact positive formula
\begin{equation}\label{eq:analytic-dirichlet-log-derivative}
       \frac{\Lambda_\chi'}{\Lambda_\chi}(3/2)
                    =m+\sum_{\gamma>0}\frac2{1+\gamma^2}.
\end{equation}

\emph{Cancellation of the conductor and the gamma differences.}
Since $L'/L=\Lambda_\chi'/\Lambda_\chi-Q_\chi'/Q_\chi$,
\[
 C_\chi=\frac{\Lambda_\chi'}{\Lambda_\chi}(3/2)
          +2+\frac12\left[
              \psi\left(\frac{1/2+a}{2}\right)
             -\psi\left(\frac{3/2+a}{2}\right)\right].
\]
The conductor terms have canceled exactly. The digamma
reflection and recurrence identities, DLMF~5.5.4 and~5.5.2,
give
\[
 \psi(1/4)-\psi(3/4)=-\pi,\qquad
 \psi(3/4)-\psi(5/4)=\pi-4.
\]
Substitution in \eqref{eq:analytic-dirichlet-log-derivative}
proves \eqref{eq:analytic-dirichlet-C-positive} and both strict
sign bounds. This uses the actual completion unit, without a
conductor-size estimate or a numerical verification.

\emph{The shifted Puiseux coefficient.}
Assume $m>0$, hence $m\ge2$. Write $x=s-1/2$ and
\[
 L(1/2+x,\chi)=A_mx^m+A_{m+1}x^{m+1}+O(x^{m+2}).
\]
Evenness of $\Lambda_\chi$ gives
$A_{m+1}/A_m=-Q_\chi'/Q_\chi(1/2)$. The shifted series of
NN-A004 is locally
\[
 L_t(s,\chi)=L(s,\chi)-sL(s+1,\chi)t+O(t^2).
\]
Put $t=u^m$ and $x=cu+du^2+O(u^3)$. The coefficient of $u^m$
gives $A_mc^m=L(3/2,\chi)/2$. Its $m$ roots are distinct.
The coefficient of $u^{m+1}$ then gives
\[
 mA_mc^{m-1}d+A_{m+1}c^{m+1}
       -\left(L(3/2,\chi)+\tfrac12L'(3/2,\chi)\right)c=0,
\]
or $d=(C_\chi/m)c^2$. Since $2m\ge m+2$, the $O(t^2)$ term
does not affect these two equations. The analytic implicit
function theorem after dividing by $u^m$ at each distinct
leading root gives all $m$ analytic branches in $u$, with the
stated error. This proves
\eqref{eq:analytic-dirichlet-central-branches}.

Finally put $R=(L(3/2,\chi)/(2A_m))^{1/m}>0$. The leading
coefficients are $c_j=Re^{2\pi ij/m}$. If $m\equiv2\pmod4$,
none is purely imaginary and exactly half have each sign of
their real part. If $m\equiv0\pmod4$, precisely two are
$\pm iR$; each has second coefficient
$-C_\chi R^2/m<0$. All other branches keep the sign of their
nonzero first real coefficient. Counting them gives the two
rows of the table, with no branch left on the line for
sufficiently small positive $t$.
\end{proof}

\paragraph{Geometric translation.}
The quadratic character is the rank-one Artin representation of
its quadratic finite etale cover away from the conductor; its
root number is the actual $+1$ in Gross's equations~(1.3)--(1.6).
At a hypothetical central zero of length $m$, the shifted
analytic divisor is finite flat of rank $m$ near the parameter
origin. The ramified base change $t=u^m$ separates its leading
sections. Complex conjugation pairs the two sections with
leading coefficients $\pm iR$, and the strictly negative real
second coefficient puts both on the left of the chosen real
normal line. This is a statement about the full local analytic
algebra and its completion unit, rather than a single label on
its reduced support. The finite-free discriminant/etale
criterion is SGA~1, Expos\'e~I, section~4, Proposition~10; no
arithmetic family realizing the shifted parameter, central
nonvanishing theorem, or positive integral cohomology is inferred.


\section{An unconditional central-order certificate and exactly two right branches}
\label{sec:analytic-central-order-certificate}

The next argument uses only the classical critical strip, not a
Riemann hypothesis. It complements NN-A018 by determining the
central multiplicity of its explicit quaternionic example.
The central order itself is treated as a classical type of
calculation, without a historical novelty claim. The new
catalogue consequence is the complete two-branch conclusion for
the specified shifted deformation, together with the fully
written rational certificate below.

For a number field $K$ of signature $(r_1,r_2)$, use the exact
pole-removed completion
\begin{equation}\label{eq:analytic-order-completion}
 \xi_K(s)=s(s-1)D_K^{s/2}\pi^{-r_1s/2}(2\pi)^{-r_2s}
        \Gamma(s/2)^{r_1}\Gamma(s)^{r_2}\zeta_K(s).
\end{equation}
It is real entire of order one, is invariant under $s\mapsto1-s$,
and its zeros are the nontrivial zeros in $0\le\Re s\le1$,
with multiplicity. A primary locator for these analytic inputs
and the completion is Hasanalizade--Shen--Wong,
\emph{Counting zeros of Dedekind zeta functions},
\emph{Math. Comp.}~91 (2022), section~2.1, equations~(2.1)--(2.2),
p.~279. Its complex-place gamma convention differs from
\eqref{eq:analytic-order-completion} by the positive constant
$2^{r_2}$, by the duplication formula; every logarithmic
derivative here is therefore unchanged. The Euler product gives
nonvanishing in $\Re s>1$, and the functional equation and the
nonzero residues at $0,1$ remove any additional completed zeros.

\begin{theorem}[NN-A033: a positive rational test for central order]
\label{thm:analytic-central-order-test}
Put $F_K(z)=\xi_K(1/2+z)$,
$G_K(z)=F_K'(z)/F_K(z)$, and
$m_K=\operatorname{ord}_{s=1/2}\zeta_K(s)$. For every real
$\delta>1/\sqrt2$,
\begin{equation}\label{eq:analytic-T2-order}
 m_K\le T_{2,K}(\delta):=
 \frac{\delta G_K(\delta)-\delta^2G_K'(\delta)}2
 =m_K+\sum_{\{\alpha,-\alpha\}}
                 \frac{2\delta^4}{(\delta^2-\alpha^2)^2}.
\end{equation}
Here $\alpha=\rho-1/2\ne0$, one representative is chosen for
each pair $\{\alpha,-\alpha\}$, multiplicities are retained,
and conjugate terms are grouped. Each grouped contribution is
strictly positive. The series is absolutely convergent.

Set $\sigma=1/2+\delta$ and
\[
 B_K(s)=\frac1s+\frac1{s-1}+\frac12\log D_K
       -\frac{r_1}2\log\pi-r_2\log(2\pi)
       +\frac{r_1}2\psi(s/2)+r_2\psi(s).
\]
Then the entirely archimedean upper bound is
\begin{equation}\label{eq:analytic-T2-arch}
 T_{2,K}(\delta)<
 U_K(\delta):=
       \frac{\delta B_K(\sigma)-\delta^2B_K'(\sigma)}2.
\end{equation}
Both assertions are unconditional.
\end{theorem}

\begin{proof}
Evenness gives $m_K$ even. Hadamard factorization and pairing
opposite zeros give
\[
 F_K(z)=c z^{m_K}\prod_{\{\alpha,-\alpha\}}
                      (1-z^2/\alpha^2),\qquad c\ne0.
\]
The product converges normally because the order-one zero
divisor satisfies $\sum_{\alpha\ne0}|\alpha|^{-2}<\infty$.
Pairing the genus-one factors cancels their exponential parts;
any remaining factor $e^{az+b}$ has $a=0$ by evenness.
The point $\delta$ is not a zero, since
$\sigma>1$. Logarithmic differentiation on a neighbourhood of
that point is consequently legitimate and gives
\[
 G_K(z)=\frac{m_K}{z}
       +\sum_{\{\alpha,-\alpha\}}
                        \frac{2z}{z^2-\alpha^2}.
\]
Differentiating once more, the contribution of a noncentral
pair to $(\delta G_K-\delta^2G_K')/2$ is
$2\delta^4/(\delta^2-\alpha^2)^2$, and the central contribution
is $m_K$. This also proves absolute convergence of the series
in \eqref{eq:analytic-T2-order}, whose tail is
$O_\delta(\sum |\alpha|^{-4})$.

It remains to check positivity without moving zeros to the
critical line. Write $\alpha=a+ib$, where $|a|\le1/2$. Put
\[
 A=\delta^2-a^2+b^2,\qquad B=2ab.
\]
Then
\[
 A-|B|=\delta^2-2a^2+(|b|-|a|)^2
                    >0.
\]
For a genuine quartet $\alpha,-\alpha,\bar\alpha,-\bar\alpha$,
the two paired contributions therefore have positive real sum,
because
\[
 \Re\frac1{(\delta^2-\alpha^2)^2}
             =\frac{A^2-B^2}{(A^2+B^2)^2}>0.
\]
If the pair is real or purely imaginary, its single contribution
is itself positive. This proves the first inequality with all
possible off-line zeros included.

For real $\sigma>1$, the ideal Euler product gives the absolutely
convergent positive sums
\[
 E_1(\sigma)=\sum_{\mathfrak p}\sum_{k\ge1}
           \frac{\log N\mathfrak p}{(N\mathfrak p)^{k\sigma}}>0,
 \qquad
 E_2(\sigma)=\sum_{\mathfrak p}\sum_{k\ge1}
           \frac{k(\log N\mathfrak p)^2}
                {(N\mathfrak p)^{k\sigma}}>0.
\]
Thus $\zeta_K'/\zeta_K=-E_1$ and
$(\zeta_K'/\zeta_K)'=E_2$. Taking the logarithmic derivatives of
\eqref{eq:analytic-order-completion} now yields exactly
\[
 T_{2,K}(\delta)=U_K(\delta)
                  -\frac\delta2E_1(\sigma)
                  -\frac{\delta^2}2E_2(\sigma)
                  <U_K(\delta),
\]
which proves \eqref{eq:analytic-T2-arch}.
\end{proof}

\begin{corollary}[The quaternionic example has central order exactly two]
\label{cor:analytic-martinet-exact-two}
Let
\[
 F=\mathbb Q(\sqrt5,\sqrt{41}),\qquad
 N=F\left(\sqrt{\frac{(5+\sqrt5)(41+\sqrt{205})}{4}}\right).
\]
For this totally real tame quaternionic field of discriminant
$205^6$ from NN-A018,
\[
               \operatorname{ord}_{s=1/2}\zeta_N(s)=2.
\]
Its unique irreducible degree-two Artin factor has central order
exactly one. These statements require neither A nor M.

For the specific shifted ideal-norm series
$\zeta_{N,t}(s)=\sum_{n\ge1}a_N(n)(n+t)^{-s}$, continued as in
NN-A004, there are exactly two zero branches near $1/2$. Both
are strictly to the right for all sufficiently small real
$t>0$. More precisely, put
\[
 a_2=\frac{\zeta_N''(1/2)}2<0,\qquad
 A_N=\left(\frac{\zeta_N(3/2)}{2|a_2|}\right)^{1/2},\qquad
 d_N=-\frac{C_N A_N^2}{2}>0,
\]
where $C_N$ is the coefficient of NN-A018. Then
\begin{equation}\label{eq:analytic-martinet-two-right}
 s_\pm(t)=\frac12\pm i A_N\sqrt t+d_Nt+O(t^{3/2}),
 \qquad
 \Re s_\pm(t)-\frac12=d_Nt+O(t^2).
\end{equation}
In particular, their real velocities are
$\frac{d}{dt}\Re s_\pm(t)=d_N+O(t)>0$ on a sufficiently small
positive interval.
\end{corollary}

\begin{proof}
The field, discriminant and root-number input are those already
verified in NN-A018: Martinet, 1971, Propositions~III.1--III.2,
Theorem~IV.1 and Example~3, pp.~404--407, together with
Fr\"ohlich, 1974, Theorem~1, p.~82. The three quadratic central
values for discriminants $5,41,205$ are strictly positive by
the saved interval certificate of NN-A018. Hence the
factorization
\[
 \zeta_N(s)=\zeta(s)L(s,\chi_5)L(s,\chi_{41})
                         L(s,\chi_{205})L(s,\psi)^2
\]
and the negative symplectic root number imply
$m_N=2\operatorname{ord}_{1/2}L(s,\psi)\equiv2\pmod4$.
In particular $m_N\ge2$ and $a_{m_N}<0$; none of these steps
uses a Riemann hypothesis.

Apply the theorem at $\delta=1$. Since $r_1=8$, $r_2=0$ and
$\log D_N=6\log205$,
\[
 B_N(3/2)=\frac83+3\log205-4\log\pi+4\psi(3/4),
 \qquad
 B_N'(3/2)=-\frac{40}9+2\psi'(3/4).
\]
The reflection and duplication formulas for $\psi$, with
$\psi(1)=-\gamma_{\!E}$, give
\[
 \psi(3/4)=-\gamma_{\!E}+\frac\pi2-3\log2.
\]
For clarity, the other special value follows directly from
DLMF~5.15.1:
\[
 \psi'(3/4)=16\sum_{n\ge0}\frac1{(4n+3)^2}
                  =\pi^2-8\mathcal G,
 \qquad
 \mathcal G=\sum_{n\ge0}\frac{(-1)^n}{(2n+1)^2}.
\]
Indeed the sum of the two odd residue-class sums is
$\pi^2/8$ and their difference is $\mathcal G$.
Consequently
\begin{equation}\label{eq:analytic-martinet-U-exact}
 U_N(1)=\frac{32}9+\frac32\log205-2\log\pi
       -2\gamma_{\!E}+\pi-6\log2-\pi^2+8\mathcal G.
\end{equation}

Only coarse rational estimates are needed:
\begin{equation}\label{eq:analytic-martinet-coarse}
 \frac{157}{50}<\pi<\frac{22}7,\quad
 \log205<\frac{16}3,\quad \log\pi>\frac{57}{50},\quad
 \gamma_{\!E}>\frac12,\quad \log2>\frac{69}{100},\quad
 \mathcal G<\frac{23}{25}.
\end{equation}
Here are elementary finite checks, included so that no
floating-point constant is required. The Machin identity
$\pi=16\arctan(1/5)-4\arctan(1/239)$ follows by the tangent
addition law and the angle ranges. Alternating-series bounds
give
\[
 \begin{split}
 16\left(\frac15-\frac1{3\cdot5^3}\right)-\frac4{239}
                         &>\frac{157}{50},\\
 16\left(\frac15-\frac1{3\cdot5^3}
                         +\frac1{5\cdot5^5}\right)
 -4\left(\frac1{239}-\frac1{3\cdot239^3}\right)
                         &<\frac{22}7.
 \end{split}
\]
The positive series $\log x=2\sum_{j\ge0}y^{2j+1}/(2j+1)$,
$y=(x-1)/(x+1)$ for $x>1$, gives
\[
 \log\pi>\log(157/50)
       >2\sum_{j=0}^3\frac{(107/207)^{2j+1}}{2j+1}
       >\frac{57}{50},\qquad
 \log2>2\left(\frac13+\frac1{81}\right)
       =\frac{56}{81}>\frac{69}{100}.
\]
For the two upper logarithm bounds used here, the positive
exponential series gives the exact rational comparisons
\[
 \begin{split}
 \sum_{j=0}^{16}\frac{(16/3)^j}{j!}-205
       &=\frac{58220464777647308}{27485885585032875}>0,\\
 \sum_{j=0}^{8}\frac{(39/20)^j}{j!}-7
       &=\frac{3130848254169}{114688000000000}>0.
 \end{split}
\]
The first proves $\log205<16/3$. The elementary integral
comparison $\gamma_{\!E}>H_6-\log7$, the second comparison and
$H_6=49/20$ give $\gamma_{\!E}>1/2$. Finally, the alternating
Catalan series gives
\[
 \mathcal G<\sum_{j=0}^6\frac{(-1)^j}{(2j+1)^2}
       =\frac{23}{25}-\frac{3085982}{2029052025}
       <\frac{23}{25}.
\]
This proves every inequality in
\eqref{eq:analytic-martinet-coarse}. Substitution in
\eqref{eq:analytic-martinet-U-exact} now gives
\[
 \begin{split}
 U_N(1)&<\frac{32}9+8-\frac{57}{25}-1+\frac{22}7
            -\frac{207}{50}-\frac{24649}{2500}
            +\frac{184}{25}\\
       &=\frac{752663}{157500}
         =5-\frac{34837}{157500}<5.
 \end{split}
\]
The theorem yields $m_N<5$. Together with
$m_N\equiv2\pmod4$, this forces $m_N=2$ and the Artin order
one.

NN-A018 already proves $C_N<-3.5197<0$ using its saved
certificate. Its full local formula, now with $m=2$, gives
the first expansion in \eqref{eq:analytic-martinet-two-right}.
There are exactly two holomorphic sections after $t=u^2$.
Real coefficients and uniqueness identify these sections both
by conjugation on real $u$ and by $u\mapsto-u$; thus the odd
coefficients of either section are imaginary and its even
coefficients are real. The real part consequently has a
convergent power series in $t$, with first coefficient
$d_N>0$. This proves the improved $O(t^2)$ remainder, permits
termwise differentiation, and establishes the strict side and
velocity assertions.
\end{proof}

\paragraph{Classification and source boundary.}
The quaternionic root-number and central-existence results are
classical. Sami Omar, \emph{On Artin $L$-Functions for Octic
Quaternion Fields}, \emph{Experimental Mathematics}~10 (2001),
237--245, DOI~10.1080/10586458.2001.10504445, studies exact
central orders in these fields. The present proof does not
claim to be the first determination of the order of this field,
nor does it import a table entry that has not been inspected.
Its order bound is fully proved above and uses the existing
NN-A018 certificate only for the quadratic nonvanishing and
the second-jet sign. The comparison with an exact Omar table
entry remains a source-location audit, not a missing hypothesis
of the deduction.

\paragraph{AG and representation translation.}
The symplectic irreducible representation of the actual
$Q_8$-extension has negative functional-equation sign and,
by the proved upper bound, a length-one central analytic zero.
Its multiplicity two in the regular representation produces
the local Dedekind Artin algebra $\mathbb C[z]/(z^2)$. For the
specified shifted deformation, the base change $t=u^2$ gives
two sections with a common positive second real normal
coefficient $d_N$. Their discriminant is
$-4A_N^2t+O(t^2)$, so the cover has simple ramification at
$t=0$, is etale off that point on a sufficiently small disc,
and has the transposition monodromy. This is the concrete
finite-free discriminant criterion of SGA~1, Expos\'e~I,
section~4, Proposition~10. Positivity of the rational test is
a statement about conjugate analytic zero orbits, not
Frobenius positivity or a claimed motive for the shifted
series. The sign of the second normal jet retains the analytic
unit data that the length and the reduced support alone do not
record.


\section{An explicit cubic compositum with two excess left branches}
\label{sec:analytic-cubic-left-imbalance}

Let $N$ be the quaternionic field of NN-A018 and NN-A033.
Put $p=751$, and let $\vartheta$ be the primitive cubic
Dirichlet character modulo $p$ specified by
\[
  \vartheta(3^j)=e^{2\pi ij/3},\qquad
  \vartheta(n)=0\quad(p\mid n).
\]
Let $E$ be the fixed field of $\ker\vartheta$ in
$\mathbb Q(\zeta_{751})$, and define $K=NE$.
This explicitly specifies the field; no search over number
fields is part of the statement.

\begin{theorem}[NN-A034: an unconditional arithmetic left imbalance]
\label{thm:analytic-cubic-left-imbalance}
The extension $K/\mathbb Q$ is totally real, has Galois group
$Q_8\times C_3$, degree $24$, and absolute discriminant
\[
                   D_K=205^{18}751^{16}.
\]
There is an integer $k\ge0$ such that
\[
 \operatorname{ord}_{s=1/2}\zeta_K(s)=m=2+4k,
 \qquad \frac{\zeta_K^{(m)}(1/2)}{m!}<0,
 \qquad C_K>0.
\]
For all sufficiently small positive $t$, the shifted ideal-norm
deformation of NN-A004 has exactly $2+2k$ central branches on
the left of $\Re s=1/2$ and exactly $2k$ on the right. Thus it
has exactly two excess left branches. Its two initially vertical
branches both bend strictly left at the second Puiseux order.
All these conclusions are unconditional; neither A nor M is
used. The integer $k$, a central order of a complex Artin
twist, is not determined by this theorem.
\end{theorem}

\begin{lemma}[Four certified nonzero central values]
\label{lem:analytic-cubic-four-values}
Write $\chi_d(n)=(d/n)$ for $d=1,5,41,205$, with
$\chi_1=1$. Each $\eta_d=\chi_d\vartheta$ is primitive and
even, of conductor $q_d=751d$. Its central value lies in the
following open rectangle. Every decimal endpoint in this table
is an exact rational number.
\[
\begin{array}{c|r|c|c}
d&q_d&\Re L(1/2,\eta_d)&\Im L(1/2,\eta_d)\\\hline
1&751&(-0.0265980,-0.0265979)&(0.2188449,0.2188450)\\
5&3755&(1.2071682,1.2071683)&(-0.5141633,-0.5141632)\\
41&30791&(0.3418888,0.3418889)&(0.0415524,0.0415525)\\
205&153955&(0.8806422,0.8806423)&(2.0675988,2.0675989)
\end{array}
\]
In particular all four values are nonzero.
\end{lemma}

\begin{proof}
Primality of $751$ follows by trial division by the primes at
most $\sqrt{751}<28$. Since $750=2\cdot3\cdot5^3$ and
\[
 3^{375}\equiv750,\qquad3^{250}\equiv678,\qquad
 3^{150}\equiv80\pmod{751},
\]
the residue class of $3$ is primitive. This defines the character
exactly by its finite discrete-log table. Its parity is even
because $\vartheta(-1)=e^{2\pi i375/3}=1$. The quadratic
characters have positive fundamental discriminants $d$, and
their conductors are coprime to $751$; the product characters
are consequently primitive of conductor $q_d$ and even.

Here is the exact analytic formula and rigorous tail used for
the finite check. For any primitive even character $\eta$ of
conductor $q$, put
\[
 \epsilon_\eta=\frac1{\sqrt q}
          \sum_{a=1}^{q-1}\eta(a)e^{2\pi ia/q},\qquad
 w_q(n)=\frac{\Gamma(1/4,\pi n^2/q)}{\Gamma(1/4)\sqrt n}.
\]
The primitive Gauss sum has modulus $\sqrt q$, so
$|\epsilon_\eta|=1$. The completed functional equation has
this sign in the convention
$\Lambda(s,\eta)=(q/\pi)^{s/2}\Gamma(s/2)L(s,\eta)$;
compare DLMF~25.15.5--25.15.6. To derive the central formula,
apply Poisson summation in the residue classes to
$\Theta_\eta(x)=\sum_{n\in\mathbb Z}\eta(n)e^{-\pi n^2x/q}$.
The finite Fourier transform is
\[
 \sum_{a\bmod q}\eta(a)e^{2\pi ian/q}
              =\bar\eta(n)\tau(\eta),
\]
so
$\Theta_\eta(x)=\epsilon_\eta x^{-1/2}
\Theta_{\bar\eta}(1/x)$. Splitting its Mellin integral at $1$
gives the absolutely convergent identity
\begin{equation}\label{eq:analytic-cubic-central-sum}
 L(1/2,\eta)=\sum_{n\ge1}
            \bigl(\eta(n)+\epsilon_\eta\bar\eta(n)\bigr)w_q(n).
\end{equation}
The change of variable $u=\pi n^2x/q$ evaluates each integral
as the upper incomplete gamma function displayed in $w_q$;
dividing by $(q/\pi)^{1/4}\Gamma(1/4)$ leaves precisely
$n^{-1/2}$. This verifies every scale and phase in
\eqref{eq:analytic-cubic-central-sum}.

For $x>0$,
\[
 \Gamma(1/4,x)=\int_x^\infty u^{-3/4}e^{-u}\,du
                  \le x^{-3/4}e^{-x}.
\]
Also
$\Gamma(1/4)>\int_0^1u^{-3/4}(1-u)\,du=16/5>1$.
Consequently the omitted tail after an integer $B\ge1$ is
bounded by
\begin{equation}\label{eq:analytic-cubic-tail}
 \begin{split}
 \left|\sum_{n>B}
       (\eta(n)+\epsilon_\eta\bar\eta(n))w_q(n)\right|
 &\le 2(q/\pi)^{3/4}
       \int_B^\infty x^{-2}e^{-\pi x^2/q}\,dx\\
 &\le\frac{(q/\pi)^{7/4}}{B^3}
                         e^{-\pi B^2/q}.
 \end{split}
\end{equation}
The first comparison uses the decreasing positive summand;
the second uses $x^{-2}\le x/B^3$ for $x\ge B$ and integrates
the resulting Gaussian derivative exactly.

Only the Gauss sum of conductor $751$ is needed. The Chinese remainder
theorem, writing a residue modulo $751d$ as $751a+db$, gives
\[
 \tau(\chi_d\vartheta)
   =\chi_d(751)\vartheta(d)\tau(\chi_d)\tau(\vartheta),
 \qquad
 \epsilon_{\eta_d}
   =\chi_d(751)\vartheta(d)\epsilon_\vartheta.
\]
Here $\tau(\chi_d)=\sqrt d$, with the positive quadratic Gauss
sign verified in Gross, section~1, equations~(1.3)--(1.6),
as used in NN-A032; for $d=1$ this is the trivial identity.
This exact formula avoids any ambiguity about conjugating a
root number.

The saved script \path{work/analytic/check_cubic_compositum.py}
uses these formulas with FLINT/Arb complex balls at $60$
decimal digits. It takes
\[
 B=5(\lfloor\sqrt q\rfloor+1),\qquad
 B\in\{140,310,880,1965\},
\]
computes the full finite sum in
\eqref{eq:analytic-cubic-central-sum}, and adds the upper bound
\eqref{eq:analytic-cubic-tail} to the error in both real and
imaginary components. The respective tail bounds are strictly
less than
\[
 1.31\cdot10^{-38},\quad9.86\cdot10^{-38},\quad
 6.87\cdot10^{-37},\quad1.29\cdot10^{-36}.
\]
It checks all eight real and imaginary interval inclusions in
the table, as well as nonvanishing, and writes the complete
balls and tail bounds to
\path{work/analytic/cubic_compositum_certificate.json}.
The stated rectangles are certified enclosures, not numerical
approximations substituted for a proof.
\end{proof}

\begin{proof}[Proof of Theorem~\ref{thm:analytic-cubic-left-imbalance}]
The cyclic field $E$ exists by the cyclotomic Galois
correspondence. It has degree three and is real because
$\vartheta(-1)=1$. The only ramified prime is $751$, where the
ramification is total and tame of index three, so
$D_E=751^2$. This also follows directly from the
conductor--discriminant formula for its two nontrivial cubic
characters. The degrees $8$ and $3$ are coprime, so
$N\cap E=\mathbb Q$. Both fields are Galois, which gives
$\operatorname{Gal}(K/\mathbb Q)=Q_8\times C_3$ and degree
$24$. The discriminants $205^6$ and $751^2$ are coprime.
The integral tensor-product discriminant formula therefore
gives
\[
 D_K=D_N^3D_E^8=205^{18}751^{16}.
\]
A precise primary exposition is B.~Conrad,
\emph{Discriminant of composite fields}, Math.~154 handout,
p.~1: its displayed formula and trace-tensor determinant
calculation apply with exactly these degree and coprimality
hypotheses. Equivalently, the tame inertia indices $4,4,3$ at
$5,41,751$ give exponents $24(1-1/4)=18$ and
$24(1-1/3)=16$ directly.

Let $\psi$ denote the unique degree-two irreducible character
of $Q_8$, inflated where appropriate. The irreducible
characters of the product group are
\[
 \chi_d\vartheta^j\quad(d=1,5,41,205;\ j=0,1,2),
 \qquad
 \psi\vartheta^j\quad(j=0,1,2).
\]
The former have degree one and the latter degree two.
Decomposing the regular representation, Artin formalism gives
the exact factorization, including ramified Euler factors,
\begin{equation}\label{eq:analytic-cubic-factorization}
 \zeta_K(s)=
  \prod_{d\in\{1,5,41,205\}}\prod_{j=0}^2
                L(s,\chi_d\vartheta^j)
  \prod_{j=0}^2 L(s,\psi\vartheta^j)^2.
\end{equation}
For clarity, the regular-representation Euler identity does
not discard bad primes: on inertia invariants, Frobenius acts
by cycles corresponding to the prime ideals over the given
rational prime, with their residue degrees. Its determinant
is therefore the product of the corresponding Dedekind Euler
factors. Direct-sum decomposition of that same invariant
space gives \eqref{eq:analytic-cubic-factorization} at every
prime.

There is no unproved Artin holomorphy input. Choose $C_4\subset
Q_8$ and a faithful linear character $\lambda$ of $C_4$.
Since its nontrivial conjugate is $\lambda^{-1}$,
$\operatorname{Ind}_{C_4}^{Q_8}\lambda=\psi$. Thus
\[
 \psi\vartheta^j=
 \operatorname{Ind}_{C_4\times C_3}^{Q_8\times C_3}
                       (\lambda\boxtimes\vartheta^j).
\]
Artin induction invariance identifies each $L$-function with
the nontrivial finite-order Hecke $L$-function over the
quadratic fixed field of $C_4\times C_3$. Such a function is
entire: the classical analytic input is, for example,
B.~Poonen, \emph{Tate's thesis}, section~5.11,
Theorem~5.22(b), p.~38, together with the Galois-character
construction in Example~5.14, p.~32. A nontrivial finite-order
character cannot be a norm power, the only exceptional
component in that theorem. The local induction identity is
the determinant identity for induced cosets, and holds at the
ramified places as well; it is the same unconditional
monomial argument already used for $\psi$ in NN-A018.

NN-A033 proves
$\operatorname{ord}_{1/2}L(s,\psi)=1$. Put
$k=\operatorname{ord}_{1/2}L(s,\psi\vartheta)\ge0$.
Conjugating the Artin Euler factors gives
\[
 L(s,\psi\bar\vartheta)
       =\overline{L(\bar s,\psi\vartheta)},
\]
because $\psi$ is real valued. This identity extends from
$\Re s>1$ to the entire continuations and shows that the
two cubic Artin twists have the same order $k$ at the center.
The four nonzero values in Lemma~\ref{lem:analytic-cubic-four-values}
and their conjugates dispose of the eight nonreal
degree-one factors. The remaining factors are $\zeta$ and
the three positive quadratic central values already certified
in NN-A018; $\zeta(1/2)<0$. Hence
\eqref{eq:analytic-cubic-factorization} gives exactly
\[
                  m=2+4k.
\]
It also determines the leading sign. On the real axis the
nonreal character factors occur as absolute squares. If
$b_1$ and $c_k$ are the first nonzero coefficients of
$L(s,\psi)$ and $L(s,\psi\vartheta)$, respectively, their
contribution to the leading coefficient is
$b_1^2|c_k|^4>0$. Multiplying by the negative nonzero product
of the degree-one central values gives
$a_m=\zeta_K^{(m)}(1/2)/m!<0$.

Finally the root discriminant is
\[
 D_K^{1/24}=205^{3/4}751^{2/3}>4371.
\]
This comparison is an exact integer inequality: raising to
the twelfth power gives
\[
 205^9 751^8-4371^{12}
    =16062670137855473202693706441001475807620284>0.
\]
The conductor criterion proved and interval-certified in
NN-A018 therefore gives $C_K>0$. Applying its full second
Puiseux formula with $a_m<0$ and $m=2+4k$ yields
\[
 \#\{\Re s_j(t)<1/2\}=\frac{m+2}2=2+2k,
 \quad
 \#\{\Re s_j(t)>1/2\}=\frac{m-2}2=2k
\]
on a common sufficiently small interval $t>0$. The two
vertical coefficients are $\pm iA$, and their second normal
coefficient is $-C_KA^2/m<0$. This proves all the assertions.
\end{proof}

\paragraph{Scope and arithmetic consequence.}
Together NN-A033 and NN-A034 supply explicit number fields
with opposite central side imbalances under the same
ideal-norm shift. The latter field is fixed before the
finite certification and requires no nonabelian GRH. The
unknown $k$ affects the total number of branches but not the
excess of two on the left. No noncentral sign-density law,
Riemann-zeta simplicity theorem or heat-collision statistic
is inferred from these central examples.

\paragraph{AG and representation translation.}
Away from $5\cdot41\cdot751$, the specified compositum gives
an actual finite etale $Q_8\times C_3$-cover of
$\operatorname{Spec}\mathbb Z[1/(5\cdot41\cdot751)]$.
The external tensor-product decomposition of its regular
representation gives the displayed Artin factorization.
Complex conjugation pairs the two cubic degree-two twists,
so their contribution to the central analytic length is
$4k$, while the old symplectic factor contributes two.
After $t=u^{2+4k}$, the local analytic Cartier divisor has
$2+4k$ normalized sections, with the two vertical first jets
having strictly negative second real normal component.
These are actual representation multiplicities and local
Artin lengths. The exterior-power/representation-ring
formalism is SGA~6, Expos\'e~V, Example~2.2(b),
Lemma~2.2.1; the local finite-cover/discriminant dictionary
is SGA~1, Expos\'e~I, section~4, Proposition~10. The
specified analytic norm shift is not asserted to arise
from an arithmetic deformation of the cover.


\section{Multiple arithmetic zeros cannot remain tangent through the second Puiseux jet}
\label{sec:analytic-multiple-zero-bending}

This section treats zeros of multiplicity $m\ge2$. The simple
case has an additional deformation term, displayed at the end;
no conclusion about simple-zero transversality is inferred.
All statements at a specified critical-line zero are
unconditional. Assumption A makes them apply to every multiple
nontrivial zero of $\zeta$ or a primitive Dirichlet $L$-function.

\begin{lemma}[A positive logarithmic derivative at one horizontal unit]
\label{lem:analytic-reflected-logderivative}
Let $\Lambda$ be an entire order-one completion with zeros in
$0\le\Re s\le1$, infinitely many zeros, and reflection
\[
 \Lambda(s)=\epsilon\overline{\Lambda(1-\bar s)},
 \qquad |\epsilon|=1.
\]
Suppose $\rho=1/2+i\gamma$ is a zero of order $m\ge1$.
Then, unconditionally,
\begin{equation}\label{eq:analytic-reflected-positive}
 \Re\frac{\Lambda'}{\Lambda}(\rho+1)
   =m+\sum_{\rho'\ne\rho}
       \frac{3/2-\Re\rho'}
       {(3/2-\Re\rho')^2+(\Im\rho'-\gamma)^2}>m.
\end{equation}
Multiplicity is retained and the displayed real series is
absolutely convergent.
\end{lemma}

\begin{proof}
Put $F(z)=\Lambda(\rho+z)$ and write its first two coefficients
as $b_mz^m+b_{m+1}z^{m+1}+O(z^{m+2})$. Since
$1-\bar\rho=\rho$, reflection becomes
$F(z)=\epsilon\overline{F(-\bar z)}$. Therefore
\[
 b_m=\epsilon(-1)^m\bar b_m,
 \qquad b_{m+1}=\epsilon(-1)^{m+1}\bar b_{m+1},
 \qquad
 \Re\frac{b_{m+1}}{b_m}=0.
\]
The genus-one Hadamard product has the form
\[
 F(z)=z^m e^{az+b}\prod_{\alpha\ne0}
           E_1(z/\alpha),\qquad
 E_1(w)=(1-w)e^w,
 \quad \alpha=\rho'-\rho.
\]
All multiplicities are included and
$\sum |\alpha|^{-2}<\infty$. Since each $E_1$ has zero
linear coefficient at the origin, the logarithmic derivative
of $F(z)/z^m$ at zero is $a=b_{m+1}/b_m$. Thus
$\Re a=0$; this is the required control of the exponential
factor, not an assumption about its sign.

At $z=1$, which is not a zero, logarithmic differentiation gives
\[
 \frac{F'(1)}{F(1)}=m+a+
     \sum_{\alpha\ne0}\left(\frac1{1-\alpha}
                                      +\frac1\alpha\right).
\]
The combined complex terms converge absolutely. Moreover
$|\Re\alpha|\le1/2$, so
$\sum|\Re(1/\alpha)|<\infty$. Reflection pairs $\alpha$ with
$-\bar\alpha$, with the same multiplicity, and
\[
 \Re\left(\frac1\alpha+\frac1{-\bar\alpha}\right)=0.
\]
For a purely imaginary zero this cancellation holds termwise.
Taking real parts therefore leaves just
$m+\sum\Re(1/(1-\alpha))$. Each summand is
\[
 \frac{1-\Re\alpha}
 {(1-\Re\alpha)^2+(\Im\alpha)^2}>0,
\]
because $1-\Re\alpha\ge1/2$. Its tail is bounded by a fixed
multiple of $\sum|\alpha|^{-2}$, proving absolute convergence.
There are noncentral zeros, so the inequality is strict.
Substituting $\alpha=\rho'-\rho$ proves
\eqref{eq:analytic-reflected-positive}.
\end{proof}

\begin{theorem}[NN-A035: strict left bending at multiple Dirichlet zeros]
\label{thm:analytic-multiple-dirichlet-bending}
Let $L(s)=L(s,\chi)$ for a nonprincipal primitive Dirichlet
character, real or complex, of parity $a\in\{0,1\}$ and
conductor $q$. Suppose $\rho=1/2+i\gamma$ is a zero of
multiplicity $m\ge2$. Write
\[
 L(\rho+z)=A_mz^m+A_{m+1}z^{m+1}+O(z^{m+2}),
 \qquad
 D_\rho=\frac1\rho+\frac{L'}L(\rho+1)
                              -\frac{A_{m+1}}{A_m}.
\]
Then
\begin{equation}\label{eq:analytic-multiple-D-positive}
 \Re D_\rho>m+h_a(\gamma)>m,
\end{equation}
where the exact positive functions are
\[
 h_1(\gamma)=\frac{\pi}{2\cosh(\pi\gamma)},\qquad
 h_0(\gamma)=\frac1{2(\gamma^2+1/4)}
                      -\frac{\pi}{2\cosh(\pi\gamma)}>0.
\]
For the shifted deformation of NN-A004, the $m$ normalized
branches after $t=u^m$ satisfy
\begin{equation}\label{eq:analytic-multiple-full-jets}
 s_j(t)=\rho+c_jt^{1/m}
                  +\frac{D_\rho}{m}c_j^2t^{2/m}
                  +O(t^{3/m}),\qquad
 c_j^m=\frac{\rho L(\rho+1)}{A_m}\ne0.
\end{equation}
Every branch has a nonzero real normal coefficient by its
second Puiseux order. Specifically, if $\Re c_j=0$ and
$c_j=i\beta_j$, then
\[
 \Re s_j(t)-\frac12
       =-\frac{\beta_j^2}{m}\Re D_\rho\,t^{2/m}
                         +O(t^{3/m})<0
\]
for all sufficiently small positive $t$.
If vertical first directions occur, the exact left/right
counts are $(m/2+1,m/2-1)$ for even $m$, and
$((m+1)/2,(m-1)/2)$ for odd $m$.
\end{theorem}

\begin{proof}
Use the exact completion
\[
 Q_\chi(s)=\left(\frac q\pi\right)^{(s+a)/2}
                         \Gamma((s+a)/2),
 \qquad \Lambda(s)=Q_\chi(s)L(s,\chi).
\]
The primitive functional equation and conjugation give
$\Lambda(s)=\epsilon_\chi
\overline{\Lambda(1-\bar s)}$ even when $\chi$ is complex;
the root number has modulus one. These are DLMF~25.15.5--25.15.6.
The theta integral as in NN-A032, now with the conjugate theta
series on the reflected half, gives the entire order-one
completion. The classical critical-strip and infinite-zero
inputs are DLMF~25.15(ii); cancellation at the endpoint uses
$L(1,\chi)\ne0$, DLMF~25.15.9, just as in NN-A032.
Thus Lemma~\ref{lem:analytic-reflected-logderivative} applies.

The coefficient ratio in the completed germ is
\[
 \frac{b_{m+1}}{b_m}
          =\frac{A_{m+1}}{A_m}+\frac{Q_\chi'}{Q_\chi}(\rho).
\]
Its real part is zero by reflection. Consequently
\begin{equation}\label{eq:analytic-multiple-D-completed}
 \Re D_\rho=
 \Re\frac{\Lambda'}\Lambda(\rho+1)
 +\Re\left(\frac1\rho+
       \frac{Q_\chi'}{Q_\chi}(\rho)
       -\frac{Q_\chi'}{Q_\chi}(\rho+1)\right).
\end{equation}
The conductor term cancels exactly. It remains to evaluate
the real gamma difference, retaining its height dependence.

Put $z=(\rho+a)/2$. Since $\Re\rho=1/2$, one has
$\overline{(\rho+1)/2}=1-\rho/2$. The digamma reflection
identity $\psi(1-z)-\psi(z)=\pi\cot(\pi z)$, DLMF~5.5.4,
therefore gives
\[
 \Re\left(\frac12\psi(\rho/2)
                   -\frac12\psi((\rho+1)/2)\right)
             =-\frac{\pi}{2\cosh(\pi\gamma)}.
\]
Indeed $\Re\cot(\pi/4+i\pi\gamma/2)=1/\cosh(\pi\gamma)$,
as follows by writing sine and cosine in exponentials.
This yields $h_0$ after adding
$\Re(1/\rho)=1/[2(\gamma^2+1/4)]$.
For odd parity, the recurrence (DLMF~5.5.2)
$\psi(1+\rho/2)=\psi(\rho/2)+2/\rho$ reverses the preceding
difference after the same $1/\rho$ term is added, giving
$h_1=\pi/(2\cosh(\pi\gamma))$.
The positivity of $h_0$ is exact: from $2<\pi<4$,
\[
 \cosh(\pi\gamma)\ge1+\frac{\pi^2\gamma^2}{2}
              >\pi(\gamma^2+1/4).
\]
This proves both positive functions. Combining them with
\eqref{eq:analytic-reflected-positive} proves
\eqref{eq:analytic-multiple-D-positive}.

For completeness, the local deformation coefficient is
derived rather than imported from a simple-zero formula.
The shifted expansion gives
\[
 L_t(\rho+z)=A_mz^m+A_{m+1}z^{m+1}
       -t(B_0+B_1z)+O(z^{m+2})+O(tz^2)+O(t^2),
\]
where $B_0=\rho L(\rho+1)$ and
$B_1=L(\rho+1)+\rho L'(\rho+1)$. Since $\Re(\rho+1)=3/2$,
the Euler product makes $B_0\ne0$. Put $t=u^m$ and
$z=c_ju+d_ju^2+O(u^3)$. The coefficient of $u^m$ gives
$A_mc_j^m=B_0$, with distinct nonzero roots. The coefficient
of $u^{m+1}$ gives
\[
 mA_mc_j^{m-1}d_j+A_{m+1}c_j^{m+1}-B_1c_j=0,
\]
because $2m\ge m+2$ for $m\ge2$. Thus
$d_j=c_j^2D_\rho/m$, proving
\eqref{eq:analytic-multiple-full-jets}. The analytic implicit
function theorem after division by $u^m$ justifies these
convergent sections and their common error estimate.

For a vertical $c_j=i\beta_j$, its second real coefficient
is strictly negative by \eqref{eq:analytic-multiple-D-positive}.
Every other branch is decided by its nonzero first real
coefficient. If $m$ is even, vertical directions come as the
opposite pair; the remaining roots are antipodal and split
equally between left and right. If $m$ is odd, there is exactly
one vertical root, and the other roots pair with opposite
real parts by their angular positions $\pi/2+2\pi j/m$ (or
the corresponding lower vertical starting angle). Adding
the left-bending vertical branches gives the stated counts.
\end{proof}

\begin{corollary}[Dedekind zeros at sufficient height, and every multiple zeta zero under A]
\label{cor:analytic-multiple-dedekind-bending}
Let $K$ have signature $(r_1,r_2)$ and degree $n$. At a
critical-line zero $\rho=1/2+i\gamma$ of $\zeta_K$ with
order $m\ge2$, the same $D_\rho$ and Puiseux expansion obey
\begin{equation}\label{eq:analytic-dedekind-D-lower}
 \Re D_\rho>m-
 \left(\frac{3}{2(\gamma^2+9/4)}
       +\frac{r_1\pi}{2\cosh(\pi\gamma)}
       +\frac{r_2}{2(\gamma^2+1/4)}\right).
\end{equation}
In particular, the sufficient condition
\[
 \gamma^2\ge\frac1m
       \left(\frac{r_2+3}{2}+\frac{r_1}{\pi}\right)
\]
already implies $\Re D_\rho>0$ and strict left bending of
every vertical first direction. If
\[
 \gamma^2\ge\frac{r_2+3}{2}+\frac{r_1}{\pi},
\]
then the stronger bound $\Re D_\rho>m-1>0$ holds. The simpler sufficient bound
$\gamma^2\ge n+2$ also works.

For $K=\mathbb Q$, the existing NN-A014 interval certificate
implies that every critical-line zero has $\gamma^2>8$.
Hence every multiple critical-line zeta zero satisfies the
stronger bound
\[
                    \Re D_\rho>m-\frac{11}{48}>0.
\]
Under A this treats all multiple nontrivial zeta zeros, without
assuming that any such zero exists.
\end{corollary}

\begin{proof}
Apply Lemma~\ref{lem:analytic-reflected-logderivative} to the
entire pole-removed Dedekind completion of NN-A033. The
local coefficient relation and
\eqref{eq:analytic-multiple-D-completed} still hold. The
pole-removal factors contribute
\[
 \Re\left(\frac1\rho+\frac1{\rho-1}
                                  -\frac1{\rho+1}\right)
                  =-\frac{3}{2(\gamma^2+9/4)}.
\]
Each real-place gamma factor contributes
$-\pi/(2\cosh(\pi\gamma))$, as just calculated. Each
complex-place factor contributes
$\Re(\psi(\rho)-\psi(\rho+1))=-\Re(1/\rho)$.
This proves \eqref{eq:analytic-dedekind-D-lower} with the
exact stated completion and no location hypothesis beyond
the specified zero being on the line.

For $\gamma\ne0$, the bracket is strictly less than
\[
 \frac1{\gamma^2}
        \left(\frac{r_2+3}{2}+\frac{r_1}{\pi}\right),
\]
using $\cosh(\pi\gamma)>\pi^2\gamma^2/2$. Comparing this
upper bound first with $m$ and then with $1$ proves both
thresholds. Since $\pi>3$ and $n=r_1+2r_2$,
the numerator is less than $n/3+3/2<n+2$, giving the
simpler sufficient threshold.

For the zeta assertion, reuse the positive theta integral
and exact normalization $H_0(x)=\xi(1/2+ix/2)/8$ of
NN-A006. The full saved Arb ball of NN-A014 is
\[
 A_0=\frac{\xi''(0)}{4\xi(1/2)}
       \in (0.0117395,0.0117396)\subset(0,1/16).
\]
No new computation is required for this wider rational
enclosure. If $c_1$ is the normalized second theta moment,
the positive coefficient sum in NN-A014 gives
$2c_1\le A_0<1/16$. The elementary inequality
$\cos v\ge1-v^2/2$ therefore gives, unconditionally,
\[
 H_0(x)\ge H_0(0)(1-c_1x^2)>0\qquad(x^2\le32).
\]
This is a statement about the integral and does not use a
real-zero product or RH. Thus a critical-line zero has
$x=2\gamma$ with $\gamma^2>8$. In the case $r_1=1,r_2=0$,
the bracket in \eqref{eq:analytic-dedekind-D-lower} is
strictly less than
$((3/2)+(1/3))/8=11/48$. The claimed lower bound and all
second-jet conclusions follow.
\end{proof}

\paragraph{The simple-zero extra term and the current remaining gate.}
For $m=1$, write $v=\rho L(\rho+1)/L'(\rho)$. Direct
coefficient matching instead gives
\begin{equation}\label{eq:analytic-simple-extra-term}
 s(t)=\rho+vt+
 \left(v^2D_\rho-
       \frac{v(\rho+1)L(\rho+2)}{2L(\rho+1)}\right)t^2
       +O(t^3).
\end{equation}
The extra term is the contribution of
$t^2\rho(\rho+1)L(\rho+2)/2$ in the shifted series.
If $v$ is purely imaginary, positivity of $\Re D_\rho$
controls only the first term in the parentheses. The real
part of the second term has not been controlled here.
Thus this theorem resolves the second normal jet at multiple
zeros, and neither proves $\Re v\ne0$ at simple zeros nor
assigns a universal bending direction to a simple tangent.
The subsequent NN-A044 controls the exterior ratio: every downward
simple tangency bends strictly left for primitive Dirichlet
heights at least three and for every positive zeta height.
It gives an exact positive threshold for upward tangencies.
Equality at that threshold, and the existence of any simple
tangency, remain separate questions.

\paragraph{AG translation.}
At a length-$m$ critical-line analytic zero, completed
reflection puts the first reduced logarithmic coefficient
in $i\mathbb R$. The genus-one exponential and reflected
zero pairs make this local Real constraint compatible with
the strictly positive real logarithmic derivative one unit
to the right. After $t=u^m$, the finite analytic Cartier
divisor normalizes into $m$ sections. For $m\ge2$ the
specified shift has no $t^2$ contribution at the second
Puiseux order, so every vertical tangent acquires the
proved nonzero left normal coefficient. The local algebra
is $\mathbb C[z]/(z^m)$; its normalization and ramified
cover have the standard differentials/discriminant
interpretation of SGA~1, Expos\'e~I, section~3,
Proposition~1 and section~4, Proposition~10. This is a
second-jet theorem for the given analytic deformation,
not a purity statement or an arithmetic realization of
the shifted denominators.


\section{The direction of bending at a simple downward tangency}
\label{sec:A044-simple-tangencies}

The extra term in NN-A035's simple-zero formula has a controlled
sign in one vertical direction. The reason is a uniform sector
bound for the exterior Euler ratio, proved here with elementary
constants. No existence of a simple tangency is assumed.

\begin{lemma}[A uniform exterior Euler sector]
\label{lem:A044-euler-sector}
For every Dirichlet character, including the conductor-one
character, and every real $t$, the absolutely convergent Euler
logarithm of
\[
              R_\chi(t)=
                 \frac{L(5/2+it,\chi)}{L(3/2+it,\chi)}
\]
satisfies
\[
                   |\arg R_\chi(t)|<\frac9{10}.
\]
Consequently, for $\gamma\ge0$,
\begin{equation}\label{eq:A044-exterior-positive}
 \Im\{(3/2+i\gamma)R_\chi(\gamma)\}
   >|R_\chi(\gamma)|\,\frac{119\gamma-270}{200}.
\end{equation}
In particular this imaginary part is positive for
$\gamma\ge3$, and also whenever $\gamma^2>8$.
\end{lemma}

\begin{proof}
A prime factorization is supplied by the absolutely convergent
Euler product, DLMF~25.15.2.
A prime with $\chi(p)=0$ contributes one. Otherwise its
Euler factor is
\[
               w=\frac{1-z}{1-z/p},
       \qquad |z|=r=p^{-3/2}.
\]
Its image circle has real center and radius
\[
 c=\frac{1-r^2/p}{1-r^2/p^2},\qquad
 R=\frac{r(1-1/p)}{1-r^2/p^2}.
\]
Indeed $|1-w|=r|1-w/p|$ gives this circle by completing
the square. Since $c>R>0$, the largest absolute argument
of any point on it is $\arcsin(R/c)$. Thus its principal
argument has absolute value at most $\arcsin u_p$, where
\[
 u_p=\frac{p^{-3/2}(1-1/p)}{1-p^{-4}}.
\]
Put $f(x)=x^{-3/2}-x^{-5/2}$. It is positive and decreasing
for $x\ge2$, because $f'(x)=(5-3x)/(2x^{7/2})<0$.
The denominator $1-x^{-4}$ is increasing, so
$u_p\le u_2=4/(15\sqrt2)$. Integrating the derivative of
arcsine on $[0,u_p]$ gives
\[
 \arcsin u_p\le\frac{15u_p}{\sqrt{217}}
                    \le\frac{16}{\sqrt{217}}f(p).
\]
Here $1-u_2^2=217/225$ and
$(1-p^{-4})^{-1}\le16/15$.

Every prime other than two is an odd integer at least three.
The monotonicity of $f$ bounds the corresponding sum by
\[
 \sum_p f(p)
 \le f(2)+f(3)+\frac12\int_3^\infty f(x)\,dx
       =\frac1{4\sqrt2}+\frac{10}{9\sqrt3}.
\]
For each odd integer at least five, use the preceding
interval of length two to justify the integral upper bound.
The elementary rational estimates
$\sqrt{217}>147/10$, $\sqrt2>7/5$, and $\sqrt3>12/7$
therefore give
\[
 \sum_p\arcsin u_p
 <\frac{160}{147}\left(\frac5{28}+\frac{35}{54}\right)
   =\frac{100000}{111132}<\frac9{10}.
\]
The last comparison is the integer inequality
$1000000<1000188$. Absolute Euler convergence permits the
sum of these local arguments and proves the sector bound,
with no choice of a winding number.

Write $R_\chi=|R_\chi|e^{i\varphi}$, where
$|\varphi|<9/10$. The elementary inequalities
$\cos x\ge1-x^2/2$ and $|\sin x|\le|x|$ give
\[
 \cos\varphi>\frac{119}{200},\qquad
 \sin\varphi>-\frac9{10}.
\]
Hence
$\gamma\Re R_\chi+(3/2)\Im R_\chi$
has the bound in \eqref{eq:A044-exterior-positive}.
The ratio is nonzero by its Euler product.
Finally $270/119<3$ and $(270/119)^2<8$, proving both
stated ranges.
\end{proof}

\begin{theorem}[NN-A044: downward simple tangencies bend strictly left]
\label{thm:A044-downward}
Let $L(s)$ be a nonprincipal primitive Dirichlet $L$-function,
possibly for a complex character, and let
$\rho=1/2+i\gamma$ be a simple zero with $\gamma\ge3$.
Use the specified denominator shift of NN-A004 and suppose its
ordinary velocity is purely imaginary:
\[
                   v_\rho=
             \frac{\rho L(\rho+1)}{L'(\rho)}=i\beta,
                  \qquad \beta\in\mathbb R\setminus\{0\}.
\]
If $\beta<0$, then the moving zero has strictly negative second
real normal coefficient. In particular, for some $d_\rho>0$,
\[
          \Re s(\tau)-\frac12=-d_\rho\tau^2+O(\tau^3)<0
\]
for all sufficiently small positive $\tau$.

The same conclusion holds for every simple positive-height
critical-line zero of $\zeta$, with no further height condition.
These local assertions do not require a global zero-location
hypothesis once the specified zero is on the line. A locates
all nontrivial zeros in the indicated families, but does not
assert that any such tangency occurs.

More precisely, in either range define
\[
 \begin{aligned}
 D_\rho&=\frac1\rho+\frac{L'}L(\rho+1)
                                  -\frac{L''(\rho)}{2L'(\rho)},\\
 d_\rho^{\,*}&=\Re D_\rho>0,\\
 P_\rho&=\Im\left\{(\rho+1)\frac{L(\rho+2)}{L(\rho+1)}\right\}>0,
       \qquad \beta_\rho^*=\frac{P_\rho}{2d_\rho^{\,*}}.
 \end{aligned}
\]
For an upward tangency, $\beta>0$, the second real coefficient
is positive if $0<\beta<\beta_\rho^*$, negative if
$\beta>\beta_\rho^*$, and zero if $\beta=\beta_\rho^*$.
The equality case still requires the higher-jet recurrence.
\end{theorem}

\begin{proof}
First justify the positivity of $\Re D_\rho$ at multiplicity
one, rather than assuming that the higher-multiplicity theorem
already states it. The reflected Hadamard-product lemma
of NN-A035 applies to an entire completed zero of any positive
order $m$, including $m=1$. It gives
$\Re(\Lambda'/\Lambda)(\rho+1)>1$ here. The coefficient
ratio in the completed local germ has zero real part, by
completed reflection, so exactly as in
\eqref{eq:analytic-multiple-D-completed},
\[
 \Re D_\rho=\Re\frac{\Lambda'}{\Lambda}(\rho+1)
   +\Re\left(\frac1\rho+
           \frac{Q'}Q(\rho)-\frac{Q'}Q(\rho+1)\right).
\]
There is no use of $m\ge2$ in this identity or in the
Hadamard lemma. For primitive nonprincipal $L$, the gamma
difference is the positive function $h_a(\gamma)$ already
proved in NN-A035. Thus
$\Re D_\rho>1+h_a(\gamma)>0$.
The exact classical completion, endpoint cancellation and gamma
identities have locators DLMF~25.15.5--25.15.6 and 25.15.9,
and 5.5.2/5.5.4, as used there.

For zeta, use the entire completion with the two pole-removal
factors. The same proof gives the NN-A035 Dedekind inequality
with $m=1$, $r_1=1$ and $r_2=0$. Its existing theta-integral
bound and saved NN-A014 certificate imply $\gamma^2>8$ for
every critical-line zero, and hence
$\Re D_\rho>1-11/48=37/48>0$.
The restriction $m\ge2$ in NN-A035 entered only when ignoring
the second parameter derivative in its Puiseux calculation,
not in these inequalities. No new numerical computation is
being used.

For a simple zero retain that additional derivative. Expanding
the shifted series to second order gives
\[
 L_\tau(s)=L(s)-\tau sL(s+1)
                +\frac{\tau^2s(s+1)}2L(s+2)+O(\tau^3)
\]
near $\rho$. Substitution of
$s(\tau)=\rho+v\tau+a_2\tau^2+O(\tau^3)$ gives
$L'(\rho)v=\rho L(\rho+1)$ and
\[
 L'(\rho)a_2+\frac{L''(\rho)}2v^2
 -v\{L(\rho+1)+\rho L'(\rho+1)\}
 +\frac{\rho(\rho+1)}2L(\rho+2)=0.
\]
Therefore
\[
               a_2=v^2D_\rho
                   -\frac{v(\rho+1)L(\rho+2)}{2L(\rho+1)}.
\]
The implicit function theorem supplies the convergent simple
section and its $O(\tau^3)$ remainder. For $v=i\beta$ this
identity becomes
\begin{equation}\label{eq:A044-real-coefficient}
                   \Re a_2
                    =-\beta^2d_\rho^{\,*}+\frac{\beta P_\rho}{2}
                    =\beta d_\rho^{\,*}(\beta_\rho^*-\beta).
\end{equation}
The Euler-sector lemma gives $P_\rho>0$ for the stated
primitive-character height range and for all the zeta heights
just verified. Thus $\beta<0$ makes both terms in the
middle expression strictly negative. For $\beta>0$, the last
expression proves each of the three threshold alternatives.
Strict sign of the quadratic coefficient gives the eventual
side for the positive parameter ray.
\end{proof}

\paragraph{The threshold uses the corrected phase derivative.}
Let $\theta_\chi$ be the root-normalized Hardy phase,
$\alpha_\chi(t)=\Im\log L(3/2+it,\chi)$, and
$\psi_\chi(t)=\theta_\chi(t)+\alpha_\chi(t)-\arctan(1/(2t))$
as in NN-A042. At every simple critical zero in the theorem,
\[
                         d_\rho^{\,*}=\psi_\chi'(\gamma).
\]
Indeed the completed local coefficient constraint gives
$\Re(L''/(2L'))(\rho)=-\Re(Q'/Q)(\rho)=-\theta_\chi'(\gamma)$.
Also $\alpha_\chi'=\Re(L'/L)(\rho+1)$, while
the derivative of $-\arctan(1/(2t))$ is
$\Re(1/\rho)$ at $t=\gamma$. Substitution gives the identity.
Thus the upward threshold is
$P_\rho/(2\psi_\chi'(\gamma))$, with a strictly positive
denominator at these zeros. It does not require a second
derivative value of $L$ once the phase is known.

\paragraph{Geometric and Grothendieck translation.}
At a reduced analytic zero, the specified divisor projects
isomorphically to the parameter disc and defines a section
$s(\tau)=\rho+i\beta\tau+a_2\tau^2+\cdots$. Its first
normal class vanishes. The proved Euler-sector constraint and
the reflected completed germ imply that the second normal
class is strictly negative whenever the tangent lowers the
height, $\beta<0$. Upward tangents have the exact real threshold
in \eqref{eq:A044-real-coefficient}; equality leaves the
second normal class zero and retains the higher formal
neighborhood. These are ordered real analytic normal jets,
not an intrinsic sign on arbitrary complex schemes.
The reduced local finite algebra has vanishing relative
differentials, in the finite-algebra dictionary of SGA~1,
Expos\'e~I, section~3, Proposition~1; the deformation and its
chosen positive real ray supply the additional orientation.
No arithmetic denominator-shift family, purity statement,
ordinary-sign density or existence of a tangent zero follows.

\paragraph{Exact source and certificate dependencies.}
The Euler product, functional equation and endpoint
nonvanishing locators are
\url{https://dlmf.nist.gov/25.15}, equations~25.15.2,
25.15.5--25.15.6 and 25.15.9.
The gamma identities are
\url{https://dlmf.nist.gov/5.5}, equations~5.5.2 and 5.5.4.
The zeta height bound reuses exactly NN-A035's proof from the
saved NN-A014 interval certificate; there is no new
$L$-value or zero computation. The circle, prime-sum sector
bound, rational comparisons and full simple-jet deduction
are all proved above.


\subsection{The same tangency direction for quadratic Dedekind functions}
\label{sec:A044-quadratic}

This is a local refinement of NN-A044, not a new finding.
No global zero-location hypothesis is used in this subsection.

\begin{lemma}[A sharper absolute Euler sector]
\label{lem:A044-sharper-sector}
For every Dirichlet character, with the same continuous Euler
argument as in Lemma~\ref{lem:A044-euler-sector},
\[
 \left|\arg\frac{L(5/2+it,\chi)}{L(3/2+it,\chi)}\right|<\frac34.
\]
For every quadratic number field \(K\),
\[
 \left|\arg\frac{\zeta_K(5/2+it)}{\zeta_K(3/2+it)}\right|<\frac32.
\]
These bounds are uniform in the character, field and real \(t\).
\end{lemma}

\begin{proof}
Keep \(f(x)=x^{-3/2}-x^{-5/2}\) and \(u_p\) from
Lemma~\ref{lem:A044-euler-sector}. The local circle proof there
gave \(\arcsin u_p\le16f(p)/\sqrt{217}\).
Every prime greater than three belongs to one of the two
progressions \(5+6j\) and \(7+6j\), \(j\ge0\).
Monotonicity of \(f\) therefore gives
\[
 \begin{aligned}
 \sum_p f(p)
 &\le f(2)+f(3)+f(5)+f(7)
        +\frac16\left(\int_5^\infty f(x)\,dx
                       +\int_7^\infty f(x)\,dx\right)\\
 &=\frac1{4\sqrt2}+\frac2{9\sqrt3}
          +\frac{106}{225\sqrt5}+\frac{194}{441\sqrt7}.
 \end{aligned}
\]
For each term after the initial term of either progression,
the preceding interval of length six gives its integral bound.
The rational inequalities
\[
 \sqrt{217}>\frac{147}{10},\quad
 \sqrt2>\frac75,\quad \sqrt3>\frac{12}{7},\quad
 \sqrt5>\frac{20}{9},\quad \sqrt7>\frac{21}{8}
\]
are all verified by squaring positive numbers. Consequently
\[
 \begin{aligned}
 \sum_p\arcsin u_p
 &<\frac{160}{147}
       \left(\frac5{28}+\frac7{54}
                     +\frac{53}{250}+\frac{1552}{9261}\right)\\
 &=\frac{8492776}{11344725}<\frac34.
 \end{aligned}
\]
The last comparison is \(33971104<34034175\).
This proves the first assertion by absolute Euler convergence.

For the second, a prime ideal \(\mathfrak p\) above \(p\) has
norm \(p^f\), \(f\ge1\). Its factor in the ratio is the same
circle factor with \(p\) replaced by \(p^f\).
The function
\[
 u(x)=\frac{x^{-3/2}-x^{-5/2}}{1-x^{-4}}
\]
is decreasing for \(x\ge2\), since its numerator is positive
and decreasing and its denominator positive and increasing.
Thus the prime-ideal factor contributes at most
\(\arcsin u_p\) in absolute argument.
There are at most two prime ideals above each rational prime:
the fundamental degree identity
\(\sum_{\mathfrak p\mid p}e_{\mathfrak p}f_{\mathfrak p}=2\)
and positivity of its integer summands give this directly,
including ramified primes. Summing the absolute arguments
over the prime ideals therefore gives a bound strictly below
\(2(3/4)=3/2\). No splitting-density assertion or factorization
of \(\zeta_K\) into Dirichlet functions is used.
\end{proof}

\begin{corollary}[NN-A044: quadratic-field simple tangencies]
\label{cor:A044-quadratic}
Let \(K\) be a quadratic number field, and let
\(\rho=1/2+i\gamma\) be a simple critical-line zero of
\(\zeta_K\) with \(\gamma\ge24\). For its specified
norm-denominator shift, suppose
\[
                  v_\rho
       =\frac{\rho\zeta_K(\rho+1)}{\zeta_K'(\rho)}
       =i\beta,\qquad \beta\in\mathbb R\setminus\{0\}.
\]
Every downward tangency, \(\beta<0\), bends strictly left at
second order. More precisely, with
\[
 \begin{aligned}
 D_\rho&=\frac1\rho+\frac{\zeta_K'}{\zeta_K}(\rho+1)
                            -\frac{\zeta_K''(\rho)}{2\zeta_K'(\rho)},\\
 d_\rho^*&=\Re D_\rho,\qquad
 P_\rho=\Im\left\{(\rho+1)\frac{\zeta_K(\rho+2)}
                                      {\zeta_K(\rho+1)}\right\},
 \end{aligned}
\]
both \(d_\rho^*\) and \(P_\rho\) are positive and
\[
       \Re a_2=\beta d_\rho^*
                   \left(\frac{P_\rho}{2d_\rho^*}-\beta\right).
\]
Thus the same upward threshold alternatives as in NN-A044 hold.
The assertion is unconditional once the specified simple zero
is on the line, and does not assert a tangent zero exists.
\end{corollary}

\begin{proof}
Put \(R=\zeta_K(5/2+i\gamma)/\zeta_K(3/2+i\gamma)\).
The sharper sector lemma gives \(|\arg R|<3/2\).
For \(0\le x\le3/2\), the alternating cosine series gives
\[
 \cos x\ge 1-\frac{x^2}{2}+\frac{x^4}{24}-\frac{x^6}{720}
       \ge\frac{359}{5120}>\frac1{15}.
\]
The displayed polynomial is decreasing on this interval:
its derivative is
\(-x(1-x^2/6+x^4/120)\le0\), since
\(1-x^2/6\ge5/8>0\).
The rational last comparison is \(5385>5120\).
Also \(\sin x\ge-1\) for every real \(x\). It follows that
\[
 P_\rho=\Im\{(3/2+i\gamma)R\}
      >|R|\left(\frac{\gamma}{15}-\frac32\right)>0
               \qquad(\gamma\ge24).
\]

For \(d_\rho^*\), retain both pole-removal factors in the
entire Dedekind completion. The proof of
\eqref{eq:analytic-dedekind-D-lower}, using the reflected
Hadamard lemma, applies at \(m=1\), as explained in NN-A044.
For the quadratic signature \(r_1+2r_2=2\), it yields
\[
 d_\rho^*>
 1-\left(\frac{3}{2(\gamma^2+9/4)}
       +\frac{r_1\pi}{2\cosh(\pi\gamma)}
       +\frac{r_2}{2(\gamma^2+1/4)}\right)>0.
\]
Indeed the bracket is strictly less than
\(\gamma^{-2}((r_2+3)/2+r_1/\pi)\), whose numerator is
less than \(2/3+3/2<4\). At \(\gamma\ge24\) it is
therefore strictly less than one. This uses the full
meromorphic completion via its entire pole-removed version,
not an entire-function assumption on the meromorphic
completion itself.

The second-parameter expansion and simple implicit-function
calculation in NN-A044 apply unchanged to the norm shift.
They give
\[
 a_2=v_\rho^2D_\rho
       -\frac{v_\rho(\rho+1)\zeta_K(\rho+2)}
                                {2\zeta_K(\rho+1)}.
\]
Substitution of \(v_\rho=i\beta\) proves the stated formula.
Both terms in its real part are strictly negative for
\(\beta<0\); the upward alternatives follow by factoring.
\end{proof}

\paragraph{Inputs and geometric scope.}
Only the prime-ideal Euler product and degree identity are
added to the already checked NN-A035 and NN-A044 interfaces.
For the latter, Milne, \emph{Algebraic Number Theory},
version~3.08 (2020), Theorem~3.34, equation~(10),
printed pp.~59--60, gives the exact degree identity;
the norm convention is defined on printed p.~68:
\url{https://www.jmilne.org/math/CourseNotes/ANT.pdf}.
The full proof of the cited degree identity was read directly,
not the complete notes.
The norm shift defines the same reduced analytic divisor
section, and the sign conclusion concerns its second real
normal jet on the chosen positive parameter ray.
The finite relative-differentials dictionary remains
SGA~1, Expos\'e~I, section~3, Proposition~1.
No global Dedekind RH, arithmetic deformation scheme,
prime-splitting distribution or ordinary-sign density is
inferred.


\section{A positive proportion of actual zeta zeros have a definite normal velocity}
\label{sec:analytic-arithmetic-transverse-proportion}

All zero counts in this section use positive ordinates and
count multiplicity unless a simple-zero condition is written
explicitly. Put
\[
 N(T)=\sum_{0<\gamma\le T}1,\qquad
 N_s(T)=\#\{\rho=1/2+i\gamma:0<\gamma\le T,
                                      \ \zeta'(\rho)\ne0\}.
\]
Thus a zero of multiplicity $m$ contributes $m$ to $N(T)$ and
nothing to $N_s(T)$ when $m>1$.

\begin{theorem}[NN-A036: a quantitative arithmetic transversality proportion]
\label{thm:analytic-arithmetic-transverse-proportion}
Assume A. At each simple zero, let
\[
 v_\rho=\frac{\rho\zeta(\rho+1)}{\zeta'(\rho)},\qquad
 r_\rho=\frac{|\Re v_\rho|}{|v_\rho|}\in[0,1],\qquad
 c=\frac1{2\zeta(3)}-\frac13>\frac1{15}.
\]
Then
\begin{equation}\label{eq:analytic-normal-mean}
 \liminf_{T\to\infty}\frac1{N(T)}
                 \sum_{\substack{0<\gamma\le T\\\rho\ {
m simple}}}
                           r_\rho\ \ge c.
\end{equation}
In particular, at least a lower proportion $c$ of all zero
occurrences are simple with $\Re v_\rho\ne0$.

For $0<\eta<3c/2$, let
\[
 G_\eta(T)=\#\{\rho\ {
m simple}:0<\gamma\le T,
                              \ |\Re v_\rho|\ge\eta|v_\rho|\}.
\]
The bounds with their two distinct denominators are
\begin{align}
 \liminf_{T\to\infty}\frac{G_\eta(T)}{N(T)}
       &\ge\frac{c-(2/3)\eta}{1-\eta},
         \label{eq:analytic-total-transverse-density}\\
 \liminf_{T\to\infty}\frac{G_\eta(T)}{N_s(T)}
       &\ge\frac{(3/2)c-\eta}{1-\eta}.
         \label{eq:analytic-simple-transverse-density}
\end{align}
For example, taking $\eta=1/30$ gives the convenient rational
consequences
\[
 \liminf\frac{G_{1/30}(T)}{N(T)}\ge\frac4{87},\qquad
 \liminf\frac{G_{1/30}(T)}{N_s(T)}\ge\frac2{29}.
\]
These are lower bounds for the magnitude of the relative
normal component. They assign neither its sign nor a half
proportion to left and right labels, and do not establish
transversality at every simple zero.
\end{theorem}

\begin{lemma}[The exterior Euler phase averaged at the zeros]
\label{lem:analytic-exterior-phase-average}
Under RH, for each fixed $\sigma>1$,
\begin{equation}\label{eq:analytic-exterior-ratio-mean}
 \lim_{T\to\infty}\frac1{N(T)}
       \sum_{0<\gamma\le T}
       \frac{\zeta(\sigma+i\gamma)}
            {\zeta(\sigma-i\gamma)}
                     =\frac1{\zeta(2\sigma)}.
\end{equation}
This is an unweighted zero average with multiplicity, not a
joint independence assertion involving the Hardy phase.
\end{lemma}

\begin{proof}
For every fixed $x>1$, Landau's formula is
\[
 \sum_{0<\gamma\le T}x^\rho
      =-\frac{T}{2\pi}\Lambda(x)+O_x(\log T),
\]
where $\Lambda(x)=\log p$ when $x=p^k$ with integer $k\ge1$,
and is zero for other real $x>1$. A primary locator is
S.~M.~Gonek, \emph{A Formula of Landau and Mean Values of
$\zeta(s)$}, author-hosted PDF, p.~1, equation~(1);
its uniform extension is Theorem~1, equation~(3).
The fixed-$x$ formula is all that is used here.
Under RH, $x^\rho=\sqrt x\,e^{i\gamma\log x}$ and
$N(T)\sim T\log T/(2\pi)$. Thus
\begin{equation}\label{eq:analytic-fixed-frequency-zero-average}
 \frac1{N(T)}\sum_{0<\gamma\le T}e^{i\gamma\log x}\longrightarrow0
                    \qquad(x>0,\ x\ne1).
\end{equation}
For $0<x<1$ use complex conjugation and $1/x>1$.

For $\sigma>1$, the following double series is absolutely and
uniformly convergent in real $t$:
\[
 \frac{\zeta(\sigma+it)}{\zeta(\sigma-it)}
 =\sum_{n,m\ge1}\frac{\mu(m)}{(nm)^\sigma}
                                  e^{it\log(m/n)}.
\]
Indeed the absolute coefficient sum is at most
$\zeta(\sigma)^2<\infty$. Each normalized zero average of a
term has modulus at most one, so dominated convergence for
this summable family applies to
\eqref{eq:analytic-fixed-frequency-zero-average}.
Only $m=n$ survives, and its sum is
$\sum_{n\ge1}\mu(n)n^{-2\sigma}=1/\zeta(2\sigma)$.
This proves \eqref{eq:analytic-exterior-ratio-mean} without
requiring uniformity of Landau's error in $x$.
\end{proof}

\begin{proof}[Proof of Theorem~\ref{thm:analytic-arithmetic-transverse-proportion}]
Let $\theta(t)$ be the usual Hardy--Riemann--Siegel phase and
use the Euler-product branch
\[
 \phi(t)=\arg\zeta(3/2+it),\qquad
 \delta(t)=\arctan\frac1{2t},\qquad
 \psi(t)=\phi(t)-\delta(t).
\]
The first necessary phase average is already a proved
classical statement, not full Gram equidistribution:
\begin{equation}\label{eq:analytic-theta-first-mode}
 \frac1{N(T)}\sum_{0<\gamma\le T}e^{2i\theta(\gamma)}
                                      \longrightarrow0.
\end{equation}
The exact locator is J.~Arias de Reyna,
\emph{On the distribution (mod 1) of the normalized zeros of
the Riemann Zeta-function}, arXiv:1403.7126v1,
Corollary~3.2, equation~(13), with $\kappa=1$; its proof gives
$O(1/\log T)$ in this case. That result sums
$\rho=1/2+i\alpha$ without assuming RH. Here A makes
$\alpha=\gamma$ real, with precisely the multiplicity
convention used in $N(T)$.

Since $e^{2i\phi(t)}=\zeta(3/2+it)/\zeta(3/2-it)$,
Lemma~\ref{lem:analytic-exterior-phase-average} gives
\[
 \frac1{N(T)}\sum_{0<\gamma\le T}e^{2i\phi(\gamma)}
                                  \longrightarrow\frac1{\zeta(3)}.
\]
The replacement of $\phi$ by $\psi$ has vanishing average
error. To see this with no interchange of divergent sums,
fix $Y>0$: the finitely many zeros below $Y$ contribute
$o_T(1)$, while for $\gamma>Y$ the difference of the two
unit phases is at most $2\delta(Y)$. Let $T\to\infty$ and
then $Y\to\infty$. Taking real parts, the two limits are
\begin{equation}\label{eq:analytic-two-separate-real-means}
 \langle\cos2\theta\rangle_T\longrightarrow0,
 \qquad
 \langle\cos2\psi\rangle_T\longrightarrow\frac1{\zeta(3)}.
\end{equation}
No joint distribution or decorrelation has been used.

Define at every zero occurrence, including multiple zeros,
the bounded phase observable
$f(\gamma)=|\cos(\theta(\gamma)+\psi(\gamma))|$.
The pointwise trigonometric identity gives
\[
 \cos2\theta+\cos2\psi
   =2\cos(\theta+\psi)\cos(\theta-\psi)
   \le2|\cos(\theta+\psi)|=2f.
\]
Consequently
\begin{equation}\label{eq:analytic-all-phase-lower}
 \liminf_{T\to\infty}\langle f\rangle_T
                              \ge\frac1{2\zeta(3)}.
\end{equation}
This deterministic inequality is why the separate first
moments suffice for the present lower bound, although they
do not prove the parity-weighted half-label law.

At a simple zero, $Z(t)=e^{i\theta(t)}\zeta(1/2+it)$ is
real, $Z'(\gamma)\ne0$, and
\[
 \zeta'(\rho)=-i e^{-i\theta(\gamma)}Z'(\gamma),\qquad
 v_\rho=\frac{i\rho e^{i\theta(\gamma)}
                   \zeta(3/2+i\gamma)}{Z'(\gamma)}.
\]
For $\gamma>0$, $\arg(i\rho)=\pi-\delta(\gamma)$.
The real sign of $Z'(\gamma)$ only adds a multiple of $\pi$
to this argument. Therefore the corrected exact velocity
formula of NN-A001 gives
\begin{equation}\label{eq:analytic-phase-is-relative-normal}
                   r_\rho=f(\gamma)
                       \qquad(\rho\text{ simple}).
\end{equation}
The Euler product makes the numerator of $v_\rho$ nonzero,
so the displayed normalization is well defined.

The remaining imported input is precisely Montgomery's
proved simple-proportion bound under RH:
\begin{equation}\label{eq:analytic-montgomery-simple-proportion}
                   N_s(T)\ge(2/3+o(1))N(T).
\end{equation}
The source is \emph{The Pair Correlation of Zeros of the
Zeta Function}, Corollary~2, printed p.~182, equation~(7).
Its proof in section~4, pp.~191--192, explicitly states
that the zero sums count multiplicity; a multiplicity-$m$
zero is entered $m$ times and contributes $m^2$ to the
equal-ordinate pair sum. Thus
\eqref{eq:analytic-montgomery-simple-proportion} has exactly
the denominator $N(T)$ used in both phase averages. This is
the already catalogued classical input NT-Z02, not a new
simplicity claim or a full pair-correlation conjecture.

Because $0\le f\le1$, removing the nonsimple occurrences
from \eqref{eq:analytic-all-phase-lower} loses at most
$(N(T)-N_s(T))/N(T)\le1/3+o(1)$. Together with
\eqref{eq:analytic-phase-is-relative-normal}, this proves
\eqref{eq:analytic-normal-mean}. Since $r_\rho\le1$ and
vanishes at a tangent simple zero, it also proves the lower
proportion $c$ of simple transverse zeros.

For the sharper threshold bounds, put
$p_T=N_s(T)/N(T)$ and $g_T=G_\eta(T)/N(T)$. The pointwise
upper bounds on the three subsets (nonsimple, simple below
threshold, and simple above threshold) give
\[
 \langle f\rangle_T
       \le (1-p_T)+\eta(p_T-g_T)+g_T
       =1-(1-\eta)p_T+(1-\eta)g_T.
\]
Hence
\[
 g_T\ge p_T-\frac{1-\langle f\rangle_T}{1-\eta}.
\]
Taking lower limits using $p_T\ge2/3+o(1)$ proves
\eqref{eq:analytic-total-transverse-density}. Dividing the
same inequality by $p_T$, which is bounded away from zero,
and using
\[
 \frac{1-\langle f\rangle_T}{p_T}
       \le\frac{1-1/(2\zeta(3))+o(1)}{2/3+o(1)}
\]
proves \eqref{eq:analytic-simple-transverse-density}.
There is no hidden conversion from a distinct-zero count.

Finally,
\[
 \zeta(3)=1+\frac18+\sum_{n\ge3}\frac1{n^3}
       <1+\frac18+\int_2^\infty x^{-3}\,dx=\frac54,
\]
so $c>2/5-1/3=1/15$. Inserting $\eta=1/30$ in the two
exact bounds gives respectively a number greater than
$4/87$ and a number greater than $2/29$, proving the stated
rational consequences without a numerical constant
calculation.
\end{proof}

\paragraph{Scope.}
The theorem supplies a quantitative lower density of actual
arithmetic simple zeros at which the normal velocity is a
fixed fraction of the full velocity. It uses A, the single
proved Hardy-phase mode, the fixed-frequency Landau formula,
and the classical $2/3$ simplicity proportion. It does not
use a normalized CLT, a joint phase law, full GUE, or an
assumption that all zeros are simple. The positive and
negative signs may remain correlated with index parity;
the half-label law and the exclusion of every simple tangent
remain open here.

\paragraph{AG and harmonic-analysis translation.}
The zero divisor is counted with its actual local lengths.
At a reduced point, the equation $\zeta_t(s)=0$ gives
$d\zeta_t=\zeta'(\rho)\,ds-\rho\zeta(\rho+1)\,dt$;
its section has tangent $\partial_t+v_\rho\partial_s$.
The observable $r_\rho$ is the absolute normalized real
normal component of this tangent relative to the critical
line. The theorem gives a positive lower counting density
on that reduced locus, retaining all multiplicities in the
ambient denominator. The K\"ahler-differential local model
is the one of SGA~1, Expos\'e~I, section~3, Proposition~1.
For the exterior phase, the characters $x\mapsto
e^{i\gamma\log x}$ are characters of the multiplicative
group $\mathbb Q_{>0}$, whose compact dual is the product
of prime circles. Landau's formula identifies the relevant
Haar means, and absolute Dirichlet convergence makes the
ratio's constant term $\sum\mu(n)n^{-3}=1/\zeta(3)$.
This is an actual harmonic-analysis computation, not an
asserted Frobenius equidistribution theorem for the zeros
or an arithmetic family realizing the denominator shift.


\section{Joint phase cancellation and primitive Dirichlet transversality}

The two marginal phase means used in NN-A036 can be replaced by a
proved mixed mean. This also answers the fixed primitive Dirichlet
analogue. Throughout this section a character and its conductor are
fixed as the height tends to infinity. No uniformity over conductors,
family average, full GUE conjecture, or simplicity assumption is used.
The character of conductor one denotes the Riemann zeta function.

\begin{theorem}[NN-A037: a joint arithmetic mode and definite normal motion]
Assume A. Let $L(s,\chi)$ be any fixed primitive Dirichlet $L$-function,
including $\zeta$. Write $a\in\{0,1\}$ for its parity and put
\[
 Q_\chi(s)=(q/\pi)^{(s+a)/2}\Gamma((s+a)/2),\qquad
 \Lambda_\chi(s)=Q_\chi(s)L(s,\chi),\qquad
 \Lambda_\chi(s)=\epsilon_\chi\overline{\Lambda_\chi(1-\bar s)}.
\]
Choose a square root of $\epsilon_\chi$ and the corresponding continuous
real phase $\theta_\chi$ for which
$Z_\chi(t)=e^{i\theta_\chi(t)}L(1/2+it,\chi)$ is real. Let
$N_\chi(T)$ count positive-ordinate nontrivial zeros with multiplicity,
and $N_{\chi,s}(T)$ count simple zeros. At a simple zero $\rho=1/2+i\gamma$
the specified norm shift has velocity
\[
 v_{\rho,\chi}=
 \frac{\rho L(\rho+1,\chi)}{L'(\rho,\chi)},\qquad
 r_{\rho,\chi}=\frac{|\Re v_{\rho,\chi}|}{|v_{\rho,\chi}|}.
\]
For every $\sigma>1$, let
$\phi_{\chi,\sigma}(t)=\arg L(\sigma+it,\chi)$ be the Euler-product
argument. Then
\begin{equation}\label{eq:A037-joint}
 \frac1{N_\chi(T)}\sum_{0<\gamma\le T}
 e^{2i\theta_\chi(\gamma)}
 \frac{L(\sigma+i\gamma,\chi)}
      {\overline{L(\sigma+i\gamma,\chi)}}\longrightarrow0.
\end{equation}
Every sum over zeros in this theorem retains multiplicity. In particular,
with $\delta_\gamma=\arctan(1/(2\gamma))$,
\begin{equation}\label{eq:A037-second}
 \frac1{N_\chi(T)}\sum_{0<\gamma\le T}
 \cos^2\bigl(\theta_\chi(\gamma)+
              \phi_{\chi,3/2}(\gamma)-\delta_\gamma\bigr)
 \longrightarrow\frac12.
\end{equation}
Consequently,
\[
 \liminf_{T\to\infty}\frac1{N_\chi(T)}
 \sum_{\substack{0<\gamma\le T\\\rho\ {
m simple}}}
 r_{\rho,\chi}^{\,2}\ge\frac16.
\]
At least a lower proportion $1/6$ of all zero occurrences are simple
and transverse; relative to the simple zeros the lower proportion is
$1/4$. More quantitatively, if $0<\eta<1/2$ and $G_{\chi,\eta}(T)$
counts simple zeros with $r_{\rho,\chi}\ge\eta$, then
\begin{equation}\label{eq:A037-counts}
 \liminf\frac{G_{\chi,\eta}(T)}{N_\chi(T)}
 \ge\frac{1-4\eta^2}{6(1-\eta^2)},\qquad
 \liminf\frac{G_{\chi,\eta}(T)}{N_{\chi,s}(T)}
 \ge\frac{1-4\eta^2}{4(1-\eta^2)}.
\end{equation}
For example, at $\eta=1/4$ these bounds are $2/15$ and $1/5$.
For each fixed $\varepsilon>0$, every zero counted at a fixed positive
threshold also satisfies
$|\Re v_{\rho,\chi}|\gg_{\chi,\eta,\varepsilon}
\gamma^{1-\varepsilon}$ as $\gamma\to\infty$.
No positive/negative sign proportion is asserted.
\end{theorem}

\begin{lemma}[Twisted first Hardy mode]\label{lem:A037-twist}
For every fixed $x>0$, under GRH for the specified primitive character,
\begin{align}\label{eq:A037-twist}
 \sum_{0<\gamma\le T}e^{2i\theta_\chi(\gamma)}x^{i\gamma}
 &=-\frac{i^a\epsilon_\chi^{-1}}{\sqrt{qx}}
    \sum_{n\le qxT/(2\pi)}\chi(n)\Lambda(n)e^{-2\pi i n/(qx)}
    +O_{\chi,x}(\sqrt T\log^2 T)\\
 &=O_{\chi,x}(T).\nonumber
\end{align}
Changing the convention at the upper endpoint changes the displayed
expression by an admissible error. For $q=1$, $x=1$, this is the
$\kappa=1$ formula of Arias de Reyna, Corollary~3.2, equation~(13).
The proof below gives the required twist and character extension,
including their constants; they are not inferred from that untwisted
statement.
\end{lemma}

\begin{proof}
Put
\[
 \Delta_\chi(s)=\epsilon_\chi Q_\chi(1-s)/Q_\chi(s).
\]
At a critical-line point,
$\Delta_\chi(1/2+it)^{-1}=e^{2i\theta_\chi(t)}$.
The gamma quotient is meromorphic with all of its zeros and poles real;
it is holomorphic and nonzero in the open upper half-plane. In every
fixed vertical strip, Stirling's formula gives, uniformly for $t\ge2\pi$,
\begin{equation}\label{eq:A037-stirling}
 \Delta_\chi(\sigma+it)^{-1}
 =\epsilon_\chi^{-1}(qt/(2\pi))^{\sigma-1/2}
 e^{i\{t\log(qt/(2\pi))-t+\pi(2a-1)/4\}}
 (1+O_\chi(t^{-1})).
\end{equation}
This follows directly by applying the usual Stirling expansion to
$\Gamma((s+a)/2)$ and $\Gamma((1-s+a)/2)$. The phase constant is
recorded because it supplies $i^a$ in \eqref{eq:A037-twist}.

First choose $T'$ in $(T,T+1)$ at distance $\gg_\chi1/\log T$ from
all zero ordinates. Such a choice is possible because the number of
zeros in this interval, or in a fixed enlargement of it, is
$O_\chi(\log T)$: see Iwaniec--Kowalski, Proposition~5.7(1),
equation~(5.27), pp.102--103. Their Theorem~5.6,
equations~(5.23)--(5.24), p.102, gives the Hadamard product and
partial-fraction formula. That formula gives
\[
 \frac{L'}L(\sigma+iT',\chi)=O_\chi(\log^2 T)
 \quad(-3/2\le\sigma\le 1+1/\log T).
\]
For clarity, this standard bound does not require simple zeros:
subtract the partial-fraction expansion at $2+iT'$; the poles at
distance at most one contribute at most $O(\log T)$ terms of size
$O(\log T)$, counting multiplicity. The remaining terms are bounded
by a constant times $\sum_{k\ge1}\log(T+k+2)/k^2=O(\log T)$.
The gamma and pole factors contribute $O(\log T)$. This proves the
claimed bound and also explains the admissible contour choice.

Integrate
\[
 \frac1{2\pi i}\frac{L'(s,\chi)}{L(s,\chi)}
          \Delta_\chi(s)^{-1}x^{s-1/2}
\]
around the rectangle with left side $\Re s=-3/2$, right side
$c=1+1/\log T'$, bottom at a fixed
$t_0\in(2\pi,2\pi+1)$ avoiding zero ordinates, and top at $T'$.
Any finite set of omitted lower zeros contributes $O_{\chi,x}(1)$.
There are no gamma poles or trivial zeros in this rectangle. The zeta
pole when $q=1$ is below it. The residue at a zero is its multiplicity
times $\Delta_\chi(\rho)^{-1}x^{\rho-1/2}$; GRH identifies this with
the summand in \eqref{eq:A037-twist}.

The bottom edge is $O_{\chi,x}(1)$. On the top edge the preceding
logarithmic-derivative bound and \eqref{eq:A037-stirling} give
$O_{\chi,x}(\sqrt T\log^2T)$ (in fact one logarithm can be saved).
On the left edge the functional equation and absolute convergence at
$5/2-it$ give $L'/L=O_\chi(\log(t+2))$, while
$|\Delta_\chi|^{-1}=O_\chi(t^{-2})$. Its integral is bounded.

Extending the right-edge integral down from $t_0$ to $2\pi$ costs
$O_{\chi,x}(1)$: on this fixed compact segment the integrand is
uniformly bounded as $c\downarrow1$. On the right edge use the absolutely convergent series
$L'/L=-\sum\chi(n)\Lambda(n)n^{-s}$ and Stirling. Its error is
\[
 O_{\chi,x}\left(\sum_{n\ge2}\frac{\Lambda(n)}{n^c}
                 \int_{2\pi}^{T'}t^{c-3/2}\,dt\right)
 =O_{\chi,x}(\sqrt T\log T).
\]
Here $T^{c-1}=O(1)$ and $\sum\Lambda(n)n^{-c}=O(\log T)$.
For the remaining integral set $u=t/(2\pi)$, $X=T'/(2\pi)$,
and $b=qx$. Its value is
\begin{equation}\label{eq:A037-stationary-integral}
 -\epsilon_\chi^{-1}e^{i\pi(2a-1)/4}
 b^{c-1/2}\sum_{n\ge2}\frac{\chi(n)\Lambda(n)}{n^c}
 \int_1^X u^{c-1/2}
 e^{2\pi i u\log(bu/(en))}\,du.
\end{equation}

Here are the full summation bounds for stationary phase. Set $r=n/b$.
Split $[1,X]$ into its intersection with $[r/2,2r]$ and the two
remaining intervals. On either remaining interval the derivative of
$u\log(u/(er))$ has absolute value at least $\log2$. The first
derivative test bounds each such integral by $O(X^{c-1/2})$.
Summing gives $O_{\chi,x}(\sqrt T\log T)$. The finitely many $r<1$
can be handled by the same test away from a bounded interval and add
at most this error.

On $[A,B]=[1,X]\cap[r/2,2r]$, the stationary-phase lemma of Arias
de Reyna, Lemma~2.5 (the underlying primary lemmas are Levinson and
Gonek), states that the integral is
\[
 {\bf1}_{A\le r\le B}\,r^c e^{i\pi/4}e^{-2\pi i r}
 +O\left(r^{c-1/2}
 +\frac{r^{c+1/2}}{|A-r|+\sqrt r}
 +\frac{r^{c+1/2}}{|B-r|+\sqrt r}\right).
\]
Its constants are uniform for $c-1/2$ in the compact interval in use.
Multiplication by the coefficient in
\eqref{eq:A037-stationary-integral} turns the main term into
\[
 -i^a\epsilon_\chi^{-1}b^{-1/2}
       \chi(n)\Lambda(n)e^{-2\pi i n/b}.
\]
The first remainder sums to $O_{\chi,x}(\sqrt T)$, since
$\sum_{n\le C T}\Lambda(n)n^{-1/2}\ll\sqrt T$.
Except for finitely many $n$, $A=r/2$, so the lower-endpoint
remainder has the same bound. If $r\le X/2$, then $B=2r$ and the
upper-endpoint remainder again has that bound. For $X/2<r\le2X$
it is at most
\[
 C_{\chi,x}\sqrt T\log(T+2)
 \sum_{bX/2<n\le2bX}
       \frac1{|n-bX|+\sqrt T}
 \ll_{\chi,x}\sqrt T\log^2T.
\]
The last bound follows by splitting the distances into unit intervals,
or comparing with $\int_0^{C T}(u+\sqrt T)^{-1}\,du$; it is valid
even when $bX$ is an integer. The constant number of small $n$ has
already been absorbed. This proves the displayed asymptotic at $T'$.
Chebyshev's bound $\sum_{n\le y}\Lambda(n)\ll y$
(Iwaniec--Kowalski, Section~2.2, equation~(2.12), p.33) gives its
$O_{\chi,x}(T)$ consequence without requiring cancellation in the
additive prime sum.

Finally, under GRH each summand on the left has modulus one, and
there are $O_\chi(\log T)$ zeros between $T$ and $T'$. The change
in the prime sum has $O_{q,x}(1)$ terms, each $O(\log T)$.
Thus the formula holds at every real cutoff $T$, with the stated
endpoint convention and error.
\end{proof}

\begin{lemma}[The precise fixed-character simplicity input]
Under GRH for a fixed primitive Dirichlet $L$-function,
\[
 \liminf_{T\to\infty}\frac{N_{\chi,s}(T)}{N_\chi(T)}\ge\frac23.
\]
This is a classical restricted-pair-correlation consequence, used here
as background rather than counted as a new simplicity theorem.
\end{lemma}
\begin{proof}
We specify the imported theorem and the multiplicity calculation to
avoid confusing a family average with a result for a fixed character.
Rudnick--Sarnak, \emph{Zeros of principal $L$-functions and random
matrix theory}, Theorem~1.2, specializes at degree $m=1$, correlation
order two, to every fixed primitive Dirichlet character. Its hypothesis
(1.7) is automatic here (and the theorem explicitly permits $m\le3$).
Their zeros are indexed with multiplicity; equation~(1.4) sums over
distinct indices, not necessarily distinct ordinates. The test functions
$f(x_1,x_2)=g(x_1-x_2)$ have Fourier support in
$(u,-u)$ with $|u|<1$, exactly the permitted
$|\xi_1|+|\xi_2|<2$. With $K(u)=\sin(\pi u)/(\pi u)$, this gives
\[
 \frac1N\sum_{j\ne k}g(\widetilde\gamma_j-
                         \widetilde\gamma_k)
 \longrightarrow\int_{\mathbb R}g(u)(1-K(u)^2)\,du.
\]
Here the indices run through the first $N$ positive zeros, and their
normalization is that of the cited theorem. Equal ordinates have
exactly equal normalized ordinates, which is all that is needed below.

Take nonnegative Schwartz tests $g$ with $g(0)=1$ and Fourier support
strictly inside $(-1,1)$. Every zero of multiplicity $m\ge2$ contributes
$m(m-1)\ge m$ to the sum on the left. Hence the fraction of occurrences
belonging to nonsimple zeros is bounded above by the displayed integral
in the limit. The same statement holds at cutoffs $T$ that include each
whole multiplicity block.

For completeness the endpoint optimization uses only admitted
Schwartz tests. For $0<\beta<1$, approximate
$\beta^{-1}{\bf1}_{[-\beta/2,\beta/2]}$ in $L^2$ and $L^1$ by
smooth real even functions $b_j$ supported inside that interval with
$\int b_j=1$. Then $g_j=|\widehat b_j|^2$ are nonnegative Schwartz,
$g_j(0)=1$, have Fourier support in $[-\beta,\beta]$, and converge in
$L^1$ to $g_\beta(u)=(\sin(\pi\beta u)/(\pi\beta u))^2$.
Since $0\le1-K^2\le1$, their integrals converge to
\[
 \int g_\beta(1-K^2)
 =\frac1\beta-1+\frac\beta3.
\]
Indeed $\widehat g_\beta(u)=\beta^{-1}(1-|u|/\beta)_+$ and
$\widehat {K^2}(u)=(1-|u|)_+$, so their product integral is
$1-\beta/3$. First take $N\to\infty$, then the smooth approximation,
then $\beta\uparrow1$. The nonsimple fraction is at most $1/3$.
Central zeros and finitely many endpoint choices do not affect this
positive-height limit.
\end{proof}

\begin{proof}[Proof of NN-A037]
For $\sigma>1$ the absolutely convergent expansion is
\[
 \frac{L(\sigma+it,\chi)}{\overline{L(\sigma+it,\chi)}}
 =\sum_{n,m\ge1}
 \frac{\chi(n)\mu(m)\overline{\chi(m)}}{(nm)^\sigma}
       e^{it\log(m/n)}.
\]
Its absolute coefficient sum is at most $\zeta(\sigma)^2$.
For every fixed pair $m,n$, Lemma~\ref{lem:A037-twist}, with
$x=m/n$, says that the normalized mixed sum tends to zero because
$N_\chi(T)\sim T\log T/(2\pi)$. For the positive-height counting
convention this is the calculation in the proof of
Iwaniec--Kowalski, Theorem~5.8, equations~(5.34)--(5.36), pp.104--105;
Rudnick--Sarnak equation~(2.11) records the corresponding two-sided
count. All such normalized sums have
absolute value at most one under GRH. Dominated convergence therefore
proves \eqref{eq:A037-joint}, without any uniform estimate as $m/n$
varies. This proves one mixed Fourier mode; it does not assert
independence or uniform distribution of the full joint phase.

Multiplication by $e^{-2i\delta_\gamma}$ changes the mean by $o(1)$.
To see this without a height-uniform error claim, fix $Y$, discard the
finitely many zeros below $Y$, and bound the remaining error by
$\sup_{\gamma>Y}2|\delta_\gamma|$; then take $T\to\infty$ and
$Y\to\infty$. Taking real parts and using
$\cos^2 u=(1+\cos2u)/2$ proves \eqref{eq:A037-second}.

At a simple zero,
$L'(\rho,\chi)=-i e^{-i\theta_\chi(\gamma)}Z_\chi'(\gamma)$,
where $Z_\chi'(\gamma)$ is real and nonzero. Since
$\arg(i\rho)=\pi-\delta_\gamma$, the exact velocity formula gives
\[
 r_{\rho,\chi}
 =\big|\cos(\theta_\chi(\gamma)+
              \phi_{\chi,3/2}(\gamma)-\delta_\gamma)\big|.
\]
The choice of square root of the root number changes the phase by
$\pi$ and therefore changes neither this expression nor the theorem.

Put $p=N_{\chi,s}(T)/N_\chi(T)$ and
$g=G_{\chi,\eta}(T)/N_\chi(T)$. The preceding simplicity lemma
gives $p\ge2/3-o(1)$. The contribution of all nonsimple occurrences
to \eqref{eq:A037-second} is at most $1-p$, so the simple contribution
has lower limit at least $1/6$. Also
\[
 \frac12+o(1)
 \le(1-p)+\eta^2(p-g)+g
 =1-(1-\eta^2)p+(1-\eta^2)g.
\]
It follows that
$g\ge p-[2(1-\eta^2)]^{-1}+o(1)$, which gives the first inequality
in \eqref{eq:A037-counts}. Dividing by $p$, bounded below by
$2/3-o(1)$, gives the second. Letting $\eta\downarrow0$ proves the
two asserted nonzero-velocity proportions. The numerical fractions at
$\eta=1/4$ are exact rational substitutions.

Finally, the Euler product gives
$|L(3/2+i\gamma,\chi)|\ge1/\zeta(3/2)$. The exact classical input
for the fixed-character Lindelof bound is Iwaniec--Kowalski,
\emph{Analytic Number Theory}, Section~5.7, Corollary~5.20,
equation~(5.60), printed p.116, following Theorem~5.19 on that page.
The hypotheses are GRH and the Ramanujan--Petersson coefficient
bound; the latter holds trivially for a Dirichlet character, since
$|\chi(n)|\le1$. The analytic conductor for fixed $\chi$ is
$\asymp_\chi 1+|t|$. Cauchy's derivative formula on a small
fixed circle, using the functional equation on its left half, gives
$|L'(1/2+i\gamma,\chi)|\ll_{\chi,\varepsilon}\gamma^\varepsilon$.
Explicitly, choose the circle radius smaller than $\varepsilon/4$;
Lindelof on the right and the gamma ratio bound on the left make the
supremum $O(\gamma^\varepsilon)$ after reducing the auxiliary exponent.
Thus $|v_{\rho,\chi}|\gg\gamma^{1-\varepsilon}$, and multiplication
by the fixed margin $\eta$ proves the final assertion.
\end{proof}

\begin{corollary}[Sharper proportions for the Riemann zeta function]
\label{cor:A037-zeta-sharpening}
For $\zeta$, under A, the lower proportion of all zero occurrences
that are simple and transverse is at least $11/54$, and the lower
proportion among simple zeros is at least $11/38$. More precisely,
for $0<\eta<\sqrt{11/38}$,
\[
 \liminf\frac{G_{\zeta,\eta}(T)}{N(T)}
       \ge\frac{19}{27}-\frac1{2(1-\eta^2)},\qquad
 \liminf\frac{G_{\zeta,\eta}(T)}{N_s(T)}
       \ge1-\frac{27}{38(1-\eta^2)}.
\]
At $\eta=1/4$ the two bounds are $23/135$ and $23/95$.
The absolute-speed clause of the theorem holds on these subsets.
\end{corollary}

\begin{proof}
Bui--Heath-Brown, \emph{On simple zeros of the Riemann
zeta-function}, Bull.\ London Math.\ Soc.\ \textbf{45} (2013),
953--961, arXiv:1302.5018v1, Section~1, the unnumbered Theorem
on p.~1, proves under RH alone that
\[
                \liminf_{T\to\infty}\frac{N_s(T)}{N(T)}
                   \ge p:=\frac{19}{27}.
\]
The same page explicitly defines $N$ with multiplicity and $N_s$
as the count of simple zeros. This is an imported classical
simplicity theorem, not an additional new finding here.
Write $M_s(T)=\sum_{\rho\ {\rm simple}}r_\rho^2$. The exact
second-moment identity \eqref{eq:A037-second}, subtracting the
nonsimple contribution which is bounded above by $N-N_s$, gives
\[
                   M_s\ge N_s-\tfrac12N+o(N).
\]
On the simple zeros the upper bound
$M_s\le G_{\zeta,\eta}+\eta^2(N_s-G_{\zeta,\eta})$ implies
\[
 G_{\zeta,\eta}\ge N_s-\frac{N}{2(1-\eta^2)}+o(N).
\]
Division first by $N$ and then by $N_s$, using $N_s/N\ge p-o(1)$,
proves both inequalities. Positivity of the lower bounds holds
for $\eta^2<1-1/(2p)=11/38$. Letting a fixed positive $\eta$
decrease to zero proves the two transverse-count bounds.
At $\eta=1/4$ the exact differences are
$19/27-8/15=23/135$ and $1-216/285=23/95$.
The previously proved speed estimate only requires a fixed positive
angular margin, so it applies without change.
\end{proof}

\paragraph{Exact scope and the Dedekind obstruction.}
This theorem is about every fixed primitive degree-one function. It
strengthens NN-A036 for $\zeta$ and proves its requested Dirichlet
analogue. A general Dedekind zeta function is not primitive. Repeated
Artin factors can produce nonsimple zeros systematically, and the
degree-one first-mode calculation cannot be imported for its product
without a new estimate. Nor do \eqref{eq:A037-joint} and
\eqref{eq:A037-second} distinguish the two signs of $\Re v$: the
sign of the real derivative $Z_\chi'(\gamma)$ is lost upon squaring.
The index-parity correlation and the magnitude-tail obstruction
recorded in Q-A036 therefore remain genuine sign-density gates.

\paragraph{Algebraic-geometric and representation translation.}
A primitive character gives its actual rank-one Artin representation
on the finite etale cyclotomic cover of
$\operatorname{Spec}\mathbb Z[1/q]$. Its root number supplies the
half-phase that makes the completed critical-line germ Real, including
for complex characters. At a reduced analytic zero the specified
deformation has differential
\[
 dL_t=L'(\rho,\chi)\,ds-\rho L(\rho+1,\chi)\,dt,
\]
so its tangent section has the stated normalized real normal
component; the ambient divisor is counted with its actual local
lengths. SGA~1, Expose~I, section~3, Proposition~1 is the local
differentials dictionary. On the compact dual of $\mathbb Q_{>0}$,
the mixed-mode theorem annihilates the first Hardy character times
each fixed prime-circle character, and absolute Euler convergence
extends that annihilation to the displayed exterior ratio. This is
a precise Fourier statement, not a full product-Haar or Frobenius
equidistribution theorem for analytic zeros. Rudnick--Sarnak's
restricted correlation theorem controls the nonreduced occurrence
mass sufficiently to obtain the explicit reduced-locus density.
No shifted arithmetic family, pure motive, or complexity separation
is supplied by this dictionary.

\paragraph{Primary locators used.}
Arias de Reyna, \emph{On the distribution (mod~1) of the normalized
zeros of the Riemann zeta-function}, arXiv:1403.7126v1,
Lemma~2.5, Corollary~3.2 equation~(13), and Section~4.4 give the
stationary-phase input and the untwisted comparison.
Rudnick--Sarnak, \emph{Zeros of principal $L$-functions and random
matrix theory}, author manuscript PDF pp.1--5, especially the
multiplicity convention on p.1, definitions (1.2)--(1.4) on p.3,
and Theorems~1.1--1.2 on pp.4--5, give precisely the fixed-character
restricted correlation input. Their distinction between primitive
and product functions is on p.6. The manuscript is available at
\url{https://empslocal.ex.ac.uk/people/staff/mrwatkin/zeta/rudnick-sarnak.pdf}.
The completion and root number use DLMF~25.15.5--25.15.6; the
Stirling expansion is DLMF~5.11.3 with its remainder formulation
in Section~5.11(ii). The absolute-speed clause uses
Iwaniec--Kowalski, \emph{Analytic Number Theory} (AMS Colloquium
Publications~53, 2004), Section~5.7, Corollary~5.20, equation~(5.60),
p.116, with Theorem~5.19's GRH and Ramanujan hypotheses checked
as above. The same book supplies the auxiliary Chebyshev bound at
Section~2.2, equation~(2.12), p.33; Hadamard and the logarithmic
derivative at Theorem~5.6, equations~(5.23)--(5.24), p.102; the
local zero count at Proposition~5.7(1), equation~(5.27), pp.102--103;
and the positive-height count in Theorem~5.8's proof,
equations~(5.34)--(5.36), pp.104--105.
The printed text was checked in an online OCR transcription
of the primary book, rather than a page image; the publisher record
is \url{https://doi.org/10.1090/coll/053}.
The new mixed-mode and transversality deductions
are proved in full above; no historical novelty claim is made.


\section{Both branch directions recur in every fixed multiplicative height window}
\label{sec:analytic-signed-recurrence}

The phase-square theorem loses the sign of an individual derivative.
Nevertheless its exact mean can be combined with the symmetric
residue trace to rule out an entire high window having only one
branch direction. Multiple zeros are retained in both ingredients.
The result below proves signed recurrence, not a positive counting
density or a half-label law.

\begin{lemma}[Primitive Dirichlet good heights and symmetric traces]
\label{lem:A039-dirichlet-trace}
Fix a primitive Dirichlet character $\chi$, allowing conductor one,
and assume its GRH. For every $e>0$ there are heights
$\tau_n\in(n,n+1)$ avoiding all zero ordinates, for all sufficiently
large integers $n$, such that
\begin{equation}\label{eq:A039-good-heights}
 \sup_{1/2\le\sigma\le2}
       |L(\sigma+i\tau_n,\chi)|^{-1}\ll_{\chi,e}\tau_n^e.
\end{equation}
These heights may depend on $e$. For every fixed integer $k\ge0$
one can choose them, with an auxiliary exponent sufficiently small,
so that
\begin{equation}\label{eq:A039-trace}
 \sum_{0<\Im\rho<\tau_n}
 \operatorname{Res}_{s=\rho}
        \frac{s^kL(s+1,\chi)}{L(s,\chi)}
 =\frac{i^k\tau_n^{k+1}}{2\pi(k+1)}
                  +O_{\chi,k,e}(\tau_n^{k+e}).
\end{equation}
Each multiple zero contributes its full residue. In particular the
real part of the $k=1$ sum is $O_{\chi,e}(\tau_n^{1+e})$.
\end{lemma}

\begin{proof}
We give the good-height construction explicitly. The fixed-character
GRH--Lindelof theorem already checked in NN-A037 and the functional
equation give, for every fixed compact real interval and every $d>0$,
\begin{equation}\label{eq:A039-growth}
 |L(\sigma+it,\chi)|\ll_{\chi,d}
           (1+|t|)^{\max(1/2-\sigma,0)+d},
\end{equation}
uniformly in that interval at large height.

Here are the elementary complex-analysis estimates used below, with
their proofs. If $F$ is holomorphic on $|z|\le R$, put
$G=F-F(0)$ and $M=\sup\Re G\ge0$. When $M>0$, the function
$G/(2M-G)$ maps the disc into the closed unit disc and vanishes at
zero, because $|G|^2\le|2M-G|^2$ is equivalent to $\Re G\le M$.
Schwarz's lemma therefore gives
$|G(z)|\le2M|z|/(R-|z|)$. If $M=0$, the maximum principle
makes $G=0$. This proves the Borel--Caratheodory inequality in
the form used here. To obtain Jensen's formula, divide a holomorphic
function with nonzero center value by its factors $(z-a)$ for the
zeros inside the circle. The logarithm of the remaining modulus is
harmonic and has its center value as circle average; for $|a|<R$,
the circle average of $\log|Re^{i\vartheta}-a|$ is $\log R$, as
follows by expanding $\log(1-ae^{-i\vartheta}/R)$ and averaging.
Subtracting the center identity gives
$\sum\log(R/|a|)=\operatorname{avg}\log|f(Re^{i\vartheta})|
-\log|f(0)|$. Boundary zeros are handled by a radius limit.
Finally the Poisson kernel of a disc at radius $s<R$ is at most
$(R+s)/(R-s)$ times its center kernel. Integrating against
nonnegative boundary values proves the stated Harnack bound.
These arguments justify all three estimates without an additional
unlocated analytic input.

We first establish the reciprocal estimate off the critical line
needed in the construction. On a disc centered at $2+it$ of radius
$3/2-a/2$, where $a>0$ is fixed and small, GRH supplies a holomorphic
branch of $\log L$, normalized at the center by the absolutely
convergent Euler logarithm. Its value at the center is uniformly
bounded. Equation~\eqref{eq:A039-growth} makes the real part of this
logarithm at most $d\log|t|+O_{\chi,d,a}(1)$. The
Borel--Caratheodory inequality on the smaller concentric disc of
radius $3/2-a$ gives
\[
 |\log L(\sigma+it,\chi)|
       \le C_a d\log|t|+O_{\chi,d,a}(1),
                 \qquad 1/2+a\le\sigma\le2.
\]
Indeed the usual inequality is
$\sup_{|z|\le r}|f(z)|\le
2r\sup_{|z|\le R}\Re f(z)/(R-r)
+(R+r)|f(0)|/(R-r)$; subtracting the bounded center value
gives the displayed estimate. Choosing $d$ after $a$ proves
\begin{equation}\label{eq:A039-away-inverse}
 \sup_{1/2+a\le\sigma\le2}|L(\sigma+it,\chi)|^{-1}
                      \ll_{\chi,a,b}|t|^b
 \quad\hbox{for every }b>0.
\end{equation}
At large height the possible zeta pole is outside all these discs.

Fix $0<r<1/32$, put $t_0=n+1/2$ and
$z_0=1/2+2r+it_0$, and write $f(w)=L(z_0+w,\chi)$.
For all sufficiently large $n$, depending on $r$ and $\chi$,
\eqref{eq:A039-growth} with $d=r$ and
\eqref{eq:A039-away-inverse} imply
\[
 \max_{|w|\le8r}|f(w)|\le n^{8r},\qquad
                         |f(0)|\ge n^{-r}.
\]
Jensen's formula on radius $8r$ shows that the number $M$ of
zeros in $|w|\le7r$, counted with multiplicity, satisfies
\[
 M\log(8/7)\le9r\log n,
                \qquad M\le72r\log n,
\]
where $\log(8/7)\ge1/8$ follows by integrating $1/x$ from $1$
to $8/7$. Jensen's inequality is also valid by limits if the
outer circle contains a zero.

Choose $R\in(6r,7r)$ with no zero on its boundary. Form the
finite Blaschke product over its zeros $a$, including multiplicity,
\[
 B(w)=\prod_{|a|<R}\frac{R(w-a)}{R^2-\bar a w},
                    \qquad g(w)=f(w)/B(w).
\]
The quotient is holomorphic and zero-free on the closed disc.
On its boundary $|B|=1$, so $|g|\le n^{8r}$ throughout the
disc by the maximum principle. Also $|B(0)|\le1$ and
$|g(0)|\ge n^{-r}$. Therefore
$u(w)=8r\log n-\log|g(w)|$ is a nonnegative harmonic function,
with $u(0)\le9r\log n$. The Poisson kernel bounds, equivalently
Harnack's inequality, give for $|w|\le3r$
\[
 u(w)\le\frac{R+3r}{R-3r}u(0)\le3u(0),
                  \qquad |g(w)|\ge n^{-19r}.
\]

For $1/2\le\sigma\le1/2+r$ and $|t-t_0|\le r$ one has
$|\sigma+it-z_0|\le\sqrt5\,r<3r$. Each Blaschke factor at
this point has modulus at least
\[
 \frac{|\sigma+it-z_0-a|}{R+|\sigma+it-z_0|}
                         \ge\frac{|t-\gamma_a|}{10r},
 \qquad\gamma_a=\Im(z_0+a).
\]
For every real $\gamma$, translation and positivity give
\[
 \frac1{2r}\int_{t_0-r}^{t_0+r}
       \log^+\frac{10r}{|t-\gamma|}\,dt
 \le\frac1{2r}\int_{-10r}^{10r}
       \log\frac{10r}{|u|}\,du=10.
\]
Consequently a height $t_*\in[t_0-r,t_0+r]$ can be chosen
away from all zeros with
\[
             \sum_{|a|<R}\log^+\frac{10r}{|t_*-\gamma_a|}
                            \le10M.
\]
For that same height and every $\sigma\in[1/2,1/2+r]$,
\[
 |L(\sigma+it_*,\chi)|
            \ge n^{-19r}e^{-10M}\ge n^{-739r}.
\]
The estimate away from the line, \eqref{eq:A039-away-inverse},
covers $\sigma\in[1/2+r,2]$. Choose $r$ so small that
$739r<e/2$, and use exponent $e/2$ in that latter estimate.
This proves \eqref{eq:A039-good-heights}. The selected height
lies strictly in $(n,n+1)$. Any zero at its ordinate and in the
critical strip would, by GRH, be at $1/2+it_*$ and would have
belonged to the Blaschke disc, so no zero ordinate was missed.

To prove the trace formula, integrate
$s^kL(s+1,\chi)/L(s,\chi)$ on the rectangle with vertical
sides $2$ and $-1/2$, a fixed positive bottom height avoiding
zeros, and top $\tau_n$. The vertical functional equation
relates the modulus at $\sigma+i\tau_n$ to that at
$1-\sigma+i\tau_n$ for the same character. Thus
\eqref{eq:A039-good-heights} extends to the top's left half
with the additional factor $\tau_n^{\sigma-1/2}$.
On the right side the absolute Euler expansion is
\[
 \frac{L(s+1,\chi)}{L(s,\chi)}
   =\sum_{m\ge1}b_\chi(m)m^{-s},\qquad
 b_\chi(m)=\frac{\chi(m)}m\prod_{p\mid m}(1-p),
 \quad b_\chi(1)=1,\quad |b_\chi(m)|\le1.
\]
The constant term contributes
$i^k\tau_n^{k+1}/(2\pi(k+1))+O_k(1+\tau_n^k)$.
For $m\ge2$, one integration by parts bounds the oscillatory
integral by $C_k(1+\tau_n)^k/\log m$, and the sum with
$|b_\chi(m)|m^{-2}$ converges. On the left vertical side,
the functional equation and Euler convergence give
$|1/L(-1/2+it,\chi)|\ll_\chi t^{-1}$, while the numerator
is $O(t^d)$ by Lindelof. Its integral is
$O_{\chi,k,d}(\tau_n^{k+d})$. The top integral is
$O(\tau_n^{k+e'+d})$ when exponent $e'$ was used in the
good-height construction. Choose $e',d$ with $e'+d<e$.
The bottom and the finitely many omitted small zeros contribute
$O(1)$. All possible trivial zeros and poles lie outside this
positive-height rectangle. The residue theorem proves
\eqref{eq:A039-trace} with every full multiple-zero residue.
\end{proof}

\begin{theorem}[NN-A039: recurrent arithmetic signs, retaining multiplicity]
\label{thm:A039-recurrence}
Assume A and fix a primitive Dirichlet character $\chi$, including
the zeta case. Let
\[
 L_t(s,\chi)=\sum_{n\ge1}\chi(n)(n+t)^{-s},\qquad t>0,
\]
with its local analytic continuation in $s$ and small $t$ as
in NN-A004. At each nontrivial zero
$\rho=1/2+i\gamma$, $\gamma>0$, consider its local Puiseux
branches at $t=0$. For every fixed $B>1$ there is $T_B$ such
that every $T\ge T_B$ has the following two witnesses with
centers $T<\gamma<BT$:

\begin{enumerate}
\item A branch whose first Puiseux coefficient has strictly
positive real part. At a simple zero this says $\Re v_\rho>0$.
\item A branch whose first Puiseux coefficient has strictly
negative real part, or a vertical branch at a double zero whose
second Puiseux coefficient has strictly negative real part.
At a simple zero this says $\Re v_\rho<0$.
\end{enumerate}

Both witnesses lie strictly on the indicated side of the critical
line for sufficiently small positive $t$, with the smallness bound
allowed to depend on the chosen center. The two centers may
coincide at a multiple zero.

Let $R_\chi(X)$ count distinct positive-ordinate centers at most
$X$ supporting a branch in the first class, and $L_\chi(X)$
count distinct centers supporting a branch in the second class.
Then
\begin{equation}\label{eq:A039-superlog}
 \frac{R_\chi(X)}{\log X}\longrightarrow+\infty,
 \qquad\frac{L_\chi(X)}{\log X}\longrightarrow+\infty.
\end{equation}
In particular both directions occur infinitely often. If all
sufficiently high zeros of the specified function are simple,
these are recurrent strictly positive and strictly negative
ordinary velocities. No eventual simplicity, positive counting
density, comparable sign counts, or common deformation-time
scale is asserted.
\end{theorem}

\begin{proof}
\emph{The exact phase probe at a vertical double zero.}
Use the root-normalized real Hardy function $Z_\chi$ from NN-A037
and put
\[
 w_\rho=\cos^2\bigl(\theta_\chi(\gamma)
          +\arg L(3/2+i\gamma,\chi)-\delta_\gamma\bigr),
                     \qquad\delta_\gamma=\arctan(1/(2\gamma)).
\]
It is defined at every zero and belongs to $[0,1]$. At a simple
zero it equals $r_\rho^2$. At a zero of order $m$, write
$A_m=L^{(m)}(\rho,\chi)/m!$. Differentiating the identity defining
$Z_\chi$ $m$ times at the zero gives exactly
\begin{equation}\label{eq:A039-real-leading}
 e^{i\theta_\chi(\gamma)}A_m
            =i^{-m}Z_\chi^{(m)}(\gamma)/m!.
\end{equation}
The leading branch coefficients are the distinct roots of
$c^m=\rho L(\rho+1,\chi)/A_m$.
For a vertical double zero these are $c=\pm ib$, $b>0$,
so $c^2=-b^2<0$. Equation~\eqref{eq:A039-real-leading} for
$m=2$ makes $e^{i\theta_\chi}A_2$ real and nonzero, and hence
$\rho e^{i\theta_\chi}L(\rho+1,\chi)$ real and nonzero.
Since $\arg\rho=\pi/2-\delta_\gamma$, its argument is
$\pi/2+\theta_\chi+\arg L(3/2+i\gamma,\chi)-\delta_\gamma$.
It follows that
\begin{equation}\label{eq:A039-double-probe}
                         w_\rho=0
               \quad\hbox{at every vertical double zero}.
\end{equation}
This is an exact identity, not a small-height error estimate.

\emph{Trace signs of all possible exceptional clusters.}
For $m\ge3$ the regular polygon of leading coefficients has
both positive and negative real parts. Indeed their sum is
zero; if all their real parts had the same weak sign, every
coefficient would lie on the imaginary axis, impossible for
three or more distinct points on one circle. A nonvertical
double pair also has both signs. A vertical double pair has
the expansion from NN-A035
\[
 s_\pm(t)=\rho\pm ib\,t^{1/2}
                    -\frac{b^2D_\rho}{2}t+O(t^{3/2}),
                         \qquad\Re D_\rho>0.
\]
For nonprincipal primitive characters the proved bound is
$\Re D_\rho>2+h_a(\gamma)>0$; for zeta it is
$\Re D_\rho>2-11/48>0$. Thus each vertical double branch
has strictly negative second real coefficient.
The two branches are exchanged by $t^{1/2}\mapsto-t^{1/2}$,
so their sum is holomorphic in $t$ and
\[
             s_+(t)+s_-(t)=2\rho-b^2D_\rho t+O(t^2).
\]
NN-A005's contour trace formula therefore gives
\begin{equation}\label{eq:A039-double-trace}
 \Re\operatorname{Res}_{s=\rho}
               \frac{sL(s+1,\chi)}{L(s,\chi)}
                       =-b^2\Re D_\rho<0.
\end{equation}
For clarity, that formula follows by integrating
$s\,\partial_s L_t/L_t$ around a small fixed circle around
the cluster. Differentiation in $t$ and integration by parts
give $-(2\pi i)^{-1}\int(\partial_tL_t/L_t)\,ds$.
Since $\partial_tL_t|_0=-sL(s+1,\chi)$, this is precisely the
displayed residue. At a simple zero the residue is $v_\rho$.

\emph{A fixed window has a large phase-square mass.}
Use the heights of Lemma~\ref{lem:A039-dirichlet-trace} with
trace-error exponent $e=1/4$. For each large real $T$ choose
\[
 \tau_- =\tau_{\lceil T\rceil},\qquad
 \tau_+ =\tau_{\lfloor BT\rfloor-1}.
\]
They lie inside $[T,BT]$, avoid all zero ordinates, and equal
$T+O(1)$ and $BT+O(1)$ respectively. Write $\mathcal J_T$
for the zero occurrences with $\tau_-<\gamma<\tau_+$.
The positive-height zero count and NN-A037 give
\begin{align}\label{eq:A039-window-mass}
 |\mathcal J_T|&\sim\frac{(B-1)T\log T}{2\pi},\\
 \sum_{\mathcal J_T}w_\rho
       &=\frac12|\mathcal J_T|+o(T\log T).
 \nonumber
\end{align}
Indeed subtract the two global second-moment asymptotics at
$\tau_+$ and $\tau_-$; their two error terms are $o(T\log T)$
because $B$ is fixed. Thus the unit-size endpoint adjustments
lose no fixed fraction of the moment mass. All multiplicities
remain present in both lines of \eqref{eq:A039-window-mass}.

At every simple zero in this window, the Euler product and the
GRH--Lindelof/Cauchy argument in NN-A037 give
\[
 |v_\rho|=\frac{|\rho L(\rho+1,\chi)|}{|L'(\rho,\chi)|}
                                      \gg_{\chi,B}T^{3/4}.
\]
Since $0\le r_\rho\le1$,
\begin{equation}\label{eq:A039-absolute-mass}
 \sum_{\substack{\rho\in\mathcal J_T\\\rho\ {\rm simple}}}
             |\Re v_\rho|
 \gg_{\chi,B}T^{3/4}
       \sum_{\substack{\rho\in\mathcal J_T\\\rho\ {\rm simple}}}
                    r_\rho^2.
\end{equation}
The trace formula at the two good endpoints, on the other hand,
gives
\begin{equation}\label{eq:A039-small-real-trace}
 \Re\sum_{\substack{\rho\ {
m distinct}\\\tau_-<\Im\rho<\tau_+}}
       \operatorname{Res}_{s=\rho}
                \frac{sL(s+1,\chi)}{L(s,\chi)}
                              =O_{\chi,B}(T^{5/4}).
\end{equation}

\emph{Excluding absence of the right witness.}
Suppose there were no branch in the first class in this window.
Every simple zero would have $\Re v_\rho\le0$.
The polygon argument excludes all multiple zeros except
vertical doubles. By \eqref{eq:A039-double-probe}, those
doubles contribute zero to \eqref{eq:A039-window-mass}.
Consequently its full positive mass comes from the simple
zeros, and \eqref{eq:A039-absolute-mass} is bounded below by
a positive constant times $T^{7/4}\log T$.
All simple residues have nonpositive real part, and every
remaining multiple residue has negative real part by
\eqref{eq:A039-double-trace}. The real trace is therefore at
most minus that lower bound. This contradicts
\eqref{eq:A039-small-real-trace} for all sufficiently large $T$.

\emph{Excluding absence of the left witness.}
Every multiple zero has a branch in the second class, by the
polygon argument and the strictly left-bending vertical double
case. Its absence would therefore make all zeros in the window
simple, with $\Re v_\rho\ge0$. The full phase-square mass again
equals the simple mass, so the real trace is now at least a
positive constant times $T^{7/4}\log T$. This contradicts the
same estimate. Both contradictions are uniform over all large
$T$ for the specified fixed $B$, proving the recurrence statement.

\emph{The counting assertion and its precise strength.}
Fix any $B>1$. Apply the result to the disjoint open windows
$(B^jT_B,B^{j+1}T_B)$. Each contains a center of each class, so
\[
 \liminf_{X\to\infty}\frac{R_\chi(X)}{\log X}
       \ge\frac1{\log B},\qquad
 \liminf_{X\to\infty}\frac{L_\chi(X)}{\log X}
       \ge\frac1{\log B}.
\]
The theorem holds for each fixed $B>1$; letting $B$ decrease
to one after this limiting assertion proves
\eqref{eq:A039-superlog}. This does not give a uniform
threshold as $B\downarrow1$, a positive proportion of the
$\asymp X\log X$ zero occurrences, or a ratio of the two counts.
\end{proof}

\paragraph{Geometric and Grothendieck translation.}
For the specified analytic norm shift, a length-$m$ local divisor
normalizes after the finite base change $t=u^m$ into $m$ analytic
sections with the displayed first and second normal jets. The
symmetric sum descends to the unramified parameter $t$, even when
the individual sections do not; its derivative is the full contour
residue. This is the finite locally free trace, which is defined
on the nonreduced algebra before passage to the normalization.
The finite differentials and discriminant dictionary is SGA~1,
Expos\'e~I, section~3, Proposition~1 and section~4, Proposition~10.
The proof combines that trace with a nonnegative analytic phase
probe; the probe vanishes identically on the only kind of multiple cluster
that can lack a right first jet. Thus signed recurrence is proved
without replacing the divisor by its support or assuming a
positive inner product on all its jets. The rank-one Artin
representation attached to $\chi$ is arithmetic, whereas this
specified denominator shift is an analytic deformation; no
arithmetic family or Frobenius distribution for its moving zeros
is asserted.

\paragraph{Inputs and the remaining signed gate.}
The only additional analytic input in the good-height lemma is the
precisely sourced fixed-character GRH--Lindelof bound already used in
NN-A037 (Iwaniec--Kowalski, Corollary~5.20, equation~(5.60),
p.~116). The phase-square theorem, multiple-jet sign and local
trace formula are NN-A037, NN-A035 and NN-A005. In the zeta
case NN-A013 already supplies the needed good-height trace;
the lemma above proves the primitive extension in full.
No simplicity-proportion theorem is required for this recurrence
argument. The unresolved half-label problem still asks for a
counting density, or equivalently sufficiently strong information
about the parity-sensitive odd phase modes and possible
exceptional clusters. Recurrence and the superlogarithmic lower
counts do not settle that problem.


\section{Quantitative windows containing both arithmetic branch directions}
\label{sec:analytic-quantitative-signed-windows}

Keeping the dependence on the rational twist in NN-A037 gives a
rate for its mixed phase mean. This yields a quantitative version
of NN-A039, still without a simplicity hypothesis or a sign-density
claim. Throughout, $\chi$ and its conductor $q$ are fixed.

\begin{lemma}[A uniform range of Hardy twists]
\label{lem:A040-uniform-twist}
Assume GRH for the fixed primitive $L(s,\chi)$. Uniformly for
$2\le M\le T^{1/8}$ and $M^{-1}\le x\le M$, at every real
cutoff $T$ sufficiently large in terms of $\chi$,
\begin{align}\label{eq:A040-uniform-twist}
 A_\chi(T,x)&:=\sum_{0<\gamma\le T}
              e^{2i\theta_\chi(\gamma)}x^{i\gamma}\\
 &=-\frac{i^a\epsilon_\chi^{-1}}{\sqrt{qx}}
       \sum_{n\le qxT/(2\pi)}
                  \chi(n)\Lambda(n)e^{-2\pi in/(qx)}
       +O_\chi(M^2\sqrt T\log^2T).\nonumber
\end{align}
In particular
\begin{equation}\label{eq:A040-uniform-bound}
 |A_\chi(T,x)|\ll_\chi
             T\sqrt x+M^2\sqrt T\log^2T.
\end{equation}
All zero sums count multiplicity.
\end{lemma}

\begin{proof}
We track every twist-dependent part of the already proved contour
in NN-A037. Use its height $T'\in(T,T+1)$ with
$L'/L=O_\chi(\log^2T)$ on the top, independent of $x$, and put
$c=1+1/\log T'$, $X=T'/(2\pi)$, $b=qx$. We may assume
$1<c\le3/2$. On the entire horizontal range
$-3/2\le\sigma\le c$,
\[
                         |x^{\sigma-1/2}|\le M^2.
\]
The top is consequently $O_\chi(M^2\sqrt T\log^2T)$;
the bottom and the left edge are $O_\chi(M^2)$.
The fixed bottom extension on the right is $O_\chi(M)$.
The right-edge Stirling error is
$O_\chi(M\sqrt T\log T)$, because
$b^{c-1/2}\ll_\chi M$ and
$\sum\Lambda(n)n^{-c}\ll\log T$.
The finitely many omitted low zeros have modulus-one summands,
so contribute $O_\chi(1)$.

It remains to make the stationary-phase summation uniform.
Write $r=n/b$ in equation~\eqref{eq:A037-stationary-integral}.
On the portions outside $[r/2,2r]$ the phase derivative has
absolute value at least $\log2$. Its first derivative estimate
is $O(X^{c-1/2})$, uniformly in $r$, and summing its coefficient
gives $O_\chi(M\sqrt T\log T)$. When $r<1$, only
$1/2\le r<1$ can have a stationary-neighborhood intersection
with $[1,X]$. Its length and amplitude are bounded, and its
weighted total is
\[
 b^{c-1/2}\sum_{b/2\le n<b}\Lambda(n)n^{-c}
                                           \ll\sqrt b.
\]
Restoring the missing main terms $n<b$ also costs at most
$b^{-1/2}\sum_{n<b}\Lambda(n)\ll\sqrt b$. Empty ranges cause
no problem when $b$ is small.

For $r\ge1$, the first stationary remainder from Arias de
Reyna's Lemma~2.5, multiplied by the exact coefficient, is
\[
 b^{c-1/2}\Lambda(n)n^{-c}r^{c-1/2}
                              =\Lambda(n)n^{-1/2}.
\]
Its sum for $n\le2bX$ is $O(\sqrt{bX})$ by Chebyshev and
partial summation. The lower-endpoint remainder has this same
bound when $r\ge2$. For $1\le r<2$, its total is
$O(\sqrt b)$ by the preceding bounded-interval argument.
The upper-endpoint remainder is of the same first-remainder
size when $r\le X/2$. In the remaining range
$X/2<r\le2X$, its exact bound is
\[
 C\Lambda(n)\frac{b^{-1/2}\sqrt r}
                       {|X-r|+\sqrt r}
 \le C\Lambda(n)\frac{\sqrt{bX}}
                       {|n-bX|+b\sqrt X}.
\]
Here $b\sqrt X\gg_\chi T^{3/8}\ge1$ and
$\log(2bX)\ll_\chi\log T$. Grouping integer distances from
the real number $bX$ gives, uniformly even at an integer endpoint,
\[
 \sum_{bX/2<n\le2bX}
       \frac1{|n-bX|+b\sqrt X}
 \ll 1+\log\left(1+\frac{bX}{b\sqrt X}\right)
 \ll\log T.
\]
Thus this remainder totals
$O_\chi(\sqrt{bT}\log^2T)$, which is within the asserted
error. The stationary-phase constants are uniform for
$c-1/2\in[1/2,1]$, as in the checked primary lemma.

The coefficient of the main term remains exactly
$-i^a\epsilon_\chi^{-1}b^{-1/2}$, as computed in NN-A037.
Replacing $T'$ by $T$ changes the zero sum by $O_\chi(\log T)$.
The prime sum changes by at most $O(b+1)$ terms, each of size
$O_\chi(\log T)$ before multiplication by $b^{-1/2}$.
Its change is therefore
$O_\chi((\sqrt b+b^{-1/2})\log T)$, again admissible uniformly
on the displayed range. This proves
\eqref{eq:A040-uniform-twist}. Chebyshev gives
$b^{-1/2}\sum_{n\le bT/(2\pi)}\Lambda(n)\ll\sqrt b\,T$,
and proves \eqref{eq:A040-uniform-bound}.
\end{proof}

\begin{theorem}[NN-A040: shrinking relative windows and quantitative sign counts]
\label{thm:A040-quantitative-signs}
Assume A and fix a primitive Dirichlet character $\chi$.
With the phase probe $w_\rho$ of NN-A039,
\begin{equation}\label{eq:A040-probe-error}
             \sum_{0<\gamma\le T}w_\rho
                =\tfrac12N_\chi(T)+O_\chi(T\log\log T).
\end{equation}
There is a constant $C_\chi>0$ such that every sufficiently
large interval
\begin{equation}\label{eq:A040-short-window}
       \left(T,T+C_\chi\frac{T\log\log T}{\log T}\right)
\end{equation}
contains a center of each of the two branch classes in NN-A039.
Consequently their distinct-center counting functions satisfy
\begin{equation}\label{eq:A040-counts}
 R_\chi(X)\gg_\chi\frac{(\log X)^2}{\log\log X},\qquad
 L_\chi(X)\gg_\chi\frac{(\log X)^2}{\log\log X}.
\end{equation}
The constants are allowed to depend on the fixed character and
are not assigned numerical values here. Under eventual simplicity,
these are the two strict ordinary-velocity signs. Neither a positive
proportion of zeros nor equality of the two sign densities follows.
\end{theorem}

\begin{proof}
The absolutely convergent Euler expansion at $3/2$ is
\[
 \frac{L(3/2+it,\chi)}{\overline{L(3/2+it,\chi)}}
   =\sum_{n,m\ge1}
    \frac{\chi(n)\mu(m)\overline{\chi(m)}}{(nm)^{3/2}}
                             e^{it\log(m/n)}.
\]
Truncation to $n,m\le M$ has a uniform error $O(M^{-1/2})$,
because $\sum_{n>M}n^{-3/2}\ll M^{-1/2}$.
Apply \eqref{eq:A040-uniform-bound} at each $x=m/n$.
The weighted main bound is
\[
 T\sum_{n,m\le M}\frac{\sqrt{m/n}}{(nm)^{3/2}}
     =T\left(\sum_{n\le M}n^{-2}\right)
         \left(\sum_{m\le M}m^{-1}\right)
     \ll T\log(2M).
\]
The weighted error is
$O_\chi(M^2\sqrt T\log^2T)$ since the total absolute
coefficient mass is bounded independently of $M$.
Under GRH every summand of $A_\chi$ has modulus one, so the
discarded tail costs $O(N_\chi(T)M^{-1/2})$.
Choose $M=(\log T)^4$ (integer parts do not affect the bounds).
For all sufficiently large $T$ this is within the lemma's range;
$\sqrt T\log^{10}T=o(T)$ and
$N_\chi(T)M^{-1/2}\ll_\chi T/\log T$. It follows that
\[
 \sum_{0<\gamma\le T}e^{2i\theta_\chi(\gamma)}
       \frac{L(3/2+i\gamma,\chi)}
            {\overline{L(3/2+i\gamma,\chi)}}
                            =O_\chi(T\log\log T).
\]
The correction $e^{-2i\delta_\gamma}$ changes this by
$O_\chi(\log^2T)$: use
$|e^{-2i\delta_\gamma}-1|\le2|\delta_\gamma|\le1/\gamma$
and partial summation with $N_\chi(u)\ll_\chi u\log(u+2)$,
discarding a fixed bounded initial range. Taking real parts in
$\cos^2v=(1+\cos2v)/2$ proves
\eqref{eq:A040-probe-error}.

Put $H(T)=C_\chi T\log\log T/\log T$ and choose the good
trace endpoints from NN-A039 inside $[T,T+H(T)]$:
$\tau_- =\tau_{\lceil T\rceil}$ and
$\tau_+ =\tau_{\lfloor T+H(T)\rfloor-1}$.
The full Riemann--von Mangoldt formula for this fixed character,
in the positive-height convention already sourced in NN-A037,
gives
\[
 N_\chi(\tau_+)-N_\chi(\tau_-)
          =\frac{H(T)}{2\pi}\log T
                       +O_\chi(H(T)+H(T)^2/T+\log T).
\]
Subtract \eqref{eq:A040-probe-error} at these endpoints.
The two error terms total at most $K_\chi T\log\log T$
for some fixed $K_\chi$, whereas half of the displayed main
zero count is asymptotic to
$C_\chi T\log\log T/(4\pi)$.
Choosing $C_\chi>8\pi(K_\chi+1)$ ensures
\begin{equation}\label{eq:A040-window-mass}
               \sum_{\tau_-<\gamma<\tau_+}w_\rho
                            \gg_\chi T\log\log T.
\end{equation}
Since $H(T)=o(T)$, the constants bounding the endpoint errors
can be fixed first, independently of the later choice of
$C_\chi$; the required lower threshold may depend on that choice.

The proof of NN-A039 now applies with this moment mass.
If the right class is absent, all multiple clusters are vertical
doubles with probe zero and negative real symmetric trace;
all simple real velocities are nonpositive. If the left class
is absent, every zero is simple with nonnegative real velocity.
In either case \eqref{eq:A040-window-mass}, the lower bound
$|v_\rho|\gg_\chi T^{3/4}$, and $r_\rho\ge r_\rho^2$
force the absolute real trace to be
$\gg_\chi T^{7/4}\log\log T$. The good-endpoint trace bound
is $O_\chi(T^{5/4})$, a contradiction. This proves
\eqref{eq:A040-short-window}, with exactly the same
first/second-jet classes as NN-A039.

To make the counting consequence explicit, consider a large
dyadic block $[2^k,2^{k+1}]$. For every left endpoint in this
block, its required window length is at most
\[
 D_k=2C_\chi\,2^k
            \frac{\log((k+1)\log2)}{k\log2}.
\]
Partition the initial part of the block into
$\lfloor2^k/D_k\rfloor$ disjoint intervals of length $D_k$.
Each contains a center of each class, in its open interior,
by \eqref{eq:A040-short-window}. The number is
$\gg_\chi k/\log k$ for all large $k$.
For $K=\lfloor\log X/\log2\rfloor$, sum the contributions
from $K/2\le k<K$. These blocks lie below $X$ and give
$\gg_\chi K^2/\log K$ distinct centers of each class.
This is \eqref{eq:A040-counts}.
\end{proof}

\paragraph{Geometric translation and precise remaining gate.}
The finite locally free trace and ramified local sections are exactly
those of NN-A039, with the differentials/discriminant dictionary in
SGA~1, Expos\'e~I, section~3, Proposition~1 and section~4,
Proposition~10. The new input quantifies the analytic phase pairing:
absolute Euler truncation reduces it to finitely many characters
of the compact dual of $\mathbb Q_{>0}$, and the proved uniform
stationary-phase estimate controls all of them. The two signed
normal-jet loci must therefore meet the explicit shrinking-relative
height windows. This remains a statement about the specified
analytic divisor, not an arithmetic Frobenius equidistribution or
a positive full-jet Hilbert realization. The $\asymp X\log X$
scale needed for positive counting densities is still far larger
than the proved lower count. The parity-sensitive half-label gate
and any uniformity over varying conductors remain unresolved.


\subsection{A conductor-only main term and square-root height windows}
\label{sec:A040-rational-refinement}

The following strengthens NN-A040 by using the full primitive
Dirichlet GRH content of A for the auxiliary moduli in the rational
twists. Put
\begin{equation}\label{eq:A040-cq}
 c_q=\frac5{4\pi}\prod_{p\mid q}\frac{p(p+1)}{p^2+1},
\end{equation}
with empty product one.

\begin{theorem}[NN-A040 refinement: exact mixed mean and square-root windows]
\label{thm:A040-square-root}
Under A, for every fixed primitive character $\chi$ of conductor
$q$, including $\zeta$ when $q=1$,
\begin{align}\label{eq:A040-exact-mixed}
 \sum_{0<\gamma\le T}e^{2i\theta_\chi(\gamma)}
   \frac{L(3/2+i\gamma,\chi)}
        {\overline{L(3/2+i\gamma,\chi)}}
     &=-c_qT+O_\chi(\sqrt T\log^3T),\\
 \sum_{0<\gamma\le T}w_\rho
     &=\tfrac12N_\chi(T)-\tfrac12c_qT
                       +O_\chi(\sqrt T\log^3T).\nonumber
\end{align}
Consequently, for some constant $C_\chi>0$, every sufficiently
large interval
\begin{equation}\label{eq:A040-square-root-window}
                    (T,T+C_\chi\sqrt T\log^2T)
\end{equation}
contains a center of each of NN-A039's two specified branch
classes. Their distinct-center counts obey
\begin{equation}\label{eq:A040-square-root-counts}
 R_\chi(X),\ L_\chi(X)\gg_\chi\frac{\sqrt X}{\log^2X}.
\end{equation}
This is still not a positive counting density or a half-label law.
No numerical value for the character-dependent window constant
or threshold is asserted.
\end{theorem}

\begin{proof}
\emph{A wider twist estimate retaining its weights.}
For $T^{-1}\le x\le T$, repeat the proof of
Lemma~\ref{lem:A040-uniform-twist} with left side
$\Re s=-1/2$ instead of $-3/2$. There are still no gamma poles
or trivial zeros at positive height. The top logarithmic
derivative has the same $O_\chi(\log^2T)$ bound, on the smaller
strip. On its high horizontal edge the modulus of the other
factor is, by Stirling,
\[
 |\Delta_\chi(s)^{-1}x^{s-1/2}|
                \ll_\chi(qxT/(2\pi))^{\sigma-1/2}
                \ll_\chi\sqrt{xT}.
\]
For the last inequality, $xT\ge1$, and the exponent ranges
from $-1$ to $1/2+1/\log T'$; the base lies between a fixed
positive constant and $O_\chi(T^2)$. The top is therefore
$O_\chi(\sqrt{xT}\log^2T)$. On the left side the functional
equation and absolute convergence at $3/2-it$ give
$L'/L=O_\chi(\log(t+2))$ and
$|\Delta_\chi|^{-1}=O_\chi(t^{-1})$. This edge costs
$O_\chi(x^{-1}\log^2T)$. The fixed bottom costs
$O_\chi(x^{-1}+\sqrt x)$, since
$x^{1/\log T'}$ is bounded for $x\le T$.

Keep $b=qx$, $X=T'/(2\pi)$ and $r=n/b$ only within this
paragraph. The right Stirling error and the nonstationary
portions of its integral are now bounded by
$O_\chi(\sqrt{bT}\log T)$: their exact prefactor
$b^{c-1/2}=\sqrt b\,b^{1/\log T'}$ is at most a fixed
multiple of $\sqrt b$ throughout $q/T\le b\le qT$.
The bounded $r<2$ contributions, including restoration of
the main terms with $r<1$, are $O(\sqrt b)$; the first
remainder sums to $O(\sqrt{bT})$, exactly as in the previous
lemma. At the upper endpoint one must now retain the nearest
integer term instead of assuming $b\sqrt X\ge1$:
\[
 \sum_{bX/2<n\le2bX}\frac1{|n-bX|+b\sqrt X}
     \ll\frac1{b\sqrt X}
                 +1+\log(1+\sqrt X).
\]
This follows by isolating the two nearest integers and comparing
the remaining monotone tails with their integrals. After the
factor $\sqrt{bX}\log(2+bX)$ is inserted, the error is
$O_\chi(b^{-1/2}\log T+\sqrt{bT}\log^2T)$.
This bound also covers $bX$ bounded, including empty ranges.
Replacing $T'$ by $T$ costs
$O_\chi((\sqrt b+b^{-1/2})\log T)$ as before. Thus, uniformly
for the entire range $T^{-1}\le x\le T$,
\begin{equation}\label{eq:A040-weighted-twist}
 A_\chi(T,x)=
 -\frac{i^{a_\chi}\epsilon_\chi^{-1}}{\sqrt{qx}}
 \sum_{n\le qxT/(2\pi)}\chi(n)\Lambda(n)e^{-2\pi in/(qx)}
 +O_\chi\bigl((\sqrt{xT}+x^{-1}+x^{-1/2})\log^2T\bigr).
\end{equation}
Here $a_\chi\in\{0,1\}$ denotes parity, to distinguish it
from the rational denominator below.

\emph{The exact induced Gauss sum.}
Let $x=b/a$ in lowest terms, with $1\le a,b\le T$, and put
$Q=qb$, $Y=qbT/(2\pi a)$. Define
\[
 S_{Q}(a)=\sum_{u\bmod Q\atop(u,Q)=1}
                   \chi(u)e^{-2\pi iau/Q}.
\]
Then
\begin{equation}\label{eq:A040-gauss}
 S_Q(a)=\mu(b)\chi(b)\overline{\chi(-a)}\,\tau(\chi).
\end{equation}
This includes zero values when $(ab,q)>1$ or $b$ is not
squarefree. To verify it, insert
${\bf1}_{(u,b)=1}=\sum_{d\mid(u,b)}\mu(d)$; coprimality
with $q$ is already enforced by $\chi$. Terms with $(d,q)>1$
vanish. In any other term put $u=dv$, $c=b/d$, and write
$v=v_0+qj$ modulo $qc$. The sum over $j$ is zero unless
$c\mid a$. Since $c\mid b$ and $(a,b)=1$, this forces
$c=1$ and $d=b$. The remaining sum is
$\mu(b)\chi(b)\sum_{v_0\bmod q}\chi(v_0)e^{-2\pi ia v_0/q}$,
which is \eqref{eq:A040-gauss}.
The primitive identity used in this last step follows by a
unit substitution when $(a,q)=1$. If $(a,q)>1$, primitivity
provides a unit $v\equiv1\pmod{q/(a,q)}$ with $\chi(v)\ne1$;
otherwise $\chi$ would factor through that proper divisor
(the reduction map on unit groups is onto by the Chinese
remainder theorem). Multiplication by $v$ leaves the additive
character unchanged and multiplies the sum by $\chi(v)$, so
the sum is zero. This proves every case directly. The classical
general statements are Montgomery--Vaughan, \emph{Multiplicative
Number Theory I}, Section~9.2, Theorems~9.7, 9.10 and 9.12,
pp.~287, 289--291, with the same Gauss convention.

\emph{Uniform prime sums at auxiliary moduli.}
The imported GRH estimate is Iwaniec--Kowalski,
\emph{Analytic Number Theory}, Section~5.7, Theorem~5.15,
equation~(5.56), printed p.~114. For a primitive degree-one
character it gives, with an absolute implied constant,
\[
 \sum_{n\le Y}\psi(n)\Lambda(n)
       =\delta_\psi Y+O\bigl(\sqrt Y\log^2(2+QY)\bigr),
\]
when its conductor divides $Q$. Its Ramanujan hypothesis is
automatic for Dirichlet characters, and A supplies their GRH.
For an imprimitive character modulo $Q$, removing prime powers
at the extra primes costs at most $\omega(Q)\log Y$ for $Y\ge2$,
within this error. The bound extends to $Y$ in any fixed compact
positive interval by enlarging the absolute constant. Here
$Y\ge q/(2\pi)$, so all uses are uniform in $a,b,T$ for fixed
$\chi$.

On $(\mathbb Z/Q\mathbb Z)^*$ expand
$f(u)=\chi(u)e^{-2\pi iau/Q}$ in multiplicative characters.
Its Fourier coefficients $\alpha_\psi$ satisfy
\[
 \alpha_{\psi_0}=S_Q(a)/\varphi(Q),\qquad
 \sum_\psi|\alpha_\psi|^2=1,
                     \qquad\sum_\psi|\alpha_\psi|\le\sqrt{\varphi(Q)}.
\]
This is the exact finite-group Parseval identity and
Cauchy--Schwarz. Applying the preceding prime bound yields
\[
 \sum_{n\le Y\atop(n,Q)=1}
      \chi(n)\Lambda(n)e^{-2\pi ian/Q}
 =\frac{S_Q(a)}{\varphi(Q)}Y
           +O\bigl(\sqrt{QY}\log^2(2+QY)\bigr).
\]
The omitted prime powers at $p\mid Q$ cost
$O(\omega(Q)\log(2Y))$. After multiplication by
$(qx)^{-1/2}$ their error is
$O_\chi(\sqrt{a/b}\log^2T)$, while the preceding error is
$O_\chi(\sqrt{bT}\log^2T)$.
Indeed $Q\le qT$, $Y\le qT^2/(2\pi)$ and
$\sqrt{QY}/\sqrt{qb/a}=\sqrt{qbT/(2\pi)}$.

Combining with \eqref{eq:A040-weighted-twist}, the root number
$\epsilon_\chi=\tau(\chi)/(i^{a_\chi}\sqrt q)$, and
$\overline{\chi(-a)}=(-1)^{a_\chi}\overline{\chi(a)}$, gives
\begin{equation}\label{eq:A040-rational-main}
 A_\chi(T,b/a)=
 -\frac{qT}{2\pi\varphi(q)}\sqrt{\frac ba}\,
        \frac{\chi(b)\overline{\chi(a)}\mu(b)}{\varphi(b)}
 +O_\chi\left(
       \left[\sqrt{bT}+\frac ab+\sqrt{\frac ab}\right]\log^2T\right).
\end{equation}
If $(b,q)>1$, both the main coefficient and Gauss sum vanish;
otherwise $\varphi(qb)=\varphi(q)\varphi(b)$. This explains
the displayed formula also in the noncoprime cases.

\emph{Weighted summation and the Euler constant.}
Truncate the Euler ratio to $n,m\le M=\lfloor T\rfloor$,
and write $n=ga$, $m=gb$, $(a,b)=1$. Every twist lies in the
range of \eqref{eq:A040-rational-main}. The three error sums,
after multiplication by $(nm)^{-3/2}$, are bounded by
\begin{align*}
 \sqrt T\sum_{g,a,b\atop ga,gb\le M}g^{-3}a^{-3/2}b^{-1}
       &\ll\sqrt T\log(2M),\\
 \sum_{n,m\le M}n^{-1/2}m^{-5/2}&\ll\sqrt M,\\
 \sum_{n,m\le M}n^{-1}m^{-2}&\ll\log(2M).
\end{align*}
Dropping the coprimality and character restrictions only
increases these absolute bounds. Their total contribution is
$O_\chi(\sqrt T\log^3T)$.
The discarded Euler tail costs
$O(N_\chi(T)/\sqrt M)=O_\chi(\sqrt T\log T)$, since the
zero summands have modulus one under A.

In the main term all character phases cancel. The remaining
absolutely convergent series is
\begin{equation}\label{eq:A040-constant-series}
 \mathcal C_q=\sum_{g,a,b\ge1\atop(a,b)=1}
 \frac{|\chi(gab)|^2\mu(gb)\mu(b)}
                         {g^3a^2b\varphi(b)}.
\end{equation}
Its truncation error on $ga,gb\le M$ is $O(M^{-1/2})$.
To see this rigorously, use
${\bf1}_{\max(ga,gb)>M}\le
M^{-1/2}((ga)^{1/2}+(gb)^{1/2})$.
Both resulting majorants converge: the $g$-sums are bounded by
$\sum g^{-5/2}$, the $a$-sums by $\sum a^{-3/2}$, and the
only larger $b$-sum is
\[
 \sum_b\frac{\mu(b)^2}{\sqrt b\,\varphi(b)}
       =\prod_p\left(1+\frac1{\sqrt p(p-1)}\right)<\infty.
\]
Thus replacement by the full series costs $O_\chi(T/\sqrt M)$.

For a prime $p\nmid q$, the nonzero local choices in
\eqref{eq:A040-constant-series} have either $p\nmid b$, giving
$(1-p^{-3})/(1-p^{-2})$, or $p\mid b$, forcing $p\nmid ga$
and giving $1/(p(p-1))$. Their sum is
\[
 \frac{1-p^{-3}}{1-p^{-2}}+\frac1{p(p-1)}
                          =\frac{p^2+1}{p^2-1}.
\]
Primes dividing $q$ have local factor one. Absolute convergence
therefore gives
\[
 \mathcal C_q=\prod_{p\nmid q}\frac{p^2+1}{p^2-1}
       =\frac{\zeta(2)^2}{\zeta(4)}
                    \prod_{p\mid q}\frac{p^2-1}{p^2+1}
       =\frac52\prod_{p\mid q}\frac{p^2-1}{p^2+1}.
\]
The elementary special values $\zeta(2)=\pi^2/6$ and
$\zeta(4)=\pi^4/90$ are recorded in DLMF~25.6.1,
\url{https://dlmf.nist.gov/25.6.E1}.
Multiplication by $q/(2\pi\varphi(q))$ is precisely $c_q$.
This proves the first asymptotic in
\eqref{eq:A040-exact-mixed}. The exact $\delta_\gamma$
correction costs only $O_\chi(\log^2T)$, as already proved,
and the cosine-square identity proves the second.

\emph{The square-root windows and counts.}
Put $H=C_\chi\sqrt T\log^2T$ and choose the good trace
endpoints inside $[T,T+H]$ as in NN-A039, now with trace
exponent $e=1/8$. Subtract the second asymptotic in
\eqref{eq:A040-exact-mixed} at these endpoints.
The main zero-count contribution is
$H\log T/(4\pi)$; the linear secondary term contributes
$-c_qH/2+O_\chi(1)$, which is smaller by one logarithm.
The two errors total $O_\chi(\sqrt T\log^3T)$.
Choosing $C_\chi$ sufficiently large yields window mass
$\gg_\chi\sqrt T\log^3T$.
If either specified sign class were absent, the exact cluster
argument of NN-A039 puts this entire mass on simple zeros
whose real velocities have one weak sign, with all remaining
multiple trace contributions reinforcing that sign.
The already proved lower bound $|v_\rho|\gg_\chi T^{3/4}$ then
makes the absolute real trace $\gg_\chi T^{5/4}\log^3T$.
The good-height residue formula bounds it by
$O_\chi(T^{9/8})$, a contradiction. This proves
\eqref{eq:A040-square-root-window}.

For large $X$, every required window with left endpoint in
$[X/2,X]$ has length at most $C_\chi\sqrt X\log^2X$.
Partition the initial part of this half-interval into disjoint
intervals of that length. There are
$\gg_\chi\sqrt X/\log^2X$ of them, and each contains a
center of each class in its open interior. This proves
\eqref{eq:A040-square-root-counts} without identifying
centers that happen to support both signs.
\end{proof}

\paragraph{Additional arithmetic and geometric input.}
This refinement uses A for the primitive characters inducing
all characters modulo $qb$ in the finite truncation. It thus
uses a larger explicitly named Dirichlet family than the
single-character GRH sufficient for the coarse mixed bound.
Finite character orthogonality is the actual Fourier decomposition
of the cyclotomic Galois quotient, and its Parseval norm controls
the aggregate prime error. The conductor-only constant comes
from the fully computed local Euler factors above. The analytic
local trace and normal-jet interpretation remain those in NN-A039
and SGA~1, Expos\'e~I, sections~3--4. No Frobenius law for
zero labels, integral model for the denominator shift, or
positive-density signed conclusion is inferred.

The uniform prime estimate was read in the primary-book OCR
at Iwaniec--Kowalski Theorem~5.15, equation~(5.56), p.~114;
no page-image reading is claimed. The Gauss-sum locators are
the author-hosted Montgomery--Vaughan chapter
\url{https://personal.science.psu.edu/rcv4/personal/Publications/MNTI/13.0_pp_282_325_Primitive_characters_and_Gauss_sums.pdf},
Section~9.2, pp.~287--291. Their required specialized identities
were also proved in full above.


\section{Almost linear counts of centers with each arithmetic branch direction}
\label{sec:analytic-signed-cells}

The full residue trace need not bound an individual velocity. It does,
however, bound the real trace of every interval between successive
good heights. Removing the intervals that contain a chosen direction
leaves traces of one weak sign. This gives a stronger count without
an upper moment for inverse derivatives or an assumption of simplicity.

\begin{theorem}[NN-A041: signed counts by removal of good-height cells]
\label{thm:A041-cell-count}
Assume A and fix a primitive Dirichlet character $\chi$, including
conductor one. Use the two classes of distinct positive-ordinate
zero centers defined in NN-A039: the right class supports a branch
with strictly positive first real Puiseux coefficient; the left
class supports a branch with strictly negative first real coefficient,
or a vertical double branch with strictly negative second real
coefficient. Write $R_\chi(I)$ and $L_\chi(I)$ for their respective
counts in an open height interval $I$.
For every fixed $0<\varepsilon<1$, as $T\longrightarrow\infty$,
\begin{equation}\label{eq:A041-dyadic}
 R_\chi((T,2T)),\ L_\chi((T,2T))
             \gg_{\chi,\varepsilon}T^{1-\varepsilon}\log T.
\end{equation}
This deduction uses GRH for the specified function only, through
the fixed-character assertions of NN-A037 and NN-A039.

Using also the all-auxiliary-modulus consequence of A proved in
the refinement of NN-A040, there is a constant $C_\chi>0$ such
that, uniformly for
\[
 C_\chi\sqrt T\log^2T\le H\le T,
\]
one has
\begin{equation}\label{eq:A041-short}
 R_\chi((T,T+H)),\ L_\chi((T,T+H))
             \gg_{\chi,\varepsilon}H T^{-\varepsilon}\log T.
\end{equation}
All counts are of distinct centers. A multiple center can belong
to both classes. Under eventual simplicity the counts instead
describe strictly positive and strictly negative ordinary normal
velocities. Neither estimate asserts a positive limiting proportion
of all zero occurrences, a half-label law, or a common positive
deformation-time interval for all the counted branches.
\end{theorem}

\begin{proof}
We first state precisely the four estimates used and account for
the endpoints. At a zero $\rho=1/2+i\gamma$ put
\[
 w_\rho=\cos^2\bigl(\theta_\chi(\gamma)
       +\arg L(3/2+i\gamma,\chi)-\arctan(1/(2\gamma))\bigr).
\]
It is a number in $[0,1]$ at every zero, repeated with the zero's
multiplicity when summed over occurrences. At a simple zero,
$w_\rho=r_\rho^2$, where
$r_\rho=|\Re v_\rho|/|v_\rho|\in[0,1]$.

\emph{Uniform speed at the simple zeros.}
For every fixed $d>0$,
\begin{equation}\label{eq:A041-speed}
 |L'(1/2+i\gamma,\chi)|\ll_{\chi,d}T^d,
 \qquad |v_\rho|\gg_{\chi,d}T^{1-d}
              \quad(T<\gamma<2T,\ \rho\text{ simple}).
\end{equation}
Here is the required Cauchy deduction for arbitrary $d$, rather
than just the previously displayed exponent $1/4$. The checked
GRH--Lindelof estimate, with exponent $d/4$, and the functional
equation give uniformly on a fixed strip around $\Re s=1/2$
\[
 |L(\sigma+it,\chi)|\ll_{\chi,d}
          (1+|t|)^{\max(1/2-\sigma,0)+d/4}.
\]
Choose $a=\min(d/4,1/8)$. On the circle $|s-\rho|=a$,
its right-hand side is $O_{\chi,d}(T^{a+d/4})$, which is
$O_{\chi,d}(T^{d/2})$. The possible zeta pole is outside this
disc at large height. Cauchy's integral formula therefore gives
$|L'(\rho,\chi)|\ll_{\chi,d}a^{-1}T^{d/2}
\ll_{\chi,d}T^d$. Finally the absolutely convergent Euler
product gives
$|L(3/2+i\gamma,\chi)|\ge\zeta(3/2)^{-1}$, and the exact
formula $v_\rho=\rho L(\rho+1,\chi)/L'(\rho,\chi)$ proves
\eqref{eq:A041-speed}.

\emph{A bounded-length interval has bounded occurrence mass.}
For every interval $J\subset[T-2,2T+2]$ of length at most four,
the local zero count, including multiplicity, gives
\begin{equation}\label{eq:A041-local-count}
 \#\{\rho:\Im\rho\in J\}\ll_\chi\log T.
\end{equation}
The same bound holds for its $w$-mass. The precise classical
locator is Iwaniec--Kowalski, Proposition~5.7(1), equation~(5.27),
pp.102--103; covering a length-four interval by a fixed number of
the source's unit neighborhoods proves the version used here.

\emph{Good-height cell traces.}
Fix $e>0$ and choose the sequence $\tau_n\in(n,n+1)$ in
Lemma~\ref{lem:A039-dirichlet-trace}, with its $k=1$ trace
error exponent $e$. Each $\tau_n$ avoids every zero ordinate.
For an interval $J$ whose endpoints belong to this sequence,
define its full trace by
\[
 V_\chi(J)=\sum_{\substack{\rho\text{ distinct}\\\Im\rho\in J}}
       \operatorname{Res}_{s=\rho}
                \frac{sL(s+1,\chi)}{L(s,\chi)}.
\]
If both endpoints lie in $[T-2,2T+2]$, the difference of the
two global trace formulas has a purely imaginary main term.
Consequently
\begin{equation}\label{eq:A041-trace-bound}
                   |\Re V_\chi(J)|\ll_{\chi,e}T^{1+e}.
\end{equation}
This applies both to a single cell $(\tau_n,\tau_{n+1})$ and
to a union of consecutive cells. There is no estimate here for
an individual residue and no division by a derivative at a
multiple zero.

\emph{The positive probe mass.}
First set $H=T$. NN-A037 and the positive-ordinate zero count give
\[
 \sum_{T<\gamma<T+H}w_\rho\gg_\chi H\log T.
\]
All occurrences are included. Indeed the global mean is
$\frac12N_\chi(X)+o(X\log X)$, and subtracting at $2T$ and
$T$ gives a positive multiple of $T\log T$.
For the shorter-interval assertion, NN-A040's proved refinement is
\[
 \sum_{0<\gamma\le X}w_\rho
 =\tfrac12N_\chi(X)-\tfrac12c_qX
                  +O_\chi(\sqrt X\log^3X),\qquad
 c_q=\frac5{4\pi}\prod_{p\mid q}\frac{p(p+1)}{p^2+1}.
\]
The fixed-conductor zero-count formula gives, uniformly for
$0<H\le T$,
\[
 N_\chi(T+H)-N_\chi(T)
       =\frac{H\log T}{2\pi}+O_\chi(H+\log T).
\]
After choosing $C_\chi$ large enough, subtraction proves
\begin{equation}\label{eq:A041-probe}
 \sum_{T<\gamma<T+H}w_\rho\ge c_\chi H\log T
 \quad\bigl(C_\chi\sqrt T\log^2T\le H\le T\bigr)
\end{equation}
for some fixed $c_\chi>0$ and all sufficiently large $T$.
The possible endpoint occurrences have mass $O_\chi(\log T)$
by \eqref{eq:A041-local-count}, so the open-interval convention
does not change this assertion after reducing $c_\chi$.

We now carry out the same counting argument in either range.
Put
\[
 n_0=\lceil T\rceil,\qquad n_1=\lfloor T+H\rfloor-1,
 \qquad U=(\tau_{n_0},\tau_{n_1}).
\]
The interval $U$ lies strictly inside $(T,T+H)$; each removed
end interval has length less than two. By
\eqref{eq:A041-local-count}, they remove only $O_\chi(\log T)$
occurrences and probe mass. Thus, for one constant $c>0$,
\begin{equation}\label{eq:A041-interior-mass}
                 Q(U):=\sum_{\Im\rho\in U}w_\rho
                                \ge c H\log T.
\end{equation}
The good-height sequence may depend on $e$, but this endpoint
argument and its mass bound are uniform over all its allowed
choices.

Fix one of the two sign classes. Call a cell
$I_n=(\tau_n,\tau_{n+1})$, $n_0\le n<n_1$, bad if it
contains a center of that class, and good otherwise. Let $J$ be
the number of bad cells. Different cells are disjoint and have
no zero on their boundary, so
\begin{equation}\label{eq:A041-cell-to-centers}
           J\le\text{number of centers of the selected class in }U.
\end{equation}
Each cell has length less than two. Equation~\eqref{eq:A041-local-count}
therefore bounds its probe mass by $K\log T$, for a constant
$K=K_\chi$. If $J\ge cH/(2K)$, we already have
$J\gg_\chi H$, and hence the desired lower bound
$\gg_{\chi,\varepsilon}HT^{-\varepsilon}\log T$ for all
sufficiently large $T$.

In the other case the union $G$ of the good cells has
\begin{equation}\label{eq:A041-good-mass}
                         Q(G)\ge\tfrac12cH\log T.
\end{equation}
If the selected class is the right class, NN-A039's complete
cluster classification says that every simple zero of $G$ has
$\Re v_\rho\le0$, and every multiple zero of $G$ is a
vertical double. Every such double has $w_\rho=0$ and full
residue with strictly negative real part. If the selected class
is the left class, $G$ contains no multiple zero and all its
simple zeros have $\Re v_\rho\ge0$. In either case all real
traces in $G$ have the same weak sign, its full probe mass comes
from simple zeros, and \eqref{eq:A041-speed} with $r\ge r^2$
gives
\begin{equation}\label{eq:A041-reinforcing}
 |\Re V_\chi(G)|
 \ge\sum_{\substack{\rho\in G\\\rho\text{ simple}}}|\Re v_\rho|
 \gg_{\chi,d}T^{1-d} Q(G)
 \gg_{\chi,d}H T^{1-d}\log T.
\end{equation}
Here $V_\chi(G)$ denotes the sum of the full traces over its
disjoint cells; the inequality for the right class only becomes
stronger if vertical doubles are present.

Full residues are additive over disjoint contours. Subtracting
the $J$ bad-cell traces from the trace on the whole interval $U$
and using \eqref{eq:A041-trace-bound} yields
\[
 |\Re V_\chi(G)|
       \le|\Re V_\chi(U)|+
                       \sum_{I_n\text{ bad}}|\Re V_\chi(I_n)|
       \ll_{\chi,e}(1+J)T^{1+e}.
\]
Together with \eqref{eq:A041-reinforcing}, this proves
\begin{equation}\label{eq:A041-j-plus-one}
              1+J\gg_{\chi,d,e}HT^{-d-e}\log T.
\end{equation}
Choose $d=e=\varepsilon/2$. When $H=T$, the right-hand
side tends to infinity because $0<\varepsilon<1$, so the
additive one can be absorbed. Equations~\eqref{eq:A041-cell-to-centers}
and \eqref{eq:A041-j-plus-one} prove \eqref{eq:A041-dyadic}.

For the shorter windows, NN-A040 already gives at least one
center of each selected class in $(T,T+H)$, after increasing
$C_\chi$ if necessary. If the right-hand side of
\eqref{eq:A041-j-plus-one} is bounded, this one center gives
\eqref{eq:A041-short}; if it is large, absorb the additive one
as before. More explicitly, if the number of selected centers
is $M$, we have $M\ge1$ and $M+1\ge c'HT^{-\varepsilon}\log T$;
therefore $M\ge(c'/2)HT^{-\varepsilon}\log T$ in both cases.
This proves the short-window estimate uniformly throughout its
stated range. No step assumes a bound for any individual large
residue or a positive proportion of simple zeros.
\end{proof}

\paragraph{Geometric and Grothendieck translation.}
For each analytic Weierstrass cluster of the specified shift,
the derivative of its finite locally free trace is the full
residue used above. It is defined on the length-$m$ algebra,
including its nilpotents, and agrees with the sum of the
ramified branch derivatives after descent. The trace is additive
over disjoint contour cells. Removing $J$ cells that meet one
normal-jet class leaves at most $J+1$ groups of consecutive cells;
their real traces reinforce the opposite weak sign, while the
removed cells are controlled by their boundary traces. The
proof thus counts how many cells must meet each signed jet locus,
without replacing full traces by traces on the reduced support.
The finite-algebra dictionary is SGA~1, Expos\'e~I, section~3,
Proposition~1 and section~4, Proposition~10; the analytic trace
identity itself is proved in NN-A005 and NN-A039. The rank-one
Artin representation defining $\chi$ supplies the arithmetic
$L$-function, but this norm-denominator deformation is analytic.
No algebraic family or Frobenius counting law for the moving
zeros is asserted.

\paragraph{Inputs and remaining quantitative gate.}
The local zero count and GRH--Lindelof inputs are precisely
Iwaniec--Kowalski, \emph{Analytic Number Theory} (2004),
Proposition~5.7(1), equation~(5.27), pp.102--103, and
Corollary~5.20, equation~(5.60), p.116, with their degree-one
hypotheses checked in NN-A037. The positive-height count is
Theorem~5.8's proof, equations~(5.34)--(5.36), pp.104--105.
All further inputs are the proved fixed-character phase mean
NN-A037, the complete jet signs NN-A035, the full good-height
trace and cluster identities NN-A039, and, only for
\eqref{eq:A041-short}, the square-root remainder in NN-A040.
The lower bound is a nearly linear count in height for every
fixed exponent loss. It still permits either class to have
density zero among the $\asymp T\log T$ zero occurrences.
Removing that loss would require stronger uniform trace/speed
control or additional signed phase information; it is not a
consequence of the present cell argument.


\section{Positive proportions from corrected phase crossings}
\label{sec:analytic-signed-positive-density}

The following argument uses the real meromorphic quotient of the
Hardy function by the corrected phase cosine. Its cancellation at
a vertical double zero is essential: counting only simple crossings
of the Hardy function would not give the needed bound.

\begin{lemma}[An interval missing a branch direction cannot be overcrowded]
\label{lem:A042-crossing}
Fix a primitive Dirichlet character $\chi$, allowing conductor one,
and assume GRH for this function. Let $R$ and $L$ be the two
classes of centers in NN-A039. There are constants $T_\chi$ and
$B_\chi$ such that the following holds. If $u\ge T_\chi$,
$0<h\le1$, and neither endpoint is a zero of the Hardy function
or of the phase cosine defined below, then absence of either
one of $R,L$ in $(u,u+h)$ implies
\begin{equation}\label{eq:A042-one-sided-increment}
                  S_\chi(u+h)-S_\chi(u)\le2+B_\chi h.
\end{equation}
All zeros in the counting argument have their actual multiplicities.
\end{lemma}

\begin{proof}
Use the root-normalized Hardy function of NN-A037,
$Z_\chi(t)=e^{i\theta_\chi(t)}L(1/2+it,\chi)\in\mathbb R$,
and set
\begin{align*}
 \alpha_\chi(t)&=\Im\log L(3/2+it,\chi),
   &\delta(t)&=\arctan(1/(2t)),\\
 \psi_\chi(t)&=\theta_\chi(t)+\alpha_\chi(t)-\delta(t),
   &c_\chi(t)&=\cos\psi_\chi(t).
\end{align*}
Here the logarithm is the absolutely convergent Euler logarithm,
so it is single-valued and real analytic in $t$. Its derivative
satisfies
\[
 |\alpha_\chi'(t)|
 =|\Re (L'/L)(3/2+it,\chi)|
 \le\sum_{n\ge2}\frac{\Lambda(n)}{n^{3/2}}<\infty.
\]
Also $\delta'(t)=-2/(4t^2+1)$, and Stirling's formula gives
$\theta_\chi'(t)=\frac12\log(qt/(2\pi))+O_\chi(t^{-1})$.
Consequently $\psi_\chi'(t)>0$ at all sufficiently large heights.
The zeros of $c_\chi$ there are simple. If $C(I)$ denotes their
number in an interval $I=(u,v)$ whose endpoints are not zeros,
\begin{equation}\label{eq:A042-cos-count}
 C(I)\le\frac{\psi_\chi(v)-\psi_\chi(u)}\pi+1.
\end{equation}

Consider the real meromorphic quotient
\[
                         g(t)=\frac{Z_\chi(t)}{c_\chi(t)}.
\]
Both numerator and denominator are real analytic in a neighborhood
of every high real point; thus their common factors can be
canceled in the ordinary local analytic sense. We count only
their real zeros and poles in $I$. Write $N_Z(I)$ for the
multiplicity-counted zeros of $Z_\chi$ there, and $N_g(I),P_g(I)$
for the corresponding orders of the real zeros and poles of $g$.
Additivity of orders at each point gives the exact identity
\begin{equation}\label{eq:A042-divisor-difference}
                  N_Z(I)-C(I)=N_g(I)-P_g(I).
\end{equation}

At a simple zero $\rho=1/2+i\gamma$ with $c_\chi(\gamma)\ne0$,
the exact velocity identity is
\begin{equation}\label{eq:A042-simple-crossing}
 \Re v_\rho=-\frac{E_\rho c_\chi(\gamma)}{Z_\chi'(\gamma)},
 \qquad E_\rho=|\rho|\,|L(3/2+i\gamma,\chi)|>0.
\end{equation}
In fact $L'(\rho,\chi)=-i e^{-i\theta_\chi(\gamma)}
Z_\chi'(\gamma)$ and
$\rho e^{i\theta_\chi(\gamma)}L(\rho+1,\chi)
=i E_\rho e^{i\psi_\chi(\gamma)}$, proving the formula with
its sign. Hence $g'(\gamma)=Z_\chi'(\gamma)/c_\chi(\gamma)$
is positive when $\Re v_\rho<0$, and negative when
$\Re v_\rho>0$. If a simple zero instead has
$c_\chi(\gamma)=0$, the two simple factors cancel completely:
$g$ is holomorphic and nonzero there. These are precisely the
simple tangencies in this argument, and no direction is assigned
to them by cancellation.

Now suppose $I$ contains no right-class center. The regular-polygon
argument in NN-A039 excludes all zeros of order at least three,
and excludes nonvertical doubles. Every remaining multiple zero
is a vertical double. At such a zero, writing
$A_2=L''(\rho,\chi)/2$, the real Hardy identity gives
$e^{i\theta_\chi}A_2=-Z_\chi''(\gamma)/2$. Its leading shift
coefficients satisfy
\[
 c^2=\frac{\rho L(\rho+1,\chi)}{A_2}
      =-\frac{2iE_\rho e^{i\psi_\chi(\gamma)}}{Z_\chi''(\gamma)}
                      <0.
\]
It follows both that $c_\chi(\gamma)=0$ and that
$Z_\chi''(\gamma)\sin\psi_\chi(\gamma)<0$.
After canceling the one denominator factor, $g$ has a simple
zero with
\begin{equation}\label{eq:A042-double-crossing}
 g'(\gamma)=
 \frac{Z_\chi''(\gamma)}{-2\sin\psi_\chi(\gamma)
                                  \psi_\chi'(\gamma)}>0.
\end{equation}
All simple zeros of $Z_\chi$ not canceled by the denominator
have strictly negative velocity, since the right class is absent
and the cosine is nonzero. Thus every real zero of $g$ is simple
and crosses from negative to positive.

If $I$ instead contains no left-class center, NN-A035 and
NN-A039 exclude every multiple zero: every multiple cluster
has a negative first real direction or a strictly negative
second direction at a vertical double. All uncanceled simple
zeros now have strictly positive velocity, so every real zero
of $g$ is simple and crosses from positive to negative.

In either case at most one zero of $g$ lies in each connected
component of $I$ minus its real poles. To see this, two successive
simple zeros of a continuous real nonvanishing function between
them must cross in opposite directions: its sign immediately
after the first equals its sign immediately before the second.
The prescribed common crossing orientation rules this out.
All poles of $g$ are simple because the denominator zeros are
simple; hence there are $P_g(I)+1$ such components and
\[
                       N_g(I)\le P_g(I)+1.
\]
By \eqref{eq:A042-divisor-difference} and
\eqref{eq:A042-cos-count},
\begin{equation}\label{eq:A042-overcrowding}
 N_Z(I)\le C(I)+1
          \le\frac{\psi_\chi(v)-\psi_\chi(u)}\pi+2.
\end{equation}

For completeness the exact counting-phase identity required here
also holds for a complex primitive character. On $\Re s>1/2$,
GRH gives the holomorphic Euler-continuation logarithm of $L$ at
positive heights; the zeta pole at height zero is irrelevant.
At a real height which is not a zero, its boundary argument is
$\pi S_\chi(t)$. Since $Z_\chi(t)$ is real and nonzero,
$\theta_\chi(t)+\pi S_\chi(t)$ is a locally constant integer
multiple of $\pi$. Across a zero of order $m$, the local factor
$(s-\rho)^m$ makes this boundary argument jump by $m\pi$:
on the right of the line, $\arg(\eta+i(t-\gamma))$ increases
from $-\pi/2$ to $\pi/2$ as $\eta\downarrow0$.
Thus for nonzero endpoints
\[
 N_Z((u,v))=\frac{\theta_\chi(v)-\theta_\chi(u)}\pi
                              +S_\chi(v)-S_\chi(u).
\]
Substituting into \eqref{eq:A042-overcrowding}, and using the
bounded derivatives of $\alpha_\chi$ and $\delta$, proves
\eqref{eq:A042-one-sided-increment}. The excluded endpoint
heights form a discrete, hence measure-zero, set, including
after translating them by a fixed $h$.
\end{proof}

\begin{lemma}[From two even moments to a positive upper tail]
\label{lem:A042-positive-tail}
For a fixed primitive character suppose that, for every fixed
$C\ge1$ with $h=C/\log T$, the second and fourth shifted
moments have the following form as $T\longrightarrow\infty$.
The constants in the $O_\chi$ terms must be independent of $C$;
the vanishing errors may depend on both $\chi$ and $C$:
\begin{align}\label{eq:A042-moments}
 \frac1T\int_T^{2T}X_{T,h}(t)^2\,dt
      &=\frac{v}{\pi^2}+O_\chi(1+v^{1/2})+o_{\chi,C}(1),\\
 \frac1T\int_T^{2T}X_{T,h}(t)^4\,dt
      &=\frac{3v^2}{\pi^4}+O_\chi(1+v^{3/2})+o_{\chi,C}(1),\nonumber\\
 X_{T,h}(t)&=S_\chi(t+h)-S_\chi(t),\qquad
                         v=\log(2+C).\nonumber
\end{align}
Suppose also $S_\chi(t)=O_\chi(\log(t+2))$.
Then there is a fixed $C_\chi>0$ such that, with
$h=C_\chi/\log T$, the set
\[
            \{t\in[T,2T]:X_{T,h}(t)>3\}
\]
has measure at least $T/1536$ for all sufficiently large $T$.
The size of $C_\chi$ is not numerically asserted.
\end{lemma}

\begin{proof}
Take $C>0$ fixed, put $h=C/\log T$, and regard normalized
Lebesgue measure on $[T,2T]$ as probability. The number
$v=\log(2+C)$ is fixed while $T\longrightarrow\infty$.
Choose it sufficiently large, using the constants in
\eqref{eq:A042-moments}, that
\[
 av\le\mathbb EX^2\le bv,\qquad
 \mathbb EX^4\le dv^2,\qquad
 a=\frac1{2\pi^2},\quad b=\frac2{\pi^2},\quad d=\frac6{\pi^4}.
\]
The mean itself tends to zero. Indeed an exact change of
integration variables gives
\[
 T\mathbb EX=\int_{2T}^{2T+h}S_\chi(t)\,dt
                         -\int_T^{T+h}S_\chi(t)\,dt
                         =O_\chi(h\log T)=O_\chi(C).
\]
H\"older's inequality in the form
$\mathbb EX^2\le(\mathbb E|X|)^{2/3}
(\mathbb EX^4)^{1/3}$ yields
\[
 \mathbb E|X|\ge\frac{(\mathbb EX^2)^{3/2}}
                         {(\mathbb EX^4)^{1/2}}
             \ge A\sqrt v,\qquad
 A=\frac{a^{3/2}}{\sqrt d}=\frac1{4\sqrt3\,\pi}.
\]
Increase $C$ if necessary so that $3\le A\sqrt v/8$.
For all sufficiently large $T$ we have
$|\mathbb EX|\le A\sqrt v/4$, whence
\[
 \mathbb EX_+=\frac{\mathbb E|X|+\mathbb EX}{2}
                       \ge\frac{3A\sqrt v}{8},
 \qquad
 \mathbb E\bigl(X\mathbf1_{X>3}\bigr)
                       \ge\frac{A\sqrt v}{4}.
\]
Cauchy's inequality now gives
\[
 \mathbb P(X>3)
 \ge\frac{A^2v}{16\mathbb EX^2}
 \ge\frac{A^2}{16b}=\frac1{1536}.
\]
Only two fixed even moments and a mean estimate have been used;
no Gaussian convergence, odd-moment assertion, or equidistribution
modulo an integer is needed.
\end{proof}

\begin{theorem}[NN-A042: positive signed-center proportions]
\label{thm:A042-positive-proportions}
Under A, for every fixed primitive character $\chi$, including
the conductor-one character, there is $b_\chi>0$ such that
\begin{equation}\label{eq:A042-density}
 R_\chi((T,2T))\ge b_\chi T\log T,
 \qquad L_\chi((T,2T))\ge b_\chi T\log T
\end{equation}
for all sufficiently large $T$. These are counts of distinct
centers of the two NN-A039 classes, with all multiple-zero
possibilities retained. Therefore both classes have positive
lower density even when the denominator counts all zero
occurrences with multiplicity.

For $\zeta$, the moment hypothesis follows unconditionally from
Tsang's Theorem~4 in \emph{Some $\Omega$-theorems for the Riemann
zeta-function}, Acta Arithmetica~46 (1986), 369--395, in its
range $0<h<1$; it is reproduced as Theorem~5.1, p.242, of
Trudgian, Acta Arithmetica~148 (2011), 225--256.
For every nonprincipal primitive character, the complete direct
proof in Theorem~\ref{thm:A042-direct-moments} below supplies
exactly the microscopic moment bounds actually required by
\eqref{eq:A042-moments}. Thus every assertion here is a deduction
from A with no additional statistical conjecture or unaudited
historical correction as an input.

If the number of nonsimple centers in $(T,2T)$ is
$o(T\log T)$, both strict ordinary normal-velocity signs also
have positive lower counting density on the simple zeros.
Without that additional assumption the theorem concerns the
specified branch classes, which can overlap at a multiple
center. Equal proportions, a half-label law on simple zeros,
and a limit for either proportion are not asserted.
\end{theorem}

\begin{proof}
Choose $C_\chi$ from Lemma~\ref{lem:A042-positive-tail} and
put $h=C_\chi/\log T$. For all sufficiently large $T$,
$B_\chi h<1$ in Lemma~\ref{lem:A042-crossing}.
Except at the measure-zero endpoint set, every height with
$X_{T,h}(t)>3$ therefore has at least one center of each
class in $(t,t+h)$. Restrict these heights to
$[T,2T-h]$; this removes length at most $h$, and leaves
measure at least $T/3072$ for large $T$.

For the right class, the resulting set is contained in
\[
       \bigcup_{\substack{\gamma\in(T,2T)\\\gamma\text{ in the right class}}}
                         (\gamma-h,\gamma).
\]
Each center covers length at most $h$, regardless of its
multiplicity or the number of right branches. The measure
bound thus gives
\[
 R_\chi((T,2T))\ge\frac{T}{3072h}
                         =\frac{T\log T}{3072C_\chi}.
\]
The identical covering argument applies to the left class.
This proves \eqref{eq:A042-density}, with
$b_\chi=1/(3072C_\chi)$. The standard count
$N_\chi(2T)-N_\chi(T)\sim T\log T/(2\pi)$ proves the
positive lower proportions under the stated multiplicity
convention. If the nonsimple centers have count $o(T\log T)$,
discarding all of them from each class changes neither positive
lower bound; at the remaining simple centers the two classes
are exactly the strict signs of $\Re v_\rho$.
\end{proof}

\begin{corollary}[One lower constant for all fixed primitive characters]
There is an absolute $b_0>0$ such that the two bounds in
\eqref{eq:A042-density} hold with $b_\chi=b_0$ for every fixed
primitive character. The threshold $T_0(\chi)$ may depend on
the character; no assertion uniform in a conductor growing
with $T$ is made.
\end{corollary}

\begin{proof}
The two $O$ constants in Theorem~\ref{thm:A042-direct-moments}
are absolute, independent of both the character and the fixed
$C$. For the single conductor-one case the two fixed Tsang
moment constants are absolute as well. Choose one larger pair
of constants covering these inputs. The choices of $v$ and
$C$ in Lemma~\ref{lem:A042-positive-tail} then depend only on
that pair and the displayed absolute numbers $a,b,d,A$.
Every remaining requirement---the vanishing moment errors,
the vanishing mean, $B_\chi C/\log T<1$, and the high-height
crossing classification---can be satisfied by increasing
$T_0(\chi)$ after this same $C$ has been fixed. The covering
proof therefore gives the single constant $b_0=1/(3072C)$.
\end{proof}

\paragraph{Geometric and Grothendieck translation.}
The real analytic divisor of $g=Z_\chi/\cos\psi_\chi$ is the
difference of the two real analytic Cartier divisors. At a simple
tangency their two local length-one factors cancel. At a vertical
double, the length-two Hardy factor loses exactly one phase
factor and leaves a simple oriented crossing. Its positive
orientation follows from the negative real leading discriminant
coefficient of the original two-branch deformation. The absence
of one normal-jet class makes all remaining crossings have a
common orientation, so the real divisor has at most one more
zero than pole on an interval. This is a real ordered-divisor
statement with an explicitly specified phase; no order or
positivity is assigned to a general complex scheme.
SGA~1, Expos\'e~I, section~3, Proposition~1 and section~4,
Proposition~10 supply the finite differentials and discriminant
dictionary for the local branch algebra. The contour trace and
ramified sections remain those of NN-A005 and NN-A035. The
moment theorem counts analytic height intervals; it is not a
Frobenius equidistribution law and does not construct an
arithmetic family realizing the denominator shift.

\paragraph{Historical source audit, not an input to the general-character proof.}
The displayed crossing and moment-to-density implications are
proved in full above; the required microscopic moments for
nonprincipal characters are proved in full in the following
section from the located interval majorants and explicit formula.
Fujii's \emph{On the zeros of Dirichlet
$L$-functions.~I}, Transactions of the AMS~196 (1974), 225--235,
Main Theorem p.225, states the needed fixed even shifted moments
for primitive characters, including complex characters and the
conductor-one case. Its later correction is
\emph{On the zeros of Dirichlet $L$-functions.~II},
Transactions of the AMS~267 (1981), 33--40,
DOI~10.1090/S0002-9947-1981-0621970-5. The exact correction
has not been directly inspected here. Nguyen's 2021 thesis
\emph{Topics in Multiplicative Number Theory}, Lemma~41,
equation~(8.1), printed p.136, reproduces exactly the required
fixed-character moments for bounded positive $h$, including
complex characters; that statement was read directly, while
its subsequent proof is not imported. These historical statements
are therefore not used to discharge the general-character input.
The zeta input uses the proved $0<h<1$ Tsang formulation,
rather than the false historical
extension of the shift range to $h\asymp T$ noted in
Heath-Brown's errata to Titchmarsh, p.3, entry for book p.246.


\section{A direct microscopic moment input for every fixed character}
\label{sec:A042-direct-moments}

This addendum supplies the two moments needed by NN-A042 without
using the inaccessible correction to the historical shifted-moment
theorem. The order of quantifiers is essential: choose a fixed $C$
first, and then let $T$ tend to infinity. No estimate for a shift
growing with $T$ is claimed here.

\begin{theorem}[NN-A042: the microscopic second and fourth moments]
\label{thm:A042-direct-moments}
Let $\chi$ be a fixed primitive nonprincipal Dirichlet character.
Assume GRH for $L(s,\chi)$ and for $\zeta(s)$, as supplied by input A.
Put $h=C/\log T$, $v=\log(2+C)$, and
$X(t)=S_\chi(t+h)-S_\chi(t)$. For every fixed $C\ge1$,
\begin{align}\label{eq:A042-direct-moment-result}
 \frac1T\int_T^{2T}X(t)^2\,dt
   &=\frac{v}{\pi^2}+O(1+\sqrt v)+o_{\chi,C}(1),\\
 \frac1T\int_T^{2T}X(t)^4\,dt
   &=\frac{3v^2}{\pi^4}+O(1+v^{3/2})+o_{\chi,C}(1).
 \nonumber
\end{align}
The constants in the two $O$ terms are independent of $C$ and
$\chi$; the thresholds for the $o(1)$ terms may depend on both.
Together with the already checked zeta moment input, this supplies
the premise actually used in NN-A042 for every fixed primitive
character, under A alone. The extra RH assumption for zeta is used
only to obtain an elementary prime asymptotic from the already
checked Iwaniec--Kowalski input; the ordinary prime number theorem
would suffice in its place.
\end{theorem}

\begin{proof}
Write $L=\log T$, fix $\eta=1/100$, and set
$N=T^\eta$ and $b=C/(2\pi)$. For functions on $[T,2T]$ use
normalized Lebesgue measure, denoted by $\mathbb E$, and its
$L^r$ norm $\|\cdot\|_r$.

\paragraph{Classical inputs and normalization.}
We import the interval majorants and minorants of
Carneiro--Chandee--Milinovich, \emph{A note on the zeros of zeta
and $L$-functions}, Lemma 2, printed p.~3, and their general
explicit formula (4.6), printed p.~7, with the analytic test
conditions of Lemma 4, printed p.~3. Section 4, pp.~5--6, and
Example 8, p.~13, include complex primitive Dirichlet characters.
The source is the fourteen-page author-hosted version at
\url{https://home.olemiss.edu/~mbmilino/new_S(t).pdf}.
These are imported classical statements, not new theorems of
this catalogue. The finite polynomial moment calculation below
is given in full.

Translate those interval functions to $[0,b]$, with bandwidth
$\eta$, and call them $F_-,F_+$. With the Fourier convention
$\widehat F(\xi)=\int F(x)e^{-2\pi ix\xi}\,dx$, their needed
properties are
\begin{gather*}
 F_-\le\mathbf1_{[0,b]}\le F_+,
 \qquad \int F_\pm=b\pm\eta^{-1},
 \qquad \operatorname{supp}\widehat F_\pm\subset[-\eta,\eta],\\
 \left|\widehat F_\pm(\xi)-\widehat{\mathbf1}_{[0,b]}(\xi)\right|
       \le\eta^{-1},\qquad
 F_\pm(x)=O_{C,\eta}((1+|x|)^{-2})\quad (|x|\longrightarrow\infty).
\end{gather*}
The Fourier error follows directly from the exact $L^1$ error.
The functions are real entire and satisfy the strip decay
required by the explicit formula; the source explains this
consequence of its growth and real-axis decay on pp.~3--4.
Endpoint normalization will not matter, since the excluded
values of $t$ have measure zero.

For a real admissible test function $f$, specializing (4.6) gives
\begin{equation}\label{eq:A042-direct-explicit}
 \sum_\gamma f(\gamma)
 =\frac1{2\pi}\int_{\mathbb R}f(u)G_\chi(u)\,du
  -\frac1\pi\Re\sum_{n\ge2}
       \frac{\chi(n)\Lambda(n)}{\sqrt n}
       \widehat f\left(\frac{\log n}{2\pi}\right),
\end{equation}
where zeros are counted with multiplicity, all $\gamma$ are real
by GRH, and
\[
 G_\chi(u)=\log(q/\pi)+
  \Re\frac{\Gamma'}{\Gamma}\left(\frac{1/2+a+iu}{2}\right),
 \qquad a=\frac{1-\chi(-1)}2\in\{0,1\}.
\]
There is no pole term because $\chi$ is primitive nonprincipal.
The two prime terms in (4.6) combine into a real part using
$\Lambda_{\bar\chi}(n)=\overline{\chi(n)}\Lambda(n)$ and
$\widehat f(-\xi)=\overline{\widehat f(\xi)}$; no self-duality
assumption is being made.

Apply this to $f(u)=F_\pm((u-t)L/(2\pi))$. Its transform is
$(2\pi/L)e^{-2\pi it\xi}\widehat F_\pm(2\pi\xi/L)$.
Consequently the prime term is the finite real polynomial
\[
 D_\pm(t)=-2\Re\sum_{n\le N}
  \frac{\chi(n)\Lambda(n)}{L\sqrt n}
  \widehat F_\pm\left(\frac{\log n}{L}\right)n^{-it}.
\]
Stirling's formula gives
$G_\chi(u)=\log(q|u|/(2\pi))+O(1/|u|)$ for $|u|\ge1$,
and $G_\chi$ is bounded on compact real intervals. Uniformly for
$T\le t\le2T$ it follows that
\begin{equation}\label{eq:A042-direct-arch}
 \frac1{2\pi}\int f(u)G_\chi(u)\,du
 =\frac{\log(qt/(2\pi))}{L}(b\pm\eta^{-1})
                      +o_{\chi,C}(1).
\end{equation}
Here is a direct error check. After $u=t+2\pi x/L$, restrict
first to $|x|\le TL/(4\pi)$, where $u\ge T/2$.
The gamma-density difference is
$O(|x|/(TL)+1/T)$; the integrals of $|F_\pm(x)|$ and of
$|xF_\pm(x)|$ up to this cutoff are respectively
$O_{C,\eta}(1)$ and $O_{C,\eta}(1+\log(TL))$.
On the complement, the tail bound gives an integral of the
absolute density difference at most
$O_{\chi,C,\eta}(\log(TL)/(TL))$.
Indeed the density is bounded by a constant times
$\log(2+T+|x|/L)+\log T$, including where $u$ is near zero.
The extra prefactor $1/L$ makes both errors tend to zero.

The count in $(t,t+h)$ equals
$\{\theta_\chi(t+h)-\theta_\chi(t)\}/\pi+X(t)$ by the
exact phase identity proved in NN-A042. Stirling gives for its
first term $b\log(qt/(2\pi))/L+o_{\chi,C}(1)$.
The pointwise majorant and minorant therefore imply, almost
everywhere and uniformly in $t$,
\begin{equation}\label{eq:A042-direct-sandwich}
 D_-(t)-\eta^{-1}+o_{\chi,C}(1)
       \le X(t)\le
 D_+(t)+\eta^{-1}+o_{\chi,C}(1).
\end{equation}
Replacing the factor $\log(qt/(2\pi))/L$ by $1$ in these
bounded error terms is legitimate for fixed $q$.

\paragraph{An elementary fourth-moment bound.}
We will repeatedly use the following finite calculation.
If $Q(t)=\sum_{n\le N}e_n n^{-it}$, then
\begin{equation}\label{eq:A042-direct-DPbound}
 \mathbb E|Q|^4
 \le \left(\sum_{n\le N}d(n)|e_n|^2\right)^2
       +\frac{2N^2}{T}\left(\sum_{n\le N}|e_n|\right)^4.
\end{equation}
To prove it, square $Q$, writing its coefficients as
$c_m=\sum_{ab=m,\ a,b\le N}e_a e_b$, with $m\le N^2$.
For distinct positive integers $m,n\le N^2$,
$|\log(m/n)|^{-1}\le N^2$, so the absolute normalized
integral of $(m/n)^{-it}$ is at most $2N^2/T$.
The off-diagonal sum is bounded by the second term above.
On the diagonal, Cauchy's inequality and
$d(ab)\le d(a)d(b)$ give
\[
 \sum_m|c_m|^2
 \le\sum_{a,b\le N}d(ab)|e_a|^2|e_b|^2
 \le\left(\sum_{n\le N}d(n)|e_n|^2\right)^2.
\]
This proves the bound without invoking a mean-value theorem.

Put $I=\mathbf1_{[0,b]}$, so that
$\widehat I(\xi)=(1-e^{-2\pi ib\xi})/(2\pi i\xi)$.
Let $D_0$ be defined as $D_\pm$, with $\widehat F_\pm$
replaced by $\widehat I$ and the same cutoff $N$.
For the coefficients of $-(D_\pm-D_0)/2$ before taking real
parts, the Fourier error gives
\[
 |e_{n,\pm}|\le\frac{\Lambda(n)}{\eta L\sqrt n},
 \qquad \sum_{n\le N}|e_{n,\pm}|=O(\sqrt N),
 \qquad \sum_{n\le N}d(n)|e_{n,\pm}|^2=O(1).
\]
The last constant is independent of $C$: primes contribute at
most $2(\eta L)^{-2}\sum_{p\le N}(\log p)^2/p=O(1)$;
prime powers of exponent at least two contribute at most
\[
 (\eta L)^{-2}\sum_p\sum_{r\ge2}
                (r+1)(\log p)^2p^{-r}=O(L^{-2}).
\]
The prime estimate follows by partial summation from
$\vartheta(x)=O(x)$. For definiteness, both that estimate and
the prime asymptotic used below follow from the already checked
Iwaniec--Kowalski, Theorem 5.15, equation (5.56), p.~114,
specialized to zeta under A: $\psi(x)=x+O(\sqrt x\log^2x)$.
Removing higher prime powers gives the same sufficient
asymptotic for $\vartheta(x)$.
Since $N^4/T\to0$, \eqref{eq:A042-direct-DPbound} yields
\[
                         \|D_\pm-D_0\|_4=O(1).
\]
Thus \eqref{eq:A042-direct-sandwich} implies
$\|X-D_0\|_4=O(1)+o_{\chi,C}(1)$, with an absolute $O$
constant (the fixed $\eta$ is suppressed).

For $n=p^r$, the coefficient in $D_0=-2\Re\sum a_n n^{-it}$ is
\[
 a_{p^r}=\frac{\chi(p)^r(1-p^{-irh})}{2\pi i r p^{r/2}}.
\]
The part with $r\ge2$ also has bounded fourth norm:
its divisor-weighted square sum is bounded by
$\pi^{-2}\sum_p\sum_{r\ge2}(r+1)/(r^2p^r)<\infty$,
and its coefficient absolute sum is $O(\sqrt N)$.
We have therefore proved the useful decomposition
\begin{equation}\label{eq:A042-direct-decomposition}
 X=P+Y,\qquad
 P(t)=-2\Re\sum_{p\le N}a_p p^{-it},\qquad
 \|Y\|_4=O(1)+o_{\chi,C}(1),
 \quad a_p=\frac{\chi(p)(1-p^{-ih})}{2\pi i\sqrt p}.
\end{equation}
Every bound for $Y$ is independent of $C$ before the final
vanishing term; this is the point of fixing the bandwidth.

\paragraph{The two prime moments.}
Set $V_T=\sum_{p\le N}|a_p|^2$.
On expanding $P^2$ and $P^4$, unique factorization gives
\begin{equation}\label{eq:A042-direct-prime-diag}
 \mathbb EP^2=2V_T+o(1),\qquad
 \mathbb EP^4=12V_T^2-6\sum_{p\le N}|a_p|^4+o(1).
\end{equation}
For clarity, if $A(t)=\sum a_p p^{-it}$, only the
$2A\bar A$ term has a diagonal in $(A+\bar A)^2$.
In the fourth power, only $6A^2\bar A^2$ has a diagonal,
whose coefficient sum is $2V_T^2-\sum|a_p|^4$.
Other splits have different numbers of prime factors.
All off-diagonal products have two positive integer arguments
at most $N^4$, so their normalized integrals are at most
$2N^4/T$. The sum of their absolute coefficients is
$O((\sum_{p\le N}p^{-1/2})^4)=O(N^2)$.
Thus the total fourth error is $O(N^6/T)=o(1)$, and the
second error is smaller. Moreover
$\sum|a_p|^4\le\pi^{-4}\sum_p p^{-2}=O(1)$.

The variance satisfies
\begin{equation}\label{eq:A042-direct-variance}
 2V_T=\frac1{\pi^2}
   \sum_{\substack{p\le T^\eta\\p\nmid q}}
       \frac{1-\cos(C\log p/L)}p
 \longrightarrow\frac1{\pi^2}
      J_\eta(C),\qquad
 J_\eta(C)=\int_0^\eta\frac{1-\cos(Cu)}u\,du.
\end{equation}
Here is justification at the lower integration endpoint.
For fixed $0<\epsilon<\eta$, partial summation using
$\vartheta(x)\sim x$ gives convergence on
$T^\epsilon<p\le T^\eta$ to the integral over
$[\epsilon,\eta]$. The contribution below $T^\epsilon$
is at most
$C^2(2L^2)^{-1}\sum_{p\le T^\epsilon}(\log p)^2/p
=O(C^2\epsilon^2)$; the omitted integral has the same bound.
Let $T\to\infty$ and then $\epsilon\downarrow0$.
The finitely many primes dividing $q$ contribute $o_{\chi,C}(1)$.

Changing variables $z=Cu$ and integrating $\cos z/z$ by
parts on $[1,\eta C]$ shows, for $C\ge\eta^{-1}$,
$J_\eta(C)=\log C+O_\eta(1)$.
For $1\le C\le\eta^{-1}$ the same assertion with
$\log(2+C)$ follows by boundedness on that fixed compact
range. Hence, for all $C\ge1$,
\[
                         J_\eta(C)=v+O(1).
\]
Equations \eqref{eq:A042-direct-prime-diag}--
\eqref{eq:A042-direct-variance} now give
\[
 \mathbb EP^2=v/\pi^2+O(1)+o_{\chi,C}(1),\qquad
 \mathbb EP^4=3v^2/\pi^4+O(1+v)+o_{\chi,C}(1).
\]
Finally use \eqref{eq:A042-direct-decomposition} and H\"older.
The second-moment change is at most
$2\|P\|_2\|Y\|_2+\|Y\|_2^2=O(1+\sqrt v)+o_{\chi,C}(1)$.
Expansion of $(P+Y)^4$ bounds the fourth-moment change by
\[
 4\|P\|_4^3\|Y\|_4+6\|P\|_4^2\|Y\|_4^2
       +4\|P\|_4\|Y\|_4^3+\|Y\|_4^4
       =O(1+v^{3/2})+o_{\chi,C}(1).
\]
This proves the theorem with exactly the stated quantifiers.
\end{proof}

\paragraph{Use in NN-A042.}
Its positive-tail proof needs only the two estimates above for
$h=C/\log T$, with the $O$ constants independent of the fixed
$C$. One chooses $C$ sufficiently large and then $T$ sufficiently
large. Thus its stronger, uniform bounded-shift premise can be
replaced by this microscopic premise. The count and crossing
proofs, the bound $T/1536$ for the upper-tail measure, and the
distinct-center lower bound $T\log T/(3072C)$ are unchanged.
The general-character conclusion now uses A and the explicitly
located classical interval/explicit-formula inputs only; it does
not depend on having inspected Fujii's correction. This closes
the source premise for branch-center positive proportions, not
the ordinary-simple-sign or half-balance questions.


\section{A signed statistic supported on simple zeros}
\label{sec:A042-simple-signed-cancellation}

This addendum to NN-A042 concerns the zeta function and uses A
for the auxiliary Dirichlet family in the explicitly cited
first-moment theorem. It removes multiple-center residues from
a signed cancellation estimate, while retaining the separate
absolute-mass gate required for ordinary sign counting.

\begin{proposition}[NN-A042: simple-zero weighted cancellation]
Assume A. For every fixed $\varepsilon>0$,
\begin{equation}\label{eq:A042-simple-cancellation}
 \sum_{T<\gamma\le2T}
 \frac{\Re\{\rho\zeta(\rho+1)\overline{\zeta'(\rho)}\}}{\gamma}
       =O_\varepsilon(T^{3/4+\varepsilon}),
 \qquad \rho=\tfrac12+i\gamma.
\end{equation}
The sum can count every zero occurrence with its multiplicity:
each summand at a multiple center is exactly zero. At every
simple zero its summand is
\[
                  \frac{|\zeta'(\rho)|^2}{\gamma}\Re v_\rho,
 \qquad v_\rho=\frac{\rho\zeta(\rho+1)}{\zeta'(\rho)}.
\]
It therefore has exactly the ordinary normal-velocity sign
whenever that sign is nonzero. No lower bound for the absolute
mass of these summands, or for the count of either sign, is
part of this proposition.
\end{proposition}

\begin{proof}
We use Benli--Elma--Ng, \emph{A discrete mean value of the
Riemann zeta function}, arXiv:2311.13554v1, Theorem~1.1,
equation~(9), PDF p.~2, with the exact length condition~(4)
on PDF p.~1, support and submultiplicativity conditions~(6)--(7)
on PDF p.~2, and the function~(10) on PDF p.~3.
Its conditional error in part~(ii) is uniform on
$|\alpha|\le1/(15\log T)$ under its
$\operatorname{GRH}(\Theta)$. Input A supplies
$\Theta=1/2$. This is an imported mean-value theorem; the
specialization, truncation, and signed-statistic deduction are
proved below.

Fix $0<\delta<1/4$, put
$M=\lfloor T^{1/2-\delta}\rfloor$, and use the two
Dirichlet polynomials
\[
 X(s)=1,\qquad
 Y_M(s)=\sum_{m\le M}\frac{m^{-s}}m.
\]
Their coefficient sequences are
$x(m)=\mathbf1_{m=1}$ and
$y(m)=m^{-1}\mathbf1_{m\le M}$. Both satisfy the source's
submultiplicativity bound with constant one. Indeed, if
$ab\le M$ then $a,b\le M$ and $y(ab)=y(a)y(b)$;
if $ab>M$ the left side is zero. The assertion for $x$
is immediate. The length exponent $1/2-\delta$ is fixed
and strictly below $1/2$. Thus the dependence of these
polynomials on $T$ is exactly within the stated hypotheses.

For completeness the source's main term here is
\begin{align*}
 \mathcal M_M(\alpha,T)
  ={}&\frac{T}{2\pi}\log\frac{T}{2\pi e}
                    \sum_{m\le M}m^{-2-\alpha}\\
   &-\frac{T}{2\pi}\sum_{m\le M}
              \frac{(\Lambda*s_{-\alpha})(m)}{m^2}
      +\sum_{h\le M}\frac{F_{\alpha,h,1}(T)}{h^2},
 \qquad s_{-\alpha}(m)=m^{-\alpha},
\end{align*}
where equation~(10) specializes to
\[
 F_{\alpha,h,1}(T)=\frac{T}{2\pi}
  \left(
   h^{-\alpha}\frac{\zeta'}{\zeta}(1+\alpha)
   -\zeta(1-\alpha)\frac{(T/(2\pi))^{-\alpha}}{1-\alpha}
  \right).
\]
The two polar parts at $\alpha=0$ cancel. Thus this expression
is holomorphic at zero and real for real $\alpha$ near zero.
In particular $\mathcal M_M'(0,T)$ is real.

Here are all the coefficient norms in the conditional error,
so no growing-length constant is hidden. With $\Theta=1/2$,
\begin{gather*}
 \|y(m)m^{-1/2}\|_1\le\zeta(3/2),\qquad
 \|m^{1/2}x(m)(1*|y|)(m)\|_1=1,\\
 \|x(m)y(m)/m\|_1
 =\|x(m)/m\|_1
 =\|x(m)/m^{3/2}\|_1
 =\|x\|_1=1,\\
 \|y(m)/m\|_1\le\zeta(2),\qquad
 \|y\|_1\le1+\log M.
\end{gather*}
Theorem~1.1(ii) therefore gives, uniformly on the stated
shift disc and for every fixed $\eta>0$,
\[
 \sum_{0<\gamma\le T}\zeta(\rho+\alpha)Y_M(1-\rho)
       =\mathcal M_M(\alpha,T)
          +O_{\delta,\eta}
                  \bigl(T^{1/2+\eta}(1+\log M)\bigr).
\]
Cauchy's formula on the circle
$|\alpha|=1/(30\log T)$ applies to the holomorphic error.
It follows that
\begin{equation}\label{eq:A042-imaginary-truncated}
 \Im\sum_{0<\gamma\le T}\zeta'(\rho)Y_M(1-\rho)
      =O_{\delta,\eta}(T^{1/2+\eta}\log^2T).
\end{equation}
All zero occurrences are included by the source contour
formula. Differentiating introduces no simplicity assumption.

Under RH, $1-\rho=\overline\rho$, and the exterior series is
absolutely convergent. In particular
\[
 \bigl|\zeta(2-\rho)-Y_M(1-\rho)\bigr|
       \le\sum_{m>M}m^{-3/2}\ll M^{-1/2}.
\]
The arbitrary-exponent derivative estimate already proved
in NN-A041 gives
$|\zeta'(\tfrac12+i\gamma)|\ll_\eta(\gamma+2)^\eta$.
Together with $N(T)\ll T\log T$, this bounds the omitted sum by
\[
              O_\eta(M^{-1/2}T^{1+\eta}\log T).
\]
Consequently, writing the actual untruncated sum as
\[
 R(T)=\sum_{0<\gamma\le T}\zeta'(\rho)\zeta(2-\rho),
\]
we have
\[
 \Im R(T)\ll_{\delta,\eta}
 T^{1/2+\eta}\log^2T+
 T^{3/4+\delta/2+\eta}\log T.
\]
For any prescribed $\varepsilon>0$, choose sufficiently small
fixed $\delta,\eta>0$ and absorb the logarithms to obtain
\[
                  \Im R(T)=O_\varepsilon(T^{3/4+\varepsilon}).
\]
The same conclusion at $2T$ and $T$ gives the identical
bound for the dyadic difference. This subtraction concerns
the untruncated sums, so their independently chosen cutoffs
cause no additional term.

Set $A_\rho=\zeta'(\rho)\zeta(2-\rho)=a_\rho+i b_\rho$.
Since $\zeta(2-\rho)=\overline{\zeta(\rho+1)}$,
direct multiplication yields
\[
 \frac{\Re\{\rho\zeta(\rho+1)\overline{\zeta'(\rho)}\}}{\gamma}
                       =b_\rho+\frac{a_\rho}{2\gamma}.
\]
The sum of the first terms is the dyadic imaginary part
just bounded. Absolute convergence gives
$|\zeta(2-\rho)|\le\zeta(3/2)$, so the second terms have
total at most
\[
 \frac{\zeta(3/2)}{2T}
       \sum_{T<\gamma\le2T}|\zeta'(\rho)|
             \ll_\eta T^\eta\log T.
\]
This proves \eqref{eq:A042-simple-cancellation}.
Finally, at multiplicity at least two, $\zeta'(\rho)=0$.
At a simple zero, multiplication of the exact velocity
formula by $|\zeta'(\rho)|^2/\gamma$ gives the claimed
identity and sign.
\end{proof}

\paragraph{The remaining count gate.}
Let $K_\rho$ denote the summand and put
$A_T=\sum|K_\rho|$, $B_T=\sum K_\rho$,
$C_T=\sum K_\rho^2$. If $C_T>0$, Cauchy's inequality gives
the exact bounds
\[
 \#\{K_\rho>0\}\ge\frac{(A_T+B_T)^2}{4C_T},
 \qquad
 \#\{K_\rho<0\}\ge\frac{(A_T-B_T)^2}{4C_T}.
\]
Indeed the positive and negative absolute masses are respectively
$(A_T+B_T)/2$ and $(A_T-B_T)/2$, and their squared sums
are at most $C_T$. Thus a common scale $a_T$ with
$A_T\gg N(T,2T)a_T$, $C_T\ll N(T,2T)a_T^2$, and
$|B_T|\le(1-\kappa)A_T$ for fixed $\kappa>0$ would prove
positive proportions of both strict ordinary signs. The
proposition controls $B_T$, but supplies neither required
absolute-mass estimate. In particular, A042's interval
overcrowding can still be carried by multiple centers.
No new sign-count theorem is registered from this criterion.
The subsequent NN-A045 proves the lower absolute mass at
the \(T^{1-\varepsilon}\log T\) scale on an isolated
phase-margin subset. Combined with the signed estimate here
and the GRH supremum bound, it gives almost linear counts of
both strict ordinary signs without a nonsimple-density
assumption. Positive proportions still require stronger,
scale-compatible absolute-mass control.
NN-A046 further proves the signed mean with error
\(O_{\chi,\varepsilon}(T^{1/2+\varepsilon})\) for every
fixed primitive character and extends the almost linear
ordinary sign counts to that scope. Its direct contour
proof requires RH alone for the conductor-one case.

\paragraph{Geometric meaning and limitations.}
The numerator, evaluated at the underlying complex point,
vanishes at each nonreduced zero center and
restricts on the reduced locus to the normal displacement
paired with the positive squared norm of the derivative
section. This norm uses complex conjugation and the specified
real critical line; it is not a regular algebraic function
on an arbitrary complex scheme. The estimate is a real
analytic weighted statistic of the specified divisor motion.
It does not construct a positive arithmetic measure on its
normal-jet classes or identify a Frobenius counting law.

\paragraph{Source locator.}
\url{https://arxiv.org/pdf/2311.13554}, v1:
Theorem~1.1(ii), equation~(9), PDF p.~2; the exact main-term
factor is equation~(10), PDF p.~3. The differentiated
statement is also Corollary~1.2, with its $m=1$ version
equation~(27), PDF p.~6, but the necessary Cauchy
differentiation is proved explicitly above.
The published primary version is \emph{Mathematika}~70 (2024),
e12281, Theorem~1.1(ii) and equations~(9)--(10), p.~3,
with the hypotheses on p.~2:
\url{https://londmathsoc.onlinelibrary.wiley.com/doi/pdf/10.1112/mtk.12281}.
The exact theorem interface was checked directly in both
versions; no complete review of the external paper's
forty-seven-page proof is asserted.


\section{A positive square-mollifier interface for ordinary sign proportions}
\label{sec:A042-square-mollifier-gate}

This is a gate addendum to NN-A042, not an assertion of a new
positive-proportion theorem.  The signed mean required by a positive
square-mollifier weight is automatic under the second-moment
hypothesis below.  The remaining input is one explicit strict
moment inequality.  No mollifier satisfying that inequality is
asserted here.

Put \(I_T=(T,2T]\), \(N_T=N(2T)-N(T)\), and \(\ell=\log T\).
Every sum in this section counts zero occurrences with multiplicity.
For a zero \(\rho=\frac12+i\gamma\), put
\[
 K_\rho=\frac{\Re\{\rho\zeta(\rho+1)
                              \overline{\zeta'(\rho)}\}}{\gamma}.
\]
Thus \(K_\rho=0\) at every multiple center, and its sign at a
simple zero is the ordinary normal-velocity sign.

\begin{lemma}[Removing a technical submultiplicativity restriction]
\label{lem:A042-bilinear-majorant}
The unconditional conclusion of Benli--Elma--Ng, Theorem~1.1(i),
extends by bilinearity to real coefficient sequences \(x_T,y_T\)
supported in \([1,T^\vartheta]\), for a fixed
\(\vartheta<1/2\), satisfying uniform bounds
\[
 |x_T(n)|+|y_T(n)|
       \le C\tau_r(n)(\log(2n))^j
\]
with fixed \(C,r,j\), where \(r\) is a positive integer and
\(j\) a nonnegative integer.  Its main term remains real for real
shifts, and is holomorphic at zero.  In particular, for every
fixed \(A>0\),
\[
 \Im\sum_{0<\gamma\le U}
        \zeta'(\rho)X_T(\rho)Y_T(1-\rho)
       =O_A\bigl(U(\log U)^{-A}\bigr)
       \qquad(U=T,\ 2T).
\]
The same two polynomials are used at both upper limits.
\end{lemma}

\begin{proof}
Here and below coefficient estimates at \(n=1\) are understood
separately as bounded constants; replacing \(\log n\) in the
source's divisor bound by \(\log(2n)\) changes nothing for
\(n\ge2\).  Choose a positive majorant
\[
 H(n)=C_0\tau_r(n)(\log(2n))^j
                  \mathbf1_{n\le T^\vartheta}
\]
with \(C_0\) fixed and sufficiently large that
\(|x_T|,|y_T|\le H\) and \(H(mn)\le H(m)H(n)\).
For the latter assertion, the integer divisor function satisfies
\(\tau_r(mn)\le\tau_r(m)\tau_r(n)\).  At each prime this follows
by decomposing any \(r\)-tuple of nonnegative integers of sum
\(a+b\) into two such tuples of sums \(a\) and \(b\);
the addition map from the two sets of tuples is surjective.
Furthermore
\[
 \log(2mn)\le
                  \frac{\log(2m)\log(2n)}{\log2}.
\]
Taking \(C_0\ge(\log2)^{-j}\) proves the majorant inequality;
the cutoff preserves it because \(mn\) inside the support
forces both \(m,n\) inside.

The two positive sequences \(2H+x_T\) and \(H\) satisfy the
source's submultiplicativity condition with a uniform constant:
\[
 H\le2H+x_T\le3H,\qquad
 (2H+x_T)(mn)\le
           3(2H+x_T)(m)(2H+x_T)(n).
\]
The same holds for \(2H+y_T\).  Write
\(x_T=(2H+x_T)-2H\), and similarly for \(y_T\), and expand
the bilinear mean into four terms.  Theorem~1.1(i) applies to
each, with the same fixed support exponent, divisor bounds,
and submultiplicativity constants.  Its main term is itself
bilinear in \(x_T,y_T\), so the four main terms recombine
to the stated one.

For real shifts the main term in equation~(9) of that theorem
is real.  Its apparent poles at zero cancel: in equation~(10),
for \(k=1\) the poles of \(\zeta'/\zeta(1+\alpha)\) and
\(\zeta(1-\alpha)\) cancel; for \(k>1\) the factor
\(\Phi(\alpha,k)=\prod_{p\mid k}(1-p^{-\alpha})\) cancels
the latter pole.  All other factors are holomorphic on a
sufficiently small disc about zero.  Cauchy's formula on
\(|\alpha|=1/(30\log U)\) differentiates the uniform
holomorphic error and costs one factor \(O(\log U)\).
The arbitrary power of logarithmic saving in part~(i)
absorbs that factor.  The differentiated main term is real.
The length bound holds for \(U=T\) and \(U=2T\), so the
last assertion follows without changing the polynomials.
\end{proof}

\begin{proposition}[NN-A042: the square-mollifier reduction]
\label{prop:A042-square-mollifier}
Assume RH.  Let
\[
 B_T(s)=\sum_{n\le T^\nu}b_T(n)n^{-s},
 \qquad 0<\nu<1/2,
\]
have real coefficients with a uniform fixed divisor bound
\(|b_T(n)|\le C\tau_r(n)(\log(2n))^j\).
Suppose first that
\begin{equation}\label{eq:A042-square-second-bound}
 \sum_{\gamma\in I_T}|\zeta'(\rho)B_T(\rho)^2|^2
                                      =O(N_T\ell^2).
\end{equation}
Then the positive weight \(w_\rho=|B_T(\rho)|^2\) satisfies
\begin{equation}\label{eq:A042-square-signed-mean}
                 \sum_{\gamma\in I_T}K_\rho w_\rho
                                      =o(N_T\ell).
\end{equation}

In addition, suppose that for constants \(a,b>0\),
\begin{align}
 \liminf_{T\to\infty}
 \frac{\left|\sum_{\gamma\in I_T}
                     \zeta'(\rho)B_T(\rho)^2\right|}
                                      {N_T\ell}&\ge a,
 \label{eq:A042-square-first}\\
 \limsup_{T\to\infty}
 \frac{\sum_{\gamma\in I_T}
                  |\zeta'(\rho)B_T(\rho)^2|^2}
                                      {N_T\ell^2}&\le b,
 \qquad a^2>b/2 .
 \label{eq:A042-square-ratio}
\end{align}
Then each strict ordinary normal-velocity sign has a positive
lower proportion among all zero occurrences in \(I_T\);
every counted center is simple.  More explicitly, choose
\(0<\eta<1\) such that
\[
 p_\eta=\frac1{2(1-\eta^2)}<\frac{a^2}{b},
 \qquad d_\eta=a-\sqrt{b p_\eta}>0.
\]
For each sign the lower proportion is at least
\begin{equation}\label{eq:A042-square-density}
                   \frac{\eta^2d_\eta^2}
                        {4\zeta(3/2)^4 b}>0.
\end{equation}
The assertion is conditional on
\eqref{eq:A042-square-first}--\eqref{eq:A042-square-ratio};
these moment inequalities are not proved by this proposition.
\end{proposition}

\begin{proof}
Choose a fixed \(\delta>0\) with \(\nu+\delta<1/2\), and put
\(M=\lfloor T^\delta\rfloor\).  Apply
Lemma~\ref{lem:A042-bilinear-majorant} to
\[
 X_T(s)=B_T(s),\qquad
 Y_T(s)=B_T(s)\sum_{k\le M}k^{-1-s}.
\]
Both polynomials have length at most \(T^{\nu+\delta}\).
The coefficients of \(Y_T\) are bounded by
\[
 C(\log(2n))^j\sum_{d\mid n}\tau_r(d)
            =C\tau_{r+1}(n)(\log(2n))^j,
\]
independently of the two truncations.  Thus the lemma and
subtraction at the two upper limits give a negligible
imaginary part for the truncated sum.

Under RH and real coefficients,
\(B_T(1-\rho)=\overline{B_T(\rho)}\).
Absolute convergence of the exterior zeta series gives,
uniformly at every zero in \(I_T\),
\[
 \left|\zeta(2-\rho)
                   -\sum_{k\le M}k^{\rho-2}\right|
                                  \le2M^{-1/2}.
\]
The error in restoring it is at most
\[
 2M^{-1/2}\sum_{\gamma\in I_T}
                        |\zeta'(\rho)|\,|B_T(\rho)|^2
 \le2M^{-1/2}
       \left(N_T\sum_{\gamma\in I_T}
                  |\zeta'(\rho)B_T(\rho)^2|^2\right)^{1/2}
                         =o(N_T\ell).
\]
Consequently
\[
 \Im\sum_{\gamma\in I_T}
       \zeta'(\rho)\zeta(2-\rho)|B_T(\rho)|^2
                                         =o(N_T\ell).
\]
Writing the product before the weight as \(u_\rho+i v_\rho\)
gives exactly \(K_\rho=v_\rho+u_\rho/(2\gamma)\).
The extra term is
\[
 O\left(T^{-1}\zeta(3/2)
          \sum_{\gamma\in I_T}
                       |\zeta'(\rho)|\,|B_T(\rho)|^2\right)
                           =O(N_T\ell/T).
\]
This proves \eqref{eq:A042-square-signed-mean}.

Use the corrected phase
\[
 \psi(t)=\theta(t)+\arg\zeta(3/2+it)
                                     -\arctan(1/(2t)).
\]
The conductor-one proof of NN-A037 gives, under RH alone,
\[
          \sum_{\gamma\in I_T}\cos^2\psi(\gamma)
                                    =(1/2+o(1))N_T.
\]
It follows that the set \(E_T=\{\gamma\in I_T:
|\cos\psi(\gamma)|<\eta\}\) has size at most
\((p_\eta+o(1))N_T\).  The counting convention retains
multiplicity at every nonsimple center, and uses no independence between the phase
and the mollifier.  Put \(A_\rho=\zeta'(\rho)B_T(\rho)^2\).
Cauchy's inequality and the moment assumptions give
\[
 \sum_{\gamma\in I_T\setminus E_T}|A_\rho|
 \ge\left|\sum_{\gamma\in I_T}A_\rho\right|
                   -\sqrt{\#E_T\sum_{\gamma\in I_T}|A_\rho|^2}
                    \ge(d_\eta-o(1))N_T\ell.
\]
For simple zeros the exact phase identity is
\[
 |K_\rho|w_\rho
  =\frac{|\rho|}{\gamma}|\zeta(\rho+1)|
                         |A_\rho|\,|\cos\psi(\gamma)|.
\]
At multiple centers both sides are zero.  The Euler product
bounds \(1/\zeta(3/2)\le|\zeta(3/2+it)|\le\zeta(3/2)\)
and \(|\rho|/\gamma=1+o(1)\), uniformly on \(I_T\), therefore
yield
\[
 \sum_{\gamma\in I_T}|K_\rho|w_\rho
       \ge\left(\frac{\eta d_\eta}{\zeta(3/2)}-o(1)\right)
                                                    N_T\ell,
 \qquad
 \sum_{\gamma\in I_T}(K_\rho w_\rho)^2
       \le(\zeta(3/2)^2 b+o(1))N_T\ell^2.
\]
The signed mean is \(o(N_T\ell)\), so each of the positive
and negative absolute masses is at least half of the first
displayed lower bound, up to \(o(N_T\ell)\).
Cauchy's inequality on each sign set now proves
\eqref{eq:A042-square-density}.  A nonzero weighted summand
has \(B_T(\rho)\ne0\), \(\zeta'(\rho)\ne0\), and the same
strict sign as \(K_\rho\).  Thus these are distinct simple
centers counted once, as asserted.
\end{proof}

\paragraph{The moment interface and its subsequent resolution.}
The usual simplicity mollifier \(Q_T\) in Bui--Heath-Brown,
\emph{On simple zeros of the Riemann zeta-function},
arXiv:1302.5018v1, section~2, equations~(1)--(2) and~(10),
has the optimized limiting first/second moment ratio \(19/27\).
Their theorem does not assert that \(Q_T=B_T^2\) for a
Dirichlet polynomial of the admissible length.  Replacing
\(Q_T\) by an analytic square changes its coefficients
and both moment constants.  Consequently the number
\(19/27\) cannot be inserted into
\eqref{eq:A042-square-ratio}.  A half-power mollifier is a
candidate for that calculation.  The present argument settles the required
positive-weight signed mean and reduces the positive sign
proportions to the displayed strict square-mollifier ratio;
NN-A047 below subsequently proves that strict ratio with
a fully localized half-power mollifier and an exact rational
kernel certificate.  It gives positive ordinary zeta sign
proportions under RH alone.  NN-A048 subsequently proves the fixed-character extension under A.
The half-label law remains open.

\paragraph{Primary interfaces and geometric meaning.}
The imported mean-value input is Benli--Elma--Ng,
\emph{A discrete mean value of the Riemann zeta function},
arXiv:2311.13554v1, Theorem~1.1(i), equations~(4)--(10),
PDF pp.~1--3:
\url{https://arxiv.org/pdf/2311.13554}.
Only the stated theorem interface is imported; no claim
of a complete independent review of its external proof is
made.  The positive weight is the Hermitian squared norm
of the analytic multiplier \(B_T\) on the reduced divisor.
It preserves the normal-velocity sign and vanishes
harmlessly at zeros of the multiplier.  This is a real
analytic norm operation, not a regular algebraic function,
a Frobenius trace, or an arithmetic polarization.


\section{An explicit square mollifier and positive ordinary sign proportions}
\label{sec:A047-positive-simple-proportions}

This section supplies the strict moment inequality left open in the
square-mollifier addendum to NN-A042.  The external inputs are the
two stated discrete mean-value theorems below.  The fractional
mollifier asymptotic, the rational kernel bound, its smoothing
argument, and the sign-count deduction are proved here.

\begin{theorem}[NN-A047: positive proportions of both ordinary zeta signs]
\label{thm:A047-positive-proportions}
Assume RH.  Let \(n_\pm(T)\) count the distinct simple zeros
\(\rho=\frac12+i\gamma\), \(T<\gamma\le2T\), for which respectively
\[
       \pm\Re\left(\frac{\rho\zeta(\rho+1)}{\zeta'(\rho)}\right)>0.
\]
With \(N_T=N(2T)-N(T)\) counting all zero occurrences with
multiplicity,
\[
                 \liminf_{T\to\infty}\frac{n_\pm(T)}{N_T}
                                                     \ge10^{-14}.
\]
This is an absolute, deliberately conservative lower bound.
It assumes neither simplicity of every zero nor a density-zero
exceptional set of multiple centers.  It does not assert half
balance, individual transversality at every simple zero, or a
uniform theorem for Dirichlet conductors.
\end{theorem}

\subsection*{The exact mean-value inputs}

The first input is Benli--Elma--Ng,
\emph{A discrete mean value of the Riemann zeta function},
arXiv:2311.13554v1, Theorem~1.1(i), equations~(4)--(11),
PDF pp.~1--3.  Its bilinear extension to arbitrary real,
uniformly divisor-bounded coefficients and its differentiated
error are proved in
Lemma~\ref{lem:A042-bilinear-majorant}.

The second is Heap--Li--Zhao,
\emph{Lower bounds for discrete negative moments of the Riemann
zeta function}, arXiv:2003.09368v1, Theorem~2, equations~(5)--(9),
PDF pp.~3--5.  That theorem and the following remark explicitly
state that the twisted mean-value formula is unconditional.
The simplicity assumption for the negative moments elsewhere
in their paper is not an assumption of this formula.
For real coefficients, write \(Q(s)=\sum_{n\le y}a(n)n^{-s}\),
\(y\le T^\vartheta\), with a fixed \(\vartheta<1/2\).  In the
notation of that source its formula is
\begin{equation}\label{eq:A047-HLZ}
 \sum_{0<\gamma\le U}\zeta(\rho+\alpha)
          \zeta(1-\rho+\beta)Q(\rho)Q(1-\rho)
 =\mathcal J(\alpha,\beta,U)+\mathcal I(\alpha,\beta,U)
       +\mathcal I(\beta,\alpha,U)+O_A(U(\log U)^{-A}),
\end{equation}
uniformly for sufficiently small complex shifts of size
\(O(1/\log U)\).  The error is available for every fixed
\(A>0\).  Here
\begin{align}
 \mathcal J={}&\frac1{2\pi}\int_0^U\log\frac t{2\pi}
 \sum_{\substack{g,h,k\ge1\\(h,k)=1}}
 \frac{a(gh)a(gk)}{ghk}
 \left\{\frac{\zeta(1+\alpha+\beta)}{h^\beta k^\alpha}
  +\left(\frac t{2\pi}\right)^{-\alpha-\beta}
       \frac{\zeta(1-\alpha-\beta)}{h^{-\alpha}k^{-\beta}}
 \right\}\,dt, \label{eq:A047-J}\\
 \mathcal I={}&\frac1{2\pi}\int_0^U
 \bigl.\partial_\gamma
 \begin{aligned}[t]
 &\sum_{\substack{g,h,k\ge1\\(h,k)=1}}
          \frac{a(gh)a(gk)}{ghk}\\
 &\quad{}\times
 \left\{\begin{aligned}
 &Z_{\alpha,\beta,\gamma,h,k}\\
 &+\left(\frac t{2\pi}\right)^{-\alpha-\beta}
          Z_{-\beta,-\alpha,\gamma,h,k}\\
 &+\left(\frac t{2\pi}\right)^{-\beta-\gamma}
          Z_{\alpha,-\gamma,-\beta,h,k}
 \end{aligned}\right\}
 \end{aligned}\bigr|_{\gamma=0}\,dt . \label{eq:A047-I}
\end{align}
Coefficients outside their support vanish.  The local factor is
\begin{equation}\label{eq:A047-Z}
\begin{aligned}
 Z_{\alpha,\beta,\gamma,h,k}
 &=\frac1{h^\beta}
  \frac{\zeta(1+\alpha+\beta)\zeta(1+\beta+\gamma)}
       {\zeta(1+\beta)}\\
 &\quad{}\times\prod_{p^v\parallel k}
 \frac{\sum_{m\ge0} f_{\alpha,\gamma}(p^{m+v})
                                  p^{-m(1+\beta)}}
      {\sum_{m\ge0} f_{\alpha,\gamma}(p^m)
                                  p^{-m(1+\beta)}} .
\end{aligned}
\end{equation}
Here \(f_{\alpha,\gamma}=\mu*(n^{-\alpha})*(n^{-\gamma})\).
Both \(\mathcal J\) and the three-term expression defining
\(\mathcal I\) have their removable values at coinciding shifts.
These are the source's actual main terms, not a ratios-conjecture
assumption.

The coefficient submultiplicativity hypothesis of that theorem
can also be removed here.  For a real divisor-bounded \(a\), take
the positive majorant \(H\) of
Lemma~\ref{lem:A042-bilinear-majorant}, and put \(U=2H+a\),
\(V=2H\).  The three sequences \(U,V,U+V\) are positive,
uniformly submultiplicative and divisor-bounded.  For any
quadratic expression arising from a bilinear form,
\[
             \mathcal F(a,a)=
          2\mathcal F(U,U)+2\mathcal F(V,V)
                                  -\mathcal F(U+V,U+V).
\]
Apply the source theorem three times.  Its quadratic main
terms recombine by this identity, and its errors remain
\(O_A(U(\log U)^{-A})\).  This proves
\eqref{eq:A047-HLZ} for the coefficients used below.

\subsection*{The fractional square and its first moment}

Let \(d_{-1/2}(n)\) denote the coefficients of
\(\zeta(s)^{-1/2}\) in \(\Re s>1\), with constant coefficient
one.  They are real, multiplicative, and have modulus at most
one: locally they are the coefficients of \((1-z)^{1/2}\),
whose nonconstant absolute coefficients are at most one.
Fix \(0<\nu<1/4\), put \(Y=T^\nu\), and let \(P\) be a fixed
real polynomial with \(P(0)=0\), \(P(1)=1\).  Define
\begin{equation}\label{eq:A047-mollifier}
 B_T(s)=\sum_{n\le Y}\frac{d_{-1/2}(n)}{n^s}
               P\left(\frac{\log(Y/n)}{\log Y}\right),
 \qquad Q_T(s)=B_T(s)^2.
\end{equation}
Thus \(Q_T\) has length \(Y^2=T^{2\nu}<T^{1/2}\);
its coefficients are bounded by \(C_P\tau_2(n)\).
All coefficient constants and the length exponent are fixed
before \(T\) tends to infinity.

\begin{lemma}[First square-mollifier constant]
\label{lem:A047-first}
Put
\[
 I_P=\int_0^1P(t)\,dt,\qquad
 J_P=\int_0^1(1-t)P(t)\,dt .
\]
Then
\begin{equation}\label{eq:A047-first-constant}
 \frac1{N_T\log T}\sum_{T<\gamma\le2T}
                       \zeta'(\rho)Q_T(\rho)
       \longrightarrow
 a(P,\nu):=\frac12+\nu I_P+
                           \frac{\nu^2}{4}(I_P^2-J_P).
\end{equation}
This mean-value assertion is unconditional when the source
sum is over its nontrivial zeros; RH is needed later to use
the conjugate square as a positive moment.
\end{lemma}

\begin{proof}
Apply the first input with \(X=Q_T\), \(Y=1\), and
differentiate its shift at zero.  Write \(L=\log T\).
The \(k=1\) term of its equation~(10) contributes
\(\frac12 L^2+O(L)\), after the factor \(T/(2\pi)\)
is removed.  For a prime \(p\), the derivative of that
factor, to the precision relevant here, is
\[
                  -L\log p+\tfrac12(\log p)^2.
\]
For distinct primes \(p,q\) it is
\(\log p\log q\).  Terms with at least three distinct
prime divisors vanish exactly: the product
\(\prod_{p\mid k}(1-p^{-\alpha})\) has order at least
three at zero, and multiplication by the simple zeta
pole followed by one derivative still leaves zero.

The Euler corrections \(p/(p-1)\) and the remaining
Laurent constants contribute \(O_P(L)\).  The same is
true for terms with a repeated prime.  For example,
the sums
\[
 \sum_{p}\sum_{j\ge2}
         \frac{(j+1)^2j^2(\log p)^2}{p^j}
\]
converge, and a second distinct prime can add at most
one factor \(O(L)\).  These bounds also cover the
divisor coefficients of \(Q_T\).

With \(u=\log p/\log Y\), the coefficient of \(p\) in
\(Q_T\) is \(-P(1-u)\), for \(0\le u\le1\).
For distinct \(p,q\), putting \(v=\log q/\log Y\),
the coefficient is
\[
 \frac12\left\{
 P(1-u-v)\mathbf1_{u+v\le1}
                   +P(1-u)P(1-v)\right\},
 \qquad 0\le u,v\le1.
\]
The prime number theorem, or its RH consequence already
used in NN-A040, implies the elementary weak convergence
\[
 \frac1{\log Y}\sum_{p\le Y}
          \frac{\log p}{p}\,\delta_{\log p/\log Y}
                         \ \longrightarrow\ du
                    \quad\hbox{on }[0,1].
\]
It follows by partial summation from
\(\sum_{p\le x}\log p/p=\log x+o(\log x)\); the latter also
follows directly from the prime-power form of the prime
number theorem.  The diagonal \(p=q\) contributes
\(o((\log Y)^2)\) to the product measure.  Thus unordered
distinct prime pairs give one half the double integral.

The prime terms contribute
\(\nu I_P-\frac12\nu^2J_P\); the pair terms contribute
\(\frac14\nu^2(J_P+I_P^2)\), since
\[
 \int_{\substack{u,v\ge0\\u+v\le1}}P(1-u-v)\,du\,dv=J_P.
\]
Combining these with \(1/2\) gives
\eqref{eq:A047-first-constant}.  The source error, after
Cauchy differentiation, is negligible.  Applying its
formula at \(2T\) and \(T\) with the same polynomial
preserves these leading constants, because
\(\log(2T)=L+O(1)\) and \(N_T\sim TL/(2\pi)\).
\end{proof}

\subsection*{The four-variable second-moment limit}

The next lemma first assumes that \(P\) vanishes to order at
least sixteen at zero.  This is a technical smoothing
condition used in the proof; it will be removed only in
the limiting kernel calculation, not silently in the
arithmetic asymptotic.

Let \(s_1,s_2,t_1,t_2\) be complex variables with real parts
\(1/\nu\), and put
\[
 F_0(s,t)=
  \frac{(s_1s_2t_1t_2)^{1/2}}
             {\prod_{i,j=1}^2(s_i+t_j)^{1/4}},
 \quad
 E_L(x)=\prod_{i=1}^2(1+x/s_i)^{1/2},\quad
 E_R(x)=\prod_{j=1}^2(1+x/t_j)^{1/2}.
\]
Use the branches equal to one at \(x=0\).
Write \(E_L(x)=\sum l_kx^k\),
\(E_R(x)=\sum r_kx^k\).  Define
\begin{align}
 j(l,r)={}&l_1r_1+l_1r_2+\tfrac12l_1+l_2r_1-l_3
                                      +\tfrac12r_1-r_3+\tfrac13,
                                                    \label{eq:A047-j}\\
 i(l,r)={}&l_1r_1^2+l_1r_1r_2+l_1r_3-\tfrac13l_1
       +l_2r_1^2-l_2r_1-l_2r_2-\tfrac12l_2
       -2l_3r_1+l_4 \nonumber\\
       &+\tfrac12r_1^2+r_1r_2-r_1r_3+\tfrac16r_1
                          +r_2^2+\tfrac12r_2-r_4-\tfrac18.
                                                    \label{eq:A047-i}
\end{align}
Let \(\mathscr B=j+2i\).  Substitute the finite expressions
\[
 l_k=\sum_{h=0}^k
       {1/2\choose h}{1/2\choose k-h}s_1^{-h}s_2^{-(k-h)},
 \qquad
 r_k=\sum_{h=0}^k
       {1/2\choose h}{1/2\choose k-h}t_1^{-h}t_2^{-(k-h)}.
\]
This gives a polynomial
\[
 \mathscr B=\sum_{\mathbf k}c_{\mathbf k}
           s_1^{-k_1}s_2^{-k_2}t_1^{-k_3}t_2^{-k_4}
\]
with rational coefficients, total degree at most four,
and exactly sixty-two nonzero monomials.

For \(k\ge0\), define on \(0\le z\le1\)
\begin{equation}\label{eq:A047-vertex}
 \Psi_k^P(z)=
   \sum_{j\ge1}p_j
             \frac{\Gamma(j+1)}{\Gamma(j+k+1/2)}
                                      z^{j+k-1/2},
       \qquad P(z)=\sum_{j\ge1}p_jz^j .
\end{equation}
Let
\begin{align*}
 \mathcal D&=\{x_{ij}\ge0:
                   \textstyle\sum_jx_{ij}\le1,\
                   \sum_ix_{ij}\le1\},\\
 z_1&=1-x_{11}-x_{12},\qquad z_2=1-x_{21}-x_{22},\\
 z_3&=1-x_{11}-x_{21},\qquad z_4=1-x_{12}-x_{22}.
\end{align*}
The following finite integral is well defined:
\begin{equation}\label{eq:A047-b-functional}
 b(P,\nu)=\frac1{\Gamma(1/4)^4}
 \int_{\mathcal D}\prod_{i,j}x_{ij}^{-3/4}
 \sum_{\mathbf k}c_{\mathbf k}\nu^{|\mathbf k|-1}
                       \prod_{\ell=1}^4\Psi_{k_\ell}^P(z_\ell)
                                      \prod_{i,j}dx_{ij}.
\end{equation}

\begin{lemma}[Second square-mollifier constant]
\label{lem:A047-second}
If \(P\) has zero of order at least sixteen at zero, then,
under RH,
\begin{equation}\label{eq:A047-second-constant}
 \frac1{N_T(\log T)^2}
       \sum_{T<\gamma\le2T}
                    |\zeta'(\rho)B_T(\rho)^2|^2
                            \longrightarrow b(P,\nu).
\end{equation}
The constant is the integral \eqref{eq:A047-b-functional};
no unproved moment formula is used.
\end{lemma}

\begin{proof}
The source formula \eqref{eq:A047-HLZ} applies to \(Q_T=B_T^2\)
by the proved coefficient extension.  Apply it at \(2T\)
and \(T\) with the same polynomial, and differentiate the
two shifts.  Its error is still negligible by Cauchy's
formula.  Its main integrals are now exactly over
\(T\le t\le2T\).

Here are the Euler products and estimates needed to evaluate
those main terms, including the uniformity in the small shifts.
For any monomial of \(P\), the elementary inverse Laplace
formula gives
\[
 P\left(\frac{\log(Y/n)}{\log Y}\right)\mathbf1_{n\le Y}
 =\frac1{2\pi i}\int_{(1/\log Y)}
 (Y/n)^u\sum_j\frac{p_jj!}{(\log Y)^j u^{j+1}}\,du .
\]
At \(n=Y\), both sides are zero.  Apply this formula in the
four coefficient variables, and scale \(u=s/\log T\).
Each vertex weight becomes
\[
             W_{P,\nu}(s)
                  =e^{\nu s}\sum_j\frac{p_jj!}{\nu^j s^{j+1}},
                  \qquad\Re s=1/\nu .
\]
Each Euler identity below is first unfolded with independent
shifts in an absolute-convergence chamber, for example positive
real coefficient variables and positive real shift variables.
It is then continued meromorphically, term by term before the
indicated substitutions, to the distinct Cauchy circles below.
The finite coefficient sums and the displayed zeta factors
times locally convergent Euler products identify that
continuation.  The initial sums and Mellin integrals converge
absolutely in the stated chamber.

For the gcd sum in \(\mathcal J\), the first-order prime
terms are
\[
 -\tfrac12\sum_i p^{-1-u_i-\beta}
 -\tfrac12\sum_j p^{-1-v_j-\alpha}
 +\tfrac14\sum_{i,j}p^{-1-u_i-v_j}.
\]
Consequently its four-variable Euler product is
\begin{equation}\label{eq:A047-F-euler}
 C(u,v;\alpha,\beta)
 \frac{\prod_{i,j}\zeta(1+u_i+v_j)^{1/4}}
      {\prod_i\zeta(1+u_i+\beta)^{1/2}
       \prod_j\zeta(1+v_j+\alpha)^{1/2}}.
\end{equation}
The factor \(C\) is an absolutely convergent Euler product
in a fixed neighborhood of the origin.  For example,
taking all real parts of its arguments of modulus less
than \(1/32\) leaves a uniform summable bound
\(O(p^{-3/2})\) after the displayed first-order terms
are removed.  Its value at the origin is one.
Indeed, the two half-power coefficient series combine
to \(\mu\) on each side, and the local gcd sum there is
\(1-1/p\), exactly the local factor of the displayed
zeta quotient.

For the \(Z\) sum in \(\mathcal I\), unfold
\(f_{\alpha,\gamma}=\mu*n^{-\alpha}*n^{-\gamma}\).
The balanced relation is \(hm=kn\), with the two half
coefficients on \(h\), the two on \(k\), the three
displayed factors on \(m\), and \(n^{-\beta}\) on \(n\).
Pairing one prime on each side gives exactly
\begin{align}
 D(u,v;\alpha,\beta,\gamma)&
 \frac{\zeta(1+\alpha+\beta)\zeta(1+\beta+\gamma)}
      {\zeta(1+\beta)}\nonumber\\
 &{}\times
 \frac{\prod_{i,j}\zeta(1+u_i+v_j)^{1/4}}
      {\prod_i\zeta(1+u_i+\beta)^{1/2}}
 \prod_j
 \frac{\zeta(1+v_j)^{1/2}}
      {\zeta(1+v_j+\alpha)^{1/2}
                         \zeta(1+v_j+\gamma)^{1/2}} .
 \label{eq:A047-G-euler}
\end{align}
The same argument bounds the higher prime terms uniformly
by \(O(p^{-3/2})\).  Thus \(D\) is also absolutely convergent
and bounded near the origin.  Its value there is one:
the full left coefficient is
\(\mu*(\mu*1*1)=\delta_1\), and the full right coefficient
is \(\mu*1=\delta_1\); hence every local balanced sum is one.
The zeta exponents in \eqref{eq:A047-G-euler} also sum to
zero there.  This proves the normalization, including all
prime powers, without an unspecified arithmetic factor.

Put \(L=\log T\), and use scaled shifts
\(\alpha=A/L,\beta=B/L,\gamma=C/L\).  Choose their Cauchy
circle radii to be distinct sufficiently small fixed
positive numbers, for example \(1/100,2/100,3/100\).
All sums or differences appearing in the three separate
terms then stay bounded away from zero on the relevant
circles.  Every zeta factor involving a coefficient variable
has real part strictly greater than one.

We give the tail estimates before taking a limit.
On the scaled Mellin lines the elementary estimates
\[
 |\zeta(1+\varepsilon+it)|\le\zeta(1+\varepsilon),
 \qquad
 |\zeta(1+\varepsilon+it)^{-1}|
                                  \le\zeta(1+\varepsilon)
\]
apply with \(\varepsilon\asymp1/L\).  The fixed-shift
prefactor of \eqref{eq:A047-G-euler} is \(O(L)\) by its
Laurent expansion, uniformly on the chosen circles.
Thus \eqref{eq:A047-F-euler} is \(O(L^3)\), and
\eqref{eq:A047-G-euler} is \(O(L^6)\), before cancellation.
The absolutely convergent factors \(C,D\) are uniformly
bounded also for large imaginary coefficient variables.
If \(J\ge16\) is the smallest power of \(P\), each vertex
weight is \(O_P((1+|\Im s|)^{-J-1})\).  Removing the
regions where any \(|\Im s|>\sqrt L\) therefore costs
\[
                   O_P(L^{3-J/2})
       \quad\hbox{and}\quad O_P(L^{6-J/2}),
\]
respectively.  These are \(o(L^{-1})\) and \(o(1)\),
the two required scales.

On the remaining box every zeta argument is \(1+O(L^{-1/2})\).
Use
\(\zeta(1+z)=z^{-1}(1+O(z))\), with the analytic branches
already fixed in \(\Re z>0\), and \(C,D=1+o(1)\).
After multiplying \eqref{eq:A047-F-euler} by \(L\),
its limit is
\[
                  F_0(s,t)E_L(B)E_R(A).
\]
The limit of \eqref{eq:A047-G-euler} is
\[
       R(A,B,C)=
       F_0(s,t)\frac{B}{(A+B)(B+C)}
                           E_L(B)E_R(A)E_R(C).
\]
The relative errors are uniform on the central box and the
chosen shift circles.  The model kernels have at most fixed
polynomial growth in each Mellin variable; their denominator
\(s_i+t_j\) has positive real part bounded away from zero.
The vertex weights with \(J\ge16\) therefore give an
integrable majorant.  Dominated convergence, followed by
Cauchy's formula in the scaled shifts, justifies the limit
and all three differentiations.  Also
\(\log(t/(2\pi))/L=1+O(1/L)\) uniformly for
\(T\le t\le2T\), so the exponential factors have exactly
the limits used next.

The normalized \(\mathcal J\) kernel is
\[
 F_0\frac{E_L(B)E_R(A)-
                    e^{-A-B}E_L(-A)E_R(-B)}{A+B}.
\]
Its \(AB\) coefficient is \(F_0j(l,r)\).
The normalized \(\mathcal I\) kernel before its \(C\)
derivative is
\[
          R(A,B,C)+e^{-A-B}R(-B,-A,C)
                                  +e^{-B-C}R(A,-C,-B).
\]
After division by \(F_0\), its common numerator is
\begin{align*}
 &B(A-C)E_L(B)E_R(A)E_R(C)\\
 &\quad-A(B+C)e^{-A-B}E_L(-A)E_R(-B)E_R(C)\\
 &\quad+C(A+B)e^{-B-C}E_L(-C)E_R(A)E_R(-B),
\end{align*}
over \((A+B)(B+C)(A-C)\).
The numerator vanishes on each of the three indicated
hyperplanes, so the quotient is holomorphic at their
intersections as well.  Expanding to total degree six
gives its \(ABC\) coefficient \(i(l,r)\) in
\eqref{eq:A047-i}.  The two copies of \(\mathcal I\)
give the same \(AB\) derivative, hence the kernel
\(F_0\mathscr B\).  These finite polynomial divisions
are also checked exactly in the accompanying certificate.

Finally, for each cross factor use
\[
 (s_i+t_j)^{-1/4}=
       \frac1{\Gamma(1/4)}
       \int_0^\infty x^{-3/4}e^{-x(s_i+t_j)}\,dx.
\]
Fubini is justified by the absolute convergence just proved.
The one-variable inverse Laplace integrals give
\[
 \frac1{2\pi i}\int_{(1/\nu)}
     W_{P,\nu}(s)s^{1/2-k}e^{-xs}\,ds
 =\nu^{k-1/2}\Psi_k^P(1-x/\nu)\mathbf1_{x\le\nu}.
\]
Scale all four cross variables by \(\nu\).  Their measure
contributes \(\nu\), and the four vertex factors contribute
\(\nu^{|\mathbf k|-2}\).  The resulting integral is precisely
\eqref{eq:A047-b-functional}.  Under RH, the differentiated
source sum is the positive second moment in
\eqref{eq:A047-second-constant}.  This completes the proof.
\end{proof}

\subsection*{An exact upper bound for a limiting polynomial}

Take \(P_0(z)=2z-z^2\) and \(\nu_0=1/4\).  We evaluate the
kernel functional, without claiming the arithmetic
asymptotic for this insufficiently smoothed \(P_0\).
Equation~\eqref{eq:A047-first-constant} gives
\begin{equation}\label{eq:A047-a0}
                         a_0=a(P_0,1/4)=1543/2304.
\end{equation}
The vertex functions are
\[
 \Psi_k^{P_0}(z)=
       \frac{2z^{k+1/2}}{\Gamma(k+3/2)}
                         \left(1-\frac{z}{k+3/2}\right).
\]
Factor \(\sqrt{z_1z_2z_3z_4}/\pi^2\) out of
the finite kernel sum in \eqref{eq:A047-b-functional}.
The remaining polynomial is \(Q(S,D,E)\), where
\[
 S=z_1+z_2=z_3+z_4,\qquad
 D=(z_1-z_2)^2,\qquad E=(z_3-z_4)^2 .
\]
It is specified without numerical approximation by
\eqref{eq:A047-j}--\eqref{eq:A047-i} and the preceding
vertex substitution.  In particular its constant and
linear-\(D\) coefficients are
\begin{align}
 q_0(S)={}&\frac8{245581875}
 \bigl(776S^8+75340S^7+287460S^6-5504100S^5
       -11151525S^4 \nonumber\\
 &\hspace{28mm}
       +151420500S^3+43659000S^2
                         -1746360000S+2619540000\bigr),
                                                    \label{eq:A047-q0}\\
 q_D(S)={}&-\frac8{245581875}
 \bigl(5207S^6-140S^5+948840S^4+1451100S^3
                    \nonumber\\
 &\hspace{28mm}
                    -49075950S^2+30492000S+261954000\bigr).
                                                    \label{eq:A047-qD}
\end{align}
The exact polynomial bounds used below are
\begin{equation}\label{eq:A047-poly-bound}
 Q(S,D,E)\le
 U(S,D,E):=q_0(S)+q_D(S)D-(2-3S/4)E
\end{equation}
for \(0\le S\le2\) and
\(0\le D,E\le\min(S,2-S)^2\), together with
\[
 q_D(S)<0,\qquad
 U(S,D,E)\ge0
 \quad\left(
 0\le D\le\min(S,2-S)^2,\quad
 0\le E\le(2-S)^2\right).
 \label{eq:A047-upper-positive}
\]

Here is an exact finite certificate for those inequalities.
Put
\[
 A(S,D)=-[D^2]Q-[D^3]Q\,D,\qquad
 B(S,D,E)=-(2-3S/4)-\frac{Q(S,D,E)-Q(S,D,0)}E .
\]
The polynomial identity
\[
                      U-Q=D^2A+EB
\]
is exact.  In each of the two substitutions
\[
 (S,D,E)=(u,u^2v,u^2w),\qquad
 (S,D,E)=(2-u,u^2v,u^2w),
                   \qquad 0\le u,v,w\le1,
\]
every tensor Bernstein coefficient is positive.  The
least coefficients, as exact rationals, are respectively
\[
 \begin{array}{c|cc}
    & S=u & S=2-u\\ \hline
 A &5060984/81860625&64264/3898125\\
 B &4133699/13097700&38706253/491163750 .
 \end{array}
\]
The univariate Bernstein coefficients of \(q_D(2u)\)
are all negative.  The least coefficients of
\[
 q_0(S)+q_D(S)\min(S,2-S)^2
                              -(2-3S/4)(2-S)^2
\]
on the same two halves are
\(167144539/6682500\) and \(1492032128/245581875\).
These prove \eqref{eq:A047-poly-bound} and the larger-domain
nonnegativity assertion.  For reproducibility the certificate
uses the formula
\[
 [B_{\mathbf k}]\ p
   =\sum_{\mathbf j\le\mathbf k}a_{\mathbf j}
           \prod_i\frac{\binom{k_i}{j_i}}{\binom{n_i}{j_i}},
 \qquad p(x)=\sum a_{\mathbf j}x^{\mathbf j},
\]
where \(n_i\) is the degree in the \(i\)-th variable.
Bernstein basis functions on the cube are nonnegative and
sum to one.  Thus this is a finite rational proof, not a
floating point quadrature.  The full coefficient construction,
identity check, all Bernstein coefficients, and the beta
integration below are reproduced by
\path{work/analytic/check_square_mollifier_kernel.py}.

To integrate the upper bound, put
\[
 r_1=x_{11}+x_{12},\quad r_2=x_{21}+x_{22},\qquad
 x_{11}=r_1v_1,\quad x_{21}=r_2v_2 .
\]
On the actual domain the column remainders are nonnegative,
and their sum is \(S=2-r_1-r_2\).  Hence
\[
 \sqrt{z_3z_4}\le S/2.
\]
Because \(Q\le U\) and \(U\ge0\), replacing
\(\sqrt{z_3z_4}Q\) by \((S/2)U\) is valid even
where \(Q\) might be negative.  The larger-domain
nonnegativity then permits dropping both column constraints.

Under the normalized angular beta measure,
\(v_1,v_2\) are independent
\(\operatorname{Beta}(1/4,1/4)\) variables, so
\[
 \mathbb E E=\frac23(r_1^2+r_2^2),\qquad
 D=(r_1-r_2)^2 .
\]
Their normalization contributes \(1/\pi\) after division
by \(\Gamma(1/4)^4\).
The remaining row weights are
\(r_i^{-1/2}(1-r_i)^{1/2}\,dr_i\); each has
integral \(\pi/2\).
Consequently, for independent
\(R_1,R_2\sim\operatorname{Beta}(1/2,3/2)\),
\begin{align}
 b(P_0,1/4)
 &\le\frac1{4\pi}\mathbb E\left[
 \frac S2\left\{
 q_0(S)+q_D(S)(R_1-R_2)^2
       -(2-3S/4)\frac23(R_1^2+R_2^2)\right\}\right]\nonumber\\
 &=\frac{1505515846753}{536481792000\,\pi},
 \qquad S=2-R_1-R_2 . \label{eq:A047-rational-upper}
\end{align}
The equality is a finite rational calculation using
\(\mathbb E R_i^k=(1/2)_k/(2)_k\); it is independently
reproduced in the certificate.

For a fully rational comparison, \(\pi>314159/100000\).
One elementary verification is Machin's identity
\(\pi=16\arctan(1/5)-4\arctan(1/239)\), followed by
the alternating lower sum through \(-1/(7\cdot5^7)\)
for the first arctangent and the upper bound \(1/239\)
for the second.  The tangent addition formula proves
Machin's identity with the angle in \((0,\pi/2)\).
Thus
\begin{equation}\label{eq:A047-strict-gap}
 2a_0^2-
 \frac{1505515846753}
      {536481792000(314159/100000)}
       =\frac{5046955580047}{1348324666343424}>0 .
\end{equation}

\subsection*{High smoothing preserves the strict gap}

For positive integers \(K\), set
\[
             P_K(z)=(2z-z^2)\bigl(1-(1-z)^K\bigr)^{15}.
\]
This is a polynomial with \(P_K(1)=1\) and a zero of
order sixteen at zero.  Uniformly on \([0,1]\),
\[
 0\le P_K'(z)\le32,\qquad
 P_K'\longrightarrow P_0'\quad\hbox{in }L^4[0,1].
\]
Indeed \(0\le P_0'\le2\), \(P_0(z)\le2z\), and
\(Kz(1-z)^{K-1}\le1\); differentiating the cutoff gives
the stated upper bound.  Pointwise convergence away
from zero and dominated convergence prove the \(L^4\)
assertion.  Also \(P_K\to P_0\) uniformly, by integrating
the \(L^1\) convergence of their derivatives.

The vertex functions admit the real integral formulas
\[
 \Psi_0^P(z)=\frac1{\sqrt\pi}
       \int_0^z\frac{P'(t)}{\sqrt{z-t}}\,dt,\qquad
 \Psi_k^P(z)=\frac1{\Gamma(k-1/2)}
       \int_0^zP(t)(z-t)^{k-3/2}\,dt\quad(k\ge1).
\]
They follow by evaluating beta integrals on each monomial.
Hölder's inequality with exponents \(4,4/3\) gives
\[
 \|\Psi_0^{P_K}-\Psi_0^{P_0}\|_\infty
       \le\frac{3^{3/4}}{\sqrt\pi}
                             \|P_K'-P_0'\|_4\longrightarrow0.
\]
The remaining finitely many vertex functions converge
uniformly by the integrable-kernel formula and uniform
convergence of \(P_K\).  The weight
\(\prod x_{ij}^{-3/4}\) has finite integral on the bounded
domain \(\mathcal D\).  Thus
\[
 a(P_K,\nu)\longrightarrow a(P_0,\nu),\qquad
 b(P_K,\nu)\longrightarrow b(P_0,\nu).
\]
The functions are also continuous in \(\nu>0\);
\eqref{eq:A047-b-functional} is a finite sum of powers
\(\nu^{|\mathbf k|-1}\).

It follows from the strict rational gap
\eqref{eq:A047-strict-gap} that there exist a fixed
rational \(0<\nu<1/4\) sufficiently close to \(1/4\),
and then a fixed sufficiently large integer \(K\),
such that the actual arithmetic moments of
\eqref{eq:A047-mollifier} obey
\begin{align*}
 a(P_K,\nu)&\ge a_*:=1543/2304-1/2000,\\
 b(P_K,\nu)&\le b_*:=
 \frac{1505515846753}
      {536481792000(314159/100000)}+1/2000 .
\end{align*}
These choices are made once and then held fixed as
\(T\to\infty\).  No height-dependent smoothing limit
or unproved uniformity in \(K\) is used.

\subsection*{The signed proportion and its scope}

Apply Proposition~\ref{prop:A042-square-mollifier}.
The triangle inequality turns the complex first-moment
asymptotic into the absolute first-moment lower bound used there.
Choose \(\eta=1/50\) and
\(p_\eta=1/(2(1-\eta^2))\).  The certificate verifies
exactly
\[
          0<a_*<1,\qquad 0<b_*<1,\qquad
                       a_*^2-b_*p_\eta>1/2000.
\]
Therefore
\[
 d_\eta=a_*-\sqrt{b_*p_\eta}
   =\frac{a_*^2-b_*p_\eta}{a_*+\sqrt{b_*p_\eta}}
                                               >1/4000.
\]
Finally \(\zeta(3/2)<3\), by comparison with
\(1+\int_1^\infty x^{-3/2}\,dx\).  The exact lower
fraction in \eqref{eq:A042-square-density} is hence
strictly greater than
\[
 \frac{(1/50)^2(1/4000)^2}{4\cdot3^4}
              =\frac1{12960000000000}>10^{-14}.
\]
The multiplier weight is nonnegative, and its nonzero
weighted summands occur only at simple transverse
centers.  The underlying normal velocity is the
specified denominator-shift velocity throughout.
This proves Theorem~\ref{thm:A047-positive-proportions}.

\paragraph{Geometric translation.}
On the reduced analytic zero divisor, the norm-shift
family has its unique local section and hence a
normal tangent relative to the critical real line.
The statistic \(K_\rho\) pairs that tangent with the
positive squared derivative norm; it vanishes at
all nonreduced centers.  The analytic multiplier
adds the nonnegative Hermitian norm \(|B_T(\rho)|^2\).
The proved moment inequality bounds its concentration,
and the signed mean forces positive counting
proportions in both oriented reduced loci.
The reduced finite-algebra differential criterion is
SGA~1, Expos\'e~I, section~3, Proposition~1:
\url{https://grothendiecksga.com/read/sga1/fr/I-3.html}.
This is a real analytic norm and counting argument;
it does not define an order on arbitrary complex
schemes, an arithmetic denominator-shift family,
a pure motive, or a Frobenius half-label law.

\paragraph{Source and certificate scope.}
Primary mean-value sources:
\url{https://arxiv.org/pdf/2311.13554}, v1,
Theorem~1.1(i), PDF pp.~1--3; and
\url{https://arxiv.org/pdf/2003.09368}, v1,
Theorem~2, equations~(5)--(9), PDF pp.~3--5.
Their exact theorem interfaces are imported.
No full independent review of both external papers
is claimed.  The new localization, limiting kernels,
smoothing continuity, and arithmetic application
are proved in this section.  The certificate
checks finite rational identities and inequalities;
the exploratory Monte Carlo file is not used
anywhere as proof.


% Full root review 2026-11-22 04:00:19 UTC; peer GCT review 03:53:18 UTC.
\section{Square mollifiers for a fixed primitive character}
\label{sec:A048-fixed-character-square}

The conductor in this section is fixed, but the auxiliary
character moduli grow with the mollifier. Their GRHs are
supplied by A. The right-edge calculation, including its
principal scalar and coefficient-averaged error, is proved
below; it is not inferred from the zeta mean-value theorem.

\begin{theorem}[NN-A048: both ordinary signs for every fixed primitive character]
\label{thm:A048-signs}
Assume A and fix a primitive Dirichlet character \(\chi\)
of conductor \(q>1\). Let \(N_{\chi,T}\) count all its
zero occurrences in \(T<\gamma\le2T\), with multiplicity,
and let \(n_{\chi,\pm}(T)\) count distinct simple centers with
\[
 \pm\Re\frac{\rho L(\rho+1,\chi)}{L'(\rho,\chi)}>0,
       \qquad \rho=\frac12+i\gamma.
\]
Then
\[
 \liminf_{T\to\infty}
       \frac{n_{\chi,\pm}(T)}{N_{\chi,T}}\ge10^{-14}.
\]
The constant is independent of the fixed character; the
eventual height threshold can depend on it. No uniformity
for a growing conductor is asserted. There is no simplicity
or density-zero multiplicity premise. For conductor one,
RH alone already suffices by NN-A047.
\end{theorem}

Put
\[
 Z_q(s)=\zeta(s)\prod_{p\mid q}(1-p^{-s}),\qquad
 E_q=\frac{\varphi(q)}q,\qquad L=\log T.
\]
Thus \(\operatorname{Res}_{s=1}Z_q(s)=E_q\).
For real coefficients \(a(n)\), write
\(A_\chi(s)=\sum_{n\le y}a(n)\chi(n)n^{-s}\).
The coefficient bounds used below are
\begin{equation}\label{eq:A048-coeff}
 |a(n)|\le C\tau_r(n)(\log(2n))^r
\end{equation}
with \(C,r\) fixed, uniformly in height and support.

\subsection*{Uniform auxiliary-character bounds proved locally}

\begin{lemma}[Uniform off-line growth]
\label{lem:A048-uniform-growth}
For each fixed \(\eta,\epsilon>0\), A implies
\begin{equation}\label{eq:A048-uniform-growth}
 |L(\sigma+it,\psi)|+|L(\sigma+it,\psi)^{-1}|
 \ll_{\eta,\epsilon}\{Q(|t|+3)\}^{\epsilon},
 \qquad \frac12+\eta\le\sigma\le2,
\end{equation}
for every nonprincipal character \(\psi\bmod Q\).
For a principal character the same assertion holds on
the line \(\sigma=1/2+\eta<1\), and on horizontal
segments in this strip with \(|t|\ge2\).
Fixed-order derivatives satisfy the corresponding bounds
on a smaller closed strip.
\end{lemma}

\begin{proof}
First let \(\psi\) be primitive and nonprincipal.
The root-number normalized completion is entire of order
one and invariant under conjugate reflection about
\(\Re s=1/2\). Its centered Hadamard linear coefficient
is purely imaginary, as proved in NN-A035. Under A every
centered zero is purely imaginary, so the real parts of
the genus-one correction constants are zero. Therefore
\begin{equation}\label{eq:A048-uniform-zero-sum}
 \sum_\rho\frac{3/2}{(3/2)^2+(t-\gamma)^2}
 =\Re\frac{\Lambda'}{\Lambda}(2+it,\psi)
 \ll\log\{Q(|t|+3)\}.
\end{equation}
The Euler series bounds \(L'/L(2+it,\psi)\) absolutely;
the gamma logarithmic derivative and
\(\frac12\log(Q/\pi)\) give the stated uniform bound.
Central zeros retain their multiplicities.

Subtract the Hadamard logarithmic derivatives at
\(\sigma+it\) and \(2+it\). Their zero contribution is
\[
 (2-\sigma)\sum_\rho
       \frac1{(\sigma+it-\rho)(2+it-\rho)}.
\]
Its absolute value is \(O_\eta\) of the sum in
\eqref{eq:A048-uniform-zero-sum}; the gamma difference
has the same bound. Hence
\begin{equation}\label{eq:A048-logderiv}
 \left|\frac{L'}L(\sigma+it,\psi)\right|
       \ll_\eta\log\{Q(|t|+3)\}.
\end{equation}

The following elementary smoothed identity supplies the
small growth and the reciprocal without importing a
uniform-conductor Lindelof theorem:
\begin{align}
 -\frac{L'}L(s,\psi)
 ={}&\sum_{n\le x}\frac{\Lambda(n)\psi(n)}{n^s}
                    \frac{\log(x/n)}{\log x}
       +\frac1{\log x}\left(\frac{L'}L\right)'(s,\psi)
       \nonumber\\
 &+\frac1{\log x}\sum_\rho
                  \frac{x^{\rho-s}}{(\rho-s)^2}
       +\frac1{\log x}\sum_{\rho_{\rm triv}}
                  \frac{x^{\rho_{\rm triv}-s}}
                       {(\rho_{\rm triv}-s)^2}.
 \label{eq:A048-smoothed-log}
\end{align}
Here \(x>1\), \(s=\sigma+it\), and the trivial zeros are
\(-a-2j\), \(j\ge0\), with \(a=(1-\psi(-1))/2\).
To prove it, integrate
\(-L'/L(s+w,\psi)x^w/w^2\) on an absolute-convergence
line. Termwise inversion gives
\(\sum_{n\le x}\Lambda(n)\psi(n)n^{-s}\log(x/n)\).
The residue at zero is
\(-L'/L(s)\log x-(L'/L)'(s)\), and every zero gives
\(-x^{\rho-s}/(\rho-s)^2\).
Shift left. The functional equation makes the new
left edge tend to zero exponentially in its distance.
Horizontal edges tend to zero along heights avoiding
zeros: \eqref{eq:A048-uniform-zero-sum} gives the local
count, one can choose distance \(\gg1/\log(QH)\), and
partial fractions then give \(O(\log^2(QH))\) against
the quadratic denominator. Both zero sums converge
absolutely. This proves the identity.

Integrate in the real coordinate from \(\sigma\) to
infinity, using the analytic Euler logarithm.
Set \({\cal C}=Q(|t|+3)\) and
\(x=\max(3,\log{\cal C})\).
The prime contribution to \(\log L(s,\psi)\) is bounded
by \(\sum_{n\le x}n^{-\sigma}\ll\sqrt x\).
The derivative term is
\(O_\eta(\log{\cal C}/\log x)\), by
\eqref{eq:A048-logderiv}. The zero term is at most
\[
 \frac{x^{-\eta}}{(\log x)^2}
       \sum_\rho\frac1{\eta^2+(t-\gamma)^2}
 \ll_\eta\frac{x^{-\eta}\log{\cal C}}{(\log x)^2}.
\]
The integrated trivial-zero term is bounded uniformly.
Thus
\[
 |\log L(s,\psi)|\ll_\eta
 \sqrt{\log{\cal C}}+
       \frac{\log{\cal C}}{\log\log({\cal C}+3)}
             =o_\eta(\log{\cal C}),
\]
uniformly in the primitive character and height.
Exponentiation proves both claimed bounds.

For an imprimitive character, the deleted Euler factors
and their reciprocals cost at most
\[
 \prod_{p\mid Q}\frac{1+p^{-1/2-\eta}}{1-p^{-1/2-\eta}}
                      \ll_{\eta,\epsilon}Q^\epsilon.
\]
After finitely many primes are removed, each logarithm
is at most \(\epsilon\log p\); the finite initial
product is an allowed constant. For principal characters
the same proof for zeta has one additional pole term at
\(1\) in \eqref{eq:A048-smoothed-log}. For \(|t|\ge2\)
its integrated absolute value is \(O(\sqrt x)\).
On the specified left line with bounded \(t\), continuity
and nonvanishing under RH apply away from the pole.
Deleted factors again cost \(Q^\epsilon\).
Cauchy's formula on a fixed smaller disc gives the
derivative bounds.
\end{proof}


\subsection*{The coefficient-averaged additive estimate}

Every shift below has modulus at most \(C_0/L\), with
\(C_0\) fixed. Let \(c\) be a convolution of at most
four factors of the forms \(\mu(n)n^{-z}\) and \(n^{-z}\),
with at most one Mobius factor, and put \(b=c*a\).
The polynomial \(a\) is supported in \([1,y]\) and
satisfies \eqref{eq:A048-coeff}. A missing polynomial
is represented by \(a=\delta_1\).

\begin{lemma}[Auxiliary character averaging]
\label{lem:A048-character-average}
Fix \(q\) and \(\theta<1/2\), and suppose \(y\le T^\theta\).
In the additive sum
\begin{equation}\label{eq:A048-additive}
 \frac{\tau(\bar\chi)}q
 \sum_{\substack{k\le y\\(k,q)=1}}
   \frac{a_2(k)\bar\chi(k)}{k^{1-\beta}}
 \sum_{n\le qUk/(2\pi)}
   \frac{\chi(n)b(n)}{n^\beta}
                      e\left(-\frac n{qk}\right),
       \qquad T\le U\le2T,
\end{equation}
all nonprincipal character components have total
\(O_{q,\epsilon}(T^{1/2+\epsilon}y)\).
The principal part is exactly
\begin{equation}\label{eq:A048-principal}
 \frac1{\varphi(q)}
 \sum_{\substack{k\le y\\(k,q)=1}}\frac{a_2(k)}{k^{1-\beta}}
 \sum_{\substack{n\le qUk/(2\pi)\\(n,q)=1}}
   \frac{b(n)}{n^\beta}
          \frac{\mu(k/(n,k))}{\varphi(k/(n,k))}.
\end{equation}
Its inner sums can be replaced by the residues of their
Dirichlet series to the right of a fixed line
\(\Re s=1/2+\eta\), with total error of the same size.
These statements are uniform on small complex shift
circles. Fixed-order shift derivatives and logarithmic
weights are allowed after enlarging \(\epsilon\).
\end{lemma}

\begin{proof}
Put \(d=(n,k)\), \(k'=k/d\), \(n=dm\), and \(Q=qk'\).
On the units modulo \(Q\), expand
\begin{equation}\label{eq:A048-parseval}
 \chi(m)e(-m/Q)=\sum_{\psi\bmod Q}u_\psi\psi(m),
       \qquad \sum_\psi|u_\psi|^2=1.
\end{equation}
The last identity is orthogonality, because the original
function has modulus one. Its principal coefficient is
\begin{equation}\label{eq:A048-induced-gauss}
 u_{\psi_0}
 =\frac{\mu(k')\chi(k')\tau(\chi)\bar\chi(-1)}
        {\varphi(q)\varphi(k')}.
\end{equation}
Chinese remaindering gives the primitive Gauss sum
\(\chi(k')\bar\chi(-1)\tau(\chi)\) and the Ramanujan
sum \(\mu(k')\). This also proves the zero value when
\(k'\) is not squarefree. Since
\(\tau(\chi)\tau(\bar\chi)=\chi(-1)q\) and
\(\chi(d)\chi(k')=\chi(k)\), multiplication by the
outer scalar leaves
\(\mu(k')/(\varphi(q)\varphi(k'))\).
This proves \eqref{eq:A048-principal} with every
complex phase retained.

For the error use the exact factorization
\begin{align}
 {\cal B}_{d,\psi}(s)
 &:=
 \sum_{(m,Q)=1}\frac{b(dm)\psi(m)}{m^s}\nonumber\\
 &=\sum_{e\mid d}
   \left(\sum_{\substack{h\le y/e\\(h,Qd/e)=1}}
             \frac{a(eh)\psi(h)}{h^s}\right)
   \left(\sum_{(m,Q)=1}
             \frac{c((d/e)m)\psi(m)}{m^s}\right).
 \label{eq:A048-gcd-factorization}
\end{align}
Indeed, for a contributing divisor \(h_0\mid dn\)
put \(e=(h_0,d)\), \(h_0=eh\). Then
\((h,d/e)=1\), \(h\mid n\), and \(n=hm\).
This is a bijection and gives all displayed unit
conditions.

The second factor is its shifted \(L\)-product
times local factors at \(p\mid d/e\).
The unshifted local generating function has the form
\[
 \frac{1-p^{-z_0}X}
      {\prod_{j=1}^{J}(1-p^{-z_j}X)},\qquad J\le4,
\]
with a missing numerator when appropriate.
At \(p^v\mid d/e\), replace the series by its tail
starting at \(v\); for \(p\mid Q\), only its \(v\)-th
coefficient remains. The coefficient bound is
\(O((v+1)^3p^{C_0v/L})\).
On \(\Re s\ge1/2+\eta\), all added denominators
\(1-p^{-s-z}\) are bounded away from zero.
Products of the remaining
\(1+O(p^{-1/2-\eta/2})\) factors are \(O(d^\epsilon)\).
Lemma~\ref{lem:A048-uniform-growth} therefore bounds
the full multiplier of the finite polynomial by
\[
 \ll_{\eta,\epsilon}\tau_4(d/e)
                 \{dQ(|\Im s|+3)\}^{\epsilon}.
\]
For a principal character this bound holds on the
new left line and high horizontal edges, away from
its recorded poles.

For that polynomial, orthogonality followed by Cauchy
in each residue class gives
\begin{align}
 &\sum_{\psi\bmod Q}
 \left|\sum_{\substack{h\le y/e\\(h,Qd/e)=1}}
                  a(eh)\psi(h)h^{-s}\right|^2\nonumber\\
 &\quad\le
 \varphi(Q)\left(1+\frac{y}{eQ}\right)
                 \sum_{h\le y/e}|a(eh)|^2h^{-2\Re s}
 \ll_{\eta,r}\tau_r(e)^2(\log(2y))^{2r}(Q+y/e).
 \label{eq:A048-polynomial-mean}
\end{align}
Here \(\tau_r(eh)\le\tau_r(e)\tau_r(h)\), and
\(\sum\tau_r(h)^2h^{-1-2\eta}<\infty\).
The elementary divisor bound used here follows by
absorbing each local polynomial \(\binom{v+r-1}{r-1}\)
in \(p^{\delta v}\) for all sufficiently large primes,
and bounding the finitely many remaining factors
separately. Combining this with
\eqref{eq:A048-parseval} bounds the weighted integrand
after the sum over \(e\mid d\) by
\begin{equation}\label{eq:A048-mean-integrand}
 \{dQ(|\Im s|+3)T\}^{\epsilon}
                 \tau_{r+4}(d)\sqrt{Q+y}.
\end{equation}

Apply the elementary truncated Perron kernel proved
in NN-A046 at \(x=qUk/(2\pi d)\asymp_q Tk/d\).
Replace \(x\) by \(\lfloor x\rfloor+1/2\), and use
height \(H=T^{10}\) and initial real part
\(1+C_1/L\), beyond every shift.
Apply the truncation kernel after the finite character combination: its original multiplier has modulus one, its principal part has modulus at most one, and its nonprincipal remainder has modulus at most two. Thus no character-count factor enters this error. The bound
\[
 |b(dm)|\ll \tau_{r+4}(d)\tau_{r+4}(m)
          (dm)^{C_2/L}(\log(2dm))^{C_2}
\]
and half-integrality make the near-cutoff truncation
error \(O(d^\epsilon T^\epsilon x/H)\).
The absolute Dirichlet norm outside that range is
a fixed power of \(L\), times \(d^\epsilon\), and
gives a negligible total error as well.
Shift to \(\sigma_0=1/2+\eta\), with \(\eta\) chosen
small relative to the final error exponent.
There are no poles for nonprincipal \(\psi\);
for the principal character record the poles of
its displayed product. Equation
\eqref{eq:A048-mean-integrand}, integrated against
\(x^s/s\), bounds the remaining weighted integral
and its horizontal edges by
\begin{equation}\label{eq:A048-perron-error}
 \ll_{q,\epsilon} T^\epsilon(Tk/d)^{1/2}
                \tau_{r+4}(d)\sqrt{qk/d+y}.
\end{equation}
Preliminary exponents are chosen smaller to absorb
\(H^\epsilon\), the logarithms and \(x^\eta\).
The factor \(H^{-1}\) on horizontal edges suffices
because \(x\ll_q Ty\).
Replacing the residue cutoff by \(x\) changes only
\(T^\epsilon d^\epsilon\) times a fixed divisor factor:
the poles have bounded order, shifts \(O(1/L)\), and
residue coefficients of fixed logarithmic growth.
This is covered by \eqref{eq:A048-perron-error}.

The outer \(1/k\), all divisors, and coefficients
then give the total bound
\[
 C_{q,\epsilon}T^{1/2+\epsilon}\sqrt y
       \sum_{k\le y}k^{-1/2}|a_2(k)|
              \sum_{d\mid k}d^{-1/2}\tau_{r+4}(d)
            \ll_{q,\epsilon}T^{1/2+\epsilon}y.
\]
The small factors \(n^{-\beta},k^\beta\) modify only
the already allowed shifts. Cauchy differentiation
on slightly larger circles and endpoint partial
summation insert the asserted derivatives and
logarithmic weights. All constants are uniform for
\(T\le U\le2T\).
\end{proof}

\subsection*{Stationary phase and the three-shift right edge}

\begin{lemma}[Stationary transform with two coefficient polynomials]
\label{lem:A048-stationary}
Unfolding
\begin{align}
 {\cal K}_\chi(\alpha,\beta,\gamma;U)
 =\frac1{2\pi i}\int_{c+it_0}^{c+iU}
   \frac{L(s+\gamma,\chi)}{L(s,\chi)}
       L(s+\alpha,\chi)L(1-s+\beta,\bar\chi)
       A_\chi(s)A_{\bar\chi}(1-s)\,ds
 \label{eq:A048-K}
\end{align}
gives \eqref{eq:A048-additive}, with
\(b=(\mu*n^{-\gamma}*n^{-\alpha}*n^\beta)*a\),
and error \(O_{q,\epsilon}(T^{1/2+\epsilon}y)\).
Here \(t_0>2\pi\) is fixed and \(c=1+C_1/L\).
Two different coefficient polynomials and at most
two additional gamma or coefficient logarithms are
allowed with the same enlarged error exponent.
\end{lemma}

\begin{proof}
Use
\[
 L(1-s+\beta,\bar\chi)
       =\Delta_\chi(s-\beta)^{-1}L(s-\beta,\chi).
\]
The series converge absolutely on the chosen line.
On setting \(v=s-\beta\), the coefficient at index
\(m\) is \(b(m)m^{-\beta}\); the outer coefficient
is \(a(k)\bar\chi(k)k^{\beta-1}\).
Moving the shifted segment back to a real line and
real endpoints creates connecting segments of
length \(O(1/L)\), bounded by \(O_q(y\sqrt T L^C)\)
using Stirling and the absolute Dirichlet norms.

The Stirling phase, already derived with its root
number in NN-A037/NN-A046, has stationary point
\(t=2\pi m/(qk)\), scalar \(\tau(\bar\chi)/q\),
and phase \(e(-m/(qk))\).
Here are the coefficient sums retaining uniformity
in \(k\le y\). The relative \(O(1/t)\) Stirling
error sums to \(O_q(y\sqrt T L^C)\).
Off the factor-two stationary interval, integration
by parts and the absolute Dirichlet norm give
\(O_q(\sqrt T L^C)\) for each outer coefficient.
On the middle interval, Arias de Reyna Lemma2.5,
with the compact-amplitude constants of Lemma2.2,
gives the three errors used in NN-A046.
The first and lower-endpoint errors are bounded by
\[
 C_q k^{-1/2}\sum_{m\le C_q kT}
             \frac{\tau_R(m)}{\sqrt m}
                 \ll_{q,\epsilon}\sqrt T(kT)^\epsilon.
\]
The far upper-endpoint errors have this bound.
The near upper-endpoint terms sum to at most
\[
 C_q\sqrt T(kT)^\epsilon
       \sum_{m\asymp qkT}
        \frac1{|m-qkU/(2\pi)|+qk\sqrt T}
              \ll_{q,\epsilon}\sqrt T(kT)^\epsilon L.
\]
This estimate includes integer cutoffs. Small shift
factors are bounded on \(m\ll kT\); on infinite
tails their absolute norms are \(O(L^C)\) because
\(c\) lies beyond every shift. Summing the outer
divisor-bounded coefficients proves the error.
Endpoint partial summation with the right line
held fixed inserts gamma logarithms; shift
derivatives insert coefficient logarithms.
\end{proof}

\begin{lemma}[The explicit three-shift main term]
\label{lem:A048-three-shift}
For \(y<T^\theta\), \(\theta<1/2\), under A,
\begin{align}
 {\cal K}_\chi(\alpha,\beta,\gamma;U)
 ={}&\frac1{2\pi}\int_0^U
 \sum_{\substack{g,h,k\ge1\\(h,k)=1\\(ghk,q)=1}}
   \frac{a(gh)a(gk)}{ghk}
   \left\{\begin{aligned}
    &Z^{(q)}_{\alpha,\beta,\gamma,h,k}\\
    &+(qt/(2\pi))^{-\alpha-\beta}
              Z^{(q)}_{-\beta,-\alpha,\gamma,h,k}\\
    &+(qt/(2\pi))^{-\beta-\gamma}
              Z^{(q)}_{\alpha,-\gamma,-\beta,h,k}
   \end{aligned}\right\}dt\nonumber\\
 &\hspace{10mm}+O_{q,\epsilon}(T^{1/2+\epsilon}y).
 \label{eq:A048-three-shift}
\end{align}
The fixed lower endpoint is included in the error.
Here \(Z^{(q)}\) is exactly \eqref{eq:A047-Z} with
each zeta function replaced by \(Z_q\); the local
factors at \(p\mid k\) are unchanged.
The combined expression has removable values at
coinciding shifts.
\end{lemma}

\begin{proof}
Apply the preceding two lemmas. Unfold \(b=c*a\)
in the principal term and write the two polynomial
indices as \(gh,gk\), with \((h,k)=1\).
The inner Dirichlet series is
\begin{equation}\label{eq:A048-principal-series}
 \sum_{\substack{n\ge1\\(n,q)=1}}
    c(n)n^{-s-\beta}
       \frac{\mu(k/(n,k))}{\varphi(k/(n,k))},
       \quad c=\mu*n^{-\gamma}*n^{-\alpha}*n^\beta.
\end{equation}
Its unrestricted product is
\[
 \frac{Z_q(s)Z_q(s+\beta+\gamma)Z_q(s+\beta+\alpha)}
      {Z_q(s+\beta)}.
\]
At \(p^v\parallel k\), the local weight is zero
through exponent \(v-2\), \(-1/(p-1)\) at \(v-1\),
and one thereafter. Multiply by the ratio of this
weighted local series to its unrestricted series.
The only possible poles to the right of the new
line are \(1,1-\alpha-\beta,1-\beta-\gamma\).

We verify the first residue explicitly. Put
\[
 F_p(X)=\sum_{j\ge0}f_{\alpha,\gamma}(p^j)X^j,
 \quad C_p(X)=\frac{F_p(X)}{1-p^\beta X},
 \quad X=p^{-1-\beta}.
\]
Writing \(C_p(X)=\sum c_jX^j\), the recurrence
\(c_j=f_{\alpha,\gamma}(p^j)+p^\beta c_{j-1}\)
gives
\[
 (1-1/p)\left(\sum_{j\ge v}c_jX^j
                  -\frac{c_{v-1}X^{v-1}}{p-1}\right)
             =\sum_{j\ge v}f_{\alpha,\gamma}(p^j)X^j.
\]
Thus the local ratio, multiplied by
\(p^{v(1+\beta)}\), is exactly the factor in
\(Z^{(q)}\). The residue of \(Z_q(s)\) is \(E_q\).
The cutoff \(qUk/(2\pi h)\) and scalar
\(1/\varphi(q)\) contribute \(q/\varphi(q)\);
this cancels \(E_q\) exactly. The remaining
coefficient is \(a(gh)a(gk)/(ghk)\) times the
claimed \(Z^{(q)}\).

Under
\((\alpha,\beta,\gamma)\mapsto(-\beta,-\alpha,\gamma)\),
the coefficients of \eqref{eq:A048-principal-series}
are multiplied by \(n^{\alpha+\beta}\).
Its pole at \(1-\alpha-\beta\) becomes the first
pole of the transformed series. The identity
\(x^{1-\alpha-\beta}/(1-\alpha-\beta)
=\int_0^x t^{-\alpha-\beta}dt\), with the outer
\(h,k\) factors restored, gives the second term.
The permutation
\((\alpha,\beta,\gamma)\mapsto(\alpha,-\gamma,-\beta)\)
similarly multiplies the coefficients by
\(n^{\beta+\gamma}\) and gives the third.
These are exact meromorphic identities, first
computed for distinct shifts. The original
finite-coefficient integral is holomorphic in
the small shifts, so the combined singularities
are removable. The stated error is uniform
on the distinct circles.
\end{proof}


\subsection*{The discrete second moment and its universal kernel}

\begin{lemma}[The fixed-character discrete twisted second moment]
\label{lem:A048-discrete-second}
For real coefficients satisfying \eqref{eq:A048-coeff}
and \(y<T^\theta\), \(\theta<1/2\), the shifted sum
\[
 \sum_{T<\gamma\le2T}
 L(\rho+\alpha,\chi)L(1-\rho+\beta,\bar\chi)
             A_\chi(\rho)A_{\bar\chi}(1-\rho)
\]
equals the difference of the three main terms
\(\mathcal J_q+\mathcal I_q(\alpha,\beta)
+\mathcal I_q(\beta,\alpha)\), at \(2T\) and \(T\),
plus \(O_{q,\epsilon}(T^{1/2+\epsilon}y)\).
These terms are exactly \eqref{eq:A047-J} and
\eqref{eq:A047-I}, with the following replacements:
every zeta function becomes \(Z_q\), every coefficient
sum is restricted by \((ghk,q)=1\), every \(t/(2\pi)\)
inside an exponential or logarithm becomes \(qt/(2\pi)\),
and the local \(Z\) factors become \(Z^{(q)}\).
The same coefficient polynomial is used at both endpoints.
The assertion is uniform in small complex shifts and
retains all zero multiplicities.
\end{lemma}

\begin{proof}
Integrate
\[
 \frac{L'}L(s,\chi)
 L(s+\alpha,\chi)L(1-s+\beta,\bar\chi)
                    A_\chi(s)A_{\bar\chi}(1-s)
\]
over the standard rectangle. Choose its two heights
within one unit of \(T,2T\), avoiding zeros at distance
\(\gg_q1/\log T\). The local count and partial-fraction
argument in NN-A037 give \(L'/L=O_q(\log^2 T)\)
on both horizontal segments. The critical-line growth used
here is the fixed-character bound of NN-A039 and its checked
Iwaniec--Kowalski Corollary5.20 interface; it is not obtained
by varying the fixed distance in Lemma~\ref{lem:A048-uniform-growth}.
Under A and the reflected growth bound, throughout the strip
\[
 |L(s+\alpha,\chi)L(1-s+\beta,\bar\chi)|
             \ll_{q,\epsilon}T^{|\Re s-1/2|+\epsilon},
 \qquad
 |A_\chi(s)A_{\bar\chi}(1-s)|\ll y^{1+\epsilon}
\]
for \(0\le\Re s\le1\), with harmless bounded-factor
extensions to the chosen edges. Hence the horizontal
contributions are \(O_{q,\epsilon}(yT^{1/2+\epsilon})\).

The right edge is
\(\partial_\gamma{\cal K}_\chi(\alpha,\beta,\gamma)|_0\).
On the left use
\[
 \frac{L'}L(s,\chi)
   =\frac{\Delta_\chi'}{\Delta_\chi}(s)
                     -\frac{L'}L(1-s,\bar\chi).
\]
The second term, after reflection and conjugation,
is the conjugate right edge with shifts
\(\bar\beta,\bar\alpha\). The main term in
Lemma~\ref{lem:A048-three-shift} has real coefficients
for real shifts; conjugation therefore gives exactly
\(\mathcal I_q(\beta,\alpha)\). This is where the
reality of the proved principal component is used;
the complex character is not declared real.

The remaining left term can be moved to the critical
line. Its gamma logarithmic derivative is
\(-\log(qt/(2\pi))+O_q(1/t)\). The resulting main
integral is
\[
 \frac1{2\pi}\int_T^{2T}\log\frac{qt}{2\pi}
    L(1/2+\alpha+it,\chi)
    L(1/2+\beta-it,\bar\chi)|A_\chi(1/2+it)|^2\,dt.
\]
The \(O(1/t)\) error and its connecting edges are
within the displayed error bound.

This continuous integral is already determined by
Lemma~\ref{lem:A048-three-shift}; no separate continuous
twisted-moment theorem is needed. Set \(\gamma=0\)
in \({\cal K}_\chi\). Its quotient becomes one,
so its line moves to \(\Re s=1/2\) with the same
horizontal bound. The first two \(Z^{(q)}\) factors
become
\[
 \frac{Z_q(1+\alpha+\beta)}{h^\beta k^\alpha},
 \qquad
 \frac{Z_q(1-\alpha-\beta)}{h^{-\alpha}k^{-\beta}}.
\]
The third is zero because it contains \(1/Z_q(1)\);
first take distinct nonzero shifts and then the
removable continuation. Endpoint integration by
parts inserts \(\log(qt/(2\pi))\), with the same
power-saving error after absorbing one logarithm.
This is exactly \(\mathcal J_q\).

Finally, restoring the original heights costs at most
\(O_{q,\epsilon}(yT^\epsilon\log T)\).
There are \(O_q(\log T)\) zero occurrences in each
unit interval, each \(L\)-factor near the line is
\(O_{q,\epsilon}(T^\epsilon)\), and
\(|A_\chi(\rho)|^2\ll y^{1+\epsilon}\).
Cauchy differentiation preserves negligible errors.
All residues count zero occurrences, so multiplicities
are retained. The proof is complete.
\end{proof}

Fix the same \(P,\nu\) as in NN-A047, with \(0<\nu<1/4\),
and define
\begin{equation}\label{eq:A048-mollifier}
 B_{\chi,T}(s)=
 \sum_{n\le T^\nu}d_{-1/2}(n)\chi(n)n^{-s}
      P\left(\frac{\log(T^\nu/n)}{\nu\log T}\right).
\end{equation}
The square \(Q_{\chi,T}=B_{\chi,T}^2\) has underlying
real coefficients bounded by \(C_P\tau_2(n)\) and
length \(y=T^{2\nu}<T^{1/2}\).

\begin{corollary}[Second moment with the same constant]
\label{cor:A048-second-constant}
If \(P\) vanishes to order at least sixteen at zero,
then under A
\[
 \frac1{N_{\chi,T}L^2}
 \sum_{T<\gamma\le2T}
       |L'(\rho,\chi)B_{\chi,T}(\rho)^2|^2
                  \longrightarrow b(P,\nu),
\]
where \(b(P,\nu)\) is precisely the finite integral
\eqref{eq:A047-b-functional}.
\end{corollary}

\begin{proof}
Apply Lemma~\ref{lem:A048-discrete-second} to \(Q_{\chi,T}\)
and differentiate both shifts. Under GRH the two factors
are conjugates. The complete four-variable Mellin proof
of NN-A047 applies with \(Z_q\) in place of zeta.
Here are all the conductor-dependent changes.
At each \(p\nmid q\) the local series and local
remainder are identical to those in
\eqref{eq:A047-F-euler} and \eqref{eq:A047-G-euler};
at \(p\mid q\) the coefficient series is one and
that Euler factor is omitted from \(Z_q\).
Thus both convergent remainder products have value
one at the origin, with no missing arithmetic factor.

For fixed \(q\),
\[
 Z_q(1+z)=\frac{E_q}{z}(1+O_q(z)).
\]
In the first Euler product the net exponent of
\(Z_q\) is \(-1\), giving \(E_q^{-1}\).
The explicit pole in \(\mathcal J_q\) contributes
\(E_q\), so these cancel. In the second Euler product
the total exponent, including its prefactor, is
zero. Therefore its \(E_q\) factors also cancel.
The scaled kernels are exactly the same
\(F_0j,F_0i\) as in NN-A047.
Also \(\log(qt/(2\pi))/L=1+O_q(1/L)\), and
\(N_{\chi,T}\sim TL/(2\pi)\).

For completeness, removing finitely many Euler factors
preserves the global bounds \(O_q(L^3)\) and
\(O_q(L^6)\) on the Mellin lines. The order-sixteen
weights give the same tails
\(O_q(L^{3-8})=o(L^{-1})\) and
\(O_q(L^{6-8})=o(1)\).
On the central box the remainders tend uniformly
to one, and the same polynomial-growth majorants
are integrable. Distinct scaled Cauchy circles
therefore justify all differentiated limits.
The four inverse Laplace integrals, their domain
and every power of \(\nu\) are unchanged.
This proves the stated constant, rather than an
analogy based only on the conductor.
\end{proof}

\subsection*{The first moment from its actual principal residues}

\begin{lemma}[First moment with the same constant]
\label{lem:A048-first-constant}
For \eqref{eq:A048-mollifier}, with \(P(0)=0,P(1)=1\),
under A,
\[
 \frac1{N_{\chi,T}L}
  \sum_{T<\gamma\le2T}
       L'(\rho,\chi)B_{\chi,T}(\rho)^2
             \longrightarrow
 a(P,\nu)=\frac12+\nu I_P+
                         \frac{\nu^2}{4}(I_P^2-J_P).
\]
The integrals \(I_P,J_P\) are those in
Lemma~\ref{lem:A047-first}.
\end{lemma}

\begin{proof}
Write \(a(k)\) for the underlying coefficients of
\(Q_{\chi,T}=B_{\chi,T}^2\).
Integrate the contour integrand
\[
                    \frac{L'(s,\chi)^2Q_{\chi,T}(s)}{L(s,\chi)}.
\]
At a simple zero the residue is the stated summand;
at a multiple zero the integrand is holomorphic
and \(L'(\rho,\chi)=0\). The right-edge Dirichlet
series has no constant term and contributes only
\(O(y^{1+\epsilon}L^C)\).
The horizontal estimates are those already proved.

The exact conjugated left-edge calculation of NN-A046
gives a negative stationary transform with outer
coefficient \(a(k)\bar\chi(k)/k\). Its inner coefficient,
after the gamma logarithms are evaluated at the
stationary point, is
\[
 c_k(n)=(\log k)^2-(\log n)^2+
          (\Lambda*\log)(n).
\]
All nonprincipal components have error
\(O_{q,\epsilon}(T^{1/2+\epsilon}y)\) by
Lemma~\ref{lem:A048-character-average}, including
its logarithmic derivatives. Its principal series
is therefore controlled by
\begin{equation}\label{eq:A048-first-residue}
 W_k(s)=\sum_{\substack{n\ge1\\(n,q)=1}}
    c_k(n)\frac{\mu(k/(n,k))}{\varphi(k/(n,k))}n^{-s},
\end{equation}
with scalar \(-a(k)/(k\varphi(q))\) and cutoff
\(x=qUk/(2\pi)\). Its principal residues are real;
conjugation hence leaves them unchanged.

We compute the only residues contributing to the
leading \(L^2\) term. Put
\[
 F_k(s)=\sum_{(n,q)=1}
       \frac{\mu(k/(n,k))}{\varphi(k/(n,k))}n^{-s}
       =Z_q(s)A_k(s),
 \quad
 A_k(s)=\prod_{p^v\parallel k}
       \frac{p^{-(v-1)s}(p^{1-s}-1)}{p-1}.
 \tag{A048.14}\label{eq:A048-Ak}
\]
The formula follows by summing the geometric local
series with the three weight ranges used above.
The exact coefficient identity
\[
 (\Lambda*\log)(n)-(\log n)^2
             =-\sum_{p^j\mid n}j(\log p)^2
\]
follows from
\(Z_q'^2/Z_q-Z_q''=-Z_q(Z_q'/Z_q)'\).
The restriction \((n,q)=1\) is understood in this
Dirichlet-series identity. Locally its insertion is
\((X^2d^2/dX^2+2X\,d/dX)/2\), since it multiplies
the \(j\)-th coefficient by \(j(j+1)/2\).
Consequently \(W_k\) is \( (\log k)^2F_k\) minus
the sum of these local insertions, weighted by
\((\log p)^2\).

Let \(z=s-1\), and write
\(Z_q(s)=E_q/z+C_q+O_q(z)\).
For \(k=1\),
\[
 W_1(s)=-Z_q(s)(Z_q'/Z_q)'(s)
             =-E_qz^{-3}+O_q(z^{-2}).
\]
For a prime \(p\nmid q\), putting \(\ell_p=\log p\),
direct differentiation of the local geometric series
in \eqref{eq:A048-Ak} gives
\begin{align}
 W_p(s)={}&
 \frac{E_q\ell_p}{p-1}z^{-2}\nonumber\\
 &+\left\{-E_q\left(
       \frac{3\ell_p^2}{2(p-1)}
            +\frac{\ell_p^2}{(p-1)^2}\right)
       +\frac{C_q\ell_p}{p-1}\right\}z^{-1}
       +O_{p,q}(1).
 \label{eq:A048-prime-pole}
\end{align}
One may verify the only additional local term as follows.
For \(v=1\) the weighted local series is
\((pX-1)/((p-1)(1-X))\), and at \(X=1/p\)
its normalized local insertion has residue
\(-p/(\ell_p(p-1))\) in the \(s\) variable.
Multiplying by \(\ell_p^2F_p(s)\), and including
the quadratic coefficient of \(A_p(s)\), yields
\eqref{eq:A048-prime-pole}.
For distinct primes \(p,r\nmid q\),
\[
 W_{pr}(s)=
 -\frac{E_q\ell_p\ell_r}{(p-1)(r-1)}z^{-1}+O_{p,r,q}(1).
\]
Every \(A_k\) has a zero of order \(\omega(k)\)
at \(1\). The unrestricted insertion adds a double
pole; a modified local insertion adds at most a
simple pole after its own local zero is divided out.
Hence \(W_k\) is holomorphic at \(1\) when
\(\omega(k)\ge3\).

Prime powers and the Euler corrections contribute
only \(O_{q,P}(L)\) after the common factor \(U/(2\pi)\).
To see this uniformly, the single-prime-power residues
are bounded by
\[
 C|a(p^v)|p^{-v}
             \{Lv^2\log p+v^3(\log p)^2\},
                    \qquad v\ge2.
\]
This follows by differentiating the displayed geometric
factor, including \(p^{-(v-1)s}\).
For two primes the bound is
\[
 C|a(p^vr^w)|p^{-v}r^{-w}
                    v^2w^2\log p\log r.
\]
If \(v\ge2\) or \(w\ge2\), one prime sum converges
absolutely and the other is \(O(L)\).
The bounds \(|a(k)|\le C_P\tau_2(k)\) suffice.
Replacing \(p-1\) by \(p\) in the simple-prime terms
also costs \(O_{q,P}(L)\), because the difference
introduces a summable extra \(1/p\).

Taking the residues of \(W_k(s)x^s/s\), with its
negative scalar, now gives the same leading expression
as for zeta:
\[
 \frac U{2\pi}\left\{\frac12(\log U)^2+
 \sum_p\frac{a(p)}p
       \left(-\log U\log p+\frac12(\log p)^2\right)
 +\sum_{p<r}\frac{a(pr)}{pr}\log p\log r
 +O_{q,P}(L)\right\}.
 \tag{A048.15}\label{eq:A048-first-leading}
\]
Here \(qE_q/\varphi(q)=1\) cancels the conductor
scalar, and \(\log(q/(2\pi))\) contributes only lower
order terms. Omitting the finitely many primes
dividing \(q\) changes this expression by \(O_{q,P}(L)\).
The prime and semiprime coefficients of the square
are exactly those computed in NN-A047:
\[
 a(p)=-P(1-u),\qquad
 a(pr)=\frac12\{P(1-u-v)\mathbf1_{u+v\le1}
                              +P(1-u)P(1-v)\}.
\]
Thus the same prime-number-theorem integrals give
\(a(P,\nu)\). Evaluate \eqref{eq:A048-first-leading}
at \(2T,T\) with the same polynomial. Replacing
\(\log U\) by \(L\) changes only \(O_{q,P}(TL)\).
Since \(N_{\chi,T}\sim TL/(2\pi)\), the result follows.
\end{proof}

\subsection*{The positive-weight signed mean}

\begin{lemma}[Reality of the bilinear derivative main term]
\label{lem:A048-real-bilinear}
Let \(X_\chi,Y_{\bar\chi}\) have real underlying
coefficients satisfying \eqref{eq:A048-coeff}, both
of length at most \(y\le T^\theta\), \(\theta<1/2\).
Under A,
\[
 \Im\sum_{T<\gamma\le2T}
       L'(\rho,\chi)X_\chi(\rho)Y_{\bar\chi}(1-\rho)
       =O_{q,\epsilon}(T^{1/2+\epsilon}y).
\]
\end{lemma}

\begin{proof}
Integrate \(L'^2/L\) times the two polynomials over
the same rectangle. The residues are exactly the
displayed sum, with every multiple-center term zero.
The horizontal contribution is
\(O_{q,\epsilon}(T^{1/2+\epsilon}y)\).
On the right edge the diagonal condition is \(nh=k\).
Its character factor is
\(\chi(n)\chi(h)\bar\chi(k)=1\) for its nonzero terms,
and every remaining coefficient is real.
The off-diagonal integral is \(O(y^{1+\epsilon}L^C)\):
combine the two right-side coefficients into index
\(m=nh\), split \(m\) outside and inside \([k/2,2k]\),
and use respectively absolute Dirichlet convergence
and \(|\log(m/k)|\gg|m-k|/k\).
The harmonic distance sum is logarithmic, uniformly
for \(k\le y\).

After left reflection and conjugation, the expression
is the exact one in NN-A046, with outer polynomial
\(X_{\bar\chi}(1-s)\) and inner polynomial \(Y_\chi(s)\).
The three terms are
\[
 \Delta_\chi(s)^{-1}
 \left\{D_\chi(s)^2L(s,\chi)
       -2D_\chi(s)L'(s,\chi)
             +L'(s,\chi)^2/L(s,\chi)\right\},
 \qquad D_\chi=\Delta_\chi'/\Delta_\chi .
\]
Replace \(D_\chi\) by \(-\log(qt/(2\pi))\);
its \(O_q(1/t)\) error is within the same bound.
Lemma~\ref{lem:A048-stationary} applies to the
remaining convolution series and gamma logarithms.
Their principal components are real by the exact
scalar \eqref{eq:A048-principal}. Indeed at a
stationary point the gamma logarithm becomes
\(\log(m/k)\), which is real, and the underlying
coefficients are real. The needed derivatives of
shifted products are uniform by
Lemma~\ref{lem:A048-character-average}; their
principal derivatives remain real.
All nonprincipal parts have the stated error.
Conjugation of this real principal term leaves it
real. Restoring the endpoints uses the same unit
interval estimate as above, with \(L'\ll T^\epsilon\).
This proves the lemma.
\end{proof}

\begin{lemma}[The square-weight cancellation]
\label{lem:A048-square-signed}
For a real-underlying-coefficient \(B_{\chi,T}\) of
fixed length \(T^\nu\), \(\nu<1/2\), suppose
\[
 \sum_{T<\gamma\le2T}
       |L'(\rho,\chi)B_{\chi,T}(\rho)^2|^2
                              =O_\chi(N_{\chi,T}L^2).
\]
Put
\[
 K_{\rho,\chi}
  =\frac{\Re\{\rho L(\rho+1,\chi)
                         \overline{L'(\rho,\chi)}\}}{\gamma}.
\]
Then under A
\[
 \sum_{T<\gamma\le2T}
       K_{\rho,\chi}|B_{\chi,T}(\rho)|^2
                                    =o_\chi(N_{\chi,T}L).
\]
\end{lemma}

\begin{proof}
Choose fixed \(\delta>0\) with \(\nu+\delta<1/2\)
and truncate the exterior series at \(M=T^\delta\).
Set
\[
 X_\chi=B_{\chi,T},\qquad
 Y_{\bar\chi}(s)=B_{\bar\chi,T}(s)
                  \sum_{m\le M}\bar\chi(m)m^{-1-s}.
\]
Its underlying coefficients are real, uniformly
divisor-bounded, and supported through \(T^{\nu+\delta}\).
Lemma~\ref{lem:A048-real-bilinear} therefore gives
a negligible imaginary part for the truncated sum.
Under GRH,
\(B_{\bar\chi,T}(1-\rho)=\overline{B_{\chi,T}(\rho)}\).
The exterior tail is at most \(2M^{-1/2}\), so restoring
it costs at most
\[
 2M^{-1/2}\sum |L'(\rho,\chi)|\,|B_{\chi,T}(\rho)|^2
 \le2M^{-1/2}
     \left(N_{\chi,T}\sum|L'(\rho,\chi)B_{\chi,T}(\rho)^2|^2\right)^{1/2}
                   =o_\chi(N_{\chi,T}L).
\]
If \(L'(\rho,\chi)\overline{L(\rho+1,\chi)}=u+iv\),
then exactly \(K_{\rho,\chi}=v+u/(2\gamma)\).
The latter term, after weighting, is
\(O_\chi(N_{\chi,T}L/T)\) by Cauchy and the exterior
Euler bound. This proves the assertion.
\end{proof}

\subsection*{Completion of the counting argument}

Take the single fixed rational \(\nu<1/4\) and fixed
high-order polynomial \(P_K\) constructed in NN-A047.
They can be chosen independently of the fixed character,
because the limiting constants just proved are exactly
\(a(P_K,\nu)\) and \(b(P_K,\nu)\). In particular,
\[
 a\ge a_*=\frac{1543}{2304}-\frac1{2000},\qquad
 b\le b_*=
 \frac{1505515846753}
      {536481792000(314159/100000)}+\frac1{2000}.
\]
The displayed rational upper bound also satisfies \(b_*<1\).
The first moment and the triangle inequality give the
required absolute first-moment lower bound. The second
moment gives square-weight cancellation by the preceding
lemma.

The corrected phase \(\psi_\chi\) of NN-A037 satisfies
\[
 \sum_{T<\gamma\le2T}\cos^2\psi_\chi(\gamma)
                     =(1/2+o_\chi(1))N_{\chi,T}.
\]
This is the proved fixed-character arithmetic mean,
not phase independence. Therefore the phase-bad set
\(|\cos\psi_\chi|<\eta\) has at most
\((p_\eta+o_\chi(1))N_{\chi,T}\) occurrences, where
\(p_\eta=1/(2(1-\eta^2))\).
The proof of Proposition~\ref{prop:A042-square-mollifier}
now applies verbatim to these actual moments and
signed weights: its only other bounds are
\[
 \zeta(3/2)^{-1}\le |L(3/2+it,\chi)|\le\zeta(3/2),
       \qquad |\rho|/\gamma=1+o(1),
\]
which hold uniformly for every Dirichlet Euler product.
At every multiple center the weighted statistic is zero.
At a nonzero weighted summand its sign is precisely the
strict ordinary velocity sign, and that center is simple.

With \(\eta=1/50\), the already checked exact rational
inequality gives
\[
 a_*^2-b_*p_\eta>\frac1{2000},\qquad
 a_*-\sqrt{b_*p_\eta}>\frac1{4000}.
\]
Each sign therefore has lower proportion greater than
\[
 \frac{(1/50)^2(1/4000)^2}{4\cdot3^4}
            =\frac1{12960000000000}>10^{-14}.
\]
The height threshold may depend on \(q,\chi\), including
all finite Euler-factor and contour constants. The
fixed polynomial and the numerical lower fraction do
not. This proves Theorem~\ref{thm:A048-signs}.

\paragraph{Geometric translation.}
The local analytic zero algebra is reduced exactly
where the specified shifted divisor has its unique
unramified local section; the finite relative-differentials
criterion is SGA1, ExposeI, section3, Proposition1.
The normal tangent is paired with
\(|L'(\rho,\chi)|^2/\gamma\), and the multiplier adds
the Hermitian norm \(|B_{\chi,T}(\rho)|^2\).
The statistic vanishes on all nonreduced centers.
The finite character decomposition on
\((\mathbb Z/(qk)\mathbb Z)^*\) is an actual
representation-theoretic step: its principal Gauss
scalar becomes real, while character orthogonality
controls the entire coefficient polynomial before
the arithmetic sum is bounded. This produces the
positive counting proportions in both oriented
reduced loci. It is a real analytic norm theorem,
not an algebraic ordering, an arithmetic deformation
scheme, a Frobenius half-label law or a pure motive.
The auxiliary GRHs are genuinely used for growing
moduli, although the counted character is fixed.

\paragraph{Source and proof scope.}
The new uniform growth, additive averaging, residues,
and weighted signed mean are proved here.
The local stationary-phase interface is Arias de Reyna,
arXiv1403.7126v1, Lemmas2.2 and2.5, with the detailed
coefficient summation supplied in this section:
\url{https://arxiv.org/html/1403.7126v1}.
The primitive completion and functional equation are
DLMF25.15.5--6; the gamma estimates are DLMF5.11.
The finite Gauss calculation is proved by Chinese
remaindering above; its classical locator is
Montgomery--Vaughan, section9.2, Theorems9.7 and9.10.
The comparison blueprint is Heap--Li--Zhao v1,
section4, equations35--52, whose relevant section
was read in full; its zeta theorem is not used to
assert the fixed-character formula.
The complete NN-A047 kernel proof and exact rational
certificate are reused after the conductor factors
have been checked, with no new numerical experiment.
The exact geometric locator is
\url{https://grothendiecksga.com/read/sga1/fr/I-3.html}.




% Full root review 2026-11-22 04:01:36 UTC; peer base/refinement reviews separately logged.
\section{A common deformation interval for positive proportions of actual zeros}
\label{sec:A049-common-time}

The sign proportions proved in NN-A047--NN-A048 concern the
first normal derivative. This section gives a single positive
parameter interval on which positive proportions of the actual
analytic branches keep their two respective sides.

\begin{theorem}[NN-A049: positive proportions on a common deformation interval]
\label{thm:A049-common-time}
Fix a primitive Dirichlet character \(\chi\), with conductor
one meaning zeta. Assume RH in that case, and A for the other
fixed characters. For every fixed \(\epsilon>0\), and every
sufficiently large \(T\), there are two disjoint sets of
simple zero centers in \(T<\Im\rho<2T\), each of size at least
\[
                         10^{-16}N_{\chi,T},
 \qquad N_{\chi,T}=N_\chi(2T)-N_\chi(T),
\]
whose branches for the specified shift
\begin{equation}\label{eq:A049-shift}
 F_\chi(s,\tau)=L(s,\chi)+
   \sum_{k\ge1}\frac{(-\tau)^k(s)_k}{k!}L(s+k,\chi)
\end{equation}
are distinct and simple throughout
\[
                 0\le\tau\le T^{-3/2-\epsilon}.
\]
Every branch in the first set satisfies
\(\Re s(\tau)>1/2\), and every branch in the second
satisfies \(\Re s(\tau)<1/2\), simultaneously for every
\(0<\tau\le T^{-3/2-\epsilon}\). Their ordinates remain
in \((T,2T)\). The same centers work for the whole interval.

More precisely, if \(\nu<1/4\) is the single fixed mollifier
exponent constructed in NN-A047, the upper endpoint can be
\(T^{-1-2\nu-\epsilon}\). The constant in the counting
proportion is absolute; the eventual height threshold can
depend on the fixed character and on \(\epsilon\).
No claim is made for a fixed nonzero deformation parameter,
a height-uniform conductor family, or a limiting half-label
law.
\end{theorem}

\subsection*{A robust two-sided sample}

Retain the fixed \(P,\nu\) of NN-A047--NN-A048 and write
\(B=B_{\chi,T}\), \(D_\rho=L'(\rho,\chi)\),
\[
 A_\rho=D_\rho B(\rho)^2,\qquad
 w_\rho=|B(\rho)|^2,\qquad
 K_\rho=\frac{\Re\{\rho L(\rho+1,\chi)\overline{D_\rho}\}}
                   {\gamma}.
\]
All moment sums below count every zero occurrence. A nonzero
\(K_\rho w_\rho\) is supported on a simple center.

\begin{lemma}[Robust signed sets]\label{lem:A049-robust}
For every fixed \(\delta>0\), and all sufficiently large \(T\),
each strict sign has at least \(10^{-16}N_{\chi,T}\)
distinct simple centers satisfying
\begin{equation}\label{eq:A049-robust-conditions}
 T+1<\gamma<2T-1,\qquad
 |D_\rho|\ge T^{-\nu-\delta},\qquad
 \frac{|\Re v_\rho|}{|v_\rho|}\ge10^{-9},
 \qquad v_\rho=\frac{\rho L(\rho+1,\chi)}{D_\rho}.
\end{equation}
The two sets retain their respective sign of \(\Re v_\rho\).
\end{lemma}

\begin{proof}
The proved moment constants obey \(b_*<1\), and with
\(\eta_0=1/50\) their phase-discard margin exceeds \(1/4000\).
The proof of Proposition~\ref{prop:A042-square-mollifier},
applied to the actual moments of NN-A047 or NN-A048, gives
\[
 \sum |K_\rho|w_\rho
     \ge\left(\frac1{600000}-o(1)\right)N_{\chi,T}\log T,
 \qquad
 \sum K_\rho w_\rho=o(N_{\chi,T}\log T).
\]
Here \(1/\zeta(3/2)>1/3\) is the same elementary bound used
there. Consequently each signed absolute mass is at least
\[
                    \frac1{2000000}N_{\chi,T}\log T
\]
eventually. Also the actual second moment gives
\begin{equation}\label{eq:A049-second-bound}
 \sum |A_\rho|^2\le N_{\chi,T}(\log T)^2
\end{equation}
for sufficiently large \(T\).

On this height interval, \(|\rho|/\gamma<2\) and
\(|L(\rho+1,\chi)|<3\). The exact phase identity therefore
gives
\[
 |K_\rho|w_\rho\le6|A_\rho|\,|\cos\psi_\chi(\gamma)|.
\]
Thus deleting the centers with
\(|\cos\psi_\chi|<10^{-9}\) removes total absolute mass at
most \(6\cdot10^{-9}N_{\chi,T}\log T\), by Cauchy and
\eqref{eq:A049-second-bound}. No independence is used.

The coefficients of \((1-z)^{1/2}\) have absolute value
at most one, so multiplicativity gives
\(|d_{-1/2}(n)|\le1\). Therefore, directly from the
defining polynomial,
\[
 |B(1/2+i\gamma)|^2
 \le \|P\|_\infty^2
        \left(\sum_{n\le T^\nu}n^{-1/2}\right)^2
 \le4\|P\|_\infty^2T^\nu .
\]
At centers with \(|D_\rho|<T^{-\nu-\delta}\), the total
absolute signed weight is accordingly
\(O_P(N_{\chi,T}T^{-\delta})\), which is
\(o(N_{\chi,T}\log T)\). This uses a pointwise polynomial
bound, not an inverse-derivative moment.

There are \(O_\chi(\log T)\) occurrences in the two discarded
unit intervals at the ends. Their total absolute weight is
at most
\[
 6\left(O_\chi(\log T)\sum|A_\rho|^2\right)^{1/2}
                         =o(N_{\chi,T}\log T).
\]
Thus each sign retains absolute weight at least
\(\kappa N_{\chi,T}\log T\), where
\(\kappa=1/4000000\), once \(T\) is sufficiently large.
For either retained set \(S\), Cauchy's inequality gives
\[
 \kappa^2 N_{\chi,T}^2(\log T)^2
 \le \#S\sum (K_\rho w_\rho)^2
 \le36\,\#S\,N_{\chi,T}(\log T)^2.
\]
It follows that
\[
 \#S\ge\frac{N_{\chi,T}}{576000000000000}
                           >10^{-16}N_{\chi,T}.
\]
All retained summands are nonzero and have the requested
strict sign. The exact phase identity makes
\(|\Re v_\rho|/|v_\rho|=|\cos\psi_\chi(\gamma)|\).
This proves every condition in \eqref{eq:A049-robust-conditions}.
\end{proof}

\subsection*{Uniform local estimates for the specified shift}

\begin{lemma}[Local deformation bounds]\label{lem:A049-local}
Fix \(\delta>0\) and the character. There is a fixed
\(r_0>0\) such that, at all sufficiently high critical-line
centers \(\rho\) with \(T<\gamma<2T\),
\[
 \max_{|s-\rho|\le2r_0}|L''(s,\chi)|\le T^\delta.
\]
Writing \(G(s)=sL(s+1,\chi)\), uniformly on these discs,
\[
 cT\le |G(\rho)|\le CT,\qquad |G'(s)|\le CT,
\]
with absolute positive constants after \(r_0<1/100\).
For \(|\tau|T\le1/10\), the series in
\eqref{eq:A049-shift} converges locally normally and satisfies
\begin{equation}\label{eq:A049-local-expansion}
 F_\chi(s,\tau)=L(s,\chi)-\tau G(s)+R(s,\tau),\qquad
 |R(s,\tau)|+|\partial_sR(s,\tau)|\le C_\delta T^2|\tau|^2
\end{equation}
on \(|s-\rho|\le r_0\).
\end{lemma}

\begin{proof}
Choose \(r_0<\min(\delta/100,1/100)\). The fixed-character
GRH growth estimate and its reflected left-half version,
proved in NN-A039 and used in NN-A041, give
\[
 \max_{|s-\rho|\le4r_0}|L(s,\chi)|
       \ll_{\chi,\delta}T^{4r_0+\delta/4}.
\]
Cauchy's formula on a fixed \(r_0\)-disc, and absorption
of its fixed constant into a further small power of \(T\),
give the displayed \(L''\) bound. The possible zeta pole
is far outside the high discs.

Since \(\Re(s+1)>1.4\), the absolute Dirichlet series for
\(L(s+1,\chi)\) and its first derivative are bounded by
absolute constants. At the center the Euler-product lower
bound is \(1/\zeta(3/2)>1/3\). These prove the \(G\) bounds.

For \(k\ge2\), the same absolute series gives
\(|L(s+k,\chi)|\le\zeta(2.4)\). If \(u=|s|\), then
\[
 |(s)_k|\le(u)_k,\qquad
 \sum_{k\ge0}\frac{(u)_k|\tau|^k}{k!}
                         =(1-|\tau|)^{-u}.
\]
Here \(u\le3T\), and \(T|\tau|\le1/10\).
The terms of order at least two are \(O(T^2|\tau|^2)\),
uniformly on the slightly larger discs. This proves local
normal convergence and the remainder bound. Cauchy's
formula in \(s\) gives the derivative bound on the smaller
disc. All constants are independent of the center.
\end{proof}

\begin{proof}[Proof of Theorem~\ref{thm:A049-common-time}]
It suffices to prove the sharper endpoint in terms of
\(\nu\). Put
\[
 \delta=\min(\epsilon/16,1/100),\qquad
 r_T=T^{-\nu-3\delta},\qquad
 \tau_T=T^{-1-2\nu-\epsilon}.
\]
Take the robust sets of Lemma~\ref{lem:A049-robust}.
For each selected center let \(D=D_\rho\). On
\(|s-\rho|\le r_T\), Taylor's formula and
Lemma~\ref{lem:A049-local} give
\[
 L(s,\chi)=D(s-\rho)+E(s),\qquad
 |E(s)|\le\tfrac12 T^\delta|s-\rho|^2.
\]
Since \(|D|\ge T^{-\nu-\delta}\), the error is \(o\) of
the linear term on the boundary. The disc contains exactly
one zero of \(L\), counted with multiplicity.

These discs are pairwise disjoint. Indeed, if two selected
centers had distance at most \(4r_T\), Taylor's formula at
one of them, evaluated at the other, would give
\[
 |D|\le2T^\delta r_T=2T^{-\nu-2\delta},
\]
contradicting \(|D|\ge T^{-\nu-\delta}\) eventually.

On each boundary, uniformly for \(|\tau|\le\tau_T\),
\[
 \frac{|F_\chi(s,\tau)-L(s,\chi)|}{|D|r_T}
 \ll \frac{\tau_T T}{T^{-2\nu-4\delta}}
                         =T^{-\epsilon+4\delta}=o(1).
\]
Rouche's theorem therefore gives exactly one shifted zero
in every disc for all such complex \(\tau\). Moreover
\[
 |\partial_sF_\chi(s,\tau)-D|
 \le T^\delta r_T+C|\tau|T+C_\delta|\tau|^2T^2=o(|D|).
\]
Thus the zero is simple throughout. Local implicit functions
patch by uniqueness to a single holomorphic branch
\(s_\rho(\tau)=\rho+w_\rho(\tau)\) on the full parameter disc.

We check its direction uniformly, rather than only its
derivative at zero. The equation
\(F_\chi(\rho+w,\tau)=0\), together with
\(|w|\le r_T\), first gives
\[
 |w|\le3|\tau|\,|G(\rho)|/|D|
                          =3|\tau v_\rho|
\]
for sufficiently large \(T\): the terms
\(\tfrac12T^\delta r_T|w|+C|\tau|T|w|\) are absorbed
into \(|D||w|\), and the residual
\(O(T^2|\tau|^2)\) is \(o(|\tau G(\rho)|)\).
Substituting this bound back into the equation yields
\[
 \frac{|w-\tau v_\rho|}{|\tau v_\rho|}
 \ll_\delta
      \frac{T^\delta|\tau|T}{|D|^2}
        +\frac{|\tau|T}{|D|}+|\tau|T
 \ll_\delta
      T^{-\epsilon+3\delta}
       +T^{-\nu-\epsilon+\delta}
       +T^{-2\nu-\epsilon}=o(1).
\]
This is uniform for \(0<|\tau|\le\tau_T\) and all selected
centers. For positive real \(\tau\), eventually the relative
error is less than \(10^{-9}/2\). Since the robust set has
\(|\Re v_\rho|\ge10^{-9}|v_\rho|\), every branch has the
same strict normal sign as its initial velocity throughout
the entire positive interval.

The center ordinates have distance at least one from both
ends, and \(r_T<1\), so all shifted ordinates stay in
\((T,2T)\). The disjoint discs retain distinct branches.
Finally \(2\nu<1/2\), so
\(T^{-3/2-\epsilon}\le\tau_T\). This proves the asserted
common interval and both counting proportions.
\end{proof}

\paragraph{Geometric translation.}
Over the parameter disc, each selected local zero divisor
is finite of degree one and has vanishing relative
differentials. The exact finite unramified criterion is
SGA1, Expose I, section3, Proposition1. Its unique section
extends over a common explicitly shrinking disc, and the
sections have disjoint analytic neighborhoods. The specified
real ray carries their two normal orientations throughout
that common interval, not merely in their first formal jets.
The positive Hermitian multiplier is used only to select
enough reduced centers; the final assertion counts the actual
unweighted branches. This remains an analytic deformation
of the specified shifted Dirichlet series. It does not
supply an arithmetic family, an ordering of complex schemes,
a Frobenius sign law, or a fixed positive time valid at every
height.

\paragraph{Inputs and scope.}
The moment and signed-weight estimates are NN-A047 for zeta
and NN-A048 for the other fixed primitive characters.
The fixed-character growth/Cauchy interface was proved in
NN-A039--NN-A041; the zero-count local bound is the exact
one retained in NN-A037. The deformation series and its
velocity were proved in NN-A004. All new disc, tail,
selection and uniform direction estimates are proved above.
No new mean-value theorem, simplicity premise, numerical
experiment or growing-conductor estimate is imported.

\subsection*{Refinement to the inverse-height scale}

\begin{lemma}[Elementary sampling of the actual mollifier]
\label{lem:A049-sampling}
For the fixed polynomial \(B=B_{\chi,T}\) above,
\[
             \sum_{T<\gamma\le2T}|B(\rho)|^2
                          \ll_{\chi,P,\nu}T(\log T)^3.
\]
The sum retains every zero occurrence and requires no
spacing or simplicity assumption.
\end{lemma}

\begin{proof}
Put \(Y=T^\nu\), and write
\(B(1/2+it)=\sum_{n\le Y}c_n n^{-it}\). The bound already
proved gives \(|c_n|\le C_Pn^{-1/2}\).
For any real interval \(J\) of length \(T+O(1)\), expand the
square and integrate each exponential. The diagonal is
\(O_P(T\log(2Y))\). The absolute off-diagonal contribution
is at most
\[
 C_P\sum_{1\le n<m\le Y}
             \frac1{\sqrt{mn}\log(m/n)}.
\]
For \(m\le2n\), use
\(\log(m/n)\ge(m-n)/m\); each term is at most
\(\sqrt2/(m-n)\), whose double sum is \(O(Y\log(2Y))\).
For \(m>2n\), the logarithm is at least \(\log2\), so the
double sum is at most
\(C(\sum_{n\le Y}n^{-1/2})^2=O(Y)\).
Thus
\[
 \int_J |B(1/2+it)|^2dt
                       \ll_P (T+Y)\log(2Y).
\]
The derivative with respect to \(t\) inserts \(-i\log n\).
Repeating the same calculation, or bounding both logarithms
in the off-diagonal sum by \(\log Y\), gives
\[
 \int_J \left|\frac{d}{dt}B(1/2+it)\right|^2dt
                       \ll_P (T+Y)(\log(2Y))^3.
\]

For every \(C^1\) complex function \(g\) on an interval
\(I\) of length one and each \(t_0\in I\),
\[
 |g(t_0)|^2\le
       \int_I|g(t)|^2dt+2\int_I|g(t)g'(t)|dt.
\]
Indeed integrate
\(|g(t_0)|^2\le |g(t)|^2+\int_I|(|g|^2)'(u)|du\)
over \(t\in I\).
Every unit interval in \([T-1,2T+1]\) contains
\(O_\chi(\log T)\) zero occurrences. Sum the last bound,
one copy per occurrence, over the covering unit intervals.
Cauchy's inequality for the two integrals, and \(Y<T\),
give
\[
 \begin{aligned}
 \sum_{T<\gamma\le2T}|B(\rho)|^2
 &\ll_\chi\log T
    \left(\int_{T-1}^{2T+1}|B|^2+
              2\left(\int|B|^2\int|B'|^2\right)^{1/2}\right)\\
 &\ll_{\chi,P,\nu}T(\log T)^3.
 \end{aligned}
\]
Assign any integer-height endpoint to just one adjacent
interval. Multiplicities are retained in the local count.
\end{proof}

\begin{corollary}[NN-A049: a common interval of inverse-height size]
\label{cor:A049-inverse-height}
Under the hypotheses of Theorem~\ref{thm:A049-common-time},
the same conclusion, with the same lower proportion
\(10^{-16}\), holds simultaneously for every
\[
                        0<\tau\le T^{-1-\epsilon}.
\]
The selected centers are allowed to depend on
\(\epsilon\) and \(T\), but not on \(\tau\) within that
whole interval.
\end{corollary}

\begin{proof}
The only use of the exponent \(\nu\) in the selection
lemma was the pointwise bound on the square weight.
For every fixed \(\delta>0\), the absolute signed weight
of centers with \(|D_\rho|<T^{-\delta}\) is instead at most
\[
 6T^{-\delta}\sum_{T<\gamma\le2T}|B(\rho)|^2
      \ll_{\chi,P,\nu}T^{1-\delta}(\log T)^3
                           =o(N_{\chi,T}\log T)
\]
by Lemma~\ref{lem:A049-sampling}. All other deletions and
constants are unchanged. Hence both robust sets still
have at least \(10^{-16}N_{\chi,T}\) centers, now with
\[
              |D_\rho|\ge T^{-\delta},\qquad
              |\Re v_\rho|\ge10^{-9}|v_\rho|.
\]

For precision, repeat the disc estimates with
\[
 \delta=\min(\epsilon/16,1/100),\qquad
 r_T=T^{-3\delta},\qquad \tau_T=T^{-1-\epsilon}.
\]
The curvature ratio on the unshifted disc is
\(T^\delta r_T/|D_\rho|\le T^{-\delta}\).
Two selected centers at distance at most \(4r_T\)
would give \(|D_\rho|\le2T^{-2\delta}\), a contradiction.
The forcing ratio is
\[
 \frac{\tau_T T}{|D_\rho|r_T}
                    \le T^{-\epsilon+4\delta}=o(1),
\]
and the perturbed derivative differs from \(D_\rho\)
by \(o(|D_\rho|)\). Thus the same disjoint discs carry
one simple holomorphic branch on the full complex
parameter disc.

Finally the already proved error formula is now bounded by
\[
 \frac{|s_\rho(\tau)-\rho-\tau v_\rho|}
      {|\tau v_\rho|}
 \ll_\delta T^{-\epsilon+3\delta}
             +T^{-\epsilon+\delta}+T^{-\epsilon}=o(1),
 \qquad 0<|\tau|\le T^{-1-\epsilon}.
\]
It is uniform in the selected centers. The positive-ray
sign, ordinate, distinctness and count conclusions follow
exactly as in the theorem. The geometric conclusion is
therefore a common analytic parameter disc of the sharper
size; no new arithmetic or probabilistic premise is used.
\end{proof}




% A049 refinement: root recorded PASS 2026-11-22 04:09:31 UTC; GCT full PASS 04:06:06 UTC.
\section{A logarithmic refinement of the common deformation interval}
\label{sec:A049-logarithmic-time}

This refinement uses the exact unweighted small-shift second
moment already available in NN-A047--NN-A048. Its source and
multiplicity interface are stated explicitly below.

\begin{lemma}[Unweighted small shifts and a local curvature mean]
\label{lem:A049-curvature-mean}
Under the hypotheses of NN-A049, put \(L=\log T\) and
\(R=(10L)^{-1}\). Uniformly for \(|z|\le R\),
\begin{equation}\label{eq:A049-small-shift-mean}
 \sum_{T<\gamma\le2T}|L(\rho+z,\chi)|^2
                                  \ll_\chi N_{\chi,T}L.
\end{equation}
If
\[
 M_\rho=\max_{|w|\le R/2}|L''(\rho+w,\chi)|,
\]
then, retaining all zero occurrences,
\begin{equation}\label{eq:A049-curvature-mean}
                     \sum M_\rho^2\ll_\chi N_{\chi,T}L^5.
\end{equation}
\end{lemma}

\begin{proof}
For \(q>1\), use Lemma~\ref{lem:A048-discrete-second}
with the coefficient sequence \(a(1)=1\), \(a(n)=0\)
for \(n>1\). This is the fully proved fixed-character
formula, whose coefficients now have length one.
For zeta use the exact imported Heap--Li--Zhao v1
Theorem2, equations5--9, with that same coefficient
sequence, as stated in NN-A047. That formula is
unconditional and counts every zero with multiplicity;
the simplicity premise of its separate negative-moment
theorem is not used. RH is used here for the conjugation
and critical-line location, as it was in NN-A047.

We bound the actual main terms, including their removable
values. Write \(Z_q(s)=\zeta(s)\prod_{p\mid q}(1-p^{-s})\),
with \(Z_1=\zeta\). For \(g=h=k=1\) the three-shift factor
in both formulas is precisely
\[
 Z^{(q)}_{\alpha,\beta,\gamma,1,1}
   =\frac{Z_q(1+\alpha+\beta)Z_q(1+\beta+\gamma)}
           {Z_q(1+\beta)}.
\]
Choose the outer shift circles
\[
                  |\alpha|=2R,\quad|\beta|=3R,\quad|\gamma|=4R.
\]
Every linear combination in the three permuted products
that is a potential pole argument has modulus between
\(R\) and \(8R\). The denominator arguments have moduli
\(2R,3R\), or \(4R\). Since
\[
 Z_q(1+u)=\frac{E_q}{u}(1+O_q(u)),
\]
each of the three products is \(O_q(L)\), uniformly on
these circles. The accompanying powers of \(qt/(2\pi)\)
are bounded on \(T\le t\le2T\), since all shifts are
\(O(1/L)\). The combined three-term expression is
holomorphic in \(\gamma\): at \(\gamma=-\beta\) the first
and third residues cancel, and at \(\gamma=\alpha\) the
second and third residues cancel, with their powers then
equal. There are no other poles in the small disc.
Cauchy's formula on the \(\gamma\) circle thus bounds
the derivative of the combined expression by \(O_q(L^2)\).

The continuous term \(\mathcal J_q\) has integrand
\[
 \log\frac{qt}{2\pi}
 \left\{Z_q(1+\alpha+\beta)
 +(qt/(2\pi))^{-\alpha-\beta}Z_q(1-\alpha-\beta)\right\},
\]
which is likewise \(O_q(L^2)\) on the outer
\(\alpha,\beta\) circles. Thus the dyadic shifted discrete
sum has modulus \(O_q(TL^2)\) there. For \(q>1\) its error is
\(O_{q,\epsilon}(T^{1/2+\epsilon})\), with \(\epsilon<1/2\).
For zeta the imported error is
\(O_A(T(\log T)^{-A})\) for every fixed \(A>0\);
this is also \(o(TL^2)\). The finite sum itself
is holomorphic in both shifts, so successive maximum-modulus
bounds extend this estimate to
\(|\alpha|,|\beta|\le R\), including coinciding shifts.
Since \(N_{\chi,T}\asymp_\chi TL\), put
\(\alpha=z,\beta=\bar z\). Under GRH,
\[
 L(1-\rho+\bar z,\bar\chi)
                         =\overline{L(\rho+z,\chi)},
\]
and \eqref{eq:A049-small-shift-mean} follows.

For \(|w|\le R/2\), Cauchy's formula on \(|z|=R\)
gives the uniform bound
\[
 |L''(\rho+w,\chi)|
 \le16R^{-2}\frac1{2\pi}\int_0^{2\pi}
                    |L(\rho+Re^{i\theta},\chi)|\,d\theta .
\]
Squaring, applying Cauchy to the angular average, and
then summing the finite zero-occurrence set gives
\[
 \sum M_\rho^2
 \le256R^{-4}\frac1{2\pi}\int_0^{2\pi}
           \sum|L(\rho+Re^{i\theta},\chi)|^2\,d\theta
 \ll_\chi N_{\chi,T}L^5.
\]
No selection of simple zeros or exchange of an infinite
sum is involved.
\end{proof}

\begin{corollary}[NN-A049: a common logarithmic inverse-height interval]
\label{cor:A049-logarithmic-time}
Under the hypotheses of NN-A049, all sufficiently large
\(T\) have two sets of distinct simple centers, each of
size at least \(10^{-16}N_{\chi,T}\), whose actual shifted
branches remain distinct and simple, stay in the ordinate
window \((T,2T)\), and keep their respective strict sides
of the critical line simultaneously for every
\[
                         0<\tau\le\frac1{T(\log T)^{10}}.
\]
The same two sets work on that whole interval. The
constant is absolute, and the eventual height threshold
can depend on the fixed character. For zeta RH alone
suffices; the other fixed characters use A as in NN-A048.
\end{corollary}

\begin{proof}
Retain the notation and signed-mass constants from
Lemma~\ref{lem:A049-robust}. Its sampling refinement gives
\[
                 \sum|B(\rho)|^2\ll_{\chi,P,\nu}TL^3
                                               \ll_{\chi,P,\nu}N_{\chi,T}L^2.
\]
Consequently the total absolute signed weight on
\(|D_\rho|<L^{-2}\) is
\[
       \sum |K_\rho|w_\rho
       \le6L^{-2}\sum|B(\rho)|^2
                         =O_{\chi,P,\nu}(N_{\chi,T})
                         =o(N_{\chi,T}L).
\]
By \eqref{eq:A049-curvature-mean}, the set with
\(M_\rho>L^4\) has at most \(O_\chi(N_{\chi,T}L^{-3})\)
occurrences. Cauchy and
\(\sum|A_\rho|^2\le N_{\chi,T}L^2\) show that its
absolute signed weight is at most
\[
 6\left(O_\chi(N_{\chi,T}L^{-3})
                       N_{\chi,T}L^2\right)^{1/2}
                              =o(N_{\chi,T}L).
\]
Delete these two sets, the same phase-margin exceptional
set and the same endpoint unit intervals as before.
Each sign still retains mass at least
\(\kappa N_{\chi,T}L\), with \(\kappa=1/4000000\),
and hence at least \(10^{-16}N_{\chi,T}\) centers.
Every retained center now has
\[
 |D_\rho|\ge L^{-2},\qquad
 M_\rho\le L^4,\qquad
 |\Re v_\rho|\ge10^{-9}|v_\rho|.
\]

Put \(r_T=L^{-7}\) and \(\tau_T=(TL^{10})^{-1}\).
Eventually \(4r_T<R/2\). Taylor's formula at every selected
center has curvature at most \(L^4\) on that whole disc.
Two selected centers at distance at most \(4r_T\) would
force
\[
                  |D_\rho|\le2L^4r_T=2L^{-3}<L^{-2},
\]
a contradiction. Thus the discs are disjoint.
The curvature and forcing ratios are, respectively,
\[
 \frac{M_\rho r_T}{|D_\rho|}\le L^{-1},\qquad
 \frac{\tau_T T}{|D_\rho|r_T}\le L^{-1}.
\]
The local deformation series and its absolute remainder
bound from Lemma~\ref{lem:A049-local} still apply: these
bounds use only absolutely convergent shifted \(L\)-series
and are uniform on the present smaller discs.
They therefore give one simple holomorphic section
throughout \(|\tau|\le\tau_T\), exactly as in NN-A049.

Finally the explicit relative-direction estimate there is
now
\[
 \frac{|s_\rho(\tau)-\rho-\tau v_\rho|}
      {|\tau v_\rho|}
 \ll \frac{M_\rho|\tau|T}{|D_\rho|^2}
             +\frac{|\tau|T}{|D_\rho|}+|\tau|T
 \ll L^{-2}+L^{-8}+L^{-10}=o(1).
\]
This is uniform in all selected centers and
\(0<|\tau|\le\tau_T\). Eventually it is smaller than
\(10^{-9}/2\). The same sign, distinctness, ordinate and
count conclusions now hold on the full stated positive
interval.
\end{proof}

\paragraph{Exact interface and geometric scope.}
For zeta the only imported mean-value input to this
refinement is Heap--Li--Zhao v1, Theorem2, equations5--9,
PDFpp3--5, with coefficient sequence \(\delta_1\);
the unconditional and multiplicity conventions were
checked in NN-A047. The official primary copy is
\url{https://pure.mpg.de/rest/items/item_3368312_2/component/file_3368314/content}.
For the other characters the required formula was proved
locally in Lemma~\ref{lem:A048-discrete-second}; no
conductor extrapolation is made. The refinement enlarges
the common base disc of the same degree-one analytic
covers in NN-A049. The SGA1 relative-differentials
dictionary and all its stated arithmetic limitations
are unchanged. A fixed positive parameter, the full
critical scale \(\tau\asymp T^{-1}\), half balance,
and conductor-uniformity remain open.



% Reviewed parameter refinement within NN-A049; root 04:09:31 UTC, GCT 04:09:37 UTC.
\subsection*{The same mean estimates give every logarithmic exponent above \(9/2\)}

\begin{corollary}[NN-A049: logarithmic interval with exponent \(9/2+\epsilon\)]
\label{cor:A049-logarithmic-sharp}
For every fixed \(\epsilon>0\), the conclusion of
Corollary~\ref{cor:A049-logarithmic-time}, with the same
absolute lower proportion \(10^{-16}\), holds
simultaneously for every
\[
                    0<\tau\le\frac1{T(\log T)^{9/2+\epsilon}}.
\]
The hypotheses and all fixed-character and geometric
limitations are unchanged.
\end{corollary}

\begin{proof}
Put \(L=\log T\) and \(\delta=\epsilon/5\).
The two established mean bounds are
\[
 \sum|B(\rho)|^2\ll_{\chi,P,\nu}N_{\chi,T}L^2,\qquad
 \sum M_\rho^2\ll_\chi N_{\chi,T}L^5.
\]
The total absolute signed weight with
\(|D_\rho|<L^{-1-\delta}\) is at most
\[
                  6L^{-1-\delta}\sum|B(\rho)|^2
                          \ll_{\chi,P,\nu}N_{\chi,T}L^{1-\delta}
                          =o(N_{\chi,T}L).
\]
The set \(M_\rho>L^{5/2+\delta}\) has at most
\(O_\chi(N_{\chi,T}L^{-2\delta})\) occurrences. Its
absolute signed weight is, by the same Cauchy bound,
\(O_\chi(N_{\chi,T}L^{1-\delta})=o(N_{\chi,T}L)\).
Deleting those sets, the phase-margin exception and the
two endpoint unit intervals retains both absolute signed
masses at least \(\kappa N_{\chi,T}L\).
The exact count deduction is unchanged. Thus each sign
has at least \(10^{-16}N_{\chi,T}\) centers with
\[
 |D_\rho|\ge L^{-1-\delta},\qquad
 M_\rho\le L^{5/2+\delta},\qquad
 |\Re v_\rho|\ge10^{-9}|v_\rho|.
\]

Use
\[
 r_T=L^{-7/2-3\delta},\qquad
 \tau_T=T^{-1}L^{-9/2-5\delta}.
\]
Eventually \(4r_T<(20L)^{-1}\), the disc on which the
curvature mean was defined. The curvature and forcing
ratios satisfy
\[
 \frac{M_\rho r_T}{|D_\rho|}\le L^{-\delta},\qquad
 \frac{\tau_TT}{|D_\rho|r_T}\le L^{-\delta}.
\]
Two selected centers within \(4r_T\) would give
\(|D_\rho|\le2M_\rho r_T< L^{-1-\delta}\)
eventually. The discs are therefore disjoint, and the
same Rouche and derivative estimates give one simple
holomorphic section in each disc throughout
\(|\tau|\le\tau_T\). The three relative-direction errors
are explicitly bounded by
\[
 \frac{M_\rho|\tau|T}{|D_\rho|^2}
        +\frac{|\tau|T}{|D_\rho|}+|\tau|T
 \le L^{-2\delta}+L^{-7/2-4\delta}+L^{-9/2-5\delta}
                                      =o(1).
\]
They are uniform in the selected centers and in the
whole punctured parameter disc. The fixed phase margin
therefore preserves both strict sides for every positive
\(\tau\) in that disc, and all ordinates remain in the
original height window. Since \(5\delta=\epsilon\), this
is the stated interval. This changes no source input
and remains a refinement of NN-A049, not a new finding.
\end{proof}



\subsection*{Using the actual mollifier coefficients improves the exponent to three}

\begin{lemma}[The fractional coefficient norm and zero sampling]
\label{lem:A049-fractionalnorm}
Let \(d_{-1/2}(n)\) be the coefficients of \(\zeta(s)^{-1/2}\).
For \(Y\ge e\),
\begin{equation}\label{eq:A049-fractionalnorm}
 \sum_{n\le Y}\frac{|d_{-1/2}(n)|^2}{n}
                         \le e^3(1+\log Y)^{1/4}.
\end{equation}
Consequently the actual fixed mollifier \(B=B_{\chi,T}\) of
NN-A047--A048 satisfies, with every zero occurrence retained,
\begin{equation}\label{eq:A049-fractionalsampling}
             \sum_{T<\gamma\le2T}|B(\rho)|^2
                  \ll_{\chi,P,\nu}N_{\chi,T}(\log T)^{5/4}.
\end{equation}
Here the zero-location and fixed-character hypotheses are the
same as in NN-A049; no spacing or simplicity is assumed.
\end{lemma}

\begin{proof}
The coefficients are multiplicative, and locally come from
\((1-z)^{1/2}\). They have absolute value at most one, with
\(d_{-1/2}(p)=-1/2\). For real \(\sigma>1\), put
\(x=p^{-\sigma}\le1/2\). The local nonnegative square series
is bounded by
\[
 \sum_{k\ge0}|d_{-1/2}(p^k)|^2x^k
       \le1+\frac x4+\frac{x^2}{1-x}
       \le\exp(x/4+2x^2).
\]
Multiplying the absolutely convergent Euler products, and using
\(\sum_p p^{-2\sigma}<1\) and
\(\sum_p p^{-\sigma}\le\log\zeta(\sigma)\), gives
\[
 \sum_{n\ge1}\frac{|d_{-1/2}(n)|^2}{n^\sigma}
                               \le e^2\zeta(\sigma)^{1/4}.
\]
Both elementary inequalities follow respectively by comparison
with \(\sum_{n\ge2}n^{-2}<\int_1^\infty u^{-2}du=1\),
and the logarithm of the Euler product for \(\zeta\).
Choose \(\sigma=1+1/\log Y\). Nonnegativity, the bound
\(Y^{\sigma-1}=e\), and
\(\zeta(\sigma)\le1+\int_1^\infty u^{-\sigma}du
=1+\log Y\) prove \eqref{eq:A049-fractionalnorm}.

For clarity we track this improvement through the full sampling
estimate. Write \(L=\log T\), \(Y=T^\nu\) with the fixed
\(0<\nu<1/4\), and
\[
 B(1/2+it)=\sum_{n\le Y}c_n n^{-it},\qquad
 |c_n|\le C_P|d_{-1/2}(n)|n^{-1/2}.
\]
The character factor has absolute value at most one. Thus
\(\sum|c_n|^2\ll_P L^{1/4}\).
On every real interval \(J\) of length \(T+O(1)\), expansion
of the square gives diagonal contribution
\(O_P(TL^{1/4})\). The absolute off-diagonal contribution is
at most a constant depending on \(P\) times
\[
 \sum_{n<m\le Y}\frac1{\sqrt{mn}\log(m/n)}
                                          \ll Y\log(2Y).
\]
Indeed for \(m\le2n\) each summand is at most
\(\sqrt2/(m-n)\), whose sum is \(O(Y\log(2Y))\).
For \(m>2n\), use \(\log(m/n)\ge\log2\) and
\((\sum_{n\le Y}n^{-1/2})^2=O(Y)\).
Differentiation with respect to \(t\) inserts \(-i\log n\),
so the differentiated diagonal is \(O_P(TL^{9/4})\) and
the absolute off-diagonal is \(O_P(YL^3)\).
Since \(\nu<1\) is fixed, both off-diagonal bounds are
eventually smaller than the corresponding diagonal bounds.
We have proved
\[
                 \int_J|B|^2\ll_P TL^{1/4},\qquad
                 \int_J|B'|^2\ll_P TL^{9/4}.
\]
Use the elementary unit-interval inequality proved in
Lemma~\ref{lem:A049-sampling},
\[
 |g(t_0)|^2\le\int_I|g|^2+2\int_I|gg'|,
\]
and the same multiplicity-aware \(O_\chi(L)\) local zero
count. Summing over unit intervals and applying Cauchy's
inequality gives
\[
 \sum_{T<\gamma\le2T}|B(\rho)|^2
       \ll_{\chi,P,\nu}
          L\left(TL^{1/4}+
                       (TL^{1/4}\cdot TL^{9/4})^{1/2}\right)
       \ll_{\chi,P,\nu}TL^{9/4}.
\]
Finally \(N_{\chi,T}\asymp_\chi TL\) proves
\eqref{eq:A049-fractionalsampling}.
\end{proof}

\begin{corollary}[NN-A049: a common interval with exponent \(3+\epsilon\)]
\label{cor:A049-logarithmic-three}
For every fixed \(\epsilon>0\), all conclusions of
Corollary~\ref{cor:A049-logarithmic-time}, including the same
absolute lower fraction \(10^{-16}\) for each strict side,
hold simultaneously for every
\[
                         0<\tau\le
                           \frac1{T(\log T)^{3+\epsilon}}.
\]
The same selected centers work throughout the interval.
Zeta requires RH alone, and every fixed primitive character
has the already stated A scope. No growing-conductor or
fixed-positive-time assertion is added.
\end{corollary}

\begin{proof}
Put \(L=\log T\), \(\delta=\epsilon/5\), and retain the
notation \(D_\rho\), \(M_\rho\) and signed weights of the
preceding NN-A049 proofs. The two mean bounds now are
\[
 \sum|B(\rho)|^2\ll N_{\chi,T}L^{5/4},\qquad
 \sum M_\rho^2\ll N_{\chi,T}L^5.
\]
By the same pointwise signed-weight bound, the total absolute
weight of centers with \(|D_\rho|<L^{-1/4-\delta}\) is at most
\[
                 6L^{-1/4-\delta}\sum|B(\rho)|^2
                       \ll N_{\chi,T}L^{1-\delta}
                         =o(N_{\chi,T}L).
\]
The curvature exceptions \(M_\rho>L^{5/2+\delta}\) have
\(O(N_{\chi,T}L^{-2\delta})\) occurrences. The previously
proved second moment of the signed weights and Cauchy's
inequality bound their absolute weight by
\(O(N_{\chi,T}L^{1-\delta})\) as well.
After these deletions, the fixed phase-margin exception and
the two endpoint unit intervals, each sign still has the
absolute mass and counting conclusion used in NN-A049.
In particular at least \(10^{-16}N_{\chi,T}\) centers of
each sign satisfy
\[
 |D_\rho|\ge L^{-1/4-\delta},\quad
 M_\rho\le L^{5/2+\delta},\quad
 |\Re v_\rho|\ge10^{-9}|v_\rho|.
\]

Choose
\[
 r_T=L^{-11/4-3\delta},\qquad
 \tau_T=T^{-1}L^{-3-5\delta}.
\]
Eventually \(4r_T<(20L)^{-1}\), inside the previously
defined curvature disk. The exact ratios are
\[
 \frac{M_\rho r_T}{|D_\rho|}\le L^{-\delta},\qquad
 \frac{\tau_TT}{|D_\rho|r_T}\le L^{-\delta}.
\]
If two selected centers were within \(4r_T\), Taylor's
formula and the curvature bound would give
\(|D_\rho|\le2M_\rho r_T< L^{-1/4-\delta}\), a contradiction
for large \(T\). Hence the disks are disjoint.
The same locally convergent denominator-shift deformation,
Rouche estimates and derivative bound give one simple
holomorphic section in each disk for all \(|\tau|\le\tau_T\).
The uniform relative-direction errors are explicitly at most
\[
 \frac{M_\rho|\tau|T}{|D_\rho|^2}
       +\frac{|\tau|T}{|D_\rho|}+|\tau|T
   \le L^{-2\delta}+L^{-11/4-4\delta}+L^{-3-5\delta}=o(1).
\]
The fixed phase margin therefore preserves both strict sides
throughout the entire positive interval, and the endpoint
deletion keeps every ordinate in \((T,2T)\).
Since \(5\delta=\epsilon\), this is the asserted scale.
\end{proof}

\paragraph{Geometric scope.}
The same positive fractions of reduced centers of the analytic
denominator-shift divisor extend to disjoint unramified sections
over the larger common parameter disk
\(|\tau|\le1/[T(\log T)^{3+\epsilon}]\).
Their normal real directions retain their signs on its positive
radius. The refinement uses an elementary coefficient norm and
the already reviewed curvature and signed-weight moments; it
does not change the underlying geometry or identify this
denominator-shift parameter with de Bruijn--Newman heat time.
The full scale \(\tau\asymp1/T\) remains a separate gate.


\subsection*{Direct rational-twist sampling}

\begin{lemma}[A rational-twist contour with all multiplicities]
\label{lem:A049-rational-landau}
Fix a primitive Dirichlet character \(\chi\), with the conductor-one
case denoting zeta, and assume its GRH. Put \(L=\log T\), and fix
\(0<\nu<1/4\). For all sufficiently large \(T\), there are
\[
 U\in[T,T+1],\qquad V\in[2T,2T+1]
\]
independent of \(x\), such that for every \(1<x\le T^\nu\),
\begin{equation}\label{eq:A049-rational-landau}
 \begin{aligned}
 \sum_{U<\gamma<V}x^\rho
   &=-\frac{V-U}{2\pi}\Lambda_\chi(x)\\
   &\quad+O_\chi\!\left(
      xL^2+\frac{x\log(2x)}{\operatorname{dist}_*(x,\mathbb Z)}
                         +\frac{L}{\log x}\right).
 \end{aligned}
\end{equation}
The sum retains every zero occurrence. Here
\(\Lambda_\chi(k)=\chi(k)\Lambda(k)\) on positive integers and
is zero otherwise, and
\[
 \operatorname{dist}_*(x,\mathbb Z)
     =\min_{\substack{k\in\mathbb Z_{\ge1}\\ k\ne x}}|x-k|.
\]
In particular this distance equals one at an integer \(x>1\).
All implied constants are uniform in the displayed range of \(x\).
The heights \(U,V\) are not zero ordinates.
\end{lemma}

\begin{proof}
We first justify the same two heights for every twist.
The fixed-character local count already proved and sourced in
NN-A037--A041 bounds the number of zero occurrences in any unit
interval near height \(T\) by \(C_\chi L\).
The exact previously audited classical locators are
Iwaniec--Kowalski, Proposition~5.7(1), equation (5.27),
pp.\ 102--103, and Theorem~5.8 with its proof (5.34)--(5.36),
pp.\ 104--105; negative ordinates follow by conjugating the
character. These fixed-object inputs do not grow with the twists.
In each of \([T,T+1]\) and \([2T,2T+1]\), remove the intervals
of radius \(\kappa_\chi/L\) about all ordinates within distance
one of that interval. Taking a sufficiently small fixed
\(\kappa_\chi>0\) makes the removed measure less than \(1/2\).
Choose \(U,V\) in the respective complements. Both are at
distance at least \(\kappa_\chi/L\) from every zero ordinate.
This uses the occurrence count, so multiple zeros cause no
exception.

For these heights and \(-1\le\sigma\le2\), we claim
\begin{equation}\label{eq:A049-good-logderivative}
 \left|\frac{L'}{L}(\sigma+it,\chi)\right|\ll_\chi L^2,
                              \qquad t\in\{U,V\}.
\end{equation}
Here and below \(L'/L\) denotes the logarithmic derivative of
the Dirichlet function, not the abbreviation \(L=\log T\).
To check the claim directly, subtract its value at \(2+it\).
The genus-one Hadamard logarithmic derivative of the entire
completion gives the absolutely convergent zero sum
\[
 \sum_\rho\left(\frac1{\sigma+it-\rho}
                         -\frac1{2+it-\rho}\right).
\]
The exponential constant and the \(1/\rho\) summands cancel.
The gamma factors and the optional zeta factors \(s(s-1)\)
contribute \(O_\chi(L)\) on this fixed strip; this follows
from the already used fixed-strip Stirling formula.
There are \(O_\chi(L)\) occurrences with \(|t-\gamma|\le1\);
each corresponding difference is \(O_\chi(L)\) at the
chosen heights. The remaining terms are
\(O((t-\gamma)^{-2})\). Summing them by unit intervals,
using \(O_\chi(\log(2+|u|))\) occurrences near height \(u\),
gives \(O_\chi(L)\); the convergent tail
\(\sum_{j\ge1}\log(j+2)/j^2\) handles \(|u|\) far from \(T\).
Finally \(L'/L(2+it,\chi)\) is bounded by its absolutely
convergent Euler series. This proves
\eqref{eq:A049-good-logderivative}.
The entire-completion/Hadamard interface is exactly the one
proved and used in NN-A035 and NN-A039, with the same
fixed-character local count; no new uniform-conductor
assertion is being imported.

Set \(c=1+1/L\) and integrate
\[
                  \frac1{2\pi i}\frac{L'}{L}(s,\chi)x^s
\]
around the positively oriented rectangle with vertical edges
\(\Re s=c,1-c\) and horizontal edges \(\Im s=U,V\).
Its poles are precisely the nontrivial zeros inside, with
their multiplicities; the possible pole of zeta and all
trivial zeros have imaginary part zero.
On either horizontal edge,
\eqref{eq:A049-good-logderivative} and
\(x^c\le e^\nu x\) give the bound \(O_\chi(xL^2)\).

On the right edge use the absolutely convergent expansion
\[
             \frac{L'}{L}(s,\chi)
                   =-\sum_{n\ge2}\frac{\chi(n)\Lambda(n)}{n^s}.
\]
If \(x\) is an integer, the term \(n=x\) contributes
\(-(V-U)\Lambda_\chi(x)/(2\pi)\); otherwise no term is
stationary. Every other term has integral bounded in
absolute value by
\[
              \frac{2}{|\log(x/n)|}\,(x/n)^c\Lambda(n).
\]
For \(n\notin[x/2,2x]\), its sum is \(O(xL^2)\):
use \(|\log(x/n)|\ge\log2\), \(x^c\ll x\), and
\[
 \sum_{n\ge2}\frac{\Lambda(n)}{n^c}
 \le \sum_{n\ge2}\frac{\log n}{n^c}\ll(c-1)^{-2}\ll L^2.
\]
The last estimate follows by the integral comparison for
\(\sum(\log n)n^{-c}\), uniformly for \(1<c\le2\).
For the remaining nonstationary terms,
\[
 (x/n)^c\le4,\qquad
 \frac1{|\log(x/n)|}\le\frac{2x}{|x-n|},\qquad
 \Lambda(n)\le\log(2x).
\]
There are at most two nearest integers to \(x\). After removing
those, the reciprocal distances are bounded by a harmonic
sum of length \(O(x)\). Thus
\[
 \sum_{\substack{x/2\le n\le2x\\ n\ne x}}\frac1{|x-n|}
       \ll \log(2x)+\operatorname{dist}_*(x,\mathbb Z)^{-1}.
\]
This proves the right-edge error in
\eqref{eq:A049-rational-landau}, since \(\log(2x)\ll L\).

It remains to bound the left edge, including its gamma term.
Write the primitive functional equation as
\[
 L(s,\chi)=\Delta_\chi(s)L(1-s,\bar\chi),\qquad
 \Delta_\chi(s)=\epsilon_\chi
    (q/\pi)^{1/2-s}
    \frac{\Gamma((1-s+e_\chi)/2)}{\Gamma((s+e_\chi)/2)}.
\]
For zeta use \(q=1,e_\chi=0,\epsilon_\chi=1\).
The convention is DLMF 25.15.5--25.15.6, already checked
for NN-A051; the root number disappears upon logarithmic
differentiation. DLMF 5.11.2 supplies the fixed-strip digamma
asymptotic. On \(s=1-c+it\), it gives
\[
 \frac{L'}{L}(s,\chi)
   =-\log\frac{qt}{2\pi}+O_\chi(t^{-1})
                   -\frac{L'}{L}(1-s,\bar\chi).
\]
The integral of the first term against \(e^{it\log x}\) is
\(O_\chi(L/\log x)\), by one integration by parts.
The \(O(t^{-1})\) term has integral \(O_\chi(1)\) on
\([U,V]\). The prefactor \(x^{1-c}\) is at most one.
For the remaining term, absolute convergence at \(1-s=c-it\)
and integration of each exponential give a bound
\[
 \sum_{n\ge2}\frac{\Lambda(n)}{n^c\log(nx)}
      \le \sum_{n\ge2}\frac1{n^c}\ll L.
\]
Indeed \(\log(nx)\ge\log n\) and \(\Lambda(n)\le\log n\).
This last bound and the \(O_\chi(1)\) term are absorbed by
\(O_\chi(xL^2)\). The orientations change no error estimate,
and the residue theorem now proves the displayed formula.
\end{proof}

\begin{lemma}[Direct sampling of the actual fractional mollifier]
\label{lem:A049-directsampling}
Under the hypotheses of NN-A049, the same actual polynomial
\(B=B_{\chi,T}\) satisfies
\begin{equation}\label{eq:A049-directsampling}
              \sum_{T<\gamma\le2T}|B(\rho)|^2
                    \ll_{\chi,P,\nu}N_{\chi,T}(\log T)^{1/4}.
\end{equation}
Every zero occurrence, including every multiple one, is retained.
No assumption of zero spacing is used.
\end{lemma}

\begin{proof}
Write \(Y=T^\nu\), \(L=\log T\), and
\[
 B(s)=\sum_{n\le Y}b_n n^{-s},\qquad
 |b_n|\le C_P a(n),\qquad a(n)=|d_{-1/2}(n)|.
\]
First sum over the good heights \(U,V\) supplied by
Lemma~\ref{lem:A049-rational-landau}.
On GRH, \(\rho=1/2+i\gamma\), so expansion of the finite
square gives the diagonal
\[
       N_\chi(U,V)\sum_{n\le Y}\frac{|b_n|^2}{n}
\]
and twice the real part of the terms with \(m>n\).
For those terms put \(x=m/n\). Dividing
\eqref{eq:A049-rational-landau} by \(\sqrt x\) gives the
sum of \(x^{i\gamma}\); reversing its sign in the exponential
just conjugates that sum. The absolute coefficient multiplying
the unscaled error in \eqref{eq:A049-rational-landau} is at
most
\[
                    \frac{C_P^2}{\sqrt{mn}\sqrt{m/n}}
                              =\frac{C_P^2}{m}.
\]
There is no unstated uniformity in this finite twist summation.
For \(x=m/n>1\) we have
\[
 \operatorname{dist}_*(m/n,\mathbb Z)\ge\frac1n,\qquad
                 \log(m/n)\ge\frac{m-n}{m}.
\]
Consequently the three summed errors are, respectively,
\begin{align*}
 \sum_{n<m\le Y}\frac{xL^2}{m}
     &\ll YL^3,\\
 \sum_{n<m\le Y}
       \frac{x\log(2x)}{m\,\operatorname{dist}_*(x,\mathbb Z)}
     &\ll Y^2L,\\
 \sum_{n<m\le Y}\frac{L}{m\log(m/n)}
     &\ll YL^2.
\end{align*}
Here the first and third use the elementary harmonic sum,
and the second uses at most \(Y^2/2\) terms each bounded
by \(\log(2Y)\).
Their sum is \(O_P(Y^2L+YL^3)=o(T)\), since the fixed
\(\nu<1/4\).

Only ratios \(m/n=r\) that are integer prime powers contribute
to the main off-diagonal term. Its absolute value is bounded by
\[
 C_P^2 T\sum_{r\le Y}\frac{\Lambda(r)}r
                     \sum_{n\le Y/r}\frac{a(nr)a(n)}n.
\]
The local coefficients \(a(p^j)\) are nonincreasing in \(j\).
Indeed \(a(1)=1,a(p)=1/2\), and for \(j\ge1\)
\[
             \frac{a(p^{j+1})}{a(p^j)}
                           =\frac{j-1/2}{j+1}<1.
\]
Thus \(a(np^k)\le a(n)\), without a coprimality restriction.
For every contributing \(r\), the inner sum is at most
\[
                 S(Y):=\sum_{n\le Y}\frac{a(n)^2}{n}
                                   \ll(\log Y)^{1/4}
\]
by Lemma~\ref{lem:A049-fractionalnorm}.
The elementary factorial identity, with \(M=\lfloor Y\rfloor\),
gives
\[
 \log(M!)=\sum_{r\le M}\Lambda(r)\lfloor M/r\rfloor
       \ge\frac M2\sum_{r\le M}\frac{\Lambda(r)}r.
\]
Since \(\log(M!)\le M\log M\), we have
\(\sum_{r\le Y}\Lambda(r)/r\le2\log Y\).
The main off-diagonal is therefore \(O_P(TL^{5/4})\).
The diagonal is \(O_{\chi,P}(N_\chi(U,V)L^{1/4})\),
of the same size because \(N_\chi(U,V)\asymp_\chi TL\).

Finally replacing \(U,V\) by \(T,2T\) changes at most
\(O_\chi(L)\) zero occurrences. On the critical line,
\[
              |B(\rho)|\le C_P\sum_{n\le Y}n^{-1/2}
                                      \le2C_P\sqrt Y.
\]
The resulting endpoint error is \(O_{\chi,P}(YL)=o(T)\).
All estimates counted occurrences, so this also covers
multiple zeros or an endpoint that is itself a zero.
This proves \eqref{eq:A049-directsampling}.
\end{proof}

\begin{corollary}[NN-A049: a common interval with exponent \(1+\epsilon\)]
\label{cor:A049-logarithmic-one}
For every fixed \(\epsilon>0\), the positive actual branch
proportions and all other conclusions of NN-A049 hold
simultaneously throughout
\[
                         0<\tau\le
                           \frac1{T(\log T)^{1+\epsilon}}.
\]
Each side has the same absolute lower proportion \(10^{-16}\),
with the selected centers fixed throughout the interval.
The scopes remain RH alone for zeta and A for every other
fixed primitive character. The full scale \(\tau\asymp1/T\)
and a fixed nonzero time at arbitrarily large heights are
not conclusions of this refinement.
\end{corollary}

\begin{proof}
Put \(L=\log T\), \(\delta=\epsilon/5\), and use the same
signed weights, derivative \(D_\rho\), and local curvature
\(M_\rho\) as in the previous NN-A049 proofs.
The new sample bound and unchanged curvature mean are
\[
 \sum|B(\rho)|^2\ll N_{\chi,T}L^{1/4},\qquad
 \sum M_\rho^2\ll N_{\chi,T}L^5.
\]
The absolute signed weight of centers with
\(|D_\rho|<L^{3/4-\delta}\) is at most
\[
                 6L^{3/4-\delta}\sum|B(\rho)|^2
                       \ll N_{\chi,T}L^{1-\delta}
                         =o(N_{\chi,T}L).
\]
As previously, the set \(M_\rho>L^{5/2+\delta}\)
has \(O(N_{\chi,T}L^{-2\delta})\) occurrences and,
by the proved signed-weight second moment and Cauchy,
absolute signed weight \(O(N_{\chi,T}L^{1-\delta})\).
Deleting these two sets, the fixed phase-margin exception
and the two endpoint unit intervals therefore retains
the same signed masses and the same counting lower bound.
Each side has at least \(10^{-16}N_{\chi,T}\) centers with
\[
 |D_\rho|\ge L^{3/4-\delta},\quad
 M_\rho\le L^{5/2+\delta},\quad
 |\Re v_\rho|\ge10^{-9}|v_\rho|.
\]

Choose
\[
 r_T=L^{-7/4-3\delta},\qquad
 \tau_T=T^{-1}L^{-1-5\delta}.
\]
The disks remain inside the previously defined curvature
neighborhood, since \(4r_T<(20L)^{-1}\) eventually.
The two Rouche ratios are exactly bounded by
\[
 \frac{M_\rho r_T}{|D_\rho|}\le L^{-\delta},\qquad
 \frac{\tau_TT}{|D_\rho|r_T}\le L^{-\delta}.
\]
Two selected centers within \(4r_T\) would, by Taylor's
formula, force
\(|D_\rho|\le2M_\rho r_T<L^{3/4-\delta}\) for large \(T\).
Thus the disks are disjoint. The already proved convergent
denominator-shift expansion and its derivative estimate
give one simple holomorphic section in every disk,
simultaneously for \(|\tau|\le\tau_T\).
Their relative-direction errors are uniformly bounded by
\[
 \frac{M_\rho|\tau|T}{|D_\rho|^2}
       +\frac{|\tau|T}{|D_\rho|}+|\tau|T
   \le L^{-2\delta}+L^{-7/4-4\delta}+L^{-1-5\delta}=o(1).
\]
The fixed phase margin keeps each section strictly on its
prescribed side for every \(0<\tau\le\tau_T\).
The endpoint deletion keeps every ordinate in \((T,2T)\).
Finally \(5\delta=\epsilon\) gives the stated interval.
\end{proof}

\paragraph{Geometric scope.}
The same positive fractions of reduced centers of the actual
denominator-shift divisor extend to disjoint unramified
sections over the larger common disk
\(|\tau|\le1/[T(\log T)^{1+\epsilon}]\).
The rational-twist contour counts the full initial divisor,
including its nonreduced points, before the signed masses
select the reduced centers. The refined sampling changes
the analytic base radius, not the arithmetic object or the
meaning of its oriented real parameter ray.
It gives no algebraization, no heat-time identification and
no inference of an exact half-label law.


\subsection*{Weighted curvature and a fixed quarter-logarithmic scale}

\begin{lemma}[The undifferentiated square-mollifier moment]
\label{lem:A049-weightedcurvature}
Use the single fixed \(P,\nu\) chosen in NN-A047--A048, and
write \(B=B_{\chi,T}\), \(L=\log T\), \(N=N_{\chi,T}\).
There are absolute constants \(C_Q,C_M>0\) such that, for
every fixed primitive character in the hypotheses of NN-A049,
and all sufficiently large \(T\) depending on that character,
\begin{equation}\label{eq:A049-weightedvalue}
 \sum_{T<\gamma\le2T}|B(\rho)|^4|L(\rho+z,\chi)|^2
                  \le C_QN,\qquad |z|\le R:=\frac1{1000L}.
\end{equation}
Consequently, for
\[
             M_\rho=\max_{|w|\le R/2}|L''(\rho+w,\chi)|,
\]
one has, retaining every zero occurrence,
\begin{equation}\label{eq:A049-weightedcurvature}
                  \sum |B(\rho)|^4M_\rho^2\le C_MNL^4.
\end{equation}
The constants are absolute because \(P,\nu\) are fixed once;
no height threshold uniform over growing conductors is asserted.
\end{lemma}

\begin{proof}
We use the already proved uniform shifted main terms with the
actual coefficient polynomial \(Q=B^2\), not the coefficient
sequence supported at one. This distinction is responsible
for the stronger bound.
For zeta the input is the exact Heap--Li--Zhao Theorem~2,
equations (5)--(9), extended to these coefficients and
fully used in NN-A047. For the other fixed characters it is
the locally proved Lemma~\ref{lem:A048-discrete-second}.
Under GRH the finite shifted sum is
\begin{align*}
 S_\chi(\alpha,\beta)
   ={}&\sum_{T<\gamma\le2T}
       L(\rho+\alpha,\chi)L(1-\rho+\beta,\bar\chi)\\
      &\hspace{15mm}{}\times
                 B_{\chi,T}(\rho)^2B_{\bar\chi,T}(1-\rho)^2.
\end{align*}
It is holomorphic in both shifts on the small disks in use.
All these sums retain multiplicities.

Here are the bounds on the actual main terms, before any
differentiation in \(\alpha,\beta\). Put \(c_0=1/1000\),
and choose
\[
 |\alpha|=2c_0/L,\qquad |\beta|=3c_0/L,\qquad
 |\gamma|=c_0/L.
\]
The gamma circle here denotes the auxiliary shift in
\(\mathcal I_q\), not a zero ordinate. It is smaller than
both other circles. Every possible pole at
\(\gamma=-\beta\) or \(\gamma=\alpha\) therefore lies
outside its closed disk. Also \(\alpha+\beta\ne0\)
on the two outer circles. Each separate permuted main
term is consequently holomorphic inside the auxiliary
gamma circle. The reciprocals of \(Z_q(1-\gamma)\)
that occur are holomorphic, with a zero at \(\gamma=0\);
they cause no pole. We do not take a Cauchy derivative
around uncancelled poles.

For precision, the four-variable Mellin proof in NN-A047
uses the fixed vertex weight
\[
 W_{P,\nu}(s)=e^{\nu s}
                 \sum_j\frac{p_jj!}{\nu^js^{j+1}},
                       \qquad \Re s=1/\nu.
\]
The smallest nonzero degree \(J\) of \(P\) is at least sixteen.
Its two actual coefficient Euler products are
\eqref{eq:A047-F-euler} and \eqref{eq:A047-G-euler}, with
\(\zeta\) replaced by
\[
                 Z_q(s)=\zeta(s)\prod_{p\mid q}(1-p^{-s})
\]
and the same coprimality restrictions for \(q>1\).
These identities are initially unfolded in the
absolute-convergence chamber and continued to the
displayed distinct shift circles, exactly as in the
checked NN-A047--A048 proof.

The global bounds before localization are \(O_q(L^3)\)
and \(O_q(L^6)\), respectively. Removing the region where
any scaled Mellin imaginary part exceeds \(\sqrt L\)
costs
\[
                 O_q(L^{3-J/2})=o(L^{-1}),\qquad
                 O_q(L^{6-J/2})=o(1).
\]
These are the full existing global tail estimates; they
are applied before the Laurent approximations.
On the remaining box use
\[
                Z_q(1+u)=\frac{E_q}{u}(1+O_q(u)),
                 \qquad E_q=\prod_{p\mid q}(1-p^{-1}),
\]
with \(E_1=1\). The convergent Euler remainder products
tend to one. The first coefficient product has net
zeta exponent \(-1\), so its integrated Mellin value is
\(O_q(L^{-1})\); the second, including its explicit
three-shift prefactor, has net exponent zero, so its
integrated value is \(O_q(1)\).
The same polynomial-growth model kernels and the
order-sixteen vertex weights give an integrable majorant.
Thus these bounds are uniform on all displayed shift
circles. They follow directly from the dominated
convergence calculation before the three differentiations
in the proof of Lemma~\ref{lem:A047-second}.

In \(\mathcal J_q\), the explicit pole factor is \(O_q(L)\)
and \(\log(qt/(2\pi))=O_q(L)\); the coefficient sum just
bounded is \(O_q(L^{-1})\). Its integrand is therefore
\(O_q(L)\), and its dyadic integral is \(O_q(TL)\).
In \(\mathcal I_q\), the undifferentiated coefficient
sum is \(O_q(1)\). The accompanying shifted powers are
bounded, and Cauchy on the smaller gamma circle bounds
its physical \(\gamma\)-derivative by \(O_q(L)\).
The resulting integral is again \(O_q(TL)\).

We also justify a constant independent of the fixed
character. The preceding localized limits are exactly
the universal scaled kernels already derived in NN-A048:
the \(E_q^{-1}\) from the first coefficient product
cancels the \(E_q\) from its explicit pole, while the
second full product has exponent zero. The Euler
remainders have value one at the origin, and
\[
        \frac{\log(qt/(2\pi))}{L}=1+o_q(1),\qquad
                      N\sim\frac{TL}{2\pi}.
\]
Consequently the normalized main terms converge uniformly
on the chosen circles to bounded functions depending only
on the already fixed \(P,\nu,c_0\), not on \(q\) or \(\chi\).
The domination and tail statements above justify this
uniform limit for each fixed \(q\); the required starting
height may depend on \(q\).

The zeta source error is \(O_A(TL^{-A})=o(N)\).
For fixed \(q>1\) the proved error is
\(O_{q,\eta}(T^{1/2+2\nu+\eta})=o(N)\), choosing
\(0<\eta<1/2-2\nu\). Hence, for an absolute \(C_Q\),
\(|S_\chi(\alpha,\beta)|\le C_QN\) on the two outer
circles once \(T\) is large enough for the fixed character.
Successive maximum-modulus bounds for the finite shifted
sum extend this inequality to \(|\alpha|,|\beta|\le R\),
including coinciding shifts. On setting
\(\alpha=z,\beta=\bar z\), GRH gives
\[
 L(1-\rho+\bar z,\bar\chi)=\overline{L(\rho+z,\chi)},
 \qquad
 B_{\bar\chi,T}(1-\rho)=\overline{B_{\chi,T}(\rho)}.
\]
This proves \eqref{eq:A049-weightedvalue}.

Finally the same Cauchy estimate as in
Lemma~\ref{lem:A049-curvature-mean} gives
\[
 M_\rho^2\le256R^{-4}\frac1{2\pi}
       \int_0^{2\pi}|L(\rho+Re^{i\theta},\chi)|^2\,d\theta .
\]
Multiply by the nonnegative number \(|B(\rho)|^4\),
sum over all occurrences, and apply
\eqref{eq:A049-weightedvalue}. The resulting bound is
\eqref{eq:A049-weightedcurvature}, with an absolute
\(C_M\). No phase independence or inverse-derivative
moment is used.
\end{proof}

\begin{corollary}[NN-A049: a fixed quarter-logarithmic common interval]
\label{cor:A049-quarter-logarithmic}
There is an absolute constant \(c_*>0\) such that, for
every fixed primitive character in the scope of NN-A049
and every sufficiently large \(T\) depending on it,
two sets of distinct simple centers in \((T,2T)\),
each of size at least \(10^{-16}N_{\chi,T}\), have
the same distinct simple shifted branches on the entire
complex parameter disk
\[
                      |\tau|\le
                         \frac{c_*}{T(\log T)^{1/4}}.
\]
For all positive real \(\tau\) in that disk, the two sets
stay on their respective strict sides of the critical
line, with ordinates in \((T,2T)\).
The hypotheses are RH alone for zeta and A for the
other fixed characters. The height threshold is not
uniform as the conductor grows.
\end{corollary}

\begin{proof}
We make the fixed deletion budgets explicit. Retain
\(c_0=1/1000\), and write
\begin{align*}
 A_\rho&=D_\rho B(\rho)^2,&
 W_\rho&=|K_\rho|\,|B(\rho)|^2,\\
 \kappa&=\frac1{4000000},&
 \lambda&=10^{-9}.
\end{align*}
The previously proved signed moment and second moment
give each strict sign mass at least \(2\kappa NL\),
\[
                       W_\rho\le6|A_\rho|,\qquad
                       \sum|A_\rho|^2\le NL^2.
\]
The phase-margin deletion costs at most
\(6\lambda NL<\kappa NL/8\), and the two endpoint unit
intervals cost \(o(NL)\), eventually at most \(\kappa NL/8\).

There is also an absolute \(C_B\) such that
\begin{equation}\label{eq:A049-absolute-sampleconstant}
                         \sum|B(\rho)|^2\le C_BNL^{1/4}
\end{equation}
for all sufficiently large \(T\) for each fixed character.
Indeed the proof of Lemma~\ref{lem:A049-directsampling}
bounds its diagonal and prime-power main sum using only
the absolute coefficient norm, \(P,\nu\), and
\(N\sim TL/(2\pi)\). Its character-dependent rational-twist
and endpoint errors are \(o(T)\), and so are eventually
absorbed by one further absolute multiple of \(NL^{1/4}\).
This gives an absolute eventual bound, not a
conductor-uniform starting height.

Choose the fixed absolute constants
\[
 d=\min\{1,\kappa/(48C_B)\},\qquad
 K=\max\{1,48\sqrt{C_M}/\kappa\}.
\]
The centers with \(|D_\rho|<dL^{3/4}\) have total absolute
signed weight at most \(6dC_BNL\le\kappa NL/8\).
For the remaining centers, discard those with
\(M_\rho/|D_\rho|>KL\). By
\eqref{eq:A049-weightedcurvature},
\[
 \sum_{\mathrm{discarded}}|A_\rho|^2
    \le\frac1{K^2L^2}\sum |B(\rho)|^4M_\rho^2
    \le\frac{C_M}{K^2}NL^2.
\]
Cauchy's inequality bounds their absolute signed weight
by \(6\sqrt{C_M}NL/K\le\kappa NL/8\).
All multiple centers have \(D_\rho=0\) and zero signed
weight; they were retained in the moment estimates and
removed by the first derivative threshold.

The four deletions cost at most \(\kappa NL/2\).
Each sign therefore retains mass greater than \(\kappa NL\).
Exactly the count calculation of
Lemma~\ref{lem:A049-robust}, namely
\[
 \kappa^2N^2L^2\le
        \#S\sum (K_\rho|B(\rho)|^2)^2
                       \le36\,\#S\,NL^2,
\]
gives more than \(10^{-16}N\) centers of each sign.
Every selected center satisfies
\[
 T+1<\gamma<2T-1,\quad
 |D_\rho|\ge dL^{3/4},\quad
 \frac{M_\rho}{|D_\rho|}\le KL,\quad
                 |\Re v_\rho|\ge\lambda|v_\rho|.
\]

We now choose constants for the actual common branches.
On a fixed disk of sufficiently small absolute radius,
the local deformation proof in NN-A049 supplies absolute
positive constants \(c_G,C_G,C_R\) such that
\[
 \begin{gathered}
 c_GT\le|G(\rho)|,\qquad |G(s)|+|G'(s)|\le C_GT,\\
 F_\chi(s,\tau)=L(s,\chi)-\tau G(s)+R_\chi(s,\tau),\\
 |R_\chi|+|\partial_sR_\chi|\le C_RT^2|\tau|^2,
                              \qquad T|\tau|\le1/10 .
 \end{gathered}
\]
These constants come from the absolutely convergent
shifted Dirichlet series and a fixed Cauchy disk; they
are independent of the character. The local curvature
here is the actual \(M_\rho\), not a new growth assumption.
Put
\begin{align*}
 h&=\min\{c_0/16,1/(100K)\},\qquad r_T=h/L,\\
 C_{\mathrm{dir}}&=\max\{9C_G/2,3C_G,C_R/c_G\},
\end{align*}
and choose
\[
 0<c_*\le
 \min\{1,\ dh/(100C_G),
                     \lambda d/(100C_{\mathrm{dir}}K)\}.
\]
All these constants are absolute.
For \(|\tau|\le\tau_T:=c_*/(TL^{1/4})\), the disk
\(|s-\rho|\le4r_T\) lies inside the defining curvature
disk when \(T\) is large, since \(4h<c_0/2\).
Two selected centers within \(4r_T\) would imply by
Taylor's formula
\[
              |D_\rho|\le2M_\rho r_T
                                  \le2Kh|D_\rho|<|D_\rho|,
\]
a contradiction. Their \(r_T\)-disks are disjoint.

On a boundary \(|s-\rho|=r_T\), the Taylor error for
\(L(s,\chi)\) is at most \((Kh/2)|D_\rho|r_T\).
The forcing term has relative bound
\[
       \frac{|\tau G(s)|}{|D_\rho|r_T}
                     \le\frac{C_Gc_*}{dh}\le\frac1{100},
\]
and the quadratic parameter remainder has relative
bound \(O(L^{-1/4})\).
Rouche's theorem therefore gives exactly one shifted
zero in every disk for every complex \(|\tau|\le\tau_T\).
Also
\[
 \frac{|\partial_sF_\chi(s,\tau)-D_\rho|}{|D_\rho|}
       \le Kh+O(L^{-1})+O(L^{-5/4})<1/2
\]
eventually. Each zero is simple, and the unique local
sections patch holomorphically over the entire disk.

For completeness the direction bound has a fixed,
controlled constant. From its equation, absorbing
the Taylor term and the \(G'\) term, the zero satisfies
\(|s_\rho(\tau)-\rho|\le3|\tau v_\rho|\) for large \(T\).
Substitution back into that equation gives
\[
 \frac{|s_\rho(\tau)-\rho-\tau v_\rho|}{|\tau v_\rho|}
 \le C_{\mathrm{dir}}\left(
       \frac{M_\rho|\tau|T}{|D_\rho|^2}
                  +\frac{|\tau|T}{|D_\rho|}+|\tau|T\right).
\]
The three terms are bounded respectively by
\[
               \frac{Kc_*}{d},\qquad
               \frac{c_*}{dL},\qquad
               c_*L^{-1/4}.
\]
The first, after multiplication by \(C_{\mathrm{dir}}\),
is at most \(\lambda/100\); the other two tend to zero.
The total is therefore less than \(\lambda/2\) eventually,
uniformly in all selected centers and
\(0<|\tau|\le\tau_T\).
The phase margin preserves the required strict normal
side for every positive \(\tau\) in the disk.
The endpoint deletion and \(r_T<1\) keep all ordinates
in \((T,2T)\). This proves the assertion.
\end{proof}

\paragraph{Geometric scope and remaining gate.}
The selected reduced analytic zero divisors now give the
same positive fractions of disjoint unramified sections
over one disk of radius \(c_*/[T(\log T)^{1/4}]\).
The Hermitian weight controls curvature relative to the
derivative, so the full initial divisor, including its
nonreduced points, is still counted before the selection.
The orientation comes from the specified real parameter
ray, not an ordering of complex schemes.
This is the denominator-shift family throughout; it is
not Newman heat time. The critical scale \(c/T\), an
arithmetic half-label law, conductor-uniform height
thresholds and analogous general Dedekind-product
ordinary-sign statements remain separate.


\subsection*{Separated sampling and a subpower logarithmic loss}

This addendum uses the already proved weighted-curvature estimate
and derives the additional logarithmic approximation locally.
No central limit theorem for derivatives, simplicity hypothesis,
unrestricted pair correlation, or negative derivative moment is used.
Write \(L=\log T\), \(N=N_{\chi,T}\), and retain the one fixed
square mollifier and all hypotheses of NN-A049.

\begin{lemma}[A smoothed logarithm with an integrable zero kernel]
\label{lem:A049-smoothed-logarithm}
Fix absolute constants \(0<\nu_0<1/100\) and \(\eta>0\), and put
\(x=T^{\nu_0}\), \(\sigma_0=1/2+\eta/L\). Define
\[
 w_x(n)=
 \begin{cases}
 1,&n\le x,\\
 \log(x^2/n)/\log x,&x<n\le x^2,\\
 0,&n>x^2,
 \end{cases}
 \qquad
 P_{\chi,T}(t)=\sum_{2\le n\le x^2}
       \frac{\Lambda(n)w_x(n)\chi(n)}{n^{\sigma_0+it}\log n}.
\]
Under GRH for the fixed character, uniformly for \(T\le t\le2T\),
\begin{equation}\label{eq:A049-log-kernel}
 \left|\log|L(\sigma_0+it,\chi)|-\Re P_{\chi,T}(t)\right|
 \le C_\eta E_T(t)+1,\qquad
 E_T(t)=\sum_{\gamma'}\frac1{1+L^2(t-\gamma')^2}.
\end{equation}
The sum is over every nontrivial zero occurrence of the same
character, at positive, negative and zero ordinates.
One may take
\[
 C_\eta=\max\{4/(\nu_0\eta),3/\nu_0^2\}.
\]
The assertion holds for all sufficiently large \(T\) for the
fixed character, also for zeta with its pole retained.
\end{lemma}

\begin{proof}
We give the contour identity, including its nonzero-height
pole terms. For \(s=\sigma+it\), \(\sigma>1/2\), start with
\[
 \frac1{2\pi i\log x}\int_{(2)}
       -\frac{L'}L(s+w,\chi)
             \frac{x^{2w}-x^w}{w^2}\,dw.
\]
The Dirichlet series and this integral converge absolutely.
The elementary Perron integral
\((2\pi i)^{-1}\int_{(c)}y^w w^{-2}\,dw=\max(\log y,0)\)
gives \(\sum\Lambda(n)w_x(n)\chi(n)n^{-s}\).
At \(w=0\) the kernel has residue one, since
\(x^{2w}-x^w=w\log x+O(w^2)\).
A zero \(\omega\), repeated according to multiplicity,
has residue minus
\((x^{2(\omega-s)}-x^{\omega-s})/((\omega-s)^2\log x)\).
For zeta the pole at one contributes this expression with
the opposite sign.

Here are sufficient details to justify shifting all the way
to the left. First shift to a fixed vertical line on which
\(\Re(s+w)\) is negative and midway between consecutive
trivial zeros. Choose the top and bottom ordinates away
from nontrivial zeros using their \(O_\chi(\log V)\) local
count. The Hadamard subtraction proved in the direct
Landau lemma bounds \(L'/L\) there by
\(O_{\chi,R}(\log^2 V)\) on each fixed strip of width \(R\).
The horizontal integrals tend to zero, because their
denominators have size \(V^2\). On the far left vertical
lines the functional equation and gamma reflection give
\[
 \left|\frac{L'}L(s+w,\chi)\right|
       \ll_\chi \log(R+|t+\Im w|+3),
\]
with an additional harmless absolute constant, as the
reflected Dirichlet series is absolutely convergent.
The choice midway between trivial zeros bounds the
cotangent term also when \(t+\Im w=0\).
Consequently the remaining vertical integral is bounded by
\[
 O_\chi\left(x^{-R}
                \frac{\log(R+|t|+3)+1}{R}\right),
\]
after increasing \(R\) along these lines, and tends to zero.
The nontrivial-zero residue sum converges absolutely by
the zero count and its squared denominator. The trivial-zero
sum converges geometrically. These arguments prove
\begin{align}
 -\frac{L'}L(s,\chi)
 ={}&\sum_n\frac{\Lambda(n)w_x(n)\chi(n)}{n^s}
 +\frac1{\log x}\sum_{\omega\ {\rm zero}}
          \frac{x^{2(\omega-s)}-x^{\omega-s}}{(\omega-s)^2}
 \nonumber\\
 &-\frac{\delta_\chi}{\log x}
          \frac{x^{2(1-s)}-x^{1-s}}{(1-s)^2},
 \label{eq:A049-smoothed-logderivative}
\end{align}
where the sum includes trivial zeros and
\(\delta_\chi=1\) for zeta and zero otherwise.
The underlying classical interfaces are the exact
Hadamard and local-count statements already audited in
Iwaniec--Kowalski, Proposition~5.7(1), equation~(5.27),
pp.~102--103, and Theorem~5.8 and its proof,
equations~(5.34)--(5.36), pp.~104--105.
The primitive functional equation and trivial-zero locations
are DLMF~25.15.5--6 and its zeros subsection, with
\(L(1,\chi)\ne0\) in 25.15.9 for the endpoint cancellation; the gamma
asymptotic/reflection interfaces are DLMF~5.11 and 5.5.
The displayed smoothed identity itself has been derived here.

Integrate \eqref{eq:A049-smoothed-logderivative} in \(\sigma\)
from \(\sigma_0\) to infinity. The logarithm is the branch
obtained from the Euler product along that horizontal ray;
GRH makes it zero-free. Absolute convergence of the integrated
zero kernels justifies this operation and the interchange.
The prime sum integrates to \(P_{\chi,T}(t)\).
For a nontrivial zero \(1/2+i\gamma'\), put
\(v=L(\sigma-1/2)\), \(y=L(t-\gamma')\). Its integrated
absolute contribution is at most
\[
 \frac1{\nu_0}\int_\eta^\infty
       \frac{e^{-\nu_0v}+e^{-2\nu_0v}}{v^2+y^2}\,dv
                    \le\frac{C_\eta}{1+y^2}.
\]
For \(|y|\le1\), use \(2\int_\eta^\infty v^{-2}dv\);
for \(|y|\ge1\), use the integral of the two exponentials
divided by \(y^2\). These give the stated constants.

Every trivial zero has real part at most zero, with the
usual spacing two and multiplicity one. Their total
integrated contribution is
\(O(x^{-\sigma_0}/(T^2(\log x)^2))\).
The zeta pole contributes at most
\[
 \frac{x^{2(1-\sigma_0)}/2+x^{1-\sigma_0}}
                         {T^2(\log x)^2}
                 \le\frac{2x}{T^2(\log x)^2}.
\]
Their sum is eventually less than one. Taking real parts
now proves \eqref{eq:A049-log-kernel}.
\end{proof}

\begin{lemma}[Two sampling bounds, with every complementary multiplicity]
\label{lem:A049-logarithmic-sampling}
For \(P=P_{\chi,T}\) above there is an absolute constant
\(C_P>0\) such that, eventually for each fixed character,
\begin{equation}\label{eq:A049-prime-square}
                  \sum_{T<\gamma\le2T}|P(\gamma)|^2
                                      \le C_P N\log L.
\end{equation}
Fix \(b>0\). If a set \(S\) of distinct zero centers in
\((T,2T)\) has ordinate separation at least \(b/L\),
then an absolute constant \(C_E(b)>0\) gives
\begin{equation}\label{eq:A049-reversed-kernel}
                         \sum_{\gamma\in S}E_T(\gamma)
                                              \le C_E(b)N.
\end{equation}
Both bounds use the full multiplicities of the ambient zeros;
the second uses separation only of the chosen centers.
\end{lemma}

\begin{proof}
Write \(Y=x^2\), put \(a_1=0\), and for \(n\ge2\) set
\[
 a_n=\frac{\Lambda(n)w_x(n)}{\log n}\,n^{-\eta/L}.
\]
These are nonnegative, bounded by one, and supported on prime
powers. Thus \(P(\gamma)=\sum_{n\le Y}a_n\chi(n)n^{-\rho}\).
Use the same good endpoints \(U,V\) and uniform Landau formula
from Lemma~\ref{lem:A049-directsampling}. The diagonal is
\(N(U,V)\sum a_n^2|\chi(n)|^2/n\). For \(m>n\), the
off-diagonal main term is zero unless \(m=rn\) with \(r\)
an integer prime power; in that case its contribution,
including its conjugate, is
\[
           -\frac{V-U}{\pi}
                   \frac{a_na_m\Lambda(r)}m\le0
\]
when \((m,q)=1\), and is zero otherwise. Indeed
\(\chi(n)\overline{\chi(m)}\chi(m/n)=1\) in the former case.
This proves the needed sign also for nonreal characters.

Every remaining contour error can be summed absolutely as
in the direct sampling proof, now using merely \(a_n\le1\).
For clarity, after multiplication by \(1/m\) the three
errors sum to
\[
 \begin{aligned}
 L^2\sum_{n<m\le Y}\frac1n&\ll YL^3,\\
 \sum_{n<m\le Y}\frac{\log(2m/n)}
                       {n\operatorname{dist}_*(m/n,\mathbb Z)}
                                  &\ll Y^2L,\\
 L\sum_{n<m\le Y}\frac1{m\log(m/n)}&\ll YL^2.
 \end{aligned}
\]
Here \(\operatorname{dist}_*\ge1/n\), including the convention
at integer ratios, and \(\log(m/n)\ge(m-n)/m\).
Since \(Y=T^{2\nu_0}\), all these errors are \(o(T)\).
The endpoint strips contain \(O_\chi(L)\) occurrences;
\(|P|\le2\sqrt Y\) bounds their cost by \(O_\chi(YL)=o(T)\).

Finally,
\[
 \sum_{n\le Y}\frac{a_n^2}n
 \le\sum_{p\le Y}\frac1p+
             \sum_p\sum_{k\ge2}\frac1{k^2p^k}
                                     \ll\log L.
\]
No prime number theorem is required for this upper bound:
for \(\epsilon=1/\log Y\),
\(\sum_{p\le Y}p^{-1}\le e\sum_p p^{-1-\epsilon}
 \le e\log\zeta(1+\epsilon)
 \le e\log(1+\log Y)\).
The higher-power series converges absolutely. The diagonal
constants are absolute; the character-dependent \(o(T)\)
errors can eventually be absorbed by one further absolute
multiple of \(N\log L\). This proves
\eqref{eq:A049-prime-square}.

For the second assertion, separation implies, for every real
\(u\), the elementary lattice bound
\[
              \sum_{\gamma\in S}
                \frac1{1+L^2(\gamma-u)^2}\le C(b).
\]
For example, group the centers into successive intervals of
length \(b/L\) to compare with the convergent sum
\(2+2\sum_{j\ge0}(1+b^2j^2)^{-1}\).
Apply this first to ambient ordinates in \([T/2,3T]\).
Their total multiplicity is \(O(N)\) by the fixed-character
zero count, with an eventual absolute constant.

The infinite complement must be estimated separately.
The zeros outside \([T/2,3T]\) with \(|\gamma'|\le4T\)
have distance at least \(T/2\) from every center and total
multiplicity \(O(TL)\). Their contribution to the double
sum is at most \(C\#S/(TL)\).
For \(2^jT<|\gamma'|\le2^{j+1}T\), \(j\ge2\), the
count is \(O(2^jT(L+j))\), and every distance is at least
\(2^{j-1}T\). These shells contribute at most
\[
 \frac{C\#S}{TL^2}\sum_{j\ge2}2^{-j}(L+j)
                         \ll\frac{\#S}{TL}.
\]
The fixed-character count for both signs of the ordinate
has these absolute eventual constants; for negative
ordinates use the conjugate character. Since \(\#S\le N\),
the complement is \(O(1)\). Tonelli's theorem is valid
because the kernels are nonnegative. This proves
\eqref{eq:A049-reversed-kernel}.
\end{proof}

\begin{corollary}[NN-A049: a common interval with subpower logarithmic loss]
\label{cor:A049-sublogarithmic}
There are absolute constants \(c,C>0\) such that the same
positive branch-proportion conclusion of NN-A049 holds on
one common complex parameter disk
\begin{equation}\label{eq:A049-sublog-time}
               |\tau|\le\frac cT
                         \exp\bigl(-C\sqrt{\log\log T}\bigr).
\end{equation}
Each strict side retains at least \(10^{-16}N_{\chi,T}\)
distinct simple centers, the same centers throughout the
disk. Their ordinates stay in \((T,2T)\), and their sides
remain strict for every positive real \(\tau\) in it.
Zeta requires RH alone; every other fixed primitive character
uses A. The starting height can depend on the character.
For every fixed \(\epsilon>0\), this interval eventually
contains \(0<\tau\le1/[T(\log T)^\epsilon]\).
It does not assert an interval \(c/T\).
\end{corollary}

\begin{proof}
Use exactly the constants \(\kappa=1/4000000\),
\(\lambda=10^{-9}\), \(C_M\), and
\(K=\max\{1,48\sqrt{C_M}/\kappa\}\) from the weighted
curvature proof. Each initial strict sign has weight at
least \(2\kappa NL\), where
\[
 W_\rho=|K_\rho|\,|B(\rho)|^2\le6|A_\rho|,
 \qquad A_\rho=D_\rho B(\rho)^2,
 \qquad\sum|A_\rho|^2\le NL^2.
\]
Remove the phase-margin exceptions, endpoint unit strips,
and nonzero-derivative centers with
\(M_\rho/|D_\rho|>KL\). Each costs at most \(\kappa NL/8\)
by that proof. Multiple centers have zero weight and are
removed without cost. Denote the remaining centers by \(S\).
No lower bound on the nonzero derivative has yet been imposed.

The points of \(S\) have separation \(b/L\), for an absolute
\(b>0\). Indeed, set \(h=\min\{c_0/16,1/(100K)\}\),
\(c_0=1/1000\). If two were within \(4h/L\), Taylor's
formula between these two original zeros would give
\[
             |D_\rho|\le\frac12M_\rho(4h/L)
                            \le2Kh|D_\rho|<|D_\rho|.
\]
The segment lies in the defining curvature disk. Thus
\(b=4h\) is valid, without a spacing hypothesis about
any complementary zero.

Choose once
\(0<\eta\le\min\{c_0/4,1/(4K)\}\). Along the positive
horizontal segment from \(\rho\) to \(\rho+\eta/L\),
Taylor's formula gives
\[
 L(\rho+\eta/L,\chi)=\frac\eta L D_\rho(1+e_\rho),
                   \qquad |e_\rho|\le K\eta/2\le1/8.
\]
In particular the shifted value is nonzero, and
\begin{equation}\label{eq:A049-log-derivative-transfer}
 \log\frac{|D_\rho|}L
     =\log|L(\rho+\eta/L,\chi)|-\log\eta
                                         -\log|1+e_\rho|.
\end{equation}

Put \(\delta=(\kappa/96)^2\). The two preceding lemmas
allow absolute choices \(A_0,B_0>0\) such that the union
of the exceptions
\[
 |P_{\chi,T}(\gamma)|>A_0\sqrt{\log L},\qquad
                         E_T(\gamma)>B_0
\]
contains at most \(\delta N\) centers of \(S\).
For instance take \(A_0^2\ge2C_P/\delta\) and
\(B_0\ge2C_E(b)/\delta\), and use Markov's inequality.
Their total signed absolute weight is at most
\[
                  6\sqrt{\delta N}\sqrt{NL^2}
                                      =\kappa NL/16.
\]
On every remaining center, \eqref{eq:A049-log-kernel} and
\eqref{eq:A049-log-derivative-transfer} give, after increasing
an absolute \(C>0\),
\begin{equation}\label{eq:A049-sublog-derivative}
                   |D_\rho|\ge L\omega_T,
            \qquad \omega_T=e^{-C\sqrt{\log L}}.
\end{equation}
Here the fixed constants \(C_\eta B_0\), \(|\log\eta|\)
and the bounded \(\log|1+e_\rho|\) are absorbed into
\(C\sqrt{\log L}\), which is legitimate for \(\log L\ge1\).
Each sign still has weight greater than \(\kappa NL\).
Cauchy's inequality as in the preceding proof gives at
least \(\kappa^2N/36>10^{-16}N\) centers of each sign.

For completeness all local continuation scales are checked.
Keep radius \(r_T=h/L\), and the absolute local constants
\(c_G,C_G,C_R,C_{\rm dir}\) from that proof. Choose
\[
 0<c\le\min\{1,h/(100C_G),
                           \lambda/(100C_{\rm dir}K)\}.
\]
For \(|\tau|\le c\omega_T/T\), the boundary forcing
relative to \(|D_\rho|r_T\) is at most \(C_Gc/h\).
The quadratic parameter remainder has relative size
\(O(\omega_T)\). The Taylor error has relative size
at most \(Kh/2\). Rouche therefore gives exactly one
zero in each disjoint disk for every such complex parameter.
The derivative error divided by \(|D_\rho|\) is at most
\[
                    Kh+O(L^{-1})+O(\omega_T/L)<1/2.
\]
The zeros are simple and their unique sections patch
holomorphically. Their displacement obeys
\(|s_\rho(\tau)-\rho|\le3|\tau v_\rho|\).
The same exact direction estimate is bounded by
\[
 C_{\rm dir}\left(
       \frac{M_\rho|\tau|T}{|D_\rho|^2}
             +\frac{|\tau|T}{|D_\rho|}+|\tau|T\right)
             \le C_{\rm dir}(Kc+c/L+c\omega_T).
\]
Its first term is at most \(\lambda/100\), and the
others tend to zero. The phase margin thus preserves
the strict normal direction on the entire positive ray
segment. The endpoint strips and \(r_T<1\) preserve the
ordinate window. Finally,
\(c\exp(-C\sqrt{\log\log T})(\log T)^\epsilon\to\infty\)
for every fixed \(\epsilon>0\), proving the last assertion.
\end{proof}

\paragraph{Geometric scope and precise remaining loss.}
The same fixed positive fractions of reduced analytic
zero divisors have disjoint unramified sections over the
common disk \eqref{eq:A049-sublog-time}. Their local
relative differentials vanish and their real normal
orientations persist, in the same analytic finite-cover
dictionary as SGA~1, Expos\'e~I, section~3, Proposition~1.
The full complementary divisor, including every
nonreduced multiplicity, enters the positive Poisson-type
kernel before its sum is reversed over the selected centers.
This is an actual denominator-shift theorem, not a heat-time
statement, a derivative central limit theorem or an arithmetic
half-label law. A positive interval comparable to \(1/T\),
conductor-uniform starting heights and general Dedekind
ordinary branch proportions remain unproved.


\subsection*{A Gaussian restriction and the full inverse-height interval}

This addendum uses A, including its auxiliary primitive Dirichlet
family even in the conductor-one case. All characters and conductors
are fixed before the height tends to infinity. The polynomial and
length parameter of the square mollifier are exactly the one fixed
choice in NN-A047--A048. In particular its two actual constants obey
\[
 a\ge a_*,\qquad b\le b_*<1,\qquad
 a_*^2-\frac{b_*}{2(1-1/2500)}>\frac1{2000}.
\tag{A049g.1}
\]
The proof below concerns a short prime polynomial and weighted finite
measures. It does not import a central limit theorem for derivatives.

\begin{theorem}[NN-A049: a common inverse-height disk]
\label{thm:A049-critical-disk}
Under A there is an absolute constant $c>0$ such that, for every
fixed primitive character $\chi$, including conductor one, and all
sufficiently large $T$ depending on $\chi$, each of the two strict
normal-velocity signs has at least $10^{-24}N_{\chi,T}$ distinct
simple centers in $(T,2T)$ whose branches in the specified
denominator-shift family are simultaneously simple and disjoint on
the entire complex parameter disk
\[
                              |\tau|\le c/T.
\tag{A049g.2}
\]
For every positive real $\tau$ in this disk, these same two sets
remain on their respective strict sides of the critical line, with
ordinates in $(T,2T)$. The fraction is deliberately conservative.
This does not assert the earlier larger numerical fraction at this
larger time scale, a fixed positive time interval, or a half-label law.
\end{theorem}

\paragraph{Notation.}
Put $L=\log T$, $V=\log L$, and, for sufficiently large $T$, set
\[
 \delta=e^{-(\log V)^2},\qquad z_0=e^{(\log V)^2},\qquad
 x_0=T^\delta,\qquad
 S=\sum_{z_0<p\le x_0\atop p\nmid q}\frac1p,\qquad
 \sigma^2=S/2.
\tag{A049g.3}
\]
The letter $\sigma$ here is a positive normalization, not a real
coordinate in the complex plane. Define the finite polynomial
\[
 P_0(s)=\sum_{z_0<p\le x_0}\chi(p)p^{-s},\qquad
 X_\rho=\frac{\Re P_0(\rho)}{\sigma},\qquad
 A_\rho=L'(\rho,\chi)B_{\chi,T}(\rho)^2.
\tag{A049g.4}
\]
As before, $K_\rho=\Re\{\rho L(\rho+1,\chi)
\overline{L'(\rho,\chi)}\}/\gamma$ and
$C_\rho=K_\rho|B_{\chi,T}(\rho)|^2$. All sums retain every zero
occurrence; at a multiple center $A_\rho=C_\rho=0$.
Let $m_k$ be the $k$th moment of a standard real Gaussian, so
$m_{2j}=(2j)!/(2^j j!)$ and $m_{2j+1}=0$.

The prime number theorem already used in NN-A047, or the degree-one
instance of the audited IK Theorem 5.15, gives by partial summation
$\sum_{p\le u}p^{-1}=\log\log u+c+o(1)$. To be explicit, the
error $O(\sqrt u\log^2u)$ for $\psi(u)-u$ has an integrable
derivative weight $1/(u^2\log u)$; the prime powers of exponent
at least two contribute a convergent constant. Consequently
\[
 S=V-(\log V)^2-2\log\log V+O_\chi(1)\sim V,
 \qquad \delta V^j\longrightarrow0\quad(j\text{ fixed}).
\tag{A049g.5}
\]

\subsubsection*{The mixed first-moment interface}

\begin{lemma}[Two coefficient polynomials in the first shifted moment]
\label{lem:A049-two-shift-first}
Let $X_\chi(s)=\sum x(n)\chi(n)n^{-s}$ and
$Y_{\bar\chi}(s)=\sum y(n)\bar\chi(n)n^{-s}$ have real,
uniformly divisor-bounded coefficients, supported through
$y\le T^\theta$ for fixed $\theta<1/2$. Let $T\le U\le2T$
and let the physical shifts have modulus at most a fixed
constant times $1/L$. Under A,
$\sum_{0<\gamma\le U}L(\rho+\alpha,\chi)
X_\chi(\rho)Y_{\bar\chi}(1-\rho)$ has, up to
$O_{q,\varepsilon}(T^{1/2+\varepsilon}y)$, the main term
\begin{align*}
 &\frac1{2\pi}\int_0^U\log\frac{qt}{2\pi}\,
       \sum_{(n,q)=1}\frac{(n^{-\alpha}*x)(n)y(n)}n\,dt\\
 &\quad-\frac U{2\pi}
       \sum_{(n,q)=1}\frac{(\Lambda*n^{-\alpha}*x)(n)y(n)}n\\
 &\quad+\frac1{2\pi}\int_0^U\left.\partial_\gamma
       \sum_{(h,k)=1\atop(ghk,q)=1}\frac{y(gh)x(gk)}{ghk}
       \left\{Z^q_{\alpha,\gamma,h,k}
       +(qt/(2\pi))^{-\alpha-\gamma}
                           Z^q_{-\gamma,-\alpha,h,k}\right\}
                                                    \right|_{\gamma=0}dt,
\end{align*}
where $Z_q(s)=\zeta(s)\prod_{p\mid q}(1-p^{-s})$ and
\[
 Z^q_{\alpha,\gamma,h,k}
 =h^{-\alpha}k^{-\gamma}
       \frac{Z_q(1+\alpha+\gamma)}{Z_q(1+\alpha)}
       \prod_{p\mid k}\frac{1-p^\gamma}{1-p^{-1-\alpha}}.
\tag{A049g.6}
\]
The combined expression has its removable values. The error is
uniform on distinct small shift circles, and fixed derivatives are
allowed after enlarging $\varepsilon$. The same assertion holds for
the dyadic difference with unchanged coefficient polynomials.
\end{lemma}

\begin{proof}
For $q=1$ the exact main term is Benli--Elma--Ng, published
\emph{Mathematika} 70 (2024), e12281, Theorem 1.1(i), equations
(9)--(11), with the equivalent two-$Z$ expression (13)--(14),
pp.~2--4. The coefficient-majorant argument of NN-A042 gives its
bilinear version. We give the fixed-character contour calculation
as well, so the conductor replacement is not an extrapolation.

Integrate $L'/L(s,\chi)L(s+\alpha,\chi)X_\chi(s)
Y_{\bar\chi}(1-s)$ around the same rectangle as NN-A048.
Its residues are the displayed zero sum, with multiplicities. The
horizontal edges are $O_{q,\varepsilon}(T^{1/2+\varepsilon}y)$.
The diagonal on the right edge is the negative convolution term
shown above: its condition is $nh=k$, so all character factors
cancel. The off-diagonal estimate is precisely the harmonic
distance estimate in Lemma~\ref{lem:A048-real-bilinear}, with one
fewer logarithm; it is within this error.

On the left reflect and conjugate as in NN-A046. Write
$D_\chi=\Delta_\chi'/\Delta_\chi$. The remaining upper-right
integral, before its outside minus sign, is the conjugate of
\[
 \frac1{2\pi i}\int
  \Delta_\chi(s-\bar\alpha)^{-1}
  \left(D_\chi(s)-\frac{L'}L(s,\chi)\right)
  L(s-\bar\alpha,\chi)
  Y_\chi(s)X_{\bar\chi}(1-s)\,ds.
\tag{A049g.7}
\]
This follows directly from $L(s,\chi)=\Delta_\chi(s)
L(1-s,\bar\chi)$ and
$D_\chi(1-s)=D_{\bar\chi}(s)$. Conjugation changes barred
shift parameters back to the unbarred ones in the final real
coefficient main term.

For completeness consider the auxiliary integral obtained by
replacing the parentheses in (A049g.7) by
$L(s+\gamma,\chi)/L(s,\chi)$. In this auxiliary computation
rename the barred shift as $\alpha$; the outside conjugation
restores the unbarred holomorphic formula. Set $v=s-\alpha$ before
stationary phase. Its inner convolution is
\[
 c=1*(n^{-\alpha-\gamma})*(\mu(n)n^{-\alpha}).
\]
There are three factors and one Mobius factor, so the complete
coefficient-uniform stationary proof and auxiliary-character
averaging of NN-A048 apply literally. They permit two different
polynomials. All nonprincipal parts and Perron remainders have
the stated error. After writing the two coefficient indices as
$gh,gk$, the unrestricted principal series is exactly
\[
                \frac{Z_q(s)Z_q(s+\alpha+\gamma)}{Z_q(s+\alpha)}.
\tag{A049g.8}
\]
At $p^v\parallel k$, the principal weight is zero below exponent
$v-1$, $-1/(p-1)$ there, and one above it. At the pole $s=1$,
the weighted-to-unweighted local ratio is
\[
       \frac{(1-p^\gamma)p^{-v(1+\alpha+\gamma)}}
                                      {1-p^{-1-\alpha}}.
\]
Indeed the remaining local numerator series has coefficients
$f_0=1$ and $f_j=(1-p^\gamma)p^{-j(\alpha+\gamma)}$ for
$j\ge1$; summing its tail proves the formula. Restoring the
outer $p^{v(1+\alpha)}$ contribution and the $h$ contribution
gives (A049g.6). As in NN-A048, the residue
$E_q=\prod_{p\mid q}(1-p^{-1})$ cancels $q/\varphi(q)$.
The other pole is $1-\alpha-\gamma$. Multiplication of the
coefficients by $n^{\alpha+\gamma}$ gives exactly the parameter
change $(\alpha,\gamma)\mapsto(-\gamma,-\alpha)$, so its
residue is the second $Z$ term.

At $\gamma=0$ this main term reduces to
$\sum (n^{-\alpha}*x)(n)y(n)/n$: in its first term only
$k=1$ survives, and the second has the reciprocal pole zero
$Z_q(1)^{-1}$. The identity holds initially at $\alpha\ne0$,
then by removal. Replacing $D_\chi(s)$ by
$-\log(q\Im s/(2\pi))$ costs the usual $O(1/t)$ Stirling
error. Endpoint partial summation inserts the first displayed
logarithm. Differentiating the auxiliary integral in $\gamma$
gives the third term, with the positive sign from (A049g.7).
This proves the formula and its signs. The fixed bottom interval
and boundary zero conventions are within the same error, as in
NN-A048. Uniform Cauchy differentiation on separated circles
proves the asserted derivative version.
\end{proof}

\subsubsection*{Five Gaussian moment limits}

\begin{lemma}[Gaussian restriction of the actual moments]
\label{lem:A049-gaussian-moments}
For every fixed integer $k\ge0$, under A,
\begin{align}
 \frac1N\sum X_\rho^k&\longrightarrow m_k,
 &\frac1{NL}\sum A_\rho X_\rho^k&\longrightarrow a m_k,
                                                        \tag{A049g.9}\\
 \frac1{NL^2}\sum |A_\rho|^2X_\rho^k&\longrightarrow b m_k,
 &\frac1N\sum e^{2i\psi_\chi(\gamma)}X_\rho^k&\longrightarrow0,
                                                        \tag{A049g.10}\\
 \frac1{NL}\sum C_\rho X_\rho^k&\longrightarrow0.&&
                                                        \tag{A049g.11}
\end{align}
Here $N=N_{\chi,T}$ and $a,b$ are precisely the already proved
square-mollifier constants, independent of the fixed character.
\end{lemma}

\begin{proof}
We give separately the direct prime sampling, the two actual
mollifier moments, and the two phase statements. No factorization
of an unproved joint random model is used.

\paragraph{Unweighted prime moments.}
Expand $P_0(\rho)^j\overline{P_0(\rho)}^{\,l}$, with $j+l$
fixed. All indices are at most $x_0^{j+l}=T^{o(1)}$.
The diagonal products are equal only if $j=l$. Unique prime
factorization gives their sum as
$j!S^j+O_j((1+S)^{j-1}/z_0)$: all terms with a repeated prime
are bounded by a fixed multiple of
$\sum_{p>z_0}p^{-2}$ times a smaller power of $1+S$.

Use the complete rational Landau estimate of
Lemma~\ref{lem:A049-rational-landau} for the off-diagonals.
An integer prime-power ratio can occur only by adding copies
of one prime to the shorter prime product. Its absolute main
term is at most
\[
 C_{j,l}T(1+S)^{j+l}
                     \sum_{p\le x_0}\frac{\log p}{p}
       \ll_{j,l}T\delta L(1+V)^{j+l}.
\]
For higher powers the prime sum is smaller. This is negligible
relative to $N\sigma^{j+l}$, by (A049g.5). All rational
distance errors are $T^{o(1)}L^{O_{j,l}(1)}=o(N)$:
one can use the crude bound of $Y^2$ index pairs, with
$Y=x_0^{j+l}$, rational distance at least $1/Y$, and
$|\log(m/n)|\gg1/Y$. The coefficient bounds and endpoint
occurrences cost further fixed powers of $Y,L$ only.
The character factors cancel on diagonals and cannot enlarge
these absolute bounds. Expanding the real part and using
$\sigma^2=S/2$ gives the first limit in (A049g.9).

\paragraph{Polynomial insertions into the exact mean-value formulas.}
Fix an integer $K\ge k$ and work through degree $K$ in a formal
variable $z$. Put $E_K(u)=\sum_{j=0}^K u^j/j!$ and use
\[
 Q_z(s)=B(s)^2E_K\left(\frac{zP_0(s)}{2\sigma}\right).
\tag{A049g.12}
\]
The degree-$k$ coefficient of the product of its two insertions
at conjugate points is $X_\rho^k/k!$. Its underlying coefficient
sequences are real for real $z$, uniformly divisor-bounded for
$z$ in any fixed bounded set, and their lengths are at most
$T^{2\nu+K\delta}$, eventually below a fixed $T^\theta$ with
$\theta<1/2$. The coefficients of a positive $j$th power of a prime
polynomial are bounded by $j!\tau_j$ (the zeroth power is one),
and $\sigma^{-1}\le1$
eventually, which proves these uniform assertions. The exact
source formulas therefore apply. If real coefficients are
required, apply them at $2K+1$ distinct bounded real values of
$z$ and interpolate; their errors remain uniform for each fixed
coefficient through degree $K$.

For the first moment use Lemma~\ref{lem:A049-two-shift-first}
with $X=Q_z$ and $Y=E_K(zP_0/(2\sigma))$; for the second use
the actual shifted discrete second moment of NN-A047--A048 with
$Q_z$ at both endpoints. Differentiate in the physical shifts
only after establishing the uniform limits below. For $q=1$
the latter input is Heap--Li--Zhao Theorem 2, equations (5)--(9);
its audited formula counts multiplicities and assumes no
simplicity. For $q>1$ it is
Lemma~\ref{lem:A048-discrete-second}.

Here is the Euler-product calculation, including localization
and tails. Retain the exact Mellin representation of each
half-power vertex from NN-A047, on scaled lines
$\Re s=1/\nu$ with weight
\[
 W(s)=e^{\nu s}\sum_j\frac{p_jj!}{\nu^j s^{j+1}},
 \qquad |W(s)|\le C(1+|\Im s|)^{-J-1},\quad J\ge16.
\tag{A049g.13}
\]
Euler products are first unfolded where the coefficient variables
and shifts ensure absolute convergence, then continued to the
separated Cauchy circles. The finite coefficient sums identify
that continuation, exactly as in NN-A047--A048.

At a prime in $(z_0,x_0]$, write $u,v$ for its two insertion
variables. Before insertion the local coefficient of $p$ on
each squared-mollifier side tends to $-1$; on the first-moment
side without a mollifier it is zero. When all scaled Mellin
variables have modulus at most $R$, the error in replacing
these coefficients and the small shifted powers by their origin
values is $O(\delta(1+R))$. At the origin put
$x=-1+u$, $y=-1+v$. The following local coefficients of
$1/p$ are exact:
\[
 \begin{aligned}
 \text{first diagonal:}&\quad v(1+x)=uv,\\
 \text{first two-}Z\text{ product:}&\quad xv+v=uv,\\
 \text{second gcd product:}&\quad xy+x+y=-1+uv,\\
 \text{second full three-}Z\text{ product:}&\quad
                                    1+xy+x+y=uv.
 \end{aligned}
\tag{A049g.14}
\]
Thus every ratio to its uninserted product has coefficient
$uv/p$. In the first two-$Z$ row the $k$-only term vanishes
at the origin because of $1-p^\gamma$. In the last row the
initial $1$ is the explicit zeta prefactor. The table also
follows directly by listing the local possibilities $g=p$,
$h=p$, $k=p$ in the displayed exact gcd sums. All higher prime
terms, with any fixed formal degree, are $O_K(p^{-2})$ uniformly
on these contours. Since $p>z_0\to\infty$, the uninserted
local factors are bounded away from zero. No division by an
uncontrolled small-prime Euler factor is involved.

For clarity the uninserted first two-$Z$ product, with its two
coefficient variables $u_1,u_2$, is
\[
 C_1(u_1,u_2;\alpha,\gamma)
 \frac{Z_q(1+\alpha+\gamma)}{Z_q(1+\alpha)}
 \prod_{i=1}^2
       \frac{Z_q(1+u_i)^{1/2}}{Z_q(1+u_i+\gamma)^{1/2}}.
\tag{A049g.15}
\]
Here $C_1$ is an absolutely convergent remainder product near
the origin and equals one there. Its first-order prime terms
are $\frac12\sum_i(p^{-1-u_i}-p^{-1-u_i-\gamma})$;
at the origin every full local factor is exactly one. This
proves (A049g.15), including its normalization. The two
second-moment products are the exact $F,G$ products
(\ref{eq:A047-F-euler})--(\ref{eq:A047-G-euler}).

Take $R=\delta^{-1/2}$. Substituting
$u=v=z/(2\sigma)$ in (A049g.14), and multiplying the marked
prime factors, now gives coefficientwise through degree $K$
\[
 \frac{\text{inserted product}}{\text{uninserted product}}
    =\exp\left(\frac{z^2S}{4\sigma^2}\right)+o(1)
    =e^{z^2/2}+o(1).
\tag{A049g.16}
\]
This is uniform on the box $|\Im s_i|\le R$ and all the shift
circles. The error is bounded by a fixed power of $1+V$ times
$\delta R V+z_0^{-1}$, which tends to zero. This follows by
expanding logarithms of $1+O(1/p)$ through the fixed formal
degree; the linear shifted errors sum to $O(\delta RV)$ and
the quadratic errors to $O(\sum_{p>z_0}p^{-2})$.

We check that the smaller box has lost no main term. On all
Mellin lines, coefficients through degree $K$ of the insertion
ratio are $O_K((1+V)^K)$, independently of their imaginary
parts: the nonconstant coefficient of each local ratio is
$O_K(1/p)$ and its constant coefficient is one. The original
global bounds are $O_q(L^3)$ and $O_q(L^6)$ for $F,G$, and
$O_q(L^2)$ for (A049g.15). Removing $|\Im s_i|>\sqrt L$
therefore costs respectively
\[
 O_{q,K}(V^K L^{3-J/2}),\quad
 O_{q,K}(V^K L^{6-J/2}),\quad
 O_{q,K}(V^K L^{2-J/2}).
\]
These are negligible at their required $L^{-1},1,1$ scales.
On $|\Im s_i|\le\sqrt L$, the Laurent approximation
$Z_q(1+w)=E_qw^{-1}(1+O_q(w))$ is uniform. After the
normalizing powers of $L$ are extracted, each model kernel
is bounded by $C\prod_i(1+|\Im s_i|)^2$: the sums of
coefficient variables have a positive real part, and the
physical shift denominators are separated. Combining with
(A049g.13), the intermediate region
$R<\max|\Im s_i|\le\sqrt L$ contributes
$O_K(V^K R^{-14})=o(1)$. The convergent Euler remainders are
uniformly bounded and tend to one on the central box.
These estimates justify localization before any shift
derivative, including the fractional-power branches.

The first diagonal in the source formula has the same factor
(A049g.16). Its uninserted value is one. The source convolution
with $\Lambda$ is zero when $z=0$, and is negligible here:
every index supported by $Y$ is a product of at most $K$ marked
primes. Hence, through degree $K$, its absolute coefficient is
$O_K(\log x_0(1+V)^K)=O_K(\delta L(1+V)^K)$. After division
by the main logarithm this tends to zero, uniformly on the
physical shift circle. Cauchy differentiation preserves this
negligibility.

All source errors are still negligible. For $q>1$ the largest
length is $T^{2\nu+K\delta}$, so choose a fixed
$\varepsilon>0$ with $1/2+2\nu+\varepsilon<1$; fixed
derivatives add logarithms only. For $q=1$ the audited
arbitrary logarithmic saving in the two source formulas
absorbs every fixed polynomial factor. The $E_q$ exponents
cancel exactly as in NN-A048: the first two-$Z$ product has
net exponent zero, $F$ has exponent $-1$ canceled by its
explicit pole, and $G$ has exponent zero. Thus the normalized
limits are the same universal $a,b$.

Choose outer physical circles with distinct sufficiently small
fixed radii after multiplication by $L$, and a smaller auxiliary
$\gamma$ circle that crosses no separate pole. The combined
finite shifted sums are holomorphic. Uniformity of
(A049g.16) and the preceding tail bounds allows Cauchy
differentiation in $\alpha$ for the first moment and in
$\alpha,\beta$ for the second. Extracting the coefficient
of $z^k/k!$ proves the remaining mollifier limits in
(A049g.9)--(A049g.10). This argument uses the actual square
coefficients throughout, rather than substituting the
coefficient sequence supported at one.

\paragraph{The even phase with prime insertions.}
Expand the absolutely convergent exterior ratio in
$e^{2i\psi_\chi(t)}$. Truncate both exterior indices at
$M=T^{1/2}$, and expand the fixed power of $P_0+\overline P_0$.
The discarded exterior tail is $O(M^{-1/2})$; the elementary
bound $|P_0(\rho)|^k\le T^{o(1)}$ makes its total $o(N)$.
The factor $e^{-2i\arctan(1/(2\gamma))}$ differs from one
by $O(1/T)$ and has the same harmless bound.

For each resulting rational frequency $x=b/a$ in lowest terms,
the already proved uniform estimate (A040) gives
\begin{align*}
 \sum_{0<\gamma\le T}e^{2i\theta_\chi(\gamma)}(b/a)^{i\gamma}
 ={}&-\frac{qT}{2\pi\varphi(q)}\sqrt{b/a}\,
       \frac{\chi(b)\bar\chi(a)\mu(b)}{\varphi(b)}\\
 &+O_\chi\left((\sqrt{bT}+a/b+\sqrt{a/b})L^2\right).
\end{align*}
The main term is zero in the noncoprime cases. All our reduced
$a,b$ are at most $MY$ with $Y=x_0^k=T^{o(1)}$, so the
proved range $T^{-1}\le b/a\le T$ applies eventually.

Here is an absolute bound for the full main-term coefficient;
it avoids an unjustified independence assertion. At a prime,
let $i,j$ be its multiplicities in the two inserted prime
products and $r,s$ those in the exterior numerator and
denominator. The frequency is $p^{j+s-i-r}$. Multiplication
of its square root by all critical and exterior weights gives
$p^{-i-2r-s}$. Including $1/\varphi(b)$ bounds it by
\[
 2p^{-\max(i+r,j+s)-r-s}
\tag{A049g.17}
\]
when the exponent in $b$ is positive; otherwise the factor
2 is unnecessary. At an unmarked prime, summing $r,s\ge0$
gives $1+O(p^{-2})$. At a marked prime it is at most
$C p^{-\max(i,j)}$. Thus the unmarked product converges
and summing the at most $k$ prime positions costs at most
$C_k(1+V)^k$. The total main term is
$O_{\chi,k}(T(1+V)^k)=o(N\sigma^k)$.

The weighted errors also sum absolutely. The $\sqrt{bT}$
term costs at most $\sqrt T\,Y^{C_k}L^{C_k}$: its exterior
sum has only one harmonic factor after the square root is
inserted. The $a/b$ term costs at most
$\sqrt M\,Y^{C_k}L^{C_k}$, and $\sqrt{a/b}$ costs at most
$Y^{C_k}L^{C_k}$. These bounds follow respectively from
$\sum_{n\le M}n^{-1}$, $\sum_{n\le M}n^{-1/2}\ll\sqrt M$,
and the convergent remaining exterior sums. For example in
the first term $\sqrt b\le\sqrt{m h}$, so the $m^{-3/2}$
weight becomes $m^{-1}$; all prime-product sums are bounded
by a fixed power of $Y$. Each error is $T^{1/2+o(1)}=o(N)$.
The dyadic difference has the same bounds. This proves the
even-phase limit in (A049g.10), including multiplicities.

\paragraph{The signed moment.}
Expand $X_\rho^k$ into its $k+1$ monomials in $P_0$ and
$\overline P_0$. For each monomial apply
Lemma~\ref{lem:A048-real-bilinear} to the two coefficient
polynomials $BP_0^j$ and $BP_0^{k-j}$, with the exterior
$L(2-\rho,\bar\chi)$ series truncated at $T^{\delta_0}$.
Choose fixed $\delta_0>0$ so that $\nu+\delta_0<1/2$.
Their lengths are $T^{\nu+\delta_0+o(1)}$ and their real
underlying coefficients satisfy the same uniform divisor
bounds. The imaginary part is consequently $o(NL)$.
The exterior restoration costs at most
\[
 O(T^{-\delta_0/2})\sum|A_\rho||X_\rho|^k
 \le O(T^{-\delta_0/2})
       \left(\sum|A_\rho|^2\sum|X_\rho|^{2k}\right)^{1/2}
                         =o(NL).
\]
The extra real term divided by $2\gamma$ in the exact identity
$K_\rho=\Im\{L'(\rho,\chi)\overline{L(\rho+1,\chi)}\}
+(2\gamma)^{-1}\Re\{L'(\rho,\chi)\overline{L(\rho+1,\chi)}\}$
is $O(NL/T)$ by the same Cauchy bound. This proves (A049g.11).
\end{proof}

\subsubsection*{From polynomial moments to a fixed positive interval}

\begin{lemma}[Interval restriction of the five measures]
\label{lem:A049-gaussian-interval}
For every fixed finite interval $I$ the five limits above also
hold with $X_\rho^k$ replaced by $\mathbf1_{X_\rho\in I}$
and $m_k$ replaced by the standard Gaussian mass $p_I$.
\end{lemma}

\begin{proof}
We include the measure argument to avoid any unproved joint
distribution premise. The first and third measures, with
weights $1/N$ and $|A_\rho|^2/(NL^2)$, are positive, have
bounded mass and second moment, and are therefore tight.
Their limiting polynomial moments are those of a Gaussian
times respectively $1,b$. Taylor's formula on the real axis
gives, for every fixed $t$,
\[
 \left|e^{itx}-\sum_{j=0}^{K-1}\frac{(itx)^j}{j!}\right|
                                      \le |tx|^K/K!.
\]
Use even $K$ and the moment limits. The limiting remainder
is at most $|t|^Km_K/K!$, which tends to zero because
$m_{2j}/(2j)!=1/(2^j j!)$. Thus the characteristic functions
converge to $e^{-t^2/2}$ and $be^{-t^2/2}$.

This implies convergence against compactly supported smooth
functions directly by Fourier inversion and dominated
convergence: their Fourier transforms are integrable and
the characteristic functions are bounded by the total
masses. Approximate an interval from above and below by
such smooth functions. The Gaussian density is continuous
and has no atoms at its endpoints, so both positive measures
converge on every fixed interval.

For the complex first measure, Cauchy bounds its total
variation on a set $E$ by the square root of the product of
the first and third positive masses there. Its Taylor
remainder is bounded by
\[
 \frac{|t|^K}{K!}
       \left(\frac1{NL^2}\sum|A_\rho|^2\right)^{1/2}
       \left(\frac1N\sum|X_\rho|^{2K}\right)^{1/2}.
\]
Its limit tends to zero: $m_{2K}^{1/2}/K!$ decays faster
than any fixed exponential by the displayed factorial
formula. Its characteristic function therefore converges
to $ae^{-t^2/2}$. The same Fourier argument and the preceding
setwise Cauchy bound give the interval conclusion.
The even-phase measure is bounded in variation by the first
positive measure. The signed measure is bounded by six
times the Cauchy bound, since $|C_\rho|\le6|A_\rho|$ at
these heights. Their moment limits are zero, and the
identical Taylor, Fourier and interval-approximation
arguments give zero limits on $I$.
\end{proof}

\subsubsection*{Derivative size on a positive Gaussian interval}

\begin{lemma}[Logarithmic approximation on separated centers]
\label{lem:A049-positive-prime-derivative}
Fix $K_0\ge1$. Among simple centers with
\[
 \max_{|z|\le c_0/(2L)}|L''(\rho+z,\chi)|
            \le K_0L|L'(\rho,\chi)|,\qquad c_0=1/1000,
\tag{A049g.18}
\]
all but $o(N)$ of those with $1\le X_\rho\le2$ satisfy
$|L'(\rho,\chi)|\ge L$. The implied exceptional count
allows every multiplicity in the complementary zero sums.
\end{lemma}

\begin{proof}
Choose $h>0$ with $h\le\min(c_0/16,1/(100K_0))$.
Taylor's theorem shows that centers satisfying (A049g.18)
are separated by $4h/L$: a second zero within that distance
would give $1\le2K_0h$, a contradiction.

Fix $0<\eta\le\min(c_0/4,1/(4K_0))$. Use the already
proved smoothed logarithm (\ref{eq:A049-log-kernel}), with
fixed $\nu_0<1/100$, at $\rho+\eta/L$. Denote its prime-power
polynomial by $P_{\chi,T}$. We claim
\[
 \sum_{T<\gamma\le2T}
       |P_{\chi,T}(\gamma)-P_0(\rho)|^2=o(NV).
\tag{A049g.19}
\]
Here are the coefficient details. The prime terms omitted below
$z_0$ or above $x_0$ have squared coefficient norm
$O_\chi((\log V)^2)$ by partial summation and (A049g.5).
At included primes, $|p^{-\eta/L}-1|\le\eta\delta$, so
their norm is $O(\delta^2V)$. The prime-power terms of
exponent at least two have bounded squared norm. Apply the
direct Landau calculation to these three pieces separately.
For a prime-supported piece, distinct indices have no
integer prime-power ratio, so every off-diagonal main term
is zero irrespective of coefficient signs. For the
prime-power piece the nonnegative coefficients make the
prime-power-ratio main terms nonpositive, as proved in the
sublogarithmic addendum. The rational errors have total
$O_\chi(T^{4\nu_0}L^C)=o(N)$ by the same explicit
pair bounds there. Thus each squared sample is at most
$O_\chi(N(1+(\log V)^2))+o(N)$, proving (A049g.19).
In particular $|P_{\chi,T}-P_0|\le\sigma/4$ fails at
only $o(N)$ occurrences.

For the separated centers, the complete reversed-kernel
bound in the sublogarithmic addendum gives
\[
 \sum_{\rho\ \mathrm{selected}}E_T(\gamma)\le C(h)N.
\]
That proof splits all positive, negative and central zero
occurrences into $[T/2,3T]$, a bounded complementary range,
and dyadic infinite tails; it does not discard nearby
multiplicities. Hence $C_\eta E_T(\gamma)\le\sigma/4$
fails at $O(N/\sigma)=o(N)$ selected centers.

On the remaining centers in $[1,2]$, $\Re P_0\ge\sigma$,
so (\ref{eq:A049-log-kernel}) gives
\[
                 \log|L(\rho+\eta/L,\chi)|\ge\sigma/2-1.
\]
Taylor's formula and (A049g.18) give
$L(\rho+\eta/L,\chi)=(\eta/L)L'(\rho,\chi)(1+e)$
with $|e|\le K_0\eta/2\le1/8$. Consequently
\[
 |L'(\rho,\chi)|\ge \frac{L}{2\eta}e^{\sigma/2-1}\ge L
\]
eventually. This proves the lemma without an inverse
derivative moment or a global derivative distribution law.
\end{proof}

\subsubsection*{Proof of the common-disk theorem}

\begin{proof}[Proof of Theorem~\ref{thm:A049-critical-disk}]
Take $I=[1,2]$ and let $p$ be its Gaussian mass. The elementary
bound $p\ge e^{-2}/\sqrt{2\pi}>1/24$ follows from $e^2<8$
and $2\pi<9$. Restrict every sum in this paragraph to $X_\rho\in I$.
By Lemma~\ref{lem:A049-gaussian-interval}, the first and second
square-mollifier moments are $apNL+o(NL)$ and
$bpNL^2+o(NL^2)$. The even-phase limit gives
$\sum\cos^2\psi=(p/2+o(1))N$. Thus the phase-bad count
at $\eta_0=1/50$ is at most
$(p/(2(1-\eta_0^2))+o(1))N$. Cauchy and (A049g.1) imply
\[
 \sum_{|\cos\psi|\ge\eta_0}|A_\rho|
     \ge\left(p\left(a_*-
                 \sqrt{b_*/(2(1-\eta_0^2))}\right)-o(1)\right)NL
     > (p/4000-o(1))NL.
\]
The exact phase identity and the Euler bound
$1/3<|L(3/2+i\gamma,\chi)|<3$ then give
$\sum|C_\rho|\ge(p/600000-o(1))NL$. The interval-restricted
signed moment is $o(NL)$. Each strict sign consequently has
absolute mass greater than $NL/40000000$ eventually.

Put $\kappa=10^{-9}$ and $\lambda=10^{-12}$. Delete the two
unit endpoint strips, which cost $o(NL)$ by the local zero
count and the global second bound $\sum|A_\rho|^2\le NL^2$.
Delete $|\cos\psi|<\lambda$, at cost at most $6\lambda NL$.
Use the existing absolute weighted-curvature bound
$\sum|B(\rho)|^4M_\rho^2\le C_MNL^4$ and choose the fixed
absolute $K_0\ge\max(1,100\sqrt{C_M}/\kappa)$. Deleting
$M_\rho>K_0L|D_\rho|$ costs at most
\[
 \frac6{K_0L}\sum |B(\rho)|^2M_\rho
                         \le\frac{6\sqrt{C_M}}{K_0}NL.
\]
Multiple centers have zero weight and may be discarded.
Apply Lemma~\ref{lem:A049-positive-prime-derivative} to the
remaining separated centers. Its $o(N)$ exceptions have
weight $o(NL)$ by Cauchy and the same global second bound.
For large $T$ each sign therefore still has mass at least
$\kappa NL$, while every retained center satisfies
\[
 |D_\rho|\ge L,\qquad M_\rho/|D_\rho|\le K_0L,
 \qquad |\Re v_\rho|/|v_\rho|\ge\lambda.
\tag{A049g.20}
\]
Cauchy gives at least $(\kappa^2/36)N>10^{-24}N$ such
centers per sign, each simple and counted once.

We verify the full complex disk, not merely pointwise
implicit-function neighborhoods. Choose absolute
$h\le\min(c_0/16,1/(100K_0))$ and put $r=h/L$.
The center disks of radius $r$ are disjoint by the preceding
separation proof. On each boundary,
\[
 |L(\rho+z,\chi)-D_\rho z|
      \le M_\rho r^2/2\le(K_0h/2)|D_\rho|r.
\]
The already proved local shift expansion is
$F_\chi(s,\tau)=L(s,\chi)-\tau G(s)+R(s,\tau)$, where
$c_G T\le|G(\rho)|\le C_G T$, $|G'|\le C_GT$, and
$|R|+|R_s|\le C_RT^2|\tau|^2$ whenever $T|\tau|\le1/10$.
All constants here can be absolute on the fixed small discs,
as in the weighted-curvature addendum. Choose an absolute
$c>0$, initially so small that
\[
 c\le1/10,\qquad C_Gc+C_Rc^2\le h/10.
\]
For every $|\tau|\le c/T$, the perturbation on the boundary
is at most $|D_\rho|r/10$ because $|D_\rho|\ge L$.
Rouch\'e gives exactly one zero inside each disk.
Moreover
\[
 |F_s-D_\rho|\le K_0h|D_\rho|+C_Gc+C_Rc^2
                                      <|D_\rho|/2
\]
eventually. That zero is simple and its local sections patch
uniquely to a holomorphic section on the whole parameter disk.

Let its displacement be $z(\tau)$. The zero equation first
gives $|z(\tau)|\le C|\tau|T/|D_\rho|$. Subtracting the
linear displacement $v_\rho\tau=\tau G(\rho)/D_\rho$ and
dividing by $|v_\rho\tau|$ gives, for $\tau\ne0$,
\[
 \left|\frac{z(\tau)}{v_\rho\tau}-1\right|
 \le C'\left(\frac{M_\rho|\tau|T}{|D_\rho|^2}
                    +\frac{|\tau|T}{|D_\rho|}+T|\tau|\right)
 \le C'((K_0+1)c+c/L).
\tag{A049g.21}
\]
Decrease the same absolute $c$ so this is less than
$\lambda/2$. For positive real $\tau$, (A049g.20) then
forces $\Re z(\tau)$ to have the strict sign of
$\Re v_\rho$. The unit endpoint strips and $r\to0$ keep
the ordinates in $(T,2T)$. This proves every assertion.
\end{proof}

\paragraph{Standalone geometric interpretation and scope.}
For each of the selected reduced points, the actual holomorphic
zero divisor over $|\tau|<c/T$ is a degree-one analytic cover:
the $s$ derivative is nowhere zero and the unique section remains
in its own disjoint disk. The real orientation is the sign of
the normal tangent to the specified critical real line, preserved
by (A049g.21). The first and second analytic measures are the
derivative section multiplied by the specified square-mollifier
section and its Hermitian norm. Their Gaussian restriction is
proved by finite arithmetic moments; it is not a declaration
of probabilistic independence of divisors or Frobenius classes.
The reduced differential interface is the same SGA 1, Expos\'e I,
section 3, Proposition 1 used in NN-A047,
\url{https://grothendiecksga.com/read/sga1/fr/I-3.html}; its analytic
implicit-function counterpart is proved by the displayed
derivative bound. No arithmetic algebraization, heat-flow
identification, fixed positive time, half-label proportion,
or growing-conductor uniformity is asserted.

\paragraph{Source and review boundary.}
The new two-pole calculation uses the fully written NN-A048
stationary/auxiliary-character interfaces. In conductor one
its exact independent source check is Benli--Elma--Ng,
\url{https://londmathsoc.onlinelibrary.wiley.com/doi/pdf/10.1112/mtk.12281},
Theorem 1.1(i), equations (9)--(14), printed pp.~2--4. The
second-moment interface remains the exact Heap--Li--Zhao
Theorem 2 already audited in NN-A047, and the fixed-character
derivation in NN-A048. The rational phase input is the complete
A040 square-root refinement, with A on every auxiliary modulus.
All new moment factorization, localization, interval restriction,
derivative selection and continuation steps are proved here.


\subsection*{A larger disk for simplicity and disjointness}

This consequence concerns the same denominator shift as NN-A049.
The strict-side assertion remains on the inverse-height interval
proved in Theorem~\ref{thm:A049-critical-disk}; it is not extended
to the larger disk below.

\begin{corollary}[NN-A049: a square-root-log-log survival disk]
\label{cor:A049-loglog-survival}
Under A, each fixed primitive character has, for all sufficiently
large $T$ depending on that character, two sets of simple centers
in $(T,2T)$, each containing at least $10^{-24}N_{\chi,T}$ centers,
with opposite strict initial normal-velocity signs. Their same
branches in the denominator-shift family remain simple and in
pairwise disjoint center disks throughout
\[
               |\tau|\le\frac{\sqrt{\log\log T}}{100T}.
\tag{A049s.1}
\]
The selected branches have ordinates in $(T,2T)$ on this complex
parameter disk. Their strict sides for positive real parameters
are guaranteed on $0<\tau\le c/T$ by the preceding theorem.
No claim about sides on the rest of (A049s.1) is made.
\end{corollary}

\begin{proof}
Retain exactly the sets chosen in the proof of
Theorem~\ref{thm:A049-critical-disk}, with $L=\log T$,
$V=\log L$, and fixed absolute $K_0,h,\eta>0$. The derivative
selection proof actually gives the stronger inequality
\[
 |D_\rho|\ge\frac{L}{2\eta}e^{\sigma/2-1},\qquad
 \sigma^2=S/2,\qquad S\sim V,
 \qquad M_\rho\le K_0L|D_\rho|.
\tag{A049s.2}
\]
Thus eventually $\sigma\ge\sqrt V/2$. No new exceptional
centers or counting argument are required.

We use the full convergent shift series, not its quadratic
truncation. Fix a small absolute $r_0<1/100$ and put
\[
 \mathcal R_\chi(s,\tau)=F_\chi(s,\tau)-L(s,\chi)
  =\sum_{k\ge1}\frac{(-1)^k(s)_k\tau^k}{k!}L(s+k,\chi).
\tag{A049s.3}
\]
For $|s-\rho|\le2r_0$, every $L(s+k,\chi)$ on the right
is bounded in modulus by $\zeta(1.4)$, using its absolutely
convergent series. Also $|s|\le3T$ and
$|(s)_k|\le(|s|)_k$. Hence, for $|\tau|\le1/2$,
\[
 |\mathcal R_\chi(s,\tau)|
 \le\zeta(1.4)\bigl((1-|\tau|)^{-|s|}-1\bigr)
 \le C e^{6T|\tau|}.
\tag{A049s.4}
\]
The same summation proves local normal convergence jointly
in $s,\tau$. Cauchy's formula on fixed $s$ circles of radius
$r_0$ gives, on $|s-\rho|\le r_0$,
\[
             |\partial_s\mathcal R_\chi(s,\tau)|
                                        \le C_1e^{6T|\tau|},
\tag{A049s.5}
\]
where $C,C_1$ are absolute. This bound includes complex
parameters throughout the displayed disk.

Set $r=h/L$ as in the preceding theorem. These center disks
are pairwise disjoint. Their unshifted Taylor remainder is
at most $(K_0h/2)|D_\rho|r$. On (A049s.1), (A049s.2)--(A049s.4)
give uniformly
\[
 \frac{|\mathcal R_\chi(s,\tau)|}{|D_\rho|r}
 \le\frac{2e\eta C}{h}
       \exp\left(\frac{3}{50}\sqrt V-\frac{\sigma}{2}\right)
 \le\frac{2e\eta C}{h}e^{-19\sqrt V/100}\longrightarrow0.
\tag{A049s.6}
\]
The final exponent uses $\sigma/2\ge\sqrt V/4$.
Similarly, (A049s.5) divided by $|D_\rho|$ tends to zero.
Rouch\'e therefore gives exactly one zero in each center
disk for every parameter in (A049s.1). Since
\[
 |\partial_sF_\chi-D_\rho|
   \le K_0h|D_\rho|+C_1e^{6T|\tau|}<|D_\rho|/2
\]
eventually, that zero is simple. The unique local sections
patch to one holomorphic section on the full parameter disk.
The discarded unit endpoint strips and $r\to0$ retain the
ordinate condition. Initial velocity signs and the smaller
strict-side interval are inherited without change.
\end{proof}

\paragraph{Geometric meaning and limitation.}
The same reduced analytic points now give disjoint degree-one
analytic covers of the larger parameter disk (A049s.1).
The nonvanishing relative derivative supplies their local
etale differential criterion, in the analytic counterpart of
SGA 1, Expos\'e I, section 3, Proposition 1,
\url{https://grothendiecksga.com/read/sga1/fr/I-3.html}.
The enlarged disk controls simplicity and separation. A real
normal orientation need not be preserved merely by analytic
continuation, so no additional label claim is inferred.
This is not a de Bruijn--Newman heat-flow assertion and does
not identify a global collision law or a fixed positive time.


\subsection*{An ordinary-count Gaussian component for both velocity signs}

The five moment limits in the preceding addendum also give an
arithmetic distribution statement for the actual normal speeds.
This is a lower bound for finite-interval counting measures, not a
full Gaussian law or a claim that the two labels have equal density.

\begin{theorem}[NN-A049: Gaussian lower components of the normal-speed counts]
\label{thm:A049-speed-gaussian-component}
Assume A and fix a primitive character $\chi$. At a simple zero
$\rho=1/2+i\gamma$ define
\[
 v_\rho=\frac{\rho L(\rho+1,\chi)}{L'(\rho,\chi)},\qquad
 u_\rho=\Re v_\rho.
\]
For $T<\gamma\le2T$ put $L=\log T$, $V=\log L$. For each
fixed finite interval $I$ of positive length, let
\[
 M_\pm(T;I)=\#\left\{\begin{array}{l}
 \rho\text{ simple},\quad T<\gamma\le2T,\quad \pm u_\rho>0,\\[1mm]
 \displaystyle
 \frac{\log|u_\rho|-\log(T/\log T)}{\sqrt{(\log\log T)/2}}\in I
 \end{array}\right\}.
\]
For either sign,
\[
 \liminf_{T\to\infty}\frac{M_\pm(T;I)}{N_{\chi,T}}
       \ge10^{-16}\,\mathbb P\{G\in I\},\qquad G\sim N(0,1).
\tag{A049v.1}
\]
The same assertion holds with $|u_\rho|$ replaced by $|v_\rho|$.
For each sign it also holds for the normalized logarithmic
derivative
\[
                 \frac{\log|L'(\rho,\chi)|-\log\log T}
                                      {\sqrt{(\log\log T)/2}}.
\tag{A049v.2}
\]
The lower constant is absolute. The starting height may depend
on the character and the fixed interval. No height-uniform claim
for growing conductors, full derivative central limit theorem,
or zero-density multiple-center assertion is made.
\end{theorem}

\begin{proof}
Retain $A_\rho,C_\rho,X_\rho,\sigma$ and the actual constants
$a,b$ from the preceding Gaussian-moment proof. First fix a
finite interval $J$, and write $p=\mathbb P\{G\in J\}>0$.
The interval-restricted first, second, even-phase and signed
limits are exactly
\[
 \begin{aligned}
 \sum_J A_\rho&=(ap+o(1))NL,&
 \sum_J|A_\rho|^2&=(bp+o(1))NL^2,\\
 \sum_J\cos^2\psi&=(p/2+o(1))N,&
 \sum_JC_\rho&=o(NL).
 \end{aligned}
\]
where the subscript means $X_\rho\in J$. The same certified
$\eta_0=1/50$ calculation as in the common-disk theorem gives
absolute mass at least $(p/600000-o(1))NL$. Each strict
sign therefore has mass at least $pNL/2000000$ eventually.

We keep careful control when $p$ is small but fixed. Delete
$|\cos\psi|<\lambda$, with the absolute $\lambda=10^{-12}$.
Using the restricted count and second moment, its mass is at
most
\[
 6\lambda\sum_J|A_\rho|
 \le6\lambda\left(\#\{X_\rho\in J\}
                           \sum_J|A_\rho|^2\right)^{1/2}
                      \le(6\lambda\sqrt b+o(1))pNL.
\]
The two unit endpoint strips cost $o(NL)$. Choose a fixed
$K_0$ sufficiently large that
$6\sqrt{C_M}/K_0<p/16000000$. The already proved weighted
curvature estimate shows that deleting
$M_\rho>K_0L|D_\rho|$ has at most this mass times $NL$.
All remaining nonzero summands are simple, and their centers
are separated by $4h/L$, with fixed $h=h(K_0)>0$.

The smoothed-log proof now applies for every bounded interval
$J$, whether positive or negative. Choose fixed
$\eta\le\min(c_0/4,1/(4K_0))$. Its full reversed kernel mean
is $O_{K_0}(N)$, while the prime-polynomial difference in
(A049g.19) has squared sample $o(NV)$. Therefore
\[
 \frac{\log|L(\rho+\eta/L,\chi)|-\Re P_0(\rho)}{\sigma}
                                                        \longrightarrow0
\]
in counting measure. Explicitly, for each fixed positive
$\varepsilon$, the count where the absolute expression
exceeds $\varepsilon$ is $o(N)$, by Chebyshev and Markov.
Taylor's formula gives
\[
 L(\rho+\eta/L,\chi)
     =(\eta/L)D_\rho(1+e_\rho),\qquad |e_\rho|\le1/8.
\]
Taking absolute logarithms is valid, and the fixed additive
terms $\log\eta+\log|1+e_\rho|$ are bounded. Since
$\sigma/\sqrt{V/2}\to1$, on $X_\rho\in J$ it follows that
\[
 Y_\rho:=\frac{\log|D_\rho|-\log L}{\sqrt{V/2}}
                          =X_\rho+o(1)
\tag{A049v.3}
\]
in counting measure on the retained set. This uses all
complementary multiplicities in the zero kernel. It does
not assert a distribution for unweighted derivatives at
the deleted centers.

Any $o(N)$ exceptions just described carry $o(NL)$ mass:
$|C_\rho|\le6|A_\rho|$ and the global second moment allow
Cauchy's inequality. Hence, after every deletion and the
transfer (A049v.3), each sign retains mass at least
$pNL/4000000$ eventually. Its squared weights sum to at
most
\[
             36\sum_J|A_\rho|^2\le72pNL^2.
\]
Cauchy now gives at least
\[
 \frac{pN}{72\cdot4000000^2}>10^{-16}pN
\tag{A049v.4}
\]
centers per sign. For a target interval $I=(a,b)$ in
(A049v.2), take $J=[a+\varepsilon,b-\varepsilon]$ and
delete the $o(N)$ failures of $|Y_\rho-X_\rho|<\varepsilon$.
Then take the lower limit in $T$, followed by
$\varepsilon\downarrow0$. The Gaussian has no endpoint
atoms, so (A049v.4) proves (A049v.2) for open, closed or
half-open finite intervals.

Finally the exact velocity formula and the exterior Euler
bounds imply uniformly on the dyadic interval
\[
 \log|v_\rho|=\log T-\log|D_\rho|+O(1).
\]
On the retained phase-margin set,
$|u_\rho|/|v_\rho|=|\cos\psi_\chi(\gamma)|\in[\lambda,1]$,
so the same equality holds with $|u_\rho|$. The normalized
logarithmic speeds are thus $-Y_\rho+o(1)$. Apply the
derivative argument to the reflected inner interval $-J$;
the standard Gaussian is symmetric. The strict sign of
$u_\rho$ is the sign of $C_\rho$ at every counted center.
This proves (A049v.1) and its full-speed version.
\end{proof}

\paragraph{Geometric translation.}
The reduced zero divisor supplies an actual derivative section
and a normal tangent section for the specified analytic shift.
The theorem bounds ordinary counting measures of the logarithmic
Hermitian norms of these sections on both oriented reduced loci.
The normal norm is compared to the full tangent norm only after
the proved fixed phase-margin deletion. This is a lower Gaussian
component of real analytic counting measures, not an equality
of measures, an order on a complex scheme, or a Galois probability
distribution. The relative reduced differential interface is
SGA 1, Expos\'e I, section 3, Proposition 1,
\url{https://grothendiecksga.com/read/sga1/fr/I-3.html}.
No multiplicity assumption or arbitrary joint independence is
added to A, and the unresolved signed half-law remains separate.


% Same NN-A049. Root-authored new continuation deduction; pending full reviews.
\subsection*{A larger region near the real parameter axis}

All statements here concern the specified denominator shift, under A.
Put $L=\log T$, $V=\log L$, $N=N_{\chi,T}$,
and $s_T=\sqrt{V/2}$.
The real parameter direction admits a stronger perturbation estimate
than a bound depending only on $|\tau|$.  This changes the region
of simplicity and separation; it does not by itself preserve a real
normal orientation throughout that region.

\begin{theorem}[NN-A049: a region of anisotropic continuation]
\label{thm:A049-anisotropic}
For every fixed primitive character and all sufficiently large $T$,
there is an absolute $c_1>0$ such that the same two sets of at least
$10^{-24}N_{\chi,T}$ centers selected in the critical inverse-height
proof have simple, pairwise disjoint branches throughout
\[
 \Omega_T=\left\{\tau\in\mathbb C:
       (1+T|\tau|)\exp(12T|\Im\tau|)
                      \le c_1\exp(s_T)\right\}.
\tag{A049r.1}
\]
The branches are holomorphic on the interior and extend locally
across each boundary point. Their ordinates remain in $(T,2T)$.
In particular the same centers survive on the real interval
\[
       |\tau|\le {c_1\over2T}\exp(s_T)
\tag{A049r.2}
\]
for sufficiently large $T$. The previously proved strict-side
guarantee remains on $0<\tau\le c/T$.

More generally, for every fixed $A_0>0$ there are constants
$c_{A_0}>0$ and
\[
 \delta_{A_0}=10^{-17}\,
       \mathbb P\{G\in[A_0+1,A_0+2]\}>0,
       \qquad G\sim N(0,1),
\tag{A049r.3}
\]
such that at least $\delta_{A_0}N_{\chi,T}$ simple centers of each
strict initial normal-velocity sign have the same simplicity,
disjointness and ordinate assertions on
\[
 (1+T|\tau|)\exp(12T|\Im\tau|)
                       \le c_{A_0}\exp(2A_0s_T).
\tag{A049r.4}
\]
These constants are independent of the fixed character, while the
starting height may depend on it and $A_0$. The sets in the second
assertion may depend on $A_0$. Each also has a strict-side interval
$0<\tau\le c'_{A_0}/T$ for some $c'_{A_0}>0$.
\end{theorem}

\begin{proof}
We first prove the perturbation estimate and then apply the two
already established center selections.

\paragraph{The normally convergent difference.}
For $\Re s>0$ and $|\tau|<1/2$, use the logarithm on
$\Re(1+\tau/n)>0$ and set
\[
 R_\chi(s,\tau)=\sum_{n\ge1}\chi(n)n^{-s}
             \bigl((1+\tau/n)^{-s}-1\bigr).
\tag{A049r.5}
\]
On each compact subset the summands are
$O(|\tau|n^{-1-\Re s})$, so the series is locally normally
convergent jointly in $s,\tau$. It is the difference
$F_\chi(s,\tau)-L(s,\chi)$: initially this follows by subtracting
the absolutely convergent series for $\Re s>1$; expanding in $\tau$
gives the same shift series, and analytic continuation gives the
identity on the stated domain. No new number-theoretic input is
needed for this identity.

Let $|s-\rho|\le h/L$, where $h\le1/10$ is fixed, and put
$\sigma_s=\Re s$, $a=\Re\tau$, $b=\Im\tau$,
$U=T|\tau|$, $v=T|b|$. For all sufficiently large $T$,
$2/5\le\sigma_s\le3/5$ and $|s|,|\Im s|\le3T$.
For $0\le u\le1$ we have
\[
 \log|1+u\tau/n|\ge-2|a|/n,\qquad
 |\arg(1+u\tau/n)|\le2|b|/n.
\tag{A049r.6}
\]
Indeed its real part is at least $1-|a|/n>1/2$, and
$\log(1-x)\ge-2x$ for $0\le x\le1/2$.
The argument bound follows from
$|\arctan y|\le|y|$ and that same real-part lower bound.
Consequently
\[
 |(1+u\tau/n)^{-s}|\le e^{1+6v},\qquad
 |(1+u\tau/n)^{-s-1}|\le e^{2+6v}.
\]
Both the trivial difference bound and the exact integral identity
\[
 (1+\tau/n)^{-s}-1
       =-{s\tau\over n}\int_0^1(1+u\tau/n)^{-s-1}\,du
\]
therefore give
\[
 |(1+\tau/n)^{-s}-1|
           \le3e^2 e^{6v}\min(1,U/n).
\tag{A049r.7}
\]

If $U\le1$, then
$\sum_{n\ge1}n^{-\sigma_s}\min(1,U/n)\le4U$.
If $U\ge1$, comparison with integrals gives
\[
 \sum_{n\ge1}n^{-\sigma_s}\min(1,U/n)
                          \le10U^{1-\sigma_s}.
\]
Here $U<T$, since $|\tau|<1/2$, and
$U^{|\sigma_s-1/2|}\le e^h<2$. Using $|\chi(n)|\le1$
proves the uniform estimate
\[
 \sup_{|s-\rho|\le h/L}|R_\chi(s,\tau)|
        \le C_0 e^{6T|\Im\tau|}\sqrt{1+T|\tau|},
                  \qquad C_0=10^4.
\tag{A049r.8}
\]
The deliberately loose constant is absolute. The derivation also
applies to negative real parameters in the stated domain.

\paragraph{The same critical-scale centers.}
Retain the constants $K_0,h,\eta$ and the exact selected sets of
Theorem~\ref{thm:A049-critical-disk}. They satisfy
\[
 M_\rho\le K_0L|D_\rho|,\quad K_0h\le1/100,\quad
 |D_\rho|\ge {L\over2\eta}\exp(\sigma/2-1),
 \qquad \sigma^2=S/2.
\]
Equation (A049g.5) gives
$S=V+O_\chi((\log V)^2)$; hence
$\sigma-s_T=O_\chi((\log V)^2/\sqrt V)=o(1)$.
For $r=h/L$, it follows eventually that
\[
        |D_\rho|r\ge d_0 e^{s_T/2},
                     \qquad d_0={h\over4e\eta}>0.
\tag{A049r.9}
\]
Choose $0<c_1\le\min(1,(d_0/(10C_0))^2)$.
The region (A049r.1) is eventually contained in $|\tau|<1/2$,
since $e^{s_T}=T^{o(1)}$. On that region, (A049r.8) is at
most $|D_\rho|r/10$. On $|s-\rho|=r$, the unperturbed Taylor
remainder is at most
$(K_0h/2)|D_\rho|r\le|D_\rho|r/200$.
Rouch\'e, with reference function $D_\rho(s-\rho)$, gives
exactly one zero counting multiplicity in each disk.
It is therefore simple. The disks are already pairwise disjoint.

The strict boundary inequality persists in a neighborhood of each
parameter point. The analytic implicit-function theorem and
uniqueness in a fixed center disk patch the local zero into a
single holomorphic branch on the whole region's interior, with
local continuation across its boundary. The discarded unit endpoint
strips and $r\to0$ preserve the ordinate window. For real parameters
in (A049r.2), $1+T|\tau|\le c_1e^{s_T}$ eventually, proving
the claimed real interval. The original strict-side guarantee is
inherited on its stated smaller interval only.

\paragraph{Larger fixed regions with smaller positive fractions.}
Apply the full proof of Theorem~\ref{thm:A049-speed-gaussian-component}
to $J=[A_0+1,A_0+2]$, of Gaussian mass $p>0$.
That proof retains a fixed phase margin, a curvature bound
$M_\rho\le K_{A_0}L|D_\rho|$, and separated simple centers,
then proves in counting measure on that retained set
\[
 {\log|D_\rho|-\log L\over s_T}=X_\rho+o(1).
\]
Delete the $o(N)$ centers where this difference has modulus
at least $1/2$. Their signed absolute weight is $o(NL)$ by the
same global second moment. The interval-restricted mass and
second-moment argument there still gives, for each sign,
\[
 {p\over72\cdot4000000^2}N>10^{-16}pN
\]
eventually, with any further fixed small slack available.
In particular at least $\delta_{A_0}N$ centers satisfy
\[
       |D_\rho|\ge L e^{A_0s_T},\qquad
       M_\rho\le K_{A_0}L|D_\rho|.
\tag{A049r.10}
\]
Choose their common $h_{A_0}>0$ as in that curvature/separation
proof, with $K_{A_0}h_{A_0}\le1/100$ and $h_{A_0}\le1/10$.
Now $|D_\rho|h_{A_0}/L\ge h_{A_0}e^{A_0s_T}$.
Taking
$c_{A_0}\le\min(1,(h_{A_0}/(10C_0))^2)$, the identical
Rouch\'e argument proves (A049r.4) with no additional deletion.
The phase margin and (A049r.10), in particular $|D_\rho|\ge L$,
are exactly the hypotheses of the critical theorem's local
direction estimate. Its bound
$C'((K_{A_0}+1)c+c/L)$ is less than half the fixed phase margin
after choosing a sufficiently small fixed $c=c'_{A_0}>0$.
This proves the asserted smaller strict-side interval.
\end{proof}

\paragraph{A useful scale consequence.}
For every fixed $B>0$, some positive fraction depending on $B$
of centers of each initial sign has simple disjoint branches on
both $|\tau|\le B\sqrt{\log\log T}/T$ and the real interval
$|\tau|\le T^{-1}\exp(B\sqrt{\log\log T})$.
For example choose $A_0=12B+1$ in (A049r.4).
In the first case its left logarithm is at most
$12B\sqrt V+O_B(\log V)$; in the second it is at most
$B\sqrt V+O(1)$, whereas the right logarithm is
$\sqrt2 A_0\sqrt V+\log c_{A_0}$.
Both inequalities follow eventually. The fraction is not asserted
uniformly as $B\to\infty$, and every displayed time still tends
to zero. No new sign-preservation claim follows at these scales.

\paragraph{Standalone geometric dictionary and exact inputs.}
The selected reduced analytic zero points give disjoint degree-one
analytic covers of the displayed parameter regions: the fixed-disk
zero count is one including multiplicity, so the relative derivative
does not vanish. Unique local sections patch without a monodromy
choice, even though the new region is not circular. This is the
analytic counterpart of the reduced relative-differential criterion
in SGA 1, Expos\'e I, section 3, Proposition 1, already audited for
NN-A049, \url{https://grothendiecksga.com/read/sga1/fr/I-3.html}.
The real interval controls the reduced fiber and its separation;
orientation is an additional real condition. The precise inputs are
the critical proof's derivative selection (A049g.19)--(A049g.20),
its local direction estimate, and the complete interval-restricted
proof (A049v.3)--(A049v.4). The normally convergent difference and
all new bounds were derived above. The statement supplies no
arithmetic descent, heat-flow identification, fixed positive-time
claim, or P versus NP/VP versus VNP consequence.


\subsection*{Uniformity over a small power of the conductor}

This addendum changes a quantifier in NN-A049. The conductor is
allowed to vary with the height. All estimates are under A, including
its auxiliary primitive Dirichlet family. No fixed-character implied
constant is used as a uniform one.

\begin{theorem}[NN-A049: uniform small-power conductors]
\label{thm:A049-uniform-conductor}
There are absolute constants $\kappa,c,T_0>0$ such that, if $T\ge T_0$
and $\chi$ is any primitive character of conductor $q\le T^\kappa$,
including conductor one, each strict normal-velocity sign has at
least $10^{-30}N_{\chi,T}$ distinct simple centers in $(T,2T)$
with the following properties. Their same denominator-shift branches
are simple, disjoint, and have ordinates in $(T,2T)$ throughout
\[
                  |\tau|\le\frac{\sqrt{\log\log T}}{100T}.
\tag{A049u.1}
\]
For every positive real $0<\tau\le c/T$ they remain on their
respective strict sides of the critical line.

For every fixed finite interval $I$ of positive length, the lower
limit as $T\to\infty$ of the infimum over all these primitive
characters of either signed counting fraction in
Theorem~\ref{thm:A049-speed-gaussian-component} is at least
\[
                      10^{-30}\mathbb P\{G\in I\},
                         \qquad G\sim N(0,1).
\tag{A049u.2}
\]
This applies separately to the normal speed, the full speed, and
the logarithmic derivative normalization stated there. Its threshold
may depend on the fixed interval, but not on the character in the
displayed conductor range. Neither a full limiting distribution
nor a half-label law is asserted.
\end{theorem}

\paragraph{Fixed parameters and error convention.}
Use the one fixed polynomial $P$ and rational $0<\nu<1/4$
chosen in NN-A047, and put
\[
 L=\log T,\quad V=\log L,\quad d=\tfrac12-2\nu>0,\quad
 \ell=\frac{\log(qT)}L,\quad N=N_{\chi,T}.
\tag{A049u.3}
\]
Coefficient bounds are exactly \eqref{eq:A048-coeff}, with their
constants, degree, and number of differentiations fixed.
Preliminary estimates hold for $2\le q\le T$ and $y\le T^\theta$,
$\theta<1/2$. Their $O_\epsilon$ constants are independent of
$q,\chi,T,y$; they may depend on the fixed coefficient and shift
parameters. We eventually choose $\kappa<d/4$ and smaller still
for a strict moment gap. For every fixed insertion degree,
$y=T^{2\nu+o(1)}$ and, after fixing a sufficiently small $\epsilon$,
\[
                     qT^{1/2+\epsilon}y=o(T)
\tag{A049u.4}
\]
uniformly. The single conductor-one case is covered by the already
proved NN-A049 source formulas and errors; its possible zeta pole
in mean-value rectangles needs no growing-conductor argument.
The zero and sampling estimates below include conductor one.

\subsubsection*{Uniform counts and horizontal contours}

\begin{lemma}[Counts, good heights, and convexity]
\label{lem:A049-uniform-counts}
Under A, uniformly for primitive characters,
\[
 \#\{\rho:|t-\Im\rho|\le1\}\ll\log(q(|t|+3)).
\tag{A049u.5}
\]
For $q\le T$,
\[
 N=\frac{T}{2\pi}\{\log(qT)+O(1)\}
                       +O(\log(qT)\log T).
\tag{A049u.6}
\]
Each of $[T,T+1]$, $[2T,2T+1]$, and $[10,11]$ contains
a height at distance at least an absolute constant divided by
$\log(q(|t|+3))$ from all zero ordinates. At these heights
$|L'/L(s,\chi)|\ll\log^2(q(|t|+3))$ on the fixed strips
used in NN-A048--A049.

For fixed $C_0$, $u=C_0/L\le1/10$,
$-u\le\sigma\le1+u$, and $T/2\le t\le3T$,
\[
 |L(\sigma+it,\chi)|\ll L(qT)^{(1+u-\sigma)/2}.
\tag{A049u.7}
\]
Shifts $O(1/L)$ and fixed derivatives are allowed at the
cost of fixed powers of $L$.
\end{lemma}

\begin{proof}
The proved uniform Hadamard identity
\eqref{eq:A048-uniform-zero-sum}, with $a=3/2$, is
\[
 \sum_\rho\frac{a}{a^2+(t-\gamma)^2}
        =\Re\frac{\Lambda'}{\Lambda}(2+it,\chi)
        \ll\log(q(|t|+3)).
\tag{A049u.8}
\]
All multiplicities and negative and central ordinates are retained.
For zeta its entire completion adds the two harmless polynomial
logarithmic derivatives. Each occurrence in a unit interval
has a fixed positive contribution, proving (A049u.5).

Integrate (A049u.8) over $[T,2T]$. The gamma and conductor
terms have integral $\frac12T\log(qT)+O(T)$. The Euler
logarithmic derivative has integral $O(1)$, by the convergent
majorant $2\sum_{n\ge2}\Lambda(n)/(n^2\log n)$ after termwise
integration. For $|\gamma|\le4T$, the Poisson-kernel integral
differs from $\pi\mathbf1_{(T,2T)}(\gamma)$ by at most
\[
 C\left(\frac1{1+|\gamma-T|}+\frac1{1+|\gamma-2T|}\right).
\]
Unit-interval counting gives $O(\log(qT)\log T)$ for their
sum. For $|\gamma|>4T$, each integral is $O(T/\gamma^2)$;
dyadic-shell use of (A049u.5) gives $O(\log(qT))$.
This proves (A049u.6), with either endpoint convention.

Removing intervals of sufficiently small radius
$c/\log(q(|t|+3))$ around the ordinates within distance one
leaves positive measure in each proposed unit interval.
At such a height subtract the genus-one logarithmic derivative
at $2+it$ from that at $\sigma+it$. There are
$O(\log(q(|t|+3)))$ near zero terms, each of that size.
The far differences are $O((t-\gamma)^{-2})$ and sum to
$O(\log(q(|t|+3)))$. The gamma difference has the latter
bound, and the Euler derivative at $2+it$ is absolutely
bounded. This proves the good-height estimate.

To prove (A049u.7), the Euler series on $1+u$ gives $CL$.
The functional equation and its two parity gamma ratios give,
on the left boundary,
\[
 |L(-u+iv,\chi)|\le CL\{q(|v|+3)\}^{1/2+u}
\]
with an absolute constant for $0<u\le1/10$. Fix the target
height $t$ and use the maximum principle on the strip for
\[
 L(s,\chi)(qT)^{(s-(1+u))/2}\exp((s-it)^2).
\]
The right-boundary bound is $O(L)$. On the left the
remaining factor is
$((|v|+3)/T)^{1/2+u}e^{-(v-t)^2}$, uniformly bounded.
The horizontal edges of exhausting rectangles tend to zero,
because the Dirichlet function has order one and the Gaussian
dominates. The exhaustion is applied to each character,
but the resulting boundary bound is uniform. For zeta
first multiply by $(s-1)/(s+1)$; this cancels its pole,
has no pole in the strip, and its inverse is bounded at
the target heights. Evaluation proves (A049u.7).
Slightly wider strips and Cauchy circles of radius $c/L$
give the final statement.
\end{proof}

\subsubsection*{Exact mean-value interfaces and every conductor loss}

\begin{lemma}[Uniform actual mean-value formulas]
\label{lem:A049-uniform-interfaces}
For $2\le q\le T$, the exact main terms in
Lemmas~\ref{lem:A048-three-shift},
\ref{lem:A048-discrete-second}, and
\ref{lem:A049-two-shift-first} have uniform error
\[
                         O_\epsilon(qT^{1/2+\epsilon}y).
\tag{A049u.9}
\]
This also bounds the imaginary part in
Lemma~\ref{lem:A048-real-bilinear}. Uniformity holds on the
fixed distinct shift circles and after fixed shift
derivatives or endpoint logarithmic weights.
\end{lemma}

\begin{proof}
The convolution, Gauss, gcd, and principal-residue identities
in NN-A048 are exact for every conductor. We expose all
size estimates in those proofs.

\paragraph{Horizontal and bounded-height terms.}
Two opposite $L$ factors on the high horizontal edges have
product at most $CL^C\sqrt{qT}$ by (A049u.7); one factor
has no larger bound. The opposite coefficient-polynomial
product is $O_\epsilon(y^{1+\epsilon}L^C)$. The good-height
logarithmic derivative costs $O(L^2)$. Thus all such terms,
including $L'^2/L=L(L'/L)^2$, are
$O_\epsilon(\sqrt{qT}\,yT^\epsilon)$.
An $O(1/t)$ gamma error on a dyadic vertical segment has
the same bound by absolute integration.

At bounded heights the same strip argument, with $T$
replaced by $3$ in its conductor factor, gives the
corresponding $q^{(1+u-\sigma)/2}L$ bound. There are
$O(\log(q+3))$ low occurrences. The good bottom in
$[10,11]$ has logarithmic derivative $O(\log^2(q+3))$.
Their contributions are $O_\epsilon(qyT^\epsilon)$,
within (A049u.9). Restoring the high endpoints costs
$O_\epsilon(\sqrt{qT}\,yT^\epsilon)$ by (A049u.5),
(A049u.7), and Cauchy, also for the shifted second
moment and multiple occurrences.

\paragraph{Stationary transform.}
Use Lemma~\ref{lem:A048-stationary}'s notation. For outer
index $k$ and inner $m$, the stationary point is
$t=2\pi m/(qk)$ and the exact scalar is
$\tau(\bar\chi)/q$, of modulus $q^{-1/2}$.
The pre-transform absolute prefactor is
$q^{c-1/2}k^{c-1}m^{-c}$, where $c=1+C_1/L$.
Small shifted powers on $m\ll qkT$ are bounded uniformly.

After this prefactor, Arias de Reyna Lemma 2.5's first
and lower-endpoint errors have factor $k^{-1/2}m^{-1/2}$:
the $q$ powers cancel against $(m/(qk))^{c-1/2}$.
Their total is at most
\[
 k^{-1/2}\sum_{m\le CqkT}\frac{\tau_R(m)}{\sqrt m}
                 \ll_\epsilon\sqrt{qT}(qkT)^\epsilon .
\tag{A049u.10}
\]
Near the upper endpoint the bound is
\[
 \sqrt{qT}(qkT)^\epsilon
 \sum_{m\asymp qkT}
 \frac1{|m-qkU/(2\pi)|+qk\sqrt T}
           \ll_\epsilon\sqrt{qT}(qkT)^\epsilon L .
\tag{A049u.11}
\]
The nearest integers are included; $qk\sqrt T\ge1$.
The remainder is bounded by a reciprocal integral.
Outside the factor-two stationary range, integration by
parts and the absolute Dirichlet norm give
$O(\sqrt{qT}L^C)$ per outer coefficient; inserting the
same prefactor in the nonstationary endpoint terms
of the cited lemma gives this bound as well.
The relative $O(1/t)$ Stirling error has absolute bound
\[
 Cq^{c-1/2}\int_{10}^{2T}t^{c-3/2}\,dt
 \sum_m|b(m)|m^{-c}
 \sum_{k\le y}|a_2(k)|k^{c-1}
       \ll_\epsilon\sqrt{qT}\,yT^\epsilon.
\]
The line $c$ is beyond every shift, so the inner norm
is $O(L^C)$. Short shift-connecting segments obey the
same bound. Restoring stationary points below the
bottom cutoff costs at most
\[
 q^{-1/2}\sum_{k\le y}\frac{|a_2(k)|}{k}
               \sum_{m\le Cqk}|b(m)|
                    \ll_\epsilon\sqrt q\,yT^\epsilon .
\]
Endpoint partial summation and fixed differentiations
insert only $\log(qkT)=O(L)$ powers. Summing over the
outer coefficients proves
$O_\epsilon(\sqrt{qT}\,yT^\epsilon)$.
The exact imported stationary interface remains the
audited Arias de Reyna Lemmas 2.2 and 2.5, with compact
amplitude constants containing no conductor parameter.

\paragraph{Auxiliary-character average.}
Write $d_1=(n,k)$, $Q=qk/d_1$, and
$x=qUk/(2\pi d_1)$. The divisor is denoted $d_1$ to
distinguish it from the fixed exponent gap $d$.
The exact factorization \eqref{eq:A048-gcd-factorization},
Parseval, and \eqref{eq:A048-polynomial-mean} are unchanged.
The previously proved uniform off-line growth
Lemma~\ref{lem:A048-uniform-growth} gives the integrand bound
\[
 (d_1Q(|\Im s|+3)T)^\epsilon
                     \tau_R(d_1)\sqrt{Q+y}.
\tag{A049u.12}
\]
Principal characters obey it on the left line and high
horizontal edges, with their poles recorded. This bound
has no hidden inducing-conductor constant.

Use Perron height $H=T^{10}$ and temporary cutoff
$\lfloor x\rfloor+1/2$, shifting to $1/2+\eta$ with
$\eta$ small relative to the final exponent.
Here $T/(2\pi)\le x\ll T^{5/2}$.
Truncation is applied after the character combination,
whose original multiplier has modulus one and whose
nonprincipal remainder has modulus at most two.
There is consequently no character-count loss.
The nearest term and absolute Dirichlet tail give
$O(d_1^\epsilon T^\epsilon x/H)$, just as in NN-A048.
The high Perron horizontal edges are negligible with
(A049u.12), since $H=T^{10}$.

The remaining vertical error is bounded by
\[
 T^\epsilon(qTk/d_1)^{1/2}\tau_R(d_1)
                                  \sqrt{qk/d_1+y}.
\tag{A049u.13}
\]
Preliminary exponents are chosen small enough to absorb
$H^\epsilon$, $x^\eta$, and fixed logarithms into the
final $T^\epsilon$. Restoring the residue cutoff costs
only $T^\epsilon d_1^\epsilon$ times fixed divisor
factors: pole order is bounded and shifts are $O(1/L)$.
It is within (A049u.13). The two complete outer sums are
\[
\begin{aligned}
&q\sqrt T\,T^\epsilon
 \sum_{k\le y}|a_2(k)|\sum_{d_1\mid k}\frac{\tau_R(d_1)}{d_1}
          \ll_\epsilon q\sqrt T\,yT^\epsilon,\\
&\sqrt{qT}\sqrt y\,T^\epsilon
 \sum_{k\le y}\frac{|a_2(k)|}{\sqrt k}
       \sum_{d_1\mid k}\frac{\tau_R(d_1)}{\sqrt{d_1}}
          \ll_\epsilon\sqrt{qT}\,yT^\epsilon .
\end{aligned}
\tag{A049u.14}
\]
We even discarded the beneficial outer $q^{-1/2}$.
The elementary divisor bounds here are exactly those
proved in NN-A048. Thus both nonprincipal contributions
and principal Perron remainders satisfy (A049u.9).

The principal expression is still
\eqref{eq:A048-principal}, with its exact real scalar.
Its residues give the same three-shift and two-shift
main terms. Conjugation and removable-pole identities
hold identically in $q$. Right-edge off-diagonal errors
are $O_\epsilon(y^{1+\epsilon}L^C)$ independently of $q$,
since character factors have modulus at most one.
For the signed derivative sum the principal term remains
real. Every other term has just been bounded.
Cauchy derivatives and gamma logarithms cost only powers
of $L$. This proves the lemma.
\end{proof}

\subsubsection*{Uniform Euler localization and main constants}

\begin{lemma}[Deleted factors and localized kernels]
\label{lem:A049-uniform-euler}
Uniformly for $q\le T$ and $|w|\le C/\sqrt L$,
\[
 Z_q(1+w)=\frac{E_q}{w}(1+O_C(|w|V)).
\tag{A049u.15}
\]
The value at zero means the residue. The Euler remainders
and all Mellin tails in the critical Gaussian proof are
uniform in $q\le T$. When (A049u.4) holds, the normalized
first and second constants are exactly
\[
 a_\ell=\ell\,a(P,\nu/\ell),\qquad
 b_\ell=\ell^2b(P,\nu/\ell).
\tag{A049u.16}
\]
The weighted value and curvature estimates consequently
have absolute constants uniformly for $q\le T^\kappa$,
once $\kappa$ is sufficiently small.
\end{lemma}

\begin{proof}
The absolute prime estimates already used in NN-A047 give
$\sum_{p\le x}p^{-1}=\log\log x+O(1)$ and
$\sum_{p\le x}(\log p)/p\ll\log x$.
Splitting primes dividing $q$ at $\log(q+3)$ gives
\[
 \sum_{p\mid q}\frac1p=O(1+\log\log\log(q+10)),\qquad
 E_q^{-1}\ll1+\log\log(q+3).
\tag{A049u.17}
\]
For the tail use $\sum_{p\mid q}\log p\le\log q$;
the small part uses the displayed absolute estimates.

For the deleted-factor logarithmic derivative in
$|w|\le C/\sqrt L$, split at $p=L^2$.
On the small part $p^{|\Re w|}\le e^{2CV/\sqrt L}$,
so its derivative sum is $O_C(V)$. On the large part
it is at most
$CL^{-2+2C/\sqrt L}\sum_{p\mid q}\log p=O_C(1)$.
All $1-p^{-1-w}$ denominators are bounded away from zero.
Integration from zero to $w$, followed by the ordinary
zeta Laurent expansion, proves (A049u.15).

The remainder products $C_1,F,G$ have local remainders
$1+O(p^{-3/2})$ on one fixed small argument neighborhood,
and every local remainder equals one at the origin.
Primes dividing $q$ simply have their factors omitted.
A common convergent dominating series therefore proves
boundedness and convergence to one uniformly over all
omitted-prime subsets, as the arguments tend to zero.

On the full Mellin lines the restricted Euler series give
$|Z_q(s)|,|Z_q(s)^{-1}|\le\zeta(\Re s)$ for $\Re s>1$.
Physical-shift pole factors are bounded by (A049u.15)
and (A049u.17). The former global bounds $L^3,L^6,L^2$
can thus all be multiplied by a fixed power $V^B$
and become uniform. The order-sixteen weights remove
$|\Im s_i|>\sqrt L$ with errors respectively
\[
 O(V^B L^{3-8})=o(L^{-1}),\quad
 O(V^B L^{6-8})=o(1),\quad O(V^B L^{2-8})=o(1).
\]
On the remaining box the relative Laurent error is
$O(V/\sqrt L)=o(1)$. All $E_q$ exponents cancel exactly:
the $F$ exponent $-1$ cancels its explicit pole, and
the full $G$ and first two-$Z$ exponents are zero.
The same polynomial-growth kernels and vertex weights
supply the integrable majorants. Cancellations precede
differentiation; the distinct outer and smaller auxiliary
circles avoid separate poles as in the curvature proof.

Set $\mathcal L=\log(qT)=\ell L$.
The very same polynomial has length
$T^\nu=\exp((\nu/\ell)\mathcal L)$, with cutoff denominator
$\nu L=(\nu/\ell)\mathcal L$.
Rescaling all coefficient and physical variables by
$\mathcal L$ instead of $L$ therefore gives exactly the
NN-A047 kernels with length parameter $\nu/\ell$.
Also $\log(qt/(2\pi))=\mathcal L+O(1)$ dyadically and
$N\sim T\mathcal L/(2\pi)$ uniformly by (A049u.6).
One physical derivative changes the normalization by
$\ell$, and two by $\ell^2$, proving (A049u.16).
All errors are uniform for $\ell$ in a compact interval
near one. The same undifferentiated bounds and Cauchy
estimate used in the curvature proof yield
\[
 \sum|B(\rho)|^4|L(\rho+z,\chi)|^2\le C_QN,\qquad
 \sum|B(\rho)|^4M_\rho^2\le C_MNL^4
\tag{A049u.18}
\]
uniformly on the fixed small shift disks.
\end{proof}

\paragraph{Choosing the conductor range.}
The finite formulas for $a(P,\xi),b(P,\xi)$ are continuous
for $\xi>0$, with the common integrable majorant already
proved in NN-A047. The certified gap at $\ell=1$ permits
a fixed $\kappa_0>0$ such that, for $1\le\ell\le1+\kappa_0$,
\[
 0<a_\ell<2,\quad0<b_\ell<2,\quad
 a_\ell^2-\frac{b_\ell}{2(1-1/2500)}>\frac1{4000}.
\tag{A049u.19}
\]
Choose $0<\kappa<\min(1/10,d/4,\kappa_0)$ once and then
a sufficiently small exponent in (A049u.9).
Arithmetic uniformity and the tail estimates were proved
before this use of continuity. This is therefore a uniform
moment gap, not an extrapolation from fixed characters.

\subsubsection*{Uniformity of the five Gaussian measures}

\begin{lemma}[Uniform marked moments and interval restrictions]
\label{lem:A049-uniform-gaussian}
Use $P_0,X_\rho,A_\rho,C_\rho,\sigma$ from the critical
Gaussian proof. Uniformly over $q\le T^\kappa$, for
every fixed integer $k\ge0$,
\[
\begin{aligned}
 N^{-1}\sum X_\rho^k&=m_k+o(1),&
 (NL)^{-1}\sum A_\rho X_\rho^k&=a_\ell m_k+o(1),\\
 (NL^2)^{-1}\sum|A_\rho|^2X_\rho^k&=b_\ell m_k+o(1),&
 N^{-1}\sum e^{2i\psi_\chi(\gamma)}X_\rho^k&=o(1),\\
 (NL)^{-1}\sum C_\rho X_\rho^k&=o(1).&&
\end{aligned}
\tag{A049u.20}
\]
For every fixed finite interval the same statements hold
with its indicator in place of $X_\rho^k$ and its Gaussian
mass in place of $m_k$.
\end{lemma}

\begin{proof}
The fixed-degree computations in the critical Gaussian
proof are exact. We check every possible conductor
dependence in their errors.

First, (A049u.17), with the same $\delta,z_0,x_0$, gives
\[
 S=V-(\log V)^2-2\log\log V+O(1+\log V)\sim V
\tag{A049u.21}
\]
uniformly. The rational Landau contour has absolute
constants for $q\le T$: (A049u.5) gives common good heights
and $O(L^2)$ logarithmic derivatives; its left gamma
logarithm is $\log(qt/(2\pi))=O(L)$. Every rational-distance
and Euler-side estimate in Lemma~\ref{lem:A049-rational-landau}
is independent of $q$. Thus its displayed error is uniform.
The prime diagonal and repeated-prime calculation are
unchanged. The off-diagonal prime-power main term is
$O_k(T\delta L(1+V)^k)$ and the rational errors are
$T^{o(1)}L^{O_k(1)}$, uniformly. This proves the first limit.

For the first and second weighted moments insert the same
degree-$K$ truncated exponentials at the endpoints. Their
lengths are $T^{2\nu+K\delta}$, their real coefficient bounds
are uniform, and $\sigma^{-1}\le1$ uniformly eventually.
Their error (A049u.9) is $o(T)$ for each fixed $K$.
The four marked local $1/p$ coefficients remain exactly
the table (A049g.14), so the insertion ratio is
$e^{z^2/2}+o(1)$ on $|\Im s_i|\le R=\delta^{-1/2}$.
Its error is bounded by a fixed power of $V$ times
$\delta RV+z_0^{-1}$, independently of the primes dividing
$q$. On all Mellin lines its fixed-degree coefficients
are $O_K((1+V)^K)$. The global tails remain negligible
after this factor. The intermediate region
$R<\max|\Im s_i|\le\sqrt L$ costs
$O(V^{B+K}R^{-14})=o(1)$; on the inner box the uniform
Laurent relative error is $O(RV/L)=o(1)$ and the
remainder products tend uniformly to one. Thus localization,
continuation, and distinct-circle differentiation are
uniform. The first $\Lambda$ convolution remains
$O_K(\delta L(1+V)^K)$ before normalization.
Equation (A049u.16) gives exactly the two weighted constants.

We expose separately the conductor factors for the even
phase. For $T^{-1}\le x\le T$, the A040 wider contour
has high horizontal factor
\[
 C(qxT/(2\pi))^{\sigma-1/2}\le C\sqrt{qxT};
\]
the smallest possible base is $1/(2\pi)$.
Its left side costs $O((qx)^{-1}L^2)$, its bounded
bottom costs $O(((qx)^{-1}+\sqrt{qx})L^2)$, and the
retained nearest-integer stationary error is
$O((qx)^{-1/2}L)$. The other stationary and Stirling
errors are $O(\sqrt{qxT}L^2)$, from the same three
Arias terms with their exact prefactors. The unaveraged
twist remainder is therefore uniformly
\[
 O\bigl((\sqrt{qxT}+(qx)^{-1}+(qx)^{-1/2})L^2\bigr).
\tag{A049u.22}
\]
For reduced $x=b/a$, $a,b\le T$, set $Q=qb$ and
$Y=qbT/(2\pi a)$. The exact A040 induced Gauss identity,
including all noncoprime zero cases, is unchanged.
The absolute GRH prime estimate of Iwaniec--Kowalski
Theorem 5.15, equation (5.56), applies to every auxiliary
character. Parseval costs $\sqrt Q$; after multiplication
by $(qx)^{-1/2}$ its error is
$O(\sqrt{qbT}\log^2(2+QY))$. Removed prime powers cost
$O(q^{-1/2}\sqrt{a/b}L^2)$. Here $Q\le T^2$,
$Y\le T^3$, and $Y\ge1/(2\pi)$.
The resulting exact rational main term is
\[
 -\frac{qT}{2\pi\varphi(q)}\sqrt{\frac ba}\,
       \frac{\chi(b)\bar\chi(a)\mu(b)}{\varphi(b)},
\]
with remainder
\[
 O\left(\left[\sqrt{qbT}+q^{-1}a/b+
                  q^{-1/2}\sqrt{a/b}\right]L^2\right).
\tag{A049u.23}
\]
Its main prefactor $q/\varphi(q)$ is $O(V)$.

Truncate the exterior ratio at $M=T^{1/2}$ and let
$Y_0=x_0^k=T^{o(1)}$ bound the marked prime products.
The fully enumerated local coefficient bound (A049g.17)
gives total main term
$O_k(T(1+V)^{k+1})=o(N\sigma^k)$.
The three weighted errors in (A049u.23) are at most
\[
 \sqrt{qT}\,Y_0^{C_k}L^{C_k},\qquad
 q^{-1}\sqrt M\,Y_0^{C_k}L^{C_k},\qquad
 q^{-1/2}Y_0^{C_k}L^{C_k},
\]
all $o(N)$ uniformly. The exterior tail and the
$O(1/T)$ phase correction remain $o(N)$ by absolute
Euler bounds. This proves the even-phase limit with
no independence premise.

For the signed moment fix $\delta_0>0$ with
$\kappa+\nu+\delta_0<1/2$. Apply the real-bilinear
part of (A049u.9) to $BP_0^j,BP_0^{k-j}$ with the
exterior series cut at $T^{\delta_0}$.
The maximum length is $T^{\nu+\delta_0+o(1)}$, so
its imaginary error is uniformly $o(NL)$.
Restoring the tail costs $T^{-\delta_0/2}$ times the
Cauchy bound from the proved weighted and unweighted
moments. The exact extra term divided by $2\gamma$
is $O(NL/T)$ by the same bound.

Finally the interval-measure argument is uniform.
Positive masses and every fixed moment are uniformly
bounded; Taylor remainders are controlled by Gaussian
moments, and the complex first and signed variations
obey the same setwise Cauchy bound. Equivalently, a
failure of a fixed-interval uniform limit gives a
sequence $T_j,\chi_j$ and a subsequence on which
$\ell_j$ converges in $[1,1+\kappa]$.
The uniform moment limits, continuity of $a_\ell,b_\ell$,
and the already proved Taylor, Fourier-inversion and
interval-approximation argument give the claimed limit
on that subsequence, a contradiction. No uniform
assertion over moving intervals is needed.
\end{proof}

\subsubsection*{Uniform derivative selection and continuation}

\begin{proof}[Proof of Theorem~\ref{thm:A049-uniform-conductor}]
We give the remaining constants and deletions explicitly.

On $I=[1,2]$, with Gaussian mass $p>1/24$, the uniform
moment limits and (A049u.19) imply
\[
 \sum_{I,\ |\cos\psi|\ge1/50}|A_\rho|
             \ge(p/20000-o(1))NL .
\]
Indeed $a_\ell-\sqrt{b_\ell/(2(1-1/2500))}>1/16000$,
since their sum is less than four.
The exact phase identity and
$1/3<|L(3/2+i\gamma,\chi)|<3$ give
$\sum_I|C_\rho|\ge(p/3000000-o(1))NL$.
The signed sum is $o(NL)$ uniformly; each sign therefore
has mass at least $pNL/10000000$ eventually.
Also $\sum|A_\rho|^2\le3NL^2$ and
$|C_\rho|\le6|A_\rho|$ eventually.

Put $\omega=10^{-11}$, $\lambda=10^{-14}$.
The two unit endpoint strips have $O(L)$ occurrences
and cost $o(NL)$ by Cauchy. Deleting
$|\cos\psi|<\lambda$ costs at most $6\sqrt3\lambda NL$.
Fix an absolute $K_0$ so large that
$6\sqrt{C_M}/K_0<\omega/10$.
Equation (A049u.18) bounds the deletion
$M_\rho>K_0L|D_\rho|$ by this latter constant times
$NL$. Every remaining nonzero weighted center is
simple. Taylor's theorem separates them by $4h/L$
for $h\le\min(c_0/16,1/(100K_0))$.

The smoothed logarithm \eqref{eq:A049-log-kernel}
is uniform here. Its residue identity uses the
uniform good-height and Hadamard bounds, whereas
its explicit constant $C_\eta$ does not involve $q$.
Trivial zeros depend only on parity; the optional
zeta pole has the same explicit vanishing error.
In the comparison with $P_0$, omitted prime terms
have squared coefficient norm
$O((\log V)^2+\log V)$ by (A049u.21), included prime
shifts cost $O(\delta^2V)$, and higher prime powers
have bounded norm. The uniform Landau estimate gives
the same zero or nonpositive prime-power off-diagonal
main terms. Remaining rational errors are
$O(T^{4\nu_0}L^C)=o(N)$ for fixed $\nu_0<1/100$.
Consequently, uniformly in the conductor range,
\[
          \sum|P_{\chi,T}(\gamma)-P_0(\rho)|^2=o(NV).
\tag{A049u.24}
\]

The full reversed zero-kernel bound is uniform as well.
The lattice sum over separated centers against one
ambient ordinate is bounded by a constant depending
only on $h$. Ambient occurrences in $[T/2,3T]$
number $O(T\log(qT))=O(N)$.
The bounded complementary range $|\gamma'|\le4T$
contributes $O(\#S/(TL))$.
On shells $2^jT<|\gamma'|\le2^{j+1}T$, $j\ge2$,
the count is $O(2^jT(L+j))$, giving total
\[
 \frac{C\#S}{TL^2}\sum_{j\ge2}2^{-j}(L+j)
                         =O(\#S/(TL)).
\]
Thus $\sum_{\rho\in S}E_T(\gamma)\le C(h)N$, including
all complementary multiplicities and infinite tails.
Chebyshev, Markov, and (A049u.24) discard only $o(N)$
centers where $|P_{\chi,T}-P_0|>\sigma/4$ or
$C_\eta E_T>\sigma/4$. They carry $o(NL)$ weight uniformly.

With fixed $\eta\le\min(c_0/4,1/(4K_0))$, the same
Taylor argument at $\rho+\eta/L$ now gives
\[
 |D_\rho|\ge\frac{L}{2\eta}e^{\sigma/2-1}\ge L,\quad
 M_\rho\le K_0L|D_\rho|,\quad
 |\Re v_\rho|/|v_\rho|\ge\lambda
\tag{A049u.25}
\]
on the retained part of $[1,2]$. Uniformity of the
threshold follows from (A049u.21).
Each sign still has mass at least $\omega NL$.
Its squared weight sum is at most $108NL^2$, giving
at least $\omega^2N/108>10^{-30}N$ centers per sign.

On their disjoint disks of radius $h/L$, the local
shift expansion has absolute $c_G,C_G,C_R$:
$c_GT\le|G(\rho)|\le C_GT$, $|G'|\le C_GT$ and
$|R|+|R_s|\le C_RT^2|\tau|^2$ for $T|\tau|\le1/10$.
These follow from absolutely convergent exterior
series and have no conductor dependence.
Choose absolute $c>0$ such that
$C_Gc+C_Rc^2\le h/10$ and the exact direction error
(A049g.21) is less than $\lambda/2$.
Rouch\'e, nonvanishing of the derivative, and unique
patching give the common $c/T$ disk and strict sides.
The terms proportional to $c,c^2$ are controlled by
this choice, not by a later height.

For the larger survival disk use the full series
(A049s.3), whose value and $s$ derivative are bounded
by absolute multiples of $e^{6T|\tau|}$.
Equation (A049u.21) gives $\sigma\ge\sqrt V/2$
uniformly eventually. The perturbation ratio to
$|D_\rho|h/L$ on (A049u.1) is
$O(e^{-19\sqrt V/100})$ uniformly.
The simple disjoint sections and ordinate condition
follow as in Corollary~\ref{cor:A049-loglog-survival}.
No additional side assertion is inferred.

For (A049u.2), fix an interval $J$ with Gaussian
mass $p>0$. The same initial calculation gives
sign mass $pNL/10000000$.
The phase deletion costs $O(\lambda pNL)$ by the
restricted moments. Choose $K_0=K_0(p)$ so its
curvature deletion is less than $pNL/10^9$.
Constants and thresholds may now depend on $p$,
but not on the character in the displayed range.
The full kernel and polynomial comparison yield,
uniformly in counting measure on the retained set,
\[
 \frac{\log|D_\rho|-\log L}{\sqrt{V/2}}=X_\rho+o(1).
\]
Every $o(N)$ failure has $o(NL)$ weight.
Retain mass $pNL/10^8$; the restricted squared
weights are at most $144pNL^2$.
Cauchy gives at least
$pN/(144\cdot10^{16})>10^{-30}pN$ centers per sign.
Use inner intervals and their increasing Gaussian
masses as in the preceding statistical theorem.
Finally the exterior Euler bounds give uniformly
$\log|v_\rho|=\log T-\log|D_\rho|+O(1)$; the phase
margin gives the same for $|u_\rho|$.
Reflect the Gaussian interval to obtain both speed
assertions. This proves every stated quantifier.
\end{proof}

\paragraph{Geometric translation and source scope.}
For every selected primitive character in the height-dependent
finite range, the reduced analytic zero divisor has disjoint
degree-one analytic covers over the stated disk. The same
absolute smaller real interval preserves its two normal
orientations. The derivative criterion is uniform in this
collection of analytic families, and is the analytic
implicit-function counterpart of SGA 1, Expos\'e I,
section 3, Proposition 1,
\url{https://grothendiecksga.com/read/sga1/fr/I-3.html}.
No common finite-type arithmetic family over the conductor
parameter, Frobenius distribution, or algebraization is asserted.

The primary interfaces remain the previously audited Arias de
Reyna Lemmas 2.2 and 2.5 for the three stationary remainders;
Iwaniec--Kowalski Theorem 5.15, equation (5.56), printed p.~114,
for the absolute auxiliary-character prime estimate; and
DLMF 25.15.5--6 and 5.11 for the functional equation and
fixed-strip gamma ratios. The coefficient-averaged auxiliary
bounds and every conductor loss have been proved here from
the full NN-A048 calculation. In the single conductor-one
case the exact Benli--Elma--Ng and Heap--Li--Zhao interfaces
retained in NN-A047--A049 suffice. Continuity is used only
after the uniform arithmetic errors and localization are proved.


\subsection*{Uniform anisotropic continuation}

\begin{corollary}[NN-A049: conductor-uniform anisotropic regions]
\label{cor:A049-uniform-anisotropic}
Retain the absolute $\kappa$ of
Theorem~\ref{thm:A049-uniform-conductor} and put
$s_T=\sqrt{(\log\log T)/2}$.
For all sufficiently large $T$, with one absolute threshold
for every primitive $q\le T^\kappa$, the same two sets of at least
$10^{-30}N_{\chi,T}$ centers selected in that theorem
have simple disjoint branches with ordinates in $(T,2T)$
throughout
\[
 (1+T|\tau|)e^{12T|\Im\tau|}\le c_1e^{s_T},
\tag{A049u.26}
\]
where $c_1>0$ is absolute. In particular they survive on
$|\tau|\le c_1e^{s_T}/(2T)$ on the real axis.
Their strict sides remain guaranteed on $0<\tau\le c/T$.

For each fixed $A_0>0$, a possibly different set of each
initial strict sign, of size at least
\[
 \delta_{A_0}N_{\chi,T},\qquad
 \delta_{A_0}=10^{-31}\mathbb P\{G\in[A_0+1,A_0+2]\}>0,
\]
has the same survival assertions on
\[
 (1+T|\tau|)e^{12T|\Im\tau|}
                     \le c_{A_0}e^{2A_0s_T}.
\tag{A049u.27}
\]
The constants and threshold may depend on $A_0$, but
not on the character in the displayed conductor range.
Each such set also retains its strict side for
$0<\tau\le c'_{A_0}/T$. For every fixed $B>0$, the scale
$B\sqrt{\log\log T}/T$ as a complex disk and every
fixed $e^{B\sqrt{\log\log T}}/T$ as a real interval
therefore have positive survival fractions depending
on $B$, uniformly over these conductors.
\end{corollary}

\begin{proof}
The direct normally convergent difference (A049r.5)
and its bound (A049r.8) use only $|\chi(n)|\le1$.
They are already uniform over all characters:
\[
 \sup_{|s-\rho|\le h/L}|F_\chi(s,\tau)-L(s,\chi)|
       \le C_0e^{6T|\Im\tau|}\sqrt{1+T|\tau|},
       \qquad C_0=10^4,\quad |\tau|<1/2 .
\]
The uniform theorem's selected centers have fixed
absolute $K_0,h,\eta$, relative curvature at most
$K_0L$, and derivative floor
$L e^{\sigma/2-1}/(2\eta)$.
Equation (A049u.21) implies
\[
 \sigma-s_T=O((\log V)^2/\sqrt V)=o(1)
\]
uniformly. Thus $|D_\rho|h/L\ge d_0e^{s_T/2}$ with
absolute $d_0>0$. Choose
$c_1\le\min(1,(d_0/(10C_0))^2)$.
On (A049u.26) the perturbation is at most
$|D_\rho|h/(10L)$; the unperturbed Taylor remainder
is at most $K_0h|D_\rho|h/(2L)$.
Exactly the count-one Rouch\'e argument of
Theorem~\ref{thm:A049-anisotropic} proves simplicity,
disjointness, and unique analytic patching. The region
is uniformly contained in $|\tau|<1/2$ eventually.

For the second assertion use the full interval-selection
proof of (A049u.2) on $J=[A_0+1,A_0+2]$.
It retains a fixed phase margin and relative curvature
$K_{A_0}L$, and its logarithmic derivative transfer is
uniform in counting measure. Delete the $o(N)$ centers
where that transfer differs from $X_\rho$ by at least
$1/2$; Cauchy's bound makes their signed mass $o(NL)$
uniformly. The count proof in (A049u.2) has fixed slack
above $10^{-30}\mathbb P(G\in J)N$.
Hence at least $\delta_{A_0}N$ centers of each sign have
\[
       |D_\rho|\ge Le^{A_0s_T},\qquad
       M_\rho\le K_{A_0}L|D_\rho|.
\]
Choose the fixed radius $h_{A_0}/L$ with
$K_{A_0}h_{A_0}\le1/100$ and set
$c_{A_0}\le\min(1,(h_{A_0}/(10C_0))^2)$.
The same perturbation and count-one argument gives
(A049u.27). The existing phase-margin direction estimate
gives $c'_{A_0}/T$, after choosing this constant small.
All thresholds are uniform by the uniform moment and
kernel arguments. Taking $A_0=12B+1$ gives the two
last scale consequences by the logarithmic comparison
already written after (A049r.4).
\end{proof}

\paragraph{Geometric scope.}
These are uniform disjoint degree-one covers of the
indicated analytic parameter regions, with no monodromy
choice because each section is the unique zero in its
fixed center disk. The reduced differential dictionary
is the same SGA 1, Expos\'e I, section 3, Proposition 1.
Orientation remains an extra real condition with its
stated smaller interval. No arithmetic family across
conductors, fixed positive time, or larger strict-side
claim is inferred.


\section{Almost linear counts for both ordinary simple-zero signs}
\label{sec:A045-simple-signs}

The weighted cancellation in NN-A042 can be combined with
isolated zeros to obtain actual counts of both ordinary
velocity signs. The derivative lower bound needed here holds
on a positive-size selected subset, rather than at every
simple zero. This fixed-radius method retains a subpower loss;
NN-A047 subsequently closes the positive-proportion endpoint
for zeta by a separate square-mollifier argument.

\begin{theorem}[NN-A045: both strict ordinary signs have almost linear counts]
\label{thm:A045-simple-signs}
Assume A. For every fixed \(\varepsilon>0\), each of
\[
 \begin{aligned}
 \#\{\rho:\ T<\Im\rho\le16T,\ \rho\text{ simple},
                                      \ \Re v_\rho>0\},\\
 \#\{\rho:\ T<\Im\rho\le16T,\ \rho\text{ simple},
                                      \ \Re v_\rho<0\}
 \end{aligned}
\]
is \(\gg_\varepsilon T^{1-\varepsilon}\log T\) as
\(T\to\infty\), where
\[
            v_\rho=\frac{\rho\zeta(\rho+1)}{\zeta'(\rho)}.
\]
The counts are of distinct simple centers. No assumption on
the density of nonsimple zeros is made.
The result does not assert a positive counting proportion,
half balance, a limiting sign frequency or transversality
at every simple zero.
\end{theorem}

\begin{lemma}[Isolated zeros retain a positive phase mass]
\label{lem:A045-isolated-phase}
Fix a primitive Dirichlet character \(\chi\), allowing
conductor one, and assume the zero-location input used in
NN-A037. There is an absolute constant \(\delta>0\) with the
following property for every sufficiently large \(T\), with
the threshold allowed to depend on \(\chi\).
Put
\[
 h_T=\frac{\delta}{\log T},\qquad
 I_T=(T,16T],\qquad N_0=N_\chi(T).
\]
Call an occurrence in \(I_T\) isolated if there is no other
zero occurrence of this same function at distance less than
\(h_T\) in ordinate. Every isolated occurrence is simple.
There are at least \(N_0/2\) isolated centers in \(I_T\) with
\[
                 |\cos\psi_\chi(\gamma)|\ge\frac14,
\]
where \(\psi_\chi\) is the exact corrected phase of NN-A042.
\end{lemma}

\begin{proof}
Use precisely the restricted pair-correlation input proved
classically by Rudnick--Sarnak, Theorem~1.2, in the degree-one,
two-point specialization already checked in NN-A037.
Its normalized ordinates are
\[
          F(\gamma)=\frac{\gamma\log\gamma}{2\pi}
                  \qquad(\gamma>0).
\]
There is a nonnegative even Schwartz function \(g\), with
\(g(0)=1\) and Fourier support strictly inside \((-1,1)\),
such that
\[
 I_g:=\int_{\mathbb R}g(u)
       \left(1-\left(\frac{\sin\pi u}{\pi u}\right)^2\right)\,du
                              <\frac{17}{50}.
\]
Indeed NN-A037 gives a complete admissible smooth
approximation of the triangular Fourier test, for which
these integrals tend to \(1/3\). Choose one test once and for
all. By continuity choose \(\delta>0\), also less than one,
such that \(g(u)\ge99/100\) for \(|u|\le\delta\).

The pair theorem, summing distinct indices with every
multiplicity retained, gives at the cutoff \(16T+1\)
\[
 \sum_{j\ne k}g(F(\gamma_j)-F(\gamma_k))
                  \le\frac7{20}N_\chi(16T+1)
\]
for all sufficiently large \(T\). In particular, repeated
ordinates are included in the sum. For ordinates within
\(h_T\) of a center in \(I_T\), both ordinates are positive
and at most \(16T+1\), once \(T\) is large. On the whole
interval \([T-1,16T+1]\),
\[
              0<F'(t)=\frac{\log t+1}{2\pi}\le\log T
\]
eventually. Thus each nonisolated occurrence in \(I_T\)
has a distinct neighboring index contributing at least
\(99/100\). Summing one such contribution per first index
counts each chosen ordered pair at most once. The number
of nonisolated occurrences in \(I_T\) is consequently at most
\[
 \frac{35}{99}N_\chi(16T+1)
                         <\frac38N_\chi(16T+1).
\]
This also excludes all multiple centers from the isolated set.

The positive-height zero count and NN-A037's exact cosine-square
mean imply
\[
 \begin{aligned}
 \# I_T&=15N_0+o(N_0),&
 N_\chi(16T+1)&=16N_0+o(N_0),\\
 \sum_{\gamma\in I_T}\cos^2\psi_\chi(\gamma)
                         &=\frac{15}{2}N_0+o(N_0).
 \end{aligned}
\]
The sums here count occurrences. Remove the nonisolated
occurrences, each contributing at most one, and the isolated
ones with absolute cosine less than \(1/4\), each contributing
at most \(1/16\). If \(G_T\) denotes the retained set, its
cardinality is at least
\[
 \frac{15}{2}N_0-\frac38(16N_0)-\frac1{16}(15N_0)
                  +o(N_0)
          =\frac9{16}N_0+o(N_0)\ge\frac12N_0.
\]
No independence of phase and spacing was used: the removal
bound is valid for the worst possible correlation.
\end{proof}

\begin{lemma}[Derivative lower bounds on most isolated centers]
\label{lem:A045-isolated-derivative}
In the setting of Lemma~\ref{lem:A045-isolated-phase}, fix
any \(a>0\). At least \(N_0/4\) of its retained centers satisfy
\[
                          |L'(\rho,\chi)|\ge T^{-a}
\]
for all sufficiently large \(T\), with the threshold allowed
to depend on \(a\) and \(\chi\).
\end{lemma}

\begin{proof}
All zeros considered below belong to this same fixed function.
Choose a fixed radius
\[
 0<r<\min\left(\frac1{32},\frac a{76},
                               \frac{\delta a}{10000}\right).
\]
The order of choices is \(\delta\), then \(a\), then \(r\),
and finally the sufficiently large height \(T\).
The GRH--Lindelof bound, its reflected left-half version and
the off-line reciprocal estimate were proved in
NN-A039, equations~\eqref{eq:A039-growth} and
\eqref{eq:A039-away-inverse}. They imply uniformly for
\(\gamma\in I_T\), at \(z_0=1/2+2r+i\gamma\),
\[
 \max_{|w|\le8r}|L(z_0+w,\chi)|\le T^{8r},
                         \qquad |L(z_0,\chi)|\ge T^{-r}.
\]
For the first estimate the leftmost real part is
\(1/2-6r\); use auxiliary growth exponent \(r\) and absorb
the fixed constants into the remaining \(T^r\).
For the second use the reciprocal estimate at the fixed
distance \(2r\), with exponent \(r/2\), and absorb constants.
The possible zeta pole is outside these high discs.

For each isolated center choose \(R\in(6r,7r)\) with no zero
on its boundary. Put \(f(w)=L(z_0+w,\chi)\), and form its
finite Blaschke product over all zeros in the disc,
with multiplicity:
\[
 B(w)=\prod_{|b|<R}\frac{R(w-b)}{R^2-\bar b w},
                              \qquad H(w)=\frac{f(w)}{B(w)}.
\]
The quotient is holomorphic and zero-free on the closed
disc. Since \(|B|=1\) on its boundary,
\(|H|\le T^{8r}\) throughout; since \(|B(0)|\le1\),
\(|H(0)|\ge T^{-r}\).
Apply the Poisson-kernel Harnack bound, locally proved in
NN-A039, to \(8r\log T-\log|H|\). On \(|w|\le3r\)
it gives
\[
                              |H(w)|\ge T^{-19r}.
\]

The selected zero is at \(b_0=-2r\) and is simple. Its own
Blaschke factor has derivative
\[
                       \frac{R}{R^2-|b_0|^2}\ge\frac1{7r}.
\]
Every other factor, evaluated at \(b_0\), has modulus at least
\[
     \frac{|b_0-b|}{R+|b_0|}
            \ge\frac{|\gamma-\gamma_b|}{9r}.
\]
Here RH puts every zero of the high disc on the critical line.
It follows that
\begin{equation}\label{eq:A045-derivative-product}
 |L'(\rho,\chi)|\ge\frac{T^{-19r}}{7r}\exp(-P_\rho),
 \qquad
 P_\rho=\sum_{\substack{\gamma'\ne\gamma\\
                        |\gamma'-\gamma|<10r}}
                         \log\frac{10r}{|\gamma'-\gamma|}.
\end{equation}
The sum retains every neighboring multiplicity. Adding zeros
outside the Blaschke disc can only enlarge this nonnegative
potential, so it remains a valid lower bound.

Let \(S_T\) be the entire isolated set in \(I_T\).
Its centers are mutually at distance at least \(h_T\), and
each is at least \(h_T\) from every other zero occurrence.
Fix a neighboring occurrence \(\gamma'\).
After omitting a coincident center if it belongs to \(S_T\),
the \(k\)-th isolated center on either side of \(\gamma'\),
ordered by distance, is at distance at least \(kh_T\).
Therefore
\[
 \begin{aligned}
 \sum_{\substack{\rho\in S_T\\\gamma\ne\gamma'}}
             \log^+\frac{10r}{|\gamma-\gamma'|}
 &\le 2\sum_{k=1}^{\lfloor10r/h_T\rfloor}
                         \log\frac{10r}{kh_T}\\
 &\le\frac{20r}{h_T}.
 \end{aligned}
\]
The last inequality follows by comparing each term with the
integral of the decreasing function
\(\log(A/x)\) on \([k-1,k]\), where \(A=10r/h_T\);
the total integral on \([0,A]\) is \(A\).

Only occurrences in \((T-10r,16T+10r)\) can contribute
to the sum of the potentials. The zero-count asymptotic gives
\(N_\chi(17T)\le18N_0\) for large \(T\). Interchange the
two finite nonnegative sums to obtain
\[
                    \sum_{\rho\in S_T}P_\rho
         \le\frac{20r}{h_T}N_\chi(17T)
         \le\frac{360r}{\delta}N_0\log T.
\]
Thus at most
\[
               \frac{720r}{\delta a}N_0
                               <\frac1{10}N_0
\]
isolated centers have \(P_\rho>(a/2)\log T\).
Discard them from the set \(G_T\) of cardinality at least
\(N_0/2\). At least \(2N_0/5\), and hence at least
\(N_0/4\), remain. At each remaining center,
\eqref{eq:A045-derivative-product} and \(19r<a/4\) give
\[
 |L'(\rho,\chi)|\ge(7r)^{-1}T^{-3a/4}\ge T^{-a}.
\]
This is an averaged lower bound on the isolated set.
It neither bounds inverse derivatives at arbitrary simple
zeros nor discards multiplicities from the neighboring sum.
\end{proof}

\begin{proof}[Proof of Theorem~\ref{thm:A045-simple-signs}]
It suffices to prove the assertion for \(0<\varepsilon<1\).
Choose \(a=\min(\varepsilon/2,1/16)\).
Use the two lemmas for the conductor-one character.
At every retained zero the absolute value of the statistic
from NN-A042 is exactly
\[
 \left|K_\rho\right|
 :=\left|\frac{\Re\{\rho\zeta(\rho+1)
                          \overline{\zeta'(\rho)}\}}{\gamma}\right|
 =\frac{|\rho|}{\gamma}|\zeta(\rho+1)|\,|\zeta'(\rho)|
                                     |\cos\psi(\gamma)|.
\]
The exact Euler lower bound
\(|\zeta(3/2+i\gamma)|\ge1/\zeta(3/2)\) therefore yields
\[
 A_T:=\sum_{T<\gamma\le16T}|K_\rho|
                              \gg_a T^{1-a}\log T.
\]
The sum may retain all occurrences: the summand at every
multiple center is zero.

Add the four dyadic estimates of NN-A042 on
\((T,2T],(2T,4T],(4T,8T],(8T,16T]\), using its exponent
parameter \(1/32\). This gives
\[
 B_T:=\sum_{T<\gamma\le16T}K_\rho
                                 =O(T^{25/32})=o(A_T),
\]
since \(a\le1/16\).
On the other hand, the arbitrary-exponent derivative bound
from NN-A041, proved by Cauchy's formula and GRH--Lindelof,
and the absolute exterior Euler product give
\[
                \max_{T<\gamma\le16T}|K_\rho|\ll_a T^a.
\]
There is no inverse-derivative estimate in this upper bound.

The positive and negative absolute masses are
\((A_T+B_T)/2\) and \((A_T-B_T)/2\), respectively.
Both are at least \(A_T/3\) eventually. Dividing each by
the displayed maximum gives
\[
       \#\{K_\rho>0\},\ \#\{K_\rho<0\}
                    \gg_a T^{1-2a}\log T
                    \ge c_\varepsilon T^{1-\varepsilon}\log T.
\]
A nonzero summand necessarily belongs to a simple center,
and its sign equals the strict ordinary velocity sign by
NN-A042. Thus these are precisely lower bounds for the
counts asserted in the theorem. Each such center is counted
once. This completes the proof.
\end{proof}

\paragraph{Exact inputs and their scope.}
The only new use of restricted pair correlation is the
small fixed isolation-radius count; the exact source is
Rudnick--Sarnak, Theorem~1.2, PDF p.~5, with the first-positive-
zero convention and normalization on PDF pp.~1--3, and
Fourier-support hypotheses in Theorem~1.1, PDF p.~4:
\url{https://empslocal.ex.ac.uk/people/staff/mrwatkin/zeta/rudnick-sarnak.pdf}.
Those statements and definitions were freshly read, while
NN-A037 contains the full admissible-test approximation used
here. The argument uses no unrestricted pair-correlation
conjecture or phase-spacing independence.
The GRH growth and reciprocal interfaces are proved in
NN-A039 from Iwaniec--Kowalski, Corollary~5.20,
equation~(5.60), printed p.~116, with all local complex-analysis
inequalities proved there. The signed-mean input is precisely
the reviewed NN-A042 specialization of Benli--Elma--Ng,
Theorem~1.1(ii), rather than an unproved complex second moment.
The ordinary zero-count asymptotic is the same one used in
NN-A037 and NN-A041.

\paragraph{Geometric and Grothendieck translation.}
At every isolated center the analytic zero algebra is reduced,
and projection of the denominator-shift divisor to the
parameter disc defines a unique local section. Vanishing
relative differentials gives the usual finite-algebra
unramified criterion, SGA~1, Expos\'e~I, section~3,
Proposition~1. Its normal tangent is \(\Re v_\rho\).
The statistic \(K_\rho\) is that normal tangent paired with
the positive derivative norm \(|\zeta'(\rho)|^2/\gamma\);
it vanishes at every nonreduced center.
The Blaschke factorization retains the complementary divisor
through its nonnegative logarithmic distance potential.
Spacing of the selected reduced centers bounds the sum of
those potentials even when the complementary divisor has
arbitrary multiplicities. The resulting lower absolute mass
and signed cancellation force the stated almost linear
counts on both real normal-orientation loci.
This is an ordered real analytic counting assertion, not an
algebraic order on complex schemes, a pure motive,
a Frobenius distribution theorem or a complexity separation.
The residual loss \(T^\varepsilon\) leaves positive ordinary
sign proportions and half balance unresolved.
The subsequent NN-A046 supplies the fixed-character signed
mean needed to extend these ordinary sign counts to every
fixed primitive Dirichlet function under the stated
fixed-modulus GRHs, and gives the zeta conclusion under RH alone.


\section{A primitive Dirichlet signed statistic and ordinary signs}
\label{sec:A046-dirichlet-signed}

The exterior value in the normal velocity permits an absolutely
convergent integer-twist expansion. After reflection, the
principal character part of each twist is real. This gives a
direct contour proof of cancellation, with no growing-length
mollifier theorem.

\begin{theorem}[NN-A046: fixed-character ordinary signed counts]
\label{thm:A046-dirichlet-signed}
Assume A and fix a primitive Dirichlet character \(\chi\) of
conductor \(q\), including conductor one. For every
\(\varepsilon>0\),
\begin{equation}\label{eq:A046-signed-mean}
 \sum_{0<\gamma\le T}
  \frac{\Re\{\rho L(\rho+1,\chi)
                    \overline{L'(\rho,\chi)}\}}{\gamma}
                       =O_{\chi,\varepsilon}(T^{1/2+\varepsilon}),
 \qquad \rho=\frac12+i\gamma.
\end{equation}
Every zero occurrence may be retained: a multiple center
contributes zero. At a simple center the summand is
\[
             \frac{|L'(\rho,\chi)|^2}{\gamma}\Re v_{\rho,\chi},
 \qquad
             v_{\rho,\chi}=\frac{\rho L(\rho+1,\chi)}{L'(\rho,\chi)}.
\]
Consequently, for every \(\varepsilon>0\), each strict
ordinary sign occurs at least
\[
                         c_{\chi,\varepsilon}
                                  T^{1-\varepsilon}\log T
\]
times at distinct simple centers in \((T,16T]\), for all
sufficiently large \(T\), with no nonsimple-density assumption.

For the signed-mean proof it is enough to assume GRH for the
primitive characters inducing the characters modulo \(q\).
In particular, for \(\zeta\) the signed estimate and ordinary
counts follow under RH alone. Conductors are fixed throughout;
neither positive ordinary sign proportions nor half balance
is asserted.
\end{theorem}

\subsection{Two uniform estimates used in the contour}

Write
\[
 a(n)=(\Lambda*\log)(n)
        =\sum_{d\mid n}\Lambda(d)\log(n/d).
\]
Then \(a(1)=0\), and
\begin{equation}\label{eq:A046-coefficients}
             0\le a(n)\le(\log n)^2\quad(n\ge2),
\end{equation}
because \(\sum_{d\mid n}\Lambda(d)=\log n\).

\begin{lemma}[Cancellation of a nonprincipal convolution]
\label{lem:A046-convolution}
Let \(\psi\ne\psi_0\) be a character modulo the fixed \(q\).
Under GRH for its inducing primitive character, for every
\(\nu>0\) and \(Y\ge1\),
\[
             \sum_{n\le Y}\psi(n)a(n)
                          \ll_{q,\nu}(1+Y)^{1/2+\nu}.
\]
Also
\[
 \sum_{n\le Y}\psi(n)=O_q(1),\qquad
 \sum_{n\le Y}\psi(n)(\log n)^2=O_q(\log^2(2Y)).
\]
\end{lemma}

\begin{proof}
The last two bounds follow from periodicity, the vanishing
sum of a nonprincipal character over a full period, and
partial summation. For the first, the exact Dirichlet series is
\[
 F_\psi(s)=\frac{L'(s,\psi)^2}{L(s,\psi)}
                     =\sum_{n\ge1}\frac{\psi(n)a(n)}{n^s}
                       \qquad(\Re s>1).
\]
Under the stated GRH it is holomorphic on \(\Re s>1/2\).
Indeed the inducing nonprincipal \(L\)-function has no pole
and has no zero there; the finitely many deleted Euler factors
of an imprimitive character have zeros only on \(\Re s=0\).
For every fixed \(\eta,b>0\), the growth, Cauchy derivative,
and reciprocal bounds of NN-A039 give
\begin{equation}\label{eq:A046-perron-growth}
        F_\psi(\sigma+it)\ll_{q,\eta,b}(1+|t|)^b
                   \qquad(1/2+\eta\le\sigma\le2).
\end{equation}
These bounds extend to compact \(t\)-intervals by continuity
and nonvanishing. Finite Euler factors preserve them.

Here are the necessary truncated Perron details. Replace \(Y\)
by \(y=\lfloor Y\rfloor+1/2\), which leaves its integer
summation range unchanged. Bounded \(y\) may be absorbed, so
assume \(y\ge3\). Put \(c=1+1/\log y\) and \(H=y^2\).
For \(x>0\), the integral kernel
\[
      \frac1{2\pi i}\int_{c-iH}^{c+iH}\frac{x^s}{s}\,ds
\]
differs from \({\bf1}_{x>1}\) by
\[
              O\!\left(x^c
                     \min\left(1,\frac1{H|\log x|}\right)\right)
                       \quad(x\ne1).
\]
The infinite vertical integral is the indicated step
function, by closing the contour on the decaying side.
Integration by parts on each omitted tail gives
\(O(x^c/(H|\log x|))\). The uniform \(O(x^c)\) alternative
follows by pairing positive and negative imaginary parts:
the cosine term has integrable denominator \(c^2+t^2\),
and the sine term is uniformly bounded after rescaling
\(u=t|\log x|\), splitting at \(u=1\), and integrating
the remaining sine integral by parts.

Apply this kernel with the absolutely convergent series on
\(\Re s=c\). For \(n\in[y/2,2y]\), half-integrality gives
\(|n-y|\ge1/2\) and
\(|\log(y/n)|\gg |n-y|/y\).
The total error from this range is
\[
               O\left(\log^2(2y)\frac yH\log(2y)\right).
\]
Outside it, \(|\log(y/n)|\) is bounded below, and
\eqref{eq:A046-coefficients} bounds the error by
\[
       O\left(\frac{y^c}{H}
                       \sum_{n\ge2}\frac{(\log n)^2}{n^c}\right)
                  =O\left(\frac{y^c}{H}\log^3y\right).
\]
Both are \(O(\log^3(2y))\). Thus the partial sum equals
the truncated Perron integral plus this admissible error.

Shift that integral to \(\sigma_0=1/2+\eta\), without
crossing any pole. The new vertical integral is
\(O_{q,\eta,b}(y^{\sigma_0}H^b)\), and the two horizontal
integrals are \(O_{q,\eta,b}(y^cH^{b-1})\), by
\eqref{eq:A046-perron-growth}.
Choose \(0<\eta<\min(\nu/4,1/8)\) and \(0<b<\nu/8\),
reducing \(b\) further if necessary to have \(b<1/4\).
Since \(H=y^2\), these estimates and the Perron error are
\(O_{q,\nu}(y^{1/2+\nu})\). This proves the lemma with
all parameters fixed before \(Y\).
\end{proof}

\begin{lemma}[A coefficient-uniform stationary transform]
\label{lem:A046-stationary}
Fix \(\chi\), and let \(t_0>2\pi\) be fixed. For large \(U\)
put \(c=1+1/\log U\), and write
\[
 \Delta_\chi(s)=\epsilon_\chi
                \frac{Q_\chi(1-s)}{Q_\chi(s)},\qquad
                  \ell(t)=\log\frac{qt}{2\pi}.
\]
If the complex coefficients satisfy
\(|d_N|\le\log^2(2N)\), then, for \(k=0,1,2\),
\begin{align}\label{eq:A046-stationary}
 &\frac1{2\pi}\int_{t_0}^{U}
       \Delta_\chi(c+it)^{-1}\ell(t)^k
                           \sum_{N\ge1}d_NN^{-c-it}\,dt\\
 &\hspace{12mm}
   =\frac{\tau(\bar\chi)}q
         \sum_{N\le qU/(2\pi)}d_N(\log N)^k e^{-2\pi iN/q}
                   +O_\chi(\sqrt U\log^7 U).\nonumber
\end{align}
The constant is independent of every coefficient and of
any support restriction on them. A convention for the
last term at an integer cutoff changes only this error.
\end{lemma}

\begin{proof}
For \(k=0\), use the exact Stirling expansion
\eqref{eq:A037-stirling}. Its relative \(O(1/t)\) error
contributes at most
\[
 C_\chi\left(\sum_N\log^2(2N)N^{-c}\right)
                    \int_{t_0}^{U}t^{c-3/2}\,dt
                         \ll_\chi\sqrt U\log^3 U.
\]
In the main integral put \(u=t/(2\pi)\), \(X=U/(2\pi)\),
and \(r=N/q\). Its oscillatory phase is
\(2\pi u\log(u/(er))\), and its power amplitude is
\(u^{c-1/2}\), exactly as in NN-A037 with \(b=q\).
The finitely many \(r\) below the fixed lower endpoint
contribute an admissible \(O_\chi(\log^3U)\) error when
their main terms are inserted.

For clarity the complete error summation does not require a
mean bound on \(d_N\). Split each interval into
\([r/2,2r]\) and its complement. On the complement the
phase derivative has modulus at least \(\log2\);
integration by parts with monotone power amplitude gives
\(O(X^{c-1/2})\). Summing with absolute coefficients gives
\(O_\chi(\sqrt U\log^3U)\).
On the middle interval Arias de Reyna's Lemma~2.5 gives the
three remainders displayed and proved summable in NN-A037.
The first and lower-endpoint terms, after multiplication
by the coefficient, sum to at most
\[
                    C_\chi\sum_{N\le C_\chi U}
                             \frac{\log^2(2N)}{\sqrt N}
                         \ll_\chi\sqrt U\log^2 U.
\]
The upper-endpoint terms away from \(r\asymp X\) have the
same bound. In the remaining range their sum is at most
\[
 C_\chi\sqrt U\log^2U
       \sum_{qX/2<N<2qX}
                   \frac1{|N-qX|+q\sqrt X}
                        \ll_\chi\sqrt U\log^3U.
\]
This estimate is uniform even if \(qX\) is an integer.
Thus the base transform has error
\(O_\chi(\sqrt U\log^5U)\), allowing spare logarithms.

The stationary main coefficient is
 \(i^a\epsilon_\chi^{-1}/\sqrt q\).
Since \(\epsilon_\chi=\tau(\chi)/(i^a\sqrt q)\) and
\(\tau(\chi)\tau(\bar\chi)=\chi(-1)q\), it is exactly
\(\tau(\bar\chi)/q\). Its phase is \(e^{-2\pi iN/q}\).
This verifies the constant, including complex characters.

For \(k=1,2\), apply ordinary partial summation in the upper
integration endpoint to the \(k=0\) formula, keeping \(c\)
fixed at its value determined by \(U\).
The same argument at any endpoint \(V\in[t_0,U]\) has
error \(O_\chi((1+\sqrt V)\log^5U)\): here
\(V^{c-1}\le U^{c-1}=e\) and the coefficient norm is still
\(O(\log^3U)\).
Since \(\ell'(V)=1/V\), the endpoint error and its
integral are \(O_\chi(\sqrt U\log^{5+k}U)\).
At the stationary endpoint \(V=2\pi N/q\), the weight is
\(\ell(V)^k=(\log N)^k\). This proves the stated formula.
\end{proof}

\subsection{The integer twists and their real principal terms}

\begin{proof}[Proof of Theorem~\ref{thm:A046-dirichlet-signed}]
Write \(L(s)=L(s,\chi)\).
Choose \(U\in(T,T+1)\) at distance
\(\gg_\chi1/\log T\) from every zero ordinate, as in
NN-A037. Put \(c=1+1/\log U\). Take a fixed bottom
\(t_0>2\pi\) avoiding zero ordinates.
The partial-fraction argument in NN-A037 gives
\[
       \frac{L'}L(\sigma+iU)=O_\chi(\log^2T)
                       \qquad(1-c\le\sigma\le c).
\]
At the fixed bottom this logarithmic derivative is uniformly
bounded on the same compact real interval for all large \(T\).

For integers \(1\le m\le T^2\), define
\[
                  S_m(U)=\sum_{0<\gamma\le U}m^\rho L'(\rho).
\]
Integrate
\[
                   \frac{L'(s)^2}{L(s)}m^s
\]
over the positively oriented rectangle with vertical sides
\(c\) and \(1-c\), bottom \(t_0\), and top \(U\).
At a simple zero its residue is \(m^\rho L'(\rho)\).
At a zero of order \(j\ge2\), its local order is \(j-2\ge0\),
so it has no residue. Thus this is the exact occurrence sum
with multiple-center terms zero, apart from finitely many
lower zeros, whose total is \(O_\chi(\sqrt m)\) under GRH.
All trivial zeros, gamma singularities and the possible zeta
pole lie below the rectangle.

For any fixed \(b>0\), the Cauchy and reflected growth bounds
used in NN-A041 give
\[
                 |L'(\sigma+iU)|
          \ll_{\chi,b} T^{\max(1/2-\sigma,0)+b}
                          \qquad(1-c\le\sigma\le c).
\]
For \(1\le m\le T^2\), this implies that the top edge is
\[
               O_{\chi,b}\bigl(T^b(\sqrt T+m)\log^2T\bigr).
\]
Indeed on the left half
 \(m^\sigma T^{1/2-\sigma}
   =\sqrt T(m/T)^\sigma\)
is bounded by \(O(\sqrt T+\sqrt m)\), and on the right
half \(m^\sigma\le m^c\ll m\), since \(m\le T^2\).
The bottom edge is \(O_\chi(m)\).

On the right vertical edge the absolutely convergent
Dirichlet series is
\[
                  \frac{L'(s)^2}{L(s)}
                        =\sum_{n\ge1}\chi(n)a(n)n^{-s}.
\]
The \(n=m\) term contributes
\((U-t_0)\chi(m)a(m)/(2\pi)\).
The remaining terms are \(O(m\log^3(2T))\), uniformly in
the displayed range of \(m\). To see this, their absolute
integrated bound is
\[
   C m^c\sum_{n\ne m}
               \frac{(\log(2n))^2 n^{-c}}{|\log(m/n)|}.
\]
For \(n\notin[m/2,2m]\), use a constant lower bound for
the logarithm and
\(\sum_n(\log(2n))^2n^{-c}\ll\log^3U\).
In the remaining range use
\(|\log(m/n)|\ge |m-n|/(2m)\) and
\((m/n)^c\ll1\); the harmonic distance sum gives
\(O(m\log^3(2m))\). Since \(m^c\ll m\) and
\(\log(2m)\ll\log T\), the assertion follows.

Let \(J_m\) denote the downward left-edge contribution.
Set \(w=c+it\) and \(D_\chi=\Delta_\chi'/\Delta_\chi\).
The functional equations
\[
 L(s,\chi)=\Delta_\chi(s)L(1-s,\bar\chi),\qquad
       \Delta_\chi(s)\Delta_{\bar\chi}(1-s)=1
\]
give the exact conjugated expression
\begin{equation}\label{eq:A046-left}
 \overline{J_m}
   =-\frac m{2\pi}\int_{t_0}^{U}\Delta_\chi(w)^{-1}
       \left(D_\chi(w)^2L(w)-2D_\chi(w)L'(w)
                         +\frac{L'(w)^2}{L(w)}\right)m^{-w}\,dt.
\end{equation}
In particular, conjugation of a complex character has not
been discarded. It changes the left-half dual back into
\(\chi\) at \(w\).

The differentiated Stirling formula is
\(D_\chi(c+it)=-\ell(t)+O_\chi(1/t)\).
Its error in \eqref{eq:A046-left} is
\(O_\chi(m\sqrt U\log^4U)\): absolute convergence on
\(\Re w=c\), the bounds for
\(\sum n^{-c}\) and \(\sum(\log n)n^{-c}\), and
\(\int t^{c-3/2}\log(t+2)\,dt\) give this directly.
In the remaining three terms apply
Lemma~\ref{lem:A046-stationary}. Relabel \(N=mn\) and
put zero coefficients at \(N\) not divisible by \(m\).
The three coefficient sequences are respectively
\(\chi(n)\), \(-\chi(n)\log n\), and \(\chi(n)a(n)\);
each is bounded by a fixed constant times \(\log^2(2N)\),
uniformly in \(m\). The lemma therefore yields
\begin{align}\label{eq:A046-left-main}
 \overline{J_m}
 &=-\frac{m\tau(\bar\chi)}q
      \sum_{n\le qU/(2\pi m)}\chi(n)e^{-2\pi imn/q}
               f_m(n)+O_\chi(m\sqrt U\log^7U),\\
 f_m(n)&=(\log m)^2-(\log n)^2+a(n).\nonumber
\end{align}
Indeed the stationary weight is \(\log(mn)\), and
\((\log(mn))^2-2\log(mn)\log n
   =(\log m)^2-(\log n)^2\).

Only integers \(m\) coprime to \(q\) will be used below.
On the finite group \((\mathbb Z/q\mathbb Z)^*\), write
\[
 \chi(n)e^{-2\pi imn/q}
                  =\sum_{\psi\bmod q}c_\psi(m)\psi(n),
 \qquad
 c_\psi(m)=\frac1{\varphi(q)}
       \sum_{r\bmod q}^{*}
          \chi(r)e^{-2\pi imr/q}\overline{\psi(r)}.
\]
Both sides vanish for nonunits. The coefficients have modulus
at most one. The exact principal coefficient is
\[
       c_{\psi_0}(m)=\frac{\tau(\chi)\bar\chi(-m)}{\varphi(q)}.
\]
The primitive Gauss identity thus gives
\[
 -\frac{m\tau(\bar\chi)}q c_{\psi_0}(m)
                            =-\frac{m\bar\chi(m)}{\varphi(q)}.
\]
Since \(f_m(n)\) and \(\psi_0(n)\) are real, this principal
contribution to \(\overline{J_m}\), after conjugation and
multiplication by \(\bar\chi(m)\), is exactly real.
The right-edge diagonal has the same reality property.
This is the required cancellation; no asymptotic formula
for either real principal term is needed.

For every nonprincipal \(\psi\), put \(Y=qU/(2\pi m)\).
Lemma~\ref{lem:A046-convolution} and its two elementary
companions imply, for fixed \(\nu>0\),
\[
   \left|\sum_{n\le Y}\psi(n)f_m(n)\right|
       \ll_{q,\nu}\log^2(2T)\left(1+Y^{1/2+\nu}\right)
                         \quad(Y\ge1).
\]
For \(Y<1\) the sum is empty. There are only
\(\varphi(q)-1\) such characters; their inducing GRH
hypotheses are supplied by A. This controls the entire
nonreal arithmetic part of \eqref{eq:A046-left-main}.

\subsection*{Absolute Euler summation and the final signs}

Take \(M=\lfloor T^2\rfloor\). The exterior value has the
absolutely convergent expansion
\[
 \overline{L(\rho+1,\chi)}
              =L(2-\rho,\bar\chi)
              =\sum_{m\ge1}\bar\chi(m)m^{\rho-2},
\]
whose tail after \(M\) is \(O(M^{-1/2})\), uniformly at
every critical-line zero. By the arbitrary-exponent
derivative bound and the zero count, this gives
\begin{equation}\label{eq:A046-euler-truncation}
 \sum_{0<\gamma\le T}L'(\rho,\chi)\overline{L(\rho+1,\chi)}
     =\sum_{m\le M}\frac{\bar\chi(m)}{m^2}S_m(T)
                          +O_{\chi,b}(T^b\log T).
\end{equation}
The finitely many low zeros are harmless here.
Changing from \(S_m(T)\) to \(S_m(U)\) costs at most
\(O_{\chi,b}(\sqrt m\,T^b\log T)\) per twist; its weighted
sum is \(O_{\chi,b}(T^b\log T)\).

All previously estimated errors are summable with these
absolute weights, uniformly as \(M\) grows. More explicitly,
the horizontal and bottom errors sum to
\[
 O_{\chi,b}\bigl(T^b(\sqrt T+\log(2M))\log^2T+\log(2M)\bigr).
\]
The right-edge errors sum to \(O(\log^4T)\), and the
stationary errors sum to \(O_\chi(\sqrt T\log^8T)\).
The nonprincipal arithmetic contributions are at most
\[
 C_{q,\nu}\log^2(2T)
    \left(\sum_{m\le M}\frac1m
          +T^{1/2+\nu}\sum_{m\ge1}m^{-3/2-\nu}\right)
       \ll_{q,\nu}T^{1/2+\nu}\log^2(2T)+O_q(\log^3T).
\]
Each principal contribution is exactly real after its
Euler weight. Choosing \(b,\nu\) sufficiently small in
terms of \(\varepsilon\), and absorbing fixed powers of
logarithms, proves
\begin{equation}\label{eq:A046-imaginary-mean}
 \Im\sum_{0<\gamma\le T}
       L'(\rho,\chi)\overline{L(\rho+1,\chi)}
                        \ll_{\chi,\varepsilon}T^{1/2+\varepsilon}.
\end{equation}

Write the individual product in this display as \(A_\rho+iB_\rho\).
Direct multiplication by \(\rho=1/2+i\gamma\) gives
\[
  \frac{\Re\{\rho L(\rho+1,\chi)\overline{L'(\rho,\chi)}\}}{\gamma}
                         =B_\rho+\frac{A_\rho}{2\gamma}.
\]
The sum of the latter terms is
\(O_{\chi,b}(1+T^b\log T)\), by Euler boundedness,
the derivative upper bound and partial summation of
\(N_\chi(t)\ll_\chi t\log(t+2)\).
This proves \eqref{eq:A046-signed-mean}, including all
endpoint conventions and multiplicities.

Finally apply the two isolated phase and derivative lemmas
of NN-A045 to this same fixed primitive character.
They provide lower absolute statistic mass
\(\gg_{\chi,a}T^{1-a}\log T\) on \((T,16T]\), for every
fixed \(a>0\). The Euler lower bound used there remains
\(|L(3/2+i\gamma,\chi)|\ge1/\zeta(3/2)\).
The supremum bound is \(O_{\chi,a}(T^a)\), while
\eqref{eq:A046-signed-mean} gives a signed sum
\(O_{\chi,\eta}(T^{1/2+\eta})\).
Choose \(a=\min(\varepsilon/2,1/16)\) and, for example,
\(\eta=1/16\). Both signed absolute masses are eventually
at least one third of their total. Dividing each by the
supremum gives the stated lower count
\(\gg_{\chi,\varepsilon}T^{1-\varepsilon}\log T\).
Every nonzero statistic is at a simple center and has the
ordinary velocity sign. For \(q=1\), the nonprincipal
character sum is absent; all inputs used in this proof and
the isolated lemmas then require RH alone.
\end{proof}

\paragraph{Exact source interfaces.}
The functional equation and Gauss normalization are
DLMF~25.15.5--25.15.6, together with the primitive Gauss
identity in Montgomery--Vaughan, \emph{Multiplicative Number
Theory I}, section~9.2, Theorem~9.7, printed p.~287,
already checked for NN-A040.
The only imported stationary estimate is Arias de Reyna,
Lemma~2.5, \url{https://arxiv.org/html/1403.7126v1};
all new coefficient-uniform summation is proved above.
The growth and reciprocal bounds are the completely written
NN-A039 interfaces, based on Iwaniec--Kowalski,
Corollary~5.20, equation~(5.60), p.~116.
Rudnick--Sarnak's restricted correlation theorem and its
exact normalization enter only through the fully written
NN-A045 lemmas.

Kobayashi, \emph{On a certain sum of the derivatives of
Dirichlet \(L\)-functions}, arXiv:2003.05629v2,
Theorem~1.1, PDF p.~2, proves a real principal asymptotic
for the unweighted derivative sum; its Lemma~2.3,
PDF p.~5, records the corresponding Gonek transform:
\url{https://arxiv.org/pdf/2003.05629}.
Those exact interfaces were read as a comparison. Neither
an unstated uniform-twist version of that theorem nor the
complete external proof is imported here. There is no
historical novelty assertion about those classical contour
ingredients.

\paragraph{Geometric and Grothendieck translation.}
The factor \(L'(\rho,\chi)\) annihilates the statistic at
every nonreduced zero center. On the reduced analytic locus,
the specified denominator-shift divisor has a unique local
section, and its real normal tangent is paired with the
positive derivative norm. SGA~1, Expos\'e~I, section~3,
Proposition~1 supplies the finite relative-differentials
criterion for this reduced local situation.
The finite character expansion in the proof is the actual
Fourier decomposition on the cyclotomic group
\((\mathbb Z/q\mathbb Z)^*\). Its principal component is
real after the primitive Gauss pairing; the remaining
components are controlled by GRH and the displayed Perron
argument. These concrete rank-one arithmetic characters
therefore control a real analytic signed statistic, without
asserting an arithmetic scheme for the denominator shift
or a Frobenius distribution of its normal signs.
The conclusion counts both ordinary reduced-center
orientations with a subpower loss; it does not give
positive proportions, purity or a complexity separation.


\section{Positive signed-center proportions for Dedekind zeta functions}
\label{sec:A043-dedekind-density}

This extension uses the nonnegative prime-ideal coefficients of
$\zeta_K$. It needs no Artin holomorphy assertion and no
independence of zeros of its factors. Every multiplicity is
retained throughout. The hypothesis here is the Riemann
hypothesis for the specified Dedekind zeta function; input M
supplies it conditionally, and is not promoted to input A.

\begin{theorem}[NN-A043: both Dedekind branch classes have positive density]
\label{thm:A043-density}
Let $K$ be a number field of degree $n$, and assume that every
nontrivial zero of $\zeta_K$ lies on $\Re s=1/2$.
Consider the specified denominator shift
\[
 (\zeta_K)_\tau(s)=\sum_{\mathfrak a\ne0}
                      (N\mathfrak a+\tau)^{-s}
\]
and its local meromorphic continuation near a nontrivial zero.
A center belongs to $R_K$ if at least one of its Puiseux
branches has positive first real normal coefficient. It belongs
to $L_K$ if a branch has negative first real normal coefficient,
or if it is a vertical double whose branches have negative
second real normal coefficient. These are precisely the
NN-A039 branch classes, now for $\zeta_K$ at sufficiently
large heights.

There is a constant $C_n\ge1$, depending only on $n$, such that
for each fixed field $K$ and all sufficiently large $T$,
\begin{equation}\label{eq:A043-counts}
 R_K((T,2T)),\ L_K((T,2T))
       \ge\frac{T\log T}{3072\,n^3 C_n}.
\end{equation}
The threshold for $T$ may depend on $K$.
In particular both distinct-center classes have positive lower
proportion even relative to the count of all zero occurrences,
whose dyadic asymptotic is $nT\log T/(2\pi)$.
No numerical value of $C_n$, uniform height threshold over
fields, ordinary-simple-sign density, or half-label law is
asserted. Multiple centers can belong to both classes.
\end{theorem}

\begin{proof}
Write the signature as $(r_1,r_2)$ and the absolute
discriminant as $D_K$. Use the meromorphic completion
\[
 Q_K(s)\zeta_K(s),\qquad
 Q_K(s)=D_K^{s/2}\Gamma_{\mathbb R}(s)^{r_1+r_2}
                         \Gamma_{\mathbb R}(s+1)^{r_2},
 \quad \Gamma_{\mathbb R}(s)=\pi^{-s/2}\Gamma(s/2).
\]
It has simple poles at zero and one; multiplication by
$s(s-1)$ makes it entire of order one, with its actual
nontrivial zeros. The functional equation has sign $+1$.
These exact completion facts are given, for example, by
Greni\'e--Molteni, \emph{An explicit Chebotarev density theorem
under GRH}, equations (2.5)--(2.8), author p.~5, specialized to
the trivial Hecke character. We also use their Theorem~1.1,
author p.~2, with the trivial extension $K/K$, which gives
\begin{equation}\label{eq:A043-PNT}
 \psi_K(x)=x+O_K\bigl(\sqrt{x}\log^2(2x)\bigr).
\end{equation}
This is an imported prime-ideal theorem, not a new finding.

\paragraph{The crossing argument, including all multiple centers.}
Let
\[
 Z_K(t)=e^{i\theta_K(t)}\zeta_K(1/2+it)\in\mathbb R,
 \qquad \theta_K(t)=\arg Q_K(1/2+it),
\]
with a continuous argument. Put
\begin{align*}
 \alpha_K(t)&=\Im\log\zeta_K(3/2+it),
   &\delta(t)&=\arctan(1/(2t)),\\
 \psi_K(t)&=\theta_K(t)+\alpha_K(t)-\delta(t),
   &c_K(t)&=\cos\psi_K(t).
\end{align*}
The Euler logarithm is single-valued on $\Re s>1$.
If $\Lambda_K(m)$ is defined by
$-\zeta_K'/\zeta_K(s)=\sum_m\Lambda_K(m)m^{-s}$, then
$0\le\Lambda_K(m)\le n\Lambda(m)$; this bound is proved
again below from residue degrees. It gives
\[
 |\alpha_K'(t)|\le n\sum_{m\ge2}\Lambda(m)m^{-3/2}<\infty.
\]
Stirling's formula gives
$\theta_K'(t)=\frac12\log D_K+\frac n2\log(t/(2\pi))+O_n(t^{-1})$.
Thus $\psi_K'>0$ at all sufficiently large heights, and
the zeros of $c_K$ there are simple.

The local shift expansion of NN-A004--NN-A005 gives, at
$\rho=1/2+i\gamma$ of order $m$,
\[
 s_j(\tau)=\rho+c_j\tau^{1/m}+O(\tau^{2/m}),\qquad
 c_j^m=\frac{m!\rho\zeta_K(\rho+1)}{\zeta_K^{(m)}(\rho)}\ne0.
\]
For $m\ge2$, NN-A035 proves that every vertical first
direction bends strictly left at second order once
$\gamma^2\ge n+2$. In particular, an interval missing
$R_K$ can contain, besides simple zeros, only vertical doubles;
an interval missing $L_K$ contains no multiple center.

Here are the local crossing calculations, so that this
application does not require an unstated simplicity assumption.
Set $g=Z_K/c_K$, canceling common real analytic factors.
At a simple zero with $c_K(\gamma)\ne0$,
\[
 \Re v_\rho=-\frac{E_\rho c_K(\gamma)}{Z_K'(\gamma)},
 \qquad E_\rho=|\rho|\,|\zeta_K(\rho+1)|>0,
 \qquad g'(\gamma)=\frac{Z_K'(\gamma)}{c_K(\gamma)}.
\]
A simple tangency has $c_K(\gamma)=0$ and its two
length-one factors cancel completely. At a vertical double,
writing $A_2=\zeta_K''(\rho)/2$, one has
\[
 e^{i\theta_K}A_2=-Z_K''(\gamma)/2,
 \qquad c^2=-\frac{2iE_\rho e^{i\psi_K(\gamma)}}{Z_K''(\gamma)}<0.
\]
Hence $c_K(\gamma)=0$, $Z_K''(\gamma)\sin\psi_K(\gamma)<0$,
and after one cancellation the quotient has a simple zero with
\[
 g'(\gamma)=\frac{Z_K''(\gamma)}{-2\sin\psi_K(\gamma)
                                       \psi_K'(\gamma)}>0.
\]
If $R_K$ is absent, all zeros of $g$ therefore cross from
negative to positive. If $L_K$ is absent, they all cross from
positive to negative. Each component between its simple real
poles contains at most one such zero. With all local orders
counted, this says
\[
 N_Z(I)-N_c(I)=N_g(I)-P_g(I)\le1,
 \qquad N_c(I)\le\frac{\psi_K(v)-\psi_K(u)}\pi+1
 \quad (I=(u,v)).
\]
On the positive-height right half-strip the assumed zero
location gives an Euler-continuation logarithm of $\zeta_K$.
The real Hardy identity and the jump of $m\pi$ at an
order-$m$ zero give exactly
\[
 N_Z((u,v))=\frac{\theta_K(v)-\theta_K(u)}\pi
                                      +S_K(v)-S_K(u),
 \qquad S_K(t)=\pi^{-1}\arg\zeta_K(1/2+it).
\]
The pole at height zero is outside this domain. Consequently,
for $0<h\le1$ and nonzero endpoints, absence of either
class in $(u,u+h)$ implies
\begin{equation}\label{eq:A043-crossing}
                 S_K(u+h)-S_K(u)\le2+B_nh
\end{equation}
with a constant $B_n$ depending only on the degree. The
excluded endpoint set is discrete and has measure zero.

\paragraph{The two pole terms in the explicit formula.}
Use exactly the interval majorants $F_-,F_+$ for
$[0,b]$, where $b=C/(2\pi)$ and bandwidth $\eta=1/100$,
from the proof of NN-A042's microscopic moment theorem.
We need their values at two nonreal arguments, and prove the
required decay rather than extrapolating a real-axis estimate.

Write the Beurling functions displayed in Section~2 of
Carneiro--Chandee--Milinovich as $H^\pm=H\pm K_0$, where
\[
 H(z)=\left(\frac{\sin\pi z}{\pi}\right)^2
       \left(\sum_{m\in\mathbb Z}
                \frac{\operatorname{sgn}m}{(z-m)^2}+\frac2z\right),
 \qquad K_0(z)=\left(\frac{\sin\pi z}{\pi z}\right)^2.
\]
For $x=\Re z>0$, the partial-fraction identity
$\sum_{m\in\mathbb Z}(z-m)^{-2}=\pi^2/\sin^2(\pi z)$
(DLMF~4.22.4) gives
\[
 H(z)=1+\left(\frac{\sin\pi z}{\pi}\right)^2
       \left(-2\sum_{m\ge1}(z+m)^{-2}-z^{-2}+2/z\right).
\]
Comparing the sum with
$\int_0^\infty(z+u)^{-2}\,du=1/z$ yields
\[
 \left|\sum_{m\ge1}(z+m)^{-2}-\frac1z\right|
 \le\int_0^\infty2|z+u|^{-3}\,du\le x^{-2}.
\]
For example, this follows by integrating
$(z+m)^{-2}-(z+u)^{-2}$ on each $m-1\le u\le m$
and bounding the derivative. Since
$|\sin\pi z|\le e^{\pi|\Im z|}$, it follows that
\[
 |H^\pm(z)-1|\le\frac4{\pi^2x^2}e^{2\pi|\Im z|}.
\]
The identity $H^\pm(-z)=-H^\mp(z)$ handles $x<0$.
In
\[
 F_\pm(z)=\frac12\{H^\pm(\eta z)+H^\pm(\eta(b-z))\},
\]
the two constants cancel outside the interval. Thus, if
$d=\operatorname{dist}(x,[0,b])>0$,
\begin{equation}\label{eq:A043-complex-tail}
 |F_\pm(x+iy)|\le\frac4{\pi^2\eta^2d^2}e^{2\pi\eta|y|}.
\end{equation}
This also directly verifies the required decay on each fixed
horizontal strip.

Put $L=\log T$ and
$f_t(u)=F_\pm((u-t)L/(2\pi))$.
The two pole contributions in the general explicit formula
are $f_t(1/(2i))+f_t(-1/(2i))$. By
\eqref{eq:A043-complex-tail}, uniformly for $T\le t\le2T$,
\begin{equation}\label{eq:A043-poles}
 |f_t(1/(2i))|+|f_t(-1/(2i))|
       =O_\eta\bigl(T^{-2+\eta/2}L^{-2}\bigr)=o(1).
\end{equation}
Indeed the normalized real part is $-tL/(2\pi)$ and the
absolute imaginary part is $L/(4\pi)$. Both poles, with their
actual order one, have been retained before being estimated.

\paragraph{Reduction to a prime polynomial.}
The general explicit formula (4.6), author p.~7, of
Carneiro--Chandee--Milinovich applies with pole order one,
degree $n$, conductor $D_K$, and spectral parameters zero
repeated $r_1+r_2$ times and one repeated $r_2$ times.
Their Example~10, p.~13, expressly includes $\zeta_K$.
For real admissible $f$ it reads
\[
 \sum_\gamma f(\gamma)
 =f(1/(2i))+f(-1/(2i))
   +\frac1{2\pi}\int_{\mathbb R}f(u)G_K(u)\,du
   -\frac1\pi\Re\sum_{m\ge2}\frac{\Lambda_K(m)}{\sqrt m}
                           \widehat f(\log m/(2\pi)),
\]
where $G_K(u)=2\Re(Q_K'/Q_K)(1/2+iu)$ and every zero
occurrence is included. In particular
\[
 G_K(u)=\log D_K+n\log(|u|/(2\pi))+O_n(1/|u|)
 \quad (|u|\ge1),
\]
and it is bounded on compact real intervals.
The archimedean-tail calculation in NN-A042 now gives
\[
 \frac1{2\pi}\int f_t(u)G_K(u)\,du
       =\frac{G_K(t)}{L}(b\pm\eta^{-1})+o_{K,C}(1).
\]
For clarity, after $u=t+2\pi x/L$, on
$|x|\le TL/(4\pi)$ the density difference is
$O_n(|x|/(TL)+1/T)$; its integral against $F_\pm$ is
$O_{n,C}(1+\log(TL))/(TL)$ before the extra $1/L$.
On the complement, the real decay in
\eqref{eq:A043-complex-tail} bounds the density integral by
$O_{K,C}(\log(TL)/(TL))$. Both vanish. Also
$G_K(t)/L=n+o_K(1)$ uniformly in the indicated dyadic range.

Set $N=T^\eta$ and $X(t)=S_K(t+C/L)-S_K(t)$. The
pointwise interval sandwich, exact phase-count identity and
\eqref{eq:A043-poles} give
\[
 D_-(t)-n\eta^{-1}+o_{K,C}(1)
       \le X(t)\le D_+(t)+n\eta^{-1}+o_{K,C}(1),
\]
where
\[
 D_\pm(t)=-2\Re\sum_{m\le N}
       \frac{\Lambda_K(m)}{L\sqrt m}
       \widehat F_\pm(\log m/L)m^{-it}.
\]
If $p\mathcal O_K=\prod_j\mathfrak p_j^{e_j}$ and
$f_j$ are the residue degrees, then
\[
 \Lambda_K(p^r)=\log p\sum_{j:f_j\mid r}f_j,
 \qquad \sum_j e_jf_j=n.
\]
This proves $0\le\Lambda_K(p^r)\le n\log p$, including
ramified primes. In particular
$a_K(p):=\Lambda_K(p)/\log p$ is the integer counting prime
ideals of norm $p$, with $0\le a_K(p)\le n$.

The exact $L^1$ Fourier error for $F_\pm$ is $\eta^{-1}$,
independent of $C$. Replace their Fourier transforms by that
of $\mathbf1_{[0,b]}$ and then remove prime powers of exponent
at least two. The elementary divisor-weighted fourth-moment
bound \eqref{eq:A042-direct-DPbound}, with coefficients bounded
by $n$ times the corresponding NN-A042 coefficients, proves
\begin{equation}\label{eq:A043-decomposition}
 \begin{aligned}
 X&=P+Y,\qquad \|Y\|_4=O_n(1)+o_{K,C}(1),\\
 P(t)&=-2\Re\sum_{p\le N}A_p p^{-it},\\
 A_p&=\frac{a_K(p)(1-p^{-iC/L})}{2\pi i\sqrt p}.
 \end{aligned}
\end{equation}
The norm uses normalized Lebesgue measure on $[T,2T]$.
For the Fourier replacement, the divisor-weighted square sum
is $O(n^2)$, since the prime part is bounded by
$2n^2(\eta L)^{-2}\sum_{p\le N}(\log p)^2/p$ and the
higher powers by a convergent series times $n^2L^{-2}$.
Here $\sum_{p\le x}\log p\ll x$ is the unconditional
Chebyshev bound, with the exact classical locator
Iwaniec--Kowalski, equation (2.12), p.~33; thus this estimate
does not add a hypothesis about the zeros of $\zeta$.
For the removed higher powers it is bounded by
$n^2\pi^{-2}\sum_p\sum_{r\ge2}(r+1)/(r^2p^r)$.
Both coefficient absolute sums are $O_n(\sqrt N)$, and
the fourth-moment error is $O_n(N^4/T)=o(1)$.
Thus the constant in $O_n(1)$ is independent of $C$ and $K$.

\paragraph{A variance lower bound from prime ideals alone.}
Removing the terms $p^r$ with $r\ge2$ from
\eqref{eq:A043-PNT} gives
\begin{equation}\label{eq:A043-degreeonePNT}
 \sum_{p\le x}a_K(p)\log p\sim x.
\end{equation}
Indeed the removed terms total at most
$n\sum_{r\ge2,p^r\le x}\log p=O_n(\sqrt{x}\log^2(2x))$.
Let $V_T=\sum_{p\le N}|A_p|^2$. Unique factorization and
the complete diagonal calculation in
\eqref{eq:A042-direct-prime-diag} give
\[
 \mathbb EP^2=2V_T+o(1),\qquad
 \mathbb EP^4=12V_T^2-6\sum_{p\le N}|A_p|^4+o(1).
\]
The off-diagonal error is $O_n(N^6/T)=o(1)$, and
$\sum|A_p|^4=O_n(1)$.
Since $a_K(p)$ is a nonnegative integer at most $n$,
\[
                   a_K(p)\le a_K(p)^2\le n a_K(p).
\]
Partial summation in \eqref{eq:A043-degreeonePNT} therefore
implies, for each fixed $C$,
\begin{equation}\label{eq:A043-variancebounds}
 \begin{aligned}
 \frac{J_\eta(C)}{\pi^2}+o_{K,C}(1)
     &\le2V_T\le\frac{nJ_\eta(C)}{\pi^2}+o_{K,C}(1),\\
 J_\eta(C)&=\int_0^\eta\frac{1-\cos(Cu)}u\,du.
 \end{aligned}
\end{equation}
At the lower endpoint, the discarded primes $p\le T^\epsilon$
contribute at most
\[
 \frac{n^2C^2}{2L^2}\sum_{p\le T^\epsilon}
                          \frac{(\log p)^2}{p}
                   =O_n(C^2\epsilon^2).
\]
Take $T\to\infty$ and then
$\epsilon\downarrow0$. Thus no uniform prime ideal estimate
as $K$ varies is being assumed. As proved in NN-A042,
$J_\eta(C)=v+O(1)$ for $v=\log(2+C)$.

Combining \eqref{eq:A043-decomposition} with H\"older gives
\begin{align*}
 \mathbb EX^2&\ge v/\pi^2-O_n(1+\sqrt v)+o_{K,C}(1),\\
 \mathbb EX^2&\le nv/\pi^2+O_n(1+\sqrt v)+o_{K,C}(1),\\
 \mathbb EX^4&\le3n^2v^2/\pi^4+O_n(1+v^{3/2})+o_{K,C}(1).
\end{align*}
Consequently one can choose $C=C_n$ sufficiently large,
depending only on the degree, so that for every fixed $K$
and sufficiently large $T$,
\[
 av\le\mathbb EX^2\le bv,\qquad \mathbb EX^4\le dv^2,
 \quad a=\frac1{2\pi^2},\quad b=\frac{2n}{\pi^2},
 \quad d=\frac{6n^2}{\pi^4}.
\]
The classical $S_K(t)=O_K(\log(t+2))$ follows, for example,
from CCM Theorem~5 and Example~10 under the assumed RH.
It implies
\[
 T\mathbb EX=\int_{2T}^{2T+h}S_K(t)\,dt
                    -\int_T^{T+h}S_K(t)\,dt=O_K(C_n),
 \qquad h=C_n/L.
\]
Hence $\mathbb EX\to0$.

\paragraph{The positive tail and the covering count.}
The full H\"older argument used in NN-A042 now gives
\[
 \mathbb E|X|\ge A_n\sqrt v,
 \qquad A_n=\frac{a^{3/2}}{\sqrt d}
                       =\frac1{4\sqrt3\,\pi n}.
\]
Increase $C_n$ so that $3\le A_n\sqrt v/8$.
For large $T$, $|\mathbb EX|\le A_n\sqrt v/4$, so
$\mathbb E(X\mathbf1_{X>3})\ge A_n\sqrt v/4$.
Cauchy's inequality yields
\[
            \mathbb P(X>3)\ge\frac{A_n^2}{16b}
                                  =\frac1{1536\,n^3}.
\]
For large $T$, \eqref{eq:A043-crossing} and $B_nh<1$
imply that every such interval contains a center of both
classes, outside the measure-zero endpoint set. Restrict the
basepoint to $[T,2T-h]$; its good subset has measure at least
$T/(3072n^3)$. A given distinct center belongs to intervals
$(t,t+h)$ for a set of basepoints of length at most $h$.
The covering bound gives \eqref{eq:A043-counts}.

Finally, the exact counting-phase identity, Stirling's formula
and $S_K(t)=O_K(\log(t+2))$ imply
$N_K(2T)-N_K(T)\sim nT\log T/(2\pi)$, with every
zero occurrence counted. This proves the stated proportions.
All thresholds are taken after fixing the field. No step
asserts that its multiple zeros disappear or that the two
classes are disjoint.
\end{proof}

\paragraph{Geometric and Grothendieck translation.}
The integers $a_K(p)$ count the degree-one closed points of
$\operatorname{Spec}\mathcal O_K$ over the rational prime $p$.
Their positivity and the prime-ideal theorem, rather than a
purported Frobenius law for zeros, supply the variance lower
bound. On the analytic side, each isolated length-$m$ zero
of the specified shifted divisor is a finite locally free
analytic algebra over the parameter disc. Its normalization
after $\tau=u^m$ gives the Puiseux normal-jet classes. The
quotient of its real Hardy section by the corrected phase
cosine retains oriented Cartier-divisor multiplicities: a
vertical double loses exactly one phase factor. The resulting
real crossing index and the prime polynomial moments prove
the positive center proportions. SGA~1, Expos\'e~I, section~3,
Proposition~1 and section~4, Proposition~10 give the finite
differentials and discriminant dictionary. No arithmetic
family with these denominator shifts, pure motive, order on
general complex schemes, or complexity separation is asserted.

\paragraph{Exact additional primary locators.}
The prime-ideal input and completion are Greni\'e--Molteni
(2019), Theorem~1.1, author p.~2, and equations~(2.5)--(2.8),
author p.~5, at
\url{https://sites.unimi.it/molteni/research/papers-pdf/44-molteni-An_explicit_Chebotarev_density_theorem_under_GRH.pdf}.
The explicit-formula inputs are CCM Lemma~2 and Lemma~4,
author p.~3, equation~(4.6), author p.~7, and Example~10,
author p.~13, at
\url{https://home.olemiss.edu/~mbmilino/new_S(t).pdf}.
The elementary partial fraction is
\url{https://dlmf.nist.gov/4.22.E4}.
These classical sources are inputs; the new deduction is the
two-sided branch-center density for the specified deformation.


\section{The Newman boundary and exact positive heat matrices}

Use the normalization of the supplied workbench22,
\[
 H_t(z)=\int_0^\infty e^{tu^2}\Phi(u)\cos(zu)\,du,
 \qquad H_0(z)=\tfrac18\xi(1/2+iz/2),
 \qquad \partial_tH_t=-\partial_z^2H_t.
\]
The kernel $\Phi$ is positive and decays doubly exponentially. The
de Bruijn--Newman threshold property is imported as a known theorem:
all zeros of $H_t$ are real exactly when $t\ge\Lambda$. Rodgers--Tao,
arXiv:1801.05914v5, Theorem 1.1, proves $\Lambda\ge0$; A gives
$\Lambda\le0$. We use their intersection $\Lambda=0$, without claiming
that either theorem supplies simplicity at the endpoint.

\begin{theorem}[NN-A006: endpoint collision or escape to infinity]
Under A, every zero of $H_t$ is real and simple for $t>0$.
The zeros of $H_0$ are real, with their possible multiplicities retained.
If $t_j<0$, $t_j\uparrow0$, and $z_j$ is a nonreal zero of $H_{t_j}$,
then every finite accumulation point of $z_j$ is a multiple real zero
of $H_0$. In particular, if all zeta zeros are simple, such nonreal zeros
escape every compact subset of $\mathbb C$.

At a real multiplicity-$m$ zero $(t_0,z_0)$ the local branches are
\[
 z_j(t)=z_0+\sqrt{2(t-t_0)}\,h_{m,j}+O(|t-t_0|),
\]
where $h_{m,j}$ are the distinct real roots of the probabilists' Hermite
polynomial $\operatorname{He}_m$. For $m\ge2$, a real collision can occur only
at $t_0\le0$ under A. Replacing this by $t_0<0$ requires simplicity at
$t=0$, which A does not state.
\end{theorem}
\begin{proof}
The integral is jointly entire in $(t,z)$ locally, by the kernel decay,
so all compact-set convergences and derivatives used here are uniform.
Write $H_{t_0}(z)=a_m(z-z_0)^m+O((z-z_0)^{m+1})$.
The PDE gives
$\partial_t^\ell H=(-1)^\ell\partial_z^{2\ell}H$.
Set $t-t_0=u^2/2$ and $z-z_0=uy$; divide by $a_mu^m$.
The resulting analytic function equals at $u=0$
\[
 \sum_{\ell=0}^{\lfloor m/2\rfloor}
 \frac{m!(-1)^\ell y^{m-2\ell}}
 {2^\ell\ell!(m-2\ell)!}=\operatorname{He}_m(y).
\]
Gaussian differentiation and Rolle's theorem give $m$ real roots;
the Hermite differential equation excludes repeated roots. The implicit
function theorem at each root, with Rouche's theorem for completeness,
gives the stated expansion. For negative real $t-t_0$, at least two
of the nonzero Hermite roots give nonreal branches. If $t_0>0$, nearby
negative increments still have $t>0$, contradicting $\Lambda=0$.
This proves simplicity at all positive times.

If $z_j$ has a finite accumulation point $z_*$, uniform convergence
implies $H_0(z_*)=0$, and A makes $z_*$ real. If that zero were simple,
the implicit function theorem would give a unique local zero for all
small real $t$. Reality of the coefficients means its conjugate is
another local zero, so uniqueness forces the branch to be real. This
contradicts the nonreal $z_j$. Compactness gives the final escape claim.
\end{proof}

The theorem gives an exact dichotomy and no height law. In particular,
$\Lambda=0$ alone does not prove an exponential collision front or an
independent-pair collision approximation. In the completed arithmetic
coordinate $s=1/2+iz/2$, the double-zero law becomes
$s-s_0=\pm i\sqrt{(t-t_0)/2}+O(|t-t_0|)$; the normalization cannot
be omitted when comparing gap-based heuristics. The geometric local
model is a ramified double cover of the parameter disc (higher
Hermite branching for $m>2$), not a family with unramified individual
zero sections at the collision.

\begin{theorem}[NN-A009: strict heat positivity and determined moments]
Let the complete arithmetic expression $Z(a)$ be the one in
workbench22, equation \path{eq:rh-connes-arithmetic-heat}, including
both pole terms, the archimedean integral, and the prime-power sum.
Under A it has the exact representation
\[
 Z(a)=\sum_\rho e^{-a\gamma_\rho^2},\qquad a>0,
\]
with nontrivial zeros counted with multiplicity. Consequently:
\begin{enumerate}
\item $(-1)^kZ^{(k)}(a)>0$ for every $k\ge0$, and
$(\log Z)''(a)>0$.
\item If $a_1,\ldots,a_d>0$ are distinct, the matrix
$[Z(a_i+a_j)]_{i,j=1}^d$ is positive definite.
\item For every $d\ge1$, the matrix
$[(-1)^{i+j}Z^{(i+j)}(a)]_{i,j=0}^{d-1}$ is positive definite.
\item The positive finite measure
$\mu_a=\sum_\rho e^{-a\gamma_\rho^2}\delta_{\gamma_\rho^2}$
is uniquely determined among positive measures on $[0,\infty)$ by its
moments $(-1)^kZ^{(k)}(a)$.
\end{enumerate}
No simplicity assumption is needed.
\end{theorem}
\begin{proof}
The supplied unrestricted explicit formula is
$Z(a)=\sum_\rho e^{a(\rho-1/2)^2}$, absolutely and locally uniformly
convergent with all derivatives. A replaces $\rho-1/2$ by $i\gamma$.
There are infinitely many distinct positive values $\gamma^2$, since
an entire function has finitely many zeros in each compact set and
$\xi$ has infinitely many zeros. None is zero: $\zeta(1/2)\ne0$,
for example from the alternating eta series.
Termwise differentiation proves strict complete monotonicity.
After dividing the weights by $Z(a)$,
$Z''/Z-(Z'/Z)^2$ is the variance of $\gamma^2$, and is strictly
positive because its support has more than one point.

For the first matrix the quadratic form is
$\sum_\rho|\sum_i c_i e^{-a_i\gamma^2}|^2$.
A nonzero real linear combination of $d$ exponentials with distinct
real exponents has at most $d-1$ distinct real zeros: divide by one
exponential, differentiate, and use Rolle's theorem inductively on
any chosen finite interval containing the zeros. If a complex linear
combination vanished at infinitely many points, its real and imaginary
parts would each vanish identically by this fact; independence of the
exponentials then makes every coefficient zero. Thus the quadratic form
vanishes only at $c=0$. The second matrix has quadratic form
$\sum_\rho e^{-a\gamma^2}|\sum_{i=0}^{d-1}c_i\gamma^{2i}|^2$.
A nonzero polynomial has finitely many roots, proving strictness.

For determinacy, $\int e^{(a/2)x}\,d\mu_a=Z(a/2)<\infty$.
Any positive measure $\nu$ with the same moments has the same finite
exponential moment by monotone convergence of its nonnegative Taylor
series. Its moment generating function therefore agrees with that of
$\mu_a$ in $|z|<a/2$, by dominated termwise integration. Both Laplace
transforms are holomorphic on $\Re z<a/2$; the identity theorem gives
equality on that half-plane, hence equality of characteristic functions
on the imaginary axis. Fourier uniqueness for finite measures yields
$\nu=\mu_a$.
\end{proof}

The finite matrices are genuine Gram matrices of the displayed functions
in $L^2(\mu)$ or $L^2(\mu_a)$. An approximate matrix
$\widetilde G$ with a certified operator-norm error $\epsilon$ obeys
$\lambda_{\min}(\widetilde G)\ge\lambda_{\min}(G)-\epsilon$.
Under A any rigorously certified negative exact value is impossible;
an observed negative value must therefore fail either its error
enclosure, its admissibility, or the retained complete formula.
Strictness gives no uniform positive lower bound as the matrix size,
parameters, or conditioning vary. This is an analytic Hilbert-space
Gram realization. It is not a finite-dimensional pure Weil object,
and the tautological diagonal operator with spectrum $\gamma^2$
does not construct a new arithmetic Hilbert--Polya operator.


\section{A total-positivity boundary and two actual character constructions}

\begin{theorem}[NN-A007: the Newman threshold is a PF boundary]
Define the entire function
\[
 F_t(w)=\frac{H_t(i\sqrt w)}{H_t(0)}=\sum_{k\ge0}c_k(t)w^k,
 \qquad
 c_k(t)=\frac{\int_0^\infty e^{tu^2}\Phi(u)u^{2k}\,du}
 {(2k)!H_t(0)}.
\]
The square-root notation denotes the entire even-power series, so it
introduces no branch choice. Every $c_k(t)$ is strictly positive for
real $t$. Under A, the infinite Toeplitz matrix
\[
 T_t=(c_{j-i}(t))_{i,j\ge0},\qquad c_k=0\quad(k<0),
\]
is totally nonnegative if and only if $t\ge0$. Thus for each $t<0$
some finite minor is negative, although every entry on or above its
diagonal is positive. At each $t\ge0$,
\[
 F_t(w)=\prod_j(1+w/x_j(t)^2),
\]
where $x_j(t)>0$ lists the positive zeros of $H_t$, with multiplicities.
\end{theorem}
\begin{proof}
Positivity and convergence follow directly from the theta integral.
The order estimate $\log\max_{|z|\le R}|H_t(z)|=O_t(R\log R)$
in workbench22, \texttt{sec:rh-newman-collision}, implies that $F_t$
has order at most $1/2$. When all zeros of $H_t$ are real, pairing
$x_j$ with $-x_j$ gives the displayed genus-zero product; the possible
zero-free exponential factor is constant because the order of $F_t$
is less than one, and $F_t(0)=1$ fixes it. Moreover
$\sum_jx_j^{-2}<\infty$.

Each factor $1+bw$, $b\ge0$, has upper-bidiagonal Toeplitz matrix
$I+bS$, which is totally nonnegative: every nonzero minor has its
order-preserving matching and nonnegative product. Multiplication of
generating series is multiplication of their Toeplitz matrices.
Cauchy--Binet therefore gives nonnegative minors for every finite
product. Local uniform convergence gives coefficient convergence and
hence convergence of each finite determinant, proving the forward
direction without invoking a classification theorem.

For the converse use the known Edrei theorem, in the precise original
form of Proposition A on pp.86--87 and its proof in Section 3, p.92,
of Edrei, \emph{Proof of a conjecture of Schoenberg}, Canad. J. Math.
5 (1953), 86--94, DOI 10.4153/CJM-1953-010-3: an entire generating
function of a totally nonnegative sequence has all its zeros on the
negative real axis. A zero $z$ of $H_t$ gives a zero $w=-z^2$ of
$F_t$. If all such $w$ are negative real, every $z$ is real. Thus PF
is equivalent to all-real zeros. The threshold property and A together
with Rodgers--Tao give exactly $t\ge0$. Failure of total nonnegativity
means, by its definition, failure at a finite minor; coefficients are
real, so that minor is strictly negative.
\end{proof}

The distinction between positive coefficients and all minors is material:
entrywise positivity persists throughout the nonreal-zero phase. The
negative certificate has finite size for each fixed negative time, but
this argument supplies neither its size nor an algorithmic rate as
$t\uparrow0$. It therefore gives no complexity upper bound.

\begin{theorem}[NN-A008: zeta parameters in symmetric, unitary, and flag geometry]
Under A, let $\gamma_j>0$ list the positive zeta ordinates with
multiplicities and put
\[
 C=\sum_j\gamma_j^{-2},\qquad
 \beta_j=\gamma_j^{-2}/C,\qquad
 E(w)=\prod_j(1+\beta_jw)=\sum_{k\ge0}e_k(\beta)w^k.
\]
Then $C$ is finite and positive, every $0<\beta_j<1$, and
$\sum\beta_j=1$. The following constructions are well defined.
\begin{enumerate}
\item A normalized central positive-definite character of $S_\infty$ is
\[
 \chi_S(\sigma)=\prod_{k\ge2}
 \left((-1)^{k-1}\sum_j\beta_j^k\right)^{m_k(\sigma)}.
\]
It has Thoma parameters $(\alpha,\beta,\gamma)=(0,\beta,0)$.
\item A normalized continuous central positive-definite character on
$U(\infty)=\bigcup_nU(n)$ is
\[
 \chi_U(U)=\prod_j\det((1-\beta_j)I_n+\beta_jU),\qquad U\in U(n).
\]
These are two different characters, not restriction of a single formula.
On a permutation cycle of length $k$, the latter has factor
\[
 \prod_j\{(1-\beta_j)^k-(-\beta_j)^k\},
\]
whereas the former has factor $(-1)^{k-1}\sum_j\beta_j^k$.
\item Every finite upper-triangular Toeplitz truncation
$T^{(n+1)}=(e_{j-i}(\beta))_{0\le i,j\le n}$ has strictly positive
minors whenever upper triangularity does not force them to vanish.
For $1\le i\le n$, put
\[
 D_i=\det T^{(n+1)}_{\{0,\ldots,i-1\},\{n+1-i,\ldots,n\}},
 \quad D_0=D_{n+1}=1,\quad
 q_i=\frac{D_{i-1}D_{i+1}}{D_i^2}>0.
\]
These are literal positive coordinates in the finite Toeplitz
parametrization associated with quantum cohomology of the complete flag
variety $\mathrm{Fl}_{n+1}$, by the specified Rietsch theorem below.
\end{enumerate}
\end{theorem}
\begin{proof}
The paired Hadamard product for $\Xi(z)=\xi(1/2+iz)$ and the zero-count
bound give $\sum\gamma_j^{-2}<\infty$; its infinitely many positive
terms give $0<\beta_j<1$.
For (1), let $D$ be the positive trace-class diagonal operator with
eigenvalues $\beta_j$ on $\ell^2$, and let $U_\sigma$ permute the
factors of $(\ell^2)^{\otimes n}$. Direct summation in the tensor basis
gives
\[
 \operatorname{Tr}(D^{\otimes n}U_\sigma)
 =\prod_{k\ge1}(\operatorname{Tr}D^k)^{m_k(\sigma)}.
\]
Multiplication by $\operatorname{sgn}\sigma$ gives exactly the stated
formula, since $\operatorname{Tr}D=1$. It is compatible with adjoining
fixed points. For any finite complex coefficients $a_i$ and permutations
$\sigma_i$, its positive-definiteness is
\[
 \sum_{i,j}\overline{a_i}a_j\chi_S(\sigma_i^{-1}\sigma_j)
 =\operatorname{Tr}(D^{\otimes n}A^*A)\ge0,
 \quad A=\sum_j a_j\operatorname{sgn}(\sigma_j)U_{\sigma_j}.
\]
The tensor density commutes with all permutations, so the character is
central. The displayed parameters coincide with the standard Thoma
formula; its extremality is the known classification theorem, with exact
modern primary locator Okounkov, arXiv:math/9803037, Section 2.3,
Theorem (Thoma), printed p.23. No classification is needed for the
explicit existence proof just given.

For (2), at a fixed $0\le b\le1$ the determinant expands as
\[
 \det((1-b)I_n+bU)
 =\sum_{k=0}^n b^k(1-b)^{n-k}\operatorname{Tr}(\bigwedge^kU).
\]
It is a nonnegative combination of representation characters, normalized
at the identity because the dimensions sum to $(b+1-b)^n=1$.
Products remain positive definite. For fixed $n$, uniform convergence
follows from $|\det(I_n+\beta_j(U-I_n))-1|\le C_n\beta_j$ and
$\sum\beta_j<\infty$; the finitely many early factors cause no issue,
including any zeros. Hence the limit is continuous and positive
definite. The embedding $U\mapsto U\oplus1$ leaves every factor
unchanged. For a permutation $k$-cycle the eigenvalues are the $k$th
roots of unity, whose product gives $(1-b)^k-(-b)^k$.

For (3), use the finite products of the bidiagonal factors in the
preceding proof. A minor with increasing row indices $r_i$ and column
indices $c_i$ is forced zero unless $r_i\le c_i$ for all $i$.
When these inequalities hold, realize a positive term in the iterated
Cauchy--Binet formula by moving the rightmost selected index from $r_d$
to $c_d$, one unit at a time, then the next rightmost to its target,
and so on. At every step the selected indices are distinct and the
corresponding bidiagonal minor is either one or a positive $\beta_j$.
After finitely many steps all targets are reached. Further factors
have a diagonal term equal to one, so every longer product retains a
fixed strictly positive lower bound for this minor. Passing to the
limit proves strictness. The $D_i$ therefore are positive.
Rietsch, \emph{Totally Positive Toeplitz Matrices: Classical and Modern},
ICM2026, DOI 10.1137/25M1806727, Theorem 2.1 and equations (2.1)--(2.3),
identifies the positive Toeplitz space with the positive diagonal torus
using $d_i=D_i/D_{i-1}$ and $d_{n+1}=1/D_n$. Its successive-ratio
coordinates are exactly the displayed $q_i$. Section 2.1 identifies
these with the quantum parameters for $\mathrm{Fl}_{n+1}$.
\end{proof}

For a fully explicit finite representation measure, let
$\rho_n(\sigma)=\operatorname{sgn}(\sigma)U_\sigma$ and let
$P_\lambda$ be its central orthogonal projection onto the $S_n$
isotypic component of type $\lambda$. Then
\[
 p_n(\lambda)=\operatorname{Tr}(D^{\otimes n}P_\lambda)\ge0,
 \qquad\sum_{\lambda\vdash n}p_n(\lambda)=1,
\]
and $\chi_S|_{S_n}$ is the mixture of normalized irreducible
characters with these weights. Orthogonal-projection positivity and
$\sum P_\lambda=I$ prove both assertions. The standard Frobenius formula
identifies these weights as $f^\lambda s_{\lambda'}(\beta)$; the
factor $f^\lambda$ is essential when using normalized irreducible
characters.

This gives actual algebraic-group flag varieties and representation
objects with explicit real parameters. The parameters are infinite
analytic zero data; they are not nonnegative integer structure constants,
not an asserted canonical basis for GCT, and not Frobenius eigenvalues
of a finite Weil object. The finite quantum-cohomology dictionary is a
specified imported parametrization; it does not create an arithmetic
model of $\operatorname{Spec}\mathbb Z$ or prove a GCT positivity
hypothesis. In particular it is compatible with the supplied GCT
counterexamples.


\section{An explicit asymptotic for the actual symmetric velocity trace}

\begin{theorem}[NN-A013: every fixed polynomial residue moment]
Under A there is a sequence $T_n\in(n,n+1]$ avoiding all zero ordinates
such that, for every fixed integer $r\ge0$ and every $\epsilon>0$,
\[
 \mathcal V_r(T_n):=
 \sum_{0<\Im\rho<T_n}
 \operatorname{Res}_{s=\rho}
 \frac{s^r\zeta(s+1)}{\zeta(s)}
 =\frac{i^rT_n^{r+1}}{2\pi(r+1)}
 +O_{r,\epsilon}(T_n^{r+\epsilon}).
\]
The sum is over distinct nontrivial zeros, with the full residue at a
multiple zero. No simplicity assumption is needed. If all these zeros
are simple, then
\[
 \mathcal V_r(T_n)=\sum_{0<\gamma<T_n}\rho^{r-1}v_\rho,
 \qquad v_\rho=\frac{\rho\zeta(\rho+1)}{\zeta'(\rho)}.
\]
In particular,
\[
 \mathcal V_0(T_n)=\frac{T_n}{2\pi}+O_\epsilon(T_n^\epsilon),
 \qquad
 \mathcal V_1(T_n)=\frac{iT_n^2}{4\pi}
 +O_\epsilon(T_n^{1+\epsilon}).
\]
The latter is the derivative of the sum of zeros in the fixed contour
under Hurwitz deformation, including ramified clusters as symmetric
traces. It is not an asymptotic for individual velocities at multiple
zeros.
\end{theorem}
\begin{proof}
We explicitly import the following standard consequence of RH: there
exist $T_n\in(n,n+1]$ and an absolute $A_0$ such that
\[
 \sup_{1/2\le\sigma\le2}|\zeta(\sigma+iT_n)|^{-1}
 \ll\exp(A_0\log T_n/\log\log T_n).
\]
A primary paper recording the precise statement is Milinovich--Ng,
\emph{A note on a conjecture of Gonek}, Section 3, equation (3.2),
printed p.5; its original locator is Titchmarsh, Theorem 14.16.
This imported bound assumes RH, not simplicity. The functional equation
and Stirling's formula extend it to $-1/2\le\sigma\le2$ as
\[
 |\zeta(\sigma+iT_n)|^{-1}
 \ll T_n^{\min(\sigma-1/2,0)+\epsilon}
\]
for every fixed $\epsilon>0$. A also gives the usual Lindelof bound
$\zeta(\sigma+it)\ll_\epsilon t^\epsilon$ uniformly on fixed compact
real intervals with $\sigma\ge1/2$, away from bounded $t$.

Integrate $f_r(s)=s^r\zeta(s+1)/\zeta(s)$ counterclockwise around
the rectangle with vertical sides $\Re s=2$ and $\Re s=-1/2$,
bottom at a fixed positive nonzero-ordinate height $T_0$, and top at
$T_n$. Its only poles are the nontrivial zeta zeros inside; the pole
of $\zeta(s+1)$ is at $s=0$, outside the rectangle, and $1/\zeta(s)$
has a zero rather than a pole at $s=1$. Trivial zeros are outside.
Adding the fixed finitely many omitted lower residues contributes $O(1)$.

On the right edge,
\[
 \frac{\zeta(s+1)}{\zeta(s)}=
 \sum_{m\ge1}b(m)m^{-s},\qquad
 b(m)=m^{-1}\prod_{p\mid m}(1-p),\quad b(1)=1.
\]
The bound $|b(m)|\le1$ gives absolute convergence. The $m=1$ term
contributes
\[
 \frac1{2\pi}\int_{T_0}^{T_n}(2+it)^r\,dt
 =\frac{i^rT_n^{r+1}}{2\pi(r+1)}+O_r(1+T_n^r).
\]
For $m\ge2$, integration by parts yields
$|\int_{T_0}^{T_n}(2+it)^re^{-it\log m}dt|
\le C_r(1+T_n)^r/\log m$.
Summing with $|b(m)|m^{-2}$ gives an $O_r(1+T_n^r)$ error.

On the left edge, the functional equation and absolute convergence at
$3/2-it$ give $|1/\zeta(-1/2+it)|\ll t^{-1}$.
The numerator is $\zeta(1/2+it)\ll t^\epsilon$, so
$|f_r(-1/2+it)|\ll t^{r-1+\epsilon}$.
Its integral is $O_{r,\epsilon}(T_n^{r+\epsilon})$, including $r=0$.
On the top edge the numerator lies in $\Re(s+1)\ge1/2$, hence is
$O(T_n^\epsilon)$, and the imported good-height reciprocal estimate
gives $f_r=O(T_n^{r+2\epsilon})$ uniformly. The horizontal length
is fixed. The bottom edge is a constant contribution. Replace
$\epsilon$ by $\epsilon/2$ at the start and apply the residue theorem.
At a simple zero the residue equals
$\rho^r\zeta(\rho+1)/\zeta'(\rho)=\rho^{r-1}v_\rho$.
At a multiple zero the exact interpretation is NN-A005.
\end{proof}

These formulas strengthen the elementary contour identity by supplying a
verified asymptotic at explicit good-height sequences, and repair the
missing simplicity qualification in a velocity interpretation. They do
not extend automatically to every real cutoff $T$: an individual residue
near a cutoff need not be bounded by the displayed total error. Nor does
the small real part of $\mathcal V_1$ imply equal proportions of positive
and negative real velocities. A sum can cancel with arbitrarily unequal
sign counts; $N-1$ entries equal to $1$ and one equal to $-(N-1)$ already
show this. The odd phase-correlation gate of NN-A002 remains a separate
arithmetic problem.


\section{Finite detection must escape to larger Toeplitz windows}

\begin{theorem}[NN-A014: an explicit finite-window obstruction near the Newman boundary]
Use the normalized coefficients $c_k(t)$ of NN-A007. Let
$T_t^{[R]}=(c_{j-i}(t))_{0\le i,j\le R}$, and for $t<0$ define
\[
 R_*(t)=\min\{R\ge1:T_t^{[R]}\text{ has a negative minor}\}.
\]
Under A this minimum exists for every $t<0$, and
\[
 \lim_{t\uparrow0}R_*(t)=\infty.
\]
No hypothesis on simplicity of zeta zeros is used. The following
explicit quantitative margin is available for every fixed $R$.
Let $\delta_R>0$ be the minimum at $t=0$ of all minors of
$T_0^{[R]}$ not forced to vanish by upper triangularity. Put
\[
 C_R=\max(1,c_1(0),\ldots,c_R(0)),\qquad
 A_R=\max_{0\le k\le R}(2k+2)(2k+1)c_{k+1}(0),
\]
\[
 B_R=(R+1)!(R+1)A_RC_R^R,\qquad
 \epsilon_R=\frac{\delta_R}{2B_R}.
\]
Then all non-forced minors of $T_t^{[R]}$ are at least
$\delta_R/2$ for $-\epsilon_R\le t\le0$, so
$R_*(t)>R$ on $-\epsilon_R\le t<0$.
\end{theorem}
\begin{proof}
NN-A007 gives a negative finite minor at every $t<0$, which is
contained in some initial window. At $t=0$ the positive-zero product
has infinitely many positive parameters. The nonintersecting-index
construction in NN-A008 applies without the normalization $\sum\beta=1$
and proves positivity of every non-forced minor. There are finitely
many minors in a fixed window, so $\delta_R>0$.

The theta integral gives the exact coefficient equation
\[
 c_k'(t)=(2k+2)(2k+1)c_{k+1}(t)-2c_1(t)c_k(t).
\]
Indeed numerator differentiation inserts $u^2$, and denominator
differentiation inserts its normalized mean $2c_1(t)$.
Under the probability measure proportional to
$e^{tu^2}\Phi(u)du$, this derivative is
$\operatorname{Cov}(u^{2k},u^2)/(2k)!$.
For increasing real functions $f,g$, covariance is nonnegative:
with independent $X,Y$ it equals
$\frac12\mathbb E[(f(X)-f(Y))(g(X)-g(Y))]\ge0$.
Thus $0<c_k(t)\le c_k(0)$ for $t\le0$, and
$0\le c_k'(t)\le(2k+2)(2k+1)c_{k+1}(0)$.

An $r$ by $r$ determinant is a sum of $r!$ signed products; its
derivative is a sum of at most $r!r$ products with one differentiated
entry. For $r\le R+1$, each product has absolute value at most
$A_RC_R^{r-1}$, which is bounded by the displayed $B_R$ after
multiplication by $r!r$. Entries below the diagonal and their
derivatives are zero. The fundamental theorem of calculus gives
$|D(t)-D(0)|\le B_R|t|$ for every such minor and $t\le0$.
This proves the stated margin. Since $R$ was arbitrary, it proves
divergence of the minimal detecting window.
\end{proof}

\begin{corollary}[An effective lower bound for the detecting window]
Under A, for every integer $R\ge1$ put
\[
 D_R=\left\lfloor\frac{(R+1)^2}{4}\right\rfloor,
 \qquad
 \underline\epsilon_R=
 \frac{40000^{-D_R}}
 {2^{R+1}(D_R!)^2(R+1)!(R+1)}.
\]
Every minor of $T_t^{[R]}$ not forced to vanish by upper triangularity
is strictly positive for $-\underline\epsilon_R\le t\le0$.
In particular, for all sufficiently small negative $t$,
\[
 R_*(t)\gg
 \sqrt{\frac{\log(1/|t|)}{\log\log(1/|t|)}}.
\]
The implicit constant and the meaning of ``sufficiently small'' are
effective; the displayed factorial bound is an explicit version for
every finite $R$. No simplicity assumption is required.
The later NN-A022 improves the asymptotic lower bound to
$R_*(-a)\gg a^{-1/2}$, with an explicit constant, by controlling
all minors of each bounded order at arbitrary indices.
\end{corollary}
\begin{proof}
First we establish a deliberately generous effective bound for zero
ordinates. Hasanalizade--Shen--Wong, \emph{Counting zeros of the Riemann
zeta function}, arXiv:2107.06506, Corollary~1.2, equation~(1.5), gives
for $T\ge e$
\[
 \left|N(T)-\frac{T}{2\pi}\log\frac{T}{2\pi e}\right|
 \le .1038\log T+.2573\log\log T+9.3675.
\]
The lower bound on the right, after subtracting $T/100$, is positive
at $100$ and strictly increasing for $T\ge100$. For example,
$2\pi<7$, $e<3$, $\log100<5$, and $\log\log100<2$ give its value
at $100$ greater than
$100/7-.104\cdot5-.258\cdot2-9.368-1>0$; its derivative is greater
than $1/7-.104/100-.258/100-.01>0$.
Consequently $N(100j)>j$ and $\gamma_j\le100j$, with ordinates
indexed with multiplicity. Since $x_j=2\gamma_j$ in our heat
normalization, the product parameters satisfy
\[
 b_j=x_j^{-2}\ge\frac{1}{40000j^2}=\frac{\kappa}{j^2}.
\]

Consider an allowed $r$ by $r$ minor with increasing row and column
indices $a_i\le d_i$ in $\{0,\ldots,R\}$. The path construction in
NN-A008 moves the rightmost index first, one unit at a time, and
requires exactly $D=\sum_i(d_i-a_i)$ moves. Here
\[
 D\le r(R+1-r)\le D_R.
\]
Use one move in each of the first $D$ bidiagonal factors
$I+b_jS$, and use diagonal steps in every other factor. Its positive
Cauchy--Binet summand is $\prod_{j=1}^D b_j$; all other summands are
nonnegative. Passing from finite products to the coefficientwise
limit therefore gives
\[
 \delta_R\ge\frac{\kappa^{D_R}}{(D_R!)^2}.
\]
This includes $D=0$, whose chosen summand is $1$.

Uniform coefficient bounds remove the remaining dependence on
unknown finite minima. Define
\[
 C=\sum_{k\ge0}c_k(0)
   =\frac{H_0(i)}{H_0(0)}=\frac{1}{2\xi(1/2)},
 \qquad
 A_0=\sum_{k\ge0}(2k+2)(2k+1)c_{k+1}(0)
   =\frac{\xi''(0)}{4\xi(1/2)}.
\]
The positive theta integral proves convergence of both sums and
$C>1$, $A_0>0$. The accompanying Arb certificate, using
$\xi(s)=(s-1)\pi^{-s/2}\Gamma(1+s/2)\zeta(s)$, rigorously gives
\[
 1<C<2,\qquad 0<A_0<1.
\]
One may retain the exact constants instead of this certificate:
$C=1.00579\ldots$ and $A_0=0.0117395\ldots$.
Since $C_R\le C$ and $A_R\le A_0$, the preceding theorem applies
with derivative bound
$(R+1)!(R+1)2^R$. Substituting the lower bound for $\delta_R$
gives exactly $\underline\epsilon_R$.

For completeness, elementary factorial bounds give, for $R\ge1$,
\[
 \log(1/\underline\epsilon_R)
 \le 10(R+1)^2\log(R+2).
\]
Indeed use $D_R\le(R+1)^2/4$, $\log(n!)\le n\log n$,
$\log40000<11$, and $\log(R+2)>1$ in the explicit expression;
the sum of the resulting coefficients is less than $10$.
Writing $L=\log(1/|t|)$, choose $R+1$ of order
$\sqrt{L/(20\log L)}$. For sufficiently large $L$ the last display
is at most $L/2$, so $|t|<\underline\epsilon_R$ and $R_*(t)>R$.
This proves the asserted effective asymptotic lower bound.
\end{proof}

This is a precise sense in which backward heat leaves the positive
infinite object while every fixed finite truncation initially remains
positive. It does not assert that the minimum \emph{order} of a negative
minor diverges: a bounded-order minor could move to unbounded row or
column indices. The theorem controls the full window parameter and
supplies computable analytic margins, not a polynomial complexity
classification. In representation language, every fixed finite family
of these Toeplitz/Schur tests stays positive for a negative-time interval,
although the full specialization has already left the PF cone. In flag
geometry, every fixed-rank point remains in its positive Toeplitz chart
on that interval; the obstruction to a compatible positive object
occurs at ranks tending to infinity.

\begin{proposition}[A004 refinement: transversality is additional data]
The conditions of a real entire function with all zeros on
$\Re s=1/2$ and functional equation $L(1-s)=L(s)$ do not imply that
the formal shifted-series deformation has a nonzero first normal jet
at every simple zero. An exact example is
\[
 L(s)=((s-1/2)^2+1)(1+\tfrac23(s-1/2)^2),\qquad \rho=1/2+i.
\]
All four zeros are simple and on the line, but
\[
 \rho'(0)=\frac{25i}{4},\qquad
 \rho''(0)=-208+\frac{4979i}{16}.
\]
It therefore moves to the left at second order for small positive
time, despite a purely tangential nonzero first jet.
\end{proposition}
\begin{proof}
In $x=s-1/2$, the zeros are $\pm i,\pm i\sqrt{3/2}$, distinct,
and the polynomial is even with real coefficients.
The series defining $L_t$ converges locally for $|t|<1$, since
$L(s+k)=O(k^4)$ on compact $s$ sets. Substitution into NN-A004 gives
the displayed rational complex numbers; the included symbolic check
uses exact rational arithmetic. Since the first real coefficient is
$\Re(\rho''(0)/2)=-104$, the sign claim follows.
This model has no asserted Euler product. It does not refute
transversality for the actual zeta function; it isolates the extra
phase-nonvanishing input needed for that assertion.
\end{proof}


\section{Forward heat, character dynamics, and a rational flag limit}

\begin{theorem}[NN-A015: a universal normalized character limit]
Retain the exact heat normalization of NN-A007 and set, for $t\ge0$,
\[
 p_m(t)=\sum_{j\ge1}x_j(t)^{-2m},\qquad
 \beta_j(t)=\frac{x_j(t)^{-2}}{p_1(t)},\qquad
 E_t(w)=F_t(w/p_1(t)).
\]
Under A these are positive zero weights, indexed in decreasing order
with multiplicity, and $\sum_j\beta_j(t)=1$. They obey the exact
moment equations
\[
 p_m'=4m\sum_{j=1}^{m}p_jp_{m+1-j}-m(4m+2)p_{m+1}
 \qquad(m\ge1).                                      \tag{A015.1}
\]
The moments are analytic even at a possible multiple zero at $t=0$;
no separate differentiable labeling of colliding roots is needed.

As $t\to+\infty$,
\[
 p_1(t)\sim\frac{(\log t)^2}{32},\qquad
 E_t(w)\longrightarrow E_\infty(w)=\cosh\sqrt{2w}
                                                        \tag{A015.2}
\]
locally uniformly on $\mathbb C$. The square-root expression is its
entire even-power series. In addition,
\[
 \beta_j(t)\longrightarrow
 \beta_j^*=\frac{8}{\pi^2(2j-1)^2},\qquad
 \sum_j|\beta_j(t)-\beta_j^*|\longrightarrow0.           \tag{A015.3}
\]
Consequently the characters constructed in NN-A008 have the limits
\[
 \chi_{S,t}(\sigma)\longrightarrow
 \prod_{k\ge2}\left[
 (-1)^{k-1}\frac{2^k(4^k-1)\zeta(2k)}{\pi^{2k}}
 \right]^{m_k(\sigma)},                                \tag{A015.4}
\]
and, for $U\in U(d)$ with eigenvalues $\lambda_1,\ldots,\lambda_d$,
\[
 \chi_{U,t}(U)\longrightarrow
 \prod_{a=1}^d\cosh\sqrt{2(\lambda_a-1)}.                \tag{A015.5}
\]
The limit in (A015.2) and the asymptotic for $p_1=c_1$ follow from
the kernel without A. A identifies the whole path $t\ge0$ with
positive zero weights and the indicated characters.
\end{theorem}

\begin{proof}
For every fixed real $t$, Fourier decay gives $H_t(x)\to0$ as
$x\to\pm\infty$ on the real line. If its zero product were finite,
the order-less-than-one product for $F_t$ would make $H_t$ a nonzero
polynomial, contradicting this decay. Under A its real zero product
therefore has infinitely many positive parameters at each $t\ge0$.
The heat PDE, in the coordinate $w$, and normalization at $w=0$ give
\[
 \partial_tF_t=4wF_t''+2F_t'-2p_1F_t.
\]
Near $w=0$, put $q=F_t'/F_t$ and use the convergent product to write
\[
 \log F_t(w)=\sum_{m\ge1}\frac{(-1)^{m-1}}m p_m(t)w^m,
 \qquad q(w)=\sum_{m\ge0}(-1)^m p_{m+1}(t)w^m.
\]
Then $\partial_t\log F_t=4w(q'+q^2)+2q-2p_1$.
Comparison of each coefficient proves (A015.1). Every $p_m$ is a
fixed polynomial in $c_1,\ldots,c_m$ by this logarithmic expansion,
so it is analytic in $t$ even if zero branches meet. In particular,
\[
 p_1'=4p_1^2-6p_2>0.                                  \tag{A015.6}
\]
Strictness also follows directly from the nonzero variance below.

We give a direct concentration proof of (A015.2), including the
tail control needed for local uniform convergence. The kernel is
\[
 \Phi(u)=\sum_{n\ge1}
 (2\pi^2n^4e^{9u}-3\pi n^2e^{5u})e^{-\pi n^2e^{4u}}
 \qquad(u\ge0).
\]
There are constants $0<b<B<\infty$ such that, on this whole half-line,
\[
 b e^{9u-\pi e^{4u}}\le\Phi(u)
       \le B e^{9u-\pi e^{4u}}.                        \tag{A015.7}
\]
The first summand proves the lower bound, since
$2\pi^2-3\pi e^{-4u}\ge2\pi^2-3\pi>0$.
For the upper bound discard the negative term and bound the remaining
sum by $2\pi^2\sum_{n\ge1}n^4e^{-\pi(n^2-1)}$ after factoring
out $e^{9u-\pi e^{4u}}$.

Let $U_t\ge0$ have density $H_t(0)^{-1}e^{tu^2}\Phi(u)$, and put
$a_t=(\log t)/4$ for large $t$. Integrating on
$[a_t,a_t+1/t]$ in (A015.7) gives
\[
 H_t(0)\ge\exp(t a_t^2-Ct)                            \tag{A015.8}
\]
for an absolute $C$ and all sufficiently large $t$.
Fix $0<\epsilon<1$, an integer $p\ge0$, and $A_0\ge0$.
On $0\le u\le(1-\epsilon)a_t$, the extra weight
$(u/a_t)^p e^{A_0u/a_t}$ is bounded. The integral over this interval,
divided by (A015.8), is at most a constant times
\[
 a_t\exp\bigl(-[2\epsilon-\epsilon^2]t a_t^2
                         +Ct+9a_t+A_0\bigr)\longrightarrow0.
\]
On $u\ge b_t=(1+\epsilon)a_t$, for all sufficiently large $t$,
\[
 tu^2+9u+A_0u/a_t+p\log(u/a_t)\le\tfrac\pi2e^{4u}.
\]
Indeed $u^2e^{-4u}$ and $ue^{-4u}$ decrease for large $u$, and
$t b_t^2/e^{4b_t}=O((\log t)^2t^{-\epsilon})\to0$;
the remaining linear terms are smaller. The weighted tail integral
is therefore at most
\[
 B\int_{b_t}^\infty e^{-(\pi/2)e^{4u}}du
 \le \frac{B}{2\pi e^{4b_t}}
                 e^{-(\pi/2)e^{4b_t}},
\]
which still tends to zero after division by (A015.8).
These estimates prove $U_t/a_t\to1$ in every fixed $L^p$, and
uniform integrability of all fixed exponential weights in $U_t/a_t$.
Since $p_1=c_1=\mathbb E U_t^2/2$, the claimed asymptotic for $p_1$
follows. The variance identity
$p_1'=\operatorname{Var}(U_t^2)/2>0$ justifies (A015.6).

Set $V_t=U_t/\sqrt{p_1(t)}$. We have $V_t\to\sqrt2$ in every
fixed $L^p$ and the same exponential integrability. Moreover
\[
 E_t(w)=\mathbb E\cosh(\sqrt w V_t).
\]
For $|w|\le R$ the difference from $\cosh\sqrt{2w}$ is bounded by
\[
 \sqrt R\,\mathbb E\left[
 |V_t-\sqrt2|e^{\sqrt R(V_t+\sqrt2)}\right],
\]
which tends to zero by Cauchy--Schwarz, the $L^2$ convergence and the
exponential-moment bound. This proves local uniform convergence.

The positive zeros of
$H_t(y/\sqrt{p_1(t)})/H_t(0)=E_t(-y^2)$ converge, in their order,
to the simple positive zeros of $\cos(\sqrt2y)$, by Rouche's theorem
and uniform nonvanishing on the intervening compact real intervals.
Under A all zeros being counted are real. Thus
\[
 \sqrt{p_1(t)}x_j(t)\to\frac{\pi(2j-1)}{2\sqrt2},
\]
which gives pointwise convergence in (A015.3). Both the finite-time
and limiting sequences sum to one; the latter uses
$\sum_{j\ge1}(2j-1)^{-2}=\pi^2/8$. Pointwise convergence of
nonnegative unit-mass sequences then gives $\ell^1$ convergence:
first truncate at $J$, pass to the limit in that finite set, and
bound the two omitted tails by their remaining masses.

For every fixed $k$, $|a^k-b^k|\le k|a-b|$ on $[0,1]$, hence
the power sums converge. Summing the odd reciprocal powers gives
\[
 \sum_j(\beta_j^*)^k
 =\frac{8^k}{\pi^{2k}}(1-2^{-2k})\zeta(2k)
 =\frac{2^k(4^k-1)\zeta(2k)}{\pi^{2k}}.
\]
The finite-cycle product gives (A015.4). Finally
$\chi_{U,t}(U)=\prod_{a=1}^d E_t(\lambda_a-1)$, so (A015.5)
is the already proved local uniform limit. These are actual
positive-definite characters, either by the constructions in NN-A008
with $\beta^*$ or by the pointwise limit of their finite positive
Gram matrices.
\end{proof}

\begin{corollary}[Rational normalized flag coordinates and the classical limit]
Fix $n\ge1$ and form the $(n+1)$ by $(n+1)$ upper Toeplitz matrix
from the coefficients of $E_t(w)$. Let $\widetilde\Delta_i(t)$
be its top-right $i$ by $i$ minor, with boundary values
$\widetilde\Delta_0=\widetilde\Delta_{n+1}=1$, and let
\[
 \widetilde q_i(t)=
 \frac{\widetilde\Delta_{i-1}(t)\widetilde\Delta_{i+1}(t)}
 {\widetilde\Delta_i(t)^2}.
\]
Every limit $\widetilde q_i^*$ is a strictly positive rational number,
computed from the Toeplitz coefficients $2^k/(2k)!$. In particular,
\[
 \widetilde q_1^*=n(2n-1)(4n-3).
\]
For the unnormalized coefficients of $F_t$, the corresponding
Rietsch quantum flag coordinates satisfy
\[
 q_i(t)=p_1(t)^{-2}\widetilde q_i(t)
 \sim\frac{1024\widetilde q_i^*}{(\log t)^4}.
\]
Thus this specific positive path approaches the locus $q_1=\cdots=q_n=0$
of the quantum parameters, whose specialization is the classical
cohomology ring in the usual quantum-cohomology presentation.
\end{corollary}
\begin{proof}
Coefficient convergence gives minor convergence. The limiting product
has infinitely many strictly positive parameters $\beta_j^*$, so
NN-A008 makes every allowed limiting minor strictly positive. Its
coefficients are rational, proving the first assertion. For $n\ge2$,
put $a_k=2^k/(2k)!$. Then
\[
 \widetilde\Delta_1^*=a_n,\qquad
 \widetilde\Delta_2^*=a_{n-1}^2-a_{n-2}a_n,
\]
and substitution gives
\[
 \widetilde q_1^*=
 n^2(2n-1)^2-n(n-1)(2n-1)(2n-3)
 =n(2n-1)(4n-3).
\]
For $n=1$, both boundary minors are one and $a_1=1$, giving the same
formula. Each term of $\Delta_i$ has coefficient degree
$i(n+1-i)$ under the scaling $c_k\mapsto p_1^{-k}c_k$.
Taking the second difference of this degree gives exactly
$\widetilde q_i=p_1^2q_i$. The asymptotic follows from (A015.2).
The coordinate identification is the imported Rietsch Theorem~2.1,
equations~(2.1)--(2.3), and Section~2.1, already checked in NN-A008;
the $q=0$ specialization is that theorem's standard quantum flag
presentation, not a new claim of an arithmetic motive.
\end{proof}

The limiting symmetric-group cycle values are rational because
$\zeta(2k)/\pi^{2k}$ is rational. This supplies actual rational finite
flag points, while the limiting Thoma parameters themselves contain
$\pi^{-2}$. Neither fact gives an integral positive basis for a GCT
multiplicity problem. The moment evolution and the universal limit
belong to this specified analytic heat family; they are not a theorem
about arbitrary $L$-function deformations.


\section{The first heat correction and mass at high zero indices}
\label{sec:analytic-heat-correction}

\begin{theorem}[NN-A019: Gaussian saddle and first character correction]
Retain the notation of NN-A015. For sufficiently large $t$, let $b_t>1$
be the unique large solution of
\[
 4\pi e^{4b_t}=2tb_t+9,
 \qquad \lambda_t=8tb_t-2t+36.
\]
Let $U_t$ have density $e^{tu^2}\Phi(u)/H_t(0)$ on $[0,\infty)$,
and put
\[
 \mu_t=\mathbb E U_t^2=2p_1(t),\qquad
 v_t=\frac{\operatorname{Var}(U_t^2)}{\mu_t^2}.
\]
The following assertions about the kernel hold without A or M:
\begin{align}
 b_t&=\tfrac14\bigl(\log t+\log\log t-\log(8\pi)+o(1)\bigr),
 &p_1(t)&=\tfrac12b_t^2+o(1),\\
 \sqrt{\lambda_t}(U_t-b_t)&\Longrightarrow\mathcal N(0,1),
 &v_t&\sim\frac{1}{2tb_t^3},\\
 p_1'(t)&\sim\frac{b_t}{4t}\sim\frac{\log t}{16t}.
 \label{eq:heat-first-correction-scale}
\end{align}
The normal limit holds with every absolute moment and every fixed
exponential weight. Locally uniformly for $w\in\mathbb C$,
\begin{equation}\label{eq:heat-first-correction-entire}
 E_t(w)=E_*(w)+\frac{v_t}{4}w\bigl(E_*(w)-E_*'(w)\bigr)+o(v_t),
 \qquad E_*(w)=\cosh\sqrt{2w}.
\end{equation}
Equivalently, with $\epsilon_t=v_t/4$,
\[
 E_t(w)=e^{\epsilon_tw}E_*((1-\epsilon_t)w)+o(\epsilon_t).
\]

Under A, the positive normalized zero weights satisfy, for every
fixed $j$,
\begin{equation}\label{eq:heat-first-correction-weights}
 \beta_j(t)=\beta_j^*(1-v_t/4)+o(v_t),\qquad
 \beta_j^*=\frac{8}{\pi^2(2j-1)^2}.
\end{equation}
For every fixed $k\ge2$, their power sums obey
\[
 \sum_j\beta_j(t)^k
 =\left(1-\frac{k v_t}{4}\right)\sum_j(\beta_j^*)^k+o(v_t).
\]
If $\sigma$ is a fixed finite permutation moving $d(\sigma)$ points,
then its Thoma character has the sharp relative expansion
\begin{equation}\label{eq:heat-first-correction-character}
 \chi_{S,t}(\sigma)=\chi_{S,*}(\sigma)
           \left(1-\frac{d(\sigma)v_t}{4}+o(v_t)\right).
\end{equation}
Nevertheless, the vectors $(\beta(t)-\beta^*)/v_t$ do not converge
in $\ell^1$, and
\[
 \liminf_{t\to\infty}
       \frac{\|\beta(t)-\beta^*\|_1}{v_t}\ge\frac12.
\]
Thus the first-order finite-character limit cannot be obtained by
differentiating the full weight sequence in $\ell^1$.
\end{theorem}

\begin{proof}
Write
\[
 f_t(u)=tu^2+9u-\pi e^{4u},\qquad
 h(u)=\Phi(u)e^{-9u+\pi e^{4u}}.
\]
The exact kernel and its bounds in NN-A015 give constants
$0<c<C<\infty$ with $c\le h(u)\le C$ for all $u\ge0$ and
$h(u)\to2\pi^2$ as $u\to\infty$. The latter follows by factoring
the $n=1$ summand; its negative correction is
$-3\pi e^{-4u}$ and the other summands tend to zero uniformly by
the same summable upper bound used in (A015.7).

For large $t$, $f_t'(1)>0$ whereas $f_t'(u)\to-\infty$.
Since $(f_t')''(u)=-64\pi e^{4u}<0$, there is exactly one root
$b_t>1$, the local maximum. Evaluating $4\pi e^{4b}-2tb-9$ at
$b=(\log t)/4$ and $b=(\log t+2\log\log t)/4$ shows
$b_t\sim(\log t)/4$. Taking logarithms of its defining equation
then gives the more precise displayed expansion. Also
$f_t''(b_t)=-\lambda_t$ and $\lambda_t\sim8tb_t$.

Here is a full tail justification for the Gaussian saddle method.
Fix a small constant $\delta>0$. For all sufficiently large $t$,
\[
 f_t''(u)\le-c_0\lambda_t\qquad (u\ge b_t-\delta)
\]
with $c_0>0$ independent of $t$: insert
$16\pi e^{4b_t}=8tb_t+36$ and use $b_t\to\infty$.
Since $f_t'(b_t)=0$, this implies
\[
 f_t(u)-f_t(b_t)\le-\tfrac12c_0\lambda_t(u-b_t)^2
               \qquad(u\ge b_t-\delta).
\]
On $[0,b_t-\delta]$, the maximum of $f_t$ is at an endpoint.
Indeed $f_t'$ is strictly concave, negative at $0$, positive at $1$,
and zero at $b_t$; its earlier zero is a local minimum.
Moreover $f_t(b_t)\sim tb_t^2$ whereas $f_t(0)=-\pi$.
Consequently
\[
 \sup_{0\le u\le b_t-\delta}f_t(u)
                       \le f_t(b_t)-c_1\lambda_t
\]
for some $c_1>0$. The integral on this interval remains negligible
even after multiplication by any fixed power of
$|\sqrt{\lambda_t}(u-b_t)|$ and by
$\exp(A|\sqrt{\lambda_t}(u-b_t)|)$, for fixed $A$:
its logarithmic extra cost is at most
$A b_t\sqrt{\lambda_t}+O(\log(b_t\lambda_t))=o(\lambda_t)$.

For bounded $y$, Taylor's formula gives
\[
 f_t(b_t+y/\sqrt{\lambda_t})-f_t(b_t)
                  =-y^2/2+O(|y|^3/\sqrt{\lambda_t}),
\]
because $f_t'''=O(\lambda_t)$ in a fixed neighbourhood of $b_t$.
The bounded positive factor $h$ tends there to $2\pi^2$.
The preceding Gaussian majorant and left-tail estimate justify
dominated convergence after this change of variables, including
every fixed polynomial and exponential weight. Hence, writing
$Y_t=\sqrt{\lambda_t}(U_t-b_t)$, one has
$\mathbb E Y_t\to0$, $\mathbb E Y_t^2\to1$, and uniformly bounded
absolute moments of every fixed order with exponential weights.

It follows that
\[
 \mu_t=b_t^2+o(b_t/\sqrt{\lambda_t})+O(1/\lambda_t)
          =b_t^2+o(1),
 \qquad
 \operatorname{Var}(U_t^2)\sim\frac{4b_t^2}{\lambda_t}.
\]
For the variance, expand
$U_t^2=b_t^2+2b_tY_t/\sqrt{\lambda_t}+Y_t^2/\lambda_t$;
the cross term is $O(b_t/\lambda_t^{3/2})$ and the last variance
is $O(1/\lambda_t^2)$, both smaller than the leading term.
The claims about $v_t$ and $p_1$ follow. Differentiating the
positive exponential tilt gives
$p_1'(t)=\operatorname{Var}(U_t^2)/2$, proving its asymptotic.

Set $V_t=U_t^2/\mu_t$, so that $\mathbb E V_t=1$ and
$\mathbb E(V_t-1)^2=v_t$. For each compact set $|w|\le R$,
the entire function of $v$
$g_w(v)=\cosh\sqrt{2wv}$ has, on $v\ge0$, third derivative bounded
by $C_R\exp(C_R\sqrt v)$. This follows either from its power
series on $0\le v\le1$ or from its explicit derivatives on
$v\ge1$. Taylor expansion at $v=1$ therefore has an expected
remainder bounded by
\[
 C_R\mathbb E\bigl[|V_t-1|^3
                         e^{C_R\sqrt{\max(V_t,1)}}\bigr].
\]
To bound it, put $\eta_t=1/(b_t\sqrt{\lambda_t})$. The exact
formula for $U_t^2$ and $\mu_t\sim b_t^2$ show
\[
 |V_t-1|\le C\eta_t(|Y_t|+\mathbb E|Y_t|)
       +C\eta_t^2(Y_t^2+\mathbb E Y_t^2),
 \qquad
 \sqrt{V_t}\le C+C\eta_t|Y_t|.
\]
The weighted moment bounds imply that the remainder is
$O_R(\eta_t^3)=o(v_t)$, since $v_t\sim4\eta_t^2$.
The linear Taylor term has zero mean. With $a=\sqrt{2w}$,
\[
 g_w''(1)=\frac{a^2\cosh a-a\sinh a}{4}
          =\frac w2\bigl(E_*(w)-E_*'(w)\bigr).
\]
Using $E_t(w)=\mathbb E g_w(V_t)$ proves
\eqref{eq:heat-first-correction-entire}, uniformly on the compact set.

Under A, the simple limiting roots
$w_j^*=-1/\beta_j^*$ correspond to the ordered finite-time roots
$w_j(t)=-1/\beta_j(t)$, as proved in NN-A015. The local root
expansion from \eqref{eq:heat-first-correction-entire} is
\[
 w_j(t)-w_j^*
 =-\frac{v_t}{4}\,
       \frac{w_j^*(E_*(w_j^*)-E_*'(w_j^*))}{E_*'(w_j^*)}
       +o(v_t)
 =\frac{v_t}{4}w_j^*+o(v_t).
\]
This proves \eqref{eq:heat-first-correction-weights}. On a fixed
small disc avoiding the zeros of $E_*$,
\[
 \log E_t(w)=\log E_*(w)+\frac{v_t}{4}
                w\left(1-\frac{E_*'(w)}{E_*(w)}\right)+o(v_t).
\]
Cauchy's formula permits coefficient comparison. The coefficient
of $w$ remains one; for $k\ge2$, comparison with
$\log E_t=\sum_{k\ge1}(-1)^{k-1}(\sum_j\beta_j(t)^k)w^k/k$
gives the power-sum formula. Multiplying the finitely many nonzero
cycle values gives \eqref{eq:heat-first-correction-character}.

Finally, if $(\beta(t)-\beta^*)/v_t$ converged in $\ell^1$, every
coordinate of its limit would be $-\beta_j^*/4$ by the proved
fixed-coordinate formula, so its coordinate sum would be $-1/4$.
Continuity of summation on $\ell^1$ instead makes that sum zero,
because each finite-time vector has total mass one. This is a
contradiction. For any fixed $J$, the zero total mass of
$\beta(t)-\beta^*$ also gives
\[
 \|\beta(t)-\beta^*\|_1
 \ge2\left|\sum_{j\le J}(\beta_j(t)-\beta_j^*)\right|.
\]
Divide by $v_t$, take a lower limit, and then let $J\to\infty$ to
obtain the lower bound $1/2$.
\end{proof}

\paragraph{Exact representation-theoretic translation.}
For every sufficiently small $\epsilon>0$,
\[
 e^{\epsilon w}E_*((1-\epsilon)w)
 =e^{\epsilon w}\prod_j(1+(1-\epsilon)\beta_j^*w)
\]
is the generating function of the genuine Thoma parameters
$\alpha=0$, $\beta=(1-\epsilon)\beta^*$ and $\gamma=\epsilon$.
Their total mass is one. Formula
\eqref{eq:heat-first-correction-entire} proves that this path gives
the first-order values of every fixed finite character observable
of the actual heat family, although that family has $\gamma=0$
at every finite time. The missing first-order mass at each fixed
zero index is compensated at indices tending to infinity.
The proved failure of an $\ell^1$ derivative makes this distinction
precise. The corresponding genuine $U(\infty)$ character is
\[
 e^{\epsilon\operatorname{Tr}(U-I)}
 \prod_j\det\bigl((1-(1-\epsilon)\beta_j^*)I
                         +(1-\epsilon)\beta_j^*U\bigr).
\]
Its exponential factor is positive definite on $U(d)$ by the
Poisson expansion in normalized tensor powers of the defining
character, and is compatible under $U\mapsto U\oplus1$.
Consequently the entire-function expansion gives, by finite
products, the first-order unitary character values as well.
These are exact infinite-group and analytic statements, not an
integral positive basis or an arithmetic motive. The only imported
classification is the same primary Thoma theorem used in NN-A008;
the saddle and correction proofs above are complete.
For the exterior-power dictionary, the concrete Grothendieck
locator is SGA~6, Expos\'e~V, Example~2.2(b) and Lemma~2.2.1.
The latter gives the exterior-power filtration for a locally
split exact sequence; the representation-ring case is explicitly
included after the lemma. It does not turn these analytic weights
into integral vector-bundle multiplicities.


\section{Negative heat requires minors of increasing order}
\label{sec:analytic-pf-order-escape}

\begin{lemma}[Uniform angular confinement for bounded negative heat]
\label{lem:analytic-heat-angular}
Fix $M>0$ and $0<\varepsilon<1$. Uniformly for $0\le a\le M$ and
$z$ with $-\operatorname{Im}z\ge\varepsilon|z|$, as $|z|\to\infty$,
\begin{equation}\label{eq:analytic-angular-asymptotic}
 H_{-a}(z)=H_0(z)
   \exp\left(a\left(\frac{H_0'}{H_0}(z)\right)^2\right)
       \left(1+O_{M,\varepsilon}
              \left(\frac{\log^2|z|}{|z|}\right)\right).
\end{equation}
In particular, for sufficiently large $|z|$ this sector contains
no zero of any $H_{-a}$ with $a\in[0,M]$. By conjugation the same
holds in the reflected upper sector. This lemma is unconditional.
\end{lemma}

\begin{proof}
Set $r=|z|$. In the half-plane $\operatorname{Im}w<-1$,
$s=1/2+iw/2$ satisfies $\operatorname{Re}s>1$ and
\[
 H_0(w)=\frac18\xi(s)
       =\frac1{16}s(s-1)\pi^{-s/2}\Gamma(s/2)\zeta(s)
\]
is nonzero. Every factor is nonzero there; in particular there are
no gamma poles and no zeta zeros in this half-plane. It is simply
connected, so choose a holomorphic logarithm $G(w)=\log H_0(w)$.
Write $L=G'(z)$. Stirling's expansion and the absolutely convergent
Euler series in $\operatorname{Re}s>1$ give
\begin{equation}\label{eq:analytic-angular-log-bounds}
 |L|=O_\varepsilon(\log r),\qquad
 G''(w)=O_\varepsilon(r^{-1})\quad
              (|w-z|\le 2r^{1/3}).
\end{equation}
For clarity, the gamma contribution to $G'$ is
$(i/4)\psi(s/2)$ and that to $G''$ is
$-(1/16)\psi'(s/2)$; the polynomial terms are $O(1/r)$ and
$O(1/r^2)$ respectively, and the logarithmic zeta derivatives
are $O(2^{-c_\varepsilon r})$. These estimates are uniform on
the indicated disc for large $r$.

We also need a global bound along the contours used below.
For $p=2aL$ and every real $u$, all points on the straight segment
from $z$ to $z+p+u$ have imaginary part at most
$-\varepsilon r+O_M(\log r)<-\varepsilon r/2$.
Their associated $s$ therefore remain in
$\operatorname{Re}s\ge c_\varepsilon r$. Uniform Stirling estimates
in the right half-plane and the Euler series give
\[
 |G'(w)|\le C_{M,\varepsilon}\log(r+|u|)
\]
along that segment, after enlarging the constant and taking $r$
large. Consequently the holomorphic-logarithm remainder
\[
 R(u)=G(z+p+u)-G(z)-L(p+u)
\]
satisfies
\begin{align}
 |R(u)|&\le
 \frac{C_{M,\varepsilon}}r(|u|+\log r)^2,
                  &&|u|\le r^{1/3},\label{eq:analytic-local-R}\\
 |R(u)|&\le C_{M,\varepsilon}
       (|u|+\log r)\log(r+|u|),
                  &&u\in\mathbb R.\label{eq:analytic-global-R}
\end{align}
The first bound is Taylor's formula with
\eqref{eq:analytic-angular-log-bounds}; the second is integration
of the global derivative bound and $|L|=O(\log r)$.

For $a>0$, the defining Fourier integral gives the exact Gaussian
convolution
\begin{equation}\label{eq:analytic-negative-heat-convolution}
 H_{-a}(z)=\frac1{\sqrt{4\pi a}}
       \int_{\mathbb R}e^{-u^2/(4a)}H_0(z+u)\,du.
\end{equation}
Indeed, integrating $\cos((z+u)v)$ against the Gaussian inserts
$e^{-av^2}$; Fubini is justified by
$|\cos((z+u)v)|\le e^{|\operatorname{Im}z|v}$ for real $u$
and the rapid decay of the theta kernel.

Shift the contour in \eqref{eq:analytic-negative-heat-convolution}
from $\mathbb R$ to $\mathbb R+p$. The integrand is entire in
$u$. Throughout the finite vertical displacement the point $z+u$
still lies in the zero-free half-plane just specified, so the same
branch of $G$ is valid; no gamma pole or zero is crossed.
The Gaussian dominates the growth of $H_0(z+u)$ on the vertical
ends of a large rectangle. More explicitly, for fixed $z,p,a$,
Stirling's bound is at most
$\exp(C_{z,p}|u|\log(2+|u|))$, whereas the Gaussian contributes
$\exp(-(\operatorname{Re}u)^2/(4a)+O_{p,a}(1))$.
Thus those ends tend to zero. With the contour parametrized by
$u+p$, completing the square now gives the exact identity
\[
 \frac{H_{-a}(z)}{H_0(z)e^{aL^2}}
 =\frac1{\sqrt{4\pi a}}
     \int_{\mathbb R}e^{-u^2/(4a)}e^{R(u)}\,du.
\]

The local error is uniform down to $a=0$. On $|u|\le r^{1/3}$,
\eqref{eq:analytic-local-R} tends uniformly to zero and
$|e^{R(u)}-1|\le2|R(u)|$ for large $r$. Its Gaussian mean is
$O_{M,\varepsilon}((a+\log^2r)/r)$.
On $|u|>r^{1/3}$, the right side of
\eqref{eq:analytic-global-R} is at most $u^2/(8M)$ for sufficiently
large $r$, uniformly throughout this tail. Since $a\le M$,
\[
 e^{-u^2/(4a)}|e^{R(u)}|
       \le e^{-u^2/(8a)}.
\]
After normalization, the tail integral is
$O_M(e^{-c_Mr^{2/3}})$, uniformly for $0<a\le M$; the same
bound holds for the omitted Gaussian mass. This proves
\eqref{eq:analytic-angular-asymptotic}. At $a=0$ the assertion is
the identity $H_0=H_0$. Finally, both displayed leading factors
are nonzero, so the relative error proves zero-freeness for
sufficiently large $r$.
\end{proof}

\begin{lemma}[A quadratic factor and the minor-width dictionary]
\label{lem:analytic-quadratic-pfr}
Let $b>0$, $0\le\theta<\pi/(r+1)$, and
\[
 P(w)=(1+be^{i\theta}w)(1+be^{-i\theta}w).
\]
Its coefficient Toeplitz matrix has all minors of order at most
$r$ nonnegative. Products, and locally uniform limits normalized
to have constant coefficient one, retain this property.
\end{lemma}

\begin{proof}
Write $a=be^{i\theta}$, $\bar a=be^{-i\theta}$, so the coefficients
are the elementary symmetric functions in these two variables.
For a $k$ by $k$ minor with increasing row indices $i_1,\ldots,i_k$
and column indices $j_1,\ldots,j_k$, the minor is zero if some
$i_\ell>j_\ell$, by triangularity. Otherwise set
\[
 \lambda'_p=j_{k+1-p}+p-i_1-k,\qquad
 \mu'_p=i_{k+1-p}+p-i_1-k,\qquad 1\le p\le k.
\]
These are partitions with $\mu'\subseteq\lambda'$. Transposing
and reversing both index lists identifies the determinant with
the dual Jacobi--Trudi determinant $s_{\lambda/\mu}(a,\bar a)$.
In particular $\lambda_1\le k$. The bound concerns the width of
the skew diagram, not its possibly arbitrarily large degree.
The dual Jacobi--Trudi/tableau identity follows by the usual
sign-reversing cancellation of intersecting paths in the
determinant expansion; the remaining paths are exactly
semistandard tableaux, with their variable weights.

Here the tableau formula can be evaluated directly, without any
Littlewood--Richardson positivity assumption. If a column of the
skew diagram has more than two boxes, the value is zero. Otherwise
every two-box column must contain $1$ above $2$, contributing
$a\bar a=b^2$. In each row the boxes paired with a box below form
an initial segment, and those paired with a box above form a final
segment, by the partition inequalities. The unpaired boxes form
one intervening interval of length $d\le\lambda_1\le k$. Their
weakly increasing entries can be chosen independently, with
generating function
\[
 h_d(a,\bar a)
   =b^d\frac{\sin((d+1)\theta)}{\sin\theta}>0.
\]
At $\theta=0$ this is $(d+1)b^d$. The full skew Schur value is
therefore a power of $b^2$ times a product of these nonnegative
quantities. This proves the assertion when $k\le r$.

Coefficient convolution is multiplication of the upper Toeplitz
matrices. For each fixed minor the Cauchy--Binet sum is finite,
since all intermediate indices lie between its smallest row and
largest column index. Hence products preserve the inequalities.
Local uniform convergence implies coefficient convergence and
therefore convergence of each fixed finite determinant.
\end{proof}

\begin{theorem}[NN-A020: the detecting minor order escapes]
Under A, for every integer $r\ge1$ there is $\delta_r>0$ such that
the coefficients of
\[
 F_t(w)=H_t(i\sqrt w)/H_t(0)
\]
belong to $PF_r$ for $-\delta_r\le t\le0$. In fact every allowed
minor of order at most $r$, at arbitrary indices, is strictly
positive on this interval. Here \emph{allowed} means
$i_\ell\le j_\ell$ for its increasing index lists.

Let $m_*(t)$ be the minimum order of a negative finite Toeplitz
minor. Under A and Rodgers--Tao, $m_*(t)$ is finite for each $t<0$,
but
\[
                    m_*(t)\longrightarrow\infty
                         \qquad(t\uparrow0).
\]
No simplicity hypothesis on $H_0$ is needed. This concerns minor
order uniformly over all indices; it strengthens the qualitative
finite-window result of NN-A014. The later NN-A022 supplies the
rate $m_*(-a)\gg a^{-1/2}$ by a different uniform product estimate.
\end{theorem}

\begin{proof}
For a zero $z$ in the right half-plane write its angle as
$|\arg z|\in[0,\pi/2)$. There are no zeros on the imaginary axis,
because $H_t(iy)=\int_0^\infty e^{tu^2}\Phi(u)\cosh(yu)\,du>0$.
We first show
\begin{equation}\label{eq:analytic-zero-angles-shrink}
 \sup_{\substack{H_t(z)=0\\\operatorname{Re}z>0}}|\arg z|
                         \longrightarrow0\qquad(t\uparrow0).
\end{equation}
Fix an angle $\eta>0$. Lemma~\ref{lem:analytic-heat-angular}, with
$M=1$ and $\varepsilon=\sin\eta$, excludes all zeros with
$|z|>R_\eta$ and distance in angle at least $\eta$ from the real
axis, uniformly for $-1\le t\le0$. On the remaining compact disc,
$H_t\to H_0$ uniformly. Under A, its finitely many zeros are real
and nonzero. Choose small neighbourhoods of them lying in the
$\eta$ sectors about the positive or negative real axes, and a
zero-free compact complement. Uniform convergence excludes zeros
on that complement for sufficiently small $|t|$. This argument
allows arbitrary multiplicity in the neighbourhoods: it requires
their locations to be real, not their splitting to remain real.
It also excludes a neighbourhood of zero, since $H_0(0)>0$.
This proves \eqref{eq:analytic-zero-angles-shrink}.

Choose $|t|$ small enough that every right-half-plane zero has
$|\arg z|<\pi/(2(r+1))$. The genus-zero product for $F_t$, whose
order is less than one, is absolutely locally convergent and has
constant coefficient one. Its zeros are $w=-z^2$. Real zeros of
$H_t$ give factors $1+bw$ with $b>0$; each nonreal quadruple gives
a quadratic factor as in Lemma~\ref{lem:analytic-quadratic-pfr},
with parameter angle $|\theta|=2|\arg z|<\pi/(r+1)$.
The real linear factors are totally nonnegative. The lemma and
product convergence prove $PF_r$ at all indices.

There are infinitely many factors. Otherwise $F_t$, and hence
$H_t$, would be a polynomial by its order-less-than-one product,
contrary to Fourier decay $H_t(x)\to0$ along the real axis and
$H_t(0)>0$. Every factor has constant term one and strictly
positive linear coefficient. For any allowed minor of order
$k\le r$, move its row configuration to its column configuration
one unit at a time, rightmost particle first, using a different
factor for each move. Each selected single-factor minor is its
positive linear coefficient, while a stationary configuration
has minor one. The particles never collide. This is a positive
term in the finite Cauchy--Binet expansion once sufficiently many
factors are used. Additional factors preserve this lower bound
by their stationary term, and the limiting determinant is
strictly positive.

At $t=0$ the same assertion is NN-A008. Shrinking the chosen
neighbourhood if necessary gives a closed interval
$[-\delta_r,0]$. For each $t<0$, NN-A007, using A and
Rodgers--Tao, says that $F_t$ fails $PF_\infty$, so some finite
minor is negative. But for any fixed $r$ none of orders at most
$r$ is negative near zero. This is precisely the claimed
divergence of $m_*(t)$.
\end{proof}

\begin{corollary}[Effective finite certificates for each fixed order]
Assume A. There is a terminating algorithm which, given $r\ge1$,
returns a positive rational $\delta_r$ and finite analytic bounds
certifying the conclusion of NN-A020. The algorithm does not
need to decide simplicity or the multiplicities of any zero.
No polynomial-time bound is asserted.
\end{corollary}

\begin{proof}
Choose a positive rational $q$, for example $q=1/(4(r+1))$, with
$\arctan q<\pi/(2(r+1))$, and set
$\varepsilon=q/\sqrt{1+q^2}$. The proof of
Lemma~\ref{lem:analytic-heat-angular} is effective: the gamma
estimates have computable remainder bounds and the zeta bounds
come from an absolutely convergent Dirichlet series. Thus its
constants and a rational threshold $R$ for zero-freeness in
$|\operatorname{Im}z|\ge\varepsilon|z|$ can be computed for
$0\le a\le1$. There is no compactness-only constant in that lemma.
For example, the only estimate uniform over an unbounded real
tail is \eqref{eq:analytic-global-R}; after division by $u^2$ its
majorant is
\[
 C\frac{\log(R+|u|)}{|u|}
          \left(1+\frac{\log R}{|u|}\right),
\]
decreasing for $|u|\ge R^{1/3}$ when $R>e$.
The resulting finite elementary inequalities, and the local
$C\log^2R/R$ error, have computable thresholds. Increasing $R$
to a rational value preserves a threshold established by these
eventually decreasing majorants. Explicit digamma and trigamma
bounds follow either from the remainder form of Stirling or
from Euler summation of their absolutely convergent series.

Consider the compact semialgebraic set
\[
 K_{R,q}=\{z:|z|\le R,\ |\operatorname{Im}z|
                                \ge q|\operatorname{Re}z|\}.
\]
Under A, $H_0$ has no zero on this set; the only real point of the
set is zero, where $H_0(0)>0$. Its positive minimum can be
certified without listing its zeros. Cover the set by rational
grid boxes, evaluate $H_0$ at their centers by complex-ball
arithmetic, and subtract a box-radius times a certified derivative
bound on $|z|\le R+1$. Refine the grid and the evaluations until
all these lower bounds are positive. Continuity and the positive
compact minimum guarantee termination. Intersection of a rational
box with $K_{R,q}$ is a real-algebraic decision, or one may use
an effective finite overcover converging to this compact set.
The Fourier integral provides a computable derivative bound.
Denote the resulting rational lower bound by $m_R>0$.

Compute a rational upper bound $B_R$ for
\[
 \int_0^\infty u^2e^{Ru}\Phi(u)\,du.
\]
Such a bound is effective from the explicit theta sum and its
superexponential tail. Since $0\le1-e^{-au^2}\le au^2$,
\[
 \sup_{|z|\le R}|H_{-a}(z)-H_0(z)|\le aB_R.
\]
It follows that
$\delta_r=\min(1,m_R/(2B_R))$ has the required property:
there are no zeros on $K_{R,q}$ for $0\le a\le\delta_r$, and
the angular lemma excludes them on the unbounded part of the
same cones. Every zero therefore lies in
$|\operatorname{Im}z|<q|\operatorname{Re}z|$.
The proved factor criterion and positive paths certify all
orders at most $r$, at arbitrary indices, simultaneously.
\end{proof}

\paragraph{Quantitative interface.}
If
$\alpha(t)=\sup_{H_t(z)=0,\,\operatorname{Re}z>0}|\arg z|$
is bounded explicitly, the same proof certifies $PF_r$ whenever
$2(r+1)\alpha(t)<\pi$. Thus a quantitative angular-confinement
bound would give a quantitative minimum-order lower bound.
The fixed-sector estimate proved here does not control the rate
at which the finitely many bounded zeros approach the real line,
nor its interaction with $R_\eta$ as $\eta\downarrow0$.
It gives neither an exponential collision-height law nor a
lower bound for all algorithms.
The corollary gives an effective fixed-order certificate. The later
NN-A022 strengthens it to the explicit interval
$-2^{-36}/(r+1)^2\le t\le0$ and the corresponding
inverse-square-root order lower bound.

\paragraph{Classical locators and exact geometric scope.}
The sector implication is classical Schoenberg theory. A precise
published statement is Adm--Garloff--Tyaglov, \emph{Total
nonnegativity of finite Hurwitz matrices and root location of
polynomials}, JMAA~467 (2018), 148--170, Theorem~5.8,
equation~(5.10), printed p.~161, and the quadratic case
(5.11)--(5.12), p.~162; DOI~10.1016/j.jmaa.2018.06.065.
The original reference given there is Schoenberg,
\emph{A note on multiply positive sequences and the Descartes
rule of signs}, Rend.\ Circ.\ Mat.\ Palermo~4 (1955), 123--131.
The proof above supplies the needed two-variable assertion
directly; no unverified converse for products is used.
Stirling estimates may be read in DLMF~5.11.1--5.11.2 and
their differentiated form; the actual Gaussian-contour
comparison needed here is proved above.

The Toeplitz minors are the skew Schur/exterior-power tests with
diagram width bounded by their order. The theorem proves that
every fixed such width, even with unbounded polynomial degree or
matrix index, retains positivity for a negative-time interval.
The full infinite Schur cone has nevertheless been left at every
negative time. This is a concrete failure of bounded-order
positivity tests to detect the boundary; no finite-dimensional
pure motive, integral GCT basis, or complexity-class separation
is asserted.
The concrete exterior-power dictionary is SGA~6, Expos\'e~V,
Example~2.2(b) and Lemma~2.2.1, with its representation-ring
example. Those operations supply the symmetric-function
formalism; the actual analytic positivity and width bound
are proved here.


\section{Negative heat remains positive on every fixed unitary group}
\label{sec:analytic-finite-rank-characters}

Keep the normalization
\[
 p_1(t)=F_t'(0)>0,\qquad
 E_t(w)=F_t(w/p_1(t)),\qquad
 \Psi_t(z)=E_t(z-1)=\sum_{k\ge0}a_k(t)z^k.
\]
Thus $E_t(0)=E_t'(0)=1$ and $\Psi_t(1)=1$ for every real $t$.
For $U\in U(d)$ with eigenvalues $z_1,\ldots,z_d$, define
\[
             \chi_{t,d}(U)=\prod_{i=1}^d\Psi_t(z_i).
\]
Here a normalized character means a continuous central function of
positive type with value one at the identity. It need not be the
character of a finite-dimensional representation.

\begin{theorem}[NN-A021: positivity at fixed unitary rank]
Assume A. For every fixed $d\ge1$ there is $\eta_d>0$ such that
$\chi_{t,d}$ is a normalized character of $U(d)$ for
$-\eta_d\le t\le0$. Its expansion is
\begin{equation}\label{eq:analytic-fixed-unitary-expansion}
 \chi_{t,d}(U)=
 \sum_{\ell(\lambda)\le d}c_{\lambda,d}(t)s_\lambda(U),
 \qquad
 c_{\lambda,d}(t)=
       \det\bigl[a_{\lambda_i-i+j}(t)\bigr]_{i,j=1}^d>0,
\end{equation}
where $a_k=0$ for $k<0$ and partitions are padded with zero parts.
The series converges absolutely and uniformly on $U(d)$, and
\[
       \sum_{\ell(\lambda)\le d}
          c_{\lambda,d}(t)s_\lambda(1^d)=1.
\]
In particular every polynomial irreducible representation has
strictly positive probability weight, with no bound on its degree.

Under A and Rodgers--Tao, for every $t<0$ the compatible functions
$\chi_{t,d}$ fail positive type at some finite rank. If
\[
 d_*(t)=\min\{d\ge1:\chi_{t,d}\text{ is not of positive type}\},
\]
then $d_*(t)<\infty$ for $t<0$ and
\[
                         d_*(t)\longrightarrow\infty
                                  \qquad(t\uparrow0).
\]
For $t\ge0$, all ranks have positive type by NN-A008. No
simplicity assumption is needed at $t=0$.
\end{theorem}

\begin{proof}
We first establish positivity of the Taylor coefficients in the
required Toeplitz sense. The order of $E_t$ is less than one, and its
genus-zero product is
\[
 E_t(w)=\prod_j(1+\beta_j(t)w),\qquad
 \sum_j|\beta_j(t)|<\infty,\qquad
 \sum_j\beta_j(t)=1.
\]
The factors are real positive or occur in conjugate pairs when
$t$ is sufficiently close to zero, with angles tending uniformly
to zero by NN-A020. More explicitly, a zero $x$ of $H_t$ in the
right half-plane gives $\beta=1/(p_1(t)x^2)$, so that
$|\arg\beta|=2|\arg x|$ near zero.

Under A, all $\beta_j(0)$ are positive. There are infinitely many,
and their sum is one, hence $\beta_{\max}(0)<1$. Choose
$b$ with $\beta_{\max}(0)<b<1$. The entire function $E_0$ is
zero-free on $|w|\le1/b$. Local uniform convergence
$E_t\to E_0$ makes this disc zero-free for all sufficiently small
$|t|$. Consequently
\begin{equation}\label{eq:analytic-unitary-beta-bound}
                |\beta_j(t)|<b\qquad\text{for every }j.
\end{equation}
This argument does not require individually labelling any multiple
zero. Define
\[
       \widetilde\beta_j(t)=\frac{\beta_j(t)}{1-\beta_j(t)}.
\]
For $\beta=\rho e^{i\theta}$, $\rho\le b<1$, with $|\theta|<\pi/2$,
\[
 |\arg\widetilde\beta|
 \le |\theta|+
       \arctan\frac{\rho|\sin\theta|}{1-\rho\cos\theta}
 \le \frac{|\theta|}{1-b}.
\]
Thus the angles of all $\widetilde\beta_j$ tend uniformly to zero,
and their absolute values are summable.

The real factors $1-\beta_j$ are positive; each conjugate pair
has product $|1-\beta_j|^2>0$. Absolute convergence therefore gives
\begin{equation}\label{eq:analytic-shifted-PF-product}
 \Psi_t(z)=E_t(-1)\prod_j(1+\widetilde\beta_j(t)z),
                  \qquad E_t(-1)>0.
\end{equation}
For a fixed $d$, make $|t|$ small enough that every conjugate-pair
angle is less than $\pi/(d+1)$. The quadratic factor lemma in
NN-A020 proves $PF_d$ for the coefficients of the normalized
product in \eqref{eq:analytic-shifted-PF-product}. Positive
constant scaling preserves every minor sign. Moreover every
allowed minor of order at most $d$ is strictly positive: the
infinite product has infinitely many factors, all with positive
linear coefficient, and the one-step particle paths in the
Cauchy--Binet proof of NN-A020 apply unchanged. This establishes
strict positivity at arbitrary indices, not merely in a bounded
coefficient window. Shrinking the neighbourhood gives a closed
interval $[-\eta_d,0]$.

We give the character expansion explicitly. For any entire
function $\Psi(z)=\sum_{k\ge0}a_kz^k$, multiply
$\prod_i\Psi(x_i)$ by
$\Delta(x)=\det[x_i^{d-j}]_{i,j=1}^d$ and expand the rows:
\[
 \det\bigl[x_i^{d-j}\Psi(x_i)\bigr]_{i,j=1}^d.
\]
Cauchy--Binet, with the exponents chosen in strictly decreasing
order $m_i=\lambda_i+d-i$, gives the coefficient of the alternant
$\det[x_i^{\lambda_j+d-j}]$ as
$\det[a_{\lambda_i-i+j}]$. Dividing by $\Delta$ gives
\begin{equation}\label{eq:analytic-alternant-formal}
 \prod_{i=1}^d\Psi(x_i)
     =\sum_{\ell(\lambda)\le d}
        \det[a_{\lambda_i-i+j}]s_\lambda(x_1,\ldots,x_d)
\end{equation}
as a degree-wise formal identity. Each degree involves only
finitely many terms. The Schur polynomials are the irreducible
polynomial characters by the Weyl character formula.

For $\Psi=\Psi_t$, the coefficient determinant is the Toeplitz
minor with row indices $i-\lambda_i$ and column indices $j$,
after adding the same sufficiently large integer to both lists.
The row indices increase strictly, and
$i-\lambda_i\le i$, so this minor is allowed. It is therefore
strictly positive. Set every $x_i=u$, $0\le u<1$, in
\eqref{eq:analytic-alternant-formal}. Nonnegative coefficients
and degree-wise convergence give
\[
 \Psi_t(u)^d
      =\sum_{\ell(\lambda)\le d}
        c_{\lambda,d}(t)u^{|\lambda|}s_\lambda(1^d).
\]
Monotone convergence as $u\uparrow1$ proves that the positive
weights $c_{\lambda,d}(t)s_\lambda(1^d)$ sum to one.
For $U\in U(d)$,
$|s_\lambda(U)|\le s_\lambda(1^d)$, since this is the trace of
a unitary representation. The series thus converges absolutely
and uniformly. The formal identity inside the unit polydisc,
followed by radial convergence, identifies its value with
$\prod_i\Psi_t(z_i)$. A positive mixture of normalized unitary
characters is of positive type, proving the first assertion.

For completeness we prove the required all-ranks converse,
including a possible zero constant coefficient, rather than
postulating an infinite-group classification. Suppose an entire
function $\Psi$ with real Taylor coefficients and $\Psi(1)=1$
has $\prod_i\Psi(z_i)$ of positive type on every $U(d)$.
At rank one its Fourier coefficients are nonnegative; these
are exactly its Taylor coefficients $a_k$, with the negative
coefficients zero. Let $m$ be the first index with $a_m>0$ and
write $\Psi(z)=z^m B(z)$. Then $B(0)>0$, $B(1)=1$, and
every product $\prod_i B(z_i)$ is again of positive type:
it is the original function multiplied by the one-dimensional
unitary character $\det(U)^{-m}$.

The coefficient of the irreducible polynomial character
$s_\lambda$ in this central function is
\begin{equation}\label{eq:analytic-necessary-Schur-coefficient}
                      \det[b_{\lambda_i-i+j}]_{i,j=1}^d,
                       \qquad B(z)=\sum_{k\ge0}b_kz^k.
\end{equation}
One may obtain this identity directly from the Weyl integration
formula: multiply the two alternants, expand their determinants,
and use
$\int_{S^1}z^n\,dz/(2\pi iz)=\mathbf1_{n=0}$.
Since a continuous central positive-type function on a compact
group has nonnegative irreducible-character coefficients,
\eqref{eq:analytic-necessary-Schur-coefficient} is nonnegative.
This last fact also follows directly from a positive-type
quadratic form integrated against matrix coefficients, followed
by Schur orthogonality. Thus the specialization
\[
             \rho(h_k)=b_k/b_0,\qquad h_0=1,
\]
is nonnegative on every ordinary Schur function, because every
partition occurs at some finite rank.

Here ordinary Schur positivity implies skew Schur positivity.
We recall a finite-dimensional proof of the needed implication.
In two disjoint alphabets $X,Y$, splitting a semistandard
tableau according to its alphabet gives
\[
        s_\lambda(X,Y)=
             \sum_{\mu\subseteq\lambda}s_\mu(X)s_{\lambda/\mu}(Y).
\]
On the other hand, restrict the polynomial irreducible
representation with character $s_\lambda$ of $U(p+q)$ to
$U(p)\times U(q)$, where $p,q$ exceed the diagram size.
Complete reducibility gives a nonnegative integer sum
\[
        s_\lambda(X,Y)=
             \sum_{\mu,\nu}N_{\mu,\nu}^\lambda
                           s_\mu(X)s_\nu(Y),
                         \qquad N_{\mu,\nu}^\lambda\ge0.
\]
The summands are polynomial representations, so only partitions
occur. Linear independence of the $s_\mu(X)$ identifies
$s_{\lambda/\mu}(Y)=\sum_\nu N_{\mu,\nu}^\lambda s_\nu(Y)$.
This is an identity of symmetric functions by the choice of
large alphabets. Applying $\rho$ proves nonnegativity of every
skew Schur function. The ordinary Jacobi--Trudi formula
identifies every allowed Toeplitz minor of $(b_k/b_0)$ with
one of these values; forbidden minors are zero. Thus $(b_k)$
is $PF_\infty$.

Edrei's theorem now says that the entire generating function
$B$ has only nonpositive real zeros: its possible denominator
factors would have positive real poles and cannot cancel its
nonpositive numerator zeros. The factor $z^m$ adds only the
zero at zero. We have proved that an entire $\Psi$ of the stated
all-ranks positive type cannot have a nonreal zero.

Apply this conclusion to $\Psi_t$. For every $t<0$, A and
Rodgers--Tao give a nonreal zero $x$ of $H_t$. There is no
zero on the imaginary axis since $H_t(iy)>0$. Hence
$x^2$ is nonreal, and
\[
                        z=1-p_1(t)x^2
\]
is a nonreal zero of $\Psi_t$. All-ranks positive type is
therefore impossible, so $d_*(t)$ is finite. Compatibility
\[
        \chi_{t,d+1}(U\oplus1)=\chi_{t,d}(U)
\]
shows that positivity at rank $d+1$ implies positivity at
rank $d$. The first part of the theorem then proves
$d_*(t)\to\infty$ as $t\uparrow0$.
\end{proof}

\begin{corollary}[The exact finite-rank dictionary]
Let $\Psi(z)=\sum_{k\ge0}a_kz^k$ be entire with real coefficients,
$a_0>0$ and $\Psi(1)=1$. For a fixed $d\ge1$,
\[
 \bigl[\,U\longmapsto\det\Psi(U)
                 \text{ is of positive type on }U(d)\,\bigr]
 \quad\Longleftrightarrow\quad
                  (a_k)_{k\ge0}\in PF_d.
\]
Consequently, for $t<0$ sufficiently close to zero, $d_*(t)$ is
exactly the minimum order of a negative Toeplitz minor of the
coefficients of $\Psi_t$. These are the coefficients of the
translated normalized function $\Psi_t(z)=E_t(z-1)$, not the
untranslated coefficients of $F_t$.
\end{corollary}

\begin{proof}
The forward construction from $PF_d$ to a positive normalized
Schur expansion was proved above and did not require strict
inequalities. Conversely positive type gives the nonnegative
coefficient determinants
\eqref{eq:analytic-necessary-Schur-coefficient} for every
partition of length at most $d$. Normalize $a_0$ to one in the
symmetric-function specialization. In the preceding skew Schur
expansion, the coefficient $N_{\mu,\nu}^\lambda$ vanishes unless
$\mu\subseteq\lambda$, by the two-alphabet tableau identity.
Taking equally large alphabets and interchanging their two
unitary blocks proves
$N_{\mu,\nu}^\lambda=N_{\nu,\mu}^\lambda$; hence it also
vanishes unless $\nu\subseteq\lambda$.

A Toeplitz minor of order $k\le d$, with increasing indices
$i_1,\ldots,i_k$ and $j_1,\ldots,j_k$, is zero when some
$i_\ell>j_\ell$. Otherwise transpose and reverse the index
lists, and put
\[
 \lambda_p=j_{k+1-p}+p-i_1-k,\qquad
 \mu_p=i_{k+1-p}+p-i_1-k,\qquad 1\le p\le k.
\]
These are partitions with $\mu\subseteq\lambda$ and
$\ell(\lambda)\le k$. The normalized determinant is
$\rho(s_{\lambda/\mu})$ by the ordinary Jacobi--Trudi identity.
Its expansion uses only shapes $\nu\subseteq\lambda$, so every
$\ell(\nu)\le k\le d$. All their specialized Schur values are
nonnegative by the character coefficients already obtained.
Thus every minor of order at most $d$ is nonnegative.
The minimum-order assertion follows from the equivalence and
$E_t(-1)>0$ near zero.
\end{proof}

\paragraph{Effectivity at fixed rank.}
The interval $\eta_d$ can be certified effectively under A.
First certify that $E_0$ is zero-free on a rational-radius disc
$|w|\le1/b$ for some rational $b<1$. Such a choice exists by
$\beta_{\max}(0)<1$, and rational grid evaluation with a
derivative bound certifies a positive compact minimum.
Uniform effective bounds for $\partial_t E_t$ on a slightly
larger compact set, obtained from the defining theta integrals
and the positive denominator $p_1(t)$, then give a rational
time interval on which \eqref{eq:analytic-unitary-beta-bound}
holds. In the effective corollary to NN-A020 choose its integer
parameter $r$ large enough that
\[
       \frac{2\arctan(1/[4(r+1)])}{1-b}
                                      <\frac{\pi}{d+1}.
\]
That corollary certifies exactly the angular confinement used
above. The smaller of the two time intervals is a certified
$\eta_d$. This asserts a terminating finite computation at
each fixed rank; it does not assert an asymptotic rate or
polynomial running time.

\paragraph{Primary character locators and geometric translation.}
The Weyl character formula used here is, for example, Etingof
et al., \emph{Introduction to Representation Theory},
Theorem~5.22.1, printed p.~132; the neighbouring
Theorem~5.23.2 gives complete reducibility of algebraic
representations. For the precise compact-group convention,
Faraut's \emph{Infinite Dimensional Spherical Analysis},
notes written by Sho Matsumoto, Kyushu (September 2007), section~4.1,
printed p.~56, states the positive normalized-character
expansion, and section~4.2, printed p.~57, gives the Weyl
integration and character formulas. The new conclusion is
the application to bounded negative heat time at every fixed
rank, not those classical character formulas.

Exterior powers and their additive filtration formula have
the concrete Grothendieck locator SGA~6, Expos\'e~V,
Example~2.2(b) and Lemma~2.2.1, including the representation-ring
case. The coefficients in
\eqref{eq:analytic-fixed-unitary-expansion} give positive
probability measures on actual irreducible polynomial
$U(d)$-representations. They are real weights, not asserted
integer multiplicities in an algebraic vector bundle.
The analytic heat path lies in each fixed-rank positive
character cone on some negative interval, while it lies
outside the compatible positive cone for all ranks at every
negative time. This is an unbounded-degree statement at
fixed rank, stronger than finitely many character tests;
it neither supplies an integral positive basis nor a GCT
obstruction or a lower bound against all algorithms.


\section{A square-root angular law and quantitative rank escape}
\label{sec:analytic-angular-heat-rate}

The following argument uses A to retain the full real zero product
of $H_0$, with multiplicities. It strengthens NN-A020's fixed-sector
argument quantitatively. Write $x_j>0$ for the positive zeros of
$H_0$, including multiplicity, and
\[
 H_0(z)=H_0(0)\prod_j(1-z^2/x_j^2),\qquad
 N(X)=\#\{j:x_j\le X\}.
\]
The product has no residual exponential factor: as in NN-A007,
$H_0(i\sqrt w)/H_0(0)$ has order at most $1/2<1$, so its genus-zero
Hadamard factorization has constant zero-free factor, fixed at one.
In particular $\sum_jx_j^{-2}<\infty$. The classical zero-count
bound supplies an effectively computable constant $C_0$ with
\[
                       N(X)\le C_0 X\log(X+3)\qquad(X>0).
\]
For the explicit primary bound one can use
Hasanalizade--Shen--Wong, Corollary~1.2, equation~(1.5),
printed p.~2, with $x_j=2\gamma_j$, as already checked in
NN-A014. A compact initial interval is covered by an effective
bound and a certified zero-free neighbourhood of zero.

\begin{lemma}[A global canonical-product remainder]
\label{lem:analytic-product-remainder}
For a nonreal $z$, set
\[
 L(z)=\frac{H_0'}{H_0}(z),\qquad
 S(z)=\sum_j\left(\frac1{|z-x_j|^2}+\frac1{|z+x_j|^2}\right)>0.
\]
The entire function
\[
 R_z(h)=\frac{H_0(z+h)}{H_0(z)e^{L(z)h}}
\]
satisfies, for every $h\in\mathbb C$,
\begin{align}
 |R_z(h)|&\le \exp\bigl(S(z)|h|^2/2\bigr),
                  \label{eq:analytic-R-global}\\
 |R_z(h)-1|&\le
     8S(z)|h|^2\exp\bigl(2S(z)|h|^2\bigr).
                  \label{eq:analytic-R-difference}
\end{align}
\end{lemma}

\begin{proof}
The paired product and logarithmic derivative give the locally
uniformly convergent canonical product
\[
 R_z(h)=\prod_{\xi\in\{x_j,-x_j\}}
           \left(1+\frac h{z-\xi}\right)
                       \exp\left(-\frac h{z-\xi}\right).
\]
The sum of $|z-\xi|^{-2}$ is finite; pairing makes the relation
with $L(z)$ explicit. No unrecorded exponential term remains.
For every complex $v$,
\[
 |1+v|^2\le\exp(2\operatorname{Re}v+|v|^2),
 \qquad |(1+v)e^{-v}|\le e^{|v|^2/2}.
\]
Multiply this inequality over finite products and take their
locally uniform limit to obtain \eqref{eq:analytic-R-global}.

Also $R_z(0)=1$ and $R_z'(0)=0$. Put $S=S(z)$ and
$Q(v)=R_z(v/\sqrt S)$. On $|v|=1$, $|Q(v)|\le e^{1/2}$.
Cauchy's coefficient estimate and the vanishing linear
coefficient give, for $|v|\le1/2$,
\[
 |Q(v)-1|
       \le e^{1/2}\sum_{k\ge2}|v|^k
       \le2e^{1/2}|v|^2<4|v|^2.
\]
For $|v|\ge1/2$, the first bound instead gives
$|Q(v)-1|\le2e^{|v|^2/2}\le8|v|^2e^{2|v|^2}$.
The same latter bound covers $|v|\le1/2$, proving
\eqref{eq:analytic-R-difference}.
\end{proof}

\begin{lemma}[A numerical zero-free test for the Gaussian convolution]
\label{lem:analytic-gaussian-nonzero}
For $a>0$ and nonreal $z$, put
\[
                A_z=aS(z),\qquad B_z=a^2|L(z)|^2S(z).
\]
If $A_z\le2^{-10}$ and $B_z\le2^{-10}$, then $H_{-a}(z)\ne0$.
More precisely,
\begin{equation}\label{eq:analytic-relative-Gaussian-bound}
 \left|\frac{H_{-a}(z)}{H_0(z)e^{aL(z)^2}}-1\right|
 \le16e^{16B_z}
 \left\{\frac{4B_z}{\sqrt{1-16A_z}}+
             \frac{2A_z}{(1-16A_z)^{3/2}}\right\}
 \le\frac7{16}<1.
\end{equation}
\end{lemma}

\begin{proof}
The theta integral gives the exact convolution
\[
 H_{-a}(z)=\frac1{\sqrt{4\pi a}}
             \int_{\mathbb R}e^{-u^2/(4a)}H_0(z+u)\,du.
\]
Let $p=2aL(z)$ and translate the integration line to
$\mathbb R+p$. The integrand is entire. On the vertical ends,
the Gaussian quadratic decay dominates the order-one bound
$|H_0(w)|\le\exp(C|w|\log(|w|+3))$ for fixed $z,p,a$.
Thus the end integrals tend to zero. A real part of $p$ only
reparametrizes the translated horizontal line. There is no
logarithm to continue and no obstruction from zeros crossed
in this deformation: $R_z$ itself is entire.
Completing the square gives
\[
 \frac{H_{-a}(z)}{H_0(z)e^{aL(z)^2}}
          =\mathbb E R_z(p+U),\qquad U\sim N(0,2a).
\]
With $S=S(z)$, \eqref{eq:analytic-R-difference} and
$|p+U|^2\le2|p|^2+2U^2$ show that the error is at most
\[
 16S e^{4S|p|^2}
       \mathbb E\bigl[(|p|^2+U^2)e^{4SU^2}\bigr].
\]
The two elementary Gaussian integrals are
\[
 \mathbb E e^{4SU^2}=(1-16aS)^{-1/2},\qquad
 \mathbb E U^2e^{4SU^2}
                       =2a(1-16aS)^{-3/2}.
\]
Since $S|p|^2=4B_z$, these give the first inequality in
\eqref{eq:analytic-relative-Gaussian-bound}.
For $A_z,B_z\le2^{-10}$, bound $e^{16B_z}<2$,
$(1-16A_z)^{-1/2}<2$, and
$(1-16A_z)^{-3/2}<3$. The right side is then at most
$256B_z+192A_z\le448/1024=7/16$.
The denominator is nonzero under A, and this strict relative
bound proves the assertion.
\end{proof}

\begin{lemma}[Uniform cone bounds with a vanishing radial tail]
\label{lem:analytic-cone-logderivative}
There is an effectively computable $C\ge1$ such that, whenever
$0<q\le1$, $r=|z|$ and $|\operatorname{Im}z|\ge qr$,
\begin{align}
 S(z)&\le Cq^{-2}\frac{\log(r+3)}{r+1},
                 \label{eq:analytic-S-cone}\\
 |L(z)|&\le Cq^{-1}\log(r+3).
                 \label{eq:analytic-L-cone}
\end{align}
Consequently
\[
 |L(z)|^2S(z)
       \le C^3q^{-4}\frac{\log^3(r+3)}{r+1}.
\]
\end{lemma}

\begin{proof}
By evenness and conjugation we may assume
$\operatorname{Re}z\ge0$. Choose a fixed positive $r_0$ with
$2r_0<x_1$. For $r\le r_0$, the product gives bounded $S(z)$
and $L(z)=O(r)$, which are absorbed by the claimed constants.
For $r\ge r_0$ split the positive zeros into
$x_j<r/2$, $r/2\le x_j\le2r$, and $x_j>2r$.

In the first range, $|z\pm x_j|\ge r/2$. The contribution to
$S$ is at most $8N(r/2)/r^2\ll\log(r+3)/r$.
Using the paired logarithmic derivative
\[
                  L(z)=-2z\sum_j\frac1{x_j^2-z^2},
\]
the same range contributes $O(N(r/2)/r)=O(\log(r+3))$.
In the middle range, $|z-x_j|\ge qr$ and, since
$\operatorname{Re}z\ge0$, $|z+x_j|\ge r$.
Its contributions are at most
$2q^{-2}N(2r)/r^2$ to $S$ and
$2q^{-1}N(2r)/r$ to $|L|$.
Finally partial summation of the count bound gives
\[
             \sum_{x_j>2r}x_j^{-2}
                          \ll\frac{\log(r+3)}r.
\]
This last range contributes at most eight times that sum to
$S$, and $O(r)$ times that sum to $|L|$.
Combine the bounds and compare $r$ with $r+1$ for $r\ge r_0$.

All constants can be made effective from the stated zero-count
bound, the computable sum $\sum_jx_j^{-2}=F_0'(0)$, and a
certified zero-free disc at zero. To obtain the last disc one
can bound $|H_0(z)-H_0(0)|$ by $|z|$ times the computable theta
integral bound for $H_0'$ on $|z|\le1$, and use $H_0(0)>0$.
No simplicity or root-separation estimate is used.
\end{proof}

\begin{theorem}[NN-A022: angular rate and detecting-order lower bounds]
Assume A and define
\[
 \alpha(-a)=
 \sup_{\substack{H_{-a}(z)=0\\\operatorname{Re}z>0}}|\arg z|,
                          \qquad a>0.
\]
There are effectively computable constants $a_0,c,c_U,C>0$ such
that for $0<a\le a_0$,
\begin{equation}\label{eq:analytic-root-angle-rate}
                         \alpha(-a)\le C\sqrt a.
\end{equation}
More precisely the high zeros have a uniformly vanishing
scaled contribution:
\begin{equation}\label{eq:analytic-angular-tail-small}
 \lim_{R\to\infty}
 \sup_{0<a\le1}\ 
 \sup_{\substack{H_{-a}(z)=0,\ \operatorname{Re}z>0\\|z|\ge R}}
                       \frac{|\arg z|}{\sqrt a}=0,
\end{equation}
where the supremum over an empty set is zero.

For every integer $r\ge1$, all allowed coefficient Toeplitz
minors of $F_t$ of order at most $r$, at arbitrary indices,
are strictly positive throughout
\[
                         -\frac c{(r+1)^2}\le t\le0.
\]
For every $d\ge1$, $\chi_{t,d}$ from NN-A021 is of positive
type, with every polynomial Schur weight strictly positive,
throughout
\[
                         -\frac {c_U}{(d+1)^2}\le t\le0.
\]
Thus the least negative-minor order and the least failing
unitary rank satisfy
\[
                 m_*(-a)\gg a^{-1/2},
                 \qquad d_*(-a)\gg a^{-1/2}\qquad(a\downarrow0).
\]
The least detecting index window of NN-A014 also satisfies
$R_*(-a)\ge m_*(-a)-1\gg a^{-1/2}$, since a minor in a window
of size $R+1$ has order at most $R+1$.
Rodgers--Tao is needed only to know that these detecting
quantities are finite at every negative time.
\end{theorem}

\begin{proof}
The functions
$\log(r+3)/(r+1)$ and $\log^3(r+3)/(r+1)$ are bounded on
$[0,\infty)$ and tend to zero at infinity. The last lemma
therefore gives absolute effective constants $C_1,C_2$ with
\[
                  S(z)\le C_1q^{-2},\qquad
                  |L(z)|^2S(z)\le C_2q^{-4}
\]
on the stated cones. Choose an effective $K\ge1$ such that
$C_1/K^2\le2^{-10}$ and $C_2/K^4\le2^{-10}$.
For $q=K\sqrt a\le1/2$, Lemma~\ref{lem:analytic-gaussian-nonzero}
excludes every zero with $|\operatorname{Im}z|\ge q|z|$.
There are no imaginary-axis zeros, since $H_{-a}(iy)>0$.
Hence
\[
                 \alpha(-a)\le\arcsin(K\sqrt a)\le2K\sqrt a
                    \qquad(0<a\le1/(4K^2)).
\]
This proves \eqref{eq:analytic-root-angle-rate}.

For the sharper tail statement, fix $0<\varepsilon<1/2$.
In the cone with $q=\varepsilon\sqrt a$, $0<a\le1$, the
two parameters in the Gaussian nonvanishing test obey
\[
 A_z\le\frac C{\varepsilon^2}
               \frac{\log(r+3)}{r+1},\qquad
 B_z\le\frac{C^3}{\varepsilon^4}
               \frac{\log^3(r+3)}{r+1}.
\]
Both are at most $2^{-10}$ for $r\ge R_\varepsilon$,
uniformly in $a$. A zero of that size must therefore have
$|\arg z|\le\arcsin(\varepsilon\sqrt a)
\le2\varepsilon\sqrt a$. Let $\varepsilon$ decrease to zero
to obtain \eqref{eq:analytic-angular-tail-small}.

The quadratic factor criterion in NN-A020 applies whenever
$2\alpha(t)<\pi/(r+1)$. The displayed angular bound supplies
this inequality on $-c/(r+1)^2\le t\le0$ after choosing a
fixed smaller positive effective $c$. The infinite positive
particle paths give strictness for every allowed minor.
For example, with the preceding $K$, taking
$c\le\pi^2/(64K^2)$ gives a factor-two margin in the angle
inequality and lies inside the small-time range.

For the unitary statement retain the effective constants
$b<1$ and $\eta_0>0$ in the proof of NN-A021, so
$|\beta_j(t)|<b$ when $-\eta_0\le t\le0$. Then
\[
          |\arg\widetilde\beta_j(t)|
               \le\frac{2\alpha(t)}{1-b}
               \le\frac{4K\sqrt{|t|}}{1-b}.
\]
Choose $c_U>0$ with $c_U\le\eta_0$ and
$c_U\le\pi^2(1-b)^2/(64K^2)$. The same quadratic-factor
argument and the exact finite-rank dictionary of NN-A021
give the conclusion for every $d$. Taking the largest
integers $r,d$ permitted by these intervals gives the two
lower bounds for the detecting order and rank.
\end{proof}

\begin{corollary}[One explicit numerical certificate]
The constant in the $PF_r$ interval may be chosen as
$c=2^{-36}$. Thus, under A, every allowed minor of order at
most $r$ is strictly positive for
\[
                       -\frac{2^{-36}}{(r+1)^2}\le t\le0.
\]
For $0<a\le2^{-38}$ this implies
\[
 m_*(-a)\ge
       \left\lfloor\frac{2^{-18}}{\sqrt a}\right\rfloor,
 \qquad
 R_*(-a)\ge
       \left\lfloor\frac{2^{-18}}{\sqrt a}\right\rfloor-1.
\]
The numerical input is the existing certified inequality
$A_0<1$ from NN-A014; no new floating-point experiment is used.
\end{corollary}

\begin{proof}
The coefficient sum from NN-A014 obeys
\[
 A_0=\sum_{k\ge0}(2k+2)(2k+1)c_{k+1}(0)
                              \ge2c_1(0)=2p_1(0).
\]
Its saved Arb certificate gives $0<p_1(0)<1/2$. For real
$|x|\le1$, the elementary inequality $\cos v\ge1-v^2/2$
therefore implies
\[
 H_0(x)\ge H_0(0)(1-p_1(0)x^2)>0.
\]
In particular $x_1>1$, and $\sum_jx_j^{-2}=p_1(0)<1/2$.

Write $N_\zeta(T)$ for the positive-ordinate zeta count.
The exact source bound, for $T\ge e$, is
\[
 \left|N_\zeta(T)-\frac{T}{2\pi}
                      \log\frac{T}{2\pi e}\right|
       \le .1038\log T+.2573\log\log T+9.3675.
\]
This is Hasanalizade--Shen--Wong, Corollary~1.2,
equation~(1.5), printed p.~2. At $T=e$ the main term is
negative, so $N_\zeta(e)<10$. For $T\ge e$, the main term
is at most $T\log(T+3)$, and the error is less than
\[
 .362\log T+9.368
       \le .362T+4T
       \le4.362T\log(T+3).
\]
Here $\log\log T\le\log T$, $e>5/2$, and
$\log(T+3)>1$ suffice. Hence
$N_\zeta(T)\le6T\log(T+3)$. Together with $x_1>1$ and
$N(X)=N_\zeta(X/2)$, these bounds give the deliberately loose
but uniform inequality
\begin{equation}\label{eq:analytic-explicit-count-100}
                         N(X)\le100X\log(X+3)\qquad(X>0).
\end{equation}
Indeed for $X<1$ the left side is zero, for
$1\le X<2e$ it is less than ten, and for $X\ge2e$
it is at most $3X\log(X+3)$.

We now make every constant in the cone lemma explicit.
For $r\ge1/2$, the range $x_j<r/2$ contributes at most
\[
       400\frac{\log(r+3)}r\quad\hbox{to }S,
       \qquad \frac{400}3\log(r+3)\quad\hbox{to }|L|.
\]
The middle range contributes at most
\[
       800q^{-2}\frac{\log(r+3)}r\quad\hbox{to }S,
       \qquad800q^{-1}\log(r+3)\quad\hbox{to }|L|,
\]
using $\log(2r+3)\le2\log(r+3)$.
For the last range, partial summation and
$\log(2r+3)>1$ give
\begin{align*}
 \sum_{x_j>2r}x_j^{-2}
 &\le200\int_{2r}^\infty\frac{\log(u+3)}{u^2}\,du\\
 &\le200\,\frac{\log(2r+3)+1}{2r}
 \le400\,\frac{\log(r+3)}r.
\end{align*}
It contributes at most
$3200\log(r+3)/r$ to $S$ and
$(3200/3)\log(r+3)$ to $|L|$.
Thus
\[
 S\le4400q^{-2}\frac{\log(r+3)}r
       \le13200q^{-2}\frac{\log(r+3)}{r+1},
       \qquad |L|\le2000q^{-1}\log(r+3).
\]
For $r\le1/2$, $x_1>1$ and $p_1(0)<1/2$ give
$S<4$ and $|L|\le4r/3\le2/3$, which satisfy the same
bounds after taking $C=2^{14}$ in the cone lemma.

For every $r\ge0$,
\[
 \frac{\log(r+3)}{r+1}\le2,\qquad
 \frac{\log^3(r+3)}{r+1}\le16.
\]
For the second inequality put $y=r+3\ge3$, note
$r+1=y-2\ge y/3$, and maximize $(\log y)^3/y$:
its maximum is $27/e^3$, so the ratio is at most
$81/e^3<81/8<16$. The first follows from
$\log(r+3)\le r+2$.
Consequently one can use $C_1=2^{15}$ and $C_2=2^{46}$
in the proof of the theorem. With $K=2^{16}$,
\[
             C_1/K^2=2^{-17},\qquad C_2/K^4=2^{-18},
\]
both smaller than the required $2^{-10}$. Thus
\[
 \alpha(-a)\le2^{17}\sqrt a\qquad(0<a\le2^{-34}).
\]
If $a\le2^{-36}/(r+1)^2$, then
$2\alpha(-a)\le1/(r+1)<\pi/(r+1)$.
The quadratic factor criterion proves the explicit interval.
Finally take
$r=\lfloor2^{-18}/\sqrt a\rfloor-1\ge1$.
All minors of order at most this $r$ are positive, so
$m_*(-a)\ge r+1$, and $R_*(-a)\ge m_*(-a)-1$.
\end{proof}

\begin{corollary}[The limiting angle detects endpoint multiplicities]
For each distinct positive zero $x$ of $H_0$, let $m_x$ be its
multiplicity, and let $h_m$ be the largest absolute value of a
root of the probabilists' Hermite polynomial $\operatorname{He}_m$;
in particular $h_1=0$. Under A,
\begin{equation}\label{eq:analytic-endpoint-angular-limit}
       \lim_{a\downarrow0}\frac{\alpha(-a)}{\sqrt a}
          =\kappa_0,\qquad
       \kappa_0=\sup_{H_0(x)=0,\ x>0}\frac{\sqrt2\,h_{m_x}}x.
\end{equation}
The supremum is finite and is attained if it is positive.
In particular all nontrivial zeta zeros are simple if and only
if $\alpha(-a)=o(\sqrt a)$. Under that additional simplicity
condition the two detection bounds improve to
\[
                     \sqrt a\,m_*(-a)\longrightarrow\infty,
              \qquad \sqrt a\,d_*(-a)\longrightarrow\infty.
\]
This does not determine the height of the first nonreal zero.
\end{corollary}

\begin{proof}
The Hermite recurrence gives its roots as the eigenvalues of
the real symmetric tridiagonal matrix with off-diagonal
entries $\sqrt1,\ldots,\sqrt{m-1}$. Gershgorin's bound yields
$h_m\le2\sqrt m$. Since
$m_x\le N(2x)\ll x\log(x+3)$, it follows that
$h_{m_x}/x\to0$ as $x\to\infty$. There are finitely many
distinct zeros in each compact set, so $\kappa_0$ is finite
and any positive supremum is a maximum.

At each fixed positive zero $x$, NN-A006 gives the complete
local splitting for negative time:
\[
             z_j(-a)=x+i\sqrt{2a}\,h_{m_x,j}+O(a).
\]
Consequently the largest angle of these local branches,
divided by $\sqrt a$, tends to $\sqrt2\,h_{m_x}/x$.
This also covers a simple zero, whose branch remains real
for small real time.

Choose a radius $R$ whose boundary contains no zero of $H_0$.
Uniform compact convergence, local root splitting and
Rouch\'e's theorem account for every zero in $|z|\le R$.
There are finitely many endpoint germs there, so the scaled
maximum of their angles tends to the maximum of their
Hermite constants. The remaining zeros contribute arbitrarily
little, uniformly for $0<a\le1$, by
\eqref{eq:analytic-angular-tail-small}. Combining the finite
maximum with this vanishing tail proves
\eqref{eq:analytic-endpoint-angular-limit}.
Since $h_m>0$ precisely when $m\ge2$, $\kappa_0=0$ is
equivalent to simplicity of every endpoint zero.

If $\alpha(-a)/\sqrt a\to0$, then for any prescribed $M>0$
the quadratic angle criterion applies for all orders
$r\le M/\sqrt a$ once $a$ is sufficiently small. The
M\"obius angle bound from NN-A021 gives the same assertion
at unitary ranks $d\le M/\sqrt a$. Since $M$ is arbitrary,
the two strengthened lower bounds follow.
\end{proof}

\paragraph{Exact geometric and representation meaning.}
The multiplicity $m_x$ is the length of the completed local
zero algebra. The parameter change $t=-u^2/2$ splits its
initial heat polynomial into the distinct Hermite directions;
the positive real direction $x$ converts their imaginary
normal displacement into the coefficient
$\sqrt2\,h_{m_x}/x$. The global theorem shows that these
finite local Artin lengths determine the limiting maximal
angle, because the high-radius contribution vanishes uniformly.
The lower bounds concern every skew Schur test of bounded
width and every polynomial irreducible at bounded unitary
rank, even when its degree is unrestricted. Their
exterior-power dictionary is SGA~6, Expos\'e~V,
Example~2.2(b) and Lemma~2.2.1, with the exact finite-rank
equivalence proved in NN-A021. The bounds are for these
specified certificate families; no general complexity lower
bound, pure arithmetic motive, or full collision-height law
is asserted.


% Unintegrated draft: no completed peer review is claimed.
\section{A positive fraction survives backward heat on the microscopic time scale}
\label{sec:A050-backward-survival}

This is a real-zero survival statement, distinct from a first
collision or complete collision-exhaustion law. It controls a
positive fraction of actual arithmetic branches uniformly on
a common negative-time interval. No simplicity premise is imposed
on the complementary divisor.

\begin{lemma}[A local backward-heat survival disc]
\label{lem:A050-local-survival}
Let \(f\) be real entire of order at most one, with only real
zeros, and let \(x\) be a simple zero. Write
\[
 \ell_x=\frac{f''(x)}{2f'(x)}\in\mathbb R,\qquad
 S_x=\sum_{\substack{f(y)=0\\y\ne x}}\frac1{(x-y)^2},
\]
with multiplicity in the sum. Suppose \(V>0\) and \(S_x\le V\).
For \(f_{-a}=e^{a\partial_z^2}f\), the zero at \(x\) has a
simple holomorphic continuation throughout
\[
                         |a|\le\frac{2^{-16}}V.
\]
It is real for real \(0\le a\le2^{-16}/V\), and
\begin{equation}\label{eq:A050-local-location}
 z_x(-a)=x-2a\ell_x+w_x(a),\qquad
                    |w_x(a)|\le\frac1{32\sqrt V}.
\end{equation}
The heat convolution is understood in its usual entire
continuation, as in NN-A023.
\end{lemma}

\begin{proof}
The centered genus-one Hadamard product is
\[
 f(x+h)=f'(x)h\,e^{\ell_xh}R_x(h),\qquad
 R_x(h)=\prod_{y\ne x}
       \left(1+\frac h{x-y}\right)e^{-h/(x-y)}.
\]
There is no exponential term of degree two or higher.
The constant and linear terms are fixed by the two displayed
derivatives. The convergence follows from \(S_x<\infty\).
Exactly the elementary factor estimate and Cauchy argument
proved in Lemma~\ref{lem:analytic-product-remainder} give
\[
 |R_x(h)|\le e^{V|h|^2/2},\qquad
 |R_x(h)-1|\le8V|h|^2e^{2V|h|^2}.
\]
This proof also covers an empty product.

Put \(\kappa=2^{-5}\), \(c_0=2^{-16}\), and
\(r=\kappa/\sqrt V\). Completing the Gaussian square,
with the same real drift removal as NN-A023, gives
\begin{equation}\label{eq:A050-drift-identity}
 \frac{f_{-a}(x-2a\ell_x+w)}
       {f'(x)e^{\ell_xw-a\ell_x^2}}
   =G_a(w):=
       \mathbb E\bigl[(w+\sqrt{2a}\,Z)
                         R_x(w+\sqrt{2a}\,Z)\bigr],
                    \qquad Z\sim N(0,1).
\end{equation}
For positive real \(a\) this is an ordinary Gaussian
translation. For complex \(a\), the expectation is even
in the chosen square root; Gaussian domination makes it
holomorphic in \(a\), including at zero, on the indicated
disc. The original order-one growth likewise defines
the heat convolution there, so the identity continues.
The denominator is nonzero.

For \(|w|=r\), put \(v=\sqrt Vw\) and
\(\sigma=\sqrt{2aV}\). Then \(|v|=\kappa\) and
\(|\sigma|\le\sqrt{2c_0}\). Since \(\mathbb E Z=0\),
\[
 \sqrt V\,|G_a(w)-w|
 \le8\mathbb E\bigl[|v+\sigma Z|^3
                             e^{2|v+\sigma Z|^2}\bigr].
\]
The inequalities
\((u+v)^3\le4(u^3+v^3)\) for nonnegative \(u,v\), and
\(2|v+\sigma Z|^2\le4\kappa^2+8c_0Z^2\), imply that
this is at most
\[
 32e^{4\kappa^2}
 \left\{\kappa^3\mathbb E e^{8c_0Z^2}
       +(2c_0)^{3/2}\mathbb E|Z|^3e^{8c_0Z^2}\right\}.
\]
Here \(8c_0<1/4\), and the Gaussian integrals give
\[
 e^{4\kappa^2}<2,\qquad
 \mathbb E e^{Z^2/4}=\sqrt2<2,\qquad
 \mathbb E|Z|^3e^{Z^2/4}=\frac{16}{\sqrt{2\pi}}<8.
\]
The error is consequently less than
\[
 128\kappa^3+512(2c_0)^{3/2}
       <2^{-8}+2^{-13}<\kappa/4.
\]
Thus \(|G_a(w)-w|<r/4\) on \(|w|=r\), uniformly for
\(|a|\le c_0/V\). Rouche's theorem gives exactly one
zero in the disc, counted with multiplicity, so it is
simple. Its local implicit sections patch by uniqueness
over the parameter disc. For real \(a\), conjugation
preserves the disc and \(G_a\), so the unique zero is
real. Equation~\eqref{eq:A050-local-location} follows.
The strict boundary estimate also gives continuation
through each point of the closed parameter disc.
\end{proof}

\subsection*{A positive arithmetic sample with bounded inverse-square interaction}

Use the normalization already fixed in NN-A006--NN-A023:
\[
 H_0(z)=\xi(1/2+iz/2)/8,\qquad
 \partial_tH_t=-\partial_z^2H_t.
\]
Under RH its positive zeros are \(x_j=2\gamma_j\), with
multiplicity. Write
\[
 N_H(X)=\#\{j:0<x_j\le X\},\qquad
 M_T=N_H(2T)-N_H(T),\qquad N_0=N_H(T).
\]
The exact positive-height counting conventions imply
\(M_T/N_0\to1\) and \(N_H(cT)/N_0\to c\) for each fixed
\(c>0\).

The local count needed below follows directly from the
positive Hadamard sum, rather than from a global asymptotic.
The centered product of \(\xi\), as in NN-A035, gives under RH
\[
 \sum_\gamma
 \frac{3/2}{(3/2)^2+(u-\gamma)^2}
                  =\Re\frac{\xi'}{\xi}(2+iu)\ll\log(u+3).
\]
Indeed its explicit logarithmic derivative is
\[
 \frac1s+\frac1{s-1}-\frac12\log\pi
       +\frac12\frac{\Gamma'}{\Gamma}(s/2)
       +\frac{\zeta'}{\zeta}(s).
\]
At \(s=2+iu\), the Euler term is absolutely bounded,
the gamma term is \(O(\log(u+3))\), and the rational
terms are bounded. Each occurrence with \(|u-\gamma|\le1\)
contributes at least \(6/13\). Thus every bounded-length
interval near height \(T\) contains \(O(\log(T+3))\)
occurrences, including multiplicities; the same holds
for the ordinates \(x_j=2\gamma_j\).

\begin{lemma}[A bounded-interaction subfamily]\label{lem:A050-sample}
There are absolute constants \(B>0\) and \(T_0\) such that,
for \(T\ge T_0\), at least \(M_T/20\) distinct simple zeros
\(x\) in
\[
                         T+\log T<x<2T-\log T
\]
satisfy
\begin{equation}\label{eq:A050-interaction}
 S_x:=\sum_{\substack{H_0(y)=0\\y\ne x}}(x-y)^{-2}
                                      \le B(\log T)^2.
\end{equation}
The sum includes the negative zeros and every multiplicity
in the complementary set.
\end{lemma}

\begin{proof}
Use the fixed nonnegative even Schwartz test \(g\) constructed
and checked in NN-A037/NN-A045: \(g(0)=1\), Fourier support
strictly inside \((-1,1)\), and sine-pair integral less
than \(17/50\). Choose a fixed \(0<\delta<1\) such that
\(g(u)\ge99/100\) for \(|u|\le\delta\).
The exact degree-one two-point input, Rudnick--Sarnak
Theorem1.2, gives at the positive cutoff \(2T+1\),
with all distinct indices and multiplicities retained,
\[
 \sum_{j\ne k}
 g\left(F(x_j/2)-F(x_k/2)\right)
                   \le\frac7{20}N_H(2T+1),
 \qquad F(u)=\frac{u\log u}{2\pi},
\]
for all sufficiently large \(T\).
This is the same input and test as NN-A045, with the
normalization \(x_j=2\gamma_j\) written explicitly.

Put \(h=\delta/\log T\). Call an occurrence in \((T,2T)\)
isolated if it has no other occurrence at distance less
than \(h\). It is then simple. On \([T-1,2T+1]\),
\[
                 \frac{d}{dx}F(x/2)
                  =\frac{\log(x/2)+1}{4\pi}\le\log T
\]
eventually. Each nonisolated occurrence has a distinct
neighboring index contributing at least \(99/100\).
Selecting one such ordered pair per first index counts
each selected pair at most once. Thus the nonisolated
count is at most
\[
                  \frac{35}{99}N_H(2T+1).
\]
Since \(N_H(2T+1)/N_0\to2\) and \(M_T/N_0\to1\),
the isolated set \(\mathcal S_T\) has at least \(N_0/4\)
members for large \(T\). Its centers are mutually
\(h\)-separated, and each has distance at least \(h\)
from every other zero occurrence.

For a fixed occurrence \(y\) in \((T/2,3T)\), the
\(k\)-th isolated center on either side of \(y\), omitting
a coincident center, is at distance at least \(kh\).
Consequently
\[
 \sum_{\substack{x\in\mathcal S_T\\x\ne y}}
          (x-y)^{-2}
       \le2h^{-2}\sum_{k\ge1}k^{-2}<4h^{-2}.
\]
Reverse the finite nonnegative sum, preserving all
multiplicities of \(y\). Since \(N_H(3T)\le4N_0\)
eventually, the total contribution from this middle
region is at most \(16N_0h^{-2}\).

For the remaining real zeros, the count bound
\(N_H(X)\ll X\log(X+3)\) and partial summation give,
uniformly for \(x\in(T,2T)\),
\[
 \sum_{\substack{H_0(y)=0\\y\notin(T/2,3T)}}(x-y)^{-2}
                                  \ll\frac{\log T}{T}.
\]
For \(0<y\le T/2\), use distance at least \(T/2\);
for \(y\ge3T\), use distance at least \(y/3\).
For negative \(y=-u\), split \(u\le3T\) and \(u>3T\)
and use \(x+u\ge\max(T,u)\). The same count and tail
bound apply. Thus
\[
                  \sum_{x\in\mathcal S_T}S_x
                         \le20N_0h^{-2}
\]
eventually. With \(B=160\delta^{-2}\), at most \(N_0/8\)
members have \(S_x>B(\log T)^2\). At least \(N_0/8\)
remain. The endpoint strips of width \(\log T\) contain
only \(O((\log T)^2)=o(N_0)\) occurrences by the local
count, so at least \(N_0/9\) remain after deleting them.
Finally \(M_T\le2N_0\) eventually, and
\(N_0/9>M_T/20\). This proves the stated sample.
\end{proof}

\begin{theorem}[NN-A050: positive-fraction backward-heat survival]
\label{thm:A050-survival}
Assume RH. There are absolute constants \(c>0,T_0\) such
that for every \(T\ge T_0\), one can select at least
\[
                         \frac1{20}\bigl(N_H(2T)-N_H(T)\bigr)
\]
distinct simple zeros of \(H_0\) whose branches remain
distinct, real and simple, and remain in \((T,2T)\),
simultaneously for every
\[
                         -\frac c{(\log T)^2}\le t\le0.
\]
The same selected centers work for the entire interval.
In particular every \(H_t\) in that interval has at least
the displayed number of real simple zeros in the window.
One possible constant is \(c=2^{-16}/B\), with the
absolute \(B\) in Lemma~\ref{lem:A050-sample}.
\end{theorem}

\begin{proof}
For every selected center apply
Lemma~\ref{lem:A050-local-survival} with
\(V=B(\log T)^2\), using the same \(V\) for the entire
sample. This gives simple analytic branches over the
common complex parameter disc and real branches for
\(t=-a\), \(0\le a\le2^{-16}/V\).

We check their heights, including the possible drift.
The even Hadamard product gives exactly
\[
 \ell_x=\frac1{2x}+
 \sum_{\substack{y>0\\H_0(y)=0,\ y\ne x}}
                  \left(\frac1{x-y}+\frac1{x+y}\right).
\]
The two terms are paired in the infinite sum. For
\(y<x/2\) and \(y>2x\), their total absolute contribution
is \(O(\log T)\), by the same zero count and partial
summation. In the middle region, the terms \(1/(x+y)\)
also give \(O(\log T)\). For the terms with
\(0<|x-y|<1\), Cauchy and the local count give
\[
 \sum_{0<|x-y|<1}|x-y|^{-1}
        \le\bigl(O(\log T)S_x\bigr)^{1/2}
                         \ll_B(\log T)^{3/2}.
\]
For the remaining middle terms, sum the
\(O(\log T)\) occurrences in each unit-distance interval
against \(1/k\). This gives \(O((\log T)^2)\).
Consequently
\[
                         |\ell_x|\ll_B(\log T)^2
\]
uniformly on the whole sample. Formula
\eqref{eq:A050-local-location} now gives
\[
 |z_x(-a)-x|
       \le2a|\ell_x|+(32\sqrt V)^{-1}=O_B(1)
\]
throughout the common interval. The selected centers
have distance at least \(\log T\) from both endpoints,
so all continued real zeros stay in \((T,2T)\).

Finally distinct selected sections cannot meet. If they
met at any parameter, their common zero would be simple;
the local implicit function would identify the two
sections in a neighborhood of that parameter. The
identity theorem on the connected common parameter disc
would then identify their distinct initial centers, a
contradiction. This proves every assertion.
\end{proof}

\paragraph{Relation to the collision gate.}
The theorem closes a positive-fraction microscopic
survival subgate: real simplicity persists for a specified
positive fraction over a common backward-heat interval
of order \((\log T)^{-2}\), with arbitrary complementary
multiplicities. It neither identifies the earliest
collision nor gives the fraction that has collided, the
full collision point process, or a global exhaustion law.
It is compatible with the nonreal zeros guaranteed at
every negative time by the de Bruijn--Newman endpoint
result in NN-A006--NN-A007.

\paragraph{Geometric translation and exact inputs.}
Each selected reduced analytic divisor has a degree-one
unramified local section over the common complex time
disc; SGA1, ExposeI, section3, Proposition1 gives the
relative-differentials dictionary. Real conjugation
fixes that section on the real interval, and the local
discriminant stays nonzero there. The reciprocal-square
interaction retains the full nonreduced complementary
divisor, while the count measures actual reduced real
sections. This is a concrete analytic stability theorem,
not an arithmetic heat scheme or a Frobenius law.
The genus-one product and Gaussian interface were proved
in NN-A022--NN-A023. The only correlation input is the
precisely checked positive-index, multiplicity-aware
Rudnick--Sarnak Theorem1.2 already used in NN-A037/NN-A045,
with its admissible test constructed there. Their equation2.11
supplies the global degree-one zero-count asymptotic;
the local bound is proved directly above.
The primary source is
\url{https://empslocal.ex.ac.uk/people/staff/mrwatkin/zeta/rudnick-sarnak.pdf}.
No unrestricted GUE law, all-simple premise, M, numerical
experiment or historical novelty claim is used.



% Unintegrated refinement within NN-A050.
\subsection*{Preserving every fraction below the two-thirds bound}

\begin{corollary}[NN-A050: survival of every fraction below two thirds]
\label{cor:A050-two-thirds}
Assume RH. For every fixed \(0<\eta<2/3\) there are constants
\(c_\eta>0,T_\eta\) such that every \(T\ge T_\eta\) has at
least
\[
                         (2/3-\eta)M_T
\]
distinct simple centers whose actual backward-heat branches
are simultaneously real, simple, distinct and inside \((T,2T)\)
throughout \(-c_\eta/(\log T)^2\le t\le0\).
The same centers work for that whole interval.

Consequently, for any positive \(a(T)\) with
\(a(T)(\log T)^2\to0\), the lower limiting proportion
of initial centers with these properties throughout
\([-a(T),0]\) is at least \(2/3\).
\end{corollary}

\begin{proof}
Choose the nonnegative even admissible Schwartz test from
NN-A037 with
\[
             g(0)=1,\qquad I_g<1/3+\eta/16.
\]
The complete smoothing construction there allows any fixed
positive error above \(1/3\). Choose \(0<\delta_\eta<1\)
such that \(g(u)\ge1-\eta/16\) for
\(|u|\le\delta_\eta\), and put
\(h=\delta_\eta/\log T\).

The important localization uses the fixed normalization
\(F(x/2)\), which is independent of the upper cutoff.
Define
\[
 P(Y)=\sum_{\substack{0<x_j,x_k\le Y\\j\ne k}}
       g\left(F(x_j/2)-F(x_k/2)\right).
\]
The same exact pair theorem gives
\(P(Y)=I_gN_H(Y)+o(N_H(Y))\). Nonnegativity makes \(P\)
monotone. Every ordered pair witnessing nonisolation
of an occurrence in \((T,2T)\) lies below \(2T+1\)
and has at least one coordinate above \(T-1\). Hence
all selected witness pairs contribute to
\[
 P(2T+1)-P(T-1)=I_g M_T+o(M_T).
\]
The two asymptotic errors are \(o(M_T)\), since the
three zero counts are comparable. The bounded-width
endpoint errors are \(O(\log T)=o(M_T)\).

Consequently the proportion of nonisolated occurrences is
at most
\[
 \frac{I_g+o(1)}{1-\eta/16}
 \le\frac13+\frac{\eta}{8}+o(1)
                      <\frac13+\frac{\eta}{4}
\]
for large \(T\). The middle inequality follows from
\[
 \frac{1/3+\eta/16}{1-\eta/16}-\frac13
       =\frac{\eta/12}{1-\eta/16}<\frac{\eta}{8},
               \qquad 0<\eta<2/3.
\]
Thus at least \((2/3-\eta/4)M_T\) isolated centers remain.
Deleting the endpoint strips of width \(\log T\)
leaves at least \((2/3-\eta/2)M_T\).

The reversed inverse-square calculation in
Lemma~\ref{lem:A050-sample}, now using
\(N_H(3T)\le4M_T\), gives
\[
            \sum_{x\in\mathcal S_T}S_x
                       \le20M_T h^{-2}
\]
for the isolated set. Choose
\(B_\eta=80/(\eta\delta_\eta^2)\).
At most \(\eta M_T/4\) centers have
\(S_x>B_\eta(\log T)^2\). The remaining number is at
least \((2/3-3\eta/4)M_T\), hence at least the stated
\((2/3-\eta)M_T\).
The complete local, drift, height and distinctness proof
of Theorem~\ref{thm:A050-survival} applies with
\(V=B_\eta(\log T)^2\). Take
\(c_\eta=2^{-16}/B_\eta\).

For the final assertion fix \(\eta>0\). The condition
\(a(T)(\log T)^2\to0\) eventually puts the entire interval
\([-a(T),0]\) inside the interval just proved. Its surviving
proportion is therefore at least \(2/3-\eta\) eventually.
Taking the lower limit and then letting \(\eta\downarrow0\)
proves the claim. This is a refinement of the same survival
finding, not a claim about a full collision law.
\end{proof}



\section{A quantitative upper height for a negative-time nonreal zero}

\paragraph{Imported interface and normalization.}
Write
\[
 \Xi_{-a}(s)=8H_{-a}(-i(2s-1)),\qquad
 \gamma(s)=\frac{s(s-1)}2\pi^{-s/2}\Gamma(s/2),
 \qquad a>0.
\]
Thus \(\Xi_0(s)=\xi(s)=\gamma(s)\zeta(s)\), and the
parameter \(-a\) is exactly the heat parameter used throughout this
chapter. Dobner, \emph{A proof of Newman's conjecture for the extended
Selberg class}, Acta Arith.\ 201 (2021), 29--62,
Theorem~4, equations (10)--(13), gives the following special case.
There are absolute constants \(C_D,Y_D>0\) such that, uniformly for
\(0<a\le1\), \(1/2\le\Re s\le2\), and \(y=\Im s\ge Y_D\),
\begin{equation}\label{eq:A051-Dobner}
 \frac{\Xi_{-a}(J_a(s))}{\gamma_a(s)}
   =F_a(s)+O\!\left(C_Dy^{-1/5}e^{40/a}\right),
\end{equation}
where the constant in the displayed \(O\) may be absorbed into
\(C_D\), and
\begin{align*}
 F_a(s)&=\sum_{n\ge1}e^{-a(\log n)^2/4}n^{-s},\\
 J_a(s)&=s+\frac a4\operatorname{Log}\frac{s}{2\pi},
 &\gamma_a(s)&=\gamma(s)
       \exp\!\left(\frac{(s-J_a(s))^2}{a}\right).
\end{align*}
The logarithm is principal in the upper half-plane. The constants in
this interface are independent of \(a\). The primary source is
\url{https://arxiv.org/html/2005.05142v2}; its Section~4 proves the
uniformity, and Section~3.1 uses qualitative almost periodicity.
The finite recurrence bound below is proved directly.

\begin{theorem}[NN-A051: a double-exponential actual heat-zero height bound]
\label{thm:A051-height}
There are absolute constants \(C,a_0>0\) such that for every
\(0<a<a_0\), the function \(H_{-a}\) has a simple nonreal zero \(z_a\)
satisfying
\[
 0<\Re z_a\le \exp\!\left(\exp\frac{C}{\sqrt a}\right),
 \qquad \Im z_a\le-\frac12.
\]
Its complex conjugate is another simple zero with the same positive
real part. This statement is unconditional.

In particular, if
\[
 T_{\rm off}(a)=\inf\{\Re z>0:H_{-a}(z)=0,\ \Im z\ne0\},
\]
then
\[
       T_{\rm off}(a)\le
            \exp\!\left(\exp\frac{C}{\sqrt a}\right).
\]
Under RH and simplicity of every zero of \(H_0\), one additionally has
\(T_{\rm off}(a)\to\infty\) as \(a\downarrow0\).
The upper bound concerns a zero at the specified time \(-a\);
it does not identify that zero's collision time or collision center.
\end{theorem}

\begin{lemma}[A simple pole-scaling zero]
\label{lem:A051-profile}
The entire function
\[
              E(w)=\int_0^\infty e^{-v^2/4-wv}\,dv
\]
has a simple nonreal zero. If \(K\) is fixed, then uniformly for
\(|w|\le K\) and sufficiently small \(a>0\),
\begin{equation}\label{eq:A051-profile}
       F_a(1+\sqrt a\,w)=a^{-1/2}E(w)+\mathcal E_a(w),
              \qquad |\mathcal E_a(w)|\le7.
\end{equation}
The displayed numerical error bound applies whenever
\(\Re(1+\sqrt a\,w)\ge1/2\),
\(|1+\sqrt a\,w|\le2\), and \(0<a\le1\).
\end{lemma}

\begin{proof}
Uniform domination on compact sets proves that \(E\) is entire, and
completing the square gives
\[
             |E(w)|\le2\sqrt\pi e^{|w|^2}.
\]
Suppose \(E\) had no zeros. It would have an entire logarithm \(g\),
because the complex plane is simply connected. The preceding bound
gives \(\Re g(w)\le\log(2\sqrt\pi)+|w|^2\).
The elementary Borel--Caratheodory inequality proved in NN-A039,
applied on disks of radii \(R\) and \(2R\), gives
\(\sup_{|w|\le R}|g(w)-g(0)|=O(R^2+1)\).
Cauchy's derivative estimate, with \(R\to\infty\), implies that
every Taylor coefficient of \(g\) of degree greater than two vanishes.
Thus \(g\) is a polynomial of degree at most two.

On the positive real axis, substitution \(u=xv\) and dominated
convergence give
\[
 xE(x)=\int_0^\infty e^{-u-u^2/(4x^2)}\,du\longrightarrow1.
\]
Consequently \(\Re g(x)=-\log x+o(1)\), which is impossible for a
real polynomial of degree at most two. Hence \(E\) has a zero.
For every real \(x\), its defining integral is strictly positive,
so all its zeros are nonreal. Integration by parts gives
\[
                         E'(w)=2wE(w)-2.
\]
At any zero the derivative equals \(-2\), so every zero is simple.

For the second assertion put
\(s=1+\sqrt a\,w\) and
\(f_s(u)=e^{-a(\log u)^2/4}u^{-s}\), for \(u\ge1\).
The elementary sum--integral comparison on each interval
\([n,n+1]\) yields
\[
 \left|\sum_{n\ge1}f_s(n)-\int_1^\infty f_s(u)\,du\right|
       \le |f_s(1)|+\int_1^\infty|f_s'(u)|\,du.
\]
Under the displayed hypotheses the last expression is at most
\[
 1+2\int_1^\infty u^{-3/2}\,du
      +\frac a2\int_1^\infty(\log u)u^{-3/2}\,du
             \le1+4+2=7.
\]
Finally \(v=\sqrt a\log u\) gives the exact identity
\[
             \int_1^\infty f_s(u)\,du=a^{-1/2}E(w).
\]
This proves \eqref{eq:A051-profile}.
\end{proof}

\begin{lemma}[A finite quantitative almost period]
\label{lem:A051-recurrence}
Fix a compact disk \(D\) in the \(w\)-plane and a number \(\mu>0\).
There are fixed constants \(R>0\) and an integer \(Q\ge2\), depending
only on \(D,\mu\), with the following property.
For all sufficiently small \(a>0\), put
\(N=\lceil e^{R/\sqrt a}\rceil\).
For every prescribed real \(Y>0\), there is a number
\[
       Y\le\tau\le YQ^{N-1}
\]
such that, uniformly for \(w\in D\),
\begin{equation}\label{eq:A051-recurrence}
 \big|F_a(1+\sqrt a\,w+i\tau)
          -F_a(1+\sqrt a\,w)\big|
                    \le\frac{\mu}{16\sqrt a}.
\end{equation}
\end{lemma}

\begin{proof}
Choose \(K\ge1\) with \(|w|\le K\) on \(D\).
For \(K\sqrt a\le1/2\) the positive function
\[
        u\longmapsto
        e^{-a(\log u)^2/4}u^{-1+K\sqrt a}
\]
is decreasing for \(u\ge1\). Its sum over the integers satisfies
\begin{equation}\label{eq:A051-mass}
 \sum_{n\ge1}e^{-a(\log n)^2/4}n^{-1+K\sqrt a}
 \le1+\frac1{\sqrt a}\int_0^\infty e^{-v^2/4+Kv}\,dv
 \le\frac{W_K}{\sqrt a},
\end{equation}
where
\(W_K=1+\int_0^\infty e^{-v^2/4+Kv}\,dv\).
Choose \(R\) so large that
\[
             \int_R^\infty e^{-v^2/4+Kv}\,dv\le\mu/64.
\]
Since \(\sqrt a\log N\ge R\), monotonicity gives
\begin{equation}\label{eq:A051-tail}
 \sum_{n>N}e^{-a(\log n)^2/4}n^{-1+K\sqrt a}
                              \le\frac\mu{64\sqrt a}.
\end{equation}

Choose an integer \(Q\ge2\) such that
\(2\pi/Q\le\mu/(32W_K)\).
Apply the pigeonhole principle to the \(Q^{N-1}+1\) points
\[
 \left(\left\{\frac{jY\log n}{2\pi}\right\}\right)_{2\le n\le N},
                 \qquad 0\le j\le Q^{N-1},
\]
in the unit cube partitioned into \(Q^{N-1}\) half-open cubes
of side \(1/Q\). The difference of two points in the same cube
provides an integer \(1\le k\le Q^{N-1}\) such that
\[
       \left\|\frac{kY\log n}{2\pi}\right\|\le\frac1Q
                   \quad(2\le n\le N).
\]
Set \(\tau=kY\). Then
\(\left|n^{-i\tau}-1\right|\le2\pi/Q\) for \(n\le N\).
The sum of the absolute differences of the first \(N\) terms is
at most \(\mu/(32\sqrt a)\) by \eqref{eq:A051-mass}.
The two remaining tails together have the same upper bound by
\eqref{eq:A051-tail}. This proves \eqref{eq:A051-recurrence}.
No qualitative recurrence theorem or unquantified first return
time is used.
\end{proof}

\begin{proof}[Proof of Theorem~\ref{thm:A051-height}]
Choose a zero \(w_0\) of \(E\), and a sufficiently small closed
disk \(D=\{|w-w_0|\le r\}\) containing no other zero and meeting
no real point. Its boundary has a positive minimum
\[
                      \mu=\min_{\partial D}|E|>0.
\]
The zero in \(D\) is simple by Lemma~\ref{lem:A051-profile}.
Fix \(K,R,Q\) as in the two lemmas. All these choices are made
before \(a\) varies.

For small \(a>0\), set
\[
                    Y=e^{250/a},\qquad
                    N=\lceil e^{R/\sqrt a}\rceil,
\]
and choose \(\tau\in[Y,YQ^{N-1}]\) from
Lemma~\ref{lem:A051-recurrence}.
For \(w\in D\) write
\[
                          s(w)=1+\sqrt a\,w+i\tau.
\]
By decreasing a fixed \(a_0>0\), we ensure throughout the disk that
\[
 3/4\le\Re s(w)\le5/4,\qquad
 \Im s(w)\ge Y/2\ge Y_D.
\]
The ratio
\[
       B_a(w)=\frac{\Xi_{-a}(J_a(s(w)))}{\gamma_a(s(w))}
\]
is holomorphic on a neighborhood of \(D\): the logarithm there
is analytic, \(\gamma\) has no zero or pole there, and the
exponential factor in \(\gamma_a\) is nonzero.
The imported estimate \eqref{eq:A051-Dobner} has error bounded by
\[
 C_D (Y/2)^{-1/5}e^{40/a}
                   \le 2^{1/5}C_De^{-10/a}
                   \le\frac\mu{16\sqrt a}
\]
for all sufficiently small \(a\).
Combining this with both elementary lemmas gives, on \(\partial D\),
\begin{align*}
 \big|B_a(w)-a^{-1/2}E(w)\big|
 &\le \frac\mu{16\sqrt a}+\frac\mu{16\sqrt a}+7\\
 &\le\frac\mu{4\sqrt a}
       <|a^{-1/2}E(w)|,
\end{align*}
after additionally requiring \(7\sqrt a\le\mu/8\).
Rouche's theorem says that \(B_a\) has exactly one zero in \(D\),
counted with multiplicity. Denote it by \(w_a\).

Put \(s_a=s(w_a)\), \(u_a=J_a(s_a)\), and
\(z_a=-i(2u_a-1)\). This is a zero of \(H_{-a}\).
Moreover,
\[
       J_a'(s_a)=1+\frac{a}{4s_a}\ne0,
\]
and \(s'(w)=\sqrt a\ne0\). Since the denominator defining \(B_a\)
is nonzero and the zero of \(B_a\) has multiplicity one,
\(z_a\) is simple.
For sufficiently small \(a\), \(|s_a|>2\pi\), and hence
\[
 \Re u_a=\Re s_a+\frac a4\log\frac{|s_a|}{2\pi}\ge\frac34.
\]
It follows that \(\Im z_a=1-2\Re u_a\le-1/2\).
Also
\[
 \Im u_a=\Im s_a+\frac a4\operatorname{Arg}s_a>0,
 \qquad 0<\Re z_a=2\Im u_a\le3\tau
\]
for small \(a\), because \(|\Im s_a-\tau|\le K\sqrt a\)
and \(0<\operatorname{Arg}s_a<\pi\).
Thus
\begin{equation}\label{eq:A051-heightbudget}
 \log\Re z_a
    \le\log3+\frac{250}{a}+(N-1)\log Q
    \le\log3+\frac{250}{a}+2(\log Q)e^{R/\sqrt a}.
\end{equation}
There is a fixed \(C>R\) such that the last expression is at
most \(e^{C/\sqrt a}\) for all sufficiently small \(a\).
Exponentiating proves the asserted height bound.
The real kernel defining \(H_{-a}\) gives
\(H_{-a}(\bar z)=\overline{H_{-a}(z)}\), proving the conjugate-pair
assertion with both multiplicities equal to one.

Finally assume RH and simplicity of all zeros of \(H_0\).
The uniform angular estimate of NN-A022 puts every zero of
\(H_{-a}\), for small \(a\), inside a fixed cone about the real axis.
Consequently a bounded positive real part would place such zeros
in a fixed compact set. Local uniform convergence to \(H_0\),
and a disjoint Rouche disk around each simple zero in that compact
set, give a unique zero there, which is real by conjugation.
The complement of those disks is zero-free for small \(a\).
Thus no nonreal zero has bounded positive real part, proving
\(T_{\rm off}(a)\to\infty\).
\end{proof}

\paragraph{Exact gate comparison.}
This gives a one-sided, quantitative upper rate for actual
negative-time nonreal-zero height. It supplies no asymptotic equality,
no lower rate beyond the stated compact-escape implication, and no
extreme-gap distribution. In particular the double exponential is
an upper bound, not the proposed typical-gap exponential law.
Locating the collision which produces the constructed pair would
require control of its continuation to time zero; that is not
inferred from the endpoint location bound.

\paragraph{Algebraic-geometric and representation-theoretic form.}
The zero locus of the holomorphic family
\(\mathcal Z=\{(t,z):H_t(z)=0\}\) is an analytic effective Cartier
divisor. At each constructed point \((-a,z_a)\), the nonvanishing
of \(\partial_zH_{-a}(z_a)\) makes its projection to the parameter
line locally an unramified section; the conjugate point provides
a distinct conjugate section. This is the analytic local form of
the unramified criterion in SGA~1, Expose~I, Section~3,
Proposition~1,
\url{https://grothendiecksga.com/read/sga1/fr/I-3.html}.
The archimedean coordinate bound in the theorem quantifies a
non-real-supported fiber of this divisor. It does not assert that
the heat family algebraizes over a number field. Under RH, the
PF/character correspondence of NN-A007--A008 interprets the
nonreal zero as failure of full total positivity at \(-a\).
No numerical upper bound for a witnessing Toeplitz order or a
unitary rank follows merely from this height bound; NN-A026
explains why geometry of the rest of the alphabet matters.


\subsection{The same upper-height scale for each fixed primitive character}

\begin{lemma}[An explicit inhomogeneous prime-phase return]
\label{lem:A051-inhomogeneous}
Let \(N\ge2\), let \(p_1,\ldots,p_d\) be distinct primes at most
\(N\), with \(d\ge1\), and prescribe arbitrary real phases
\(b_1,\ldots,b_d\). Fix \(0<\delta\le1/4\), and put
\[
 L_0=\log(4d/\delta),\qquad
 M=\left\lceil4\delta^{-2}dL_0\right\rceil,
 \qquad \mathcal L=8N^{Md}(2M+1)^d.
\]
For every real \(Y\), some \(\tau\in[Y,Y+\mathcal L]\) satisfies
\[
       \left\|\frac{\tau\log p_j}{2\pi}-b_j\right\|<\delta
                          \quad(1\le j\le d).
\]
Here \(\|\cdot\|\) is distance to the nearest integer. In particular,
if \(d\le N\) and \(\delta^{-1}=O(\log N)\), with the implicit
constant fixed, then
\[
                       \log\mathcal L=O(N^2(\log N)^4).
\]
\end{lemma}

\begin{proof}
For \(x\in\mathbb R/\mathbb Z\) define
\(k_M(x)=\cos^{2M}(\pi x)\).
It is a nonnegative trigonometric polynomial of degree \(M\),
bounded by one. All of its Fourier coefficients are nonnegative,
their sum is one, and its constant coefficient is
\[
       c_M=4^{-M}\binom{2M}{M}\ge\frac1{2M+1}.
\]
The last inequality follows because the middle binomial coefficient
is the largest of the \(2M+1\) coefficients whose sum is \(4^M\).
If \(\|x\|\ge\delta\), concavity of sine on \([0,\pi/2]\) gives
\(\sin(\pi\|x\|)\ge2\|x\|\ge2\delta\), and therefore
\[
                    k_M(x)\le e^{-4M\delta^2}.
\]

Set
\[
 P(\tau)=\prod_{j=1}^d
       k_M\!\left(\frac{\tau\log p_j}{2\pi}-b_j\right).
\]
The sum of the absolute values of its Fourier coefficients, before
evaluation on the line, is one. Unique factorization shows that
the only frequency vector \(h\in\{-M,\ldots,M\}^d\) with
\(\sum_jh_j\log p_j=0\) is \(h=0\).
For every nonzero vector, put
\[
 A=\prod_{h_j>0}p_j^{h_j},\qquad
 B=\prod_{h_j<0}p_j^{-h_j}.
\]
The integers \(A,B\) are distinct and at most \(N^{Md}\).
If \(A>B\), then
\(\log(A/B)=\int_B^Adu/u\ge1/A\); the opposite ordering is
identical. Thus every nonzero frequency has absolute value at least
\(N^{-Md}\). Integration of each exponential over any interval
of length \(\mathcal L\) gives
\begin{equation}\label{eq:A051-phaseaverage}
 \left|\frac1{\mathcal L}\int_Y^{Y+\mathcal L}P(\tau)\,d\tau
                        -c_M^d\right|
       \le\frac{2N^{Md}}{\mathcal L}
       =\frac14(2M+1)^{-d}\le\frac14c_M^d.
\end{equation}

For completeness the displayed choice of \(M\) also ensures
\[
                    e^{-4M\delta^2}
                              <\frac14(2M+1)^{-d}.
\]
Indeed let \(B_0=4d/\delta\ge16\). Since \(L_0=\log B_0\),
\(M\le5\delta^{-2}dL_0\), so
\[
 2M+1\le11\delta^{-2}dL_0
          \le44d^2\delta^{-3}<B_0^4.
\]
Consequently \(\log(2M+1)\le4L_0\), whereas
\(4M\delta^2\ge16dL_0>4dL_0+\log4\).
If the asserted phase box were missed for every \(\tau\) in the
interval, at least one factor in \(P(\tau)\) would be bounded by
\(e^{-4M\delta^2}\), and the other factors are at most one.
This would contradict the lower average
\(3c_M^d/4\) in \eqref{eq:A051-phaseaverage}.

Finally \(\delta^{-1}=O(\log N)\) and \(d\le N\) give
\(M=O(N(\log N)^3)\). Taking logarithms in the definition of
\(\mathcal L\) proves the final bound.
\end{proof}

\begin{corollary}[NN-A051: fixed-character negative-heat height bound]
\label{cor:A051-characterheight}
Fix a primitive nonprincipal Dirichlet character \(\chi\) of conductor
\(q\). Let \(e_\chi=(1-\chi(-1))/2\) and use the completed function
\[
 \Lambda_\chi(s)=
       (q/\pi)^{(s+e_\chi)/2}\Gamma((s+e_\chi)/2)L(s,\chi).
\]
Choose one square root of its root number \(\epsilon_\chi\), and set
\[
 h_{\chi,0}(z)=\epsilon_\chi^{-1/2}
                     \Lambda_\chi((1+iz)/2).
\]
Here \(\epsilon_\chi=i^{-e_\chi}\tau(\chi)/\sqrt q\), with
\(\tau(\chi)=\sum_{r\bmod q}\chi(r)e^{2\pi ir/q}\).
The functional equation and Gauss-sum convention are
DLMF~25.15.5--25.15.6, \url{https://dlmf.nist.gov/25.15}.
Let \(h_{\chi,t}\) be its de Bruijn--Newman heat family, with
\(\partial_t h_{\chi,t}=-\partial_z^2h_{\chi,t}\).
There are constants \(C_\chi,a_\chi>0\) such that for every
\(0<a<a_\chi\), the function \(h_{\chi,-a}\) has a simple conjugate
nonreal pair, one member \(z_{\chi,a}\) satisfying
\[
 0<\Re z_{\chi,a}\le
        \exp\!\left(\exp\frac{C_\chi}{\sqrt a}\right),
          \qquad \Im z_{\chi,a}\le-\frac12.
\]
This is unconditional, also for nonreal \(\chi\). The character is
fixed: no conductor-uniform bound is asserted. The zeta case is
Theorem~\ref{thm:A051-height}, with its separate pole completion.
\end{corollary}

\begin{proof}
The usual functional equation implies
\[
 \epsilon_\chi^{-1/2}\Lambda_\chi(s)
   =\overline{\epsilon_\chi^{-1/2}
                           \Lambda_\chi(1-\bar s)}.
\]
Thus \(h_{\chi,0}\) is real on the real line. Its Fourier kernel
\(\Phi_\chi\), in Dobner's convention, satisfies
\(\Phi_\chi(-u)=\overline{\Phi_\chi(u)}\).
Dobner's Section~2 proves its super-exponential decay, so
\[
       h_{\chi,t}(z)=\int_{\mathbb R}
                       e^{tu^2}\Phi_\chi(u)e^{izu}\,du
\]
is entire jointly in \(t,z\), and remains real on the real axis
for real \(t\). Neither an even kernel nor an even function is
assumed for a nonreal character.

In the notation of the imported Theorem~4, take
\begin{align*}
 F_{\chi,a}(s)&=\sum_{n\ge1}
             e^{-a(\log n)^2/4}\chi(n)n^{-s},\\
 \gamma_\chi(s)&=\epsilon_\chi^{-1/2}
                 (q/\pi)^{(s+e_\chi)/2}\Gamma((s+e_\chi)/2),\\
 J_{\chi,a}(s)&=s+\frac a4\operatorname{Log}\frac{qs}{2\pi},
 &\gamma_{\chi,a}(s)&=\gamma_\chi(s)
              e^{(s-J_{\chi,a}(s))^2/a}.
\end{align*}
Uniformly on \(1/2\le\Re s\le2\), at sufficiently large positive
height \(y\) and for \(0<a\le1\), it gives
\begin{equation}\label{eq:A051-character-Dobner}
 \frac{h_{\chi,-a}(-i(2J_{\chi,a}(s)-1))}
                       {\gamma_{\chi,a}(s)}
      =F_{\chi,a}(s)+O_\chi(y^{-1/5}e^{40/a}).
\end{equation}
The parity shift is in \(\gamma_\chi\), while the displayed
\(J_{\chi,a}\) is exactly the one in the source formula.

We first obtain a pole profile for the principal coprime series
\[
              F_{q,a}(s)=\sum_{(n,q)=1}
                              e^{-a(\log n)^2/4}n^{-s}.
\]
Put \(r_q=\varphi(q)/q\). Periodicity gives
\(A_q(u)=\#\{n\le u:(n,q)=1\}=r_qu+R_q(u)\), with
\(|R_q(u)|\le q\).
For the function \(f_s\) from Lemma~\ref{lem:A051-profile},
partial summation gives
\[
 F_{q,a}(s)-r_q\int_1^\infty f_s(u)\,du
             =r_q f_s(1)-\int_1^\infty R_q(u)f_s'(u)\,du.
\]
Hence, for \(s=1+\sqrt a\,w\) in a fixed compact \(w\)-disk and
small \(a\),
\begin{equation}\label{eq:A051-coprimeprofile}
          F_{q,a}(1+\sqrt a\,w)
                 =\frac{r_q}{\sqrt a}E(w)+\mathcal R_{q,a}(w),
                     \qquad |\mathcal R_{q,a}(w)|\le7q.
\end{equation}
The endpoint term and derivative integral are both included.

Take the simple profile-zero disk \(D\), and let
\(\mu_q=r_q\min_{\partial D}|E|>0\).
Fix \(K\ge1\) with \(|w|\le K\) in \(D\), and the constant
\(W_K\) in \eqref{eq:A051-mass}.
Choose fixed \(R\) such that
\[
           \int_R^\infty e^{-v^2/4+Kv}\,dv\le\mu_q/64,
 \qquad N=\lceil e^{R/\sqrt a}\rceil.
\]
Both absolute tails, of \(F_{q,a}\) and \(F_{\chi,a}\), beyond
\(N\) are at most \(\mu_q/(64\sqrt a)\).

Let \(p_1,\ldots,p_d\) be the primes at most \(N\) not dividing
\(q\). For all sufficiently small \(a\) this set is nonempty.
For each such prime choose \(b_j\in\mathbb R/\mathbb Z\) with
\(e^{-2\pi i b_j}=\overline{\chi(p_j)}\).
Set
\[
 \varepsilon=\min\{1,\mu_q/(32W_K)\},\quad
 k_N=\lceil\log N/\log2\rceil,\quad
 \delta=\min\{1/4,\varepsilon/(2\pi k_N)\}.
\]
Apply Lemma~\ref{lem:A051-inhomogeneous} with
\(Y=e^{250/a}\). We obtain a \(\tau\in[Y,Y+\mathcal L]\) for
which \(|\chi(p)p^{-i\tau}-1|\le2\pi\delta\) at all these primes.
Every \(n\le N\) coprime to \(q\) has at most \(k_N\) prime
factors, counted with repetitions. Telescoping a product of unit
complex numbers therefore gives
\[
                  |\chi(n)n^{-i\tau}-1|\le\varepsilon.
\]
Terms with \((n,q)>1\) vanish in both series. By
\eqref{eq:A051-mass} and the two tail bounds, uniformly for \(w\in D\),
\begin{equation}\label{eq:A051-targetprofile}
 \big|F_{\chi,a}(1+\sqrt a\,w+i\tau)
             -F_{q,a}(1+\sqrt a\,w)\big|
                            \le\frac{\mu_q}{16\sqrt a}.
\end{equation}

Apply \eqref{eq:A051-character-Dobner} at
\(s=1+\sqrt a\,w+i\tau\).
Its error is at most \(\mu_q/(16\sqrt a)\) for all sufficiently
small \(a\), by the same \(e^{-10/a}\) estimate as in the zeta case.
Together with \eqref{eq:A051-coprimeprofile} and
\eqref{eq:A051-targetprofile}, and after requiring
\(7q\sqrt a\le\mu_q/8\), Rouche's theorem gives exactly one zero
counting multiplicity in \(D\).
The gamma factor is nonzero there and
\(J_{\chi,a}'(s)=1+a/(4s)\ne0\). The corresponding heat zero is
therefore simple.

Its real-part and imaginary-part estimates are the same as before:
\(\Re J_{\chi,a}(s)\ge3/4\), since \(q|s|>2\pi\), and its
positive real heat coordinate is at most \(3\tau\).
Here \(\delta^{-1}=O_q(\log N)\), so the finite phase lemma gives
\[
 \log\mathcal L=O_q(N^2(\log N)^4),\qquad
 \log\bigl(3(Y+\mathcal L)\bigr)
           \le\exp(C_\chi/\sqrt a)
\]
for a fixed sufficiently large \(C_\chi\) and all sufficiently
small \(a\). This proves the height bound.
The Hermitian-kernel identity proved at the start gives
\(h_{\chi,-a}(\bar z)=\overline{h_{\chi,-a}(z)}\), and hence the
second simple member of the pair. No exchange of \(\chi\) with
\(\bar\chi\), and no evenness assertion, is needed for this final
conjugation statement.
\end{proof}

\paragraph{Geometric meaning and limits.}
The root-number half-phase gives the completed analytic divisor a
specified real structure. The two constructed reduced points are
conjugate unramified local sections, as in the base theorem and
SGA~1, Expose~I, Section~3, Proposition~1. The finite phase argument
uses the free abelian group on finitely many rational primes and
its compact character torus; unique factorization provides every
nonzero-frequency separation used in the proof. It is a
deterministic recurrence argument, not a probabilistic independence
or Frobenius equidistribution assertion. For nonreal \(\chi\) the
heat kernel is Hermitian rather than necessarily real and even,
so an even-product PF or Thoma dictionary is not added by this
corollary. There is still no collision-origin or growing-conductor
conclusion.


\subsection{Positive-coefficient pole profiles and Dedekind heat families}

\begin{lemma}[A smoothed pole without a Tauberian input]
\label{lem:A051-generalpole}
Suppose \(F(s)=\sum_{n\ge1}b_nn^{-s}\) converges absolutely for
\(\Re s>1\). Suppose that
\[
                  G(s)=F(s)-\frac\kappa{s-1}
\]
is holomorphic on a strip \(1-\delta<\Re s<c+\delta\), for some
\(c>1\), \(\delta>0\), and has polynomial growth on its closed
substrips. For \(a>0\) write
\[
                  F_a(s)=\sum_{n\ge1}
                      b_ne^{-a(\log n)^2/4}n^{-s}.
\]
Assume this series is entire, as holds if \(b_n=O(n^B)\) for some
fixed \(B\). Then uniformly for \(z\) in every fixed compact disk,
\begin{equation}\label{eq:A051-generalpole}
           F_a(1+\sqrt a\,z)
                =\frac\kappa{\sqrt a}E(z)+O_{F,z\text{-disk}}(1)
                           \qquad(a\downarrow0).
\end{equation}
If also \(b_n\ge0\) and \(\kappa>0\), then for each fixed \(K>0\)
there are constants \(W_{F,K},D_{F,K}>0\) such that for all
sufficiently small \(a>0\),
\begin{align}
 \sum_{n\ge1}b_ne^{-a(\log n)^2/4}n^{-1+K\sqrt a}
                     &\le\frac{W_{F,K}}{\sqrt a},
                          \label{eq:A051-general-mass}\\
 \sum_{n>\exp(R/\sqrt a)}b_ne^{-a(\log n)^2/4}n^{-1+K\sqrt a}
                     &\le\frac{D_{F,K}}{\sqrt a}e^{-R^2/8}
                    \quad(R>0).\label{eq:A051-general-tail}
\end{align}
The smallness threshold for \(a\) in these bounds is independent
of \(R\).
\end{lemma}

\begin{proof}
On a line \(\Re w=c>1\), absolute convergence and Gaussian decay
justify interchanging the Dirichlet series and integral to obtain
\begin{equation}\label{eq:A051-vertical-gaussian}
 F_a(s)=\frac1{i\sqrt{\pi a}}
                  \int_{c-i\infty}^{c+i\infty}
                      F(w)e^{(s-w)^2/a}\,dw.
\end{equation}
For each individual term this follows by completing the square:
\[
 \frac{(s-w)^2}{a}-w\log n
    =\frac{(w-s-\frac a2\log n)^2}{a}
                                  -s\log n-\frac a4(\log n)^2.
\]
The remaining entire Gaussian integral along a vertical line is
\(i\sqrt{\pi a}\).

The contribution of the pole in \eqref{eq:A051-vertical-gaussian}
is exactly
\[
 \frac\kappa{i\sqrt{\pi a}}
     \int_{c-i\infty}^{c+i\infty}
            \frac{e^{(s-w)^2/a}}{w-1}\,dw
       =\kappa\int_0^\infty
               e^{-a u^2/4-(s-1)u}\,du.
\]
Indeed insert \((w-1)^{-1}=\int_0^\infty e^{-(w-1)u}\,du\);
absolute convergence on \(\Re w=c\) justifies Fubini, and the
same Gaussian evaluation then gives the displayed result.
At \(s=1+\sqrt a\,z\), it is \(\kappa E(z)/\sqrt a\).

For the \(G\) term shift the vertical contour to \(\Re w=\Re s\).
For small \(a\) and bounded \(z\), both lines lie in the stated
holomorphic strip. Polynomial vertical growth and Gaussian decay
make the horizontal boundary integrals tend to zero, uniformly
on each compact \(z\)-disk for a fixed \(a>0\).
Parameterizing the shifted line by \(w=s+i\sqrt a\,v\) gives
\begin{equation}\label{eq:A051-G-Gaussian}
 \frac1{i\sqrt{\pi a}}\int_{\Re w=\Re s}
                  G(w)e^{(s-w)^2/a}\,dw
         =\frac1{\sqrt\pi}\int_{\mathbb R}
                             e^{-v^2}G(s+i\sqrt a\,v)\,dv.
\end{equation}
For \(0<a\le1\), the polynomial growth bound dominates the
absolute integrand by
\(C_{F,K}e^{-v^2}(1+|v|)^{B_{F,K}}\), uniformly for \(|z|\le K\).
Its integral is finite. This proves \eqref{eq:A051-generalpole}.

When \(b_n\ge0\), the left side of
\eqref{eq:A051-general-mass} is \(F_a(1-K\sqrt a)\).
Apply \eqref{eq:A051-generalpole} at \(z=-K\) to obtain its bound.
For the tail note, term by term when \(\log n>R/\sqrt a\),
\[
 e^{-a(\log n)^2/4}
       \le e^{-R^2/8}e^{-(a/2)(\log n)^2/4}.
\]
Hence the tail is at most
\[
         e^{-R^2/8}F_{a/2}(1-K\sqrt a)
            \le \frac{D_{F,K}}{\sqrt a}e^{-R^2/8},
\]
where the last inequality is the same pole profile at the fixed
scaled coordinate \(-K\sqrt2\).
The bound on \(F_{a/2}\) and its smallness threshold do not depend
on \(R\), proving the final assertion.
\end{proof}

\begin{corollary}[NN-A051: unconditional Dedekind heat-zero height]
\label{cor:A051-dedekindheight}
Let \(K\) be any fixed number field of degree \(n=r_1+2r_2\),
with absolute discriminant \(D_K\), and put
\begin{align*}
 Q_K&=D_K^{1/2}\pi^{-r_1/2}(2\pi)^{-r_2},\\
 \xi_K(s)&=s(s-1)Q_K^s\Gamma(s/2)^{r_1}
                               \Gamma(s)^{r_2}\zeta_K(s),\\
 h_{K,0}(z)&=\xi_K((1+iz)/2).
\end{align*}
An immaterial constant factor in the completion does not change
the following assertion. Let \(h_{K,t}\) be this completed
de Bruijn--Newman family with \(\partial_t h_{K,t}=-h_{K,t}''\).
There are constants \(C_K,a_K>0\) such that for every
\(0<a<a_K\), it has a simple conjugate nonreal pair, one member
\(z_{K,a}\) satisfying
\[
 0<\Re z_{K,a}\le\exp\!\left(\exp\frac{C_K}{\sqrt a}\right),
                  \qquad \Im z_{K,a}\le-\frac12.
\]
This is unconditional. It does not use M, Artin holomorphy or
any assertion about the simplicity of the zeros of \(\zeta_K\).
No degree- or discriminant-uniform constants are claimed.
\end{corollary}

\begin{proof}
We recall exactly the classical interface needed here.
The Dedekind series has nonnegative coefficients counting integral
ideals, is absolutely convergent for \(\Re s>1\), and has a unique
simple pole at \(s=1\) with positive residue \(\kappa_K\).
Its indicated completion is entire after multiplication by
\(s(s-1)\), satisfies \(\xi_K(s)=\xi_K(1-s)\), and has finite
order. The ideal series, positive residue and root number \(+1\)
are stated in Voight, \emph{Quaternion Algebras} (2021),
Section~26.2, Theorem~26.2.3 and (26.2.4), and
Section~26.8, (26.8.3)--(26.8.5), respectively:
\url{https://link.springer.com/chapter/10.1007/978-3-030-56694-4_26}.
The source's completion differs from ours by the positive
constant \(2^{r_2}\) before multiplication by \(s(s-1)\).
For the entire order-one completion, a second precise source is
Grenie--Molteni (2019),
equations (2.5)--(2.8), author p.\ 5,
\url{https://sites.unimi.it/molteni/research/papers-pdf/44-molteni-An_explicit_Chebotarev_density_theorem_under_GRH.pdf}.
No GRH assumption is involved in these displayed classical
identities. As explained in Dobner's Section~2 after the functional
equation (4), the completed functional equation, absolute Dirichlet
convergence and finite order give polynomial growth of the
uncompleted series in fixed vertical strips by
Phragmen--Lindelof. Consequently
\(\zeta_K(s)-\kappa_K/(s-1)\) satisfies the local strip hypotheses
of Lemma~\ref{lem:A051-generalpole}.

Define the smoothed series
\[
 F_{K,a}(s)=\sum_{m\ge1}b_K(m)e^{-a(\log m)^2/4}m^{-s},
 \qquad b_K(m)=\#\{\mathfrak a:N\mathfrak a=m\}.
\]
Its coefficients have polynomial growth. For example, each Euler
factor is coefficientwise dominated by \((1-p^{-s})^{-n}\),
because there are at most \(n\) prime ideals above \(p\), with
positive integer residue degrees. Thus \(b_K(m)\le d_n(m)\),
which is a sufficient polynomial bound. The smoothed series is
entire, and Lemma~\ref{lem:A051-generalpole} gives
\[
             F_{K,a}(1+\sqrt a\,w)
                   =\frac{\kappa_K}{\sqrt a}E(w)+O_K(1)
\]
uniformly on the fixed simple-zero disk \(D\) used in the base
theorem.

Put \(\mu_K=\kappa_K\min_{\partial D}|E|\), and choose \(K_0\)
so that \(|w|\le K_0\) on \(D\).
The positive-coefficient mass and tail bounds of the lemma allow
fixed choices of \(R_K>0\) and an integer \(Q'_K\ge2\) such that,
with \(N=\lceil e^{R_K/\sqrt a}\rceil\),
\[
 \sum_{m>N}b_K(m)e^{-a(\log m)^2/4}m^{-1+K_0\sqrt a}
          \le\frac{\mu_K}{64\sqrt a},\qquad
 \frac{2\pi W_{K,K_0}}{Q'_K}\le\frac{\mu_K}{32}.
\]
Apply the direct homogeneous pigeonhole construction in
Lemma~\ref{lem:A051-recurrence} to the frequencies
\(\log m/(2\pi)\), for \(2\le m\le N\), starting at
\(Y=e^{250/a}\). The same proof, with these nonnegative coefficient
weights, gives
\[
 Y\le\tau\le Y(Q'_K)^{N-1},\qquad
 |F_{K,a}(1+\sqrt a\,w+i\tau)-F_{K,a}(1+\sqrt a\,w)|
                        \le\frac{\mu_K}{16\sqrt a}.
\]

Dobner's Theorem~4 applies with
\[
 \gamma_K(s)=s(s-1)Q_K^s\Gamma(s/2)^{r_1}\Gamma(s)^{r_2},
 \qquad
 J_{K,a}(s)=s+\frac{an}{4}\operatorname{Log}s+a c_K,
\]
where \(c_K=\tfrac12\log Q_K-\tfrac{r_1}{4}\log2\) is real.
On the shifted disk it gives the approximation to
\(F_{K,a}\) with error \(O_K(y^{-1/5}e^{40/a})\), divided by
the same nonzero gamma-exponential factor as in the base theorem.
The starting height makes this error at most
\(\mu_K/(16\sqrt a)\) for sufficiently small \(a\).
The profile error is also at most \(\mu_K/(8\sqrt a)\) eventually.
Rouche's theorem therefore gives exactly one zero, with
multiplicity one, in the profile disk.

Here \(J_{K,a}'(s)=1+an/(4s)\ne0\) at the high shifted point.
The transformed zero is simple. Since \(\Re s\ge3/4\),
\(\Im s\ge e^{250/a}/2\), and \(K\) is fixed, for small \(a\)
we have
\[
 \frac{an}{4}\log|s|+a c_K\ge0,
 \qquad \Re J_{K,a}(s)\ge3/4.
\]
Also \(\Im J_{K,a}(s)=\Im s+(an/4)\operatorname{Arg}s>0\)
and the resulting real heat coordinate is at most \(3\tau\).
The height budget \eqref{eq:A051-heightbudget}, with the fixed
\(R_K,Q'_K\), proves the asserted double exponential.
The completed Dedekind function is real under conjugation and
even about its critical center, so its heat family is real on the
real axis. Conjugation supplies the other simple zero at the same
positive height.
\end{proof}

\paragraph{Geometric scope.}
The coefficients \(b_K(m)\) count actual effective ideal divisors
on \(\operatorname{Spec}\mathcal O_K\). Their nonnegativity supplies
the tail control, and the simple pole supplies the universal
analytic profile; no independence of prime ideals or Artin factors
is used. The resulting nonreal reduced fiber points are local
unramified sections of the completed analytic heat divisor, as in
the base NN-A051 theorem. The construction does not identify an
arithmetic scheme for this heat family or a collision path to
time zero. In particular it does not discharge M in any separate
Dedekind zero-location or denominator-shift sign theorem.


\subsection{The first nonreal height and the least multiple endpoint}

\begin{corollary}[NN-A051: the exact finite-height alternative]
\label{cor:A051-firstlimit}
Assume RH, and use the normalization
\(H_0(z)=\xi(\tfrac12+iz/2)/8\) and
\(\partial_tH_t=-H_t''\). Define
\[
 T_{\mathrm{off}}(a)=
 \inf\{\Re z>0:H_{-a}(z)=0,\ \Im z\ne0\},\qquad a>0,
\]
and let \(x_*\) be the least positive multiple zero of \(H_0\),
with \(x_*=+\infty\) if every zero is simple.
Then, as an equality in the extended positive real line,
\begin{equation}\label{eq:A051-firstlimit}
                       \lim_{a\downarrow0}T_{\mathrm{off}}(a)=x_*.
\end{equation}
For all sufficiently small \(a>0\), the defining infimum is an
attained positive finite minimum. If \(x_*<\infty\), put
\[
 m=\operatorname{ord}_{x_*}H_0\ge2,\qquad
 \ell_*=\left.\frac{d}{dz}
       \log\!\left(\frac{H_0(z)}{(z-x_*)^m}\right)\right|_{z=x_*}.
\]
The logarithmic derivative here is independent of the local
choice of logarithm. It is real, and
\begin{equation}\label{eq:A051-firstdrift}
                T_{\mathrm{off}}(a)=x_*-2\ell_*a+O_{H_0,x_*}(a^2).
\end{equation}
Thus divergence of the first nonreal height is equivalent, under
RH, to simplicity of all endpoint zeros. This is a characterization,
not an assertion that the endpoint zeros are all simple or that
a multiple endpoint exists.
\end{corollary}

\begin{proof}
The positive theta kernel gives, for real \(y\),
\[
 H_{-a}(iy)=\int_0^\infty e^{-au^2}\Phi(u)\cosh(yu)\,du>0.
\]
Consequently there are no zeros on the imaginary axis, including
at the origin. NN-A022, whose cone proof uses RH alone, gives
a constant \(C\) such that every zero with positive real part
satisfies \(|\arg z|\le C\sqrt a\) for sufficiently small \(a\).
Choose the smallness threshold so that this angle is below
\(\pi/4\). Then every zero with \(0<\Re z\le R\) lies in
\(|z|\le\sqrt2 R\). In particular bounded positive real
coordinates cannot hide unbounded imaginary coordinates.

The base NN-A051 theorem gives a nonreal zero at every sufficiently
small \(a>0\). Fixing any such zero supplies a finite \(R\).
All possible zeros with smaller positive real coordinate lie in
the compact disk just described. There are only finitely many,
and none has real coordinate zero. Therefore the positive infimum
is attained and is finite.

We next use local uniform convergence \(H_{-a}\to H_0\).
In a fixed compact region, choose disjoint small disks centered
at the finitely many zeros of \(H_0\). On the complement of
these disks the limiting function has a positive minimum modulus,
and hence the approximating functions have no zero for small \(a\).
At a simple real center, Rouche's theorem gives exactly one zero
in a disk symmetric under conjugation. That zero is real and
simple. These arguments, together with the preceding cone bound,
show that no nonreal zero can have bounded positive real
coordinate unless it approaches a multiple real endpoint.
If all endpoints are simple, this proves
\(T_{\mathrm{off}}(a)\to+\infty\).

Suppose now \(x_*<\infty\). Choose \(\varepsilon_*>0\) so that
\((x_*-2\varepsilon_*,x_*+2\varepsilon_*)\) contains no other
zero of \(H_0\), and \(x_*-2\varepsilon_*>0\).
Every endpoint below \(x_*\) is simple. Applying the compact
argument to all zeros with positive real coordinate at most
\(x_*+\varepsilon_*\) shows that, for small \(a\), any nonreal
zero in this region must be in the chosen disk about \(x_*\).
There are exactly \(m\) zeros in that disk, counted with
multiplicity. We calculate their locations, including the term
needed to distinguish real and imaginary displacement.

Write, with real \(A_0\ne0\), \(\ell=\ell_*\) and real \(c_2\),
\[
 H_0(x_*+w)=A_0w^m(1+\ell w+c_2w^2+O(w^3)).
\]
Put \(q=\sqrt a>0\), and define the heat polynomials
\[
 P_j(u)=e^{D_u^2}u^j
       =\sum_{k=0}^{\lfloor j/2\rfloor}
          \frac{j!}{k!(j-2k)!}u^{j-2k}.
\]
Joint analyticity of the heat family near \((0,x_*)\), together
with \(H_{-a}=e^{aD_z^2}H_0\), gives the convergent local expansion
\begin{equation}\label{eq:A051-hermitejet}
 \frac{H_{-q^2}(x_*+qu)}{A_0q^m}
     =P_m(u)+q\ell P_{m+1}(u)+q^2c_2P_{m+2}(u)+O(q^3).
\end{equation}
The quotient extends holomorphically to \(q=0\), and the error
is uniform for \(u\) in a fixed compact set. To see the
divisibility explicitly, a joint Taylor monomial \(a^kw^j\)
has coefficient proportional to \(H_0^{(j+2k)}(x_*)\), so it
vanishes when \(j+2k<m\); substituting \(a=q^2,w=qu\)
therefore leaves no power below \(q^m\). Collecting the next
three powers gives exactly the displayed polynomials.

As in the Hermite calculation of NN-A022,
\[
       P_m(u)=(i\sqrt2)^m
                     \operatorname{He}_m(u/(i\sqrt2)).
\]
Its \(m\) roots \(r\) are distinct and purely imaginary;
zero is a root precisely when \(m\) is odd.
The analytic implicit-function theorem applied to
\eqref{eq:A051-hermitejet} gives one analytic branch \(u_r(q)\)
at each root. These \(m\) branches exhaust the disk's zeros
for small positive \(q\), and all are simple.

Here are the two expansion coefficients without a symmetry
shortcut. The identities
\[
 P_j'=jP_{j-1},\qquad
 P_{j+1}=uP_j+2jP_{j-1},\qquad
 2P_m''+uP_m'-mP_m=0
\]
follow directly from the displayed finite sum. At a root \(r\),
they give
\begin{align*}
 P_{m+1}(r)&=2P_m'(r),& P_{m+1}'(r)&=0,\\
 P_{m+2}(r)&=2rP_m'(r),&
 \frac{P_m''(r)}{P_m'(r)}&=-\frac r2.
\end{align*}
Writing \(u_r(q)=r+bq+dq^2+O(q^3)\), the successive
coefficients in \eqref{eq:A051-hermitejet} yield
\[
                 b=-2\ell,\qquad d=(\ell^2-2c_2)r.
\]
Therefore each resulting zero has the expansion
\begin{equation}\label{eq:A051-localfirstdrift}
 z_r(a)=x_*+r\sqrt a-2\ell_*a
                   +(\ell_*^2-2c_2)r a^{3/2}+O(a^2).
\end{equation}
For every nonzero \(r\), this is nonreal for small \(a>0\).
Since \(r\) is purely imaginary and all displayed coefficients
are real, its real part is \(x_*-2\ell_*a+O(a^2)\).
If \(m\) is odd, the branch at \(r=0\) is real: the equation
has real coefficients for real \(q\), and uniqueness in a small
conjugation-invariant neighborhood of zero forces
\(u_0(q)=\overline{u_0(q)}\).

There are thus \(2\lfloor m/2\rfloor\) nonreal simple zeros
near \(x_*\). For small \(a\) they lie below
\(x_*+\varepsilon_*\); the earlier compact argument excludes
all competing nonreal zeros of smaller real coordinate outside
this cluster. The minimum of finitely many quantities sharing
the real expansion \eqref{eq:A051-localfirstdrift} satisfies
\eqref{eq:A051-firstdrift}, and hence
\eqref{eq:A051-firstlimit}. This proves both alternatives.
\end{proof}

\paragraph{Geometric interpretation and remaining gate.}
At a finite multiple endpoint the local analytic zero divisor
has special fiber \(\mathbb C[w]/(w^m)\).
The base change \(a=q^2\), followed by normalization, separates
the \(m\) branches in \eqref{eq:A051-hermitejet}; the deck
involution pairs the nonzero Hermite roots \(r,-r\) and fixes
the real branch when \(m\) is odd. Thus local monodromy has
\(\lfloor m/2\rfloor\) disjoint transpositions. Every branch
away from \(q=0\) is reduced and unramified, with the same
first real drift. The compactness argument identifies this
particular finite cluster as the first nonreal-height cluster.
It does not compute any endpoint multiplicity, any lower escape
rate when all endpoints are simple, or the collision history of
the high pair constructed in the base theorem.


\section{An isolated pair: collision time and the surrounding zeros}
\label{sec:analytic-isolated-collision}

\begin{theorem}[NN-A023: an isolated-pair collision with its first interaction correction]
Let $f$ be a real entire function of order at most one, with
only real zeros, and let $b-\delta/2,b+\delta/2$ be adjacent
simple zeros, $\delta>0$. Put
\[
 G(z)=\frac{f(z)}{(z-b)^2-\delta^2/4},\qquad
 L=\frac{G'}G(b)\in\mathbb R,\qquad
 S=\sum_{G(\xi)=0}\frac1{(b-\xi)^2},\qquad
                       \kappa=\delta^2S.
\]
Zeros in $S$ are counted with multiplicity. Assume the
order-one Hadamard product for $G$; its possible linear
exponential factor is allowed and is retained in $L$.
The case $S=0$ is included.

For all sufficiently small $\kappa$, with an absolute
threshold, the two zeros continued from this pair under
$f_{-a}=e^{a\partial_z^2}f$ undergo a unique double collision
on the local time scale $0\le a\le\delta^2/4$. Its time and
real position satisfy
\begin{align}
 a_c&=\delta^2\left(\frac18+\frac{\kappa}{32}
                                      +O(\kappa^2)\right),
                        \label{eq:analytic-isolated-time}\\
 z_c&=b-2a_cL+O(\delta\kappa^{3/2}).
                        \label{eq:analytic-isolated-center}
\end{align}
The $O$ constants are absolute. Before the collision the
tracked zeros are real, and after it they are a nonreal
conjugate pair. The local neighbourhood moves with the real
center $b-2aL$; no smallness assumption on $\delta L$ is needed.

Under A this applies to every adjacent simple pair of $H_0$.
In particular any sequence of such arithmetic pairs with
$\kappa\to0$ has actual collision times
$t_c=-\delta^2(1+\kappa/4+O(\kappa^2))/8$.
This is a conditional statement about the stated pair
parameter, not a proof of how often these pairs occur.
\end{theorem}

\begin{proof}
Hadamard factorization has no exponential term of degree two
or higher, by the order assumption. Therefore
\[
 R(h)=\frac{G(b+h)}{G(b)e^{Lh}}
       =\prod_{G(\xi)=0}
         \left(1+\frac h{b-\xi}\right)
                           \exp\left(-\frac h{b-\xi}\right)
\]
converges locally uniformly, and the proof of
Lemma~\ref{lem:analytic-product-remainder} gives
$|R(h)|\le e^{S|h|^2/2}$. Its first Taylor coefficients are
\begin{equation}\label{eq:analytic-pair-R-expansion}
 R(h)=1-\frac S2h^2+\frac{T_3}{3}h^3+\mathcal E_4(h),
          \qquad T_3=\sum_{G(\xi)=0}(b-\xi)^{-3},
          \qquad |T_3|\le S^{3/2}.
\end{equation}
There is an absolute constant $C$ such that
\begin{equation}\label{eq:analytic-pair-R-fourth}
 |\mathcal E_4(h)|\le
               CS^2|h|^4e^{2S|h|^2}\qquad(h\in\mathbb C).
\end{equation}
Indeed, after putting $v=h\sqrt S$, Cauchy's estimate on
$|v|=1$ bounds the remainder by $2e^{1/2}|v|^4$ for
$|v|\le1/2$. For $|v|\ge1/2$, use the global growth bound,
the explicit quadratic coefficient and
$|T_3|/(3S^{3/2})\le1/3$; each of the four terms is bounded
by a fixed multiple of $|v|^4e^{2|v|^2}$.
If $S=0$, the product is empty and $R=1$, so the claimed
collision and formulas are exact.

Write $a=\delta^2\tau$ and use the moving coordinate
\[
                         z=b-2aL+\delta w.
\]
Completing the Gaussian square, or conjugating the heat
operator by $e^{L(z-b)}$, gives the exact identity
\begin{equation}\label{eq:analytic-pair-scaled-family}
 \frac{f_{-\delta^2\tau}(b-2\delta^2\tau L+\delta w)}
      {\delta^2G(b)\exp(\delta wL-\delta^2\tau L^2)}
  =J(\tau,w)
  =\mathbb E\left[(Y^2-\tfrac14)R(\delta Y)\right],
                \qquad Y=w+\sqrt{2\tau}\,V,\quad V\sim N(0,1).
\end{equation}
For real $\tau\ge0$ this is an ordinary Gaussian convolution
with a real shift $2aL$. Both sides extend holomorphically
in $\tau$ to a complex neighbourhood of $[0,1/4]$:
one may choose a square-root branch locally away from zero;
the expectation is even in that square root and extends
across zero. The order-one growth and Gaussian decay
justify the continuation. The denominator is nonzero.

For bounded complex $w,\tau$, the growth bound for $R$
is dominated by $\exp(C\kappa(|V|^2+1))$.
Consequently \eqref{eq:analytic-pair-R-fourth}, integrated
against the Gaussian, gives a uniform holomorphic expansion
with an $O(\kappa^2)$ remainder:
\begin{align}
 J(\tau,w)
  &=w^2+2\tau-\tfrac14-\frac{\kappa}{2}P_4(\tau,w)
                  +\eta_3 P_5(\tau,w)+O(\kappa^2),
                    \label{eq:analytic-pair-J-expansion}\\
 \eta_3&=\delta^3T_3/3,\qquad |\eta_3|\le\kappa^{3/2}/3,\nonumber\\
 P_4(\tau,w)
  &=w^4+(12\tau-\tfrac14)w^2+12\tau^2-\tfrac12\tau,\nonumber\\
 P_5(\tau,w)
  &=w^5+(20\tau-\tfrac14)w^3
                         +(60\tau^2-\tfrac32\tau)w.\nonumber
\end{align}
All derivatives on smaller compact sets have the same
remainder order by Cauchy's formula. In particular the
quadratic first correction is even in $w$, while the
cubic correction is odd. Uniform constants depend only
on the fixed compact sets, not on $f,\delta,L,S$.

For $0\le\tau\le1/4$, the limit polynomial
$w^2+2\tau-1/4$ has exactly two roots in $|w|<2$ and
none on the boundary. Rouch\'e's theorem therefore tracks
exactly two roots of $J$ in this moving disc for small
$\kappa$, starting with the given pair at $\tau=0$.
Its derivative is $2w+O(\kappa)$ uniformly there and has
a single zero near zero. The implicit function theorem
gives this zero as a real analytic function $w_*(\tau)$
for real $\tau$, with
$w_*(\tau)=O(\kappa^{3/2})$: evaluate
$\partial_wJ(\tau,0)$ using the parity in
\eqref{eq:analytic-pair-J-expansion}, and use
$\partial_w^2J=2+O(\kappa)$.

Set $D(\tau)=J(\tau,w_*(\tau))$. The same expansion yields
\[
 D(\tau)=2\tau-\tfrac14
              -\frac{\kappa}{2}(12\tau^2-\tfrac12\tau)
              +O(\kappa^2),\qquad D'(\tau)=2+O(\kappa).
\]
It has one real zero $\tau_c$ near $1/8$. First
$\tau_c=1/8+O(\kappa)$; evaluating the correction at $1/8$,
where $12\tau^2-\tau/2=1/8$, then gives
\[
                       \tau_c=\frac18+\frac{\kappa}{32}
                                                  +O(\kappa^2).
\]
The vanishing conditions $J=J_w=0$, together with
$J_{ww}=2+O(\kappa)\ne0$, give a double zero, not a higher
collision. Since $D'$ and $J_{ww}$ are positive for real
arguments here, its two branches are real for
$\tau<\tau_c$ and nonreal conjugate for $\tau>\tau_c$.
No other local collision is possible on $[0,1/4]$:
the derivative has only the one critical point and $D$
has only the one zero. Returning to $a,z$ proves
\eqref{eq:analytic-isolated-time} and
\eqref{eq:analytic-isolated-center}. This proof also
accounts for the full pair throughout its local interval;
it does not merely locate a nonreal zero at one time.
\end{proof}

\paragraph{Relation to the classical Lehmer criterion.}
Csordas--Smith--Varga, \emph{Lehmer Pairs of Zeros, the
de Bruijn--Newman Constant $\Lambda$, and the Riemann
Hypothesis}, Constructive Approximation~10 (1994), 107--129,
Definition~1, equations~(1.10)--(1.12), printed pp.~109--110,
uses the endpoint sum
\[
 g=\sum_{\xi\ne b-\delta/2,b+\delta/2}
       \left\{(b-\delta/2-\xi)^{-2}
                       +(b+\delta/2-\xi)^{-2}\right\}.
\]
Under A this is positive. For $\kappa\to0$,
$g=2S+O(\delta^2S^2)$, by expanding the two summands
symmetrically and using $\sum(b-\xi)^{-4}\le S^2$.
Thus $\delta^2g=2\kappa+O(\kappa^2)$, and the pairs above
eventually satisfy their $\delta^2g<4/5$ criterion.
Their Theorem~1, equations~(1.13)--(1.14), printed p.~110,
states the lower bound
\[
 \frac{(1-\tfrac54\delta^2g)^{4/5}-1}{8g}\le\Lambda.
\]
It does not assert that this displayed lower bound is an
exact collision time. Its first two expansion terms agree
with \eqref{eq:analytic-isolated-time} after changing the
sign of time. The argument here identifies an actual
unique local double collision and its interaction correction.
The original printed pages were inspected directly.

\paragraph{Geometric and representation meaning.}
The chosen two-point divisor determines a finite local
rank-two zero family in the moving coordinate. Its
discriminant has a simple zero at $\tau_c$, so after a
quadratic base change the two branches are separated.
For the resulting finite locally free rank-two algebra,
the exact \'{e}tale/discriminant criterion is SGA~1,
Expos\'e~I, section~4, Proposition~10; section~3,
Proposition~1 gives the equivalent vanishing of relative
differentials. This criterion is applied to the actual
local finite algebra, not to all of the infinite divisor.
The first correction is the inverse-square trace of the
complementary zero divisor, while its logarithmic first
trace produces only the displayed real translation.
This is a concrete local discriminant computation; it
does not make the complementary analytic divisor into a
finite pure Weil object. In the Toeplitz/character
dictionary of NN-A020--A022, each such collision creates
one conjugate pair of complex product parameters.
The relevant exterior-power operations have the locator
SGA~6, Expos\'e~V, Example~2.2(b) and Lemma~2.2.1.
No integral canonical-basis statement follows from
the analytic inverse-square parameter.


\section{The first CUE heat collision is an extreme-gap event}
\label{sec:analytic-cue-first-collision}

Let $0\le\theta_1<\cdots<\theta_N<2\pi$ be the eigenangles
of Haar $U(N)$, and extend them periodically by
$\theta_{j+N}=\theta_j+2\pi$. The real characteristic
trigonometric polynomial is
\[
 Z_N(z)=\prod_{j=1}^N2\sin\frac{z-\theta_j}{2},\qquad
 Z_{N,-a}=e^{a\partial_z^2}Z_N .
\]
For odd $N$ it is antiperiodic, and its zero divisor is
still $2\pi$-periodic. Negative heat time has exactly the
same sign convention as $H_{-a}$ in the arithmetic family.
Define
\[
 \delta_N=\min_j(\theta_{j+1}-\theta_j),\qquad
 A_N=\inf\{a>0:Z_{N,-a}\text{ has a nonreal zero}\}.
\]
This section concerns this specified CUE model, with no
arithmetic transfer assumption.

\begin{theorem}[NN-A024: first CUE collision law]
For $N\ge2$, $A_N$ is positive and finite almost surely.
As $N\to\infty$,
\[
                    \frac{8A_N}{\delta_N^2}
                                  \longrightarrow1
                                  \quad\hbox{in probability}.
\]
More precisely the ratio is $1+O_{\mathbb P}(N^{-1/3})$.
With probability tending to one, the first collision is
the one continued from the initially smallest cyclic gap.
For $u\ge0$,
\begin{equation}\label{eq:analytic-cue-first-time-law}
 \lim_{N\to\infty}
 \mathbb P\left\{
       \frac{8N^{8/3}A_N}{(72\pi)^{2/3}}>u\right\}
                              =e^{-u^{3/2}}.
\end{equation}
Equivalently, after unfolding the mean spacing to one,
the first negative heat time is of order $N^{-2/3}$,
smaller than the mean-spacing-squared scale, which is one.
\end{theorem}

\begin{proof}
We first prove a deterministic global lower bound for
the first loss of reality. Write
\[
 Z_N(z)=\sum_{k=0}^Nc_ke^{i(k-N/2)z}.
\]
The coefficients $c_0,c_N$ are nonzero. Under
$e^{a\partial_z^2}$ the $k$th coefficient is multiplied
by $e^{-a(k-N/2)^2}$. After multiplication by
$e^{iNz/2}$ this is a degree-$N$ polynomial in $w=e^{iz}$.
Its leading and constant coefficients remain nonzero,
and on a compact $a$ interval the normalized polynomial
coefficients and their reciprocals at the two ends stay
bounded. The $N$ roots therefore remain in a compact
annulus in $\mathbb C^*$.

While all these roots are on the unit circle and simple,
their angles have real analytic continuations $y_j(a)$.
There is no loss of a zero at imaginary infinity at a
finite time; and a simple real zero remains real by the
implicit function theorem and conjugation. Thus a first
loss of reality must occur at a finite real collision.
Before that collision the factorization into real sine
factors gives
\begin{equation}\label{eq:analytic-circular-backward-ode}
 y_i'(a)=-\sum_{j\ne i}\cot\frac{y_i(a)-y_j(a)}2.
\end{equation}
For an empty adjacent arc of length
$g_i=y_{i+1}-y_i\in(0,2\pi)$, the two mutual terms give
$-2\cot(g_i/2)$ to $g_i'$. Every remaining cotangent
difference is negative before the minus sign: along
this empty arc the cotangent is strictly decreasing
and crosses no pole. Therefore
\[
                  g_i'(a)\ge-2\cot(g_i(a)/2),
                  \qquad (g_i^2)'(a)\ge-8.
\]
There is a slightly sharper integrated form:
\[
       \frac d{da}\cos(g_i(a)/2)
                       \le\cos(g_i(a)/2).
\]
If this gap closes at $A$, its left side integrates to
$e^{-A}\le\cos(g_i(0)/2)$. Since
$g_i(0)\ge\delta_N$ and $\delta_N<\pi$ almost surely,
\begin{equation}\label{eq:analytic-first-time-lower}
       A_N\ge-\log\cos(\delta_N/2)\ge\delta_N^2/8.
\end{equation}
This rules out an earlier collision anywhere on the
circle, not merely near the smallest pair.

For completeness, the first collision is almost surely
finite. When $N$ is even the constant Fourier coefficient
is almost surely nonzero and eventually dominates all
others, so for large $a$ there are no real zeros.
When $N$ is odd, the almost surely nonzero frequencies
$\pm1/2$ eventually dominate, leaving one simple real
zero modulo $2\pi$, fewer than $N$ for $N\ge3$.
Uniform $C^1$ convergence after rescaling gives this
last zero count. Each asserted Fourier coefficient is
a nonzero real analytic function of the eigenangles:
one can check nonvanishing when the angles are all
close to zero, where the limiting function is
$(2\sin(z/2))^N$. Its vanishing set has Lebesgue
measure zero, and the CUE law has an absolutely
continuous density. The preceding annulus argument
then supplies a finite first collision.

We next prove a uniform isolation estimate for the
smallest pair. Put $\ell_N=N^{-7/6}$. With probability
tending to one, every circular interval of length
$\ell_N$ contains at most two eigenangles.
Here is a direct determinantal proof. Let
\[
 v_N(x)=(1,e^{ix},\ldots,e^{i(N-1)x})/\sqrt{2\pi}.
\]
The three-point intensity is
$\rho_3(x_1,x_2,x_3)=
 \|v_N(x_1)\wedge v_N(x_2)\wedge v_N(x_3)\|^2$.
Successive divided differences give
\[
 \rho_3(x_1,x_2,x_3)
       \le C N^9\prod_{i<j}|x_i-x_j|^2.
\]
Indeed the norms of $v_N,v_N',v_N''$ are respectively
$O(N^{1/2}),O(N^{3/2}),O(N^{5/2})$. The determinant
of the three divided-difference columns extracts the
Vandermonde product, and their norms are bounded by
these derivative norms (with a factor $1/2$ in the
second divided difference). This applies to real
lifts of any short arc, uniformly around the circle.
The expected number of triples lying in a forward
arc of length $\ell_N$ is therefore at most
\[
 C N^9\ell_N^8=O(N^{-1/3}).
\]
If such an arc contains three points, one of its points
can be chosen as the start of a forward arc containing
the other two. Markov's inequality proves the claim.

The classical extreme-gap theorem of Ben Arous--Bourgade
gives
\begin{equation}\label{eq:analytic-CUE-min-gap}
 \frac{N^{4/3}\delta_N}{(72\pi)^{1/3}}
          \ \Longrightarrow\ Y,\qquad
                    \mathbb P\{Y>y\}=e^{-y^3}\quad(y\ge0).
\end{equation}
In particular $\delta_N=O_{\mathbb P}(N^{-4/3})$ and
$\delta_N/\ell_N\to0$ in probability.
Choose the minimum pair, using real lifts
$b-\delta_N/2,b+\delta_N/2$. On the preceding isolation
event, for $\delta_N\le\ell_N/2$ every other zero on the
periodic real line has distance at least $\ell_N/2$
from $b$. Each interval of length $\ell_N/2$ contains
at most two zeros. Summing successive distance shells
therefore gives
\[
 S_N=\sum_{\substack{Z_N(\xi)=0\\
                   \xi\ne b-\delta_N/2,b+\delta_N/2}}
                     (b-\xi)^{-2}
                          \le C\ell_N^{-2}.
\]
The periodic copies are included in this sum. Thus the
pair parameter satisfies
\[
                    \kappa_N=\delta_N^2S_N
                             =O_{\mathbb P}(N^{-1/3}).
\]
This includes the exceptional probability
$O(N^{-1/3})$ of the isolation event failing:
outside that event no bound is asserted, while its
probability tends to zero, as required for the stated
bound in probability.

The real entire function $Z_N$ has order one and only
real zeros. Removing the selected pair leaves an
order-one function with the required canonical product.
NN-A023 applies, including its unrestricted real drift.
It gives an actual collision of that pair at
\[
 a_{\min}=\frac{\delta_N^2}{8}
                    \{1+\kappa_N/4+O(\kappa_N^2)\}.
\]
The first global collision can occur no later than this
one. Combine this upper bound with
\eqref{eq:analytic-first-time-lower} to obtain
$8A_N/\delta_N^2=1+O_{\mathbb P}(N^{-1/3})$.

To identify the colliding pair, let $\delta_N^{(2)}$
be the second-smallest initial gap. The same source's
joint extreme-gap theorem implies that
$\delta_N^{(2)}/\delta_N$ converges in distribution
to a variable strictly larger than one almost surely.
Since $\kappa_N\to0$ in probability, with probability
tending to one
\[
                    (\delta_N^{(2)})^2/8>a_{\min}.
\]
Every other initially labelled cyclic gap has its
first possible closing time at least its initial
squared length divided by eight, by the universal
gap inequality. Hence none can collide first before
$a_{\min}$, proving the assertion about the smallest pair.
Finally combine \eqref{eq:analytic-CUE-min-gap} with
$8A_N/\delta_N^2\to1$ and the continuous mapping theorem:
the variable in \eqref{eq:analytic-cue-first-time-law}
converges to $Y^2$, whose survival function is
$e^{-u^{3/2}}$.
\end{proof}

\paragraph{Exact two-angle normalization check.}
For just two angles $b\pm\delta/2$,
\[
 Z_2(z)=2\cos(\delta/2)-2\cos(z-b),\qquad
 Z_{2,-a}(z)=2\cos(\delta/2)-2e^{-a}\cos(z-b).
\]
Consequently
\[
              A_2=-\log\cos(\delta/2)
                  =\delta^2/8+\delta^4/192+O(\delta^6)
\]
when $0<\delta<\pi$. The complementary periodic zero
sum in NN-A023 is
\[
          S=\frac1{2\sin^2(\delta/4)}-\frac8{\delta^2}
                                  =\frac16+O(\delta^2).
\]
Its correction $\delta^2\kappa/32$ is exactly
$\delta^4/192$ to first order. This checks both the
heat-sign normalization and the interaction coefficient
by a closed-form example.

\paragraph{Primary locator and arithmetic scope.}
The only imported random-matrix asymptotic is
Ben Arous--Bourgade, \emph{Extreme gaps between eigenvalues
of random matrices}, Annals of Probability~41 (2013),
2648--2681, DOI~10.1214/11-AOP710, Theorem~1.1
printed p.~2650 and Corollary~1.2 p.~2651.
Their full-circle gap intensity is
$u^2\,du/(24\pi)$ for the scale $N^{4/3}\delta$,
giving the constant $72\pi$ in
\eqref{eq:analytic-CUE-min-gap}. The no-triples estimate,
heat dynamics, global earliest-collision bound and
passage from gaps to collision times were proved above.
The theorem uses neither A nor M.

The model also distinguishes two different heuristic
questions in the supplied consultation. Matching
$\sqrt a$ to the mean arithmetic gap gives a
\emph{typical-gap} or bulk scale
$\log X\asymp a^{-1/2}$. It is not the scale of the
first rare collision in the CUE model. If one
additionally postulates the analogous cubic small-gap
Poisson law for arithmetic zeros, together with
isolation and exhaustion of early collisions by
those pairs, the heat-coordinate density
$D(x)\sim\log x/(4\pi)$ predicts a count
\[
 \frac{\pi^2}{9}(8a)^{3/2}\int_0^X D(x)^4\,dx
           \sim\frac{\sqrt2}{144\pi^2}
                              a^{3/2}X(\log X)^4.
\]
The first-event scale in that stronger model is
$X\asymp a^{-3/2}(\log(1/a))^{-4}$, rather than an
exponential scale. This displayed arithmetic conversion
is explicitly a conditional model prediction.
Neither A, Montgomery's restricted pair-correlation
theorem, nor a fixed-scale sine-kernel limit supplies
the required extreme-gap process or the exhaustion
of all early arithmetic collisions. NN-A023 proves
the local dynamics once the isolation parameter is small.

\paragraph{An actual algebraic family and its discriminant.}
Put $q=e^{-a/4}$ and $w=e^{iz}$. The preceding Fourier
formula gives the polynomial
\[
              P_N(q,w)=
                    \sum_{k=0}^Nc_kq^{(2k-N)^2}w^k
                 \in\mathbb C[q,q^{-1},w,w^{-1}].
\]
Its leading and constant coefficients are units over
$\mathbb C[q,q^{-1}]$. Dividing by the leading
coefficient exhibits
\[
 \operatorname{Spec}
       \mathbb C[q,q^{-1},w,w^{-1}]/(P_N)
                      \longrightarrow\mathbb G_{m,q}
\]
as a finite free family of degree $N$; inversion of $w$
is automatic because its constant coefficient is a unit.
It is \'{e}tale precisely off its polynomial
discriminant, by the derivative or trace-pairing
criterion. The exact Grothendieck locator is
SGA~1, Expos\'e~I, section~4, Proposition~10
(invertibility of the discriminant for a finite
locally free cover), with Definition~1 in that
section and section~3, Proposition~1 for the
flat/unramified and differential formulations.
At the first collision identified above with
probability tending to one, the local rank-two
component has a simple discriminant zero and
ramified double-cover normalization as in NN-A023.
The circle is the fixed locus of
$w\mapsto1/\overline w$ on the real parameter ray.
This is a genuine complex algebraic zero family
for each CUE sample; it is not an arithmetic
Frobenius construction or a canonical-basis model.


\section{The complete early CUE collision process}
\label{sec:analytic-cue-collision-process}

The first-event result NN-A024 does not by itself control the
other zeros after its first pair becomes nonreal. The following
uniform local argument supplies that missing exhaustion step.
The heat normalization is unchanged:
$Z_{N,-a}=e^{a\partial_z^2}Z_N$ and
$Z_N(z)=\prod_{j=1}^N2\sin((z-\theta_j)/2)$.

\begin{lemma}[Uniform one-root and two-root tracking]
\label{lem:analytic-uniform-clusters}
Let $f$ be real entire of order at most one with all zeros real.
Remove either a simple zero $b$, or two adjacent simple zeros
$b\pm\delta/2$, and call the remaining entire factor $G$.
Put $L=G'(b)/G(b)$ and $S=\sum_{G(\xi)=0}(b-\xi)^{-2}$.
Possible linear exponential factors are included in $L$.
There are absolute positive constants $c,C$ with the following
properties for $T>0$.

If one zero was removed and $TS<c$, then for every $a\in[0,T]$
exactly one zero of $e^{a\partial_z^2}f$ lies in the moving disc
\[
                  |z-(b-2aL)|<2\sqrt T.
\]
It is real and simple throughout the interval.

If two zeros were removed, put $R^2=\delta^2+T$. If
$R^2S<c$, exactly two zeros lie in
$|z-(b-2aL)|<2R$ for every $a\in[0,T]$. They have at most
one collision there. The collision is double and crosses from
two real zeros to a nonreal conjugate pair; the pair cannot
return to the real axis during $[0,T]$. Its collision time,
when in this interval, is the time supplied by NN-A023:
\[
             a_c=\frac{\delta^2}{8}
                \bigl(1+\delta^2S/4+O((\delta^2S)^2)\bigr).
\]
The constants are independent of $L$, and of the ratio
$\delta^2/T$.
\end{lemma}

\begin{proof}
Use the canonical remainder
\[
 \mathcal R(h)=G(b+h)/(G(b)e^{Lh}),\qquad
 |\mathcal R(h)|\le e^{S|h|^2/2},\qquad
 |\mathcal R(h)-1|\le 8S|h|^2e^{2S|h|^2},
\]
proved in NN-A022. Conjugating the Gaussian heat operator by
$e^{L(z-b)}$ removes the common drift exactly. In the one-root
case, with $z=b-2aL+\sqrt T\,u$ and $\tau=a/T$, division by
$\sqrt T G(b)e^{\sqrt T uL-aL^2}$ gives
\[
 I(\tau,u)=\mathbb E\bigl[Y\mathcal R(\sqrt T Y)\bigr],
                    \qquad Y=u+\sqrt{2\tau}\,V,
                    \qquad V\sim N(0,1).
\]
On any fixed bounded complex $u,\tau$ set this is
$I=u+O(TS)$, with every derivative on smaller compact sets
bounded in the same way. To justify this including complex
$\tau$, the Gaussian expectation is even in its square-root
parameter. Its absolute integrand is dominated by
$C(1+|V|)^C e^{C TS|V|^2}$; for small $TS$ the usual Gaussian
density dominates it uniformly. The expression therefore
extends holomorphically through $\tau=0$, and the stated
bounds and their derivatives follow by Cauchy's formula.
Rouch\'e on $|u|=2$ gives exactly one zero. The derivative
$I_u=1+O(TS)$ there is nonzero at that zero. Real conjugation
and uniqueness make the zero real for real $\tau$.

For the two-root case scale by $R$ instead, put
$d=\delta/R\in(0,1]$, $\tau=a/R^2\in[0,1]$, and divide by
$R^2G(b)e^{RuL-aL^2}$. The resulting exact family is
\[
 J(\tau,u)=\mathbb E\bigl[(Y^2-d^2/4)\mathcal R(RY)\bigr]
            =u^2+2\tau-d^2/4+O(R^2S).
\]
The same Gaussian argument is uniform in $0<d\le1$ and
controls all derivatives on fixed compact sets. On $|u|=2$
for real $0\le\tau\le1$, the leading polynomial has modulus
at least $4-2-1/4=7/4$. Thus exactly two zeros remain inside.
Its derivative has exactly one zero there, near zero, and
$J_{uu}=2+O(R^2S)>0$ on the real portion of this disc.
The critical point $u_*(\tau)$ is real analytic. Define
$D(\tau)=J(\tau,u_*(\tau))$. Because $J_u=0$ at $u_*$,
\[
                         D'(\tau)=J_\tau(\tau,u_*)
                                      =2+O(R^2S)>0.
\]
At $\tau=0$ the two roots are exactly $\pm d/2$; strict
convexity gives $D(0)<0$, even when $d$ is very small.
Consequently $D$ has at most one zero in this time interval.
At that zero $J_{uu}$ and $D'$ are positive and nonzero, so
it is a simple discriminant zero with the stated real-to-complex
crossing. The two roots cannot return without another zero of
$D$. Applying NN-A023 on its own local time interval identifies
the collision time; whenever both intervals apply they track
the same two zeros by their common initial sections and local
uniqueness. No condition on $\delta L$ or $\sqrt T L$ was used.
\end{proof}

\begin{theorem}[NN-A025: Poisson process of early CUE collisions]
\label{thm:analytic-cue-collision-process}
For Haar CUE, let $(a_{N,j},b_{N,j})$ be the real double-zero
collisions of $Z_{N,-a}$, with $a_{N,j}>0$ and real position
$b_{N,j}$ taken modulo $2\pi$. On any fixed rescaled time
window, with probability tending to one, all multiple-zero
events are such simple double collisions; each creates one
nonreal conjugate pair, and these pairs remain nonreal on that
window. The marked point process
\[
 \sum_j\delta_{(v_{N,j},\,b_{N,j}\bmod2\pi)},
 \qquad v_{N,j}=\frac{8N^{8/3}a_{N,j}}{(72\pi)^{2/3}},
\]
converges on bounded time windows to a Poisson point process
with intensity
\begin{equation}\label{eq:analytic-collision-ppp-intensity}
                   \frac32\sqrt v\,dv\,\frac{db}{2\pi}.
\end{equation}
In particular the number of nonreal conjugate pairs at time
$a=(72\pi)^{2/3}v/(8N^{8/3})$ converges in distribution to
$\operatorname{Poisson}(v^{3/2})$. The pair-count processes
converge in finite-dimensional distributions to the corresponding
inhomogeneous Poisson counting process. If $V_k$ is the limiting
$k$th collision time, then
\[
 V_k^{3/2}\sim\operatorname{Gamma}(k,1),\qquad
 f_{V_k}(v)=\frac{3}{2(k-1)!}
                       v^{3k/2-1}e^{-v^{3/2}}\quad(v>0).
\]
The limiting collision positions are independent uniform circular
marks, independent of the times. This is a CUE theorem; it
asserts no transfer to arithmetic zero statistics.
\end{theorem}

\begin{proof}
Fix $A>0$ and put
\[
 s_N=N^{-4/3},\quad T_N=A s_N^2,\quad
 \ell_N=N^{-7/6},\quad B_N=N^{1/24},\quad h_N=B_Ns_N.
\]
The three-point divided-difference estimate in NN-A024 gives
\[
 \mathbb P(\text{some circular arc of length }\ell_N
                        \text{ contains three zeros})
       \le C N^9\ell_N^8=O(N^{-1/3}).
\]
Work on the complementary event. Call each adjacent gap
$\delta\le h_N$ a small-pair cluster and each remaining root
a singleton. Since $2h_N<\ell_N$ eventually, no root belongs
to two small pairs. Every other root lies at distance at least
$\ell_N-h_N\ge\ell_N/2$ from a small pair. Two singleton
roots are at least $h_N$ apart. These assertions apply also
across the circular cut, by using periodic real lifts.

We need uniform control of the other periodic zeros. The
no-triples condition gives at most two zeros in each interval
of length $\ell_N$ on the periodic real line. Shell summation
therefore gives, for a pair center with its two selected roots
removed,
\begin{equation}\label{eq:analytic-cue-uniform-pair-SL}
 S_{\rm pair}\le C\ell_N^{-2},\qquad
 |L_{\rm pair}|\le
                  C\ell_N^{-1}\log(2/\ell_N).
\end{equation}
For the inverse-square sum, bound the shells at distances
$k\ell_N/2$ by $C(k\ell_N)^{-2}$ and sum $k^{-2}$,
including all periodic copies. For the logarithmic drift,
the original trigonometric product gives exactly
\[
 L_{\rm pair}=\frac12\sum_{j\ {
m outside\ pair}}
                         \cot((b-\theta_j)/2).
\]
The selected two cotangents cancel at their midpoint, as do
the two linear-factor logarithmic derivatives. On circular
distances $d\in(0,\pi]$, $|\cot(d/2)|\le C/d$; summing the
finitely many shells gives the logarithm in
\eqref{eq:analytic-cue-uniform-pair-SL}. This avoids assigning
a nonconvergent unsigned first trace to all periodic lifts.

For a singleton at $b$, its other zeros are at least $h_N$
away. At most a fixed number can be closer than $\ell_N$;
the remaining shells give
\begin{equation}\label{eq:analytic-cue-uniform-single-SL}
 S_{\rm single}\le C(h_N^{-2}+\ell_N^{-2}),\qquad
 |L_{\rm single}|\le
               C(h_N^{-1}+\ell_N^{-1}\log(2/\ell_N)).
\end{equation}
Here $L_{\rm single}=\frac12\sum_{j\ne i}
\cot((b-\theta_j)/2)$; the regular term from its own sine
factor vanishes at $b$.

Uniformly over all pair clusters,
\[
 (\delta^2+T_N)S_{\rm pair}
       \le C(h_N^2+T_N)\ell_N^{-2}=O(N^{-1/4})=o(1).
\]
Uniformly over all singletons,
\[
 T_NS_{\rm single}
       \le C_A(B_N^{-2}+N^{-1/3})=o(1).
\]
Lemma~\ref{lem:analytic-uniform-clusters} therefore applies
simultaneously to every cluster on $[0,T_N]$. Their moving
discs remain pairwise disjoint. Indeed pair discs have radii
$2\sqrt{\delta^2+T_N}=O_A(h_N)=o(\ell_N)$ and their center
shifts are at most $C_AT_N\ell_N^{-1}\log(2/\ell_N)
=o(\ell_N)$. Singleton discs have radius $2\sqrt{T_N}=O_A(s_N)$;
their shifts, by \eqref{eq:analytic-cue-uniform-single-SL}, are
\[
 O_A(T_N/h_N+T_N\ell_N^{-1}\log(2/\ell_N))=o(s_N).
\]
They are separated from one another by at least $h_N=B_Ns_N$,
and from any pair by at least $\ell_N/2$ before these vanishing
relative displacements. Each disc has diameter tending to zero,
so disjointness also holds modulo $2\pi$.

The finite polynomial in $w=e^{iz}$ from NN-A024 has degree
$N$, with nonzero constant and leading coefficients at every
finite $a$. The disjoint discs contain a total of exactly $N$
zeros modulo $2\pi$, counted with multiplicity. They therefore
exhaust the entire zero divisor throughout $[0,T_N]$. This is
the step that remains valid after some pairs become nonreal;
no sign inequality for the later global root flow is required.
The singleton roots stay real and simple. Each pair has exactly
its unique collision, if that time is in the window, and cannot
return during the window. In particular there are no other
early collisions or undetected nonreal roots.

For every small pair the uniform collision expansion gives
\[
 a_c=\frac{\delta^2}{8}(1+O(N^{-1/4})).
\]
Only gaps with $\delta=O_A(s_N)$ can thus collide before $T_N$.
For those gaps the sharper bound $\delta^2S=O_A(N^{-1/3})$
is available. The collision mark differs from the left endpoint
of its initial gap by at most
\[
 \delta/2+2T_N|L_{\rm pair}|+O(\delta(\delta^2S)^{3/2})=o(1).
\]
Consequently on each bounded rescaled time window the collision
process is, up to a uniformly vanishing displacement, the image
of the initial extreme-gap process under
\[
                  (y,b)\longmapsto(y^2,b),\qquad
                  y=N^{4/3}\delta/(72\pi)^{1/3}.
\]
Ben Arous--Bourgade, Theorem~1.1, printed p.~2650, states
that the process in the unnormalized coordinate
$u=N^{4/3}\delta$ converges to the Poisson process of intensity
$u^2du/(24\pi)\,db/(2\pi)$. Thus the $y$ coordinate has
intensity $3y^2dy\,db/(2\pi)$, and squaring gives
\eqref{eq:analytic-collision-ppp-intensity}.
Its mass on $[0,v]\times\mathbb R/(2\pi\mathbb Z)$ is
$v^{3/2}$. The limit has no atoms at deterministic time or
spatial boundaries. Tightness of the number of extreme gaps
in bounded windows, supplied by that same point-process
convergence, justifies the vanishing-displacement comparison
and all asserted count limits. The formula for $V_k$ follows
by the deterministic time change $v\mapsto v^{3/2}$ to a
rate-one Poisson process. The product intensity yields the
independent uniform marks.
\end{proof}

\paragraph{Primary source and scope.}
The imported random-matrix result is exactly Ben Arous--Bourgade,
\emph{Extreme gaps between eigenvalues of random matrices},
Annals of Probability~41 (2013), 2648--2681,
Theorem~1.1 p.~2650 and Corollary~1.2 p.~2651,
\url{https://math.nyu.edu/~bourgade/papers/ExtremeGaps.pdf},
DOI~10.1214/11-AOP710. The source's full marked intensity,
including $db/(2\pi)$, was checked. The new argument supplies
the heat evolution, exhaustion after the first collision,
and absence of recollisions on bounded extreme-time windows.
A, M and ordinary fixed-scale pair correlation do not supply
an arithmetic version of the imported extreme-gap process.
There is no expectation convergence or error rate for the
limiting Poisson law claimed here.

\paragraph{Geometric and representation meaning.}
The actual finite free family
\[
 \operatorname{Spec}\mathbb C[q,q^{-1},w,w^{-1}]/
       \left(\sum_{k=0}^Nc_kq^{(2k-N)^2}w^k\right)
                           \longrightarrow\mathbb G_{m,q},
                  \qquad q=e^{-a/4},
\]
is the family already identified in NN-A024. On each bounded
extreme-time interval, with probability tending to one, its
nearby analytic zero algebra splits into the tracked singleton
and rank-two cluster factors. Each rank-two factor has at
most one simple trace-discriminant zero there, with a quadratic
local normalization; each singleton remains etale. SGA~1,
Expos\'e~I, section~4, Proposition~10 is the exact finite-locally-free
trace-discriminant criterion, and section~3, Proposition~1 gives
the vanishing-differentials criterion. The resulting simple
branch points obey the proved marked Poisson law. Analytic
monodromy can be taken on a common small tubular parameter
neighborhood of the real interval, preserving the disjoint
cluster factors. Around the separate early branch points it
exchanges disjoint pairs of sheets, so these local transpositions
commute.
This is a concrete random complex algebraic discriminant family,
not an arithmetic Frobenius equidistribution theorem, a pure
motive, or a positive integral basis.


\section{Why an angular bound alone has no converse in finite order}
\label{sec:analytic-angle-only-obstruction}

\begin{theorem}[NN-A026: fixed nonreal angles with arbitrarily high finite positivity]
For every integer $n\ge1$ put
\[
 \vartheta_n=\frac{\pi}{3n},\qquad
 g_n(z)=\prod_{j=0}^{n}
                (1+e^{i\vartheta_n(n-2j)}z),\qquad
 Q(z)=\cosh\sqrt z,\qquad F_n(z)=g_n(z)Q(z).
\]
Then $F_n$ is real entire of order $1/2$, $F_n(0)=1$, and all
its Taylor coefficients are strictly positive. Every allowed
Toeplitz minor of order at most $2n$ is strictly positive.
Nevertheless $F_n$ has the two fixed nonreal zeros
$-e^{i\pi/3}$ and $-e^{-i\pi/3}$, and its maximal zero-angle
deviation from the negative real ray is exactly $\pi/3$.
Its least negative-minor order is finite and exceeds $2n$.

Rescaling to $\widehat F_n(z)=F_n(z/F_n'(0))$ also imposes
$\widehat F_n'(0)=1$ and preserves the same angular and finite
positivity conclusions. Separately, put
$\Psi_n(z)=F_n(z)/F_n(1)$. Then
$U\mapsto\det\Psi_n(U)$ is a normalized central positive-type
function on every $U(d)$ with $d\le2n$, with strictly positive
coefficients for every polynomial irreducible, yet fails to be
positive type at some larger finite rank.

Thus the maximal zero angle, the entire growth order $1/2$,
and the displayed coefficient normalizations alone give no
upper bound on necessary detecting order or unitary rank.
The positive lower bounds in NN-A022 use an angular sufficient
condition; a matching upper bound for the actual heat family
requires additional information about that family.
\end{theorem}

\begin{proof}
Let $m=n+1$, $k=2n$, and more generally consider the centered
geometric alphabet
\[
 x_j=e^{i\vartheta(m+1-2j)},\quad 1\le j\le m,
                \qquad 0<\vartheta\le\frac{\pi}{m+k-1}.
\]
For a partition $\lambda$ with at most $m$ rows, the Weyl
alternant evaluated at this geometric progression gives
\begin{equation}\label{eq:analytic-geometric-schur-sines}
 s_\lambda(x_1,\ldots,x_m)
     =\prod_{1\le i<j\le m}
        \frac{\sin(\vartheta(\lambda_i-\lambda_j+j-i))}
             {\sin(\vartheta(j-i))}.
\end{equation}
For completeness, write the alternant columns as geometric
powers with exponents $\lambda_j+m-j$. Both determinants
become Vandermonde products. Replacing
$e^{iu}-e^{-iu}$ by $2i\sin u$ gives the quotient above;
the overall phases cancel because the alphabet is centered.
The denominator is nonzero, since $0<\vartheta(j-i)<\pi$.
The formula also follows directly from the Weyl character
formula, Etingof et al., Theorem~5.22.1, printed p.~132.

If $\lambda_1\le k$, every numerator argument lies in $[0,\pi]$:
\[
 1\le\lambda_i-\lambda_j+j-i\le k+m-1.
\]
Hence all ordinary Schur values of width at most $k$ are
nonnegative. Values with more than $m$ rows are zero by the
alternating definition. Every skew Schur function of width
at most $k$ is a nonnegative integer linear combination of
ordinary Schur functions of that width: the restriction
multiplicities used and proved in NN-A021 are nonnegative,
and a nonzero coefficient entails containment of the relevant
partition in the outer partition. The exact minor/width
correspondence in NN-A020 now proves that the coefficients
of $\prod_j(1+x_jz)$ are $PF_k$.

Taking $m=n+1,k=2n$ and
$\vartheta=\pi/(m+k-1)=\pi/(3n)$ gives the stated $g_n$.
Its roots are $-x_j^{-1}$; the endpoint phases are exactly
$\pm\pi/3$, with all other phases between them. Its
coefficients are nonnegative by the already proved order-one
minors. Conjugate pairing of the alphabet makes $g_n$ real.

The entire function $Q(z)=\cosh\sqrt z$ has Taylor coefficients
$1/(2r)!$, and the genus-zero product
\[
 Q(z)=\prod_{j\ge1}\left(1+
                       \frac{4z}{\pi^2(2j-1)^2}\right).
\]
Thus $Q$ is $PF_\infty$ and has strictly positive allowed
Toeplitz minors, by the explicit positive-path proof in
NN-A008. Its coefficients are strictly positive. Multiplying
by $g_n$ preserves $PF_{2n}$ by Cauchy--Binet. It actually
preserves strict positivity of every allowed minor: in
$T_QT_{g_n}$, select the intermediate index set equal to the
column index set of the target minor. The selected
$T_{g_n}$ minor is upper triangular with diagonal one,
while the selected $T_Q$ minor is the strictly positive
allowed minor. Every other Cauchy--Binet term is nonnegative.
The sum is finite for each fixed index set, by upper
triangularity. The coefficients of the product are positive
because the constant coefficient of $g_n$ is one.

A polynomial factor changes neither the order $1/2$ of $Q$
nor its negative-real infinite zero set. The fixed nonreal
zeros of $g_n$ remain, so the maximal angle is exactly
$\pi/3$. Edrei's entire-function converse, with the exact
original locator in NN-A007, rules out $PF_\infty$;
therefore some finite minor is negative. None has order
at most $2n$. Positive scaling of the variable multiplies
any Toeplitz minor by a positive number and preserves angles,
which proves the derivative-normalized assertion.

Finally $F_n(1)>0$, so $\Psi_n(0)>0$ and $\Psi_n(1)=1$.
The exact $PF_d$ versus positive-$U(d)$ criterion in NN-A021
proves the asserted characters, including strict positivity
of all polynomial-irreducible coefficients. The same
criterion and Edrei give failure at some finite larger rank.
This character normalization and the derivative normalization
are separate operations; no preservation under a negative
translation of the variable is assumed.
\end{proof}

\paragraph{Classical locator and the new use.}
The finite sharp family is the Schoenberg example recorded in
Adm--Garloff--Tyaglov, \emph{Total nonnegativity of finite
Hurwitz matrices and root location of polynomials},
Theorem~5.7 and equation~(5.9), printed p.~161, immediately before the sector
sufficiency Theorem~5.8, equation~(5.10). The displayed Weyl
calculation supplies a complete proof of precisely the
finite-order assertion needed here. The order-$1/2$ infinite
extension, strict-minor observation, normalizations and
finite-unitary-rank interpretation are given above. This is
an obstruction to an unrestricted converse, not a counterexample
to a heat-specific upper estimate.

\paragraph{An explicit correction to the source's product inference.}
Immediately after equations~(5.11)--(5.12), printed p.~162,
the same published source asserts that a product of quadratics
whose angles are exterior to the quadratic $PF_k$ interval
cannot belong to $PF_k$. That inference is false, even when
every factor has the same outside angle. Set
\[
 q(z)=1+\frac9{10}z+z^2,\qquad
 \theta=\arccos\frac9{20}>\frac\pi3.
\]
Then $q=(1+e^{i\theta}z)(1+e^{-i\theta}z)$ is not $PF_2$,
since its central order-two minor is $81/100-1<0$.
On the other hand
\[
 q(z)^2=1+\frac95z+\frac{281}{100}z^2+\frac95z^3+z^4
\]
is $PF_2$. Its three nontrivial adjacent log-concavity
differences are $43/100$, $46561/10000$, and $43/100$,
all positive. For a positive finite consecutive coefficient
sequence, these differences make the successive ratios
nonincreasing. Therefore every order-two Toeplitz minor
\[
       a_{h+d}a_{h+e}-a_ha_{h+d+e}\quad(d,e>0)
\]
is nonnegative by multiplying the corresponding ratio
inequalities; when an index is outside the support,
upper triangularity and the consecutive support give the
same conclusion directly. This proves the required $PF_2$
assertion without a numerical test. The exact published
source is
\url{https://www-home.htwg-konstanz.de/~garloff/JMAA-3_Adm_Garloff_Tyaglov.pdf},
DOI~10.1016/j.jmaa.2018.06.065. The valid sector sufficiency
and valid single-quadratic criterion used in NN-A020 are
unaffected; they were also proved there independently.

\paragraph{Geometric and representation meaning.}
The values in \eqref{eq:analytic-geometric-schur-sines} are
literal Weyl characters evaluated at a centered diagonal
unitary matrix. In the stated rectangle of partition widths,
they are nonnegative real numbers despite its nonreal
eigenvalues. Their skew restrictions give the positive
Toeplitz tests, and the independent positive $Q$ specialization
makes all allowed tests strict. The representation-ring
exterior-power and exact-sequence formalism has the exact
locator SGA~6, Expos\'e~V, Example~2.2(b) and Lemma~2.2.1.
The functions on each $U(d)$ are actual mixtures of its
polynomial irreducible characters; they fail to extend
positively through all ranks for the fixed analytic function.
Their real coefficients are not integral representation
multiplicities. The example separates finite character-cone
positivity from a constraint depending only on the support
angle of the analytic zero divisor.


\section{A sharp detecting scale when a double endpoint controls the angle}
\label{sec:analytic-dominant-double}

The following conditional case closes a matching upper-bound gate for
the actual heat family. The extra multiplicity hypothesis is displayed;
it is not part of A. We first prove the infinite-alphabet estimate that
makes the negative minor uniform near a coalescing pair.

\begin{lemma}[A coalescing dominant pair in rectangular Schur minors]
\label{lem:analytic-rectangular-pair}
Let an $\ell^1$ complex alphabet, closed under conjugation, depend on
a positive parameter tending to zero. Suppose its entries consist of
\[
 b_1>\cdots>b_p>b,\qquad
 \beta_+=be^{i\theta},\quad\beta_-=be^{-i\theta},\qquad
                         (\eta_j)_{j\ge1},
\]
where $\theta\downarrow0$. Assume $b_i\to b_i^0$,
$b\to b^0>0$, with strictly separated limits
$b_1^0>\cdots>b_p^0>b^0$ when $p>0$ (vacuous for $p=0$), and
$\sup_j|\eta_j|\le c_0<b^0$ uniformly. Assume also
$\sum_j|\eta_j|$ is uniformly bounded, and that the tail products
\[
             C(z)=\prod_j(1-\eta_jz)^{-1}
\]
converge locally uniformly on $|z|<1/c_0$ to the corresponding product of an endpoint
alphabet of positive real entries less than $b^0$. Write $k=p+1$,
$P=\prod_{i=1}^p b_i$, and $h_r$ for the complete symmetric functions
of the full alphabet. Then there are constants $C_1>0$, $0<q<1$,
independent of $n$ and of the sufficiently small parameter, such that
\begin{equation}\label{eq:analytic-rectangular-asymptotic}
 \frac{s_{(n^k)}}{(Pb)^n}
  =\frac{|D|\sin((n+1)\theta+\arg D)}{\sin\theta}
                         +O(C_1\theta^{-k}q^n),\qquad n\ge k,
\end{equation}
where
\[
 D=\left\{\prod_{i=1}^p
              \frac{C(1/b_i)}{1-\beta_-/b_i}\right\}C(1/\beta_+)
                             \longrightarrow D_0>0.
\]
In particular if $n\theta\to u\in(\pi,2\pi)$, the displayed
rectangular Schur value is strictly negative eventually.
The same value is the elementary Toeplitz minor
\begin{equation}\label{eq:analytic-rectangular-elementary-minor}
                  s_{(n^k)}=det[e_{k-i+j}]_{i,j=1}^n.
\end{equation}
It has order $n$ and fits in the index window $0,\ldots,k+n-1$.
\end{lemma}

\begin{proof}
Choose $c$ strictly between $c_0$ and $b^0$, and smaller than
all the limiting $b_i$. For sufficiently small parameter, the generating
function
\[
 H(z)=\sum_{r\ge0}h_rz^r
       =\prod_{i=1}^p(1-b_i z)^{-1}
           (1-\beta_+z)^{-1}(1-\beta_-z)^{-1} C(z)
\]
has exactly these $p+2$ simple poles inside $|z|=1/c$.
Its modulus on that circle is uniformly bounded. Remove the simple
partial fractions there:
\[
 H(z)=\sum_{\alpha\in\mathcal B}
          \frac{A_\alpha}{1-\alpha z}+G(z),\qquad
 \mathcal B=\{b_1,\ldots,b_p,\beta_+,\beta_-\},
\]
where
\[
 A_\alpha=\prod_{\gamma\ne\alpha}(1-\gamma/\alpha)^{-1}.
\]
The product includes the whole tail. The tail product converges
absolutely and is bounded above and away from zero at the required
points, by its uniform $\ell^1$ bound and strict modulus separation.
The only vanishing denominator is
$1-\beta_-/\beta_+$ or its conjugate. Hence
$|A_\alpha|\le C\theta^{-1}$, and Cauchy's formula gives
\begin{equation}\label{eq:analytic-h-poles-remainder}
 h_r=\sum_{\alpha\in\mathcal B}A_\alpha\alpha^r+g_r,
                         \qquad |g_r|\le C\theta^{-1}c^r.
\end{equation}
The remainder $G$ is analytic through all removed poles. Its circle
bound follows from those of $H$ and the partial fractions, so no
uniform separation of the two coalescing poles was assumed.

Use the ordinary Jacobi--Trudi determinant
$s_{(n^k)}=\det[h_{n-i+j}]_{i,j=0}^{k-1}$.
The pole part of its matrix is a sum of rank-one matrices
\[
        A_\alpha\alpha^n
              (\alpha^{-i})_{i=0}^{k-1}
              (\alpha^{j})_{j=0}^{k-1}.
\]
Cauchy--Binet for that finite sum gives one term for each
$k$-element subset $J\subset\mathcal B$. The two Vandermonde
determinants cancel all internal factors in
$\prod_{\alpha\in J}A_\alpha$, leaving exactly
\begin{equation}\label{eq:analytic-rectangular-subset-coefficient}
 \left(\prod_{\alpha\in J}\alpha\right)^n
       \prod_{\substack{\alpha\in J\\\gamma\notin J}}
                                      (1-\gamma/\alpha)^{-1}.
\end{equation}
Here the complement includes the infinite tail, only through its
absolutely convergent products. To see the cancellation directly,
each pair $\alpha,\alpha'$ within $J$ contributes
$-(\alpha\alpha')/(\alpha-\alpha')^2$ from the two partial
fractions, the inverse of its contribution from the two Vandermondes.

The two subsets containing every $b_i$ and exactly one of
$\beta_\pm$ have modulus weight $(Pb)^n$. Their coefficients are
$D/(1-e^{-2i\theta})$ and its conjugate. Since
\[
             (1-e^{-2i\theta})^{-1}
                         =e^{i\theta}/(2i\sin\theta),
\]
their sum is precisely the first term of
\eqref{eq:analytic-rectangular-asymptotic}. The endpoint tail is
positive real, all $b_i^0>b^0$, and the displayed products therefore
give $D\to D_0>0$.

Every other $k$-element subset omits at least one $b_i$ and
contains both members of the pair. Its modulus weight divided by
$Pb$ is $b/b_i$, bounded by a constant less than one. Its coefficient
is uniformly bounded: the coalescing pair is internal to $J$, so its
potential singular factors already cancelled in
\eqref{eq:analytic-rectangular-subset-coefficient}.

It remains to control determinants that use remainder columns from
\eqref{eq:analytic-h-poles-remainder}. This cannot be done by a
naive bound using the largest pole to the $k$th power. Instead use
multilinearity. A term with $m\ge1$ remainder columns has $k-m$
remaining rank-one pole columns, whose pole labels must be distinct
or the determinant is zero. Its modulus is at most a fixed constant
 times
\[
 \theta^{-k}\left(c^m\prod_{\alpha\in J}|\alpha|\right)^n,
                              \qquad |J|=k-m.
\]
Fixed shifts $j-i$ affect only the constant. The largest product of
$k-m\le p$ pole moduli uses that many of the $b_i$. Because each
remaining $b_i$ is larger than $b$,
\[
          \frac{c^m\prod_{\alpha\in J}|\alpha|}{Pb}
                                      \le(c/b)^m\le q_1<1.
\]
There are finitely many such multilinear terms, their number depending
only on $k$. Combining this bound with the nonleading subsets proves
\eqref{eq:analytic-rectangular-asymptotic}, after increasing $q<1$.
If $n\theta\to u\in(\pi,2\pi)$, its first term is asymptotic to
$D_0\sin u/\theta<0$, and its exponential remainder is smaller than
any power of $\theta$. Finally dual Jacobi--Trudi transposes
$(n^k)$ to $(k^n)$ and gives
\eqref{eq:analytic-rectangular-elementary-minor}. Explicitly use
row indices $0,\ldots,n-1$ and column indices $k,\ldots,k+n-1$
in the upper-triangular Toeplitz matrix $(e_{j-i})$.
\end{proof}

\begin{theorem}[NN-A027: exact detecting constants in the dominant-double case]
\label{thm:analytic-dominant-double}
Assume A. For a positive zero $x$ of $H_0$, let $m_x$ be its
multiplicity and $h_{m_x}$ the largest absolute root of the
probabilists' Hermite polynomial $\operatorname{He}_{m_x}$,
with $h_1=0$. Make the following additional hypothesis:
there is a positive \emph{double} zero $x_0$ such that
\begin{equation}\label{eq:analytic-dominant-double-hypothesis}
                    \frac{h_{m_x}}x\le\frac1{x_0}
                    \qquad\hbox{for every positive zero }x.
\end{equation}
This is a precise conditional case; neither this multiplicity pattern
nor existence of any multiple zero is inferred from A.

Let $m_*(-a)$ be the least order of a negative Toeplitz minor of
$F_{-a}$, and $R_*(-a)$ its least detecting index window as in
NN-A014 and NN-A020. Then
\begin{equation}\label{eq:analytic-sharp-order-window}
       \lim_{a\downarrow0}\sqrt a\,m_*(-a)
        =\lim_{a\downarrow0}\sqrt a\,R_*(-a)
                       =\frac{\pi x_0}{2\sqrt2}.
\end{equation}
For the normalized unitary characters of NN-A021, put
\[
 p_0=F_0'(0)=\sum_{x>0}x^{-2},\qquad
 \beta_0=\frac1{p_0x_0^2}\in(0,1),
\]
and let $d_*(-a)$ be the least rank at which
$U\mapsto\det E_{-a}(U-I)$ fails to be positive type. Then
\begin{equation}\label{eq:analytic-sharp-unitary-rank}
       \lim_{a\downarrow0}\sqrt a\,d_*(-a)
                       =\frac{\pi x_0(1-\beta_0)}{2\sqrt2}.
\end{equation}
There are explicit negative rectangular witnesses: if $p$ is the
number of positive endpoint zeros below $x_0$ and $k=p+1$, take
$n\theta(a)\to u\in(\pi,2\pi)$, where $\theta(a)$ is the positive
argument of the corresponding complex product parameter. The minor
$\det[e_{k-i+j}]_{i,j=1}^n$ is negative. For the translated unitary
alphabet this is the negative Fourier coefficient of the actual
one-dimensional representation $\det(U)^k$ on $U(n)$.
\end{theorem}

\begin{proof}
For $m\ge2$, $h_m\ge h_2=1$, by the principal Jacobi-matrix
interlacing of the Hermite recurrence. Therefore
\eqref{eq:analytic-dominant-double-hypothesis} forces every endpoint
zero below $x_0$ to be simple. There are finitely many of them,
$x_1<\cdots<x_p<x_0$. Their continuations at time $-a$ remain
simple and real for small $a$. The double zero at $x_0$ splits by
NN-A006 into
\[
             z_\pm(a)=x_0\pm i\sqrt{2a}+O(a).
\]
The two parameters $z_\pm(a)^{-2}$ are conjugate and have arguments
$\pm\theta(a)$ with
\begin{equation}\label{eq:analytic-pair-theta-sharp}
                         \frac{\theta(a)}{\sqrt a}
                                      \longrightarrow\frac{2\sqrt2}{x_0}.
\end{equation}
The $p$ real parameters from lower zeros have strictly larger moduli.
All remaining parameters have uniformly smaller modulus than
$x_0^{-2}$: choose a circle between $x_0$ and the next larger positive
endpoint zero, and use compact zero convergence. Zeros outside a
large fixed circle contribute parameters tending uniformly to zero.

Here and below the infinite products and their uniform bounds are
legitimate without simplicity of the other zeros. For $t$ in a compact
interval the theta integral gives
$\log\max_{|z|\le R}|H_t(z)|\le C R\log(R+3)$, with $C$ uniform,
and $H_t(0)$ is bounded away from zero near $t=0$. Jensen's formula
therefore gives a uniform zero count $N_t(R)\le C R\log(R+3)$.
Partial summation bounds uniformly the tails of
$\sum|z_j(t)|^{-2}$ and makes them tend to zero as the radius tends
to infinity. The genus-zero products of $F_t$, after removing the
fixed finite cluster, consequently converge locally uniformly to the
endpoint tail product. In particular the hypotheses of
Lemma~\ref{lem:analytic-rectangular-pair} hold for the alphabet
$z_j(-a)^{-2}$, with $k=p+1$.

Take $n=\lceil(\pi+\varepsilon)/\theta(a)\rceil$ for any fixed
$0<\varepsilon<\pi$. The lemma gives a negative minor of order $n$
in window $k+n-1$ for all sufficiently small $a$. Hence
\[
 m_*(-a)\le n,\qquad R_*(-a)\le k+n-1,
\]
and \eqref{eq:analytic-pair-theta-sharp} gives the corresponding
upper limits with $\pi+\varepsilon$ in place of $\pi$.

For the lower limits, NN-A022 and the extra hypothesis give
\[
 \frac{\alpha_H(-a)}{\sqrt a}\longrightarrow\frac{\sqrt2}{x_0}.
\]
The maximal parameter angle is exactly $2\alpha_H(-a)$ when these
angles are small. The strict quadratic sector proof in NN-A020
shows that every allowed minor of order at most $r$ is positive
whenever $(r+1)2\alpha_H(-a)<\pi$. Thus for any fixed
$\varepsilon>0$ one may take
$r=\lfloor(\pi-\varepsilon)/(2\alpha_H(-a))\rfloor-1$.
It follows that
\[
 \liminf_{a\downarrow0}\sqrt a\,m_*(-a)
                           \ge\frac{\pi x_0}{2\sqrt2}.
\]
Any minor in the window $0,\ldots,R$ has order at most $R+1$,
so $R_*(-a)\ge m_*(-a)-1$. Letting $\varepsilon\downarrow0$
in the upper bounds proves \eqref{eq:analytic-sharp-order-window}.

For the unitary statement put $p_a=F_{-a}'(0)>0$. The translated
alphabet in NN-A021 is
\[
 \widetilde\beta_j(a)
   =\frac{\beta_j(a)}{1-\beta_j(a)}
   =\frac1{p_a z_j(-a)^2-1},\qquad
                    \beta_j(a)=\frac1{p_a z_j(-a)^2}.
\]
At the endpoint all its entries are positive and are strictly decreasing
as functions of the positive zero $x$. Also $\sup_j|\beta_j(a)|<b<1$
for small $a$, by the compact zero separation already proved in
NN-A021. Its $\ell^1$ tail and local convergence therefore satisfy
the same lemma. At an endpoint zero $x$ of multiplicity $m_x$,
the largest local angle coefficient becomes
\begin{equation}\label{eq:analytic-translated-angle-coefficient}
 \frac{2\sqrt2\,h_{m_x}}{x(1-1/(p_0x^2))}.
\end{equation}
This follows by substituting the local Hermite expansion into
$(p_a z^2-1)^{-1}$; $p_a=p_0+O(a)$ is real.

The maximum of these coefficients is attained at $x_0$.
Indeed the nonzero coefficients occur only at $x\ge x_0$,
$h_{m_x}/x\le1/x_0$, and
$x\mapsto(1-1/(p_0x^2))^{-1}$ decreases there. The uniform
radial-tail estimate in NN-A022 also applies after this transformation:
for high zeros $|\beta_j(a)|\le C/|z_j(-a)|^2$, while
\[
 \left|\arg\frac{\beta}{1-\beta}\right|
                         \le\frac{|\arg\beta|}{1-|\beta|}
                         \qquad(|\beta|<1).
\]
One proves this last bound by integrating
$\operatorname{Re}(1-re^{i\theta})^{-1}\le(1-r)^{-1}$
along the argument, starting at a positive real number.
Consequently the maximal translated angle, divided by $\sqrt a$,
has the exact limit
$2\sqrt2/[x_0(1-\beta_0)]$.

The sector lower bound and the same rectangular upper bound now give
\eqref{eq:analytic-sharp-unitary-rank}, using the exact $PF_d$
versus $U(d)$ equivalence of NN-A021. Multiplication by the positive
constant $\Psi_a(0)=E_{-a}(-1)$ does not affect minor signs.
For the upper witness, the coefficient formula from that equivalence
is
\[
 c_{(k^n),n}=\det[a_{k-i+j}]_{i,j=1}^n
             =\Psi_a(0)^n s_{(n^k)}(\widetilde\beta(a))<0.
\]
The highest weight $(k,\ldots,k)$ is precisely $\det(U)^k$,
so this is the claimed concrete irreducible witness.
\end{proof}

\paragraph{Geometric and representation meaning.}
The additional hypothesis states that a local Artin algebra of length
two controls the maximal scaled angular motion of the endpoint divisor.
Its quadratic normalization produces the two conjugate dominant
parameters. The rectangular Schur computation cancels the internal
Vandermonde factors and isolates their sine oscillation with a uniformly
controlled infinite complementary divisor. Dual Jacobi--Trudi turns
this into a finite minor whose order and index window have the exact
asymptotic above. After the actual unitary translation, that same
rectangle is a negative coefficient of $\det^k$ at rank $n$.
The exterior-power/representation-ring dictionary is SGA~6, Expos\'e~V,
Example~2.2(b) and Lemma~2.2.1; the rank-two discriminant criterion is
SGA~1, Expos\'e~I, section~4, Proposition~10. These are sharp bounds
for the specified analytic and representation certificate families
under the displayed multiplicity hypothesis. They are not a proof
of that hypothesis, a pure motive, or a lower bound against all
algorithms.


\section{Endpoint simplicity and the sharp order of finite detection}
\label{sec:analytic-multiplicity-detection-equivalence}

The dominant-double hypothesis in NN-A027 is needed for its displayed
\emph{exact constants}, but not for a matching order of growth when
some endpoint is multiple. We prove this using an actual compact
unitary orbit and an explicit Schur isolation argument.

\begin{lemma}[A negative Hermitian orbital Fourier value]
\label{lem:analytic-negative-orbital}
Let $D\ne0$ be a Hermitian $m$ by $m$ matrix with trace zero and
spectrum invariant under multiplication by $-1$. Then
\[
 \phi_D(X)=\int_{U(m)}e^{i\operatorname{Tr}(DUXU^*)}\,dU,
                           \qquad X=X^*,
\]
is real and takes a strictly negative value. One may choose a
negative value at a diagonal matrix with distinct positive eigenvalues.
\end{lemma}

\begin{proof}
The orbit of $D$ equals that of $-D$. Haar invariance therefore makes
$\phi_D$ real. It is continuous, conjugation-invariant, and
$\phi_D(0)=1$. Give the real vector space $\mathcal H_m$ of Hermitian
matrices the Euclidean inner product $\operatorname{Tr}(XY)$ and
Lebesgue measure for an orthonormal basis; its dimension is $d=m^2$.
Fubini and the elementary Gaussian Fourier integral give, for every
$t>0$,
\begin{equation}\label{eq:analytic-orbital-Gaussian}
 \int_{\mathcal H_m}e^{-t\|X\|^2/2}\phi_D(X)\,dX
       =\left(\frac{2\pi}{t}\right)^{d/2}
                         e^{-\|D\|^2/(2t)}.
\end{equation}
The right side tends to zero as $t\downarrow0$, because $D\ne0$.
If $\phi_D$ were nonnegative everywhere, a fixed ball about zero on
which $\phi_D\ge1/2$ would give a positive lower limit for the left
side, a contradiction. Thus some $X$ has $\phi_D(X)<0$.
Diagonalize it, perturb its eigenvalues to be distinct while preserving
strict negativity, and add a scalar matrix large enough to make every
eigenvalue positive. This last operation does not change $\phi_D$
because $\operatorname{Tr}D=0$.
\end{proof}

\paragraph{Exact compact integral locator.}
We use the classical Harish--Chandra--Itzykson--Zuber formula with Haar
probability measure. If $A=\operatorname{diag}(a_1,\ldots,a_m)$ and
$B=\operatorname{diag}(b_1,\ldots,b_m)$ have distinct eigenvalues, then
\begin{equation}\label{eq:analytic-HCIZ}
 \mathcal I(A,B):=\int_{U(m)}e^{\operatorname{Tr}(AUBU^*)}\,dU
       =\left(\prod_{j=1}^{m-1}j!\right)
                    \frac{\det[e^{a_i b_j}]}{\Delta(a)\Delta(b)}.
\end{equation}
It extends to complex diagonal $A$ and coincident eigenvalues by
holomorphic continuation; the integral is entire in the matrix entries
on compact sets. A precise primary locator is J.-B.~Zuber,
\emph{Horn's problem and Harish-Chandra's integrals. Probability
density functions}, Ann. Inst. Henri Poincar\'e Comb. Phys. Interact.~5
(2018), 309--338, orbital definition (6), p.~311, formula (11), p.~312,
and constant (12), p.~313, DOI~10.4171/AIHPD/56,
\url{https://www.numdam.org/item/10.4171/aihpd/56.pdf}.
Only this exact classical identity is imported; the sign and isolation
arguments needed below are proved here.

\begin{lemma}[A whole coalescing cluster gives a bounded scaled-width witness]
\label{lem:analytic-cluster-Schur-witness}
Let an $\ell^1$ alphabet, closed under conjugation, have the following
form as $\varepsilon\downarrow0$:
\[
 b_1(\varepsilon)>\cdots>b_p(\varepsilon)>0,\quad
 \xi_j(\varepsilon)=b(\varepsilon)
                  e^{\varepsilon a_j(\varepsilon)}\quad(1\le j\le m),
                       \quad (\eta_j(\varepsilon))_{j\ge1}.
\]
Suppose $b>0$, $b\to b^0>0$,
$b_1^0>\cdots>b_p^0>b^0$, and
\[
 a_j(\varepsilon)=i d_j+O(\varepsilon),\qquad
 \sum_j a_j(\varepsilon)=0,
\]
where the $d_j$ are distinct real numbers, not all zero, with a
sign-symmetric multiset. Suppose the tail has uniformly bounded
$\ell^1$ norm, $\sup_j|\eta_j|\le c_0<b^0$, and its product
$C_\varepsilon(z)=\prod_j(1-\eta_jz)^{-1}$ converges locally uniformly
on $|z|<1/c_0$ to a product of positive real endpoint parameters.
The assertions about $p$ are vacuous when $p=0$.

Then there is a constant $C>0$ and a family of partitions
$\lambda(\varepsilon)$ of fixed length at most $q=p+m$, with
$\lambda_1\le C/\varepsilon+1$, such that
\[
                        s_{\lambda(\varepsilon)}<0
\]
for all sufficiently small $\varepsilon$. This is a negative
Toeplitz minor of order at most $C/\varepsilon+1$ in an index
window of size at most $C/\varepsilon+q$.
\end{lemma}

\begin{proof}
Put $D=\operatorname{diag}(d_1,\ldots,d_m)$. The preceding lemma
provides distinct numbers
$L_1>\cdots>L_m>0$ with
$\phi_D(\operatorname{diag}(L))<0$. Choose $C>L_1$ and $\eta>0$
so small that $L_m>2\eta$ and $C-L_1>2\eta$. Put
\[
 \mu_i=\lfloor L_i/\varepsilon\rfloor,\qquad
 n=\lceil C/\varepsilon\rceil,\qquad
 \lambda=(\underbrace{n,\ldots,n}_{p\text{ times}},
                                  \mu_1,\ldots,\mu_m).
\]
These are partitions for small $\varepsilon$, with
$\mu_m\ge\eta/\varepsilon$ and
$n-\mu_1\ge\eta/\varepsilon$. All row lengths are $O(1/\varepsilon)$.

First consider just the cluster, divided by its positive common scale
$b$. Write $A_\varepsilon=\operatorname{diag}(a_j(\varepsilon))$,
$\rho=(m-1,m-2,\ldots,0)$,
$U_\varepsilon=\operatorname{diag}(\varepsilon(\mu+\rho))$ and
$V_\varepsilon=\operatorname{diag}(\varepsilon\rho)$.
The Weyl alternant and \eqref{eq:analytic-HCIZ} give the \emph{exact}
identity
\begin{equation}\label{eq:analytic-Schur-orbital-limit}
 \frac{s_\mu(e^{\varepsilon a_1},\ldots,e^{\varepsilon a_m})}
      {s_\mu(1,\ldots,1)}
   =\frac{\mathcal I(A_\varepsilon,U_\varepsilon)}
          {\mathcal I(A_\varepsilon,V_\varepsilon)}
      \longrightarrow\phi_D(\operatorname{diag}(L))<0.
\end{equation}
Indeed the Vandermonde quotient
$\Delta(\mu+\rho)/\Delta(\rho)$ is the positive Weyl dimension,
and the remaining determinants give the displayed ratio. Its denominator
tends to one. Thus the cluster Schur value is negative, of magnitude
at least $c\,b^{|\mu|}s_\mu(1^m)$ for a fixed $c>0$.
This conclusion is stable under the full $O(\varepsilon)$ perturbation
of $A_\varepsilon$; no exact geometric spacing of the cluster is used.

We now show that the finite positive prefix and infinite tail change
this sign only by a positive factor and an exponentially small error.
All constants below may depend on the fixed $p,m,L,C$ and the stated
spectral separations. Choose $c_1$ with $c_0<c_1<b^0$.
The complete-symmetric generating function has exactly $q=p+m$
simple top poles $\alpha\in\mathcal A=\{b_i,\xi_j\}$ inside
$|z|=1/c_1$. The cluster pole differences are bounded below by a
constant times $\varepsilon$, since the $d_j$ are distinct. The
partial-fraction argument of NN-A027 therefore gives
\begin{equation}\label{eq:analytic-full-cluster-poles}
 h_r=\sum_{\alpha\in\mathcal A}A_\alpha\alpha^r+g_r,
 \qquad |A_\alpha|\le C_2\varepsilon^{-(m-1)},\qquad
 |g_r|\le C_2\varepsilon^{-(m-1)}c_1^r.
\end{equation}
The analytic remainder bound follows on the same fixed circle after
removing all top poles. Uniform tail product bounds ensure its constant
is independent of $r$ and $\varepsilon$.

Apply \eqref{eq:analytic-full-cluster-poles} to the ordinary
Jacobi--Trudi determinant $\det[h_{\lambda_i-i+j}]_{i,j=1}^q$.
The part using all $q$ top poles is exactly
\begin{equation}\label{eq:analytic-full-cluster-leading}
                  s_\lambda(\mathcal A)
                              \prod_{\alpha\in\mathcal A}C_\varepsilon(1/\alpha).
\end{equation}
This is the same internal Vandermonde cancellation as in NN-A027,
now with a general partition instead of a rectangle. The product on
the right is real and tends to a positive number.

Here is a uniform error estimate for all terms with remainder columns.
Let $P=\prod_{i=1}^p b_i$ and $W=P^n b^{|\mu|}>0$.
A term with $r\ge1$ remainder columns can use only $q-r$ distinct
top pole labels; repeated labels give zero. After determinant expansion,
its exponential weight assigns the corresponding moduli, and $r$ copies
of $c_1$, to the nonincreasing row exponents $\lambda_i$.
The rearrangement inequality maximizes that weight by matching their
orders. Compared with the full set of top moduli, replacing any one
by $c_1$ loses a factor at most
$\widetilde q^{\mu_m}$ for a fixed $\widetilde q<1$:
all exponents are at least $\mu_m$, and every top modulus is uniformly
larger than $c_1$. More explicitly the sorted logarithms decrease
coordinatewise, their total decrease is at least
$\log(\min|\alpha|/c_1)$, and multiplication by exponents at least
$\mu_m$ gives this loss. The full top weight is bounded by a constant
times $W$, because $|\xi_j|/b=e^{O(\varepsilon^2)}$ and
$|\lambda|=O(1/\varepsilon)$. Fixed shifts $j-i$ in the matrix entries
affect only a constant. There are finitely many terms, and the residue
bounds in \eqref{eq:analytic-full-cluster-poles} introduce only a fixed
power of $\varepsilon^{-1}$. Thus the total remainder is
\begin{equation}\label{eq:analytic-tail-exponential-error}
                  O\bigl(W\varepsilon^{-K}
                              \widetilde q^{\mu_m}\bigr)
\end{equation}
for some fixed $K$. This explicitly avoids a determinant estimate that
would repeat the largest pole in all rows.

To isolate the prefix inside $s_\lambda(\mathcal A)$, expand its Weyl
numerator along the first $p$ rows. The leading choice assigns the
$p$ positive prefix columns to those rows. Dividing by the full
Vandermonde gives exactly
\begin{equation}\label{eq:analytic-prefix-factorization}
 P^n s_\mu(\xi_1,\ldots,\xi_m)
            \prod_{i=1}^p\prod_{j=1}^m(1-\xi_j/b_i)^{-1}.
\end{equation}
For example the prefix numerator block is
$P^{n+m}\Delta(b_1,\ldots,b_p)$, and the cluster block is
$s_\mu(\xi)\Delta(\xi)$; these cancel the corresponding full
Vandermonde factors. The remaining product is real and tends to a
positive number.

Every other prefix assignment puts at least one larger $b_i$ into a
cluster row, whose length is at most $\mu_1$, and a cluster modulus
into a row of length $n$. It therefore loses a factor at most
$q_2^{n-\mu_1}$, with $q_2<1$ fixed. Pair the misplaced prefix and
cluster columns; the ratio $|\xi|/b_i$ is uniformly less than one,
and apply the rearrangement inequality to any other columns.
The cluster Vandermonde in the denominator has magnitude bounded
below by a constant times $\varepsilon^{m(m-1)/2}$, while all other
Vandermonde factors stay away from zero. Consequently the total wrong
assignment contribution is
\begin{equation}\label{eq:analytic-prefix-exponential-error}
                   O\bigl(W\varepsilon^{-K'}q_2^{n-\mu_1}\bigr).
\end{equation}
For $p=0$ this error is absent. Combining
\eqref{eq:analytic-full-cluster-leading}--\eqref{eq:analytic-prefix-exponential-error},
the full Schur value is a positive factor, bounded away from zero,
times $P^n s_\mu(\xi)$, plus $W$ times an exponentially small quantity
in $1/\varepsilon$. Equation~\eqref{eq:analytic-Schur-orbital-limit}
therefore makes it negative eventually.

Finally put $r=\lambda_1\le n$; equality need not hold when $p=0$.
Dual Jacobi--Trudi gives a minor of order $r$:
\[
                     s_\lambda=\det[e_{\lambda'_i-i+j}]_{i,j=1}^r.
\]
Transpose this determinant. Use row indices $r-i$ and column indices
$\lambda'_j+r-j$, both in decreasing order, and then reverse both.
The resulting increasing index sets lie in $0,\ldots,r+q-1$,
which is contained in $0,\ldots,n+q-1$.
This proves the stated order and window bounds.
\end{proof}

\begin{theorem}[NN-A028: a finite-detection characterization of endpoint simplicity]
\label{thm:analytic-simplicity-detection-equivalence}
Assume A. If $H_0$ has any multiple zero, then there are positive
constants $c,C,a_0$ such that for $0<a<a_0$,
\begin{equation}\label{eq:analytic-any-multiple-detection-scale}
 c a^{-1/2}\le m_*(-a),\ R_*(-a),\ d_*(-a)\le C a^{-1/2}.
\end{equation}
Here the three quantities are respectively the least negative-minor
order, least detecting index window for $F_{-a}$, and least failing
rank of the actual translated unitary characters from NN-A021.
The chained notation means that \emph{each} of the three quantities
lies between the displayed bounds.

Consequently the following four statements are equivalent:
\begin{enumerate}
\item Every nontrivial zero of $\zeta$ is simple.
\item $\sqrt a\,m_*(-a)\longrightarrow+\infty$ as $a\downarrow0$.
\item $\sqrt a\,R_*(-a)\longrightarrow+\infty$ as $a\downarrow0$.
\item $\sqrt a\,d_*(-a)\longrightarrow+\infty$ as $a\downarrow0$.
\end{enumerate}
The equivalence proves none of these equivalent assertions.
When a multiple endpoint exists, the upper witnesses have partition
height bounded by a fixed number depending on its first occurrence;
only their width and corresponding unitary rank must grow.
NN-A027 supplies sharper exact constants in its dominant-double case.
\end{theorem}

\begin{proof}
Suppose a multiple positive endpoint exists and choose the smallest
one, $x_0$, of multiplicity $m\ge2$. Such a first one exists because
the endpoint divisor is discrete and bounded away from zero. All
$p$ positive endpoint zeros below it are simple, so their inverse-square
parameters are a strictly separated positive prefix for sufficiently
small negative time. The local Hermite expansion from NN-A006 gives
\[
             z_j(-a)=x_0+i\sqrt{2a}\,h_{m,j}+O(a),
                             \qquad 1\le j\le m.
\]
These local roots are distinct for $a>0$ small. Their inverse squares
form a conjugation-stable cluster. Take its positive geometric mean
$b(a)=(\prod_j z_j(-a)^{-2})^{1/m}$ and put $\varepsilon=\sqrt a$.
Branches of the logarithm near the positive endpoint give
\[
 z_j(-a)^{-2}=b(a)e^{\varepsilon a_j(\varepsilon)},\qquad
 a_j(\varepsilon)=-\frac{2i\sqrt2}{x_0}h_{m,j}+O(\varepsilon),
                         \qquad\sum_j a_j(\varepsilon)=0.
\]
The limiting Hermite diagonal is nonzero, traceless, sign-symmetric
and has distinct entries. The uniform zero-count and tail argument
in NN-A027 gives the required smaller-modulus infinite tail and its
locally uniform positive endpoint product. Therefore
Lemma~\ref{lem:analytic-cluster-Schur-witness} supplies a negative
minor of order at most $C/\sqrt a+1$ and window at most
$C/\sqrt a+p+m$. This proves the upper bounds for $m_*$ and $R_*$.

For the actual unitary character use, as in NN-A027,
\[
       \widetilde\beta_j(a)=\frac1{p_a z_j(-a)^2-1},
                             \qquad p_a=F_{-a}'(0)>0.
\]
The positive endpoint ordering is unchanged. The cluster has the
same Hermite diagonal multiplied by the positive factor
$(1-1/(p_0x_0^2))^{-1}$. Its tail is uniformly $\ell^1$ and separated,
since the normalized original parameters have modulus bounded by
some fixed number less than one near $a=0$. The same lemma thus
gives a negative Toeplitz minor of order at most $C'/\sqrt a+1$
for the translated coefficient sequence. The exact finite-rank
criterion in NN-A021 turns it into failure of positive type at that
rank. The dual partition specifies the actual negative irreducible
coefficient; its entries are bounded by the fixed height $p+m$.

The positive lower bounds of inverse-square-root order are already
proved in NN-A022, including its explicit order/window bound and
its effective unitary-rank bound. Decreasing $c$ and $a_0$ absorbs
the fixed additive constants, proving
\eqref{eq:analytic-any-multiple-detection-scale}.

If every endpoint is simple, NN-A022 instead proves
$\sqrt a\,m_*(-a)\to\infty$ and
$\sqrt a\,d_*(-a)\to\infty$; the elementary inequality
$R_*\ge m_*-1$ gives the window statement. Conversely, the existence
of a multiple endpoint gives the bounded upper quantities just proved,
so none of these three limits can be infinite. A identifies the
positive endpoint zeros with the positive nontrivial zeta zeros,
including their multiplicities. This proves the four-way equivalence.
Rodgers--Tao and A ensure the least witnesses are finite at every
negative time even in the all-simple case; they do not decide which
case occurs.
\end{proof}

\paragraph{Geometric and representation meaning.}
The first nonreduced endpoint local algebra has finite length $m$.
Its universal leading heat unfolding is the Hermite cluster, whose
nonzero traceless Hermitian diagonal has the genuine compact orbit
$U(m)/U(1)^m$. This orbit is a full flag manifold with its usual
coadjoint-orbit symplectic form. The Fourier transform of its Haar
orbital measure is forced to become negative by
\eqref{eq:analytic-orbital-Gaussian}. The exact Weyl--HCIZ identity
then produces negative values of actual integral highest-weight
Schur characters. The finite positive prefix and infinite analytic
tail have been isolated with explicit uniform exponentially small
errors, so the same sign occurs in the actual heat coefficient minor.
The exterior-power and representation-ring operations have the exact
locator SGA~6, Expos\'e~V, Example~2.2(b) and Lemma~2.2.1.
No integrality or prequantization of the Hermite orbit's real symplectic
class is asserted. The result is an exact relation between reducedness
of the endpoint analytic divisor and growth of three concrete finite
certificate families; it is not a proof of reducedness, an arithmetic
purity realization, or a complexity-class separation.


\section{Exact detecting constants for every endpoint multiplicity pattern}
\label{sec:analytic-orbital-detection-constants}

\begin{definition}[The Hermite orbital positivity radius]
For $m\ge2$ let
\[
 D_m=\operatorname{diag}(h_{m,1},\ldots,h_{m,m}),
\]
where the entries are all roots of the probabilists' Hermite polynomial
$\operatorname{He}_m$. Define
\[
 \phi_m(X)=\int_{U(m)}e^{i\operatorname{Tr}(D_mUXU^*)}\,dU,
                 \qquad X=X^*,
\]
and
\begin{equation}\label{eq:analytic-orbital-radius}
 \mathfrak r_m=
   \inf\{\lambda_{\max}(X)-\lambda_{\min}(X):\phi_m(X)<0\}.
\end{equation}
Put $\mathfrak r_1=+\infty$ and $h_m=\max_j|h_{m,j}|$.
The orbital measure is the concrete Haar measure on the compact
Hermite orbit, not a hypothesized arithmetic measure.
\end{definition}

\begin{lemma}[Strictness, bounds and the first two radii]
\label{lem:analytic-orbital-radius-properties}
For $m\ge2$,
\[
                  0<\frac\pi{h_m}\le\mathfrak r_m<\infty.
\]
For each $c<\mathfrak r_m$, $\phi_m$ has a strictly positive minimum
on the compact set of diagonal matrices with entries in $[0,c]$.
Moreover
\begin{equation}\label{eq:analytic-first-orbital-radii}
                   \mathfrak r_2=\pi,\qquad
                   \mathfrak r_3=\frac{2\pi}{\sqrt3}.
\end{equation}
\end{lemma}

\begin{proof}
Finiteness follows from Lemma~\ref{lem:analytic-negative-orbital}.
The trace of $D_m$ is zero, so scalar translation of $X$ does not
change $\phi_m$. For small spectral diameter, subtracting a scalar
makes $\|X\|$ uniformly small. Since $\phi_m(0)=1$, this also proves
that the radius is positive.

We first rule out interior zero values in a nonnegative region.
On the Euclidean space of Hermitian matrices, differentiating under
the compact integral gives the exact elliptic equation
\begin{equation}\label{eq:analytic-orbital-Helmholtz}
                         \Delta\phi_m=-\|D_m\|^2\phi_m.
\end{equation}
The function is real analytic and real valued. If it were nonnegative
near a zero $X_0$, let $P$ be the first nonzero homogeneous term in
its Taylor series there. Such a term exists since $\phi_m$ is not
identically zero. Nonnegativity makes $P$ nonnegative in every real
direction. Equation~\eqref{eq:analytic-orbital-Helmholtz} forces its
lowest-degree part to satisfy $\Delta P=0$. The mean-value property
of harmonic polynomials gives average $P= P(0)=0$ on a sphere, so
nonnegative $P$ must vanish identically, a contradiction. Thus every
zero is approached by negative values. For $c<\mathfrak r_m$, a
neighborhood of every matrix of spectral diameter at most $c$ lies
inside the region with no negative values. No zero can occur there,
and compactness gives the stated strictly positive minimum.

To obtain the explicit lower bound, fix a diagonal spectrum $L$ of
diameter less than $\pi/h_m$, and shift its minimum to zero.
For small $\delta>0$, the finite alphabet
$e^{i\delta h_{m,j}}$ has conjugate pairs of angles at most
$\delta h_m$. The quadratic sector proof in NN-A020 makes it
$PF_r$ whenever $(r+1)\delta h_m<\pi$. Choose the partition
$\mu_i=\lfloor L_i/\delta\rfloor$, in descending order.
Its width satisfies this inequality for sufficiently small $\delta$.
Its Schur value is nonnegative. The exact normalized Weyl--HCIZ
identity \eqref{eq:analytic-Schur-orbital-limit} then gives
$\phi_m(L)\ge0$. Hence no negative value has diameter less than
$\pi/h_m$, proving the bound.

For $m=2$, the HCIZ determinant with Hermite eigenvalues $-1,1$
gives
\[
                \phi_2(X)=\frac{\sin(\lambda_1-\lambda_2)}
                                      {\lambda_1-\lambda_2},
\]
with its continuous value at equality. Its first sign change gives
$\mathfrak r_2=\pi$. For $m=3$, the eigenvalues are
$-\sqrt3,0,\sqrt3$. Factoring the exponential Vandermonde in the
same exact formula gives
\[
 \phi_3(X)=\prod_{1\le i<j\le3}
            \operatorname{sinc}\left(
                          \frac{\sqrt3}{2}(\lambda_i-\lambda_j)\right),
                    \qquad\operatorname{sinc}(u)=\frac{\sin u}{u}.
\]
All factors are positive when the diameter is less than
$2\pi/\sqrt3$. For eigenvalues $0,L/2,L$ with
$2\pi/\sqrt3<L<4\pi/\sqrt3$, the largest-difference factor is
negative and the other two are positive. This proves the second
formula in \eqref{eq:analytic-first-orbital-radii}.
\end{proof}

\begin{lemma}[Uniform finite-cluster positivity]
\label{lem:analytic-uniform-cluster-PF}
Let $m$ be fixed and let a conjugation-stable finite alphabet have
entries $b_\varepsilon e^{\varepsilon a_j(\varepsilon)}$, where
$b_\varepsilon>0$ and
$a_j(\varepsilon)=i\kappa h_{m,j}+O(\varepsilon)$ with fixed
$\kappa>0$. For each $c<\mathfrak r_m/\kappa$, its coefficient
sequence is $PF_{r(\varepsilon)}$ for all sufficiently small
$\varepsilon$, uniformly over integers $r(\varepsilon)\varepsilon\le c$.
The assertion includes every partition width in that range, with no
bound on the degree of the Toeplitz minor.
\end{lemma}

\begin{proof}
The positive common scale does not affect signs. For every partition
$\mu$ of at most $m$ rows and width at most $r$, the exact formula is
\[
 \frac{s_\mu(e^{\varepsilon a})}{s_\mu(1^m)}
       =\frac{\mathcal I(A_\varepsilon,
                       \operatorname{diag}(\varepsilon(\mu+\rho)))}
              {\mathcal I(A_\varepsilon,
                       \operatorname{diag}(\varepsilon\rho))},
            \qquad A_\varepsilon=\operatorname{diag}(a_j),
            \quad \rho=(m-1,\ldots,0).
\]
The second argument in the numerator has entries in
$[0,c+(m-1)\varepsilon]$. Choose $c'$ strictly between $c$ and
$\mathfrak r_m/\kappa$. The preceding lemma gives a positive
minimum for $\phi_m(\kappa X)$ over entries in $[0,c']$.
The Haar integral is uniformly continuous on these compact matrix
sets, while $A_\varepsilon=i\kappa D_m+O(\varepsilon)$.
The numerator is therefore uniformly close to a positive real value,
and the denominator tends to one. The character ratio itself is real,
by conjugation stability. It follows that every such ordinary Schur
value is strictly positive, uniformly in this full partition range.
Partitions with more than $m$ rows have value zero.
Every skew Schur value of width at most $r$ is a nonnegative integer
combination of ordinary Schur values of that width, by the exact
restriction-multiplicity argument in NN-A021. Thus all Toeplitz minors
of order at most $r$ are nonnegative. This proof retains the full
$O(\varepsilon)$ perturbation and the possible coincident limiting
scaled row lengths in the compact set.
\end{proof}

\begin{lemma}[Cluster isolation with earlier multiple endpoints]
\label{lem:analytic-coalescent-prefix-isolation}
The negative-witness conclusion of
Lemma~\ref{lem:analytic-cluster-Schur-witness} remains valid when its
finite prefix is a conjugation-stable multiset of possibly complex
parameters $\zeta_1(\varepsilon),\ldots,\zeta_p(\varepsilon)$ with
positive real limits strictly larger than $b^0$, even when some limits
coincide, provided the parameters are distinct for $\varepsilon>0$ and
\[
                 |\zeta_i-\zeta_j|\ge c\varepsilon
                         \quad(i\ne j)
\]
for a fixed $c>0$. Real prefix entries are positive for small
$\varepsilon$, so $P=\prod_i\zeta_i$ is positive. For any strictly
negative orbital value at spectral diameter $d$, the witness width
can be chosen with
\[
                   \limsup_{\varepsilon\downarrow0}
                               \varepsilon\lambda_1\le d+\delta
\]
for any prescribed $\delta>0$.
\end{lemma}

\begin{proof}
By continuity, first perturb the negative orbital spectrum to have
distinct eigenvalues while preserving strict negativity and increasing
its diameter by less than $\delta/4$. Shift that spectrum to have minimum in $(0,\delta/4)$,
and choose the bounding first-row length constant $C$ less than
$d+\delta$ while keeping a fixed positive gap above its maximum.
Choose the cluster row lengths and the prefix rows of length
$n=\lceil C/\varepsilon\rceil$ exactly as in NN-A028. The gaps
$\mu_m\ge\eta/\varepsilon$ and
$n-\mu_1\ge\eta/\varepsilon$ still hold for some fixed $\eta>0$.
Strict negativity is preserved by the scalar shift because the
limiting Hermite diagonal is traceless.

Here are the changes to the complete determinant estimates.
All $q=p+m$ retained poles now have pairwise separation at least a
constant times $\varepsilon$; the prefix-to-cluster separation remains
bounded below by a positive constant. Thus every partial-fraction
coefficient and the fixed-circle analytic remainder are
$O(\varepsilon^{-(q-1)})$, rather than
$O(\varepsilon^{-(m-1)})$. In a term with remainder columns, the same
sorted-log-modulus argument applies. The entire prefix is assigned
exponent $n$ in the maximal full weight. Put
\[
                  W=P^n b^{|\mu|};
\]
the maximal full weight is $W e^{O(\varepsilon)}$.
The positivity of $P$ makes this the same positive reference scale
as before. Replacing one retained modulus by the smaller tail modulus
still costs a fixed factor to the power $\mu_m$. Consequently the
total tail error is
\[
                     O(W\varepsilon^{-K}q_1^{\mu_m})
\]
for fixed $K$ and $q_1<1$; only the harmless power $K$ has changed.

In the Weyl numerator the leading prefix block is still exactly
$P^{n+m}\Delta(\zeta_1,\ldots,\zeta_p)$. It cancels its entire
prefix Vandermonde, regardless of how small that Vandermonde is.
The leading cross factor is
$\prod_{i,j}(1-\xi_j/\zeta_i)^{-1}$; it is real and tends to a
positive number. Every wrong assignment still exchanges a larger
prefix modulus with a smaller cluster modulus between rows separated
by at least $n-\mu_1$. Its exponential loss is therefore
$q_2^{n-\mu_1}$ with $q_2<1$. Dividing by the full Vandermonde now
costs at most $O(\varepsilon^{-q(q-1)/2})$, by the stated separation.
The wrong-prefix error is again
\[
                      O(W\varepsilon^{-K'}q_2^{n-\mu_1}).
\]
Both errors are smaller than any power of $\varepsilon$ times $W$.
The leading cluster Schur value, divided by its positive dimension
and $b^{|\mu|}$, tends to the fixed strictly negative orbital value.
Its sign is therefore unchanged. The dual-Jacobi--Trudi minor has
order $\lambda_1\le n$, including the empty-prefix case, and fits
in window $0,\ldots,\lambda_1+q-1$. Taking $\varepsilon\to0$
gives the asserted width bound.
\end{proof}

\begin{theorem}[NN-A029: exact orbital formulas for detecting order, window and rank]
\label{thm:analytic-exact-orbital-detection}
Assume A and put $p_0=F_0'(0)>0$. For each positive endpoint zero
$x$ of multiplicity $m_x$, put $\beta_x=1/(p_0x^2)$. Define
\begin{align}
 \mathcal L_F&=\min_{x:m_x\ge2}
                       \frac{x\mathfrak r_{m_x}}{2\sqrt2},
                          \label{eq:analytic-exact-LF}\\
 \mathcal L_U&=\min_{x:m_x\ge2}
                       \frac{x(1-\beta_x)\mathfrak r_{m_x}}{2\sqrt2},
                          \label{eq:analytic-exact-LU}
\end{align}
where a minimum over the empty set means $+\infty$. If there is any
multiple zero, these positive finite minima are attained. Then, in
the extended positive real line,
\begin{equation}\label{eq:analytic-exact-orbital-limits}
 \lim_{a\downarrow0}\sqrt a\,m_*(-a)
 =\lim_{a\downarrow0}\sqrt a\,R_*(-a)=\mathcal L_F,
       \qquad
 \lim_{a\downarrow0}\sqrt a\,d_*(-a)=\mathcal L_U.
\end{equation}
These formulas contain no unstated endpoint multiplicity or simplicity
hypothesis. They express the constants in terms of the actual endpoint
multiplicities and explicit finite-dimensional orbital integrals;
they do not determine those multiplicities.

A double endpoint contributes $\pi x/(2\sqrt2)$ to
\eqref{eq:analytic-exact-LF}; a triple endpoint contributes
$\pi x/\sqrt6$. The corresponding unitary contributions are
multiplied by $1-\beta_x$. Higher-multiplicity contributions are
specified exactly by \eqref{eq:analytic-orbital-radius}, without
asserting a closed formula for every $\mathfrak r_m$.
\end{theorem}

\begin{proof}
The large-radius multiplicity estimate in NN-A022 gives
$h_{m_x}/x\to0$ along positive zeros tending to infinity.
Since $\mathfrak r_m\ge\pi/h_m$, the quantities
$x\mathfrak r_{m_x}$ tend to infinity along multiple zeros escaping
to infinity. Also $1-\beta_x\to1$. Thus each nonempty minimum is
positive, finite and attained among a finite set of endpoint zeros.

Write $\varepsilon=\sqrt a$. The local parameters near an endpoint
$x$ have the expansion, after removing their positive geometric mean,
\[
 z_j(-a)^{-2}=b_x(a)e^{\varepsilon a_{x,j}(\varepsilon)},\qquad
 a_{x,j}(\varepsilon)=-\frac{2i\sqrt2}{x}h_{m_x,j}
                                             +O(\varepsilon).
\]
Changing the sign of the Hermite diagonal only permutes it. At a
simple endpoint the parameter stays positive and real. Uniform
$\ell^1$ tails, compact zero convergence and positivity of the endpoint
products were proved in NN-A027--A028.

Fix any finite $c<\mathcal L_F$. Choose a large circle, avoiding
endpoint zeros, so that the uniform high-radius angular estimate in
NN-A022 gives, for all zeros outside it and all sufficiently small $a$,
\[
                      |\arg(z_j(-a)^{-2})|/\varepsilon
                                      <\frac{\pi}{2(c+1)}.
\]
Every high conjugate quadratic factor is consequently $PF_r$ when
$r=\lfloor c/\varepsilon\rfloor$, for small $\varepsilon$; positive
linear factors are $PF_\infty$. Inside the circle there are finitely
many endpoint clusters. For each multiple cluster,
$c<x\mathfrak r_{m_x}/(2\sqrt2)$, so
Lemma~\ref{lem:analytic-uniform-cluster-PF} gives $PF_r$ uniformly
for its actual perturbed factors. Simple clusters are positive linear
factors. Multiplication and locally uniform convergence now give
$PF_r$ for the full entire coefficient sequence.

In fact every allowed minor of order at most $r$ is strictly positive.
There are infinitely many high factors with constant coefficient one
and positive first coefficient, each $PF_r$. In a Cauchy--Binet path,
move one index at a time one step to its required target without
colliding with another index; each such elementary factor minor is
its positive first coefficient. A finite number of high factors
realizes this path, and all remaining factors may contribute their
positive diagonal minors. Every other term is nonnegative. This is
the explicit strict-path argument used in NN-A020, and it applies
equally to the conjugate quadratic factors here. Thus
$m_*(-a)>r$, and $R_*(-a)\ge m_*(-a)-1$. Since $c<\mathcal L_F$
was arbitrary, this proves the lower limits, also when
$\mathcal L_F=+\infty$.

For the upper limit when a multiple endpoint exists, choose any such
$x$ and a negative orbital value of the scaled Hermite diagonal
$D=(2\sqrt2/x)D_{m_x}$ with spectral diameter less than
$x\mathfrak r_{m_x}/(2\sqrt2)+\delta$.
Its existence follows directly from the infimum defining
$\mathfrak r_{m_x}$. All endpoint parameters below $x$ form a finite
prefix. They need not stay real: earlier multiple endpoints may split
into conjugate clusters. Their moduli, however, remain uniformly
larger than that of the chosen cluster; all limits are positive real.
Within each earlier cluster the Hermite roots give separation at least
a constant times $\varepsilon$, and between distinct endpoint clusters
the separation stays bounded below. Their product is positive.
Hence Lemma~\ref{lem:analytic-coalescent-prefix-isolation} applies.
It gives negative minors of order $n$ and window at most $n+O_x(1)$
with
\[
 \limsup_{a\downarrow0}\sqrt a\,n
                    \le\frac{x\mathfrak r_{m_x}}{2\sqrt2}+2\delta.
\]
Choose a minimizing endpoint and let $\delta\downarrow0$. Together
with the lower limits this proves both limits involving $\mathcal L_F$.

For the unitary statement use the actual translated alphabet
\[
                \widetilde\beta_j(a)=\frac1{p_a z_j(-a)^2-1},
                                \qquad p_a=F_{-a}'(0).
\]
At endpoint $x$ its scaled Hermite diagonal is
\[
                    \frac{2\sqrt2}{x(1-\beta_x)}D_{m_x}.
\]
This follows from the full local expansion in NN-A027; it retains the
$O(\varepsilon)$ perturbation after scaling. The positive endpoint
ordering, the finite coalescent-prefix separation and the uniform
$\ell^1$ tail remain valid. The high-radius angular tail remains
uniformly negligible after division by $\varepsilon$, since
$|\beta_j(a)|\le C/|z_j(-a)|^2$ and
$|\arg(\beta/(1-\beta))|\le|\arg\beta|/(1-|\beta|)$.
The same lower and upper arguments therefore give the constant
\eqref{eq:analytic-exact-LU} for the least negative-minor order of
the translated coefficients. NN-A021 identifies that order exactly
with the least failing unitary rank. Its positive constant
$E_{-a}(-1)$ does not change any minor sign. This proves the remaining
limit. The double and triple values follow from
\eqref{eq:analytic-first-orbital-radii}.
\end{proof}

\paragraph{Geometric and representation meaning.}
For each endpoint local Artin length $m$, the leading heat unfolding
specifies the actual full-flag coadjoint orbit of the Hermite diagonal.
Its Fourier transform has the exact finite-dimensional positivity
radius \eqref{eq:analytic-orbital-radius}. The normalized Weyl--HCIZ
identity controls every integral highest-weight partition in a whole
scaled-width compact set, while its first negative orbital value gives
a negative Schur coefficient after the proved coalescent-prefix and
infinite-tail isolation. These are actual exterior-power tests, with
the representation-ring formalism in SGA~6, Expos\'e~V, Example~2.2(b)
and Lemma~2.2.1. The minima in
\eqref{eq:analytic-exact-LF}--\eqref{eq:analytic-exact-LU} are the exact
asymptotic boundaries of the corresponding analytic positive cones.
The primary HCIZ normalization is Zuber2018 equations (6), (11)--(12),
pp.~311--313, cited in NN-A028. No arithmetic Frobenius object,
integrality of the Hermite symplectic form, or complexity-class
separation is inferred.


\section{A finite Wronskian formula for every Hermite orbital radius}
\label{sec:analytic-orbital-wronskians}

The constants in NN-A029 admit a one-variable description. The
critical-length theorem used below is a classical theorem of extended
Chebyshev spaces; the identification with the Hermite orbital radius
and the resulting detecting constants are the deductions here.

For $m\ge2$ put
\begin{equation}\label{eq:analytic-hermite-ode}
 P_m(z)=i^m\operatorname{He}_m(z/i)
 =\sum_{j=0}^{\lfloor m/2\rfloor}
       \frac{m!}{2^j j!(m-2j)!}\,z^{m-2j}.
\end{equation}
Its distinct roots are $\alpha_j=ih_{m,j}$. Define the real entire
function
\begin{equation}\label{eq:analytic-fundamental-solution}
 S_m(x)=\sum_{j=1}^m\frac{e^{\alpha_jx}}{P_m'(\alpha_j)}
\end{equation}
and the Hankel Wronskians
\begin{equation}\label{eq:analytic-finite-wronskians}
 W_{m,k}(x)=\det\bigl[S_m^{(p+q)}(x)\bigr]_{p,q=0}^k,
 \qquad 0\le k\le\left\lfloor\frac{m-2}{2}\right\rfloor.
\end{equation}
The first positive zero of a function with no positive zero is, by
convention, $+\infty$.

\begin{theorem}[NN-A030: exact Wronskian reduction and computability]
\label{thm:analytic-orbital-wronskians}
The orbital radius of NN-A029 is
\begin{equation}\label{eq:analytic-wronskian-radius}
 \boxed{\displaystyle
 \mathfrak r_m=
 \min_{0\le k\le\lfloor(m-2)/2\rfloor}
        \inf\{x>0:W_{m,k}(x)=0\}.}
\end{equation}
Equivalently, it is the critical interval length for Hermite
interpolation in the real solution space of $P_m(D)f=0$.
In particular,
\begin{equation}\label{eq:analytic-radius-elementary-bounds}
       \frac\pi{h_m}\le\mathfrak r_m
                         \le\frac{(m-1)\pi}{h_m}.
\end{equation}
There is an algorithm, uniform in $m$, which returns rational bounds
$c<\mathfrak r_m<C$ of any prescribed positive rational width.
This is a computability assertion, without a polynomial-time or
uniform practical running-time bound. Consequently every local
constant in NN-A029 is specified by finitely many explicit
one-variable entire functions. Determining the actual endpoint
multiplicities remains a separate problem.
\end{theorem}

\begin{proof}
\emph{The differential equation and its canonical basis.}
The coefficient formula in \eqref{eq:analytic-hermite-ode} follows
directly from the defining formula for the probabilists' Hermite
polynomial. Thus $P_m$ is a real monic polynomial of parity $m$.
Complex conjugation permutes the terms in
\eqref{eq:analytic-fundamental-solution}, so $S_m$ is real on the
real axis. The elementary Lagrange interpolation identities give
\[
 \sum_j\frac{\alpha_j^r}{P_m'(\alpha_j)}=0
       \quad(0\le r<m-1),\qquad
 \sum_j\frac{\alpha_j^{m-1}}{P_m'(\alpha_j)}=1.
\]
Hence $S_m$ is the unique solution of $P_m(D)S_m=0$ with
\[
 S_m^{(r)}(0)=0\quad(0\le r<m-1),\qquad
 S_m^{(m-1)}(0)=1.
\]
Let $E_m$ be the real solution space. It has dimension $m$, and is
closed under translations and reflections. The functions
\[
 u_j=S_m^{(m-1-j)},\qquad 0\le j\le m-1,
\]
form a canonical basis: $u_j^{(r)}(0)=0$ for $r<j$ and
$u_j^{(j)}(0)=1$. Their complex span equals the span of
$e^{\alpha_1x},\ldots,e^{\alpha_mx}$.

\emph{Orbital determinants and arbitrary Hermite nodes.}
Write $Q=m(m-1)/2$ and $C_m=\prod_{q=1}^{m-1}q!$. The normalized
HCIZ identity, with both Vandermondes in the same order, says
\begin{equation}\label{eq:analytic-confluent-hciz}
 \phi_m(\operatorname{diag}(x_1,\ldots,x_m))
 =C_m\frac{\det[e^{\alpha_jx_k}]_{j,k=1}^m}
               {\Delta(\alpha)\Delta(x)},
 \qquad \Delta(\alpha)=i^Q\Delta(h).
\end{equation}
The identity extends continuously to repeated $x_k$. For distinct
nodes $t_1<\cdots<t_s$ with multiplicities $\nu_1+\cdots+\nu_s=m$,
the confluent numerator has columns
\[
 \left(\frac{\alpha_j^r e^{\alpha_jt_\ell}}{r!}\right)_{j=1}^m,
 \qquad 1\le\ell\le s,\quad0\le r<\nu_\ell,
\]
and the remaining denominator is the nonzero constant
$\Delta(\alpha)\prod_{\ell<v}(t_v-t_\ell)^{\nu_\ell\nu_v}$.
Thus the orbital value vanishes exactly when this Hermite
evaluation matrix is singular. Changing from the complex
exponential basis to a real basis of $E_m$ shows that singularity
is equivalent to a nonzero real $f\in E_m$ with zeros of the
prescribed multiplicities. In particular, this includes all
nonreduced node configurations, rather than just distinct nodes.

An $m$-dimensional real space is an extended Chebyshev space on an
interval if every nonzero member has at most $m-1$ zeros there,
counting multiplicities. By NN-A029, $\phi_m$ is strictly positive
for every spectrum of diameter less than $\mathfrak r_m$.
Consequently $E_m$ is an extended Chebyshev space on every
$[0,c]$ with $c<\mathfrak r_m$.

Conversely, if $c>\mathfrak r_m$, choose a negative orbital value
with spectral diameter less than $c$. Subtracting a scalar matrix
does not change it, since $\operatorname{Tr}D_m=0$. Along the
radial path from the zero matrix to this spectrum, the orbital
value starts at $1$ and ends negative. It vanishes at an
intermediate spectrum lying in $[0,c]$. The confluent determinant
criterion therefore makes $E_m$ fail the extended Chebyshev
property there. At the endpoint $c=\mathfrak r_m$, take negative
spectra with diameters decreasing to $\mathfrak r_m$, normalize
their minimum to zero, and extract a convergent subsequence. The
limit has orbital value zero: a negative limit could be radially
contracted to contradict the infimum. The endpoint also fails.
This proves that the critical interval length is exactly
$\mathfrak r_m$.

\emph{The finite Wronskian theorem and its precise indices.}
Carnicer--Mainar--Pe\~na, \emph{Spherical Bessel functions and
critical lengths}, RACSAM \textbf{117} (2023), article~29,
section~3, Proposition~3, p.~6, states that for a
reflection-invariant finite-dimensional space with a canonical
basis $(u_0,\ldots,u_n)$ and finite critical length, that length is
the least positive zero among
\[
                 \det W(u_j,\ldots,u_n),\qquad j>n/2.
\]
All hypotheses hold here; finiteness was proved in NN-A029 and
also follows from the elementary upper bound below. Put $n=m-1$
and $k=m-1-j$. Reversing these Wronskian columns changes only a
constant sign and gives $W_{m,k}$. The condition $j>(m-1)/2$
becomes $0\le k\le\lfloor(m-2)/2\rfloor$. This proves
\eqref{eq:analytic-wronskian-radius}. There are exactly
$\lfloor m/2\rfloor$ functions to test.

The lower bound in \eqref{eq:analytic-radius-elementary-bounds}
is NN-A029. For the upper bound, the nonzero function
$\sin(h_mx)$ belongs to $E_m$ and has the $m$ distinct zeros
\[
                0,\ \pi/h_m,\ \ldots,\ (m-1)\pi/h_m.
\]
The extended Chebyshev property therefore fails on their span.

\emph{Effective two-sided approximation, including tangential
Wronskian zeros.}
All roots of $\operatorname{He}_m$ are distinct real algebraic
numbers; isolating them and evaluating the finite exponential
sums with rigorous errors are effective. At zero,
\begin{equation}\label{eq:analytic-wronskian-leading}
 W_{m,k}(x)=(-1)^{k(k+1)/2}
 \left(\prod_{j=0}^k\frac{j!}{(m-1-j)!}\right)
 x^{(k+1)(m-1-k)}(1+O(x)).
\end{equation}
This follows by inserting the leading monomials
$S_m^{(j)}(x)=x^{m-1-j}/(m-1-j)!+O(x^{m-j})$ and evaluating
their Vandermonde Wronskian. Divide by the displayed power of $x$
and its nonzero signed leading coefficient. The resulting entire
functions have value $1$ at zero, computable Taylor coefficients
and effective uniform approximation bounds on every bounded
interval. Here are explicit bounds. Set $q=k+1$,
$d=q(m-1-k)$, and let $c_{m,k}$ denote the signed nonzero leading
coefficient in \eqref{eq:analytic-wronskian-leading}.
Cauchy--Binet gives the finite exponential sum
\[
 W_{m,k}(x)=\sum_{|J|=q} A_J e^{b_Jx},\qquad
 A_J=\frac{\Delta(\alpha_J)^2}{\prod_{j\in J}P_m'(\alpha_j)},
 \quad b_J=\sum_{j\in J}\alpha_j.
\]
Put $B=\sum_J|A_J|$ and $M=qh_m$. All $b_J$ are purely imaginary
and $|b_J|\le M$. For $T(x)=W_{m,k}(x)/(c_{m,k}x^d)$, its degree-$N$
Taylor polynomial $T_N$ therefore satisfies, on $|x|\le R$,
\[
 |T-T_N|\le\frac{BM^d}{|c_{m,k}|}
      \frac{(MR)^{N+1}e^{MR}}{(N+d+1)!},\qquad
 \sup_{|x|\le R}|T'(x)|
             \le\frac{BM^{d+1}}{|c_{m,k}|}e^{MR}.
\]
Indeed the coefficients of $T$ are
$W_{m,k}^{(n+d)}(0)/(c_{m,k}(n+d)!)$ and their absolute values are at
most $BM^{n+d}/(|c_{m,k}|(n+d)!)$. Summing this majorant gives both
bounds. The algebraic numbers $A_J,b_J$ and the elementary
exponentials allow outward rational bounds for every displayed
quantity.

For any rational $c<\mathfrak r_m$, all these finitely many
normalized Wronskians are strictly positive on $[0,c]$. Their
positive minima can therefore be certified by successive
uniform approximations and a finite mesh with rigorous
derivative bounds. This provides a semidecision procedure for
valid rational lower bounds. A certificate cannot pass at or
beyond the first zero.

Independently enumerate distinct rational spectra and certify
negative orbital values using \eqref{eq:analytic-confluent-hciz}.
Every rational $C>\mathfrak r_m$ admits such a negative spectrum
of diameter less than $C$: a negative spectrum exists by the
definition of the infimum, and negativity persists under a
sufficiently small perturbation to distinct rational entries.
This supplies certified upper bounds. Dovetail the lower and
upper searches until their width is smaller than the requested
error. Both searches terminate for bounds chosen sufficiently
close to the radius. No decision about whether an individual
Wronskian root is multiple, and no unjustified sign-change test
at that root, is used. This proves the algorithmic claim.
\end{proof}

\paragraph{The first examples.}
Here $S_2(x)=\sin x$ and
$S_3(x)=(1-\cos(\sqrt3x))/3$, so the single-Wronskian cases give
$\mathfrak r_2=\pi$ and $\mathfrak r_3=2\pi/\sqrt3$ again.
For $m=4$, put
\[
 a=\sqrt{3-\sqrt6},\qquad b=\sqrt{3+\sqrt6},\qquad
 S_4(x)=\frac{\sin(ax)/a-\sin(bx)/b}{b^2-a^2}.
\]
Then $\mathfrak r_4$ is exactly the smaller of the first positive
zeros of $S_4$ and $S_4S_4''-(S_4')^2$. This is a finite scalar
root problem, without a numerical value being asserted here.

\paragraph{A checked source correction.}
Beccari--Casciola--Mazure, \emph{Critical length: an alternative
approach}, arXiv:1904.09253v1, p.~12, Theorem~4.3 and
Corollary~4.4, equations~(19)--(20), print $W\ne0$ inside the
infima. The original page image was checked: this is a printed
sign error, not an extraction artifact. Those conditions must
read $W=0$, as the preceding sentence's first-zero description
and the correct Proposition~3 above specify. For $P_2(D)=D^2+1$
the printed nonzero condition gives infimum $0$ from $S_2=\sin x$,
whereas the critical length is $\pi$. No claim is made here
about whether later versions corrected that display.

\paragraph{Geometric translation.}
The algebra $\mathbb Q[z]/(P_m)$ is finite etale because $P_m$
has distinct roots; the finite-free discriminant criterion is
SGA~1, Expos\'e~I, section~4, Proposition~10. A companion matrix
$B_m$ of $P_m$ gives the actual algebraic connection $d-B_m\,dx$ on a
rank-$m$ trivial bundle on $\mathbb A^1$. Its analytic horizontal
sections are the solution space above. Evaluation on the finite
nonreduced scheme
\[
              \coprod_{\ell=1}^s
                  \operatorname{Spec}\mathbb C[x]/(x-t_\ell)^{\nu_\ell}
\]
is exactly the Hermite jet map, and its determinant vanishes
precisely when interpolation fails. A zero of $W_{m,k}(x)$ gives
a nonzero solution vanishing to order at least $m-1-k$ at $0$
and at least $k+1$ at $x$: use the span of
$u_{m-1-k},\ldots,u_{m-1}$. Thus the Wronskian formula reduces
this real interval boundary to finitely many two-point jet tests.
It does not turn the exponential connection into an arithmetic
Frobenius object, a pure motive, or a polynomial-time algorithm.


\section{Uniform growth of the Hermite orbital positivity radius}
\label{sec:analytic-hermite-radius-growth}

The one-variable formula in NN-A030 determines each fixed radius.
A uniform matrix-integral theorem supplies the missing growth rate
when the endpoint multiplicity increases. The matrix-integral
nonvanishing theorem is an imported result; the scaling to the
Hermite radii and to the heat-detection constants is the deduction.

\paragraph{The exact imported theorem and normalization.}
Novak, \emph{Topological Expansion of Oscillatory BGW and HCIZ
Integrals at Strong Coupling}, arXiv:2203.10746v1 (2022), equation
(1.2), p.~2, defines
\[
 I_N(z,A,B)=\int_{U(N)}
          \exp\bigl(zN\operatorname{Tr}(AUBU^*)\bigr)\,dU,
\]
with probability Haar measure and ordinary, unnormalized trace.
The closed polydisc $D_N(\varepsilon)$ is defined on p.~3 by
$|z|\le\varepsilon$ and modulus at most one for each eigenvalue
of $A$ and $B$. The first assertion of Theorem~1.2, p.~6, gives
an absolute $\varepsilon_0\in(0,2/27)$ such that $I_N$ is nonzero
throughout this polydisc for every $N\ge1$. It applies in
particular to the Hermitian matrices used here. The conclusion
is uniform in the dimension, not just a statement for each fixed
dimension. We do not assert a numerical value for $\varepsilon_0$.

\begin{theorem}[NN-A038: square-root growth in multiplicity]
\label{thm:analytic-hermite-radius-growth}
For the Hermite orbital radii $\mathfrak r_m$ of NN-A029 and
$h_m=\max_j|h_{m,j}|$, one has, for every integer $m\ge2$,
\begin{equation}\label{eq:analytic-radius-stable-bounds}
 \frac{2\varepsilon_0 m}{h_m}
       \le\mathfrak r_m\le\frac{\pi(m-1)}{h_m}.
\end{equation}
Consequently
\begin{equation}\label{eq:analytic-radius-root-growth}
 \varepsilon_0\sqrt m\le\mathfrak r_m
                         \le\pi\sqrt{m-1},
 \qquad\mathfrak r_m=\Theta(\sqrt m).
\end{equation}
More precisely,
\begin{equation}\label{eq:analytic-radius-liminf-limsup}
 \varepsilon_0\le
 \liminf_{m\to\infty}\frac{\mathfrak r_m}{\sqrt m}
 \le\limsup_{m\to\infty}\frac{\mathfrak r_m}{\sqrt m}
 \le\frac\pi2.
\end{equation}
These are unconditional finite-dimensional statements. Existence
of a limit in the last display, its value, an explicit admissible
$\varepsilon_0$, and a closed formula for each $m$ are not claimed.
\end{theorem}

\begin{proof}
Let $X=X^*$ have spectral diameter $R>0$, and put
\[
 c=\tfrac12\bigl(\lambda_{\max}(X)+\lambda_{\min}(X)\bigr),
 \quad A=\frac{2(X-cI)}R,\quad B=\frac{D_m}{h_m},
 \quad z=\frac{iRh_m}{2m}.
\]
Both $A$ and $B$ are Hermitian with spectra in $[-1,1]$.
Since $\operatorname{Tr}D_m=0$, invariance of Haar measure and
cyclicity of the trace give exactly
\[
 I_m(z,A,B)
   =\int_{U(m)}e^{i\operatorname{Tr}((X-cI)UD_mU^*)}\,dU
   =\phi_m(X).
\]
If $R\le2\varepsilon_0m/h_m$, the imported theorem makes this
number nonzero. The Hermite spectrum is invariant under negation,
so a unitary permutation conjugates $D_m$ to $-D_m$; hence
$\phi_m(X)$ is real. Along $uX$, $0\le u\le1$, the same
nonvanishing statement applies and the value starts at $1$.
It is therefore strictly positive. A scalar $X$, the case $R=0$,
also has value $1$. The definition of $\mathfrak r_m$ proves the
lower bound in \eqref{eq:analytic-radius-stable-bounds}. The upper
bound is the already proved extended-Chebyshev bound in NN-A030.

For completeness, all required Hermite size estimates follow
from a finite symmetric matrix. Let $J_m$ be the real tridiagonal
$m\times m$ matrix with zero diagonal and off-diagonal entries
$\sqrt1,\sqrt2,\ldots,\sqrt{m-1}$. Expanding its characteristic
determinant along the last row gives
\[
 \det(xI_m-J_m)=x\det(xI_{m-1}-J_{m-1})
                  -(m-1)\det(xI_{m-2}-J_{m-2}),
\]
with initial values $1,x$. This is the defining recurrence of
$\operatorname{He}_m$, so the eigenvalues of $J_m$ are the
$h_{m,j}$. Alternating the signs of the coordinate vectors
conjugates $J_m$ to $-J_m$, and therefore its top eigenvalue is
$h_m=\|J_m\|$. Write $J_m=T_m+T_m^*$, where $T_m$ is the
weighted one-step shift. Then
\[
 h_m\le2\|T_m\|=2\sqrt{m-1},\qquad
 h_m^2\ge\frac1m\operatorname{Tr}(J_m^2)=m-1.
\]
Substitution gives
\[
 \mathfrak r_m\ge\frac{\varepsilon_0m}{\sqrt{m-1}}
                      \ge\varepsilon_0\sqrt m,
 \qquad
 \mathfrak r_m\le\pi\sqrt{m-1},
\]
which proves \eqref{eq:analytic-radius-root-growth}.

We also prove the precise scale of $h_m$, rather than importing
a Hermite asymptotic. For $2\le k<m$, restrict to the last $k$
coordinates. Its off-diagonal weights are at least $\sqrt{m-k}$.
The positive unit eigenvector of the unweighted path on $k$
vertices has entries proportional to $\sin(j\pi/(k+1))$ and
eigenvalue $2\cos(\pi/(k+1))$. Its Rayleigh quotient in this
principal submatrix is thus at least
\[
                  2\sqrt{m-k}\cos\frac\pi{k+1}.
\]
The variational principle gives the same lower bound for $h_m$.
Choose any integer sequence $k\to\infty$ with $k/m\to0$,
and use the preceding upper bound. It follows that
$h_m/(2\sqrt m)\to1$. Dividing
\eqref{eq:analytic-radius-stable-bounds} by $\sqrt m$ proves
\eqref{eq:analytic-radius-liminf-limsup}.
\end{proof}

\paragraph{The resulting heat and representation constants.}
Under A and the exact heat normalization of NN-A029, a positive
endpoint zero $x$ of multiplicity $m\ge2$ contributes
\[
 C_F(x,m)=\frac{x\mathfrak r_m}{2\sqrt2},\qquad
 C_U(x,m)=\frac{x(1-\beta_x)\mathfrak r_m}{2\sqrt2},
 \qquad\beta_x=\frac1{p_0x^2},\quad p_0=F_0'(0).
\]
Thus, uniformly in the abstract multiplicity,
\[
 \frac{\varepsilon_0x\sqrt m}{2\sqrt2}
       \le C_F(x,m)\le\frac{\pi x\sqrt{m-1}}{2\sqrt2},
\]
and the same bounds multiplied by $1-\beta_x>0$ hold for $C_U$.
The least Toeplitz order and index window have the minimum of
the $C_F$ as their exact $\sqrt a$-scaled limit; the least failing
unitary rank has the minimum of the $C_U$, as proved in NN-A029.
These estimates describe the dependence of each candidate on
multiplicity. They do not assert that arithmetic zeros of any
prescribed multiplicity exist or that their multiplicities tend
to infinity.

\paragraph{Geometric and Grothendieck translation.}
The Hermite eigenvalues are distinct, so their conjugacy orbit
is the compact full-flag manifold $U(m)/U(1)^m$, canonically the
complex flag variety of $\mathbb C^m$ as a homogeneous space.
Its specified Haar probability measure has Fourier transform
$\phi_m$; the theorem proves a dimension-uniform positive region
for this concrete transform. The Weyl--HCIZ and Schur dictionary
proved in NN-A028--NN-A030 identifies its spectral-diameter
bound with a bound on the width of every allowed highest-weight
partition in the scaled heat alphabet. Schur and exterior powers
are the actual representation-ring operations of SGA~6,
Expos\'e~V, Example~2.2(b) and Lemma~2.2.1. The displayed
$\Theta(\sqrt m)$ bounds therefore quantify how the local
Artin length $m$ changes these positivity tests. This is a
finite-dimensional flag-orbit statement; no integral polarization
for the nonintegral Hermite weights or arithmetic cohomology
realization is asserted.

\paragraph{Source status and remaining gate.}
The input is the stated theorem of the cited 2022 primary
preprint, which explicitly supersedes arXiv:2006.04304v1.
Its definitions and complete theorem statement were checked,
as was the uniform nonvanishing step in Theorem~4.7,
equations~(4.40)--(4.42), p.~34. That step first treats all large
$N$ and then shrinks the disc to handle the finitely many smaller
dimensions. We have not claimed to reproduce or independently
verify all of its combinatorial proof. A sharper remaining
question is whether $\mathfrak r_m/\sqrt m$ converges, and, if
so, its value and a quantitative error. The growth exponent is
resolved by the displayed deduction.


\section{Exact interfaces: arithmetic tests, fluid profiles, and source corrections}

\begin{proposition}[NN-A011: the specified fluid and branch channels preserve the A constraint]
Retain the exact source maps in workbench24 and workbench29i.
Under A, every finite matrix $[W(h_i,h_j)]$ of admissible scalar
Weil tests $h_i\in C_c^\infty(\mathbb R)$ is positive semidefinite,
including every polynomial channel $h_i=L_{P_i}U_i$ constructed in
workbench24. The branch-trace form on its retained two-sheet module
is a different form and can be negative without contradicting this.

In the released-witness arithmetic realization of workbench29i, take
$H=\Xi$ and
\[
 \Psi_{\nu,j}(t,q,\zeta)
 =\Xi(S_j(X_{\nu,t}(q))+\zeta).
\]
For every preterminal time, and almost every terminal label in the
proved terminal construction, all its zeros in the variable $\zeta$
are real under A, with unchanged multiplicity. This holds together
with the source's proved divergence of the specified fixed-label
partial-derivative norm. In this particular map the derivative
singularity cannot yield an off-critical zeta zero.
\end{proposition}
\begin{proof}
For A's real ordinates the exact Weil form is
$W(f,g)=\sum_\rho m_\rho\overline{\widehat f(\gamma_\rho)}
\widehat g(\gamma_\rho)$, absolutely convergent on the specified
test class. Hence
$c^*[W(h_i,h_j)]c=\sum_\rho m_\rho|
\widehat{\sum c_ih_i}(\gamma_\rho)|^2\ge0$.
The polynomial channel preserves that test class and compact support,
as proved by its explicit differential-operator formula in workbench24.
By contrast its branch trace has fibre matrix
$\operatorname{diag}(2,-4\lambda)$ at a real $\lambda>0$; it is
not the componentwise Weil form or its scalar pullback.

The shifts $a=S_j(X_{\nu,t}(q))$ are real at all the stated parameters.
The translation $\zeta\mapsto a+\zeta$ has inverse $z\mapsto z-a$,
and transforms a germ $(z-z_0)^mA(z)$ to
$(\zeta-(z_0-a))^mA(a+\zeta)$. It preserves multiplicity and
imaginary part exactly. A makes the zeros $z_0$ of $\Xi$ real.
The completed arithmetic argument is $s=1/2+i(a+\zeta)$,
so $\Re s=1/2-\Im\zeta=1/2$ at those zeros.
No estimate on $\partial_ta$ enters this identity, so its possible
divergence does not affect the conclusion.
\end{proof}

This closes the particular proposed counterexample route after A; it does
not assert a theorem about all hypothetical fluid-to-arithmetic maps.
The underlying categories are a concrete analytic test-function module,
its Hermitian Gram form, and translated analytic divisors. The branch
trace is an algebra trace in a quadratic algebra. The equality of those
forms was never available. Neither a PDE singularity nor this translation
creates a finite-dimensional pure motive, a negative Weil matrix, or a
complexity-class separation.

The Bost--Connes formulas in source22 also retain their domains. Its
physical Gibbs partition function is $\zeta(\beta)>0$ for real
$\beta>1$. A locates the nontrivial zeros of its meromorphic continuation
on $\Re s=1/2$. The source's strong-limit projection $P_1$ has continued
expectation $1/\zeta(s)$ and therefore poles at those zeros; finite
Euler-projector expectations are entire finite products and the strong
operator limit is not automatically a meromorphic limit through those
poles. A supplies zero location, not a positive Gibbs state at complex
temperature. Any ``Lee--Yang'' terminology must refer to this stated
analytic continuation and coordinate, not silently to a unit-circle
theorem for a finite statistical-mechanical partition polynomial.

\begin{proposition}[NN-A012: two exact source corrections in paper16]
The finite-gerbe path integral printed in
\texttt{thm:finite-gerbe-anomaly} satisfies
$\mathcal A_n(S^4)=n$. Consequently its later assertion
$\mathcal A(q_n)\otimes\mathcal A(-q_n)\simeq\mathbf1$
in \path{prop:anomaly-cancellation} fails for every $n>1$ with
that normalization. Complex conjugation is not the inverse of an
arbitrary non-unit-modulus partition value.

Also, the measure algebra of an arbitrary proper adelic curve is a
Boolean sigma-algebra, but need not be a complete Boolean algebra.
The completeness assertion in \texttt{thm:final-synthesis-v11}(3)
requires an additional measure-space hypothesis such as localizability.
\end{proposition}
\begin{proof}
For $S^4$ and $\Phi_n=\mathbb Z/n\mathbb Z$, one has
$|H^0|=n$, $|H^1|=1$, and $H^2=0$. The sole quadratic action is zero,
so the displayed source formula gives $n$ for both $q_n$ and $-q_n$.
Their product is $n^2$, whereas the tensor unit has value one. On this
example a compensating Euler factor $n^{-\chi(S^4)/2}=n^{-1}$
normalizes each value; a general inverse-theory assertion must specify
the full Euler normalization, not merely replace $q$ by $-q$.

For the second assertion, let $\Omega$ be uncountable and equip it
with the countable-or-cocountable sigma-algebra and counting measure.
Assign the trivial absolute value of any field to every point.
Every principal logarithm is identically zero, so this is a proper
adelic curve under the source's stated definition. The only null set
is empty. The displayed definition at the start of
\path{sec:adelic-support-defect} imposes integrability of principal
logarithms and the product formula, but no sigma-finiteness or
localizability; the counterexample is directed at that exact definition.
Thus the measure algebra
is exactly that sigma-algebra.
Choose an uncountable subset $A$ with uncountable complement. The
family of measurable singleton sets $\{\{a\}:a\in A\}$ has no
least measurable upper bound: any upper bound is cocountable and
contains some $x\notin A$, and deleting $x$ gives a smaller measurable
upper bound. Thus arbitrary suprema need not exist, whereas countable
suprema always do.
\end{proof}

These are finite or explicit set-theoretic counterexamples, not new
consequences of E or A. They are retained as source repairs because
otherwise the proposed geometry would rely on false normalization or
completeness assertions. They supply no computational lower bound.
